\documentclass[lang=cn,11pt]{elegantbook}

\title{概率论与数理统计}
\subtitle{学习笔记}

\author{夏同}
\institute{华南理工大学}
\date{\today}
\version{1.0}
\bioinfo{自定义}{信息}

\extrainfo{居敬持志，守正出奇。—— 张红兵}

\setcounter{tocdepth}{3}

\logo{logo-blue.png}
\cover{1.jpg}

% 本文档命令
\usepackage{array}

\makeatletter
\@ifundefined{setCJKfamilyfont}{}{%
  \IfFontExistsTF{KaiTi}{%
    \setCJKfamilyfont{codexkaiti}[AutoFakeBold=2.5]{KaiTi}%
    \renewcommand{\kaishu}{\CJKfamily{codexkaiti}}%
  }{}%
  \IfFontExistsTF{FangSong}{%
    \setCJKfamilyfont{codexfangsong}[AutoFakeBold=2.5]{FangSong}%
    \renewcommand{\fangsong}{\CJKfamily{codexfangsong}}%
  }{}%
}
\makeatother

\newcommand{\ccr}[1]{\makecell{{\color{#1}\rule{1cm}{1cm}}}}
\newcommand{\answerpage}[2]{%
  \begin{center}
    \textbf{#1}\par\vspace{0.6em}
    \includegraphics[width=\textwidth,height=0.82\textheight,keepaspectratio]{#2}
  \end{center}
  \clearpage
}

% 修改标题页的橙色带
% \definecolor{customcolor}{RGB}{32,178,170}
% \colorlet{coverlinecolor}{customcolor}

\begin{document}

\maketitle
\frontmatter

\tableofcontents
\makeenvindex

\mainmatter

\chapter{随机事件与概率}
{\scriptsize
\[
\begin{cases}
\text{基本概念} & 
\begin{cases}
\text{随机试验} \ E: \text{可重复性、结果不唯一、所有结果已知、不确定性} \\
\text{样本空间} \ \Omega: \text{所有可能结果构成的集合} \\
\text{随机事件} \ A: \text{样本空间的子集} 
\begin{cases}
\text{关系：包含、相等、互斥、对立} \\
\text{表示：和事件、积事件、差事件、对立事件} \\
\text{运算：交换律、结合律、分配律、德摩根律}
\end{cases}
\end{cases} \\
\text{概率概念} & 
\begin{cases}
\text{统计概念：频率} \\
\text{公理化概念：非负性、规范性和可列可加性} \\
\text{古典模型(等可能概型)：有限元素、可能性相同} \\
\text{计算公式}
\begin{cases}
\text{加法公式}
\begin{cases}
P(A \cup B) = P(A) + P(B) - P(AB) \\
P(A \cup B) = P(A) + P(B) \ (\text{若事件} \ A \ \text{与} \ B \ \text{互斥})
\end{cases} \\
\text{减法公式}
\begin{cases}
P(A - B) = P(A) - P(AB) = P(A\overline{B}) \\
P(A - B) = P(A) - P(B) \ (\text{若事件} \ A \supset B)
\end{cases} \\
\text{求逆公式：} \ P(\overline{A}) = 1 - P(A)
\end{cases}
\end{cases} \\
\text{条件概率概念} & 
\begin{cases}
 \text{定义：当} \ P(A) > 0 \ \text{时}, \ P(B|A) = \dfrac{P(AB)}{P(A)} \\
 \text{且满足概率性质与上述公式} \\
 \quad (\text{非负性、规范性和可列可加性}) \\
 \text{计算：定义法、缩小样本空间法} \\
 \text{应用}
 \begin{cases}
 \text{乘法公式：当} \ P(A) > 0 \ \text{时}, \ P(AB) = P(B|A)P(A) \\
 \text{全概率公式(由因求果)：} \\
 \qquad P(A) = \sum_{i=1}^{n} P(A|B_i)P(B_i) \\
 \text{贝叶斯公式(执果寻因)：} \\
 \qquad P(B_i|A) = \dfrac{P(A|B_i)P(B_i)}{\sum_{j=1}^{n} P(A|B_j)P(B_j)} \\
 \quad (\text{其中} \ B_1,B_2,\dots,B_n \ \text{是一组完备事件组}, \ P(B_i) > 0)
 \end{cases}
\end{cases} \\
\text{独立性} & 
\begin{cases}
 \text{定义}
 \begin{cases}
 \text{两个事件:} \ P(AB) = P(A)P(B) \\
 \text{等价地，当} \ 0 < P(A) < 1 \ \text{时}, \ P(B|A) = P(B|\overline{A}) \\
 \text{多个事件：相互独立} \Rightarrow \text{两两独立}, \\
 \text{但两两独立} \nRightarrow \text{相互独立}
 \end{cases} \\
\text{性质}
\begin{cases}
\text{相互独立的事件中部分事件间相互独立} \\
\text{相互独立的事件中部分事件的逆事件与其余事件仍相互独立}
\end{cases} \\
 \text{应用}
 \begin{cases}
 \text{计算多个事件的和事件概率转化为逆事件概率乘积} \\
 \text{伯努利试验：只有两个结果} \ A \ \text{与} \ \overline{A}, \\
 n \ \text{次试验相互独立}
 \end{cases}
\end{cases}
\end{cases}
\]
}
\newpage

\section{基本概念}

\begin{definition}[随机试验] \label{def:random-experiment}
如果一个试验满足以下三个条件：
\begin{enumerate}
    \item 试验可以在相同的条件下重复进行；
    \item 试验所有可能结果明确可知，且不止一个；
    \item 每一次试验会出现哪一个结果，事先并不能确定，
\end{enumerate}
则称此试验为\textbf{随机试验}。随机试验简称试验，用字母 \(E\) 或 \(E_1, E_2, \dots\) 表示。
\end{definition}

\begin{definition}[样本空间] \label{def:sample-space}
随机试验的每一个可能结果称为\textbf{样本点}，通常记为 \(\omega\)。全体样本点构成的集合称为\textbf{样本空间}（或基本事件空间），记为 \(\Omega\)（或 \(S\)），可写作
\(
\Omega=\{\omega:\omega \text{ 是该随机试验的一个可能结果}\}.
\)
\end{definition}

\begin{definition}[基本事件，随机事件] \label{def:random-event}
由一个样本点构成的事件称为\textbf{基本事件}。在一次试验中可能出现，也可能不出现的结果称为\textbf{随机事件}，简称事件，用大写字母 \(A, B, C\) 等表示。随机事件是样本空间 \(\Omega\) 的子集。在每次试验中，当且仅当这一子集中的一个样本点出现时，称这一事件发生。
\end{definition}

\section{事件的关系与运算}

事件是一个集合，因此事件之间的关系和运算与集合论中集合之间的关系和运算完全类似，可以利用集合论的知识来理解事件的关系与运算。

\subsection{事件的关系}

\begin{enumerate}
    \item \textbf{包含}——如果事件 \(A\) 发生必导致事件 \(B\) 发生，则称事件 \(B\) 包含事件 \(A\)（或 \(A\) 被 \(B\) 包含），记为 \(A \subset B\)。也称 \(A\) 是 \(B\) 的子事件。
    \item \textbf{相等}——如果 \(A \subset B\) 且 \(B \subset A\)，则称事件 \(A\) 与 \(B\) 相等，记为 \(A = B\)。\(A\) 与 \(B\) 相等，事实上也就是说，\(A\) 与 \(B\) 由一些完全相同的试验结果构成，它不过是同一事件表面上看起来不同的两个说法而已。
    \item \textbf{相容}——若 \(AB \neq \varnothing\)，则称事件 \(A\) 和 \(B\) 相容。
    \item \textbf{互斥（互不相容）}——若 \(AB = \varnothing\)，则称事件 \(A\) 与 \(B\) 互不相容，也叫互斥。如果一些事件中任意两个事件都互斥，则称这些事件是两两互斥的。任意一对基本事件都是互斥的。
    \item \textbf{对立}——事件 \(A, B\) 不能同时发生，且至少一个发生，称事件 \(A, B\) 互为对立事件，或互为逆事件，即 \(A \cap B = \varnothing\)，\(A \cup B = \Omega\)。（互斥与对立的区别：互斥只要求 \(AB = \varnothing\)，而对立还要求 \(A \cup B = \Omega\)。因此对立一定是互斥的，但互斥不一定是对立的。）
\end{enumerate}

\begin{definition}[完备事件组] \label{def:complete-event-group}
如果 \(\bigcup_{i=1}^{n} A_i = \Omega\)（或 \(\bigcup_{i=1}^{\infty} A_i = \Omega\)），且 \(A_i A_j = \varnothing\)（对一切 \(i \neq j\)；\(i, j = 1, 2, \dots, n(\dots)\)），称有限个（或可列个）事件 \(A_1, A_2, \dots, A_n(\dots)\) 构成一个\textbf{完备事件组}。
\end{definition}

\begin{enumerate}
    \item \textbf{积事件（交）}——称"事件 \(A\) 与 \(B\) 同时发生"的事件为事件 \(A\) 与 \(B\) 的\textbf{积事件}（或交事件），记为 \(A \cap B\) 或 \(AB\)。称"有限个（或可列个）事件 \(A_1, A_2, \dots, A_n(\dots)\) 同时发生"的事件为这些事件的积事件，记为 \(\bigcap_{i=1}^{n} A_i\)（或 \(\bigcap_{i=1}^{\infty} A_i\)）。
    \item \textbf{和事件（并）}——称"事件 \(A\) 与 \(B\) 至少有一个发生"的事件为事件 \(A\) 与 \(B\) 的\textbf{和事件}（或并事件），记为 \(A \cup B\)。称"有限个（或可列个）事件 \(A_1, A_2, \dots, A_n(\dots)\) 至少有一个发生"的事件为这些事件的和事件，记为 \(\bigcup_{i=1}^{n} A_i\)（或 \(\bigcup_{i=1}^{\infty} A_i\)）。
    \item \textbf{差事件}——称"事件 \(A\) 发生而事件 \(B\) 不发生"的事件为事件 \(A\) 与 \(B\) 的\textbf{差事件}，记为 \(A - B\)。
    \item \textbf{逆事件（对立事件）}——称"事件 \(A\) 不发生"的事件为事件 \(A\) 的\textbf{逆事件}或\textbf{对立事件}，记为 \(\overline{A}\)。
\end{enumerate}

由定义易知
\[
A - B = A - AB = A\overline{B}, \quad B = \overline{A} \Leftrightarrow AB = \varnothing \text{ 且 } A \cup B = \Omega.
\]

\subsection{事件的运算律}

\begin{enumerate}
    \item \textbf{吸收律}——若 \(A \subset B\)，则 \(A \cup B = B\)，\(A \cap B = A\)。
    \item \textbf{交换律}——\(A \cup B = B \cup A\)，\(A \cap B = B \cap A\)。
    \item \textbf{结合律}——\((A \cup B) \cup C = A \cup (B \cup C)\)，\((A \cap B) \cap C = A \cap (B \cap C)\)。
    \item \textbf{分配律}——\(A \cap (B \cup C) = (A \cap B) \cup (A \cap C)\)，\(A \cup (B \cap C) = (A \cup B) \cap (A \cup C)\)。特别地，
    \(
    A \cap (B - C) = (A \cap B) - (A \cap C).
    \)
    \item \textbf{对偶律（德·摩根律）}——\(\overline{A \cup B} = \overline{A} \cap \overline{B}\)，\(\overline{A \cap B} = \overline{A} \cup \overline{B}\)。一般地，对于 \(n \ge 1\)，
    \[
    \overline{\bigcup_{i=1}^{n} A_i} = \bigcap_{i=1}^{n} \overline{A_i}; \quad \overline{\bigcap_{i=1}^{n} A_i} = \bigcup_{i=1}^{n} \overline{A_i}.
    \]
\end{enumerate}

\begin{exercise}\label{exer:1-30}
设随机变量 \(X,Y\)，事件 \(A=\{X\ge c\}\)，事件 \(B=\{Y\ge c\}\)。用 \(A,B\) 表示下列事件：\par
\(D=\{\min(X,Y)<c\}\)；\(E=\{\max(X,Y)\ge c\}\)；\(F=\{\min(X,Y)\ge c\}\)；\(G=\{\max(X,Y)<c\}\)。
\end{exercise}

\begin{solution}
把 \(A=\{X\ge c\}\), \(B=\{Y\ge c\}\) 代回去即可。
\begin{itemize}
    \item \(D=\{\min(X,Y)<c\}\)：至少有一个变量小于 \(c\)，所以 \(D=\overline{A}\cup\overline{B}\)；
    \item \(E=\{\max(X,Y)\ge c\}\)：至少有一个变量不小于 \(c\)，所以 \(E=A\cup B\)；
    \item \(F=\{\min(X,Y)\ge c\}\)：两个都要不小于 \(c\)，所以 \(F=A\cap B\)；
    \item \(G=\{\max(X,Y)<c\}\)：两个都要小于 \(c\)，所以 \(G=\overline{A}\cap\overline{B}=\overline{A\cup B}\)。
\end{itemize}
\end{solution}

\begin{exercise}\label{exer:1-32}
已知事件 \(A,B\) 满足 \(P(AB)=P(\overline{A}\,\overline{B})\)，记 \(P(A)=p\)，则 \(P(B)=\underline{\hspace{2cm}}\)。
\end{exercise}

\begin{solution}
\[
P(\overline{A}\overline{B})=P(1-\overline{A\cup B})=1-P(A)-P(B)+P(AB)\implies P(A)+P(B)=1
\]
\end{solution}

\begin{exercise}\label{exer:1-34}
设 \(A,B\) 为两个随机事件，且 \(P(A)=0.4\)，\(P(A\cup B)=0.7\)。若 \(A,B\) 互不相容，则 \(P(B)=\underline{\hspace{2cm}}\)；若 \(A,B\) 相互独立，则 \(P(B)=\underline{\hspace{2cm}}\)；若 \(A\) 发生则 \(B\) 必发生，则 \(P(B)=\underline{\hspace{2cm}}\)。
\end{exercise}

\begin{solution}
应填 \(0.3\)；\(0.5\)；\(0.7\)。三种情形都先用同一个加法公式
\(
P(A\cup B)=P(A)+P(B)-P(AB).
\)
\begin{itemize}
    \item 互不相容时 \(P(AB)=0\)，所以 \(P(B)=0.7-0.4=0.3\)；
    \item 相互独立时 \(P(AB)=P(A)P(B)\)，代入得 \(0.7=0.4+P(B)-0.4P(B)\)，所以 \(P(B)=0.5\)；
    \item “\(A\) 发生则 \(B\) 必发生”就是 \(A\subset B\)，于是 \(A\cup B=B\)，故 \(P(B)=0.7\)。
\end{itemize}
\end{solution}



\section{概率的定义}

\begin{definition}[概率的公理化定义] \label{def:probability-axiom}
设随机试验的样本空间为 \(\Omega\)，如果对每一个事件 \(A\) 都有一个确定的实数 \(P(A)\)，且事件函数 \(P(\cdot)\) 满足：
\begin{enumerate}
    \item \textbf{非负性}：\(P(A) \ge 0\)；
    \item \textbf{规范性}：\(P(\Omega) = 1\)；
    \item \textbf{可列可加性}：对任意可列个两两互不相容事件 \(A_1, A_2, \dots, A_n, \dots\)（即 \(A_i A_j = \varnothing\)，\(i \neq j\)；\(i, j = 1, 2, \dots\)），有
    \(
    P\!\left(\bigcup_{i=1}^{\infty} A_i\right) = \sum_{i=1}^{\infty} P(A_i),
    \)
\end{enumerate}
则称 \(P(\cdot)\) 为\textbf{概率}，\(P(A)\) 为事件 \(A\) 的\textbf{概率}。
\end{definition}

\begin{property}\label{prop:probability-basic}
概率的基本性质：
\begin{enumerate}
    \item \(P(\varnothing) = 0\)；
    \item \textbf{有限可加性}：若 \(A_1, A_2, \dots, A_n\) 两两互不相容，则
    \(
    P\!\left(\bigcup_{i=1}^{n} A_i\right) = \sum_{i=1}^{n} P(A_i);
    \)
    \item \textbf{减法公式}：若 \(B \subset A\)，则 \(P(A - B) = P(A) - P(B)\)。更一般地，
    \(
    P(A - B) = P(A) - P(AB);
    \)
    \item \textbf{单调性}：若 \(B \subset A\)，则 \(P(B) \le P(A)\)；
    \item \textbf{有界性}：对任意事件 \(A\)，\(0 \le P(A) \le 1\)；
    \item \textbf{求逆公式}：\(P(\overline{A}) = 1 - P(A)\)；
    \item \textbf{加法公式}：\(P(A \cup B) = P(A) + P(B) - P(AB)\)。当 \(A\)、\(B\) 互斥时，\(P(A \cup B) = P(A) + P(B)\)。
    \item 设 \(A_1, A_2, A_3\) 为任意三个事件，则有
    \[
    P(A_1 \cup A_2 \cup A_3) = P(A_1) + P(A_2) + P(A_3) - P(A_1A_2) - P(A_1A_3) - P(A_2A_3) + P(A_1A_2A_3).
    \]
    \item 若 \(A_1, A_2, \dots, A_n\) 是两两互不相容的事件，则有
    \[
    P(A_1 \cup A_2 \cup \dots \cup A_n) = P(A_1) + P(A_2) + \dots + P(A_n).
    \]
\end{enumerate}
\end{property}



\begin{note}
概率为 0 并不意味着为空集。即 \(P(A) = 0 \nRightarrow A = \varnothing\)，而 \(A = \varnothing \Rightarrow P(A) = 0\)。同样的，$P(A) = 0$不能断言$A = \varnothing$，而$P(A) = 1$不能断言$A = \Omega$.
\end{note}


\begin{exercise}\label{exer:1-35}
考虑一元二次方程 \(x^2+Bx+C=0\)，其中 \(B,C\) 分别是将一枚骰子接连抛两次先后出现的点数。则该方程有实根的概率为（\quad）。~~
A. \(\dfrac{19}{36}\)\quad
B. \(\dfrac{17}{36}\)\quad
C. \(\dfrac{15}{36}\)\quad
D. \(\dfrac{13}{36}\)
\end{exercise}

\begin{solution}
应选 A。方程有实根的条件为 \(B^2-4C\ge 0\)。总样本数 \(6^2=36\)。列表统计满足条件的样本点数：
\[
\begin{array}{c|cccccc}
B\setminus C & 1 & 2 & 3 & 4 & 5 & 6 \\
\midrule
1 & - & - & - & - & - & - \\
2 & 0 & - & - & - & - & - \\
3 & + & + & - & - & - & - \\
4 & + & + & + & + & - & - \\
5 & + & + & + & + & + & + \\
6 & + & + & + & + & + & + \\
\end{array}
\]
``+''和``0''对应的样本点共 \(1+2+4+6+6=19\) 个，故 \(P=\dfrac{19}{36}\)。
\end{solution}

\begin{exercise}\label{exer:1-45}
已知 \(P(B)=0.5\)，\(P(A\cup B)=0.7\)，则 \(P(A\overline{B})=\underline{\hspace{2cm}}\)，\(P(\overline{A}\overline{B})=\underline{\hspace{2cm}}\)。
\end{exercise}

\begin{solution}
由加法公式：
\(
P(A\overline{B}) = P(A) - P(AB) = P(A\cup B) - P(B) = 0.7 - 0.5 = 0.2,
\)
\[
P(\overline{A}\overline{B}) = P(\overline{A\cup B}) = 1 - P(A\cup B) = 1 - 0.7 = 0.3.
\]
\end{solution}

\begin{exercise}[提高]\label{exer:1-49}
设事件 \(A,B\) 相互独立，事件 \(B,C\) 互不相容，事件 \(A\) 与 \(C\) 不能同时发生，且 \(P(A)=P(B)=0.5\)，\(P(C)=0.2\)。则事件 \(A,B\) 和 \(C\) 中仅 \(C\) 发生或仅 \(C\) 不发生的概率为 \underline{\hspace{2cm}}。
\end{exercise}

\begin{solution}
应填 \(0.45\)。“仅 \(C\) 发生”即 \(C\overline{A}\,\overline{B}\)。由于 \(AC=\varnothing\) 且 \(BC=\varnothing\)，故 \(C\subset\overline{A}\) 且 \(C\subset\overline{B}\)，即 \(C\subset\overline{A}\,\overline{B}\)，从而 \(C\overline{A}\,\overline{B}=C\)，\(P(C)=0.2\)。

“仅 \(C\) 不发生”即 \(\overline{C}AB\)。由独立性 \(P(AB)=P(A)P(B)=0.25\)，且 \(AB\subset\overline{C}\)（因为 \(AC=\varnothing\)），故 \(\overline{C}AB=AB\)，\(P(\overline{C}AB)=0.25\)。

又 \(C\) 与 \(AB\) 互斥（\(C\subset\overline{A}\)，\(AB\subset A\)），故
\[
P(C\overline{A}\,\overline{B}\cup\overline{C}AB) = P(C)+P(AB) = 0.2+0.25 = 0.45.
\]
\end{solution}


\section{古典概型}

\begin{definition}[古典概型] \label{def:classical-prob}
如果随机试验的样本空间满足：
\begin{enumerate}
    \item 只有有限个样本点（基本事件）；
    \item 每个样本点（基本事件）发生的可能性都一样，
\end{enumerate}
则称随机试验的概率模型为\textbf{古典概型}（等可能概型）。如果古典概型的基本事件总数为 \(n\)，事件 \(A\) 包含 \(k\) 个基本事件，则 \(A\) 的概率为
\[
P(A) = \frac{k}{n} = \frac{\text{事件} \ A \ \text{所含基本事件的个数}}{\text{基本事件总数}}.
\]
由上式计算得出的概率称为 \(A\) 的\textbf{古典概率}。
\end{definition}



\begin{note}
计算古典概率的关键是基本事件和样本空间的选定以及基本事件数的计算：
\begin{enumerate}
    \item \textbf{列举法}（直接数数法）：基本事件数不多时常用。
    \item \textbf{集合对应法}：
    \begin{enumerate}
        \item \textbf{加法原理}——完成一件事有 \(n\) 类办法，第一类办法中有 \(m_1\) 种方法，第二类办法中有 \(m_2\) 种方法，\(\dots\)，第 \(n\) 类办法中有 \(m_n\) 种方法，则完成此事共有 \(\sum_{i=1}^{n} m_i\) 种方法。
        \item \textbf{乘法原理}——完成一件事有 \(n\) 个步骤，第一步有 \(m_1\) 种方法，第二步有 \(m_2\) 种方法，\(\dots\)，第 \(n\) 步有 \(m_n\) 种方法，则完成此事共有 \(\prod_{i=1}^{n} m_i\) 种方法。
        \item \textbf{排列}——从 \(n\) 个不同的元素中取出 \(m(\le n)\) 个元素，并按照一定顺序排成一列，叫作排列。所有排列的个数叫作排列数，记作
        \(
        \mathrm{P}_n^m = n(n-1)(n-2)\cdots(n-m+1) = \frac{n!}{(n-m)!}.
        \)
        当 \(m=n\) 时，\(\mathrm{P}_n^n = \frac{n!}{0!} = n!\)，叫作全排列。
        \item \textbf{组合}——从 \(n\) 个不同的元素中取出 \(m(\le n)\) 个元素，并成一组，叫作组合。所有组合的个数叫作组合数，记作
        \[
        \mathrm{C}_n^m = \frac{\mathrm{P}_n^m}{m!} = \frac{n!}{m!(n-m)!}, \quad \mathrm{C}_n^m = \mathrm{C}_n^{n-m}, \quad 0!=1, \quad \mathrm{C}_n^0=1.
        \]
    \end{enumerate}
    \item \textbf{逆数法}：先求 \(\overline{A}\) 中的基本事件数 \(n_{\overline{A}}\)，用基本事件总数 \(n\) 减去 \(n_{\overline{A}}\) 便得 \(A\) 中的基本事件数，这种方法常用于计算含有"至少"字样的事件的概率。
\end{enumerate}
\end{note}


\begin{exercise}\label{exer:1-4}\label{exer:src-2017-a-ii-02}
将 \(n\) 个球随机放入 \(N(n\le N)\) 个盒子中，每个盒子可以放任意多个球。求下列事件的概率：\par
(1) \(A=\{\text{某指定} \ n \ \text{个盒子各有一球}\}\)；\par
(2) \(B=\{\text{恰有} \ n \ \text{个盒子各有一球}\}\)；\par
(3) \(C=\{\text{指定} \ k(k\le n) \ \text{个盒子各有一球}\}\)。
\end{exercise}

\begin{solution}
设 \(n\) 个球、\(N\) 个盒子是可分辨的（例如编号），由于每个盒子可以放任意多个球，因此每个球都有 \(N\) 种不同的放置方法。将 \(n\) 个球随机放入 \(N\) 个盒子的一种放法作为基本事件，则基本事件总数为 \(N^n\)。\\
(1) 事件 \(A\) 所含的基本事件是 \(n\) 个不同球的一种排列，故
\(
n_A = n!, \quad P(A) = \frac{n!}{N^n}.
\)
(2) 事件 \(B\) 中的基本事件可以设想为先从 \(N\) 个盒子中选出 \(n\) 个（共有 \(\mathrm{C}_N^n\) 种不同选法），而后把 \(n\) 个球随机放入这 \(n\) 个盒子中，每盒一球（共有 \(n!\) 种不同放法），因此
\[
n_B = \mathrm{C}_N^n \cdot n!, \quad P(B) = \frac{\mathrm{C}_N^n \cdot n!}{N^n}.
\]
(3) 事件 \(C\) 的基本事件可以设想为先从 \(n\) 个球中选出 \(k\) 个球（有 \(\mathrm{C}_n^k\) 种不同选法），再将这 \(k\) 个球随机放入指定的 \(k\) 个盒子中，每盒一球（有 \(k!\) 种不同放法），最后将余下的 \(n-k\) 个球随机放入其余的 \(N-k\) 个盒子中，每个球都有 \(N-k\) 种放置方法，因此共有 \((N-k)^{n-k}\) 种不同放法，故
\[
n_C = \mathrm{C}_n^k \cdot k! \cdot (N-k)^{n-k}, \quad P(C) = \frac{\mathrm{C}_n^k k! (N-k)^{n-k}}{N^n} = \frac{n! (N-k)^{n-k}}{(n-k)! N^n}.
\]
\end{solution}


\begin{exercise}\label{exer:1-6}
袋中有 100 个球，40 个白球，60 个黑球。\par
(1) 先后有放回取 20 个球，求取出 15 个白球、5 个黑球的概率；\par
(2) 先后无放回取 20 个球，求取出 15 个白球、5 个黑球的概率；\par
(3) 先后有放回取 20 个球，求第 20 次取到白球的概率；\par
(4) 先后无放回取 20 个球，求第 20 次取到白球的概率。
\end{exercise}
\begin{solution}
\begin{enumerate}
    \item \(p_1 = \frac{\mathrm{C}_{20}^{15} \cdot 40^{15} \cdot 60^{5}}{100^{20}}.\)
    \item 按照"任取 20 个球"来计算，即
    \(p_2 = \frac{\mathrm{C}_{40}^{15} \mathrm{C}_{60}^{5}}{\mathrm{C}_{100}^{20}}.\)
    \item 有放回取球，每次抽取的样本空间没有变化，故每次取到白球的概率始终为 \(\dfrac{40}{100} = \dfrac{2}{5}\)
    \item 看作随机占位问题：设有 100 个盒子，每个盒子中放 1 个球，则要求第 20 个盒子中放入白球即可。故
    \(
    p_4 = \frac{\mathrm{C}_{40}^1 \cdot 99!}{100!} = \frac{40}{100} = \frac{2}{5}.
    \)
\end{enumerate}
\end{solution}

\begin{exercise}\label{exer:1-26}
从一副扑克牌（52 张）中任取 3 张（不重复），则取出 3 张牌中至少有 2 张花色相同的概率为 \underline{\hspace{2cm}}。
\end{exercise}

\begin{solution}
求“至少有 2 张花色相同”，直接分类会比较繁琐，更适合用对立事件。对立事件是\(\text{3 张牌花色各不相同}\)。总取法数为 \(\mathrm C_{52}^{3}\)。下面计算“花色各不相同”的取法数。先从 4 种花色中选出 3 种：\(\mathrm C_4^3\)。然后在每一种被选中的花色中各取 1 张牌，每种花色都有 13 张可选，所以共有\(13^3\)种取法。因此
\[
P(\text{花色各不相同})
=\frac{\mathrm C_4^3\cdot 13^3}{\mathrm C_{52}^{3}}
=\frac{4\cdot 2197}{22100}
=\frac{8788}{22100},~
P(\text{至少 2 张同花色})
=1-\frac{8788}{22100}
=\frac{13312}{22100}
=\frac{256}{425}.
\]
\end{solution}


\begin{exercise}[提高]\label{exer:1-51}
从 1 到 9 的 9 个整数中有放回地随机取 3 次，每次取一个数，则取出的 3 个数之积能被 10 整除的概率为 \underline{\hspace{2cm}}。
\end{exercise}

\begin{solution}
应填 \(\dfrac{52}{243}\)。积能被 10 整除当且仅当积同时含有因子 2 和因子 5。从 \(\{1,\dots,9\}\) 中，5 的倍数仅有 \(\{5\}\)，2 的倍数为 \(\{2,4,6,8\}\)。设 \(A\) 为“至少取出一个 5 的倍数”，\(B\) 为“至少取出一个 2 的倍数”。则所求就是 \(P(AB)\)。先算三个补事件：\[
P(\overline A)=\left(\frac89\right)^3,\qquad P(\overline B)=\left(\frac59\right)^3,\qquad P(\overline A\,\overline B)=\left(\frac49\right)^3.
\]
由容斥原理：
\begin{align*}
P(AB) &= 1 - P(\overline{A}) - P(\overline{B}) + P(\overline{A}\,\overline{B}) = 1 - \frac{8^3}{9^3} - \frac{5^3}{9^3} + \frac{4^3}{9^3} \\
&= 1 - \frac{512}{729} - \frac{125}{729} + \frac{64}{729} = \frac{156}{729} = \frac{52}{243}.
\end{align*}
\end{solution}


\begin{exercise}\label{exer:1-9}
设 \(A, B\) 为两个随机事件，则（\quad）。\\
\begin{tabular}{ll}
A. \(P(AB)+P(\overline{A}\,\overline{B})\le 1\) & B. \(P(AB)+P(\overline{A}\,\overline{B})\ge 1\) \\
C. \(P(A-B)\le P(A)-P(B)\) & D. \(P(A\cup B)+P(\overline{A}\cup\overline{B})\le 1\)
\end{tabular}
\end{exercise}

\begin{solution}
应选 A。由于 \(AB \subset A \cup B\)，即 \(P(AB) \le P(A \cup B)\)，从而
\[
P(AB) \le P(A \cup B) = 1 - P(\overline{A \cup B}) = 1 - P(\overline{A}\,\overline{B}),
\]
即 \(P(AB) + P(\overline{A}\,\overline{B}) \le 1\)，故选 A。

同时有
\[
P(A \cup B) + P(\overline{A} \cup \overline{B}) - 1 = P(A \cup B) + P(\overline{A\,B}) - 1 \ge 0,
\]
故选项 B、D 不正确。又 \(AB \subset B\)，有 \(P(AB) \le P(B)\)，从而
\(
P(A-B) = P(A) - P(AB) \ge P(A) - P(B),
\)
知选项 C 不正确。
\end{solution}

\subsection{几何概型}

\begin{definition}[几何概型] \label{def:geometric-prob}
如果随机试验的样本空间满足：
\begin{enumerate}
    \item 样本空间（基本事件空间）\(\Omega\) 是一个可度量的有界区域；
    \item 每个样本点（基本事件）发生的可能性都一样，即样本点落入 \(\Omega\) 的某一可度量的子区域 \(S\) 的可能性大小与 \(S\) 的几何度量成正比，而与 \(S\) 的位置及形状无关，
\end{enumerate}
则称随机试验的概率模型为\textbf{几何概型}。
\end{definition}

\begin{exercise}\label{exer:1-7}
在区间 \((0,1)\) 中随机地取两个数，则这两个数之差的绝对值小于 \(0.5\) 的概率为 \underline{\hspace{2cm}}。
\end{exercise}

\begin{solution}
把在 \((0,1)\) 中随机取两个数记为 \(x,y\)，则样本点 \((x,y)\) 在单位正方形
\(
D=\{(x,y)\mid 0<x<1,\ 0<y<1\}
\)
内均匀分布，所以问题转化为一个几何概型。条件 \(|x-y|<0.5\) 表示点 \((x,y)\) 必须落在两条直线
\(
y=x+0.5,\qquad y=x-0.5
\)
之间的带状区域里。故所求概率为 \(1-\frac14=\frac34\)。
\end{solution}

\begin{exercise}\label{exer:1-37}
从 \([0,1]\) 中随机地取两个数，其积大于 \(\dfrac{1}{4}\)，其和小于 \(\dfrac{5}{4}\) 的概率为 \underline{\hspace{2cm}}。
\end{exercise}

\begin{solution}
应填 \(\dfrac{15}{32}-\dfrac{1}{2}\ln 2\)。设两个数分别为 \(x,y\)，则 \((x,y)\) 在正方形 \(\Omega=[0,1]^2\) 内均匀分布。所求事件为
\(A = \left\{(x,y)\,\middle|\, x+y<\frac{5}{4},\ xy>\frac{1}{4}\right\}.\)
曲线 \(xy=\dfrac{1}{4}\) 与 \(x+y=\dfrac{5}{4}\) 交于 \(x=\dfrac{1}{4}\) 和 \(x=1\)。对 \(x\in\left[\dfrac{1}{4},1\right]\)，\(y\) 的范围从 \(\dfrac{1}{4x}\) 到 \(\dfrac{5}{4}-x\)：
\begin{align*}
P(A) &= \int_{1/4}^{1}\left(\frac{5}{4}-x-\frac{1}{4x}\right)\mathrm{d}x = \left[\frac{5}{4}x-\frac{1}{2}x^2-\frac{1}{4}\ln x\right]_{1/4}^{1} = \frac{3}{4}-\left(\frac{9}{32}+\frac{1}{2}\ln 4\right) = \frac{15}{32}-\frac{1}{2}\ln 2.
\end{align*}
\end{solution}

\section{条件概率}

\subsection{条件概率的定义}

\begin{definition}[条件概率] \label{def:conditional-prob}
设 \(A, B\) 是两个事件，且 \(P(A) > 0\)，称
\(
P(B|A) = \frac{P(AB)}{P(A)}
\)
为在事件 \(A\) 发生的条件下事件 \(B\) 发生的\textbf{条件概率}。
\end{definition}

\begin{property}\label{prop:conditional-prob}
条件概率 \(P(\cdot|A)\) 满足概率的三条公理：
\begin{enumerate}
    \item \textbf{非负性}：对于每一事件 \(B\)，有 \(P(B|A) \ge 0\)；
    \item \textbf{规范性}：\(P(\Omega|A) = 1\)；
    \item \textbf{可列可加性}：设 \(B_1, B_2, \dots, B_n, \dots\) 是两两互斥的事件，则
    \[
    P\!\left(\bigcup_{i=1}^{\infty} B_i \Big| A\right) = \sum_{i=1}^{\infty} P(B_i|A).
    \]
\end{enumerate}
因此，概率的一切性质和公式对条件概率都适用。
\end{property}

\begin{theorem}[乘法公式] \label{thm:multiplication}
设 \(P(A) > 0\)，则
\(
P(AB) = P(B|A)\,P(A).
\)
一般地，设 \(A_1, A_2, \dots, A_n\) 为 \(n\) 个事件（\(n \ge 2\)），且 \(P(A_1 A_2 \cdots A_{n-1}) > 0\)，则
\[
P(A_1 A_2 \cdots A_n) = P(A_1)\,P(A_2|A_1)\,P(A_3|A_1 A_2) \cdots P(A_n|A_1 A_2 \cdots A_{n-1}).
\]
\end{theorem}

\begin{theorem}[全概率公式] \label{thm:total-probability}
设 \(B_1, B_2, \dots, B_n\) 是一组完备事件组（两两互斥且并为全集），\(A\) 为任一事件，且 \(P(B_i) > 0\ (i=1,2,\dots,n)\)，则
\(
P(A) = \sum_{i=1}^{n} P(B_i)\,P(A|B_i).
\)
\end{theorem}

\begin{theorem}[贝叶斯公式] \label{thm:bayes}
设 \(B_1, B_2, \dots, B_n\) 是一组完备事件组，\(A\) 为任一事件，且 \(P(A) > 0\)，\(P(B_k) > 0\ (k=1,2,\dots,n)\)，则
\[
P(B_k|A) = \frac{P(B_k)\,P(A|B_k)}{P(A)} = \frac{P(B_k)\,P(A|B_k)}{\sum_{i=1}^{n} P(B_i)\,P(A|B_i)}, \quad k=1,2,\dots,n.
\]
\end{theorem}

\begin{note}
全概率公式与贝叶斯公式的记忆技巧：
\begin{itemize}
    \item \textbf{全概率公式}——由因求果。将 \(A\) 视为结果，\(B_i\) 视为导致 \(A\) 发生的原因，已知各原因的概率 \(P(B_i)\) 及各原因导致结果 \(A\) 的概率 \(P(A|B_i)\)，求结果 \(A\) 发生的概率 \(P(A)\)。
    \item \textbf{贝叶斯公式}——执果寻因。已知结果 \(A\) 已经发生，反推是由第 \(k\) 个原因 \(B_k\) 导致的概率 \(P(B_k|A)\)。
\end{itemize}
\end{note}

\begin{exercise}\label{exer:1-8}
某人外出可以乘坐飞机、火车、轮船、汽车四种交通工具，其概率依次为 \(0.05, 0.15, 0.30, 0.50\)，而乘坐这几种交通工具能如期到达的概率依次为 \(0.80, 0.70, 0.60, 0.90\)。求：
\par
(1) 该人如期到达的概率；\par
(2) 已知该人误期到达，求他是乘坐火车的概率。
\end{exercise}

\begin{solution}
设 \(A_i\ (i=1,2,3,4)\) 分别表示乘坐飞机、火车、轮船、汽车，\(B\) 表示如期到达。\\
(1) 利用全概率公式：
\[
P(B) = \sum_{i=1}^{4} P(A_i)\,P(B|A_i) = 0.05 \times 0.80 + 0.15 \times 0.70 + 0.30 \times 0.60 + 0.50 \times 0.90 = 0.775.
\]
(2) 利用贝叶斯公式：
\[
P(A_2|\overline{B}) = \frac{P(A_2)\,P(\overline{B}|A_2)}{P(\overline{B})} = \frac{0.15 \times (1-0.70)}{1-0.775} = \frac{0.045}{0.225} = 0.2.
\]
\end{solution}


\begin{exercise}\label{exer:1-10}
设随机事件 \(A, B\) 满足 \(P(A)=P(B)=\dfrac{1}{2}\) 和 \(P(A \cup B)=1\)，则（）。\par
A. \(A \cup B = \Omega\)\quad
B. \(AB = \varnothing\)\quad
C. \(P(\overline{A} \cup \overline{B})=1\)\quad
D. \(P(A-B)=0\)
\end{exercise}

\begin{solution}
\[
P(A\cup B)=P(A)+P(B)-P(AB)\implies
1=\frac12+\frac12-P(AB)\implies P(AB)=0.
\]
对 A，\(P(A\cup B)=1\) 只说明事件 \(A\cup B\) 的概率为 1，并不能推出集合恒等式 \(A\cup B=\Omega\)。概率为 1 的事件不一定就等于整个样本空间，所以 A 不对。

对 B，同理，\(P(AB)=0\) 只说明交事件是零概率事件，不能必然推出 \(AB=\varnothing\)。零概率事件在连续型情形下完全可能非空，所以 B 不对。

对 D，由
\(
P(A-B)=P(A)-P(AB)=\frac12-0=\frac12\neq 0,
\)
故 D 不对。

最后看 C。利用德摩根律，\(\overline{A}\cup\overline{B}=\overline{AB}\)，于是
\[
P(\overline{A}\cup\overline{B})=P(\overline{AB})=1-P(AB)=1-0=1.
\]
因此正确选项是 C。
\end{solution}


\begin{exercise}\label{exer:1-13}
某人给四位亲友各写一封信，然后随机地装入 4 个写好地址的信封中，且每个信封装一封信。求：
\par
(1) 4 封信都装对的概率；\par
(2) 4 封信都装错的概率。
\end{exercise}
\begin{solution}
设事件 \(A_i\ (i=1,2,3,4)\) 为"第 \(i\) 封信装对了"，\(A\) 为"4 封信都装对了"，\(B\) 为"4 封信都装错了"。\\
(1) 利用乘法公式：
\[
\begin{aligned}
P(A) &= P(A_1 A_2 A_3 A_4) = P(A_1)\,P(A_2|A_1)\,P(A_3|A_1 A_2)\,P(A_4|A_1 A_2 A_3) = \frac{1}{4} \times \frac{1}{3} \times \frac{1}{2} \times \frac{1}{1} = \frac{1}{24}.
\end{aligned}
\]
(2) 先递推求装错信封的情况总数。记 \(D_n\) 为 \(n\) 封信全部装错的方案数。看第 \(n\) 封信，若它装入第 \(k\) 个信封（\(k\ne n\)），则第 \(k\) 封信有两类可能：装入第 \(n\) 个信封，余下 \(n-2\) 封形成错排；不装入第 \(n\) 个信封，则把第 \(n\) 个信封看作第 \(k\) 封信的禁位，余下形成 \(n-1\) 个对象的错排。故
\[
  D_n=(n-1)(D_{n-1}+D_{n-2}),~~D_1=0,D_2=1,D_3=2(D_2+D_1)=2,D_4=3(D_3+D_2)=9.
\]
所有装法共有 \(4!\) 种，所以
\[
  P(B)=\frac{D_4}{4!}=\frac{9}{24}=\frac38.
\]
\end{solution}

\begin{exercise}\label{exer:1-15}
一批产品有 10 个正品和 2 个次品，任意抽取两次，每次抽一个，抽出后不再放回，则第二次抽出的是次品的概率为 \(\underline{\hspace{2cm}}\)。
\end{exercise}

\begin{solution}
记
\(
A_1=\{\text{第一次抽到次品}\},\qquad A_2=\{\text{第二次抽到次品}\}.
\)
按第一次的结果分类，有
\[
P(A_2)=P(A_1)P(A_2\mid A_1)+P(\overline{A_1})P(A_2\mid \overline{A_1})
=\frac{2}{12}\cdot\frac{1}{11}+\frac{10}{12}\cdot\frac{2}{11}
=\frac16.
\]
\end{solution}

\begin{exercise}\label{exer:1-16}
设有两批数量相同的零件，已知有一批产品全部合格，另一批产品有 25\% 不合格。从两批产品中任取 1 只，经检验是正品，放回原处，并在原所在批次再取 1 只，求这只产品是次品的概率。
\end{exercise}

\begin{solution}
对于全部合格品而言，它来自第一批（全合格）的概率为 \(\frac{4}{7}\)，来自不全合格的一批（有 25\% 不合格）的概率为 \(\frac{3}{7}\)。次品概率为$\frac{3}{7}\times\frac{1}{4}=\frac{3}{28}$
\end{solution}


\begin{exercise}\label{exer:1-40}
口袋中有 1 个球，不知其颜色是黑还是白。现再往口袋中放入 1 个白球，然后从中任意取出 1 个，发现取出的是白球。求口袋中原来那个球是白球的概率。
\end{exercise}

\begin{solution}
设
\(
A=\{\text{取出的球是白球}\},\qquad B=\{\text{原来那个球是白球}\}.
\)
由于题目没有给出“原球更可能是黑还是白”的额外信息，默认黑、白两种可能等可能，即
\[
P(B|A)=\frac{P(B)P(A|B)}{P(B)P(A|B)+P(\overline B)P(A|\overline B)}
=\frac{\frac12\times 1}{\frac12\times 1+\frac12\times\frac12}
=\frac{\frac12}{\frac12+\frac14}
=\frac{\frac12}{\frac34}
=\frac23.
\]
也可以这样说，看成是有两个口袋，一个装一白一黑，另一个装两个白球，随机挑一个袋子取出球且为白球，取出的白球有$\frac23$的概率来自于装两个白球的袋子。
\end{solution}

\begin{exercise}\label{exer:1-41}
已知装有同种零件的产品两箱，第一箱内装 50 件，其中一等品 10 件；第二箱内装 30 件，其中一等品 18 件。现从两箱中任意挑选一箱，从中先后取出两件产品（取后不放回）。已知先取出的产品是一等品，则第二次取出的产品仍然是一等品的概率。
\end{exercise}

\begin{solution}
先取出的那个一等品有$\frac14$的可能来自于第一个箱子，有$\frac34$的可能来自于第二个箱子。故
\[p = \frac{1}{4}\frac{9}{49}+\frac{3}{4}\frac{17}{29}= \dfrac{690}{1421}\]
\end{solution}


\begin{exercise}\label{exer:src-paper1-01}\label{prob:p1-01}
某卡车运送防“非典”用品下乡，顶层装 10 个纸箱，其中 5 箱民用口罩、2 箱医用口罩、3 箱消毒棉花。到目的地后发现丢失 1 箱，不知丢失哪一箱。现从剩下的 9 箱中任意打开 2 箱，结果都是民用口罩，求丢失的一箱也是民用口罩的概率。
\end{exercise}
\begin{solution}
已知看见了 2 箱民用口罩，就相当于把 5 箱民用口罩先排除掉 2 箱，剩下 8 箱里还有 3 箱民用口罩，其中一箱丢了，所以答案是 \(\frac38\)。
\end{solution}

\begin{exercise}[提高]\label{exer:1-55}
设某同学的手机在一天内收到短信数服从泊松分布 \(P(\lambda)\)，每个短信是否为垃圾短信与其到达时间独立，也与其他短信是否为垃圾短信相互独立。假设每个短信是垃圾短信的概率为 \(p\)。\par
(1) 已知一天内收到了 \(n\) 条短信，求其中恰有 \(k\) 条垃圾短信的概率（\(0\le k\le n\)）。\par
(2) 求一天内收到 \(k\) 条垃圾短信的概率（\(k=0,1,2,\dots\)）。
\end{exercise}

\begin{solution}
(1) 记 \(X\) 为一天内收到短信的总数，\(Y\) 为其中垃圾短信的条数。已知 \(X=n\) 时，每条短信独立地以概率 \(p\) 为垃圾短信，故
\(
P(Y=k\mid X=n) = \mathrm{C}_n^k\,p^k\,(1-p)^{n-k}.
\)

(2) 由全概率公式：
\begin{align*}
P(Y=k) &= \sum_{n=k}^{\infty} P(Y=k\mid X=n)\,P(X=n) = \sum_{n=k}^{\infty} \mathrm{C}_n^k\,p^k\,(1-p)^{n-k}\,\frac{\lambda^n}{n!}\,e^{-\lambda} \\
&= \frac{p^k\,e^{-\lambda}\,\lambda^k}{k!}\sum_{n=k}^{\infty}\frac{[\lambda(1-p)]^{n-k}}{(n-k)!} = \frac{(\lambda p)^k}{k!}\,e^{-\lambda}\,e^{\lambda(1-p)} = \frac{(\lambda p)^k}{k!}\,e^{-\lambda p}.
\end{align*}
\end{solution}

\begin{exercise}[提高]
某工厂进行产品质量盲测。现有外观相同的 \(A,B\) 两个初始样本盒，\(A\) 盒装有质量指标 \(1,5,10\) 的产品各一件，\(B\) 盒装有指标 \(2,5,8\) 的产品各一件。现新增 \(n\) 个均含 3 件产品的新样本盒，设指标 5 的产品为标准件。甲先从这 \(n+2\) 个盒中随机选择一个，乙再从剩下的 \(n+1\) 个盒中随机选择一个抽取。记事件 \(E_1\) 为"甲抽到标准件"，事件 \(E_2\) 为"乙抽到标准件"。设事件 \(E\) 为"甲选择 \(A\) 盒"。证明："事件 \(E_1,E_2,E\) 两两独立"当且仅当"每个新盒中恰有 1 件标准件"。
\end{exercise}

\begin{solution}
设第 \(i\) 个新盒中标准件个数为 \(s_i\)，其中 \(s_i\in\{0,1,2,3\}\)。两个旧盒 \(A,B\) 中标准件个数均为 \(1\)。记 \(S=s_1+s_2+\cdots+s_n\)。从一个随机选出的盒中再随机抽一件，抽到标准件的概率等于“平均每盒标准件数”除以 \(3\)，故 \(P(E_1)=\dfrac{2+S}{3(n+2)}\)。若已知事件 \(E\) 发生，即甲选择 \(A\) 盒，则 \(P(E_1|E)=\dfrac13\)。
由于 \(E_1\) 与 \(E\) 独立，必须有 \(P(E_1|E)=P(E_1)\)，因此
 \(\dfrac13=\dfrac{2+S}{3(n+2)}\)，从而 \(S=n\)。这说明所有新盒中标准件总数为 \(n\)，即平均每个新盒有 \(1\) 件标准件。

下面进一步说明不能只要求平均为 \(1\)，而必须每盒恰有 \(1\) 件。把 \(n+2\) 个盒中的标准件个数统一记为 \(c_1,c_2,\dots,c_{n+2}\)，其中两个旧盒对应的 \(c_i\) 均为 \(1\)，新盒对应 \(s_1,\dots,s_n\)。由上面已得 \(\sum_{i=1}^{n+2}c_i=2+S=n+2\)。
若甲已抽到标准件，则甲选中第 \(i\) 个盒的概率与该盒标准件个数 \(c_i\) 成正比，即
\[
P(\text{甲选中第 }i\text{ 个盒}\mid E_1)
=\frac{c_i}{\sum_{j=1}^{n+2}c_j}
=\frac{c_i}{n+2}.
\]
在甲选中第 \(i\) 个盒后，乙只能在剩下的 \(n+1\) 个盒中选，剩余标准件总数为 \(n+2-c_i\)，所以
\[
P(E_2|E_1)
=\sum_{i=1}^{n+2}\frac{c_i}{n+2}\cdot
\frac{n+2-c_i}{3(n+1)}
=\frac{(n+2)^2-\sum_{i=1}^{n+2}c_i^2}{3(n+2)(n+1)}.
\]
又由 \(\sum c_i=n+2\) 知 \(P(E_2)=\dfrac{n+2}{3(n+2)}=\dfrac13\)。
若 \(E_1,E_2\) 独立，则 \(P(E_2|E_1)=P(E_2)\)，从而
\[
\frac{(n+2)^2-\sum_{i=1}^{n+2}c_i^2}{3(n+2)(n+1)}
=\frac13,
\]
化简得 \(\sum_{i=1}^{n+2}c_i^2=n+2\)。
另一方面 \(\sum c_i=n+2\)。由于每个 \(c_i\in\{0,1,2,3\}\)，有 \(c_i^2\ge c_i\)，且等号仅当 \(c_i=0\) 或 \(1\) 时成立。于是所有 \(c_i\) 只能取 \(0\) 或 \(1\)。又因为共有 \(n+2\) 个盒且 \(\sum c_i=n+2\)，故每个 \(c_i=1\)。因此每个新盒中恰有 \(1\) 件标准件。

反过来，若每个新盒中恰有 \(1\) 件标准件，则所有 \(n+2\) 个盒都具有相同结构：每盒 \(3\) 件，其中 \(1\) 件为标准件。于是 \(P(E_1)=P(E_2)=\dfrac13,\ P(E)=\dfrac1{n+2}\)。
知道甲是否选择 \(A\) 盒，不会改变甲抽到标准件的概率，故 \(E_1\) 与 \(E\) 独立。若甲选择 \(A\) 盒，则乙面对的剩余 \(n+1\) 个盒仍然都是每盒 \(1\) 件标准件，故 \(P(E_2|E)=\frac13=P(E_2)\)，即 \(E_2\) 与 \(E\) 独立。又甲抽到标准件不会改变剩余盒子的结构，故 \(P(E_2|E_1)=\frac13=P(E_2)\)，即 \(E_1\) 与 \(E_2\) 独立。因此
\[
P(E_1E)=P(E_1)P(E),\qquad
P(E_2E)=P(E_2)P(E),\qquad
P(E_1E_2)=P(E_1)P(E_2).
\]
故 \(E_1,E_2,E\) 两两独立。
\end{solution}


\section{事件的独立性}

\subsection{独立性的定义}

特别地，对于三个事件 \(A, B, C\)，相互独立要求同时满足以下四个等式：
\[
\begin{cases}
P(AB) = P(A)\,P(B), \\
P(AC) = P(A)\,P(C), \\
P(BC) = P(B)\,P(C), \\
P(ABC) = P(A)\,P(B)\,P(C).
\end{cases}
\]
若仅满足前三个等式（即任意两个事件独立），则称 \(A, B, C\) \textbf{两两独立}。

\begin{property}\label{prop:independence-equiv}
\(A\) 与 \(B\) 独立（\(P(AB)=P(A)\,P(B)\)）的等价条件（设 \(0 < P(A) < 1\)）：
\begin{enumerate}
    \item \(A\) 与 \(B\) 独立 \(\Leftrightarrow\) \(A\) 与 \(\overline{B}\) 独立 \(\Leftrightarrow\) \(\overline{A}\) 与 \(B\) 独立 \(\Leftrightarrow\) \(\overline{A}\) 与 \(\overline{B}\) 独立；
    \item \(A\) 与 \(B\) 独立 \(\Leftrightarrow\) \(P(B|A) = P(B|\overline{A}) = P(B)\)；
    \item \(A\) 与 \(B\) 独立 \(\Leftrightarrow\) \(P(B|A) + P(\overline{B}|\overline{A}) = 1\)。
\end{enumerate}
\end{property}

\begin{note}
两两独立 \(\nRightarrow\) 相互独立。三个事件两两独立只保证 \(P(AB)=P(A)\,P(B)\)、\(P(AC)=P(A)\,P(C)\)、\(P(BC)=P(B)\,P(C)\)，但不保证 \(P(ABC)=P(A)\,P(B)\,P(C)\)。
\end{note}

\begin{exercise}\label{exer:src-2015-spring-a-ii-01}
设 \(A,B,C\) 三个事件两两独立，则 \(A,B,C\) 相互独立的充分必要条件是（\quad）。\par
A. \(A\) 与 \(BC\) 独立\quad
B. \(AB\) 与 \(A\cup C\) 独立\par
C. \(AB\) 与 \(AC\) 独立\quad
D. \(A\cup B\) 与 \(A\cup C\) 独立
\end{exercise}
\begin{solution}
由于 \(A,B,C\) 两两独立，已经有
\[
P(AB)=P(A)P(B),\quad P(AC)=P(A)P(C),\quad P(BC)=P(B)P(C).
\]
要使三事件相互独立，还需补上\(P(ABC)=P(A)P(B)P(C)\)，事实上
而 \(A\) 与 \(BC\) 独立正是
\[
P(A)P(B)P(C)=P(A)P(BC)=P(A\cdot BC)
\]
所以 A 是充分必要条件。
\end{solution}


\begin{exercise}\label{exer:1-18}
设 \(A, B, C\) 为三个两两独立的随机事件，\(P(A)=P(B)=P(C)\)，\(P(A \cup B \cup C)=\dfrac{23}{25}\)，\(P(\overline{A} \cup \overline{B} \cup \overline{C})=\dfrac{14}{25}\)，求 \(P(A)\)。
\end{exercise}

\begin{solution}
\[
\begin{aligned}
P(\overline{A} \cup \overline{B} \cup \overline{C}) &= P(\overline{A})+P(\overline{B})+P(\overline{C}) - P(\overline{A}\,\overline{B}) - P(\overline{B}\,\overline{C}) - P(\overline{A}\,\overline{C}) + P(\overline{A}\,\overline{B}\,\overline{C}) \\
&= 3P(\overline{A}) - 3[P(\overline{A})]^2 + P(\overline{A}\,\overline{B}\,\overline{C})=3P(\overline{A}) - 3[P(\overline{A})]^2 +1-P(A \cup B \cup C)
\end{aligned}
\]
解得 \(P(\overline{A}) = \dfrac{1}{5}\) 或 \(\dfrac{4}{5}\)。又 \(P(\overline{A}) \le P(\overline{A} \cup \overline{B} \cup \overline{C}) = \dfrac{14}{25}\)，故 \(P(\overline{A}) = \dfrac{1}{5}\)，即 \(P(A) = \dfrac{4}{5}\)。
\end{solution}

\begin{exercise}[提高]\label{exer:1-58}
设 \(A,B,C\) 相互独立，且 \(P(A)\neq 0\)，\(0<P(C)<1\)。则下面四个事件中不独立的是（\quad）。\par
A. \(\overline{A}\cup B\) 与 \(C\)\quad
B. \(\overline{AC}\) 与 \(\overline{C}\)\quad
C. \(\overline{A-B}\) 与 \(\overline{C}\)\quad
D. \(\overline{AB}\) 与 \(\overline{C}\)
\end{exercise}

\begin{solution}
应选 B。A、C、D 中左边的事件均由 \(A,B\) 构成，而 \(A,B\) 各自与 \(C\) 独立，因此由 \(A,B\) 经过并、交、差运算构成的新事件也与 \(C\) 独立。
\end{solution}

\section{独立重复试验与伯努利概型}

\begin{definition}[\(n\) 重伯努利概型] \label{def:bernoulli}
如果随机试验满足以下条件：
\begin{enumerate}
    \item 每次试验只有两个可能结果，记为 \(A\)（成功）和 \(\overline{A}\)（失败）；
    \item 每次试验中 \(P(A)=p\)（\(0 < p < 1\)）保持不变；
    \item 各次试验相互独立，
\end{enumerate}
则将这种试验独立重复进行 \(n\) 次的概率模型称为 \(\mathbf{n}\) \textbf{重伯努利概型}。
在 \(n\) 重伯努利概型中，事件 \(A\) 恰好发生 \(k\) 次的概率为
\[
P\{A \text{ 恰好发生 } k \text{ 次}\} = \mathrm{C}_n^k\, p^k\, (1-p)^{n-k}, \quad k=0,1,\dots,n.
\]
如果用 \(X\) 表示 \(n\) 重伯努利概型中事件 \(A\) 发生的次数，则 \(X\) 服从参数为 \(n, p\) 的\textbf{二项分布}，记为 \(X \sim B(n,p)\)。
\end{definition}


\begin{exercise}\label{exer:1-59}
设甲、乙两人轮流射击同一目标，直到射中为止，每次命中率分别为 0.3 和 0.4，甲先射击。\par
(1) 目标是由甲射中的概率 \(p_1\)；\par
(2) 击中目标者射击了 \(k\) 次的概率 \(p_2\)。
\end{exercise}

\begin{solution}
(1) 甲在第一轮命中的概率为 0.3；若第一轮都未命中（概率 \(0.7\times 0.6=0.42\)），则回到初始状态。故
\[
p_1 = 0.3 + 0.42\times p_1 \implies p_1 = \frac{0.3}{1-0.42} = \frac{15}{29} 
\]

(2) "击中者射击了 \(k\) 次"意味着该射手在前 \(k-1\) 轮均未命中，第 \(k\) 次命中。若甲击中：前 \(k-1\) 轮甲乙均未命中（概率 \(0.42^{k-1}\)），甲第 \(k\) 次命中（概率 0.3）。若乙击中：前 \(k-1\) 轮均未命中（概率 \(0.42^{k-1}\)），甲第 \(k\) 次未命中（概率 0.7），乙第 \(k\) 次命中（概率 0.4）。故
\[
p_2 = 0.42^{k-1}\times 0.3 + 0.42^{k-1}\times 0.7\times 0.4 = 0.42^{k-1}\times 0.58 \quad (k=1,2,\dots).
\]
\end{solution}

\begin{exercise}\label{exer:src-2007-a-09}
房间里有 \(10\) 个人，分别佩戴 \(1\) 到 \(10\) 号纪念章。任选 \(3\) 人记录其纪念章号码，求下列事件的概率：\par
(1) \(A=\{\text{最小号码为 }6\}\)；\par
(2) \(B=\{\text{所选三个号码不同时含有 }4\text{ 和 }6\}\)。
\end{exercise}
\begin{solution}
总取法数为
\(
\binom{10}{3}=120.
\)

(1) 最小号码为 \(6\)，说明必须取到 \(6\)，且另外两个号码从 \(7,8,9,10\) 中取。因此
\(P(A)=\frac{\binom42}{\binom{10}{3}}=\frac6{120}=\frac1{20}.\)

(2) “不同时含有 \(4\) 和 \(6\)”的对立事件是“三个号码中既含 \(4\) 又含 \(6\)”。若同时含 \(4\) 和 \(6\)，第三个号码可从其余 \(8\) 个号码中任取，故
\(
P(B)=1-\frac8{\binom{10}{3}}
=1-\frac8{120}
=\frac{14}{15}.
\)
\end{solution}


\begin{exercise}\label{exer:src-2008-probstat-2a-ii-01}
设 \(X,Y\) 为随机变量，且
\[
P\{X\ge0,Y\ge0\}=\frac37,\qquad
P\{X\ge0\}=P\{Y\ge0\}=\frac47.
\]
求 \(P\{\max(X,Y)\ge0\}\)。
\end{exercise}
\begin{solution}
事件 \(\{\max(X,Y)\ge0\}\) 表示 \(X,Y\) 中至少有一个不小于 \(0\)，即
\(
\{X\ge0\}\cup\{Y\ge0\}.
\)
由加法公式，
\[
\begin{aligned}
P\{\max(X,Y)\ge0\}
&=P\{X\ge0\}+P\{Y\ge0\}-P\{X\ge0,Y\ge0\}=\frac47+\frac47-\frac37
=\frac57.
\end{aligned}
\]
\end{solution}


\begin{exercise}\label{exer:src-2008-probstat-2a-ix-01}\label{exer:src-paper3-proof-01}\label{prob:p3-proof-01}
设事件 \(A,B,C\) 同时发生必导致事件 \(D\) 发生，证明：
\(
P(A)+P(B)+P(C)\le 2+P(D).
\)
\end{exercise}
\begin{solution}
连续使用\(P(A)+P(B)\le1+P(AB).\)
\[
\begin{aligned}
P(A)+P(B)+P(C)
&\le 1+P(AB)+P(C)\le 2+P(ABC)\le 2+P(D).
\end{aligned}
\]
\end{solution}


\begin{exercise}\label{exer:src-2009-a-03}
某安检系统检查时，非危险人物被误认为危险人物的概率是 \(0.02\)，而危险人物被误认为非危险人物的概率是 \(0.05\)。假设过关人中有 \(96\%\) 是非危险人物。问：\par
(1) 在被检查后认为是非危险人物的条件下，确实是非危险人物的概率；\par
(2) 如果要求对危险人物的检出率超过 \(0.999\)，至少需安设多少道相互独立的检查关卡？
\end{exercise}
\begin{solution}
（1）危险人物占比$4\%$，被认为是非危险人物的危险人物占比$4\%\times 5\%$，非危险人物占比$96\%$，被认为是非危险人物的非危险人物占比$96\%\times 98\%$，则
\begin{align*}
p=\frac{96\%\times 98\%}{96\%\times 98\%+4\%\times 5\%}=\frac{2352}{2357}\approx 0.9979
\end{align*}
(2) 对危险人物而言，单道关卡漏检概率为 \(0.05\)。若设置 \(n\) 道相互独立的关卡，则全被漏检的概率为 \(0.05^n\)，检出率为\(1-0.05^n\)，题面要求\(1-0.05^n>0.999\)，即\(0.05^n<0.001\)，因为\(0.05^2=0.0025>0.001,\qquad 0.05^3=0.000125<0.001\)，所以至少需要3道关卡。
\end{solution}
\begin{exercise}\label{exer:src-2011-a-08-2credit-polya}
袋中有 \(a\) 只红球、\(b\) 只白球、\(c\) 只黑球。从袋中任取 \(1\) 个球，观察颜色后放回，并加入 \(d\) 只与之同色的球。如此操作，求第 \(k\) 次取到红球的概率。
\end{exercise}
\begin{solution}
把白球与黑球合并看作“非红球”，初始时红球数为 \(a\)，总球数为
\(
T=a+b+c.
\)
记 \(M_j\) 为第 \(j\) 次取球后袋中红球所占比例，并令 \(M_0=a/T\)。若第 \(j\) 次取球前红球比例为 \(M_{j-1}\)，则第 \(j\) 次取到红球的条件概率就是 \(M_{j-1}\)。取到红球后比例变为 \(\dfrac{R+d}{T_j+d}\)，取到非红球后比例变为 \(\dfrac{R}{T_j+d}\)，其中 \(R/T_j=M_{j-1}\)。于是
\(
E(M_j\mid M_{j-1})
=M_{j-1}\frac{R+d}{T_j+d}+(1-M_{j-1})\frac{R}{T_j+d}
=M_{j-1}.
\)
因此 \(E(M_{j})=E(M_{j-1})=\cdots=M_0=a/T\)。第 \(k\) 次取到红球的概率为第 \(k\) 次取球前红球比例的期望，所以
\(
P(\text{第 }k\text{ 次取到红球})=E(M_{k-1})=\frac{a}{a+b+c}.
\)
\end{solution}


\begin{exercise}\label{exer:src-2014-spring-a-01}
设有 \(n\) 个人排成一行，甲与乙是其中的两个人。求这 \(n\) 个人的任意排列中，甲与乙之间恰有 \(r\) 个人的概率。若这 \(n\) 个人围成一圈，并规定从甲到乙按顺时针方向计数，证明甲与乙之间恰有 \(r\) 个人的概率与 \(r\) 无关。
\end{exercise}
\begin{solution}
先看排成一行的情形。若甲、乙之间恰有 \(r\) 个人，则二者的位置相差 \(r+1\)。这样的无序位置对共有
\(
n-r-1
\)
对；甲、乙在这两个位置上可交换，剩余 \(n-2\) 人任意排列。因此
\(
P_1=\frac{2(n-r-1)(n-2)!}{n!}
=\frac{2(n-r-1)}{n(n-1)},\qquad r=0,1,\dots,n-2.
\)

再看围成一圈的情形。固定甲的位置，只需安排其余 \(n-1\) 人。若从甲沿顺时针方向到乙之间恰有 \(r\) 个人，则相当于在其余 \(n-2\) 人中选 \(r\) 人放在甲到乙之间并排列，其余 \(n-r-2\) 人放在乙到甲之间并排列。故
\(
P_2
=\frac{\binom{n-2}{r}r!(n-r-2)!}{(n-1)!}
=\frac{(n-2)!}{(n-1)!}
=\frac1{n-1}.
\)
所以在圆排列中，该概率与 \(r\) 无关。
\end{solution}

\begin{exercise}\label{exer:src-2008-01-3c-iv-01}\label{exer:src-paper4-iv-01}
一盒乒乓球有 \(6\) 个新球、\(4\) 个旧球。不放回抽取，每次任取一个，共取两次。发现其中之一是新球时，求另一个也是新球的概率。
\end{exercise}
\begin{solution}
设 \(A=\{\text{两次中至少有一个新球}\}\)，\(B=\{\text{两次都是新球}\}\)。则 \(B\subset A\)，所求为 \(P(B|A)\)。
\(
P(B)=\frac{\binom62}{\binom{10}2}=\frac{15}{45}=\frac13.
\)
又
\[
P(A)=1-P(\text{两次都是旧球})
=1-\frac{\binom42}{\binom{10}2}
=1-\frac6{45}
=\frac{13}{15}\implies
P(B|A)=\frac{P(B)}{P(A)}=\frac5{13}.
\]
\end{solution}


\begin{exercise}\label{exer:src-2008-01-3c-v-01}\label{exer:src-paper4-proof-01}
设 \(A,B,C\) 是不能同时发生但两两独立的随机事件，且
\(
P(A)=P(B)=P(C)=\rho.
\)
证明 \(\rho\) 可取的最大值为 \(1/2\)。
\end{exercise}
\begin{solution}
因为 \(A,B,C\) 两两独立，所以
\(
P(AB)=P(A)P(B)=\rho^2,\qquad
P(AC)=P(A)P(C)=\rho^2.
\)
又因为 \(A,B,C\) 不能同时发生，即 \(ABC=\varnothing\)，所以 \(AB\) 与 \(AC\) 互不相容。并且
\(
AB\cup AC\subset A.
\)
因此
\(
P(AB)+P(AC)=P(AB\cup AC)\le P(A),
\)
即
\(
2\rho^2\le \rho.
\)
当 \(\rho>0\) 时，有
\(
\rho\le \frac12.
\)
若 \(\rho=0\) 结论显然也成立。

下面说明 \(\rho=1/2\) 可以达到。取四个等可能样本点
\(
\Omega=\{1,2,3,4\},
\)
令
\(
A=\{1,2\},\qquad B=\{1,3\},\qquad C=\{2,3\}.
\)
则 \(P(A)=P(B)=P(C)=1/2\)，且
\(
P(AB)=P(\{1\})=\frac14=P(A)P(B),
\)
\(
P(AC)=P(\{2\})=\frac14=P(A)P(C),
\)
\(
P(BC)=P(\{3\})=\frac14=P(B)P(C),
\)
所以三事件两两独立；同时 \(ABC=\varnothing\)。因此最大值确为
\(
\rho_{\max}=\frac12.
\)
\end{solution}

\chapter{一维随机变量及其分布}

{\small
\[
\begin{cases}
\text{随机变量分布函数 } F(x) = P(X \le x) \text{ 的性质}
\begin{cases}
(1)\ \text{单调不减} \\
(2)\ 0 \le F(x) \le 1,\ F(-\infty)=0,\ F(+\infty)=1 \\
(3)\ \text{右连续: } F(x+0) = F(x)
\end{cases} \\[1em]
\text{随机变量分类及概率分布}
\begin{cases}
\text{离散型: 分布律} \\
\text{连续型: 概率密度} \\
\text{混合型: 分布函数}
\end{cases}
\ \text{满足概率分布的性质} \\[1em]
\text{常见随机变量及其分布}
\begin{cases}
\text{离散型}
\begin{cases}
(0-1)\text{分布: } P(X=k) = p^k(1-p)^{1-k}\ (k=0,1) \\
\text{二项分布: } P(X=k) = C_n^k p^k(1-p)^{n-k}\ (k=0,1,2,\dots,n) \\
\text{泊松分布: } P(X=k) = \frac{\lambda^k}{k!}e^{-\lambda}\ (k=0,1,2\cdots)
\end{cases} \\[2em]
\text{连续型}
\begin{cases}
\text{均匀分布}
\begin{cases}
X \sim U(a,b) \\
\text{性质: 区间概率等于该区间占整个区间的比例}
\end{cases} \\
\text{指数分布}
\begin{cases}
X \sim E(\lambda) \\
\text{性质: 无记忆性}
\end{cases} \\
\text{正态分布}
\begin{cases}
X \sim N(\mu,\sigma^2) \\
\text{性质}
\begin{cases}
\text{参数含义: } \mu \text{ 为位置参数(注意对称性),} \\
\sigma^2 \text{ 为形状参数} \\
\text{利用标准化求解概率}
\end{cases}
\end{cases}
\end{cases}
\end{cases} \\[1em]
\text{随机变量函数的概率分布}
\begin{cases}
\text{离散型: 关键取值及对应概率} \\
\text{连续型}
\begin{cases}
\text{定义法: 求解函数的分布函数} \\
\text{公式法: 适用单调的随机变量}
\end{cases}
\end{cases}
\end{cases}
\]
}

\section{随机变量与分布函数}

\subsection{随机变量与分布函数}

\begin{definition}[随机变量] \label{def:random-variable}
设随机试验 \(E\) 的样本空间为 \(\Omega\)。如果对每一个 \(\omega \in \Omega\)，都有唯一的实数 \(X(\omega)\) 与之对应，并且对任意实数 \(x\)，集合
\[
\{\omega \mid X(\omega) \le x,\ \omega \in \Omega\}
\]
都是随机事件，则称定义在 \(\Omega\) 上的实值单值函数 \(X(\omega)\) 为\textbf{随机变量}，简记为 \(X\)。一般用大写字母 \(X, Y, Z, \dots\) 或希腊字母 \(\xi, \eta, \zeta, \dots\) 来表示随机变量。
\end{definition}

\begin{definition}[分布函数] \label{def:distribution-function}
设 \(X\) 是随机变量，\(x\) 是任意实数，称函数
\(
F(x) = P\{X \le x\} \quad (-\infty < x < +\infty)
\)
为随机变量 \(X\) 的\textbf{分布函数}，也称 \(X\) 服从 \(F(x)\) 分布，记为 \(X \sim F(x)\)。它满足以下三条性质：
\begin{enumerate}
    \item \textbf{单调不减}：对任意实数 \(x_1 < x_2\)，有 \(F(x_1) \le F(x_2)\)；
    \item \textbf{有界性}：\(0 \le F(x) \le 1\)，且
    \[
    F(-\infty) = \lim_{x \to -\infty} F(x) = 0, \quad F(+\infty) = \lim_{x \to +\infty} F(x) = 1;
    \]
    \item \textbf{右连续}：对任意 \(x_0 \in \mathbf{R}\)，有 \(\lim_{x \to x_0^+} F(x) = F(x_0)\)。当分布函数是不连续的分段函数时，为保证右连续性，分段小区间除第一个外都应是左闭右开的（等号跟在"大于等于"端）。
\end{enumerate}
记 \(F(a^-) = \lim_{x \to a^-} F(x)\) 为 \(F\) 在点 \(a\) 处的左极限值，则：
\begin{enumerate}
    \item \(P\{X \le a\} = F(a)\)；
    \item \(P\{X < a\} = F(a^-)\)；
    \item \(P\{X = a\} = F(a) - F(a^-)\)；
    \item \(P\{a < X < b\} = F(b^-) - F(a)\)；
    \item \(P\{a < X \le b\} = F(b) - F(a)\)；
    \item \(P\{a \le X < b\} = F(b^-) - F(a^-)\)；
    \item \(P\{a \le X \le b\} = F(b) - F(a^-)\)；
    \item \(P\{X > a\} = 1 - F(a)\)。
\end{enumerate}
\end{definition}



\begin{exercise}\label{exer:2-1}
设随机变量 \(X\) 的分布函数为
\[
F(x) = \begin{cases}
0, & x < 0, \\
\dfrac{1}{2}, & 0 \le x < 1, \\
1 - e^{-x}, & x \ge 1.
\end{cases}
\]
求 \(P\{X = 0\}\)、\(P\{0 < X \le 2\}\) 和 \(P\{X > 1\}\)。
\end{exercise}

\begin{solution}
\begin{enumerate}
    \item \(P\{X = 0\} = F(0) - F(0^-) = \dfrac{1}{2} - 0 = \dfrac{1}{2}\)。
    \item \(P\{0 < X \le 2\} = F(2) - F(0) = (1 - e^{-2}) - \dfrac{1}{2} = \dfrac{1}{2} - e^{-2}\)。
    注意左端是严格不等号"$>$"，所以用 $F(2) - F(0)$；若左端是"$\geqslant$"，则应写 $F(2) - F(0^-)$。
    \item \(P\{X > 1\} = 1 - F(1) = 1 - (1 - e^{-1}) = e^{-1}\)。
\end{enumerate}
\end{solution}

\begin{exercise}\label{exer:2-2}
下列函数中，可以作为某个随机变量的分布函数的是（\quad）。\par
\begin{tabular}{ll}
A. \(F(x) = \dfrac{1}{\pi}\arctan x + \dfrac{1}{2},\ x \in \mathbf{R}\) & B. \(F(x) = \dfrac{1}{1+x^2},\ x \in \mathbf{R}\) \\
C. \(F(x) = \begin{cases} 0, & x < 0 \\ 1-x, & 0 \le x < 1 \\ 1, & x \ge 1 \end{cases}\) & D. \(F(x) = \begin{cases} \sin x, & 0 \le x \le \pi \\ 0, & \text{其他} \end{cases}\)
\end{tabular}
\end{exercise}

\begin{solution}
逐条验证分布函数的三条充要条件：
\begin{itemize}
    \item A：\(\arctan x\) 单调递增，故 \(F(x)\) 单调不减；\(F(-\infty)=0\)，\(F(+\infty)=1\)；\(\arctan x\) 连续，故右连续。\textbf{三条均满足}。
    \item B：\(F(x)\) 在 \((-\infty,+\infty)\) 上\textbf{不是}单调不减的（例如 \(F(0)=1\)，\(F(1)=\frac{1}{2}\)），排除。
    \item C：在 \([0,1)\) 上 \(F(x)=1-x\) 单调\textbf{递减}，不满足单调不减，排除。
    \item D：\(\sin x\) 在 \([0,\pi]\) 上不是单调函数，不满足单调不减，排除。
\end{itemize}
\end{solution}

\begin{exercise}\label{exer:2-3}
设随机变量 \(X\) 的分布函数为
\[
F(x) = \begin{cases}
a + b e^{-x}, & x > 0, \\
0, & x \le 0.
\end{cases}
\]
求常数 \(a, b\) 及概率 \(P\{|X| < 2\}\)。
\end{exercise}

\begin{solution}
由分布函数的有界性，有
\[
F(+\infty) = \lim_{x \to +\infty} F(x) = \lim_{x \to +\infty} (a + b e^{-x}) = a = 1,
\]
得 \(a = 1\)。分布函数右连续，在 \(x=0\) 处有
\(
F(0) = \lim_{x \to 0^+} F(x) = a + b = 0,
\)
得 \(b = -1\)。所以
\[
F(x) = \begin{cases}
1 - e^{-x}, & x > 0, \\
0, & x \le 0.
\end{cases}
\implies
P\{|X| < 2\} = P\{-2 < X < 2\} = F(2^-) - F(-2) = 1 - e^{-2}.
\]
\end{solution}

\section{离散型随机变量和连续型随机变量}

\begin{definition}[离散型随机变量] \label{def:discrete-rv}
如果随机变量 $X$ 只可能取有限个或可列无限个值 $x_1, x_2, \dots$，则称 $X$ 为\textbf{离散型随机变量}。称
\(
P\{X = x_i\} = p_i, \quad i = 1, 2, \dots
\)
为 $X$ 的\textbf{分布律}（或分布列、概率分布），记为 $X \sim p_i$。常用表格形式表示：
\[
\begin{array}{c|cccc}
X & x_1 & x_2 & \cdots \\
\hline
P & p_1 & p_2 & \cdots
\end{array}
\]
\end{definition}

\begin{property}\label{prop:discrete-rv-condition}
数列 $\{p_i\}\ (i=1,2,\dots)$ 是某离散型随机变量的分布律的\textbf{充要条件}：
\begin{enumerate}
    \item $p_i \ge 0\ (i = 1, 2, \dots)$；
    \item $\displaystyle\sum_{i=1}^{\infty} p_i = 1$。
\end{enumerate}
设离散型随机变量 $X$ 的分布律为 $P\{X = x_i\} = p_i\ (i=1,2,\dots)$，则：
\begin{enumerate}
    \item $X$ 的分布函数为 $\displaystyle F(x) = \sum_{x_i \le x} p_i$，它是\textbf{阶梯函数}，在 $x = x_i$ 处有跳跃，跳跃高度为 $p_i$；
    \item $P\{X = x_i\} = F(x_i) - F(x_i^-)$；
    \item 对实数轴上的任一集合 $B$，有 $\displaystyle P\{X \in B\} = \sum_{x_i \in B} p_i$。
\end{enumerate}
\end{property}

\begin{definition}[连续型随机变量] \label{def:continuous-rv}
如果随机变量 $X$ 的分布函数 $F(x)$ 可以表示为
\[
F(x) = \int_{-\infty}^{x} f(t)\,\mathrm{d}t \quad (x \in \mathbf{R}),
\]
其中 $f(x)$ 是非负可积函数，则称 $X$ 为\textbf{连续型随机变量}，称 $f(x)$ 为 $X$ 的\textbf{概率密度函数}（简称概率密度），记为 $X \sim f(x)$。
\end{definition}


\begin{property}\label{prop:pdf-property}
函数 $f(x)$ 为某一随机变量的概率密度的\textbf{充要条件}：$f(x) \ge 0$且 $\displaystyle\int_{-\infty}^{+\infty} f(x)\,\mathrm{d}x = 1$。设连续型随机变量 $X \sim f(x)$，则：
\begin{enumerate}
    \item 对任意实数 $a < b$，$\displaystyle P\{a < X \le b\} = \int_a^b f(x)\,\mathrm{d}x = F(b) - F(a)$；
    \item \textbf{连续型随机变量取一点的概率为零}：$P\{X = c\} = 0$，从而
    \(
    P\{a < X < b\} = P\{a \le X < b\} = P\{a < X \le b\} = P\{a \le X \le b\};
    \)
    \item 若 $f(x)$ 在点 $x$ 处连续，则 $F'(x) = f(x)$。更一般地，
    \[
    f(x) = \begin{cases}
    0, & x \text{ 为 } F(x) \text{ 的不可导点}, \\
    F'(x), & x \text{ 为 } F(x) \text{ 的可导点}.
    \end{cases}
    \]
    \item 连续型随机变量的分布函数一定是\textbf{连续}函数（但反之不成立：$F(x)$ 连续的随机变量不一定是连续型的）。同一个分布函数的概率密度函数并不唯一——改变有限个点的值不影响积分结果。
\end{enumerate}
\end{property}

\begin{exercise} 
设 $F(x)$ 为某连续型随机变量的分布函数，$f(x)$ 为相应的密度函数，则（\quad）。\\
A.$F(x)$ 连续但不一定可导；~~B.$F(x)$ 可导；~~C.$f(x)$ 连续；~~D.$0\leqslant f(x)\leqslant1$
\end{exercise}
\begin{solution}
答案选\textbf{(A)}。分析：
\begin{itemize}
    \item 连续型随机变量的分布函数 $F(x)$ 一定是连续的；
    \item 但 $F(x)$ 仅在密度 $f(x)$ 的连续点处可导；
    \item 密度函数 $f(x)$ 可以不连续（例如分段均匀分布）；
    \item 密度函数 $f(x)$ 可以大于1（例如 $U(0,0.5)$ 的密度为2）。
\end{itemize}
\end{solution}

\begin{exercise}\label{exer:2-4}
设连续型随机变量 $X$ 的分布函数为
\[
F(x) = \begin{cases}
0, & x < 0, \\
Ax^2 + \dfrac{2}{3}x, & 0 \le x < 1, \\
1, & x \ge 1.
\end{cases}
\]
求常数 $A$、$X$ 的概率密度函数 $f(x)$ 以及 $P\{|X| < 0.5\}$。
\end{exercise}

\begin{solution}
由连续型随机变量分布函数的连续性，有
\(
\lim_{x \to 1^-} F(x) = A + \frac{2}{3} = 1 \implies A = \frac{1}{3}.
\)
在 $[0,1)$ 上，$F(x) = \dfrac{x^2}{3} + \dfrac{2x}{3}$，故
\(
f(x) = F'(x) = \frac{2x}{3} + \frac{2}{3}, \quad x \in [0,1),
\)
即
\[
f(x) = \begin{cases}
\dfrac{2x}{3} + \dfrac{2}{3}, & 0 \le x < 1, \\
0, & \text{其他}.
\end{cases}
\]
\(
P\{|X| < 0.5\} = P\{-0.5 < X < 0.5\} = F(0.5) - F(-0.5) = \left(\frac{1}{12} + \frac{1}{3}\right) - 0 = \frac{5}{12}.
\)
\end{solution}


\begin{exercise}\label{exer:2-26}
设随机变量 $X$ 的分布函数为
\[
F(x) = \begin{cases}
0, & x<-1, \\
0.4, & -1 \le x<1, \\
0.8, & 1 \le x<3, \\
1, & x \ge 3.
\end{cases}
\]
求：(1) $X$ 的分布律；(2) $P\{X<2 \mid X \neq 1\}$。
\end{exercise}

\begin{solution}
(1) 先判断出来是离散型随机变量：$P\{X=k\} = F(k) - F(k^-)$。
\begin{align*}
P\{X=-1\} &= F(-1) - F(-1^-) = 0.4 - 0 = 0.4, \\
P\{X=1\} &= F(1) - F(1^-) = 0.8 - 0.4 = 0.4, \\
P\{X=3\} &= F(3) - F(3^-) = 1 - 0.8 = 0.2.
\end{align*}
(2) 由条件概率公式：
\[
P\{X<2 \mid X \neq 1\} = \frac{P\{X<2,\, X \neq 1\}}{P\{X \neq 1\}} = \frac{P\{X=-1\}}{1-P\{X=1\}} = \frac{0.4}{0.6} = \frac{2}{3}.
\]
\end{solution}




\begin{exercise}\label{exer:2-6}
已知 $X$ 的分布函数为 $F(x)$，概率密度为 $f(x)$。当 $x \le 0$ 时，$f(x)$ 连续且 $f(x) = F(x)$。若 $F(0) = 1$，求 $f(x)$。
\end{exercise}

\begin{solution}
对 \(x>0\)，密度应为 0。下面处理 \(x\le 0\) 的部分。连续型随机变量满足 \(F'(x)=f(x)\)，解得
\(
F(x)=Ce^x,C=1,F(x)=e^x\quad(x\le 0).
\)
于是分布函数为
\[
F(x) = \begin{cases} e^x, & x < 0, \\ 1, & x \ge 0, \end{cases} \quad f(x) = \begin{cases} e^x, & x \le 0, \\ 0, & x > 0. \end{cases}
\]
\end{solution}


\begin{exercise}\label{exer:2-7}
已知随机变量 $X$ 的概率密度为
\[
f(x) = \begin{cases}
2\left(1 - \dfrac{1}{x^2}\right), & 1 \le x \le 2, \\
0, & \text{其他}.
\end{cases}
\]
求 $X$ 的分布函数 $F(x)$。
\end{exercise}

\begin{solution}
分段积分：
\begin{itemize}
    \item 当 $x < 1$ 时，$F(x) = \displaystyle\int_{-\infty}^x 0\,\mathrm{d}t = 0$；
    \item 当 $1 \le x < 2$ 时，
    \(
    F(x) = \int_1^x 2\left(1 - \frac{1}{t^2}\right)\mathrm{d}t = 2\left[t + \frac{1}{t}\right]_1^x = 2\left(x + \frac{1}{x} - 2\right);
    \)
    \item 当 $x \ge 2$ 时，$F(x) = \displaystyle\int_1^2 2\left(1 - \frac{1}{t^2}\right)\mathrm{d}t = 2\left[\frac{3}{2} - 1\right] = 1$。
\end{itemize}
\end{solution}

\begin{exercise}\label{exer:2-23}
下列函数中，可以作为某个随机变量的概率密度函数的是（\quad）。
\begin{align*}
&\text{A. } f(x) = \begin{cases} 2(1-|x|), & |x| \le 1 \\ 0, & |x|>1 \end{cases}
&&\text{B. } f(x) = \begin{cases} \frac{1}{\sqrt{2\pi}\sigma} e^{-\frac{(x-\mu)^2}{2\sigma^2}}, & x \ge 0 \\ 0, & x<0 \end{cases} \\
&\text{C. } f(x) = \begin{cases} x, & -1<x<0 \\ \frac{3}{4}x, & 0 \le x<2 \\ 0, & \text{其他} \end{cases}
&&\text{D. } f(x) = \begin{cases} \lambda e^{-\lambda x}, & x>0 \\ 0, & x \le 0 \end{cases} \ (\lambda>0)
\end{align*}
\end{exercise}

\begin{solution}
\begin{itemize}
    \item A：$\displaystyle\int_{-\infty}^{\infty} f(x)\,\mathrm{d}x = 2\int_0^1 2(1-x)\,\mathrm{d}x = 4\left[x-\frac{x^2}{2}\right]_0^1 = 2 \neq 1$，错误；
    \item B：正态分布的概率密度在 $x<0$ 时截断为0，但 $\displaystyle\int_0^{+\infty} f(x)\,\mathrm{d}x = \frac{1}{2} \neq 1$，错误；
    \item C：当 $-1<x<0$ 时，$x<0$，故 $f(x)=x<0$，不满足非负性，错误；
    \item D：指数分布的标准形式，满足 $\displaystyle\int_0^{+\infty} \lambda e^{-\lambda x}\,\mathrm{d}x = 1$，正确。
\end{itemize}
\end{solution}

\begin{exercise}\label{exer:2-24}
设随机变量 $X$ 的概率密度为
\[
f(x) = \begin{cases} A\sin x, & 0 \le x \le \pi \\ 0, & \text{其他} \end{cases}
\]
求：(1) 系数 $A$；(2) $X$ 的分布函数 $F(x)$；(3) $P\{0<X<1\}$。
\end{exercise}
\begin{solution}
(1) 由归一性条件：
\[
\int_{0}^{\pi} A\sin x\,\mathrm{d}x = A[-\cos x]_{0}^{\pi} = A(1+1) = 2A = 1,~~A = \dfrac{1}{2}
\]
(2) 当 $x < 0$ 时，$F(x) = 0$；当 $0 \le x < \pi$ 时，
\[
F(x) = \int_{0}^{x} \frac{1}{2}\sin t\,\mathrm{d}t = \frac{1}{2}[-\cos t]_{0}^{x} = \frac{1}{2}(1-\cos x);
\]
当 $x \ge \pi$ 时，$F(x) = 1$。故
\[
F(x) = \begin{cases}
0, & x<0, \\
\dfrac{1}{2}(1-\cos x), & 0 \le x < \pi, \\
1, & x \ge \pi.
\end{cases}
\]
(3)
\(
P\{0<X<1\} = \int_{0}^{1} \frac{1}{2}\sin t\,\mathrm{d}t = \frac{1}{2}(1-\cos 1).
\)
\end{solution}


\section{常见离散型随机变量及分布}

\begin{definition}[(0-1)分布]\label{def:0-1-dist}
随机变量 $X$ 只可能取 0 与 1 两个值，它的分布律为
\(
P\{X = k\} = p^k (1-p)^{1-k},\quad k=0,1\ (0 < p < 1).
\)
则称 $X$ 服从参数为 $p$ 的(0-1)分布，记为 $X \sim B(1,p)$。(0-1)分布是描述伯努利试验结果的计数变量——事件发生计为1，不发生计为0。0-1分布（伯努利分布）的分布函数为：
\[
F(x)=
\begin{cases}
0, & x<0, \\
1-p, & 0\leqslant x<1, \\
1, & x\geqslant1.
\end{cases}
\]
解释：
\begin{itemize}
    \item 当 $x<0$ 时，$F(x)=P\{X\leqslant x\}=0$；
    \item 当 $0\leqslant x<1$ 时，$F(x)=P\{X=0\}=1-p$；
    \item 当 $x\geqslant1$ 时，$F(x)=P\{X=0\}+P\{X=1\}=1-p+p=1$。
\end{itemize}
\end{definition}

\begin{definition}[二项分布]\label{def:binomial-dist}
若随机变量 $X$ 的分布律为
\(
P\{X = k\} = \mathrm{C}_n^k p^k (1-p)^{n-k},\quad k=0,1,\dots,n,
\)
其中 $0 < p < 1$，则称 $X$ 服从参数为 $(n,p)$ 的二项分布，记为 $X \sim B(n,p)$。若 $X \sim B(n,p)$，则令 $Y = n - X$，有 $Y \sim B(n,1-p)$。
\end{definition}

\begin{exercise}[概率密度与二项分布结合]\label{exer:2-41}\label{exer:src-2015-spring-a-i-05}
设随机变量 $X$ 的概率密度为 $f(x)=\begin{cases}2x, & 0<x<1, \\ 0, & \text{其他},\end{cases}$，$Y$ 表示对 $X$ 的 3 次独立重复试验中事件 $\left\{X\leqslant\dfrac{1}{2}\right\}$ 出现的次数，则 $P\{Y=2\}=\underline{\hspace{2cm}}$。
\end{exercise}
\begin{solution}
\[
p = P\left\{X\leqslant\frac{1}{2}\right\} = \int_0^{\frac{1}{2}}2x\,\mathrm{d}x = \left[x^2\right]_0^{\frac{1}{2}} = \frac{1}{4}.
\]
$Y$ 服从二项分布 $B\left(3,\dfrac{1}{4}\right)$，故
\[
P\{Y=2\} = \mathrm{C}_3^2\left(\frac{1}{4}\right)^2\left(\frac{3}{4}\right) = 3\cdot\frac{1}{16}\cdot\frac{3}{4} = \frac{9}{64}.
\]
\end{solution}


\begin{definition}[泊松分布] \label{def:poisson-dist}
若随机变量 $X$ 的分布律为
\[
P\{X = k\} = \frac{\lambda^k}{k!} e^{-\lambda},\quad \lambda > 0;\ k=0,1,2,\dots,
\]
则称 $X$ 服从参数为 $\lambda$ 的泊松分布，记为 $X \sim P(\lambda)$。设 $\lambda > 0$ 是一个常数，$n$ 是任意正整数，设 $\lim_{n \to \infty} np_n = \lambda$，则对于任一固定的非负整数 $k$，有
\[
\lim_{n \to \infty} \mathrm{C}_n^k p_n^k (1-p_n)^{n-k} = \frac{\lambda^k e^{-\lambda}}{k!}.
\]当$n$很大，$p$很小，且$\lambda = np$适中时，二项分布可用泊松分布近似，即
\[
\mathrm{C}_n^k p^k (1-p)^{n-k} \approx \frac{\lambda^k}{k!} e^{-\lambda}.
\]
一般地，当$n \ge 20,\ p \le 0.05$时，用泊松近似公式逼近二项分布效果比较好，特别当$n \ge 100,\ np \le 10$时，逼近效果更佳。这一结论在计算机出现前极大简化了二项分布的计算。
\end{definition}

\begin{exercise}[泊松分布]\label{exer:2-27}
假设某段时间内来百货公司的顾客数服从参数为 $\lambda$ 的泊松分布，而每位顾客购买电视机的概率均为 $p$，且顾客之间是否购买电视机相互独立。求该段时间内百货公司售出 $k$ 台电视机的概率（假设每位顾客至多买一台电视机）。
\end{exercise}

\begin{solution}
设 $A_i$ 表示这段时间内到达百货公司的顾客数为 $i\ (i=0,1,2,\dots)$，则 $P(A_i) = \dfrac{\lambda^i}{i!}e^{-\lambda}$。
当顾客数为 $i$ 时，售出的电视机数 $Y|A_i \sim B(i,p)$，故
\(
P\{Y=k \mid A_i\} = \mathrm{C}_i^k\, p^k(1-p)^{i-k}, \quad k \le i.
\)
由全概率公式：
\begin{align*}
P\{Y=k\} &= \sum_{i=0}^{\infty} P(A_i)\,P\{Y=k \mid A_i\} = \sum_{i=k}^{\infty} \frac{\lambda^i}{i!}e^{-\lambda} \cdot \frac{i!}{(i-k)!k!}\, p^k(1-p)^{i-k} \\
&= \frac{e^{-\lambda}(\lambda p)^k}{k!} \sum_{i=k}^{\infty} \frac{[\lambda(1-p)]^{i-k}}{(i-k)!}= \frac{e^{-\lambda}(\lambda p)^k}{k!} \cdot e^{\lambda(1-p)}=\frac{(\lambda p)^k}{k!} e^{-\lambda p},  k=0,1,2,\dots.
\end{align*}
即 $Y \sim P(\lambda p)$。这表明：若顾客数服从泊松分布，每位顾客独立地以概率 $p$ 购买，则购买总数仍服从泊松分布，参数变为 $\lambda p$。
\end{solution}

\begin{exercise}\label{exer:src-2013-a-01-01}
设 \(X\) 服从参数为 \(\lambda\) 的泊松分布，且
\(
E[(X-1)(X-2)]=1.
\)
求 \(\lambda\)。
\end{exercise}
\begin{solution}
\[
E(X^2-3X+2)=EX^2-3EX+2=DX+(EX)^2-3EX+2
=\lambda^2-2\lambda+2=1,\lambda=1
\]
\end{solution}

\begin{definition}[超几何分布] \label{def:hypergeo-dist}
若随机变量 $X$ 的分布律为
\[
P\{X=k\} = \frac{\mathrm{C}_M^k \mathrm{C}_{N-M}^{n-k}}{\mathrm{C}_N^n},\quad \max\{0,n-N+M\} \le k \le \min\{M,n\},
\]
其中 $M,N,n$ 为正整数且 $M \le N,\ n \le N$，则称 $X$ 服从参数为 $(n,N,M)$ 的超几何分布，记为 $X \sim H(n,N,M)$。
\end{definition}

\begin{exercise}[超几何分布] \label{exer:2-45}
20 个产品中有 5 个不合格品，若从中随机取出 8 个，求其中不合格品数 $X$ 的概率分布。
\end{exercise}
\begin{solution}
$X$ 服从超几何分布 $H(N=20, M=5, n=8)$，其分布律为
\[
P\{X=k\} = \frac{\mathrm{C}_5^k \mathrm{C}_{15}^{8-k}}{\mathrm{C}_{20}^8}, \quad k=0,1,2,3,4,5.
\]
具体计算：
\[
\begin{array}{c|cccccc}
\hline
X & 0 & 1 & 2 & 3 & 4 & 5 \\
\hline
P & \dfrac{\mathrm{C}_5^0\mathrm{C}_{15}^8}{\mathrm{C}_{20}^8} & \dfrac{\mathrm{C}_5^1\mathrm{C}_{15}^7}{\mathrm{C}_{20}^8} & \dfrac{\mathrm{C}_5^2\mathrm{C}_{15}^6}{\mathrm{C}_{20}^8} & \dfrac{\mathrm{C}_5^3\mathrm{C}_{15}^5}{\mathrm{C}_{20}^8} & \dfrac{\mathrm{C}_5^4\mathrm{C}_{15}^4}{\mathrm{C}_{20}^8} & \dfrac{\mathrm{C}_5^5\mathrm{C}_{15}^3}{\mathrm{C}_{20}^8} \\
\hline
\text{数值} & 0.0511 & 0.2554 & 0.3973 & 0.2384 & 0.0542 & 0.0036 \\
\hline
\end{array}
\]
其中 $\mathrm{C}_{20}^8 = 125970$。答案：$X$ 的超几何分布如上表所示。
\end{solution}

\begin{exercise}\label{exer:src-2011-a-02}
已知甲、乙两箱中装有同种产品，其中甲箱中装有 \(3\) 件合格品和 \(3\) 件次品，乙箱中仅装有 \(3\) 件合格品。从甲箱中任取 \(3\) 件产品放入乙箱后，求：\par
(1) 乙箱中次品件数 \(X\) 的数学期望；\par
(2) 从乙箱中任取一件产品是次品的概率。
\end{exercise}
\begin{solution}
(1) 从甲箱的 \(3\) 件次品、\(3\) 件合格品中不放回取 \(3\) 件，故 \(X\) 服从超几何分布：
\[
P(X=k)=\frac{\binom3k\binom3{3-k}}{\binom63},
\quad k=0,1,2,3.\implies
\begin{array}{c|cccc}
X & 0 & 1 & 2 & 3\\
\hline
P & \dfrac1{20} & \dfrac9{20} & \dfrac9{20} & \dfrac1{20}
\end{array}
\]
于是
\[
EX=0\cdot\frac1{20}+1\cdot\frac9{20}+2\cdot\frac9{20}+3\cdot\frac1{20}
=\frac32.
\]
(2) 记 \(A=\{\text{从乙箱中任取一件产品是次品}\}\)。放入后乙箱共有 \(6\) 件产品；若乙箱中次品数为 \(X=k\)，则
\(
P(A|X=k)=\frac{k}{6}.
\)
由全概率公式，
\(
P(A)=\sum_{k=0}^3P(X=k)P(A|X=k)
=\frac{EX}{6}
=\frac{1}{4}.
\)
\end{solution}

\begin{exercise}\label{exer:src-2011-ab-a-02}
袋中有 \(15\) 个球，其中 \(10\) 个红球、\(5\) 个黄球。不放回地分两次从袋中逐个取球，第一次取 \(5\) 个球，第二次取 \(6\) 个球。求：\par
(1) 第二次 \(6\) 个球中的第 \(5\) 个是红球的概率；\par
(2) 第一次 \(5\) 个球中有 \(2\) 个黄球且第二次 \(6\) 个球中有 \(4\) 个红球的概率；\par
(3) 第一次 \(5\) 个球中有 \(3\) 个红球或第二次 \(6\) 个球中有 \(2\) 个黄球的概率。
\end{exercise}
\begin{solution}
(1) 不放回抽取时每一个固定抽取位置为红球的概率都等于初始红球比例，故
\(
P=\frac{10}{15}=\frac23.
\)

(2) 设 \(A=\{\text{第一次有 }2\text{ 个黄球}\}\)，\(B=\{\text{第二次有 }4\text{ 个红球}\}\)。把 \(15\) 个位置分成第一次 \(5\) 个、第二次 \(6\) 个和剩余 \(4\) 个位置，则
\[
P(AB)
=\frac{\binom{10}{3}\binom52\binom74\binom32}
{\binom{15}{5}\binom{10}{6}}
=\frac{200}{1001}.
\]

(3) 设 \(C=\{\text{第一次有 }3\text{ 个红球}\}\)，\(D=\{\text{第二次有 }2\text{ 个黄球}\}\)。则
\[
P(C)=\frac{\binom{10}{3}\binom52}{\binom{15}{5}}=\frac{400}{1001},
\qquad
P(D)=\frac{\binom52\binom{10}{4}}{\binom{15}{6}}=\frac{420}{1001}.
\]
而 \(CD\) 正是第 (2) 问事件，所以
\(
P(C\cup D)=P(C)+P(D)-P(CD)
=\frac{400+420-200}{1001}
=\frac{620}{1001}.
\)
\end{solution}


\begin{definition}[几何分布] \label{def:geometric-dist}
若随机变量 $X$ 的分布律为
\(
P\{X = k\} = (1-p)^{k-1} p,\quad k=1,2,\dots;\ 0 < p < 1,
\)
则称 $X$ 服从参数为 $p$ 的几何分布，记为 $X \sim G(p)$。几何分布描述了在$n$重伯努利试验中，事件$A$首次出现所需的试验次数。它具有与指数分布类似的无记忆性。
\end{definition}


\begin{exercise}[几何分布的无记忆性] \label{exer:2-58}
设 $X$ 为取正整数值的离散型随机变量，证明：$X$ 是服从参数为 $p$ 的几何分布的充要条件是 $X$ 具有无记忆性，即 $P(X=k+n\mid X>k)$ 与 $k$ 无关（$k,n$ 为任意正整数）。
\end{exercise}

\begin{solution}
\textbf{必要性证明}：若 $X$ 服从几何分布，则 $P(X=k)=p(1-p)^{k-1},\ k=1,2,\dots$。对任意正整数 $k,n$：
\[
\begin{aligned}
P(X=k+n\mid X>k) &= \frac{P(X=k+n, X>k)}{P(X>k)} = \frac{P(X=k+n)}{P(X>k)} \\
&= \frac{p(1-p)^{k+n-1}}{\sum_{j=k+1}^{\infty}p(1-p)^{j-1}} = \frac{p(1-p)^{k+n-1}}{(1-p)^k} = p(1-p)^{n-1}.
\end{aligned}
\]
结果与 $k$ 无关，故几何分布具有无记忆性。\\
\textbf{充分性证明}：设 $P(X=k+1\mid X>k)=p$ 与 $k$ 无关。记 $q=1-p$，则
\(
P(X>k) = q\cdot P(X>k-1) = q^k\cdot P(X>0).
\)
由 $P(X>0)=1$ 可得
\(
P(X=k) = P(X>k-1) - P(X>k) = q^{k-1} - q^k = p(1-p)^{k-1}.
\)
这正是参数为 $p$ 的几何分布。
\end{solution}


\begin{exercise}\label{exer:src-2012-a-08-2credit}
某游戏的转盘被等分为 \(10\) 个区域，区域上分别标有 \(1,2,\dots,10\) 分。每次转动转盘后，若指针所指分数小于 \(8\)，游戏停止，并以该分数作为本次游戏的最后得分；若指针所指分数不小于 \(8\)，则继续转动转盘。设 \(X\) 为一次游戏中转动转盘的总次数，\(Y\) 为游戏停止时指针所指的分数。求：\par
(1) \((X,Y)\) 的联合分布律；\par
(2) \(X\) 与 \(Y\) 是否相互独立；\par
(3) 若每次继续转动可得 \(1\) 分奖励，总得分为 \(S=(X-1)Y\)，求 \(ES\)。
\end{exercise}
\begin{solution}
每次转动后继续的概率为
\(
P\{\text{继续}\}=P\{8,9,10\}=\frac3{10},
\)
停止且最后得分为 \(j\ (j=1,2,\dots,7)\) 的概率为 \(1/10\)。因此
\[
P\{X=i,Y=j\}=\left(\frac3{10}\right)^{i-1}\frac1{10},
\qquad i=1,2,\dots,\quad j=1,2,\dots,7.
\]

由上式求边缘分布，
\[
P\{X=i\}=\sum_{j=1}^7\left(\frac3{10}\right)^{i-1}\frac1{10}
=\frac7{10}\left(\frac3{10}\right)^{i-1},~
P\{Y=j\}=\sum_{i=1}^{\infty}\left(\frac3{10}\right)^{i-1}\frac1{10}
=\frac1{10}\cdot\frac1{1-3/10}=\frac17.
\]
于是
\[
P\{X=i\}P\{Y=j\}
=\frac7{10}\left(\frac3{10}\right)^{i-1}\cdot\frac17
=\left(\frac3{10}\right)^{i-1}\frac1{10}
=P\{X=i,Y=j\},
\]
故 \(X\) 与 \(Y\) 相互独立。

又 \(X\sim G(7/10)\)，所以
\(
E(X-1)=EX-1=\frac{10}{7}-1=\frac37.
\)
而 \(Y\) 在 \(1,2,\dots,7\) 上均匀分布，故
\(
EY=\frac{1+2+\cdots+7}{7}=4.
\)
由独立性，
\(
ES=E[(X-1)Y]=E(X-1)\,EY=\frac37\cdot4=\frac{12}{7}.
\)
\end{solution}

\begin{exercise}\label{exer:src-2012-b-08-2credit}
游园晚会推出有奖游戏。参与者交替用左手和右手抛掷套圈，第一次用左手抛掷；若第 \(i\) 次套中标杆，则记 \(i\) 分并换手继续，若未套中则游戏结束。记参与者结束时总抛掷次数为 \(X\)，所得总分为 \(Y\)。假定左手、右手套中的概率分别为 \(0.2\) 和 \(0.5\)，且各次抛掷相互独立。求：\par
(1) \(X\) 的分布律；\par
(2) 参与者的平均得分 \(EY\)。
\end{exercise}
\begin{solution}
(1) 若 \(X=2m-1\)，则前 \(2m-2\) 次均套中，第 \(2m-1\) 次左手未中，故
\[
P\{X=2m-1\}=0.8(0.2)^{m-1}(0.5)^{m-1}
=0.8(0.1)^{m-1},\qquad m=1,2,\dots.
\]
若 \(X=2m\)，则前 \(2m-1\) 次均套中，第 \(2m\) 次右手未中，故
\[
P\{X=2m\}=(0.2)^m(0.5)^{m}
=(0.1)^m,\qquad m=1,2,\dots.
\]
因此
\[
P\{X=k\}=
\begin{cases}
0.8(0.1)^{m-1}, & k=2m-1,\ m=1,2,\dots,\\
(0.1)^m, & k=2m,\ m=1,2,\dots.
\end{cases}
\]
(2) 若游戏在第 \(X\) 次抛掷后结束，则前 \(X-1\) 次套中，总分为
\(
Y=1+2+\cdots+(X-1)=\frac{X(X-1)}2.
\)
于是
\[
\begin{aligned}
EY
&=\sum_{m=1}^{\infty}(2m-1)(m-1)\,0.8(0.1)^{m-1}
 +\sum_{m=1}^{\infty}m(2m-1)(0.1)^m.
\end{aligned}
\]
令 \(r=0.1\)，利用
\[
\sum_{m=1}^{\infty}m r^m=\frac{r}{(1-r)^2},\qquad
\sum_{m=1}^{\infty}m^2 r^m=\frac{r(1+r)}{(1-r)^3},
\]
可得
\(
EY=\frac{14}{27}.
\)
\end{solution}

\begin{definition}[极值分布]\label{prop:extreme-value}
对于 $n$ 个相互独立的随机变量 $\xi_1,\xi_2,\dots,\xi_n$，其分布函数分别为 $F_1(x),F_2(x),\dots,F_n(x)$，则：
\begin{itemize}
    \item 最大值 $\eta_1 = \max\{\xi_1,\xi_2,\dots,\xi_n\}$ 的分布函数为
    \[ P\{\eta_1 \le x\} = P\left(\bigcap_{k=1}^n \xi_k \le x\right) = \prod_{k=1}^n F_k(x);~~P(\eta_1>x)=1-\prod_{k=1}^n F_k(x).\]
    \item 最小值 $\eta_2 = \min\{\xi_1,\xi_2,\dots,\xi_n\}$ 的分布函数为
    \[ P(\eta_2>x)=\prod_{k=1}^n (1-F_k(x)).;~~P\{\eta_2 \le x\} = 1 - P\{\eta_2 > x\}  = 1 - \prod_{k=1}^n (1-F_k(x))\]
\end{itemize}
\end{definition}


\begin{exercise}[最小值分布（系统寿命）] \label{exer:2-61}
一辆新车装有 4 个轮胎并配有 1 个备胎，当使用的轮胎损坏时就用备胎更换。设所有轮胎的寿命均服从参数为 $0.2$ 的指数分布且相互独立，记 $T$ 为开始使用备胎的时间，求 $T$ 的分布函数。
\end{exercise}

\begin{solution}
设四个在用轮胎的寿命分别为 $X_1,X_2,X_3,X_4$，开始使用备胎的时间为
\(T = \min\{X_1,X_2,X_3,X_4\}.\)单个轮胎的分布函数：\(
F_X(x) = 1-e^{-0.2x},\quad x\geqslant0.
\)，$T$ 的分布函数：
\[
\begin{aligned}
F_T(t) &= P\{T\leqslant t\} = 1-P\{T>t\} = 1-P\{X_1>t, X_2>t, X_3>t, X_4>t\} \\
&= 1-\left[P\{X_1>t\}\right]^4 = 1-\left[1-F_X(t)\right]^4 = 1-\left[e^{-0.2t}\right]^4 = 1-e^{-0.8t},\quad t\geqslant0.
\end{aligned}
\]
当 $t<0$ 时，$F_T(t)=0$。故
\[
F_T(t)=
\begin{cases}
0, & t<0, \\
1-e^{-0.8t}, & t\geqslant0.
\end{cases}
\]
即 $T\sim E(0.8)$。
\end{solution}



\section{常见连续型随机变量及分布}

\begin{definition}[均匀分布]\label{def:uniform-dist}
若随机变量 $X$ 的概率密度为
\(
f(x) = \begin{cases}
\frac{1}{b-a}, & a < x < b, \\
0, & \text{其他},
\end{cases}
\)
则称 $X$ 在区间 $(a,b)$ 内服从均匀分布，记为 $X \sim U(a,b)$。其分布函数为
\[
F(x) = \begin{cases}
0, & x < a, \\
\frac{x-a}{b-a}, & a \le x < b, \\
1, & x \ge b.
\end{cases}
\]
\end{definition}

\begin{exercise}\label{exer:2-13}
设随机变量 $X$ 满足 $1 \le X \le 4$，且 $P\{X=1\} = \frac{1}{4}$，$P\{X=4\} = \frac{1}{4}$，当 $X$ 在区间 $(1,4)$ 内取值时，$X$ 服从均匀分布，求 $X$ 的分布函数 $F(x)$。
\end{exercise}
\begin{solution}
当 $x < 1$ 时，$F(x) = 0$；
当 $1 \le x < 4$ 时，
\[
\begin{aligned}
F(x) = P\{X \le 1\} + P\{1 < X \le x\} = \frac{1}{4} + \frac{x-1}{4-1} \cdot \frac{1}{2} = \frac{5 + 2x}{12};
\end{aligned}
\]
当 $x \ge 4$ 时，$F(x) = 1$。
故
\[
F(x) = \begin{cases}
0, & x < 1, \\
\frac{5 + 2x}{12}, & 1 \le x < 4, \\
1, & x \ge 4.
\end{cases}
\]
\end{solution}

\begin{exercise}\label{exer:2-16}
设 $X \sim U(a,b)$（$a>0$），且 $P\{X>4\} = \frac{1}{2}$，$P\{0<X<3\} = \frac{1}{4}$，求 $a,b$ 及 $P\{1<X<5\}$。
\end{exercise}
\begin{solution}
在区间 \((a,b)\) 中，落在 \(4\) 右边的长度占整个区间长度的一半。 \(a>0\)，所以事件 \(\{0<X<3\}\) 在支持区间上等价于 \(\{a<X<3\}\)。于是
\[\begin{cases}
P\{X>4\}=\frac{b-4}{b-a}=\frac12\implies
2(b-4)=b-a \quad \Longrightarrow \quad a+b=8.\\
P\{a<X<3\}=\frac{3-a}{b-a}=\frac14\implies
4(3-a)=b-a \quad \Longrightarrow \quad b=12-3a.
\end{cases}
\]
得 \(a=2,\ b=6\)。因为此时 \(a=2>1\)，所以事件 \(\{1<X<5\}\) 等价于 \(\{2<X<5\}\)。因此
\(
P\{1<X<5\}=P\{2<X<5\}=\frac{5-2}{6-2}=\frac34.
\)
\end{solution}

\begin{definition}[指数分布] \label{def:exp-dist}
若随机变量 $X$ 的概率密度为
\[
f(x) = \begin{cases}
\lambda e^{-\lambda x}, & x > 0, \\
0, & x \le 0,
\end{cases},~
F(x) =P(X\leqslant x)= \begin{cases}
1 - e^{-\lambda x}, & x \ge 0, \\
0, & x < 0.
\end{cases},~~ P(X>x)=e^{-\lambda x}
\]
其中 $\lambda > 0$，则称 $X$ 服从参数为 $\lambda$ 的指数分布，记为 $X \sim E(\lambda)$。其分布函数为
\end{definition}

\begin{exercise}[指数分布与系统更换]
设一批元件的寿命独立同分布，其概率密度函数为
\[
f(x)=
\begin{cases}
e^{-x}, & x\geqslant0, \\
0, & x<0,
\end{cases}
\]
系统开始工作时只有一个元件工作，在它损坏后立即更换一个新元件接替工作，求系统从开始工作到时刻 $T(T>0)$ 为止，恰更换了一个元件的概率。
\end{exercise}
\begin{solution}
设两个元件的寿命分别为 $X_1$ 和 $X_2$，则 $X_1, X_2 \stackrel{\text{i.i.d.}}{\sim} E(1)$，"恰更换了一个元件"等价于事件：
\(
\{X_1 < T,\ X_1+X_2 > T\}.
\)
所求概率：
\[
\begin{aligned}
P\{X_1 < T,\ X_1+X_2 > T\} &= \int_0^T f_{X_1}(x)\cdot P\{X_2 > T-x\}\mathrm dx = \int_0^T e^{-x}\cdot e^{-(T-x)}\mathrm dx \\
&= \int_0^T e^{-T}\mathrm dx = Te^{-T}.
\end{aligned}
\]
\end{solution}



\begin{exercise}\label{exer:src-paper3-i-04}\label{prob:p3-i-04}
设 \(X\) 与 \(Y\) 相互独立且都服从参数为 \(\lambda\) 的指数分布 \(E(\lambda)\)，则求\(X+Y\) 的分布函数。
\end{exercise}
\begin{solution}
令 \(Z=X+Y\)。用分布函数法求 \(Z\) 的分布函数。当 \(z<0\) 时，由于 \(X,Y\ge0\)，有 \(F_Z(z)=P(Z\le z)=0\)。
当 \(z\ge0\) 时，
\[
\begin{aligned}
F_Z(z)
&=P(X+Y\le z)=\int_0^z P(Y\le z-x)f_X(x)\,\mathrm{d}x=\int_0^z \left(1-e^{-\lambda(z-x)}\right)\lambda e^{-\lambda x}\,\mathrm{d}x\\
&=\int_0^z \lambda e^{-\lambda x}\,\mathrm{d}x-\int_0^z \lambda e^{-\lambda z}\,\mathrm{d}x=1-e^{-\lambda z}-\lambda z e^{-\lambda z}.
\end{aligned}
\]
因此
\[
F_Z(z)=
\begin{cases}
0, & z<0,\\
1-e^{-\lambda z}(1+\lambda z), & z\ge0.
\end{cases}
\]
\end{solution}

\begin{definition}[正态分布]\label{def:normal-dist}
若随机变量 $X$ 的概率密度为
\[
f(x) = \frac{1}{\sqrt{2\pi}\sigma} \exp\left\{-\frac{(x-\mu)^2}{2\sigma^2}\right\},\quad -\infty < x < +\infty,
\]
其中 $\mu$ 为常数，$\sigma > 0$ 为常数，则称 $X$ 服从参数为 $(\mu,\sigma^2)$ 的正态分布，记为 $X \sim N(\mu,\sigma^2)$。当 $\mu = 0,\ \sigma = 1$ 时，称 $X$ 服从标准正态分布，记为 $X \sim N(0,1)$。其概率密度和分布函数分别记为
\[
\varphi(x) = \frac{1}{\sqrt{2\pi}} e^{-\frac{x^2}{2}},\quad \Phi(x) = \int_{-\infty}^x \varphi(t)\,\mathrm{d}t.
\]
\end{definition}



\begin{property}[正态分布的常用性质]\label{prop:normal-properties}
设 \(X\sim N(\mu,\sigma^2)\)，其中 \(\sigma>0\)，则：
\begin{enumerate}
    \item \textbf{标准化：}
    \(
    \frac{X-\mu}{\sigma}\sim N(0,1),\quad
    F_X(x)=P\{X\le x\}=\Phi\left(\frac{x-\mu}{\sigma}\right).
    \)
    因此
    \(
    P\{a<X<b\}=\Phi\left(\frac{b-\mu}{\sigma}\right)-\Phi\left(\frac{a-\mu}{\sigma}\right).
    \)

    \item \textbf{对称性：}
    \(
    \Phi(0)=\frac12,\quad \Phi(-x)=1-\Phi(x),
    \)
    且
    \(
    P\{X\le\mu\}=P\{X>\mu\}=\frac12.
    \)
    特别地，当 \(a>0\) 时，
    \(
    P\{|X-\mu|<a\}=2\Phi\left(\frac{a}{\sigma}\right)-1.
    \)

    \item \textbf{单个正态变量的线性变换：} 若 \(a\ne0\)，则
    \(
    aX+b\sim N(a\mu+b,a^2\sigma^2).
    \)

    \item \textbf{独立正态变量的线性组合：} 若 \(X_1,\dots,X_n\) 相互独立，且 \(X_i\sim N(\mu_i,\sigma_i^2)\ (i=1,2,\dots,n)\)，则对任意常数 \(a_1,\dots,a_n\)，有
    \[
    \sum_{i=1}^n a_iX_i
    \sim
    N\left(\sum_{i=1}^n a_i\mu_i,\ \sum_{i=1}^n a_i^2\sigma_i^2\right).
    \]
    特别地，若 \(X,Y\) 相互独立，\(X\sim N(\mu_1,\sigma_1^2)\)，\(Y\sim N(\mu_2,\sigma_2^2)\)，则
    \(
    aX+bY+c\sim N(a\mu_1+b\mu_2+c,\ a^2\sigma_1^2+b^2\sigma_2^2).
    \)
\end{enumerate}

注意：仅知道 \(X,Y\) 各自服从正态分布，不能推出 \(aX+bY\) 一定服从正态分布；上面的线性组合结论需要相互独立，或更一般地需要联合正态。
\end{property}

\begin{exercise}\label{exer:src-2015-spring-a-i-03}
设随机变量 \(X,Y\) 均服从正态分布 \(N(2,0.2)\)，且相互独立。求
\(
Z=X-2Y+1
\)
的概率密度函数。
\end{exercise}
\begin{solution}
独立正态随机变量的线性组合仍服从正态分布。由
\(
EZ=EX-2EY+1=2-4+1=-1,
\)
且
\(
DZ=DX+4DY=0.2+4\times0.2=1,
\)
得
\(
Z\sim N(-1,1).
\)
因此 \(Z\) 的概率密度函数为
\[
f_Z(z)=\frac1{\sqrt{2\pi}}\exp\!\left\{-\frac{(z+1)^2}{2}\right\},
\qquad -\infty<z<+\infty.
\]
\end{solution}


\begin{exercise}\label{exer:2-12}\label{exer:src-2017-a-ii-05}
设随机变量 $X$ 服从正态分布 $N(0,1)$，对给定的 $\alpha(0<\alpha<1)$，数 $u_\alpha$ 满足 $P\{X>u_\alpha\} = \alpha$，若 $P\{|X|<x\} = \alpha$，则 $x$ 等于（\quad）。
(A) $u_{\frac{\alpha}{2}}$ \quad (B) $u_{1-\frac{\alpha}{2}}$ \quad (C) $u_{\frac{1-\alpha}{2}}$ \quad (D) $u_{1-\alpha}$
\end{exercise}

\begin{solution}
题目给的是中间概率
\(
P\{|X|<x\}=P\{-x<X<x\}=\alpha.
\)
由于正态分布关于 0 对称，所以区间外剩余的总概率为 \(1-\alpha\)，左右两边各占 \(\dfrac{1-\alpha}{2}\)。
而 \(u_\beta\) 的定义是满足 \(P\{X>u_\beta\}=\beta\) 的点，与上式对照可知 \(x=u_{\frac{1-\alpha}{2}}\)。
故答案选 \textbf{(C)}。
\end{solution}


\begin{exercise}[正态分布标准差比较]\label{exer:src-2015-spring-a-ii-05}
设 $X\sim N(\mu_1,\sigma_1^2), Y\sim N(\mu_2,\sigma_2^2)$，且 $P(|X-\mu_1|<1)>P(|Y-\mu_2|<1)$，则（\quad）。
(A)$\sigma_1<\sigma_2$\quad (B)$\sigma_1>\sigma_2$\quad (C) $\mu_1<\mu_2$\quad (D) $\mu_1>\mu_2$。
\end{exercise}
\begin{solution}
答案选\textbf{(A)}。标准化得：
\[
P\left\{\frac{|X-\mu_1|}{\sigma_1}<\frac{1}{\sigma_1}\right\}>P\left\{\frac{|Y-\mu_2|}{\sigma_2}<\frac{1}{\sigma_2}\right\}
\Rightarrow 2\Phi\left(\frac{1}{\sigma_1}\right)-1 > 2\Phi\left(\frac{1}{\sigma_2}\right)-1
\]
\[
\Rightarrow \Phi\left(\frac{1}{\sigma_1}\right)>\Phi\left(\frac{1}{\sigma_2}\right)\implies
\frac{1}{\sigma_1}>\frac{1}{\sigma_2} \Rightarrow \sigma_1<\sigma_2.
\]
\end{solution}


\section{常见分布的可加性，期望与方差}
\begin{property}[常见分布的可加性]\label{prop:additivity}
有些相互独立且服从同类型分布的随机变量，其和的分布也是同类型的。设 $X$ 与 $Y$ 相互独立，则
\begin{enumerate}
    \item 二项分布：若 $X \sim B(n,p)$，$Y \sim B(m,p)$，则 $X + Y \sim B(n+m,p)$（注意 $p$ 必须相同）。
    \item 泊松分布：若 $X \sim P(\lambda_1)$，$Y \sim P(\lambda_2)$，则 $X + Y \sim P(\lambda_1 + \lambda_2)$。
    \item 正态分布：若 $X \sim N(\mu_1, \sigma_1^2)$，$Y \sim N(\mu_2, \sigma_2^2)$，则 $X + Y \sim N(\mu_1 + \mu_2, \sigma_1^2 + \sigma_2^2)$。更一般地，独立正态变量的线性组合
    \[
    \sum_{i=1}^{n}c_i X_i \sim N\!\left(\sum_{i=1}^{n}c_i\mu_i,\ \sum_{i=1}^{n}c_i^2\sigma_i^2\right).
    \]
    \item $\chi^2$ 分布：若 $X \sim \chi^2(n_1)$，$Y \sim \chi^2(n_2)$，则 $X + Y \sim \chi^2(n_1 + n_2)$（$\chi^2$ 分布详见第六章）。
    \item Gamma分布：若 $X \sim \Gamma(\alpha_1, \beta)$，$Y \sim \Gamma(\alpha_2, \beta)$，则 $X + Y \sim \Gamma(\alpha_1 + \alpha_2, \beta)$。特别地，$X_i \sim E(\lambda)$ 独立时，$\sum_{i=1}^{n}X_i \sim \Gamma(n, \lambda)$。
\end{enumerate}
以上结果对 $n$ 个相互独立的随机变量也成立。
\end{property}
\begin{table}[h]
\centering
\caption{常见分布的数学期望与方差} \label{tab:dist-expect-var}
\begin{tabular}{c|c|c|c}
\toprule
名称与记号 & 分布列或概率密度 & 数学期望 & 方差 \\
\midrule
(0-1) 分布 $B(1,p)$ & $P\{X=1\}=p,\ P\{X=0\}=1-p$ & $p$ & $p(1-p)$ \\
二项分布 $B(n,p)$ & $P\{X=k\}=\mathrm{C}_n^k p^k (1-p)^{n-k}$ & $np$ & $np(1-p)$ \\
泊松分布 $P(\lambda)$ & $P\{X=k\}=e^{-\lambda}\frac{\lambda^k}{k!}$ & $\lambda$ & $\lambda$ \\
几何分布 $G(p)$ & $P\{X=k\}=(1-p)^{k-1}p$ & $\frac{1}{p}$ & $\frac{1-p}{p^2}$ \\
超几何分布 $H(n,N,M)$ & $P\{X=k\}=\frac{\mathrm{C}_M^k \mathrm{C}_{N-M}^{n-k}}{\mathrm{C}_N^n}$ & $\frac{nM}{N}$ & $\frac{nM}{N}\left(1-\frac{M}{N}\right)\frac{N-n}{N-1}$ \\
\midrule
均匀分布 $U(a,b)$ & $f(x)=\frac{1}{b-a},\ a < x < b$ & $\frac{a+b}{2}$ & $\frac{(b-a)^2}{12}$ \\
指数分布 $E(\lambda)$ & $f(x)=\lambda e^{-\lambda x},\ x > 0$ & $\frac{1}{\lambda}$ & $\frac{1}{\lambda^2}$ \\
正态分布 $N(\mu,\sigma^2)$ & $f(x)=\frac{1}{\sqrt{2\pi}\sigma}e^{-\frac{(x-\mu)^2}{2\sigma^2}}$ & $\mu$ & $\sigma^2$ \\
\bottomrule
\end{tabular}
\end{table}


\begin{proof}[二项分布 $B(n,p)$ 的期望与方差推导]
利用组合数性质 $k \mathrm{C}_n^k = n \mathrm{C}_{n-1}^{k-1}$，计算期望：
\[
\begin{aligned}
E(X) &= \sum_{k=0}^n k \mathrm{C}_n^k p^k (1-p)^{n-k} = \sum_{k=1}^n n \mathrm{C}_{n-1}^{k-1} p^k (1-p)^{n-k} = np \sum_{k=1}^n \mathrm{C}_{n-1}^{k-1} p^{k-1} (1-p)^{(n-1)-(k-1)}=np
\end{aligned}
\]
计算 $E[X(X-1)]$，利用性质 $k(k-1)\mathrm{C}_n^k = n(n-1)\mathrm{C}_{n-2}^{k-2}$：
\[
\begin{aligned}
E[X(X-1)] &= \sum_{k=0}^n k(k-1) \mathrm{C}_n^k p^k (1-p)^{n-k} = \sum_{k=2}^n n(n-1) \mathrm{C}_{n-2}^{k-2} p^k (1-p)^{n-k} \\
&= n(n-1)p^2 \sum_{j=0}^{n-2} \mathrm{C}_{n-2}^j p^j (1-p)^{(n-2)-j} = n(n-1)p^2.
\end{aligned}
\]
由于 $D(X) = E(X^2) - [E(X)]^2 = E[X(X-1)] + E(X) - [E(X)]^2$，代入得：
\(
D(X) = n(n-1)p^2 + np - n^2p^2 = -np^2 + np = np(1-p).
\)
\end{proof}

\begin{proof}[泊松分布 $P(\lambda)$ 的期望与方差推导]
根据泰勒展开式 $\sum_{k=0}^\infty \frac{\lambda^k}{k!} = e^\lambda$。计算期望：
\[
E(X) = \sum_{k=0}^\infty k e^{-\lambda} \frac{\lambda^k}{k!} = \lambda e^{-\lambda} \sum_{k=1}^\infty \frac{\lambda^{k-1}}{(k-1)!} = \lambda e^{-\lambda} e^\lambda = \lambda.
\]
为求方差，先计算 $E[X(X-1)]$：
\[
E[X(X-1)] = \sum_{k=0}^\infty k(k-1) e^{-\lambda} \frac{\lambda^k}{k!} = \lambda^2 e^{-\lambda} \sum_{k=2}^\infty \frac{\lambda^{k-2}}{(k-2)!} = \lambda^2 e^{-\lambda} e^\lambda = \lambda^2.
\]
故方差为：
\[
D(X) = E[X(X-1)] + E(X) - [E(X)]^2 = \lambda^2 + \lambda - \lambda^2 = \lambda.
\]
\end{proof}

\begin{proof}[几何分布 $G(p)$ 的期望与方差推导]
利用幂级数求和
\(
\sum_{k=1}^\infty (1-p)^k=\frac{1-p}{p}.
\)
两边对 \(p\) 求导，得
\(
-\sum_{k=1}^\infty k(1-p)^{k-1}=-\frac{1}{p^2},
\)
即
\(
\sum_{k=1}^\infty k(1-p)^{k-1}=\frac{1}{p^2}.
\)
计算期望：
\(
E(X) = \sum_{k=1}^\infty k(1-p)^{k-1} p = p \cdot \frac{1}{p^2} = \frac{1}{p}.
\)
对
\(
\sum_{k=1}^\infty k(1-p)^{k-1}=\frac{1}{p^2}
\)
再对 \(p\) 求导，得
\(
\sum_{k=2}^\infty k(k-1)(1-p)^{k-2}=\frac{2}{p^3}.
\)
计算 \(E[X(X-1)]\)：
\[
E[X(X-1)] = \sum_{k=1}^\infty k(k-1)(1-p)^{k-1}p
=p(1-p)\sum_{k=2}^\infty k(k-1)(1-p)^{k-2}
=p(1-p)\frac{2}{p^3}
=\frac{2(1-p)}{p^2}.
\]
故方差为：
\(
D(X) = E[X(X-1)] + E(X) - [E(X)]^2 = \frac{2(1-p)}{p^2} + \frac{1}{p} - \frac{1}{p^2} = \frac{1-p}{p^2}.
\)
\end{proof}

\begin{proof}[超几何分布 $H(n,N,M)$ 的期望与方差推导]
采用引入示性变量（Indicator variables）的求法最为简便。设 $X_i$ 为第 $i$ 次抽出的物品是指定物品（次品/正品）的示性变量，若抽出则 $X_i=1$，否则 $X_i=0$。
易知 $P(X_i=1) = \frac{M}{N}$，故 $X_i \sim B(1, \frac{M}{N})$，则 $X = \sum_{i=1}^n X_i$。
期望为：
\(
E(X) = E\left(\sum_{i=1}^n X_i\right) = \sum_{i=1}^n E(X_i) = n \frac{M}{N}.
\)
对于方差，需考虑 $X_i$ 之间的协方差。当 $i \neq j$ 时：
\(
E(X_i X_j) = P(X_i=1, X_j=1) = \frac{M}{N} \cdot \frac{M-1}{N-1}.
\)
\(
\mathrm{Cov}(X_i, X_j) = E(X_i X_j) - E(X_i)E(X_j) = \frac{M(M-1)}{N(N-1)} - \frac{M^2}{N^2} = -\frac{M(N-M)}{N^2(N-1)}.
\)
由于 $D(X_i) = \frac{M}{N}\left(1-\frac{M}{N}\right)$，故总方差为：
\[
\begin{aligned}
D(X) &= \sum_{i=1}^n D(X_i) + \sum_{i \neq j} \mathrm{Cov}(X_i, X_j) \\
&= n \frac{M}{N}\left(1-\frac{M}{N}\right) + n(n-1)\left[-\frac{M(N-M)}{N^2(N-1)}\right] \\
&= n \frac{M}{N}\left(1-\frac{M}{N}\right) \left[1 - \frac{n-1}{N-1}\right] = \frac{nM}{N}\left(1-\frac{M}{N}\right)\frac{N-n}{N-1}.
\end{aligned}
\]
\end{proof}

\begin{proof}[均匀分布 $U(a,b)$ 的期望与方差推导]
根据定义，计算期望：
\(
E(X) = \int_{a}^{b} x \frac{1}{b-a} \,\mathrm{d}x = \frac{1}{b-a} \left[ \frac{1}{2}x^2 \right]_{a}^{b} = \frac{b^2-a^2}{2(b-a)} = \frac{a+b}{2}.
\)
计算二阶矩：
\[
E(X^2) = \int_{a}^{b} x^2 \frac{1}{b-a} \,\mathrm{d}x = \frac{1}{b-a} \left[ \frac{1}{3}x^3 \right]_{a}^{b} = \frac{b^3-a^3}{3(b-a)} = \frac{b^2+ab+a^2}{3}.
\]
故方差为：
\[
\begin{aligned}
D(X) = E(X^2) - [E(X)]^2 &= \frac{b^2+ab+a^2}{3} - \frac{a^2+2ab+b^2}{4} = \frac{(b-a)^2}{12}.
\end{aligned}
\]
\end{proof}

\begin{proof}[指数分布 $E(\lambda)$ 的期望与方差推导]
利用分部积分法，计算期望：
\[
E(X) = \int_{0}^{+\infty} x \lambda e^{-\lambda x} \,\mathrm{d}x = \left[ -x e^{-\lambda x} \right]_{0}^{+\infty} + \int_{0}^{+\infty} e^{-\lambda x} \,\mathrm{d}x = 0 + \left[ -\frac{1}{\lambda} e^{-\lambda x} \right]_{0}^{+\infty} = \frac{1}{\lambda}.
\]
计算二阶矩（再次利用分部积分）：
\[
\begin{aligned}
E(X^2) &= \int_{0}^{+\infty} x^2 \lambda e^{-\lambda x} \,\mathrm{d}x = \left[ -x^2 e^{-\lambda x} \right]_{0}^{+\infty} + \int_{0}^{+\infty} 2x e^{-\lambda x} \,\mathrm{d}x \\
&= 0 + \frac{2}{\lambda} \int_{0}^{+\infty} x \lambda e^{-\lambda x} \,\mathrm{d}x = \frac{2}{\lambda} E(X) = \frac{2}{\lambda^2}.
\end{aligned}
\]
故方差为：
\(
D(X) = E(X^2) - [E(X)]^2 = \frac{2}{\lambda^2} - \frac{1}{\lambda^2} = \frac{1}{\lambda^2}.
\)
\end{proof}


\begin{exercise}[拉普拉斯分布变换] 
设随机变量 $X$ 的概率密度为
\[
f(x)=\frac{1}{2\lambda}e^{-\frac{|x|}{\lambda}},\quad -\infty<x<+\infty,\ \lambda>0.
\]
\begin{enumerate}
    \item[(1)] 求方程 $y^2+2Xy+\lambda^2=0$ 有实根的概率；
    \item[(2)] 求 $Z=e^{-|X|}$ 的概率密度。
\end{enumerate}
\end{exercise}
\begin{solution}
(1) 方程有实根的条件是判别式
\[
\Delta = (2X)^2 - 4\lambda^2 = 4(X^2-\lambda^2) \geqslant 0 \Leftrightarrow |X|\geqslant\lambda.
\]
所求概率：
\[
\begin{aligned}
P\{|X|\geqslant\lambda\} &= \int_{-\infty}^{-\lambda} f(x)\,\mathrm{d}x + \int_{\lambda}^{\infty} f(x)\,\mathrm{d}x = 2\int_{\lambda}^{\infty} \frac{1}{2\lambda}e^{-\frac{x}{\lambda}}\,\mathrm{d}x = \frac{1}{\lambda}\int_{\lambda}^{\infty} e^{-\frac{x}{\lambda}}\,\mathrm{d}x = e^{-1}.
\end{aligned}
\]
(2) $Z=e^{-|X|}$ 的取值范围是 $(0,1]$。当 $0<z<1$ 时：
\[
\begin{aligned}
F_Z(z) &= P\{Z\leqslant z\} = P\{e^{-|X|}\leqslant z\} = P\{|X|\geqslant -\ln z\} = 2\int_{-\ln z}^{\infty} \frac{1}{2\lambda}e^{-\frac{x}{\lambda}}\,dx = z^{\frac{1}{\lambda}}.
\end{aligned}
\]
当 $z\leqslant0$ 时，$F_Z(z)=0$；当 $z\geqslant1$ 时，$F_Z(z)=1$。
求导得密度函数：
\[
f_Z(z) =
\begin{cases}
\dfrac{1}{\lambda}z^{\frac{1}{\lambda}-1}, & 0<z<1, \\
0, & \text{其他}.
\end{cases}
\]
\end{solution}



\begin{exercise}[分段概率密度的分布函数] \label{exer:2-47}
随机变量 $X$ 的概率密度为
\[
f(x)=
\begin{cases}
x, & 0\leqslant x<1, \\
2-x, & 1\leqslant x<2, \\
0, & \text{其他}.
\end{cases}
\]
求 $X$ 的分布函数 $F(x)$。
\end{exercise}
\begin{solution}
分段计算分布函数 $F(x) = \int_{-\infty}^{x} f(t)\,\mathrm{d}t$：
\begin{enumerate}
    \item 当 $x<0$ 时：$F(x)=0$；
    \item 当 $0\leqslant x<1$ 时：
    \(
    F(x) = \int_0^x t\,\mathrm{d}t = \frac{x^2}{2};
    \)
    \item 当 $1\leqslant x<2$ 时：
    \(
    F(x) = \int_0^1 t\,\mathrm{d}t + \int_1^x (2-t)\,\mathrm{d}t = \frac{1}{2} + \left[2t - \frac{t^2}{2}\right]_1^x = -\frac{x^2}{2} + 2x - 1;
    \)    
    \item 当 $x\geqslant2$ 时：$F(x)=1$。
\end{enumerate}
\end{solution}


\section{一维随机变量函数的分布} \label{sec:random-variable-function}

\subsection{离散型随机变量函数的分布} \label{subsec:discrete-to-discrete}

设 $X$ 为离散型随机变量，其概率分布为
\(
P\{X = x_i\} = p_i,\quad i=1,2,\dots
\)
则 $X$ 的函数 $Y = g(X)$ 也是离散型随机变量。若 $g(x_i)$ 的值互不相同，则 $Y$ 的分布律为
\(
P\{Y = g(x_i)\} = p_i,\quad i=1,2,\dots
\)
即
\[
Y \sim \begin{pmatrix} g(x_1) & g(x_2) & \cdots \\ p_1 & p_2 & \cdots \end{pmatrix}.
\]

如果有若干个 $g(x_i)$ 的值相同，则需要将这些相同值对应的概率相加合并。具体来说，若 $g(x_{i_1}) = g(x_{i_2}) = \cdots = g(x_{i_k}) = y_j$，则
\(
P\{Y = y_j\} = \sum_{t=1}^{k} P\{X = x_{i_t}\} = \sum_{t=1}^{k} p_{i_t}.
\)

\begin{property}\label{prop:discrete-function}
离散型随机变量函数的分布求解步骤：
\begin{enumerate}
    \item 确定 $Y = g(X)$ 的所有可能取值 $y_j$；
    \item 对每个 $y_j$，找出所有满足 $g(x_i) = y_j$ 的 $x_i$；
    \item 将相应的概率 $p_i$ 相加，得到 $P\{Y = y_j\}$。
\end{enumerate}
\end{property}

\subsection{连续型随机变量函数的分布} \label{subsec:continuous-function}

设 $X$ 为连续型随机变量，其概率密度函数为 $f_X(x)$，分布函数为 $F_X(x)$。对于连续型随机变量 $X$ 的函数 $Y = g(X)$，需要根据 $Y$ 的类型（离散型或连续型）采用不同的方法。

\subsubsection{当 \texorpdfstring{$Y$}{Y} 为离散型时}

若 $Y = g(X)$ 只能取有限个或可列个值，则 $Y$ 是离散型随机变量。此时应先确定 $Y$ 的所有可能取值，然后通过 $X$ 的概率分布计算 $Y$ 取各个值的概率。具体来说，若 $Y$ 的可能取值为 $y_1, y_2, \dots, y_n, \dots$，则
\[
P\{Y = y_j\} = P\{g(X) = y_j\} = \int_{\{x: g(x) = y_j\}} f_X(x) \,\mathrm{d}x.
\]
这通常涉及到对 $X$ 的取值空间按 $g(x)$ 的取值进行划分。

\subsubsection{当 \texorpdfstring{$Y$}{Y} 为连续型时} 

若 $Y = g(X)$ 的取值充满一个区间（或若干区间的并），则 $Y$ 是连续型随机变量。
根据分布函数的定义，先求 $Y$ 的分布函数 $F_Y(y)$，再对 $y$ 求导得到概率密度 $f_Y(y)$。
具体步骤如下：
\begin{enumerate}
    \item 对任意实数 $y$，求 $Y$ 的分布函数：
    \[
    F_Y(y) = P\{Y \le y\} = P\{g(X) \le y\} = \int_{\{x: g(x) \le y\}} f_X(x) \,\mathrm{d}x.
    \]
    这一步的关键是解不等式 $g(x) \le y$，确定积分区域 $\{x: g(x) \le y\}$。
    
    \item 对 $F_Y(y)$ 关于 $y$ 求导（在可导点处）：
    \(
    f_Y(y) = F_Y'(y).
    \)
    
    \item 检查 $f_Y(y)$ 是否满足概率密度函数的基本条件：非负性和归一性。
\end{enumerate}


\begin{exercise}[线性变换与平方变换] \label{exer:2-18}
设连续型随机变量 $X$ 的概率密度函数为
\(
f_X(x) = \begin{cases}
1, & -\dfrac{1}{2} < x < \dfrac{1}{2}, \\
0, & \text{其他},
\end{cases}
\)
求下列随机变量 $Y$ 的概率密度函数：
\begin{enumerate}
    \item[(1)] $Y = a + bX$, 其中 $b \neq 0$；
    \item[(2)] $Y = X^2$；
    \item[(3)] $Y = |X|$。
\end{enumerate}
\end{exercise}
\begin{solution}
(1) \textbf{线性变换}：使用分布函数法求 $Y=a+bX$ 的概率密度。设 $Y$ 的分布函数为 $F_Y(y)$，由于 $b \neq 0$，需分情况讨论：
\begin{itemize}
\item 当 $b > 0$ 时：
\begin{align*}
F_Y(y) &= P\{Y \le y\} = P\{a + bX \le y\} = P\left\{X \le \frac{y-a}{b}\right\} = F_X\left(\frac{y-a}{b}\right). \\
\intertext{对 $y$ 求导可得：}
f_Y(y) &= F_Y'(y) = f_X\left(\frac{y-a}{b}\right) \cdot \frac{1}{b},~Y=a+bx\in\left(a-\dfrac{b}{2},\ a+\dfrac{b}{2}\right).
\end{align*}
\item 当 $b < 0$ 时：
\begin{align*}
F_Y(y) &= P\{Y \le y\} = P\{a + bX \le y\} = P\left\{X \ge \frac{y-a}{b}\right\} = 1 - F_X\left(\frac{y-a}{b}\right). \\
f_Y(y) &= F_Y'(y) = -f_X\left(\frac{y-a}{b}\right) \cdot \frac{1}{b} = f_X\left(\frac{y-a}{b}\right) \cdot \frac{1}{-b},~~Y\in\left(a+\dfrac{b}{2}, a-\dfrac{b}{2}\right)
\end{align*}
\end{itemize}
综合以上两种情况，无论是 $b>0$（此时 $b=|b|$）还是 $b<0$（此时 $-b=|b|$），都可以统一写为：
\(
f_Y(y) = f_X\left(\frac{y-a}{b}\right) \cdot \frac{1}{|b|}.
\)
取值区间统一为 $\left(a-\dfrac{|b|}{2}, a+\dfrac{|b|}{2}\right)$。在该区间内，$f_X(\cdot) = 1$。因此 $Y$ 的概率密度函数为：
\[
f_Y(y) = \begin{cases}
\dfrac{1}{|b|}, & a - \dfrac{|b|}{2} < y < a + \dfrac{|b|}{2}, \\[8pt]
0, & \text{其他}.
\end{cases}
\]
(2) \textbf{平方变换}：函数 $y = x^2$ 在 $(-\dfrac{1}{2}, \dfrac{1}{2})$ 上不是单调函数，必须使用分布函数法。由于 $Y = X^2 \ge 0$，当 $y \le 0$ 时，$F_Y(y) = 0$。当 $0 < y < \dfrac{1}{4}$ 时：
\[
F_Y(y) = P\{X^2 \le y\} = P\{-\sqrt{y} \le X \le \sqrt{y}\} = \int_{-\sqrt{y}}^{\sqrt{y}} 1 \, \mathrm{d}x = 2\sqrt{y}.
\]
当 $y \ge \dfrac{1}{4}$ 时，$F_Y(y) = 1$。
求导得概率密度：
\[
f_Y(y) = \begin{cases}
\dfrac{1}{\sqrt{y}}, & 0 < y < \dfrac{1}{4}, \\
0, & \text{其他}.
\end{cases}
\]
(3) \textbf{绝对值变换}：同样使用分布函数法。由于 $Y = |X| \ge 0$，当 $y < 0$ 时，$F_Y(y) = 0$。当 $0 \le y < \dfrac{1}{2}$ 时：
\[
F_Y(y) = P\{|X| \le y\} = P\{-y \le X \le y\} = \int_{-y}^{y} 1 \, \mathrm{d}x = 2y.
\]
当 $y \ge \dfrac{1}{2}$ 时，$F_Y(y) = 1$。求导得：
\[
f_Y(y) = \begin{cases}
2, & 0 \le y < \dfrac{1}{2}, \\
0, & \text{其他}.
\end{cases}
\]
\end{solution}


\begin{exercise}[概率积分变换] \label{exer:2-20}
设随机变量 $X$ 的分布函数 $F_X(x)$ 是严格单调增加函数，其反函数 $F_X^{-1}(y)$ 存在，$Y = F_X(X)$。证明：$Y$ 服从区间 $(0,1)$ 上的均匀分布。
\end{exercise}

\begin{solution}
要证明 $Y$ 服从均匀分布，最直接的方法就是求它的分布函数 $F_Y(y)=P\{Y\le y\}$。

由于 $F_X(x)\in[0,1]$，所以 $Y=F_X(X)$ 的取值只可能落在 $[0,1]$ 内。下面分段讨论。

当 $y < 0$ 时，$F_Y(y) = P\{Y \le y\} = 0$。

当 $0 \le y < 1$ 时，因为 $F_X$ 严格单调增加，
\[
F_Y(y)=P\{Y\le y\}=P\{F_X(X)\le y\}=P\{X\le F_X^{-1}(y)\}=F_X\!\bigl(F_X^{-1}(y)\bigr)=y.
\]

当 $y \ge 1$ 时，$F_Y(y) = 1$。

综上，$Y$ 的分布函数为
\[
F_Y(y) = \begin{cases}
0, & y < 0, \\
y, & 0 \le y < 1, \\
1, & y \ge 1,
\end{cases}
\]
这正是均匀分布 $U(0,1)$ 的分布函数，因此 \(Y\sim U(0,1)\)。换句话说，\textbf{任何连续型随机变量经过自己的分布函数变换后，都会变成标准均匀分布}。
\end{solution}


\begin{exercise}[分段随机变量的函数] \label{exer:2-21}
在区间 $(0,2)$ 上随机取一点，将该区间分成两段，较短一段的长度记为 $X$，较长一段的长度记为 $Y$。令 $Z = \dfrac{Y}{X}$。(1) 求 $X$ 的概率密度；(2) 求 $Z$ 的概率密度。
\end{exercise}
\begin{solution}
(1) 设随机点的坐标为 $V \sim U(0,2)$，由于区间总长为 $2$，较短一段的长度 $X = \min\{V, 2-V\}$。显然 $X$ 的取值范围为 $(0,1)$。
先求 $X$ 的分布函数 $F_X(x) = P\{X \le x\}$：
当 $x \le 0$ 时，$F_X(x) = 0$；
当 $x \ge 1$ 时，$F_X(x) = 1$。
当 $0 < x < 1$ 时，利用对立事件有：
\[
\begin{aligned}
F_X(x) &= P\{\min\{V, 2-V\} \le x\} = 1 - P\{\min\{V, 2-V\} > x\} \\
&= 1 - P\{V > x \text{ 且 } 2-V > x\} = 1 - P\{x < V < 2-x\}=1-\int_{x}^{2-x} \frac{1}{2} \,\mathrm{d}v=x
\end{aligned}
\]
对分布函数求导，得到 $X$ 的概率密度为：
\[
f_X(x) = F'_X(x) = \begin{cases}
1, & 0 < x < 1, \\
0, & \text{其他}.
\end{cases}
\]
（即 $X \sim U(0,1)$）\\
(2) 由于 $Y = 2 - X$，故 $Z = \dfrac{Y}{X} = \dfrac{2-X}{X} = \dfrac{2}{X} - 1$。
因为 $X \in (0,1)$，所以 $Z$ 的取值范围为 $(1, +\infty)$。
现在用分布函数法求 $Z$ 的概率密度。设 $Z$ 的分布函数为 $F_Z(z) = P\{Z \le z\}$：
当 $z \le 1$ 时，$F_Z(z) = 0$。
当 $z > 1$ 时：
\[
\begin{aligned}
F_Z(z) &= P\left\{\frac{2}{X} - 1 \le z\right\} = P\left\{\frac{2}{X} \le z + 1\right\} = P\left\{X \ge \frac{2}{z+1}\right\}
\end{aligned}
\]
由于 $X \sim U(0,1)$，当 $z > 1$ 时，$\dfrac{2}{z+1} \in (0, 1)$，因此：
\[
F_Z(z) = \int_{\frac{2}{z+1}}^{1} f_X(x) \,\mathrm{d}x = \int_{\frac{2}{z+1}}^{1} 1 \,\mathrm{d}x = 1 - \frac{2}{z+1} = \frac{z-1}{z+1}
\]
对分布函数求导，得到 $Z$ 的概率密度为：
\(
f_Z(z) = F'_Z(z) = \left( 1 - \frac{2}{z+1} \right)' = \frac{2}{(z+1)^2}, \quad z > 1.
\)
故 $Z$ 的概率密度为：
\[
f_Z(z) = \begin{cases}
\dfrac{2}{(1+z)^2}, & z > 1, \\
0, & \text{其他}.
\end{cases}
\]
\end{solution}

\begin{exercise}[分段密度函数的变换] \label{exer:2-22}
设随机变量 $X$ 的概率密度为
\[
f_X(x) = \begin{cases}
\dfrac{1}{2}, & -1 < x < 0, \\
\dfrac{1}{4}, & 0 \le x < 2, \\
0, & \text{其他},
\end{cases}
\]
令 $Y = X^2$，求 $Y$ 的概率密度 $f_Y(y)$。
\end{exercise}

\begin{solution}
使用分布函数法。由于 $Y = X^2 \ge 0$，当 $y < 0$ 时，$F_Y(y) = 0$。
当 $0 \le y < 1$ 时：
\[
\begin{aligned}
F_Y(y) &= P\{-\sqrt{y} \le X \le \sqrt{y}\} = P\{-\sqrt{y} \le X < 0\} + P\{0 \le X \le \sqrt{y}\} = \frac{1}{2}\sqrt{y} + \frac{1}{4}\sqrt{y} = \frac{3}{4}\sqrt{y}.
\end{aligned}
\]

当 $1 \le y < 4$ 时（注意 $-1 < x < 0$ 对应 $0 < y < 1$，而 $0 \le x < 2$ 对应 $0 \le y < 4$）：
\(
F_Y(y) = P\{-1 < X < 0\} + P\{0 \le X \le \sqrt{y}\} = \frac{1}{2} + \frac{1}{4}\sqrt{y}.
\)

当 $y \ge 4$ 时，$F_Y(y) = 1$。

求导得概率密度：
\[
f_Y(y) = \begin{cases}
\dfrac{3}{8\sqrt{y}}, & 0 < y < 1, \\
\dfrac{1}{8\sqrt{y}}, & 1 \le y < 4, \\
0, & \text{其他}.
\end{cases}
\]
\end{solution}


\begin{exercise}\label{exer:2-31}
设随机变量 $X \sim N(4,16)$，求 $Y=|X-4|$ 的概率密度 $f_Y(y)$。
\end{exercise}

\begin{solution}
令 $Z = X-4$，则 $Z \sim N(0,16)$，概率密度为 $f_Z(z) = \dfrac{1}{\sqrt{32\pi}} e^{-\frac{z^2}{32}}$。

函数 $y = |z|$ 在 $(-\infty,+\infty)$ 上不是严格单调函数，使用分布函数法。

由于 $Y = |Z| \ge 0$，当 $y \le 0$ 时，$F_Y(y) = 0$。

当 $y > 0$ 时：
\begin{align*}
F_Y(y) &= P\{|Z| \le y\} = P\{-y \le Z \le y\} = F_Z(y) - F_Z(-y).
\end{align*}
求导得：
\[
f_Y(y) = F_Z'(y) - F_Z'(-y) \cdot (-1) = f_Z(y) + f_Z(-y) = \frac{2}{\sqrt{32\pi}} e^{-\frac{y^2}{32}}.
\]
故
\[
f_Y(y) = \begin{cases}
\dfrac{1}{\sqrt{8\pi}} e^{-\frac{y^2}{32}}, & y > 0, \\
0, & y \le 0.
\end{cases}
\]
\end{solution}


\begin{exercise}[正态分布平方变换] 
设随机变量 $X\sim N(0,1)$，求 $Y=2X^2+1$ 的概率密度函数。
\end{exercise}
\begin{solution}
当 $y\leqslant1$ 时，$F_Y(y)=0$，故 $f_Y(y)=0$。\\
当 $y>1$ 时：
\[
\begin{aligned}
F_Y(y) &= P\{2X^2+1\leqslant y\} = P\left\{-\sqrt{\frac{y-1}{2}}\leqslant X\leqslant\sqrt{\frac{y-1}{2}}\right\} = 2\Phi\left(\sqrt{\frac{y-1}{2}}\right)-1,
\end{aligned}
\]
其中 $\Phi$ 为标准正态分布函数。求导得：
\[
f_Y(y) = \frac{\mathrm{d}}{\mathrm{d}y}F_Y(y) = 2\varphi\left(\sqrt{\frac{y-1}{2}}\right)\cdot\frac{1}{2\sqrt{2(y-1)}},
\]
其中 $\varphi(x)=\dfrac{1}{\sqrt{2\pi}}e^{-x^2/2}$ 为标准正态密度。化简得：
\[
f_Y(y) =
\begin{cases}
\dfrac{1}{2\sqrt{\pi(y-1)}}\exp\left\{-\dfrac{y-1}{4}\right\}, & y>1, \\[8pt]
0, & y\leqslant1.
\end{cases}
\]
\end{solution}

\begin{exercise}[均匀分布绝对值变换] \label{exer:2-51}
设 $X\sim U(-1,2)$，求 $Y=|X|$ 的概率密度。
\end{exercise}
\begin{solution}
函数 $y=|x|$ 在区间 $(-1,2)$ 上不是单调函数，使用分布函数法。由于 $Y=|X|\in[0,2)$：
\begin{itemize}
    \item 当 $0<y<1$ 时：
    \[
    F_Y(y) = P\{|X|\leqslant y\} = P\{-y\leqslant X\leqslant y\} = \frac{y-(-y)}{3} = \frac{2y}{3},~f_Y(y)=\dfrac{2}{3}
    \]
    \item 当 $1\leqslant y<2$ 时：
    \(
    F_Y(y) = P\{-1\leqslant X\leqslant y\} = \frac{y-(-1)}{3} = \frac{y+1}{3},
    \)
    故 $f_Y(y)=\dfrac{1}{3}$；
    \item 当 $y\leqslant0$ 时，$F_Y(y)=0$；当 $y\geqslant2$ 时，$F_Y(y)=1$。
\end{itemize}
\end{solution}

\begin{exercise}\label{exer:src-2004-2a-12}
设随机变量 \(X\) 的概率密度函数为
\(
f(x)=
\begin{cases}
1-|x|, & -1<x<1,\\
0, & \text{其他}.
\end{cases}
\)
求随机变量 \(Y=X^2+1\) 的分布函数与概率密度函数。
\end{exercise}
\begin{solution}
先确定 \(Y\) 的取值范围。由于 \(-1<X<1\)，所以 \(0<X^2<1\)，从而
\(
1<Y<2.
\)

当 \(y\le 1\) 时，\(F_Y(y)=0\)。当 \(1<y<2\) 时，
\[
\begin{aligned}
F_Y(y)
&=P(X^2+1\le y)
=P(|X|\le \sqrt{y-1})=\int_{-\sqrt{y-1}}^{\sqrt{y-1}}(1-|x|)\,\mathrm{d}x\\
&=2\int_0^{\sqrt{y-1}}(1-x)\,\mathrm{d}x=2\sqrt{y-1}-(y-1).
\end{aligned}
\]
当 \(y\ge 2\) 时，\(F_Y(y)=1\)。故
\[
F_Y(y)=
\begin{cases}
0, & y\le 1,\\
2\sqrt{y-1}-(y-1), & 1<y<2,\\
1, & y\ge 2.
\end{cases}
\]
对 \(1<y<2\) 求导，得到
\(
f_Y(y)=\frac{1}{\sqrt{y-1}}-1.
\)
因此
\[
f_Y(y)=
\begin{cases}
\dfrac{1}{\sqrt{y-1}}-1, & 1<y<2,\\
0, & \text{其他}.
\end{cases}
\]
\end{solution}


\begin{exercise}\label{exer:src-2009-07-3a-i-02}\label{exer:src-probstat-test1-i-02}
设 \(X_1,X_2,\dots,X_n\) 独立同分布，且
\[
X_1\sim B(1,p),\qquad \overline X=\frac1n\sum_{i=1}^n X_i.
\]
则 \(P\{\overline X=k/n\}=\)（\quad）。\par
A. \(p\)\quad
B. \(1-p\)\quad
C. \(\binom nk p^k(1-p)^{n-k}\)\quad
D. \(\binom nk(1-p)^k p^{n-k}\)
\end{exercise}
\begin{solution}
\[
P\left\{\overline X=\frac{k}{n}\right\}
=P\left\{\sum_{i=1}^{n}X_i=k\right\}
=\binom nk p^k(1-p)^{n-k}.
\]
\end{solution}


\begin{exercise}\label{exer:src-2010-a-04}
已知某家电在 \(t=0\) 时刻正常运行。它在时刻 \(t\) 还正常运行的条件下，在 \((t,t+\Delta t)\) 这段时间损坏的概率等于
\(
\lambda\Delta t+o(\Delta t).
\)
求它正常运行时间大于 \(t\) 的概率。
\end{exercise}
\begin{solution}
设 \(\xi\) 表示正常运行时间，\(F(t)=P(\xi\le t)\)。题设给出瞬时失效率为常数 \(\lambda\)，即
\[
P(t<\xi\le t+\Delta t\,|\,\xi>t)=\lambda\Delta t+o(\Delta t).
\]
于是
\(
\frac{F(t+\Delta t)-F(t)}{1-F(t)}
=\lambda\Delta t+o(\Delta t).
\)
两边除以 \(\Delta t\)，令 \(\Delta t\to0\)，得
\(
F'(t)=\lambda[1-F(t)].
\)
令 \(S(t)=1-F(t)=P(\xi>t)\)，则
\(
S'(t)=-\lambda S(t),\qquad S(0)=1.
\)
解得
\(
S(t)=e^{-\lambda t}.
\)
因此
\(
P(\xi>t)=e^{-\lambda t},\qquad t\ge0.
\)
\end{solution}

\begin{exercise}\label{exer:src-2010-a-05}
假定某地区离婚率是 \(p\ (0<p<1)\)。为了某研究需要，决定从此地区逐个随机抽取调查对象（假设每次抽取的概率相等，并相互独立），直到抽取 \(m\) 个离婚人士为止，共抽取了 \(\xi\) 位人调查。求：\par
(1) \(\xi\) 的分布律；\par
(2) \(\xi\) 的数学期望。
\end{exercise}
\begin{solution}
\(\xi\) 表示得到第 \(m\) 个“成功”（离婚人士）时的总抽取人数，因此 \(\xi\) 服从负二项型分布，取值为
\(
\xi=m,m+1,m+2,\dots.
\)
若 \(\xi=k\)，则前 \(k-1\) 次中恰有 \(m-1\) 次抽到离婚人士，且第 \(k\) 次抽到离婚人士。因此
\(
P(\xi=k)=\binom{k-1}{m-1}p^m(1-p)^{k-m},
\qquad k=m,m+1,\dots.
\)

把 \(\xi\) 看成取得 \(m\) 次成功所需等待时间之和：
\(
\xi=\xi_1+\xi_2+\cdots+\xi_m,
\)
其中各 \(\xi_i\) 独立同分布，且均服从参数为 \(p\) 的几何分布，\(E\xi_i=1/p\)。故
\(
E\xi=\sum_{i=1}^m E\xi_i=\frac mp.
\)
\end{solution}


\begin{exercise}\label{exer:src-2008-01-3c-iii-02}\label{exer:src-paper4-iii-02}
设随机变量 \(X,Y\) 相互独立，且都服从 \([0,1]\) 上的均匀分布。令
\(
Z=|X-Y|.
\)
求 \(Z\) 的分布函数 \(F_Z(z)\)。
\end{exercise}
\begin{solution}
由于 \(X,Y\) 相互独立且都服从 \([0,1]\) 上的均匀分布，故联合密度为
\[
f_{X,Y}(x,y)=
\begin{cases}
1, & 0\le x\le1,\ 0\le y\le1,\\
0, & \text{其他}.
\end{cases}
\]

令 \(Z=|X-Y|\)。用分布函数法：
\[
F_Z(z)=P(Z\le z)=P(|X-Y|\le z)
=\iint_{\{(x,y):0\le x,y\le1,\ |x-y|\le z\}}1\,\mathrm dx\mathrm dy.
\]
当 \(z<0\) 时，事件 \(\{|X-Y|\le z\}\) 不可能发生，故 \(F_Z(z)=0\)。当 \(z\ge1\) 时，由于 \(|X-Y|\le1\) 恒成立，故 \(F_Z(z)=1\)。
当 \(0\le z<1\) 时，在单位正方形内，事件 \(\{|X-Y|>z\}\) 由两个全等三角形组成
\(
\{Y>X+z\},\qquad \{X>Y+z\}.
\)，
其中
\[
P(Y>X+z)=\int_0^{1-z}\int_{x+z}^{1}1\,\mathrm dy\mathrm dx
=\int_0^{1-z}(1-z-x)\,\mathrm dx
=\frac{(1-z)^2}{2},~~P(X>Y+z)=\dfrac{(1-z)^2}{2}
\]
因此
\(
F_Z(z)=1-P(|X-Y|>z)
=1-(1-z)^2
=2z-z^2.
\)，
综上，
\[
F_Z(z)=
\begin{cases}
0, & z<0,\\
2z-z^2, & 0\le z<1,\\
1, & z\ge1.
\end{cases}
\]
\end{solution}


\chapter{二维随机变量及分布}
\[
\begin{cases}
\text{二维随机变量的概念：二维平面中随机点的坐标} \\
\text{二维随机变量的概率分布定义及性质}
\begin{cases}
(1) \text{ 单调不减性} \\
(2) \text{ 规范性} \\
(3) \text{ 右连续性}
\end{cases}
\text{ 及 }
\begin{cases}
\text{边缘分布} \\
\text{条件分布}
\end{cases} \\

\text{二维离散型}
\begin{cases}
\text{联合分布律：} P\{X=x_i,Y=y_j\}=p_{ij};\ p_{ij}\geqslant0,\ \sum\limits_{i=1}^{\infty}\sum\limits_{j=1}^{\infty}p_{ij}=1 \\[8pt]
\text{边缘分布律：}
\begin{cases}
P\{X=x_i\}=\sum\limits_{j=1}^{\infty}p_{ij}=p_{i\cdot} \\[8pt]
P\{Y=y_j\}=\sum\limits_{i=1}^{\infty}p_{ij}=p_{\cdot j}
\end{cases} \\[12pt]
\text{条件分布律：}
\begin{cases}
P\{X=x_i|Y=y_j\}=\dfrac{p_{ij}}{p_{\cdot j}} \\[8pt]
P\{Y=y_j|X=x_i\}=\dfrac{p_{ij}}{p_{i\cdot}}
\end{cases}
\end{cases} \\[18pt]

\text{二维连续型}
\begin{cases}
\text{联合概率密度：} f(x,y)\geqslant0,\ \displaystyle\int_{-\infty}^{+\infty}\int_{-\infty}^{+\infty}f(u,v)\mathrm{d}u\mathrm{d}v=1 \\[12pt]
\text{边缘分布：}
\begin{cases}
f_X(x)=\displaystyle\int_{-\infty}^{+\infty}f(x,y)\mathrm{d}y,\ F_X(x)=\displaystyle\int_{-\infty}^{x}f_X(t)\mathrm{d}t \\[12pt]
f_Y(y)=\displaystyle\int_{-\infty}^{+\infty}f(x,y)\mathrm{d}x,\ F_Y(y)=\displaystyle\int_{-\infty}^{y}f_Y(s)\mathrm{d}s
\end{cases} \\[18pt]
\text{条件概率密度：}
\begin{cases}
f_{X|Y}(x|y)=\dfrac{f(x,y)}{f_Y(y)} \\[8pt]
f_{Y|X}(y|x)=\dfrac{f(x,y)}{f_X(x)}
\end{cases}
\end{cases} \\[18pt]

\text{随机变量的独立性：}
\begin{cases}
\text{离散型：} p_{ij}=p_{i\cdot}p_{\cdot j} \\
\text{连续型：} f(x,y)=f_X(x)\cdot f_Y(y)
\end{cases}
\iff F(x,y)=F_X(x)\cdot F_Y(y) \\[12pt]

\text{二维随机变量函数的概率分布}
\begin{cases}
\text{离散型：逐点法或表格法} \\
\text{连续型：卷积公式法，分布函数法}
\end{cases}
\end{cases}
\]

\section{二维随机变量}

\begin{definition}[二维随机变量] \label{def:2d-random-variable}
设 \(E\) 是一个随机试验，其样本空间为 \(\Omega\)。若 \(X=X(\omega)\) 和 \(Y=Y(\omega)\) 是定义在同一样本空间 \(\Omega\) 上的两个随机变量，则由它们构成的向量 \((X,Y)\) 称为\textbf{二维随机向量}（或\textbf{二维随机变量}）。
\end{definition}

\begin{definition}[联合分布函数] \label{def:joint-df}
设 $(X,Y)$ 是二维随机变量，对于任意实数 $x,y$，称二元函数
\(
F(x,y) = P\{X \leqslant x, Y \leqslant y\}
\)
为二维随机变量 $(X,Y)$ 的\textbf{分布函数}，或称为随机变量 $X$ 和 $Y$ 的\textbf{联合分布函数}，记为 $(X,Y) \sim F(x,y)$。
\end{definition}
\begin{property}\label{prop:joint-df}
联合分布函数 $F(x,y)$ 具有以下基本性质：
\begin{enumerate}
    \item \textbf{单调性}：$F(x,y)$ 关于每个变量单调不减。即对任意固定的 $y$，当 $x_1 < x_2$ 时，$F(x_1,y) \leqslant F(x_2,y)$；对任意固定的 $x$，当 $y_1 < y_2$ 时，$F(x,y_1) \leqslant F(x,y_2)$。
    \item \textbf{有界性}：$0 \leqslant F(x,y) \leqslant 1$，且
    \[
    F(-\infty, y) = 0, \quad F(x, -\infty) = 0, \quad F(-\infty, -\infty) = 0, \quad F(+\infty, +\infty) = 1.
    \]
    \item \textbf{右连续性}：$F(x,y)$ 关于每个变量右连续，即
    \(
    F(x+0, y) = F(x,y), \quad F(x, y+0) = F(x,y).
    \)
    \item \textbf{非负性}：对于任意的 $x_1 < x_2$，$y_1 < y_2$，有
    \[
    P\{x_1 < X \leqslant x_2,\ y_1 < Y \leqslant y_2\} = F(x_2,y_2) - F(x_2,y_1) - F(x_1,y_2) + F(x_1,y_1) \geqslant 0.
    \]
\end{enumerate}
性质 (4) 是二维分布函数特有的性质，也是区别于一维情形的关键。仅满足前三个性质的函数不一定能成为某个二维随机变量的联合分布函数，必须同时满足性质 (4)。例如，设
\[
F(x,y) = \begin{cases} 1, & x+y \geqslant 0, \\ 0, & x+y < 0, \end{cases}
\]
满足性质 (1)--(3)，但不满足性质 (4)，因此不是联合分布函数。
\end{property}


\begin{definition}[边缘分布函数] \label{def:marginal-df}
设二维随机变量 $(X,Y)$ 的联合分布函数为 $F(x,y)$，分别称
\begin{align*}
F_X(x) = P\{X \leqslant x\} = F(x, +\infty) = \lim_{y \to +\infty} F(x,y)\\
 F_Y(y) = P\{Y \leqslant y\} = F(+\infty, y) = \lim_{x \to +\infty} F(x,y)
\end{align*}
为 $(X,Y)$ 关于 $X$ 和关于 $Y$ 的\textbf{边缘分布函数}。
\end{definition}

\begin{exercise}\label{exer:3-1}
设 $(X,Y)$ 的联合分布函数为 $F(x,y) = a\left(b + \arctan\dfrac{x}{2}\right)\left(c + \arctan\dfrac{y}{2}\right)$，求常数 $a, b, c$ 及 $(X,Y)$ 的概率密度。
\end{exercise}
\begin{solution}
由联合分布函数的性质：\\
(1) $F(-\infty, y) = a\left(b - \dfrac{\pi}{2}\right)\left(c + \arctan\dfrac{y}{2}\right) = 0$，由于此式对所有 $y$ 成立，故 $b = \dfrac{\pi}{2}$。\\
(2) $F(x, -\infty) = a\left(b + \arctan\dfrac{x}{2}\right)\left(c - \dfrac{\pi}{2}\right) = 0$，同理得 $c = \dfrac{\pi}{2}$。\\
(3) $F(+\infty, +\infty) = a\left(\dfrac{\pi}{2} + \dfrac{\pi}{2}\right)\left(\dfrac{\pi}{2} + \dfrac{\pi}{2}\right) = a\pi^2 = 1$，故 $a = \dfrac{1}{\pi^2}$。\\
概率密度为联合分布函数的二阶混合偏导数：
\[
f(x,y) = \frac{\partial^2 F(x,y)}{\partial x\,\partial y} = \frac{1}{\pi^2} \cdot \frac{1/2}{1 + (x/2)^2} \cdot \frac{1/2}{1 + (y/2)^2} = \frac{4}{\pi^2(4+x^2)(4+y^2)}.
\]
\end{solution}

\begin{exercise}[边缘分布函数] \label{exer:3-c4}
如果二维随机变量 $(X,Y)$ 的分布函数为
\[
F(x,y) = \begin{cases} (1-e^{-\lambda_1 x})(1-e^{-\lambda_2 y} + e^{-\lambda_3 x-\lambda_2 y}), & x > 0,\ y > 0, \\ 0, & \text{其他}, \end{cases}
\]
其中 $\lambda_1, \lambda_2, \lambda_3 > 0$，求 $X$ 和 $Y$ 的边缘分布函数。
\end{exercise}
\begin{solution}
\[
F_X(x) = \lim_{y\to+\infty}F(x,y) = \begin{cases} (1-e^{-\lambda_1 x})\cdot 1 = 1-e^{-\lambda_1 x}, & x > 0, \\ 0, & \text{其他}. \end{cases}
\]
\[
F_Y(y) = \lim_{x\to+\infty}F(x,y) = \begin{cases} 1\cdot(1-e^{-\lambda_2 y}) = 1-e^{-\lambda_2 y}, & y > 0, \\ 0, & \text{其他}. \end{cases}
\]
\begin{note}
$X \sim E(\lambda_1)$，$Y \sim E(\lambda_2)$。注意 $F(x,y) \neq F_X(x)F_Y(y)$，故 $X$ 与 $Y$ 不独立（除非 $\lambda_3 = 0$）。
\end{note}
\end{solution}



\begin{exercise}[边缘密度不保证联合为正态] \label{exer:3-c5}
随机变量 $(X,Y)$ 的概率密度为
\[
f(x,y) = \frac{1+\sin x\sin y}{2\pi}\,e^{-\frac{x^2+y^2}{2}}, \quad -\infty < x, y < +\infty,
\]
求边缘概率密度 $f_X(x)$ 和 $f_Y(y)$。
\end{exercise}
\begin{solution}
$\sin y \cdot e^{-y^2/2}$ 是奇函数，在 $\mathbb{R}$ 上积分为 $0$，故
\[
f_X(x) = \int_{-\infty}^{+\infty}\frac{1+\sin x\sin y}{2\pi}\,e^{-\frac{x^2+y^2}{2}}\,\mathrm{d}y = \frac{e^{-x^2/2}}{2\pi}\left[\int_{-\infty}^{+\infty}e^{-y^2/2}\,\mathrm{d}y + \sin x \cdot 0\right] = \frac{1}{\sqrt{2\pi}}\,e^{-x^2/2}.
\]
同理 $f_Y(y) = \dfrac{1}{\sqrt{2\pi}}\,e^{-y^2/2}$。
\begin{note}
$X \sim N(0,1)$，$Y \sim N(0,1)$，但 $(X,Y)$ \textbf{不服从}二维正态分布（密度中含 $\sin x\sin y$ 项）。这说明：\textbf{边缘分布为正态不能推出联合分布为二维正态}。
\end{note}
\end{solution}


\section{n维随机变量}

\begin{definition}[$n$维随机变量] \label{def:n-d-random-variable}
如果 $X_1, X_2, \dots, X_n$ 是定义在同一个样本空间 $\Omega$ 上的 $n$ 个随机变量，则称 $(X_1, X_2, \dots, X_n)$ 为\textbf{$n$维随机变量}（或\textbf{$n$维随机向量}），$X_i$（$i=1,2,\dots,n$）称为第 $i$ 个\textbf{分量}。
\end{definition}

\begin{definition}[$n$维联合分布函数] \label{def:n-joint-df}
对任意的 $n$ 个实数 $x_1, x_2, \dots, x_n$，称 $n$ 元函数
\[
F(x_1, x_2, \dots, x_n) = P\{X_1 \leqslant x_1, X_2 \leqslant x_2, \dots, X_n \leqslant x_n\}
\]
为 $n$ 维随机变量 $(X_1, X_2, \dots, X_n)$ 的\textbf{联合分布函数}（或\textbf{分布函数}）。
\end{definition}

\begin{property}[$n$维联合分布函数的性质]\label{prop:n-joint-df-property}
$n$ 维联合分布函数 $F(x_1, x_2, \dots, x_n)$ 具有以下性质：
\begin{enumerate}
    \item \textbf{单调性}：$F$ 关于每个变量 $x_i$ 单调不减；
    \item \textbf{右连续性}：$F$ 关于每个变量 $x_i$ 右连续；
    \item \textbf{有界性}：
    \[
    F(\underbrace{-\infty, \dots, -\infty}_{k\text{个}}, x_{k+1}, \dots, x_n) = 0 \quad (k=1,2,\dots,n-1), \quad F(+\infty, +\infty, \dots, +\infty) = 1;
    \]
    \item \textbf{非负性}：对任意的 $a_i < b_i$（$i=1,2,\dots,n$），以 $F$ 计算的 $n$ 维矩形区域的概率
    \[
    P\{a_1 < X_1 \leqslant b_1, \dots, a_n < X_n \leqslant b_n\} \geqslant 0.
    \]
\end{enumerate}
\end{property}

\begin{definition}[$n$维边缘分布函数]\label{def:n-marginal-df}
$n$ 维随机变量 $(X_1, X_2, \dots, X_n)$ 中任意 $k$ 个（$1 \leqslant k < n$）分量所构成的 $k$ 维随机变量的分布函数，称为原 $n$ 维随机变量的\textbf{$k$维边缘分布函数}。特别地，关于第 $i$ 个分量 $X_i$ 的边缘分布函数为
\[
F_{X_i}(x_i) = F(+\infty, \dots, +\infty, x_i, +\infty, \dots, +\infty).
\]
\end{definition}

\begin{note}
$n$ 维边缘分布就是“只保留感兴趣的那几个分量，其余分量全部积分掉或放到 $+\infty$”。因此边缘分布反映的是局部信息，而联合分布保存的是整体联动信息。
\end{note}

类似地，对于 $n$ 维离散型随机变量，可以定义\textbf{联合分布律} \(p_{i_1 i_2 \cdots i_n} = P\{X_1 = x_{i_1}, X_2 = x_{i_2}, \dots, X_n = x_{i_n}\}\) 及相应的边缘分布律和条件分布律。

对于 $n$ 维连续型随机变量，可以定义\textbf{联合概率密度} $f(x_1, x_2, \dots, x_n)$，满足
\[
F(x_1, x_2, \dots, x_n) = \int_{-\infty}^{x_n}\cdots\int_{-\infty}^{x_1}f(t_1, t_2, \dots, t_n)\,\mathrm{d}t_1\cdots\mathrm{d}t_n,
\]
以及边缘概率密度、条件概率密度和概率密度乘法公式等，这些概念与二维情形完全类似。

\section{随机变量的独立性}

\begin{definition}[$n$ 个随机变量的相互独立性] \label{def:n-independence}
设 $n$ 维随机变量 $(X_1, X_2, \dots, X_n)$ 的联合分布函数为 $F(x_1, x_2, \dots, x_n)$，$X_i$ 的边缘分布函数为 $F_{X_i}(x_i)$。若对任意实数 $x_1, x_2, \dots, x_n$，都有
\(
F(x_1, x_2, \dots, x_n) = \prod_{i=1}^{n}F_{X_i}(x_i),
\)
则称 $X_1, X_2, \dots, X_n$ 是\textbf{相互独立}的。

进一步地，两个多维随机变量 $(X_1, \dots, X_n)$ 与 $(Y_1, \dots, Y_m)$ 相互独立，是指对任意实数 $x_i$（$i=1,\dots,n$）与 $y_j$（$j=1,\dots,m$），都有
\[
F(x_1, \dots, x_n, y_1, \dots, y_m) = F_X(x_1, \dots, x_n) \cdot F_Y(y_1, \dots, y_m),
\]
其中 $F_X$ 和 $F_Y$ 分别是 $(X_1, \dots, X_n)$ 和 $(Y_1, \dots, Y_m)$ 的分布函数。
\end{definition}

\begin{theorem}[离散型随机变量独立的判定]\label{thm:discrete-independence}
设 $(X,Y)$ 为二维离散型随机变量，其联合分布律为 $P\{X = x_i, Y = y_j\} = p_{ij}$，则 $X$ 与 $Y$ 相互独立的\textbf{充分必要条件}是对一切 $i, j$，都有
\(
p_{ij} = p_{i\cdot} \cdot p_{\cdot j}.
\)
更一般地，$n$ 个离散型随机变量 $X_1, X_2, \dots, X_n$ 相互独立的充分必要条件是对任意的 $x_i \in D_i$（$X_i$ 的一切可能取值），有
\(
P\{X_1 = x_1, \dots, X_n = x_n\} = \prod_{i=1}^{n}P\{X_i = x_i\}.
\)
\end{theorem}


\begin{theorem}[连续型随机变量独立的判定]\label{thm:continuous-independence}
设 $(X,Y)$ 为二维连续型随机变量，联合概率密度为 $f(x,y)$，边缘概率密度为 $f_X(x)$ 和 $f_Y(y)$，则 $X$ 与 $Y$ 相互独立的\textbf{充分必要条件}是
\(
f(x,y) = f_X(x) \cdot f_Y(y)
\)
几乎处处成立（即除了面积为零的集合外处处成立）。

更一般地，$n$ 维连续型随机变量 $(X_1, \dots, X_n)$ 相互独立的充分必要条件是
\(
f(x_1, x_2, \dots, x_n) = \prod_{i=1}^{n}f_i(x_i),
\)
其中 $f_i(x_i)$ 为 $X_i$ 的边缘概率密度。
\end{theorem}

\begin{property}\label{prop:independence-implications}
随机变量的独立性具有以下重要性质：
\begin{enumerate}
    \item 设 $(X,Y)$ 为二维离散型随机变量，$X$ 与 $Y$ 独立，则条件分布等于边缘分布：
    \begin{align*}
    P\{X = x_i \mid Y = y_j\} = P\{X = x_i\}\ (P\{Y = y_j\} > 0)\\
     P\{Y = y_j \mid X = x_i\} = P\{Y = y_j\}\ (P\{X = x_i\} > 0)
    \end{align*}
    \item 设 $(X,Y)$ 为二维连续型随机变量，$X$ 与 $Y$ 独立，则条件密度等于边缘密度：
    \(
    f_{X|Y}(x|y) = \frac{f(x,y)}{f_Y(y)} = f_X(x)\ (f_Y(y) > 0), \quad f_{Y|X}(y|x) = f_Y(y)\ (f_X(x) > 0).
    \)
    \item 若 $(X,Y)$ 服从二维正态分布，则 $X$ 与 $Y$ 相互独立当且仅当 $\rho = 0$。但一般而言，$X$ 与 $Y$ 不相关（$\mathrm{Cov}(X,Y) = 0$）\textbf{不能}推出 $X$ 与 $Y$ 独立。
    \item 独立一定不相关，但不相关不一定独立（除非联合分布为正态分布）。
\end{enumerate}
\end{property}

\begin{property}[不独立的判定]\label{prop:dependence-criterion}
$X$ 与 $Y$ 不独立 $\iff$ 存在 $x_0, y_0$，使得
\(
F(x_0, y_0) \neq F_X(x_0) \cdot F_Y(y_0).
\)
也可通过取特定值使 $P\{X \leqslant x_0, Y \leqslant y_0\} \neq P\{X \leqslant x_0\} \cdot P\{Y \leqslant y_0\}$ 来证明不独立。对于离散型，寻找一对 $(i,j)$ 使 $p_{ij} \neq p_{i\cdot} \cdot p_{\cdot j}$；对于连续型，验证 $f(x,y) \neq f_X(x)\,f_Y(y)$。
\end{property}

\begin{exercise}[证明 $X$ 与 $|X|$ 不独立] \label{exer:3-5c}
设随机变量 $X$ 的概率密度为
\(
f(x) = \frac{1}{2}e^{-|x|}, \quad x \in (-\infty, +\infty),
\)
证明：$X$ 与 $|X|$ 不独立。
\end{exercise}
\begin{solution}
\[
\begin{aligned}
P\{X\le 1\}
&=\int_{-\infty}^1 \frac{1}{2}e^{-|x|}\,\mathrm{d}x =\frac{1}{2}\int_{-\infty}^{0} e^x\,\mathrm{d}x+\frac{1}{2}\int_{0}^{1} e^{-x}\,\mathrm{d}x =\frac{1}{2}+\frac{1}{2}(1-e^{-1})
=1-\frac{1}{2}e^{-1},
\end{aligned}
\]
\[
\begin{aligned}
P\{|X|\le 1\}
&=\frac{1}{2}\int_{-1}^{1}e^{-|x|}\,\mathrm{d}x
=\int_{0}^{1}e^{-x}\,\mathrm{d}x
=1-e^{-1}.
\end{aligned}
\]
再看联合事件。因为 $\{|X|\le 1\}$ 已经蕴含 $X\le 1$，所以
\(
P\{X\le 1,\ |X|\le 1\}=P\{|X|\le 1\}=1-e^{-1}.
\)，
但实际上
\(
\left(1-\frac{1}{2}e^{-1}\right)(1-e^{-1})
=1-\frac{3}{2}e^{-1}+\frac{1}{2}e^{-2}
\neq 1-e^{-1}.
\)
于是
\(
P\{X\le 1,\ |X|\le 1\}\neq P\{X\le 1\}\,P\{|X|\le 1\},
\)
所以 $X$ 与 $|X|$ 不独立。
\end{solution}

\section{二维离散型随机变量}

\begin{definition}[联合分布律]  \label{def:joint-pm}
设二维离散型随机变量 $(X,Y)$ 所有可能取的值为 $(x_i, y_j)$ ($i,j = 1,2,\dots$)，则称
\(
p_{ij} = P\{X = x_i, Y = y_j\}, \quad i,j = 1,2,\dots
\)
为 $(X,Y)$ 的\textbf{联合分布律}，或称为随机变量 $X$ 和 $Y$ 的\textbf{联合概率分布}。
\end{definition}

联合分布律也可用表格形式表示：

\begin{center}
\begin{tabular}{c|ccccc|c}
\toprule
$X \backslash Y$ & $y_1$ & $y_2$ & $\cdots$ & $y_j$ & $\cdots$ & $p_{i\cdot}$ \\
\midrule
$x_1$ & $p_{11}$ & $p_{12}$ & $\cdots$ & $p_{1j}$ & $\cdots$ & $p_{1\cdot}$ \\
$x_2$ & $p_{21}$ & $p_{22}$ & $\cdots$ & $p_{2j}$ & $\cdots$ & $p_{2\cdot}$ \\
$\vdots$ & $\vdots$ & $\vdots$ & & $\vdots$ & & $\vdots$ \\
$x_i$ & $p_{i1}$ & $p_{i2}$ & $\cdots$ & $p_{ij}$ & $\cdots$ & $p_{i\cdot}$ \\
$\vdots$ & $\vdots$ & $\vdots$ & & $\vdots$ & & $\vdots$ \\
\midrule
$p_{\cdot j}$ & $p_{\cdot1}$ & $p_{\cdot2}$ & $\cdots$ & $p_{\cdot j}$ & $\cdots$ & $1$ \\
\bottomrule
\end{tabular}
\end{center}


\begin{definition}[边缘分布律] \label{def:marginal-pm}
设 $(X,Y)$ 的联合分布律为 $P\{X = x_i, Y = y_j\} = p_{ij}$，则分别称
\[
p_{i\cdot} = P\{X = x_i\} = \sum_{j=1}^{\infty}p_{ij}, \quad i = 1,2,\dots
\]
\[
p_{\cdot j} = P\{Y = y_j\} = \sum_{i=1}^{\infty}p_{ij}, \quad j = 1,2,\dots
\]
为 $(X,Y)$ 关于 $X$ 和关于 $Y$ 的\textbf{边缘分布律}。
\end{definition}


\begin{definition}[条件分布律] \label{def:conditional-pm}
设 $(X,Y)$ 为二维离散型随机变量，其联合分布律为 $P\{X = x_i, Y = y_j\} = p_{ij}$ ($i,j = 1,2,\dots$)。

若对固定的 $j$，$p_{\cdot j} = P\{Y = y_j\} > 0$，则称
\[
P\{X = x_i \mid Y = y_j\} = \frac{P\{X = x_i, Y = y_j\}}{P\{Y = y_j\}} = \frac{p_{ij}}{p_{\cdot j}}, \quad i = 1,2,\dots
\]
为在 $\{Y = y_j\}$ 条件下 $X$ 的\textbf{条件分布律}。

同理，若对固定的 $i$，$p_{i\cdot} = P\{X = x_i\} > 0$，则称
\[
P\{Y = y_j \mid X = x_i\} = \frac{p_{ij}}{p_{i\cdot}}, \quad j = 1,2,\dots
\]
为在 $\{X = x_i\}$ 条件下 $Y$ 的\textbf{条件分布律}。
\end{definition}

\begin{exercise}\label{exer:src-2008-probstat-2a-i-01}
设随机变量 \(X_1,X_2\) 的分布律均为
\[
\begin{array}{c|ccc}
X_i&-1&0&1\\
\hline
P&\dfrac14&\dfrac12&\dfrac14
\end{array}
\qquad (i=1,2),
\]
且满足 \(P\{X_1X_2=0\}=1\)，则 \(P\{X_1=X_2\}=\)（\quad）。~~
A. \(0\)\quad
B. \(\dfrac14\)\quad
C. \(\dfrac12\)\quad
D. \(1\)
\end{exercise}
\begin{solution}
\(P\{X_1X_2=0\}=1\)的意思是事件“\(X_1X_2=0\)”必然发生，概率为 1。
\begin{center}
\begin{tabular}{c|ccc|c}
\toprule
$X_1 \backslash X_2$ & $-1$ & $0$ & $1$ & $p_{i\cdot}$ \\
\midrule
$-1$ & $0$ & $\frac{1}{4}$ & $0$ & $\frac{1}{4}$ \\
$0$ & $\frac{1}{4}$ & $0$ & $\frac{1}{4}$ & $\frac{1}{2}$ \\
$1$ & $0$ & $\frac{1}{4}$ & $0$ & $\frac{1}{4}$ \\
\midrule
$p_{\cdot j}$ & $\frac{1}{4}$ & $\frac{1}{2}$ & $\frac{1}{4}$ & $1$ \\
\bottomrule
\end{tabular}
\end{center}
\end{solution}


\begin{exercise}[行列式分布] \label{exer:3-c6}
设随机变量 $X_1, X_2, X_3, X_4$ 相互独立且同分布，$P\{X_i=0\} = 0.6$，$P\{X_i=1\} = 0.4$（$i=1,2,3,4$），求行列式 $Z = \begin{vmatrix} X_1 & X_2 \\ X_3 & X_4 \end{vmatrix}$ 的概率分布。
\end{exercise}
\begin{solution}
$Z = X_1X_4 - X_2X_3$，可能取值为 $\{-1, 0, 1\}$。
\begin{align*}
P\{Z=1\} &= P\{X_1X_4=1\}\cdot P\{X_2X_3=0\}\\
& = P\{X_1=1,X_4=1\}\cdot P\{(X_2=0)\cup(X_3=0)\}\\
&= 0.4^2 \cdot (1-0.4^2) = 0.16 \times 0.84 = 0.1344.\\
P\{Z=-1\} &= P\{X_1X_4=0\}\cdot P\{X_2X_3=1\} \\
&= (1-0.4^2)\cdot 0.4^2 = 0.84 \times 0.16 = 0.1344\\
P\{Z=0\} &= 1 - 0.1344 - 0.1344 = 0.7312.
\end{align*}
所以
\[
Z \sim \begin{pmatrix} -1 & 0 & 1 \\ 0.1344 & 0.7312 & 0.1344 \end{pmatrix}.
\]
\end{solution}

\begin{exercise}[离散型函数的分布] \label{exer:3-c9}
已知 $(X,Y)$ 的概率分布为
\begin{center}
\begin{tabular}{c|cc}
\toprule
$X \backslash Y$ & $-1$ & $1$ \\
\midrule
$-1$ & $1/3$ & $1/6$ \\
$0$ & $1/3$ & $1/6$ \\
\bottomrule
\end{tabular}
\end{center}
求：(1) $Z = X - Y$ 的概率分布；(2) 设 $U = \max\{X,Y\}$，$V = \min\{X,Y\}$，求 $(U,V)$ 及 $UV$ 的分布。
\end{exercise}
\begin{solution}
(1) $Z$ 的可能取值：$Z(-1,-1)=0$，$Z(-1,1)=-2$，$Z(0,-1)=1$，$Z(0,1)=-1$。
\[
P\{Z=-2\} = 1/6,\quad P\{Z=-1\} = 1/6,\quad P\{Z=0\} = 1/3,\quad P\{Z=1\} = 1/3.
\]
\[
Z \sim \begin{pmatrix} -2 & -1 & 0 & 1 \\ 1/6 & 1/6 & 1/3 & 1/3 \end{pmatrix}.
\]
(2) $(U,V)$ 联合分布律：
\begin{center}
\begin{tabular}{c|ccc}
\toprule
$V \backslash U$ & $-1$ & $0$ & $1$ \\
\midrule
$-1$ & $1/3$ & $1/3$ & $1/6$ \\
$0$ & $0$ & $0$ & $1/6$ \\
\bottomrule
\end{tabular}
\end{center}
$UV$ 的分布：$UV \in \{-1, 0, 1\}$。
\[
UV \sim \begin{pmatrix} -1 & 0 & 1 \\ 1/6 & 1/2 & 1/3 \end{pmatrix}.
\]
\end{solution}


\begin{exercise}[独立性条件求参数] \label{exer:3-d2}
设 $X,Y$ 相互独立，$P(X=1)=P(Y=1)=p$，$P(X=0)=P(Y=0)=1-p$（$0<p<1$）。令
\[
Z = \begin{cases} 1, & X+Y \text{为偶数}, \\ 0, & X+Y \text{为奇数}. \end{cases}
\]
要使 $X$ 与 $Z$ 独立，则 $p = $（\quad）
\end{exercise}
\begin{solution}
$X+Y$ 为偶数当且仅当 $(X,Y) = (0,0)$ 或 $(1,1)$，故
\(
P(Z=1) = (1-p)^2+p^2,\quad P(Z=0) = 2p(1-p).
\)
又 $P(X=1,Z=1) = P(X=1,Y=1) = p^2$。由独立性 $P(X=1,Z=1) = P(X=1)\,P(Z=1)$：
\(
p^2 = p\left[(1-p)^2+p^2\right].
\)
若 $p=0$ 或 $p=1$，$Z$ 退化为常数，平凡独立。对 $0<p<1$，两边除以 $p$：
\[
p = 1-2p+2p^2 \implies 2p^2-3p+1 = 0 \implies (2p-1)(p-1) = 0 \implies p = 1/2.
\]
\end{solution}


\begin{exercise}\label{exer:src-2017-201207-a-09-2credit}
今有两封信投入编号为 I、II、III 的 \(3\) 个邮箱，设 \(X,Y\) 分别表示投入第 I 号和第 II 号邮箱信的数目。求：\par
(1) \((X,Y)\) 的联合分布；\par
(2) \(X\) 和 \(Y\) 是否独立；\par
(3) 随机变量 \(U=\max(X,Y)\) 及 \(V=\min(X,Y)\) 的分布律。
\end{exercise}
\begin{solution}
每封信独立且等可能地投入三个邮箱。若 \(i,j\ge0,\ i+j\le2\)，则
\(
P(X=i,Y=j)=\frac{2!}{i!\,j!\,(2-i-j)!}
\left(\frac13\right)^2.
\)
因此 \((X,Y)\) 的联合分布列为
\[
\begin{array}{c|ccc}
\toprule
X\backslash Y & 0 & 1 & 2\\
\midrule
0 & \dfrac19 & \dfrac29 & \dfrac19\\
1 & \dfrac29 & \dfrac29 & 0\\
2 & \dfrac19 & 0 & 0\\
\bottomrule
\end{array}
\]
由边缘分布可得
\(
P(X=0)=P(Y=0)=\frac49.
\)
但
\(
P(X=0,Y=0)=\frac19\ne\frac49\cdot\frac49,
\)
故 \(X,Y\) 不相互独立。

由联合分布逐项合并，得
\[
\begin{array}{c|ccc}
\toprule
u & 0 & 1 & 2\\
\midrule
P(U=u) & \dfrac19 & \dfrac69 & \dfrac29\\
\bottomrule
\end{array}
\qquad
\begin{array}{c|ccc}
\toprule
v & 0 & 1 & 2\\
\midrule
P(V=v) & \dfrac79 & \dfrac29 & 0\\
\bottomrule
\end{array}
\]
\end{solution}



\begin{exercise}[泊松变量的复合函数] \label{exer:3-d4}
设 $X_1,X_2,X_3$ 相互独立且同服从 $P(1)$，记
\[
X = \begin{cases} 1, & X_1+X_2=1, \\ 0, & X_1+X_2\neq 1, \end{cases} \quad Y = \begin{cases} 1, & X_2+X_3=1, \\ 0, & X_2+X_3\neq 1. \end{cases}
\]
求：(1) $(X,Y)$ 的联合分布律；(2) 边缘分布律；(3) $P(X+Y \leqslant 1)$。
\end{exercise}
\begin{solution}
(1) $\{X=1,Y=1\}$ 要求 $X_1+X_2=1$ 且 $X_2+X_3=1$。分两种情况：\\
\textbf{情形 $X_2=0$}：需 $X_1=1$，$X_3=1$。$P = e^{-1}\cdot e^{-1}\cdot e^{-1} = e^{-3}$。\\
\textbf{情形 $X_2=1$}：需 $X_1=0$，$X_3=0$。$P = e^{-1}\cdot e^{-1}\cdot e^{-1} = e^{-3}$。\\
当 $X_2 \geqslant 2$ 时 $X_1+X_2 \geqslant 2 \neq 1$，不贡献。
\(
P(X=1,Y=1) = 2e^{-3}.
\)
边缘概率：
\(
P(X=1) = P(X_1+X_2=1) = P(1,0)+P(0,1) = 2e^{-2}, \qquad P(Y=1) = P(X_2+X_3=1) = 2e^{-2}.
\)
\begin{align*}
P(X=0,Y=1) &= P(Y=1)-P(X=1,Y=1) = 2e^{-2}-2e^{-3}, \\
P(X=1,Y=0) &= P(X=1)-P(X=1,Y=1) = 2e^{-2}-2e^{-3}, \\
P(X=0,Y=0) &= 1-2e^{-3}-(2e^{-2}-2e^{-3})-(2e^{-2}-2e^{-3}) = 1-4e^{-2}+2e^{-3}.
\end{align*}
\begin{center}
\begin{tabular}{c|cc|c}
\toprule
$X \backslash Y$ & $1$ & $0$ & $P\{X=x_i\}$ \\
\midrule
$1$ & $2e^{-3}$ & $2e^{-2}-2e^{-3}$ & $2e^{-2}$ \\
$0$ & $2e^{-2}-2e^{-3}$ & $1-4e^{-2}+2e^{-3}$ & $1-2e^{-2}$ \\
\midrule
$P\{Y=y_j\}$ & $2e^{-2}$ & $1-2e^{-2}$ & $1$ \\
\bottomrule
\end{tabular}
\end{center}
(2) 边缘分布律见上表右侧和下方。\\
(3) $P(X+Y \leqslant 1) = 1-P(X=1,Y=1) = 1-2e^{-3}$。
\begin{note}
$X$ 与 $Y$ 不独立：$P(X=1,Y=1)=2e^{-3}$，而 $P(X=1)P(Y=1)=4e^{-4} \neq 2e^{-3}$。这是因为 $X$ 和 $Y$ 共同依赖 $X_2$。
\end{note}
\end{solution}

\section{二维连续型随机变量}

\begin{definition}[二维连续型随机变量] \label{def:2d-continuous}
设 $(X,Y)$ 为二维随机变量，若存在非负可积函数 $f(x,y)$，使得对于任意实数 $x,y$，都有
\[
F(x,y) = \int_{-\infty}^{y}\int_{-\infty}^{x}f(u,v)\,\mathrm{d}u\,\mathrm{d}v,
\]
则称 $(X,Y)$ 为\textbf{二维连续型随机变量}，称 $f(x,y)$ 为 $(X,Y)$ 的\textbf{联合概率密度函数}（简称\textbf{概率密度}），记为 $(X,Y) \sim f(x,y)$。
\end{definition}

\begin{exercise}\label{exer:src-2017-201207-a-07-2credit}
设 \((X,Y)\) 的联合密度为
\(
f(x,y)=24y(1-x),\qquad 0\le x\le1,\ 0\le y\le x,
\)
在其他地方为 \(0\)。求\((X,Y)\) 的联合分布函数。
\end{exercise}
\begin{solution}
联合分布函数的定义为
\[
F(x,y) = \int_{-\infty}^{y}\int_{-\infty}^{x}f(u,v)\,\mathrm{d}u\,\mathrm{d}v,
\]
所以$(x,y)$可能取到平面所有点，需要分如图所示5种情况讨论。
\begin{center}
\includegraphics[width=0.2\textwidth]{figure/2.png}\\
{\small 图：区域划分示意图}
\end{center}
\[
F(x,y)=
\begin{cases}
0, & x<0\ \text{或}\ y<0,\\[3pt]
3y^4-8y^3+12\left(x-\dfrac{x^2}{2}\right)y^2, & 0\le x<1,\ 0\le y<x,\\
3y^4-8y^3+6y^2, & x\ge1,\ 0\le y<1,\\
4x^3-3x^4, & 0\le x<1,\ y\ge x,\\[3pt]
1, & x\ge1,\ y\ge1.
\end{cases}
\]
\end{solution}

\begin{exercise}[概率计算] \label{exer:3-c12}
$f(x,y) = \begin{cases} Ae^{-y}, & 0 < x < y, \\ 0, & \text{其他}. \end{cases}$，求 $A$，$P\{X+Y \geqslant 1\}$，$P\{X/Y \leqslant 1/2\}$。
\end{exercise}
\begin{solution}
(1) \textbf{求解常数 $\mathbf{A}$}：由联合概率密度函数的归一性（即全空间积分为 1）可知：
\[
\iint\limits_{\mathbb{R}^2} f(x,y) \,\mathrm{d}x\,\mathrm{d}y = 1
\]
根据题目给定的非零区域 $0 < x < y < +\infty$，我们先对 $x$ 积分，再对 $y$ 积分，确定上下限如下：
\[
\int_{0}^{+\infty} \left( \int_{0}^{y} A e^{-y} \,\mathrm{d}x \right) \mathrm{d}y = 1
\]
计算内层积分（此时 $y$ 视为常数）：
\[
A \int_{0}^{+\infty} e^{-y} \Big[ x \Big]_{0}^{y} \,\mathrm{d}y = A \int_{0}^{+\infty} y e^{-y} \,\mathrm{d}y = 1
\]
利用分部积分法或 Gamma 函数性质 $\Gamma(2) = \int_0^\infty y e^{-y} \mathrm{d}y = 1! = 1$，易得：
\(
A \times 1 = 1 \implies A = 1
\)\\
(2) \textbf{求解 $\mathbf{P\{X+Y \geqslant 1\}}$}：若直接求解 $X+Y \geqslant 1$，积分区域是一个向右上角无限延伸的无界区域，计算较为繁琐。因此，我们利用对立事件（补集）来简化计算：计算图形封闭的 $P\{X+Y < 1\}$。
\(
P\{X+Y \geqslant 1\} = 1 - P\{X+Y < 1\}
\)
事件 $\{X+Y < 1\}$ 对应的概率本质上是联合密度函数在特定二维区域 $D$ 上的二重积分。该区域 $D$ 由以下三个不等式共同围成：
\[
\begin{cases}
0 < x < y & \text{(密度函数非零的固有条件)} \\
x + y < 1 & \text{(题目要求的事件条件)}
\end{cases}
\]
在平面直角坐标系中，直线 $y = x$ 与直线 $x + y = 1$ 的交点为 $(1/2, 1/2)$。观察该三角形区域 $D$，我们可以采用 $X$-型区域（先对 $y$ 积分，后对 $x$ 积分）来设置上下限：
外层 $x$ 的投影范围是 $(0, 1/2)$；对于固定的 $x$，内层 $y$ 的下界是直线 $y=x$，上界是直线 $y=1-x$。因此：
\begin{align*}
P\{X+Y < 1\} &= \iint\limits_{D} f(x,y) \,\mathrm{d}x\,\mathrm{d}y = \int_{0}^{1/2} \mathrm{d}x \int_{x}^{1-x} e^{-y} \,\mathrm{d}y \\
&= \int_{0}^{1/2} \Big[ -e^{-y} \Big]_{x}^{1-x} \mathrm{d}x = \int_{0}^{1/2} \left( e^{-x} - e^{-(1-x)} \right) \mathrm{d}x \\
&= \Big[ -e^{-x} - e^{x-1} \Big]_{0}^{1/2} = \left( -e^{-1/2} - e^{-1/2} \right) - \left( -1 - e^{-1} \right)= 1 + e^{-1} - 2e^{-1/2}
\end{align*}
最后，代回对立事件公式中：
\[
P\{X+Y \geqslant 1\} = 1 - \left( 1 + e^{-1} - 2e^{-1/2} \right) = 2e^{-1/2} - e^{-1}
\]
(3) \textbf{求解 $\mathbf{P\{X/Y \leqslant 1/2\}}$}：由于在非零区域内 $y > 0$，不等式 $X/Y \leqslant 1/2$ 可直接等价变形为 $X \leqslant Y/2$。同样地，我们需要寻找联合密度函数非零区域 $\{0 < x < y\}$ 与目标事件区域 $\{x \leqslant y/2\}$ 的交集。综合这两个条件，得到真正的积分区域：$0 < x \leqslant y/2 < +\infty$。这次我们采用 $Y$-型区域（先对 $x$ 积分，后对 $y$ 积分）来设置上下限更为简便：外层 $y$ 的范围是 $(0, +\infty)$；对于固定的 $y$，内层 $x$ 的下界是 $0$，上界是 $y/2$。
\begin{align*}
P\{X \leqslant Y/2\} &= \int_{0}^{+\infty} \mathrm{d}y \int_{0}^{y/2} e^{-y} \,\mathrm{d}x = \int_{0}^{+\infty} e^{-y} \Big[ x \Big]_{0}^{y/2} \mathrm{d}y = \int_{0}^{+\infty} \frac{y}{2} e^{-y} \,\mathrm{d}y = \frac{1}{2} \int_{0}^{+\infty} y e^{-y} \,\mathrm{d}y
\end{align*}
由第(1)问已知的结论 $\int_{0}^{+\infty} y e^{-y} \,\mathrm{d}y = 1$，故：
\(
P\{X/Y \leqslant 1/2\} = \frac{1}{2} \times 1 = \frac{1}{2}
\)
\end{solution}

\begin{exercise}[混合型：连续+离散] \label{exer:3-c13}
$X \sim E(4)$，$Y$ 为离散型随机变量：$P\{Y=-1\} = 1/4$，$P\{Y=1\} = 3/4$，且 $X$ 与 $Y$ 相互独立。求 $P\{X-Y \leqslant 1\}$ 和 $P\{XY \leqslant 2\}$。
\end{exercise}
\begin{solution}
已知 $X \sim E(4)$，即 $X$ 服从参数 $\lambda = 4$ 的指数分布。根据指数分布的定义，其概率密度函数 $f_X(x)$ 和累积分布函数 $F_X(x)$ 分别为：
\[
f_X(x) = \begin{cases} 4e^{-4x}, & x > 0 \\ 0, & x \leqslant 0 \end{cases}
\quad \text{与} \quad
F_X(x) = P\{X \leqslant x\} = \begin{cases} 1 - e^{-4x}, & x > 0 \\ 0, & x \leqslant 0 \end{cases}
\]
\textbf{（1）求解 $\mathbf{P\{X-Y \leqslant 1\}}$}:根据 $Y$ 的取值将目标事件 $\{X-Y \leqslant 1\}$ 分解为两个互斥子事件的和：
\begin{align*}
P\{X-Y \leqslant 1\} &= P\big( \{X-Y \leqslant 1\} \cap \{Y=-1\} \big) + P\big( \{X-Y \leqslant 1\} \cap \{Y=1\} \big) \\
&= P\big( \{X-(-1) \leqslant 1\} \cap \{Y=-1\} \big) + P\big( \{X-1 \leqslant 1\} \cap \{Y=1\} \big) \\
&= P\big( \{X \leqslant 0\} \cap \{Y=-1\} \big) + P\big( \{X \leqslant 2\} \cap \{Y=1\} \big)
\end{align*}
利用 $X, Y$ 的相互独立性，将联合概率拆分为边缘概率的乘积：
\begin{align*}
P\{X-Y \leqslant 1\} &= P\{X \leqslant 0\} \cdot P\{Y=-1\} + P\{X \leqslant 2\} \cdot P\{Y=1\}\\
P\{X \leqslant 0\} &= F_X(0) = 0 \quad (\text{因为 } X \text{ 恒大于 } 0) \\
P\{X \leqslant 2\} &= F_X(2) = 1 - e^{-4 \times 2} = 1 - e^{-8}
\end{align*}
将其代回主式：
\begin{align*}
P\{X-Y \leqslant 1\} &= 0 \times \frac{1}{4} + (1 - e^{-8}) \times \frac{3}{4} = \frac{3}{4}(1 - e^{-8})
\end{align*}
\textbf{（2）求解 $\mathbf{P\{XY \leqslant 2\}}$}:采用相同的全概率分解逻辑，按 $Y$ 的取值将事件 $\{XY \leqslant 2\}$ 拆分：
\begin{align*}
P\{XY \leqslant 2\} &= P\big( \{XY \leqslant 2\} \cap \{Y=-1\} \big) + P\big( \{XY \leqslant 2\} \cap \{Y=1\} \big) \\
&= P\big( \{X(-1) \leqslant 2\} \cap \{Y=-1\} \big) + P\big( \{X(1) \leqslant 2\} \cap \{Y=1\} \big) \\
&= P\big( \{-X \leqslant 2\} \cap \{Y=-1\} \big) + P\big( \{X \leqslant 2\} \cap \{Y=1\} \big)
\end{align*}
由于 $X, Y$ 独立，原式转化为：
\begin{align*}
P\{XY \leqslant 2\} &= P\{-X \leqslant 2\} \cdot P\{Y=-1\} + P\{X \leqslant 2\} \cdot P\{Y=1\} \\
&= P\{X \geqslant -2\} \cdot \frac{1}{4} + P\{X \leqslant 2\} \cdot \frac{3}{4}
\end{align*}
由于 $X \sim E(4)$，其取值范围严格为 $X > 0$。因此，事件 $\{X \geqslant -2\}$ 是一个\textbf{必然事件}（即 $X$ 的取值永远大于 0，当然也就大于 -2），故：
\(
P\{X \geqslant -2\} = 1
\)
再代入前面算出的 $P\{X \leqslant 2\} = 1 - e^{-8}$，得到：
\begin{align*}
P\{XY \leqslant 2\} &= 1 \times \frac{1}{4} + (1 - e^{-8}) \times \frac{3}{4} = \frac{1}{4} + \frac{3}{4} - \frac{3}{4} e^{-8} = 1 - \frac{3}{4} e^{-8}
\end{align*}
\end{solution}

\subsection{边缘概率密度}

\begin{definition}[边缘概率密度，分布函数] \label{def:marginal-pdf}
设 $(X,Y) \sim f(x,y)$，则分别称
\[
f_X(x) = \int_{-\infty}^{+\infty}f(x,y)\,\mathrm{d}y, \quad f_Y(y) = \int_{-\infty}^{+\infty}f(x,y)\,\mathrm{d}x
\]
为 $(X,Y)$ 关于 $X$ 和关于 $Y$ 的\textbf{边缘概率密度}。相应的边缘分布函数分别为
\[
F_X(x) = \int_{-\infty}^{x}f_X(t)\,\mathrm{d}t = \int_{-\infty}^{x}\left[\int_{-\infty}^{+\infty}f(t,y)\,\mathrm{d}y\right]\mathrm{d}t,
\]
\[
F_Y(y) = \int_{-\infty}^{y}f_Y(t)\,\mathrm{d}t = \int_{-\infty}^{y}\left[\int_{-\infty}^{+\infty}f(x,t)\,\mathrm{d}x\right]\mathrm{d}t.
\]
\end{definition}

\begin{exercise}[两种边缘密度计算] \label{exer:3-c10}
求下列联合密度的边缘概率密度：\\
（1）$f(x,y) = \begin{cases} e^{-y}, & 0 < x < y, \\ 0, & \text{其他}; \end{cases}$~~（2）$f_2(x,y) = \begin{cases} 5(x^2+y), & 0 < y < 1-x^2, \\ 0, & \text{其他}. \end{cases}$
\end{exercise}
\begin{solution}
(1)
\[
f_X(x) = \int_x^{+\infty}e^{-y}\,\mathrm{d}y = e^{-x}\ (x > 0), \qquad f_Y(y) = \int_0^y e^{-y}\,\mathrm{d}x = ye^{-y}\ (y > 0).
\]
即 $X \sim E(1)$，$Y \sim \Gamma(2,1)$。\\
(2) $f_2(x,y) \neq 0$ 要求 $0 < y < 1-x^2$，即 $|x| < 1$。
\[
f_X(x) = \int_0^{1-x^2}5(x^2+y)\,\mathrm{d}y = 5\left[x^2y + \frac{y^2}{2}\right]_0^{1-x^2} = 5\left[x^2(1-x^2) + \frac{(1-x^2)^2}{2}\right] = \frac{5}{2}(1-x^4),\quad |x|<1.
\]
\[
f_Y(y) = \int_{-\sqrt{1-y}}^{\sqrt{1-y}}5(x^2+y)\,\mathrm{d}x = 10\sqrt{1-y}\cdot\frac{1+2y}{3} = \frac{10}{3}(1+2y)\sqrt{1-y},\quad 0<y<1.
\]
\end{solution}

\begin{exercise}\label{exer:3-4}
设 $(X,Y)$ 是二维随机变量，$X$ 的边缘概率密度为
\(
f_X(x) = \begin{cases} 3x^2, & 0 < x < 1, \\ 0, & \text{其他}, \end{cases}
\)
在给定 $X = x$（$0 < x < 1$）的条件下，$Y$ 的条件概率密度为
\[
f_{Y|X}(y|x) = \begin{cases} \dfrac{3y^2}{x^3}, & 0 < y < x, \\ 0, & \text{其他}. \end{cases}
\]
\begin{enumerate}
    \item[(1)] 求 $(X,Y)$ 的联合概率密度 $f(x,y)$；
    \item[(2)] 求 $Y$ 的边缘概率密度 $f_Y(y)$；
    \item[(3)] 求 $P\{X > 2Y\}$。
\end{enumerate}
\end{exercise}
\begin{solution}
(1) 由概率密度乘法公式：
\[
f(x,y) = f_{Y|X}(y|x)\,f_X(x) = \begin{cases} \dfrac{3y^2}{x^3} \cdot 3x^2 = \dfrac{9y^2}{x}, & 0 < y < x < 1, \\ 0, & \text{其他}. \end{cases}
\]
(2) 当 $0 < y \leqslant 1$ 时，$f_Y(y) = \displaystyle\int_y^1 \frac{9y^2}{x}\,\mathrm{d}x = 9y^2\bigl[-\ln y\bigr]_y^1 = -9y^2\ln y$；其余范围 $f_Y(y) = 0$。\\
(3) $\{X > 2Y\}$ 等价于 $\{Y < X/2\}$，积分区域为 $\{(x,y): 0 < x < 1,\ 0 < y < x/2\}$：
\[
P\{X > 2Y\} = \int_0^1\mathrm{d}x\int_0^{x/2}\frac{9y^2}{x}\,\mathrm{d}y = \int_0^1 \frac{3x^2}{8}\,\mathrm{d}x = \left[\frac{x^3}{8}\right]_0^1 = \frac{1}{8}.
\]
\end{solution}

\subsection{条件概率密度}

\begin{definition}[条件概率密度] \label{def:conditional-pdf}
设 $(X,Y) \sim f(x,y)$，若边缘概率密度 $f_X(x) > 0$，则称
\(
f_{Y|X}(y|x) = \frac{f(x,y)}{f_X(x)}
\)
为在 $\{X = x\}$ 条件下 $Y$ 的\textbf{条件概率密度}。

同理，若 $f_Y(y) > 0$，则称
\(
f_{X|Y}(x|y) = \frac{f(x,y)}{f_Y(y)}
\)
为在 $\{Y = y\}$ 条件下 $X$ 的\textbf{条件概率密度}。
\end{definition}

\begin{property}[概率密度乘法公式]\label{prop:pdf-multiplication}
若 $f_X(x) > 0$，$f_Y(y) > 0$，则有
\(
f(x,y) = f_X(x)\,f_{Y|X}(y|x) = f_Y(y)\,f_{X|Y}(x|y).
\)
\end{property}


\begin{exercise}[判断独立性与条件概率] \label{exer:3-5d}
设随机变量 $X, Y$ 的联合概率密度函数为
\[
f(x,y) = \begin{cases} e^{-y}, & x > 0,\ y > x, \\ 0, & \text{其他}, \end{cases}
\]
(1) $X, Y$ 是否相互独立？(2) 求 $P\{X > 2 \mid Y < 4\}$。
\end{exercise}
\begin{solution}
(1) 求边缘概率密度：
\[
f_X(x) = \int_x^{+\infty} e^{-y}\,\mathrm{d}y = e^{-x}\ (x > 0), \qquad f_Y(y) = \int_0^y e^{-y}\,\mathrm{d}x = ye^{-y}\ (y > 0).
\]
由于 $f_X(x)\,f_Y(y) = e^{-x} \cdot ye^{-y} = ye^{-(x+y)} \neq e^{-y} = f(x,y)$（当 $y > x > 0$ 时），故 $X$ 与 $Y$ 不相互独立。

(2)
\[
P\{X > 2,\ Y < 4\} = \int_2^4\mathrm{d}y\int_2^y e^{-y}\,\mathrm{d}x = \int_2^4(y-2)e^{-y}\,\mathrm{d}y = e^{-2} - 3e^{-4},
\]
\(
P\{Y < 4\} = \int_0^4 ye^{-y}\,\mathrm{d}y = 1 - 5e^{-4},
\)
故
\(
P\{X > 2 \mid Y < 4\} = \frac{P\{X > 2,\ Y < 4\}}{P\{Y < 4\}} = \frac{e^{-2} - 3e^{-4}}{1 - 5e^{-4}}.
\)
\end{solution}


\begin{definition}[条件分布函数]\label{def:conditional-df} 
在 $\{X=x\}$ 条件下 $Y$ 的\textbf{条件分布函数}为
\[
F_{Y|X}(y|x) = \int_{-\infty}^{y}f_{Y|X}(v|x)\,\mathrm{d}v = \int_{-\infty}^{y}\frac{f(x,v)}{f_X(x)}\,\mathrm{d}v.
\]
同理，在 $\{Y=y\}$ 条件下 $X$ 的\textbf{条件分布函数}为
\[
F_{X|Y}(x|y) = \int_{-\infty}^{x}f_{X|Y}(u|y)\,\mathrm{d}u = \int_{-\infty}^{x}\frac{f(u,y)}{f_Y(y)}\,\mathrm{d}u.
\]
\end{definition}

\begin{property}[区域概率公式]\label{prop:region-prob}
对于联合分布函数 $F(x,y)$ 及边缘分布函数 $F_X(x), F_Y(y)$，有
\(
P\{X > x_0, Y > y_0\} = 1 - F_X(x_0) - F_Y(y_0) + F(x_0, y_0).
\)
\textbf{推导}：将全空间分为四个区域，
\[
P\{X > x_0, Y > y_0\} = 1 - P\{X \leqslant x_0\} - P\{Y \leqslant y_0\} + P\{X \leqslant x_0, Y \leqslant y_0\}.
\]
注意：只有当 $X$ 与 $Y$ 独立时，$P\{X > x_0, Y > y_0\} = [1-F_X(x_0)][1-F_Y(y_0)]$。
\end{property}

\begin{exercise}[混合条件分布求 $F_Y$] \label{exer:3-c17}
$P\{X=1\} = P\{X=2\} = 1/2$，在 $X=i$ 的条件下 $Y \sim U(0,i)$。求 $Y$ 的分布函数 $F_Y(y)$ 和密度 $f_Y(y)$。
\end{exercise}
\begin{solution}
由全概率公式：$F_Y(y) = P\{X=1\}\,F_{Y|X=1}(y) + P\{X=2\}\,F_{Y|X=2}(y) = \frac{1}{2}F_1(y) + \frac{1}{2}F_2(y)$。\\
$Y|X=1 \sim U(0,1)$：$F_1(y) = \begin{cases} 0, & y \leqslant 0, \\ y, & 0 < y < 1, \\ 1, & y \geqslant 1. \end{cases}$，
$Y|X=2 \sim U(0,2)$：$F_2(y) = \begin{cases} 0, & y \leqslant 0, \\ y/2, & 0 < y < 2, \\ 1, & y \geqslant 2. \end{cases}$
\begin{itemize}
    \item $y \leqslant 0$：$F_Y(y) = 0$。
    \item $0 < y \leqslant 1$：$F_Y(y) = \dfrac{1}{2}y + \dfrac{1}{2}\cdot\dfrac{y}{2} = \dfrac{3y}{4}$。
    \item $1 < y \leqslant 2$：$F_Y(y) = \dfrac{1}{2}\cdot 1 + \dfrac{1}{2}\cdot\dfrac{y}{2} = \dfrac{1}{2} + \dfrac{y}{4}$。
    \item $y \geqslant 2$：$F_Y(y) = 1$。
\end{itemize}
\[
F_Y(y) = \begin{cases} 0, & y \leqslant 0, \\ \dfrac{3y}{4}, & 0 < y \leqslant 1, \\ \dfrac{1}{2} + \dfrac{y}{4}, & 1 < y \leqslant 2, \\ 1, & y \geqslant 2. \end{cases}
\qquad
f_Y(y) = \begin{cases} \dfrac{3}{4}, & 0 < y \leqslant 1, \\ \dfrac{1}{4}, & 1 < y \leqslant 2, \\ 0, & \text{其他}. \end{cases}
\]
\end{solution}


\section{常见的二维分布}

\subsection{二维均匀分布}

\begin{definition}[二维均匀分布] \label{def:2d-uniform}
设 $D$ 为 $xOy$ 平面上的有界区域，其面积为 $S_D$。如果 $(X,Y)$ 的联合概率密度为
\[
f(x,y) = \begin{cases} \dfrac{1}{S_D}, & (x,y) \in D, \\ 0, & \text{其他}, \end{cases}
\]
则称 $(X,Y)$ 在区域 $D$ 上服从\textbf{二维均匀分布}，记为 $(X,Y) \sim U(D)$。
\end{definition}

二维均匀分布是区间上均匀分布在二维情形的推广。若 $(X,Y) \sim U(D)$，则对 $D$ 的任何子区域 $G$（其面积为 $S_G$），有
\(
P\{(X,Y) \in G\} = \frac{S_G}{S_D}.
\)


\begin{exercise}[三角形均匀分布求概率] \label{exer:3-c14}
\[f(x,y) = \begin{cases} 2-x-y, & 0 < x < 1,\ 0 < y < 1, \\ 0, & \text{其他}. \end{cases}\]
求：(1) $P\{X > 2Y\}$；(2) $Z = X+Y$ 的密度。
\end{exercise}
\begin{solution}
验证归一性：$\int_0^1\int_0^1(2-x-y)\,\mathrm{d}y\,\mathrm{d}x = 1$。\\
(1) $\{X > 2Y\}$ 等价于 $\{0 < y < x/2\}$（$0 < x < 1$）：
\[
P\{X > 2Y\} = \int_0^1\int_0^{x/2}(2-x-y)\,\mathrm{d}y\,\mathrm{d}x = \int_0^1\left[2y-xy-\frac{y^2}{2}\right]_0^{x/2}\mathrm{d}x
\]
\[
= \int_0^1\left(x - \frac{x^2}{2} - \frac{x^2}{8}\right)\mathrm{d}x = \int_0^1\left(x - \frac{5x^2}{8}\right)\mathrm{d}x = \left[\frac{x^2}{2} - \frac{5x^3}{24}\right]_0^1 = \frac{7}{24}.
\]
(2) 积分：
\[
F_Z(z)=P(X+Y\le z)
=\iint_{x+y\le z} f_{X,Y}(x,y)\,\mathrm{d}x\,\mathrm{d}y.
\]
由于联合密度只在 $0<x<1,\ 0<y<1$ 上非零，所以要按直线 $x+y=z$ 穿过单位正方形的位置分段。若 $0<z<1$，则积分区域为
\[
0<x<z,\qquad 0<y<z-x.
\]
于是
\[
F_Z(z)=\int_0^z\int_0^{z-x}(2-x-y)\,\mathrm{d}y\,\mathrm{d}x.
\]
对 $z$ 求导（变上限积分求导公式$\displaystyle\frac{\mathrm d}{\mathrm dz}\int_0^z G(x,z)\,\mathrm dx=\int_0^z \frac{\partial G(x,z)}{\partial z}\,\mathrm dx+G(z,z)$），得到
\begin{align*}
&\frac{\mathrm{d}}{\mathrm{d}z}
\int_0^z\left[\int_0^{z-x}(2-x-y)\,\mathrm{d}y\right]\mathrm{d}x=\int_0^z
\frac{\partial}{\partial z}\left[\int_0^{z-x}(2-x-y)\,\mathrm{d}y\right]\mathrm{d}x
+\left[\int_0^{z-x}(2-x-y)\,\mathrm{d}y\right]_{x=z}\\
&=\int_0^z (2-x-(z-x))\,\mathrm{d}x+0=\int_0^z(2-z)\,\mathrm{d}x=z(2-z),\qquad 0<z<1.
\end{align*}
若 $1\leqslant z<2$，则积分区域为
\[
z-1<x<1,\qquad 0<y<z-x.
\]
同理由边界线 $x+y=z$ 上的密度积分得
\[
f_Z(z)=\int_{z-1}^1 f_{X,Y}(x,z-x)\,\mathrm{d}x
=\int_{z-1}^1(2-z)\,\mathrm{d}x
=(2-z)^2,\qquad 1\leqslant z<2.
\]
综上：
\[
f_Z(z) = \begin{cases} z(2-z), & 0 < z < 1, \\ (2-z)^2, & 1 \leqslant z < 2, \\ 0, & \text{其他}. \end{cases}
\]
\end{solution}


\begin{exercise}[三角形区域上的均匀分布] \label{exer:3-13}
已知二维随机变量 $(X,Y)$ 在区域 $G$ 上服从均匀分布，$G$ 由直线 $x-y=0$、$x+y=2$ 与 $y=0$ 围成。求：\\
(1)边缘概率密度 $f_X(x)$ 和 $f_Y(y)$；~~(2)条件概率密度 $f_{X|Y}(x|y)$。
\end{exercise}
\begin{solution}
$(X,Y)$ 的正概率密度区域 $G$ 是以 $(0,0)$、$(2,0)$、$(1,1)$ 为顶点的三角形，面积 $S_G = \dfrac{1}{2} \times 2 \times 1 = 1$。
由于面积恰好为 $1$，联合概率密度为
\[
f(x,y) = \begin{cases} 1, & 0 \leqslant y \leqslant 1,\ y \leqslant x \leqslant 2-y, \\ 0, & \text{其他}. \end{cases}
\]
(1) 边缘概率密度。$f_X(x)$：固定 $x$，对 $y$ 积分。
\begin{itemize}
    \item 当 $0 < x < 1$ 时，$y$ 从 $0$ 到 $x$：$f_X(x) = \displaystyle\int_0^x 1\,\mathrm{d}y = x$；
    \item 当 $1 < x < 2$ 时，$y$ 从 $0$ 到 $2-x$：$f_X(x) = \displaystyle\int_0^{2-x} 1\,\mathrm{d}y = 2-x$。
\end{itemize}
\[
f_X(x) = \begin{cases} x, & 0 < x < 1, \\ 2-x, & 1 < x < 2, \\ 0, & \text{其他}. \end{cases}
\]
$f_Y(y)$：固定 $y$，对 $x$ 积分。
\begin{itemize}
    \item 当 $0 < y < 1$ 时，$x$ 从 $y$ 到 $2-y$：$f_Y(y) = \displaystyle\int_y^{2-y} 1\,\mathrm{d}x = 2-2y$。
\end{itemize}
\(
f_Y(y) = \begin{cases} 2-2y, & 0 < y < 1, \\ 0, & \text{其他}. \end{cases}
\)
(2) 在 $0 < y < 1$ 的条件下：
\(
f_{X|Y}(x|y) = \frac{f(x,y)}{f_Y(y)} = \frac{1}{2-2y}, \quad y \leqslant x \leqslant 2-y.
\)
即
\[
f_{X|Y}(x|y) = \begin{cases} \dfrac{1}{2-2y}, & y \leqslant x \leqslant 2-y, \\ 0, & \text{其他}. \end{cases}
\]
\end{solution}
\begin{note}
求边缘概率密度的技巧（口诀）：\textbf{求谁不积谁，不积先定限，限内画条线，先交写下限，后交写上限。}即对 $f(x,y)$ 求关于 $y$ 的积分得 $f_X(x)$，对 $f(x,y)$ 求关于 $x$ 的积分得 $f_Y(y)$。定积分限时，画出正概率密度区域的图形，用平行于积分轴的直线穿过区域，先交到的位置为下限，后交到的位置为上限。
\end{note}

\begin{exercise}[条件均匀分布] \label{exer:3-14}
已知随机变量 $X$ 在区间 $[0,1]$ 上服从均匀分布，在 $X=x$（$0<x<1$）的条件下，随机变量 $Y$ 在区间 $[0,x]$ 上服从均匀分布。求：\\
(1) $(X,Y)$ 的联合概率密度；~~(2)$Y$ 的概率密度；~~(3)概率 $P\{X+Y>1\}$。
\end{exercise}
\begin{solution}
(1) 由题设，$X$ 的边缘概率密度为
\(
f_X(x) = \begin{cases} 1, & 0 < x < 1, \\ 0, & \text{其他}. \end{cases}
\)
在 $X=x$（$0<x<1$）的条件下，$Y$ 在 $[0,x]$ 上的条件概率密度为
\(
f_{Y|X}(y|x) = \begin{cases} \dfrac{1}{x}, & 0 < y < x, \\ 0, & \text{其他}. \end{cases}
\)
由乘法公式 $f(x,y) = f_X(x)\,f_{Y|X}(y|x)$，当 $0<y<x<1$ 时 $f(x,y) = 1\cdot\dfrac{1}{x} = \dfrac{1}{x}$。故联合概率密度为
\(
f(x,y) = \begin{cases} \dfrac{1}{x}, & 0 < y < x < 1, \\ 0, & \text{其他}. \end{cases}
\)
(2) 当 $0<y<1$ 时，固定 $y$ 对 $x$ 积分（$x$ 从 $y$ 到 $1$）：
\[
f_Y(y) = \int_y^1\frac{1}{x}\,\mathrm{d}x = \bigl[-\ln y\bigr]_y^1 = -\ln y.
\]
当 $y \leqslant 0$ 或 $y \geqslant 1$ 时，$f_Y(y) = 0$。因此
\[
f_Y(y) = \begin{cases} -\ln y, & 0 < y < 1, \\ 0, & \text{其他}. \end{cases}
\]
(3) $\{X+Y>1\}$ 的积分区域为 $\{(x,y): 0<y<x<1,\ x+y>1\}$，即 $x>\dfrac{1}{2}$ 且 $1-x<y<x$：
\[
P\{X+Y>1\} = \int_{1/2}^1\mathrm{d}x\int_{1-x}^x\frac{1}{x}\,\mathrm{d}y = \int_{1/2}^1\frac{x-(1-x)}{x}\,\mathrm{d}x = \int_{1/2}^1\frac{2x-1}{x}\,\mathrm{d}x = \int_{1/2}^1\left(2-\frac{1}{x}\right)\mathrm{d}x
\]
\(
= \bigl[2x - \ln x\bigr]_{1/2}^1 = (2-0) - (1-\ln 2) = 1 - \ln 2.
\)
\end{solution}
\begin{note}
第 (1) 问中，乘法公式 $f(x,y) = f_X(x)\,f_{Y|X}(y|x)$ 仅在 $f_X(x)>0$ 时成立，因此最终结果需要加上 $f_X(x)=0$ 的部分，在分段表达式的``其他''中自然已包含。
\end{note}


\subsection{二维正态分布}

\begin{definition}[二维正态分布] \label{def:2d-normal}
如果 $(X,Y)$ 的联合概率密度为
\[
f(x,y) = \frac{1}{2\pi\sigma_1\sigma_2\sqrt{1-\rho^2}}\exp\left\{-\frac{1}{2(1-\rho^2)}\left[\frac{(x-\mu_1)^2}{\sigma_1^2} - 2\rho\frac{(x-\mu_1)(y-\mu_2)}{\sigma_1\sigma_2} + \frac{(y-\mu_2)^2}{\sigma_2^2}\right]\right\},
\]
其中 $\mu_1, \mu_2 \in \mathbb{R}$，$\sigma_1 > 0$，$\sigma_2 > 0$，$-1 < \rho < 1$，则称 $(X,Y)$ 服从参数为 $\mu_1, \mu_2, \sigma_1^2, \sigma_2^2, \rho$ 的\textbf{二维正态分布}，记为 $(X,Y) \sim N(\mu_1, \mu_2, \sigma_1^2, \sigma_2^2; \rho)$。
\end{definition}

\begin{property}\label{prop:2d-normal}
二维正态分布具有以下重要性质：
\begin{enumerate}
    \item \textbf{边缘分布仍为正态分布}：若 $(X,Y) \sim N(\mu_1, \mu_2, \sigma_1^2, \sigma_2^2; \rho)$，则
    \[
    X \sim N(\mu_1, \sigma_1^2), \quad Y \sim N(\mu_2, \sigma_2^2).
    \]
    \item \textbf{独立性的充要条件}：若 $(X,Y)$ 服从二维正态分布，则 $X$ 与 $Y$ 相互独立的充分必要条件是 $\rho = 0$。
    \item \textbf{独立正态变量联合分布}：若 $X \sim N(\mu_1, \sigma_1^2)$，$Y \sim N(\mu_2, \sigma_2^2)$，且 $X$ 与 $Y$ 相互独立，则
    \(
    (X,Y) \sim N(\mu_1, \mu_2, \sigma_1^2, \sigma_2^2; 0).
    \)
\end{enumerate}
\end{property}

\begin{note}
需要注意以下几点：
\begin{enumerate}
    \item 边缘分布为正态分布\textbf{不能}推出联合分布为二维正态分布。例如，设 $(X,Y)$ 的联合密度为
    \(
    f(x,y) = \frac{1}{2\pi}e^{-(x^2+y^2)/2}\bigl[1 + \sin x\sin y\bigr],
    \)
    则 $X$ 和 $Y$ 都服从标准正态分布，但 $(X,Y)$ 并不服从二维正态分布。
    \item 性质 (2) 仅对二维正态分布成立。一般地，由 $X$ 与 $Y$ 的边缘分布和 $\rho_{XY} = 0$ 不能推出 $X$ 与 $Y$ 独立。
    \item 二维正态分布的条件分布仍为正态分布：在 $X = x$ 的条件下，$Y$ 的条件分布为
    \[
    Y \mid X=x \sim N\!\left(\mu_2 + \rho\frac{\sigma_2}{\sigma_1}(x-\mu_1),\ \sigma_2^2(1-\rho^2)\right).
    \]
\end{enumerate}
\end{note}

\begin{exercise}[二维分布的性质判断] \label{exer:3-c5b}
下列叙述中正确的是（\quad）
\begin{enumerate}
    \item[(1)] 二维正态分布的边缘分布是正态分布，条件分布也是正态分布；
    \item[(2)] 二维正态分布的边缘分布是正态分布，但条件分布不是正态分布；
    \item[(3)] 二维均匀分布的边缘分布是均匀分布，条件分布是均匀分布；
    \item[(4)] 二维均匀分布的边缘分布不一定是均匀分布，条件分布也不一定为均匀分布。
\end{enumerate}
A. (1)(3) \quad B. (2)(3) \quad C. (1)(4) \quad D. (2)(4)
\end{exercise}
\begin{solution}
\textbf{答案选 C。}
(1) 正确：二维正态分布的边缘分布和条件分布均为正态分布。\\
(4) 正确：二维均匀分布的边缘不一定均匀。例如
\[
f(x,y) = \begin{cases} 2, & 0 < x < 1,\ 0 < y < 1-x, \\ 0, & \text{其他}, \end{cases}
\]
则 $f_X(x) = 2(1-x)$（$0<x<1$），$f_Y(y) = 2(1-y)$（$0<y<1$），均非均匀分布；条件密度 $f_{Y|X}(y|x) = 1/(1-x)$（$0<y<1-x$），在 $[0,1]$ 上也不均匀。
\end{solution}

\begin{exercise}\label{exer:3-5}
设随机变量 $X$ 与 $Y$ 相互独立，且均服从参数为 $1$ 的指数分布，即 $X, Y \sim E(1)$，求 $P\{X < Y\}$。
\end{exercise}

\begin{solution}
由独立性，联合概率密度为
\[
f(x,y) = f_X(x)\,f_Y(y) = \begin{cases} e^{-(x+y)}, & x > 0,\ y > 0, \\ 0, & \text{其他}. \end{cases}
\]
\[
P\{X < Y\} = \iint_{x < y}f(x,y)\,\mathrm{d}x\,\mathrm{d}y = \int_0^{+\infty}\mathrm{d}y\int_0^y e^{-(x+y)}\,\mathrm{d}x = \int_0^{+\infty}e^{-y}(1-e^{-y})\,\mathrm{d}y
\]
\[
= \int_0^{+\infty}(e^{-y} - e^{-2y})\,\mathrm{d}y = \left[-e^{-y} + \frac{1}{2}e^{-2y}\right]_0^{+\infty} = 1 - \frac{1}{2} = \frac{1}{2}.
\]
\end{solution}

\begin{exercise}[条件正态联合密度] \label{exer:3-5b}
设随机变量 $X \sim N(0,1)$，在 $X = x$ 的条件下，随机变量 $Y \sim N(x, 1)$，求 $(X,Y)$ 的概率密度 $f(x,y)$。
\end{exercise}

\begin{solution}
由题意知
\[
f_X(x) = \frac{1}{\sqrt{2\pi}}e^{-x^2/2}, \quad f_{Y|X}(y|x) = \frac{1}{\sqrt{2\pi}}e^{-(y-x)^2/2},
\]
故由乘法公式
\(
f(x,y) = f_X(x)\,f_{Y|X}(y|x) = \frac{1}{2\pi}e^{-\frac{x^2 + (y-x)^2}{2}} = \frac{1}{2\pi}e^{-\frac{2x^2 - 2xy + y^2}{2}}.
\)
\end{solution}


\section{二维随机变量函数的分布}

\subsection{离散型随机变量函数的分布}
\begin{note}
若 $X \sim B(n_1, p)$，$Y \sim B(n_2, p)$，且 $X$ 与 $Y$ 独立，则 $X + Y \sim B(n_1 + n_2, p)$。类似地，若 $X \sim P(\lambda_1)$，$Y \sim P(\lambda_2)$，且 $X$ 与 $Y$ 独立，则 $X + Y \sim P(\lambda_1 + \lambda_2)$。
\end{note}

\begin{exercise}[和与差的分布] \label{exer:3-6b}
将两封信投入 3 个信箱，设 $X_1, X_2$ 分别表示第一个和第二个信箱投进的信的数量，求随机变量 $Y_1 = X_1 + X_2$ 与 $Y_2 = X_1 - X_2$ 的分布。
\end{exercise}

\begin{solution}
 $(X_1, X_2)$ 的联合分布律为
\begin{center}
\begin{tabular}{c|ccc|c}
\toprule
$X_1 \backslash X_2$ & $0$ & $1$ & $2$ & $P\{X_1=x_i\}$ \\
\midrule
$0$ & $\frac19$ & $\frac29$ & $\frac19$ & $\frac49$ \\
$1$ & $\frac29$ & $\frac29$ & $0$ & $\frac49$ \\
$2$ & $\frac19$ & $0$ & $0$ & $\frac19$ \\
\midrule
$P\{X_2=x_j\}$ & $\frac49$ & $\frac49$ & $\frac19$ & $1$ \\
\bottomrule
\end{tabular}
\end{center}
下面求 $Y_1,Y_2$ 的分布时，就是把表中满足相应关系的格子概率相加。

$Y_1 = X_1 + X_2$ 的可能取值为 $0, 1, 2$：
\[
P\{Y_1 = 0\} = p_{00} = \frac{1}{9}, \quad P\{Y_1 = 1\} = p_{01} + p_{10} = \frac{4}{9}, \quad P\{Y_1 = 2\} = p_{02} + p_{11} + p_{20} = \frac{4}{9}.
\]

$Y_2 = X_1 - X_2$ 的可能取值为 $-2, -1, 0, 1, 2$：
\(
P\{Y_2 = -2\} = p_{02} = \frac{1}{9}, \quad P\{Y_2 = -1\} = p_{01} = \frac{2}{9}, \quad P\{Y_2 = 0\} = p_{00} + p_{11} = \frac{3}{9} = \frac{1}{3},
\)
\(
P\{Y_2 = 1\} = p_{10} = \frac{2}{9}, \quad P\{Y_2 = 2\} = p_{20} = \frac{1}{9}.
\)

综上：
\[
Y_1 \sim \begin{pmatrix} 0 & 1 & 2 \\ 1/9 & 4/9 & 4/9 \end{pmatrix}, \qquad Y_2 \sim \begin{pmatrix} -2 & -1 & 0 & 1 & 2 \\ 1/9 & 2/9 & 1/3 & 2/9 & 1/9 \end{pmatrix}.
\]
\end{solution}

\subsection{连续型随机变量函数的分布}

\subsubsection{分布函数法}

\begin{property}[分布函数法]\label{prop:df-method}
求 $Z = g(X,Y)$ 的分布，步骤如下：
\begin{enumerate}
    \item 先求 $Z$ 的分布函数：
    \[
    F_Z(z) = P\{Z \leqslant z\} = P\{g(X,Y) \leqslant z\} = \iint_{\{(x,y): g(x,y) \leqslant z\}}f(x,y)\,\mathrm{d}x\,\mathrm{d}y;
    \]
    \item 再对 $F_Z(z)$ 求导得概率密度：$f_Z(z) = F_Z'(z)$。有时前一步的二重积分直接算出 \(F_Z(z)\) 比较麻烦，但最后只要求密度，可以先保留积分形式，再用变上限积分求导。常见情形如
    \[
    F_Z(z)=\int_{a(z)}^{b(z)}\int_{c(x)}^{d(x,z)}f(x,y)\,\mathrm{d}y\,\mathrm{d}x,
    \]
    对 \(z\) 求导时，内层上限 \(d(x,z)\) 会给出边界上的密度：
    \[
    \begin{aligned}
    \frac{\mathrm{d}}{\mathrm{d}z}\int_{c(x)}^{d(x,z)}f(x,y)\,\mathrm{d}y=\frac{\mathrm{d}}{\mathrm{d}z}
    \left(\int_{c(x)}^t f(x,y)\,\mathrm{d}y\right)\bigg|_{t=d(x,z)}=f(x,d(x,z))\frac{\partial d(x,z)}{\partial z}.
    \end{aligned}
    \]
    同时还要看外层上下限 \(a(z),b(z)\) 是否随 \(z\) 变化；若端点处内层积分为 \(0\)，对应的端点项就会消失。比如 \(Z=X+Y\) 且积分区域可写成
    \[
    F_Z(z)=\int_0^z\int_0^{z-x}f(x,y)\,\mathrm{d}y\,\mathrm{d}x,
    \]
    则
    \[
    f_Z(z)=F_Z'(z)=\int_0^z f(x,z-x)\,\mathrm{d}x,
    \]
    这就是把密度沿直线边界 \(x+y=z\) 积起来，比先完整求出 \(F_Z(z)\) 再求导更简洁。
\end{enumerate}
\end{property}

\begin{property}\label{prop:diff-prod-quot-dist}
设 $(X,Y) \sim f(x,y)$，则
\begin{enumerate}
    \item \textbf{和的分布}（$Z = X + Y$）：
    \[
    f_Z(z) = \int_{-\infty}^{+\infty}f(x, z-x)\,\mathrm{d}x = \int_{-\infty}^{+\infty}f(z-y, y)\,\mathrm{d}y.
    \]
    当 $X$ 与 $Y$ 独立时，$f_Z(z) = \displaystyle\int_{-\infty}^{+\infty}f_X(x)\,f_Y(z-x)\,\mathrm{d}x$。
    \item \textbf{差的分布}（$Z = X - Y$）：
    \[
    f_Z(z) = \int_{-\infty}^{+\infty}f(x, x-z)\,\mathrm{d}x = \int_{-\infty}^{+\infty}f(z+y, y)\,\mathrm{d}y.
    \]
    当 $X$ 与 $Y$ 独立时，$f_Z(z) = \displaystyle\int_{-\infty}^{+\infty}f_X(x)\,f_Y(x-z)\,\mathrm{d}x$。

    \item \textbf{积的分布}（$Z = XY$）：
    \[
    f_Z(z) = \int_{-\infty}^{+\infty}\frac{1}{|x|}\,f\!\left(x, \frac{z}{x}\right)\,\mathrm{d}x = \int_{-\infty}^{+\infty}\frac{1}{|y|}\,f\!\left(\frac{z}{y}, y\right)\,\mathrm{d}y.
    \]
    当 $X$ 与 $Y$ 独立时，$f_Z(z) = \displaystyle\int_{-\infty}^{+\infty}\frac{1}{|x|}\,f_X(x)\,f_Y\!\left(\frac{z}{x}\right)\,\mathrm{d}x$。

    \item \textbf{商的分布}（$Z = X / Y$）：
    \(
    f_Z(z) = \int_{-\infty}^{+\infty}|y|\,f(yz, y)\,\mathrm{d}y.
    \)
    当 $X$ 与 $Y$ 独立时，$f_Z(z) = \displaystyle\int_{-\infty}^{+\infty}|y|\,f_X(yz)\,f_Y(y)\,\mathrm{d}y$。
\end{enumerate}
\end{property}


\begin{proof}
\begin{enumerate}
    \item 差的分布 $Z=X-Y$ 的证明：
    先求 $Z$ 的分布函数 $F_Z(z) = P(X-Y \leqslant z) = P(X \leqslant z+Y)$。
    由联合密度积分可得：
    \[
    F_Z(z) = \iint_{x-y \leqslant z} f(x,y)\,\mathrm{d}x\,\mathrm{d}y = \int_{-\infty}^{+\infty} \left( \int_{-\infty}^{z+y} f(x,y)\,\mathrm{d}x \right) \mathrm{d}y.
    \]
    对内层积分做变量代换，令 $u = x-y$（即固定 $y$，有 $x = u+y, \mathrm{d}x = \mathrm{d}u$），当 $x$ 从 $-\infty$ 积到 $z+y$ 时，$u$ 从 $-\infty$ 积到 $z$：
    \[
    F_Z(z) = \int_{-\infty}^{+\infty} \left( \int_{-\infty}^{z} f(u+y, y)\,\mathrm{d}u \right) \mathrm{d}y = \int_{-\infty}^{z} \left[ \int_{-\infty}^{+\infty} f(u+y, y)\,\mathrm{d}y \right] \mathrm{d}u.
    \]
    两边对 $z$ 求导，即得 $f_Z(z) = \displaystyle\int_{-\infty}^{+\infty}f(z+y, y)\,\mathrm{d}y$。
    同理，若按 $\mathrm{d}x$ 在外层积分，可证得另一种形式 $f_Z(z) = \displaystyle\int_{-\infty}^{+\infty}f(x, x-z)\,\mathrm{d}x$。当 $X, Y$ 独立时，$f(x,y)=f_X(x)f_Y(y)$，代入即得结论。
    \item 积的分布 $Z=XY$ 的证明：
    求 $Z$ 的分布函数 $F_Z(z) = P(XY \leqslant z)$，积分区域按 $x$ 的正负分为两部分：
    \begin{align*}
        F_Z(z) &= \iint_{xy \leqslant z} f(x,y)\,\mathrm{d}x\,\mathrm{d}y \\
        &= \int_{0}^{+\infty} \left( \int_{-\infty}^{\frac{z}{x}} f(x,y)\,\mathrm{d}y \right) \mathrm{d}x + \int_{-\infty}^{0} \left( \int_{\frac{z}{x}}^{+\infty} f(x,y)\,\mathrm{d}y \right) \mathrm{d}x.
    \end{align*}
    对内层积分做代换，令 $u = xy$（固定 $x$，则 $y = \frac{u}{x}$，$\mathrm{d}y = \frac{1}{x}\mathrm{d}u$）。
    当 $x > 0$ 时，$y$ 的区间 $(-\infty, \frac{z}{x}]$ 对应 $u$ 的区间 $(-\infty, z]$；
    当 $x < 0$ 时，$y$ 的区间 $[\frac{z}{x}, +\infty)$ 对应 $u$ 的区间 $[z, -\infty)$，此时积分限反转带来负号，即 $\mathrm{d}y = \frac{1}{x}\mathrm{d}u = -\frac{1}{|x|}\mathrm{d}u$，吸收反转积分限的负号后，两部分合并为：
    \begin{align*}
    F_Z(z) &= \int_{0}^{+\infty} \left( \int_{-\infty}^{z} f\!\left(x,\frac{u}{x}\right)\frac{1}{|x|}\,\mathrm{d}u \right) \mathrm{d}x + \int_{-\infty}^{0} \left( \int_{-\infty}^{z} f\!\left(x,\frac{u}{x}\right)\frac{1}{|x|}\,\mathrm{d}u \right) \mathrm{d}x \\
    &= \int_{-\infty}^{z} \left[ \int_{-\infty}^{+\infty} \frac{1}{|x|}f\!\left(x,\frac{u}{x}\right)\,\mathrm{d}x \right] \mathrm{d}u.
    \end{align*}
    对 $z$ 求导即得 $f_Z(z) = \displaystyle\int_{-\infty}^{+\infty}\frac{1}{|x|}\,f\!\left(x, \frac{z}{x}\right)\,\mathrm{d}x$。独立时展开联合密度即得特例。
    \item \textbf{分布函数法}：
    求 $F_Z(z) = P(X/Y \leqslant z)$，按 $y$ 的正负分为两部分：
    \begin{align*}
        F_Z(z) &= P(X \leqslant zY, Y > 0) + P(X \geqslant zY, Y < 0) \\
        &= \int_{0}^{+\infty} \left( \int_{-\infty}^{zy} f(x,y)\,\mathrm{d}x \right) \mathrm{d}y + \int_{-\infty}^{0} \left( \int_{zy}^{+\infty} f(x,y)\,\mathrm{d}x \right) \mathrm{d}y.
    \end{align*}
    对内层积分做代换，令 $u = x/y$（固定 $y$，则 $x = uy$，$\mathrm{d}x = y\,\mathrm{d}u$）。
    当 $y > 0$ 时，$x$ 的区间 $(-\infty, zy]$ 对应 $u$ 的区间 $(-\infty, z]$；
    当 $y < 0$ 时，$x$ 的区间 $[zy, +\infty)$ 对应 $u$ 的区间 $[z, -\infty)$，此时 $y = -|y|$，积分限反转吸收负号后得：
    \begin{align*}
    F_Z(z)& = \int_{0}^{+\infty} \left( \int_{-\infty}^{z} f(uy, y)y\,\mathrm{d}u \right) \mathrm{d}y + \int_{-\infty}^{0} \left( \int_{-\infty}^{z} f(uy, y)(-y)\,\mathrm{d}u \right) \mathrm{d}y\\
    & = \int_{-\infty}^{z} \left[ \int_{-\infty}^{+\infty} |y|f(uy, y)\,\mathrm{d}y \right] \mathrm{d}u.
    \end{align*}
    对 $z$ 求导即得 $f_Z(z) = \displaystyle\int_{-\infty}^{+\infty}|y|\,f(yz, y)\,\mathrm{d}y$。独立时展开联合密度即得特例。
\end{enumerate}
\end{proof}

\begin{exercise}\label{exer:3-7}
设随机变量 $X$ 和 $Y$ 相互独立，$X \sim U(0, 1)$，$Y \sim E(1)$，求 $Z = \max\{X, Y\}$ 的分布函数 $F_Z(z)$。
\end{exercise}

\begin{solution}
$X$ 的分布函数为 $F_X(x) = x$（$0 < x < 1$），$Y$ 的分布函数为 $F_Y(y) = 1 - e^{-y}$（$y > 0$）。

对于独立的 $X$ 和 $Y$：
\[
\max\{X, Y\} \leqslant z \iff X \leqslant z \text{ 且 } Y \leqslant z,~~
F_Z(z) = P\{Z \leqslant z\} = P\{X \leqslant z\}\,P\{Y \leqslant z\} = F_X(z)\,F_Y(z).
\]

当 $z \leqslant 0$ 时，$F_Z(z) = 0$；

当 $0 < z < 1$ 时，$F_Z(z) = z(1 - e^{-z})$；

当 $z \geqslant 1$ 时，$F_Z(z) = 1 \cdot (1-e^{-z}) = 1 - e^{-z}$。

综上：
\[
F_Z(z) = \begin{cases} 0, & z \leqslant 0, \\ z(1 - e^{-z}), & 0 < z < 1, \\ 1 - e^{-z}, & z \geqslant 1. \end{cases}
\]
\end{solution}

\begin{exercise}\label{exer:src-2010-0630-a-07-2credit}
设 \(X,Y\) 是相互独立的随机变量，其概率密度分别为
\[
f_X(x)=
\begin{cases}
1, & 0\le x\le1,\\
0, & \text{其他},
\end{cases}
\qquad
f_Y(y)=
\begin{cases}
e^{-y}, & y>0,\\
0, & \text{其他}.
\end{cases}
\]
定义
\(
Z=g(X,Y)=
\begin{cases}
1, & X\le Y,\\
0, & X>Y.
\end{cases}
\)
求 \(EZ\) 与 \(DZ\)。
\end{exercise}
\begin{solution}
\(Z\) 是事件 \(\{X\le Y\}\) 的示性变量，因此
\begin{align*}
EZ
&=P(Z=1)=P(X\le Y)=\iint_{\{0\le x\le1,\ y\ge x\}} f_X(x)f_Y(y)\,\mathrm{d}y\,\mathrm{d}x\\
&=\int_0^1\int_x^{+\infty} e^{-y}\,\mathrm{d}y\,\mathrm{d}x=\int_0^1 e^{-x}\,\mathrm{d}x=1-e^{-1}.
\end{align*}
又 \(Z\sim B(1,1-e^{-1})\)，所以
\(
DZ=(1-e^{-1})e^{-1}.
\)
\end{solution}

\begin{exercise}\label{exer:src-2012-a-07-2credit}
设随机变量 \(X\) 服从参数为 \(\lambda\) 的指数分布，其概率密度为
\[
f_X(x)=
\begin{cases}
\lambda e^{-\lambda x}, & x>0,\\
0, & x\le0.
\end{cases}
\]
令
\(
Y=\max\{X^2,4\}.
\)
求 \(Y\) 的分布函数。
\end{exercise}
\begin{solution}
因为 \(Y\ge4\)，故当 \(y<4\) 时，
\(
F_Y(y)=P\{Y\le y\}=0.
\)
当 \(y\ge4\) 时，
\(
\{Y\le y\}=\{\max(X^2,4)\le y\}=\{X^2\le y,\ 4\le y\}=\{X^2\le y\}.
\)
又 \(X>0\)，所以
\[
F_Y(y)=P\{0<X\le\sqrt y\}
=\int_0^{\sqrt y}\lambda e^{-\lambda x}\,\mathrm dx
=1-e^{-\lambda\sqrt y}.
\]
因此
\[
F_Y(y)=
\begin{cases}
0, & y<4,\\
1-e^{-\lambda\sqrt y}, & y\ge4.
\end{cases}
\]
特别地，\(Y\) 在 \(4\) 处有跳跃，且
\(
P\{Y=4\}=F_Y(4)-F_Y(4-)=1-e^{-2\lambda}.
\)
\end{solution}

\begin{exercise}\label{exer:src-2012-b-07-2credit}
设 \(X\) 服从参数为 \(\lambda\) 的指数分布，其概率密度为
\[
f_X(x)=
\begin{cases}
\lambda e^{-\lambda x}, & x>0,\\
0, & x\le0.
\end{cases}
\]
随机变量
\(
Y=g(X)=
\begin{cases}
X, & X>1,\\
2X, & X\le1.
\end{cases}
\)
求 \(Y\) 的分布函数 \(F_Y(y)\) 在 \(0<y\le2\) 上的表达式。
\end{exercise}
\begin{solution}
当 \(0<y\le1\) 时，若 \(X>1\)，则 \(Y=X>1\ge y\)，不会贡献；只需考虑 \(X\le1\) 且 \(2X\le y\)，故
\(
F_Y(y)=P\{X\le y/2\}=1-e^{-\lambda y/2}.
\)

当 \(1<y\le2\) 时，
\[
\begin{aligned}
F_Y(y)
&=P\{2X\le y,\ X\le1\}+P\{X\le y,\ X>1\}\\
&=P\{X\le y/2\}+P\{1<X\le y\}\\
&=(1-e^{-\lambda y/2})+(e^{-\lambda}-e^{-\lambda y}).
\end{aligned}
\]
因此
\[
F_Y(y)=
\begin{cases}
1-e^{-\lambda y/2}, & 0<y\le1,\\
1+e^{-\lambda}-e^{-\lambda y/2}-e^{-\lambda y}, & 1<y\le2.
\end{cases}
\]
\end{solution}

\begin{exercise}\label{exer:src-2011-a-06-2credit}
设随机变量 \(X\) 与 \(Y\) 独立，其中 \(X\) 的概率分布为
\(
X\sim
\begin{pmatrix}
1 & 2\\
0.3 & 0.7
\end{pmatrix},
\)
而 \(Y\) 的概率密度为 \(f(y)\)。求随机变量 \(U=X+Y\) 的概率密度 \(g(u)\)。
\end{exercise}
\begin{solution}
设 \(F(y)\) 为 \(Y\) 的分布函数，\(G(u)\) 为 \(U\) 的分布函数。由全概率公式，
\[
\begin{aligned}
G(u)
&=P(X+Y\le u)=0.3P(Y\le u-1\mid X=1)+0.7P(Y\le u-2\mid X=2)\\
&=0.3F(u-1)+0.7F(u-2)\\
g(u)&=0.3f(u-1)+0.7f(u-2).
\end{aligned}
\]
\end{solution}


\begin{exercise}[卷积公式分段积分] \label{exer:3-8b}
设随机变量 $X, Y$ 相互独立，其概率密度分别为
\[
f_X(x) = \begin{cases} 2e^{-2x}, & x > 0, \\ 0, & \text{其他}, \end{cases} \quad f_Y(y) = \begin{cases} 3y^2, & 0 < y < 1, \\ 0, & \text{其他}. \end{cases}
\]
(1) 求 $(X,Y)$ 的概率密度；(2) 求 $Z = X + Y$ 的概率密度。
\end{exercise}

\begin{solution}
(1) 由独立性，联合概率密度为：
\[
f(x,y) = f_X(x)\,f_Y(y) = \begin{cases} 6y^2 e^{-2x}, & x > 0,\ 0 < y < 1 \\ 0, & \text{其他} \end{cases}
\]

(2) 下求 $Z=X+Y$ 的分布函数 $F_Z(z) = P(X+Y \leqslant z)$。
分三种情况讨论：
\begin{enumerate}
    \item $z \leqslant 0$ 时，$F_Z(z) = 0$
    \item $0 < z < 1$ 时，积分区域为 $x>0, 0<y<1, x+y \leqslant z$，即 $0 < y < z, 0 < x < z-y$。
\begin{align*}
    F_Z(z) &= \iint_{x+y\leqslant z} f(x,y) \,\mathrm{d}x\,\mathrm{d}y = \int_0^z \int_0^{z-y} 6y^2 e^{-2x} \,\mathrm{d}x\,\mathrm{d}y \\
    &= 3 \int_0^z y^2 \left( \left. -e^{-2x} \right|_0^{z-y} \right) \,\mathrm{d}y = 3 \int_0^z y^2 \left( 1 - e^{-2(z-y)} \right) \,\mathrm{d}y \\
    &= 3 \int_0^z y^2 \,\mathrm{d}y - 3e^{-2z} \int_0^z y^2 e^{2y} \,\mathrm{d}y \\
    &= z^3 - 3e^{-2z} \left[ \left( \frac{1}{2}y^2 - \frac{1}{2}y + \frac{1}{4} \right) e^{2y} \right]_0^z \\
    &= z^3 - 3e^{-2z} \left( \left( \frac{1}{2}z^2 - \frac{1}{2}z + \frac{1}{4} \right) e^{2z} - \frac{1}{4} \right) \\
    &= z^3 - \frac{3}{2}z^2 + \frac{3}{2}z - \frac{3}{4} + \frac{3}{4}e^{-2z}
\end{align*}
求导：
\begin{align*}
    f_Z(z) &= \frac{\mathrm{d}}{\mathrm{d}z} F_Z(z) = 3z^2 - 3z + \frac{3}{2} - \frac{3}{2}e^{-2z}
\end{align*}

    \item $z \geqslant 1$ 时，积分区域为 $x>0, 0<y<1, x+y \leqslant z$，即 $0 < y < 1, 0 < x < z-y$。
\begin{align*}
    F_Z(z) &= \iint_{x+y\leqslant z} f(x,y) \,\mathrm{d}x\,\mathrm{d}y = \int_0^1 \int_0^{z-y} 6y^2 e^{-2x} \,\mathrm{d}x\,\mathrm{d}y \\
    &= 3 \int_0^1 y^2 \left( 1 - e^{-2(z-y)} \right) \,\mathrm{d}y = 3 \int_0^1 y^2 \,\mathrm{d}y - 3e^{-2z} \int_0^1 y^2 e^{2y} \,\mathrm{d}y \\
    &= 1 - 3e^{-2z} \left[ \left( \frac{1}{2}y^2 - \frac{1}{2}y + \frac{1}{4} \right) e^{2y} \right]_0^1 \\
    &= 1 - 3e^{-2z} \left( \frac{1}{4}e^2 - \frac{1}{4} \right) = 1 - \frac{3(e^2-1)}{4}e^{-2z}
\end{align*}
求导：
\begin{align*}
    f_Z(z) &= \frac{\mathrm{d}}{\mathrm{d}z} F_Z(z) = \frac{3(e^2-1)}{2}e^{-2z}
\end{align*}
\end{enumerate}
综上，$Z = X + Y$ 的分布密度函数为：
\[
f_Z(z) = \begin{cases} 3z^2 - 3z + \dfrac{3}{2} - \dfrac{3}{2}e^{-2z}, & 0 < z < 1 \\ \dfrac{3(e^2-1)}{2}e^{-2z}, & z \geqslant 1 \\ 0, & \text{其他} \end{cases}
\]
\end{solution}

\begin{exercise}\label{exer:src-paper1-app-03}\label{prob:p1-app-03}
设随机变量 \(X\sim U(0,1)\)，\(Y\sim E(1)\)，且它们相互独立。试求
\(
Z=2X-Y
\)
的密度函数 \(f_Z(z)\)。
\end{exercise}
\begin{solution}
由 \(Z=2X-Y\) 得 \(Y=2X-Z\)。在给定 \(z\) 时，需要满足
\(
0<x<1,\qquad 2x-z\ge0.
\)
因此
\[
f_Z(z)=\int f_X(x)f_Y(2x-z)\,\mathrm dx
=\int e^{-(2x-z)}\,\mathrm dx,
\]
积分区间由 \(0<x<1\) 与 \(x\ge z/2\) 共同决定。

当 \(z\le0\) 时，\(z/2\le0\)，故 \(0<x<1\)，
\(
f_Z(z)=\int_0^1 e^{-(2x-z)}\,\mathrm dx
=\frac{e^z(1-e^{-2})}{2}.
\)
当 \(0<z<2\) 时，\(z/2<x<1\)，
\(
f_Z(z)=\int_{z/2}^{1} e^{-(2x-z)}\,\mathrm dx
=\frac{1-e^{z-2}}{2}.
\)
当 \(z\ge2\) 时，没有满足条件的 \(x\)，故密度为 \(0\)。于是
\[
f_Z(z)=
\begin{cases}
\dfrac{e^z(1-e^{-2})}{2}, & z\le0,\\
\dfrac{1-e^{z-2}}{2}, & 0<z<2,\\
0, & z\ge2.
\end{cases}
\]
\end{solution}

\begin{exercise}[积的分布：矩形面积] \label{exer:3-8c}
设二维随机变量 $(X,Y)$ 在矩形区域 $D = \{(x,y) : 0 \leqslant x \leqslant 2,\ 0 \leqslant y \leqslant 1\}$ 上服从均匀分布，求边长为 $X$ 和 $Y$ 的矩形面积 $Z = XY$ 的概率密度。
\end{exercise}

\begin{solution}
由题设，联合概率密度为：
\[
f(x,y) = \begin{cases} \dfrac{1}{2}, & 0 \leqslant x \leqslant 2,\ 0 \leqslant y \leqslant 1 \\ 0, & \text{其他} \end{cases}
\]
下求 $Z=XY$ 的分布函数 $F_Z(z) = P(XY \leqslant z)$。
分三种情况讨论：
\begin{enumerate}
    \item $z \leqslant 0$ 时，$F_Z(z) = 0$
    \item $0 < z < 2$ 时，积分区域为矩形 $D$ 内满足 $xy \leqslant z$ 的部分。边界曲线为 $y = \frac{z}{x}$，与 $y=1$ 交于 $x=z$。我们将积分区域按 $x$ 划分：
\begin{align*}
    F_Z(z) &= \iint_{xy \leqslant z} f(x,y) \,\mathrm{d}x\,\mathrm{d}y = \int_0^z \left( \int_0^1 \frac{1}{2} \,\mathrm{d}y \right) \mathrm{d}x + \int_z^2 \left( \int_0^{\frac{z}{x}} \frac{1}{2} \,\mathrm{d}y \right) \mathrm{d}x \\
    &= \frac{1}{2} \int_0^z 1 \,\mathrm{d}x + \frac{1}{2} \int_z^2 \frac{z}{x} \,\mathrm{d}x = \frac{1}{2}z + \frac{z}{2} \left( \left. \ln x \right|_z^2 \right) = \frac{z}{2} + \frac{z}{2} (\ln 2 - \ln z) = \frac{z}{2}(1 + \ln 2 - \ln z)
\end{align*}
求导：
\begin{align*}
    f_Z(z) &= \frac{\mathrm{d}}{\mathrm{d}z} F_Z(z) = \frac{\mathrm{d}}{\mathrm{d}z} \left[ \frac{z}{2}(1 + \ln 2 - \ln z) \right] = \frac{1}{2}(1 + \ln 2 - \ln z) + \frac{z}{2} \left( -\frac{1}{z} \right) = \frac{1}{2}(\ln 2 - \ln z)
\end{align*}
    \item $z \geqslant 2$ 时，矩形区域完全包含在 $xy \leqslant z$ 内，$F_Z(z) = 1$。求导得 $f_Z(z) = 0$
\end{enumerate}
综上，$Z = XY$ 的分布密度函数为：
\[
f_Z(z) = \begin{cases} \dfrac{1}{2}(\ln 2 - \ln z), & 0 < z < 2 \\ 0, & \text{其他} \end{cases}
\]
\end{solution}

\begin{exercise}[混合型：全概率公式] \label{exer:3-8d}
设随机变量 $X_1, X_2, X_3$ 相互独立，其中 $X_1$ 与 $X_2$ 均服从标准正态分布，$X_3$ 的概率分布为 $P\{X_3 = 0\} = P\{X_3 = 1\} = 1/2$。令  $Y = X_1 X_3 + X_2 (1-X_3)$。(1) 求 $(X_1, Y)$ 的分布函数（结果用 $\Phi(x)$ 表示）；(2) 证明 $Y$ 服从标准正态分布。
\end{exercise}
\begin{solution}
(1) 记 $(X_1, Y)$ 的联合分布函数为 $F(x,y) = P(X_1 \leqslant x, Y \leqslant y)$。
由全概率公式，按 $X_3$ 的取值进行分解：
\begin{align*}
    F(x,y) &= P(X_3 = 0) P(X_1 \leqslant x, Y \leqslant y \mid X_3 = 0) + P(X_3 = 1) P(X_1 \leqslant x, Y \leqslant y \mid X_3 = 1) \\
    &= \frac{1}{2} P(X_1 \leqslant x, X_2 \leqslant y \mid X_3 = 0) + \frac{1}{2} P(X_1 \leqslant x, X_1 \leqslant y \mid X_3 = 1)
\end{align*}
由于 $X_1, X_2, X_3$ 相互独立，条件概率与 $X_3$ 无关，可直接转化为无条件概率：
\begin{align*}
    F(x,y) &= \frac{1}{2} P(X_1 \leqslant x, X_2 \leqslant y) + \frac{1}{2} P(X_1 \leqslant x, X_1 \leqslant y) \\
    &= \frac{1}{2} P(X_1 \leqslant x)P(X_2 \leqslant y) + \frac{1}{2} P(X_1 \leqslant \min(x,y)) \\
    &= \frac{1}{2} \Phi(x)\Phi(y) + \frac{1}{2} \Phi(\min(x,y))
\end{align*}
(2) 下求 $Y$ 的边缘分布函数 $F_Y(y)$：
\begin{align*}
    F_Y(y) &= \lim_{x \to +\infty} F(x,y) = \lim_{x \to +\infty} \left[ \frac{1}{2} \Phi(x)\Phi(y) + \frac{1}{2} \Phi(\min(x,y)) \right] = \frac{1}{2} \cdot 1 \cdot \Phi(y) + \frac{1}{2} \Phi(y) = \Phi(y)
\end{align*}
由于 $F_Y(y) = \Phi(y)$，因此严格证明了 $Y$ 服从标准正态分布 $N(0,1)$。
\end{solution}

\subsubsection{\texorpdfstring{$\max$ 和 $\min$ 的分布}{max 和 min 的分布}}

\begin{property}[$\max$ 和 $\min$ 的分布]\label{prop:max-min-dist}
设 $(X,Y) \sim f(x,y)$，分布函数为 $F(x,y)$。

\textbf{(1) 一般情形}：
\begin{enumerate}
    \item $Z = \max\{X, Y\}$ 的分布函数为
    \(
    F_{\max}(z) = P\{X \leqslant z, Y \leqslant z\} = F(z, z).
    \)
    \item $Z = \min\{X, Y\}$ 的分布函数为
    \[
    F_{\min}(z) = P\{X \leqslant z\} + P\{Y \leqslant z\} - P\{X \leqslant z, Y \leqslant z\} = F_X(z) + F_Y(z) - F(z, z).
    \]
\end{enumerate}

\textbf{(2) 独立情形}：
\begin{enumerate}
    \item 当 $X$ 与 $Y$ 独立时，$F_{\max}(z) = F_X(z) \cdot F_Y(z)$，$F_{\min}(z) = 1 - [1 - F_X(z)][1 - F_Y(z)]$。
    \item 推广到 $n$ 个相互独立的随机变量 $X_1, X_2, \dots, X_n$（分布函数为 $F_{X_i}(x)$）：
    \[
    F_{\max}(z) = \prod_{i=1}^{n}F_{X_i}(z), \qquad F_{\min}(z) = 1 - \prod_{i=1}^{n}[1 - F_{X_i}(z)].
    \]
    \item 当 $X_1, \dots, X_n$ 独立同分布（分布函数为 $F(x)$，概率密度为 $f(x)$）时：
    \(
    F_{\max}(z) = [F(z)]^n, \quad f_{\max}(z) = n[F(z)]^{n-1}f(z),
    \)
    \(
    F_{\min}(z) = 1 - [1 - F(z)]^n, \quad f_{\min}(z) = n[1 - F(z)]^{n-1}f(z).
    \)
    这些结果在数理统计部分有重要应用。
\end{enumerate}
\end{property}

\begin{exercise}\label{exer:3-9}
一辆新车装有 4 个轮胎并配有 1 个备胎，当使用的轮胎损坏时就用备胎更换。设所有轮胎的寿命均服从参数为 $\lambda = 0.2$（单位：万小时）的指数分布且相互独立，记 $T$ 为开始使用备胎的时间，求 $T$ 的分布函数。
\end{exercise}

\begin{solution}
开始使用备胎的时刻 $T$ 等于 4 个轮胎中寿命最短的那个的寿命，即 $T = \min\{X_1, X_2, X_3, X_4\}$。

每个轮胎的寿命 $X_i \sim E(0.2)$，分布函数为 $F(x) = 1 - e^{-0.2x}$（$x \geqslant 0$）。

由最小值的分布公式：
\(
F_T(t) = 1 - [1 - F(t)]^4 = 1 - [e^{-0.2t}]^4 = 1 - e^{-0.8t}, \quad t \geqslant 0.
\)

即 $T \sim E(0.8)$，备胎在参数为 $0.8$ 的指数分布时刻开始被使用。

\begin{note}
更一般地，若 $X_1, X_2, \dots, X_n$ 独立同分布于 $E(\lambda)$，则 $\min\{X_1, \dots, X_n\} \sim E(n\lambda)$。
\end{note}
\end{solution}


\begin{exercise}[均匀分布最大值密度] \label{exer:3-10b}
设随机变量 $X, Y$ 相互独立，且都在区间 $[a, b]$ 上服从均匀分布，求 $Z = \max\{X, Y\}$ 的概率密度。
\end{exercise}

\begin{solution}
由独立性，$(X,Y)$ 的概率密度为
\[
f(x,y) = \begin{cases} \dfrac{1}{(b-a)^2}, & a \leqslant x \leqslant b,\ a \leqslant y \leqslant b, \\ 0, & \text{其他}. \end{cases}
\]

$Z = \max\{X, Y\}$ 的取值范围为 $[a, b]$。当 $a \leqslant z \leqslant b$ 时，
\[
F_Z(z) = P\{Z \leqslant z\} = P\{X \leqslant z,\ Y \leqslant z\} = \int_a^z\mathrm{d}x\int_a^z \frac{1}{(b-a)^2}\,\mathrm{d}y = \frac{(z-a)^2}{(b-a)^2}.
\]

求导得概率密度：
\[
f_Z(z) = F_Z'(z) = \begin{cases} \dfrac{2(z-a)}{(b-a)^2}, & a \leqslant z \leqslant b, \\ 0, & \text{其他}. \end{cases}
\]

\begin{note}
这与独立同分布时 $f_{\max}(z) = n[F(z)]^{n-1}f(z)$ 的公式一致：$n=2$，$F(z) = \dfrac{z-a}{b-a}$，$f(z) = \dfrac{1}{b-a}$，故 $f_{\max}(z) = 2 \cdot \dfrac{z-a}{b-a} \cdot \dfrac{1}{b-a} = \dfrac{2(z-a)}{(b-a)^2}$。
\end{note}
\end{solution}



\begin{exercise}[指数分布条件区域] \label{exer:3-12b}
设二维随机变量 $(X,Y)$ 的概率密度函数为
\[
f(x,y) = \begin{cases} \lambda e^{-\lambda y}, & 0 < x < y, \\ 0, & \text{其他}, \end{cases}
\]
其中 $\lambda > 0$。求：(1) $Z = X + Y$ 的概率密度；(2) $N = \min(X,Y)$ 的概率密度。
\end{exercise}

\begin{solution}
验证归一性：$\displaystyle\int_0^{+\infty}\mathrm{d}y\int_0^y \lambda e^{-\lambda y}\,\mathrm{d}x = \int_0^{+\infty}y\lambda e^{-\lambda y}\,\mathrm{d}y = 1$（$Y \sim \Gamma(2,\lambda)$）。

边缘密度：
\[
f_X(x) = \int_x^{+\infty}\lambda e^{-\lambda y}\,\mathrm{d}y = e^{-\lambda x}, \quad x > 0. \qquad \text{即 } X \sim E(\lambda).
\]
\[
f_Y(y) = \int_0^y \lambda e^{-\lambda y}\,\mathrm{d}x = \lambda y\,e^{-\lambda y}, \quad y > 0. \qquad \text{即 } Y \sim \Gamma(2,\lambda).
\]

(1) 当 $z > 0$ 时，$\{X+Y \leqslant z\}$ 等价于 $\{0 < x < y,\ x+y \leqslant z\}$，即 $0 < x < z/2$，$x < y \leqslant z-x$：
\begin{align*}
F_Z(z) &= \int_0^{z/2}\mathrm{d}x\int_x^{z-x}\lambda e^{-\lambda y}\,\mathrm{d}y = \int_0^{z/2}(e^{-\lambda x} - e^{-\lambda(z-x)})\,\mathrm{d}x\\
&= \left[-\frac{1}{\lambda}e^{-\lambda x} - \frac{1}{\lambda}e^{-\lambda z}e^{\lambda x}\right]_0^{z/2} = \frac{1}{\lambda}\bigl(1 - e^{-\lambda z/2}\bigr) - \frac{1}{\lambda}e^{-\lambda z}(e^{\lambda z/2} - 1)\\
&= \frac{1}{\lambda}\bigl(1 - e^{-\lambda z/2} - e^{-\lambda z/2} + e^{-\lambda z}\bigr) = \frac{1}{\lambda}\bigl(1 - 2e^{-\lambda z/2} + e^{-\lambda z}\bigr).
\end{align*}
当 $z \leqslant 0$ 时，$F_Z(z) = 0$。求导：
\[
f_Z(z) = \frac{1}{\lambda}\bigl(\lambda e^{-\lambda z/2} - \lambda e^{-\lambda z}\bigr) = e^{-\lambda z/2} - e^{-\lambda z}, \quad z > 0.
\]
所以
\[
f_Z(z) = \begin{cases} e^{-\lambda z/2} - e^{-\lambda z}, & z > 0, \\ 0, & z \leqslant 0. \end{cases}
\]
(2) 在区域 $\{0 < x < y\}$ 中，$\min(X,Y) = X$。因此 $N = X$，其密度为
\[
f_N(n) = f_X(n) = \begin{cases} e^{-\lambda n}, & n > 0, \\ 0, & n \leqslant 0. \end{cases}
\]
即 $N \sim E(\lambda)$。
\end{solution}


\begin{exercise}[综合性：常数、独立性、和的分布、min的概率] \label{exer:3-12e}
设二维随机变量 $(X,Y)$ 具有联合密度函数
\[
f(x,y) = \begin{cases} cxe^{-y}, & 0 < x < y < +\infty, \\ 0, & \text{其他}. \end{cases}
\]
求：
\begin{enumerate}
    \item[(1)] 常数 $c$；
    \item[(2)] $X$ 与 $Y$ 是否独立？
    \item[(3)] $Z = X + Y$ 的密度函数；
    \item[(4)] $P\{\min(X,Y) < 1\}$。
\end{enumerate}
\end{exercise}
\begin{solution}
(1) 由概率密度的归一性：
\[
1 = \int_{-\infty}^{+\infty}\int_{-\infty}^{+\infty}f(x,y)\,\mathrm{d}x\,\mathrm{d}y = \int_0^{+\infty}\int_0^y cxe^{-y}\,\mathrm{d}x\,\mathrm{d}y.
\]
内层积分 $\int_0^y x\,\mathrm{d}x = \dfrac{y^2}{2}$，因此
\[
\int_0^{+\infty}c \cdot \frac{y^2}{2}e^{-y}\,\mathrm{d}y = \frac{c}{2}\Gamma(3) = \frac{c}{2} \cdot 2! = c = 1.
\]
(2) 求边缘概率密度：
\[
f_X(x) = \int_x^{+\infty}xe^{-y}\,\mathrm{d}y = xe^{-x}, \quad x > 0. \qquad f_Y(y) = \int_0^y xe^{-y}\,\mathrm{d}x = \frac{y^2}{2}e^{-y}, \quad y > 0.
\]
由于 $f_X(x)\,f_Y(y) = \dfrac{xy^2}{2}e^{-(x+y)} \neq xe^{-y} = f(x,y)$（取 $x=1, y=2$ 可直接验证），故 $X$ 与 $Y$ \textbf{不独立}。\\
(3) 当 $z < 0$ 时，$F_Z(z) = 0$；当 $z \geqslant 0$ 时，积分区域为 $\{0 < x < y,\ x+y \leqslant z\}$，即 $0 < x < z/2$，$x < y \leqslant z-x$：
\[
F_Z(z) = \int_0^{z/2}\mathrm{d}x\int_x^{z-x}xe^{-y}\,\mathrm{d}y = \int_0^{z/2}x(e^{-x} - e^{-(z-x)})\,\mathrm{d}x
\]
\[
= \left[1 - \left(\frac{z}{2}+1\right)e^{-z/2}\right] - e^{-z}\left[\left(\frac{z}{2}-1\right)e^{z/2} + 1\right] = 1 - ze^{-z/2} - e^{-z}.
\]
综上：
\[
f_Z(z) = \begin{cases} \left(\dfrac{z}{2}-1\right)e^{-z/2} + e^{-z}, & z \geqslant 0, \\ 0, & z < 0. \end{cases}
\]
(4)在 $\{0 < x < y\}$ 的条件下 $\min(X,Y) = X$，因此 $P\{X \geqslant 1, Y \geqslant 1\}$ 的积分区域为 $\{1 \leqslant x < y\}$。
\[
P\{\min(X,Y) < 1\} = 1 - P\{\min(X,Y) \geqslant 1\} = 1 - P\{X \geqslant 1, Y \geqslant 1\}
\]
\[
= 1 - \int_1^{+\infty}\mathrm{d}y\int_1^y xe^{-y}\,\mathrm{d}x = 1 - \int_1^{+\infty}\frac{y^2-1}{2}e^{-y}\,\mathrm{d}y
\]
\[
= 1 - \frac{1}{2}\left[\int_1^{+\infty}y^2 e^{-y}\,\mathrm{d}y - \int_1^{+\infty}e^{-y}\,\mathrm{d}y\right] = 1 - \frac{1}{2}(5e^{-1} - e^{-1}) = 1 - 2e^{-1}.
\]
\end{solution}





\begin{exercise}[min+max=X+Y] \label{exer:3-c10b}
$(X,Y)$ 服从区域 $G=\{(x,y):0\leqslant x\leqslant 1,\ 0\leqslant y\leqslant 2\}$ 上的均匀分布，令 $Z=\min(X,Y)$，$W=\max(X,Y)$，求 $P(Z+W \geqslant 1)$。
\end{exercise}
\begin{solution}
先注意到
\(
Z=\min(X,Y),\qquad W=\max(X,Y),
\)
无论 \(X\) 与 \(Y\) 谁大谁小，总有
\(
Z+W=X+Y.
\)
因此题目要求的概率就是
\(
P(Z+W\geqslant 1)=P(X+Y\geqslant 1).
\)
\[
P(X+Y \geqslant 1) = 1 - P(X+Y < 1) = 1 - \int_0^1\mathrm{d}x\int_0^{1-x}\frac{1}{2}\,\mathrm{d}y = 1 - \frac{1}{2}\int_0^1(1-x)\,\mathrm{d}x = 1-\frac{1}{4} = \frac{3}{4}.
\]
\end{solution}


\begin{exercise}\label{exer:src-2004-2a-15}
设某班车起点站上车人数 \(X\) 服从参数为 \(\lambda\ (\lambda>0)\) 的泊松分布，并且中途不再有人上车。车上每位乘客在中途下车的概率为 \(p\ (0<p<1)\)，且是否中途下车相互独立。以 \(Y\) 表示在中途下车的人数，试求：\par
(1) \((X,Y)\) 的联合概率分布律；\par
(2) \(Y\) 的分布律。
\end{exercise}
\begin{solution}
在给定 \(X=n\) 的条件下，\(Y\) 表示 \(n\) 名乘客中中途下车的人数，因此
\(
Y|X=n\sim B(n,p).
\)
所以当 \(n=0,1,2,\dots\)，\(k=0,1,\dots,n\) 时，
\[
\begin{aligned}
P(X=n,Y=k)
&=P(X=n)P(Y=k|X=n)=e^{-\lambda}\frac{\lambda^n}{n!}\binom nkp^k(1-p)^{n-k}.
\end{aligned}
\]
若 \(k>n\)，则 \(P(X=n,Y=k)=0\)。
再求 \(Y\) 的分布。对 \(k=0,1,2,\dots\)，
\[
\begin{aligned}
P(Y=k)
&=\sum_{n=k}^{\infty}P(X=n,Y=k)=\sum_{n=k}^{\infty}e^{-\lambda}\frac{\lambda^n}{n!}\binom nkp^k(1-p)^{n-k}\\
&=e^{-\lambda}\frac{(\lambda p)^k}{k!}
\sum_{m=0}^{\infty}\frac{[\lambda(1-p)]^m}{m!}=e^{-\lambda}\frac{(\lambda p)^k}{k!}e^{\lambda(1-p)}
=e^{-\lambda p}\frac{(\lambda p)^k}{k!}.
\end{aligned}
\]
因此
\(
Y\sim P(\lambda p).
\)
，原来总人数服从 \(P(\lambda)\)，每人独立保留概率为 \(p\)，则保留下来的人数服从 \(P(\lambda p)\)。
\end{solution}

\begin{exercise}\label{exer:src-2008-probstat-2a-v-01}\label{exer:src-paper3-iv-03}\label{prob:p3-iv-03}
设二维随机变量 \((X,Y)\) 的联合密度函数为
\[
f(x,y)=
\begin{cases}
2e^{-2x-y},&x>0,\ y>0,\\
0,&\text{其他}.
\end{cases}
\]
求 \(Z=\max\{X,Y\}\) 的密度函数。
\end{exercise}
\begin{solution}
由联合密度可分解为
\[
f(x,y)=\bigl(2e^{-2x}\bigr)\bigl(e^{-y}\bigr),\qquad x>0,\ y>0,
\]
故 \(X,Y\) 相互独立，且
\(
F_X(z)=1-e^{-2z},\qquad F_Y(z)=1-e^{-z}\qquad (z>0).
\)

当 \(z\le0\) 时，\(F_Z(z)=0\)。当 \(z>0\) 时，
\[
\begin{aligned}
F_Z(z)
&=P\{\max(X,Y)\le z\}\\
&=P(X\le z,Y\le z)\\
&=F_X(z)F_Y(z)
=(1-e^{-2z})(1-e^{-z}).
\end{aligned}
\]
对 \(z>0\) 求导，得
\(
f_Z(z)=2e^{-2z}(1-e^{-z})+(1-e^{-2z})e^{-z}
=e^{-z}+2e^{-2z}-3e^{-3z}.
\)
所以
\[
f_Z(z)=
\begin{cases}
e^{-z}+2e^{-2z}-3e^{-3z},&z>0,\\
0,&\text{其他}.
\end{cases}
\]
\end{solution}

\begin{exercise}\label{exer:src-2008-probstat-2a-vii-01}
随机变量 \((X,Y)\) 服从区域
\(
\{(x,y):0<x<2,\ 0<y<1\}
\)
上的均匀分布。\par
(1) 求 \((X,Y)\) 的概率密度函数及分布函数；\par
(2) 设 \(\xi=X+Y,\ \eta=aX+bY\)，且 \(\xi,\eta\) 不相关，\(D\eta=1\)，求 \(a,b\)。
\end{exercise}
\begin{solution}
(1) 矩形区域面积为 \(2\)，故联合密度为
\[
f(x,y)=
\begin{cases}
\dfrac12,&0<x<2,\ 0<y<1,\\
0,&\text{其他}.
\end{cases}
\]
边缘密度为
\[
f_X(x)=
\begin{cases}
\dfrac12,&0<x<2,\\
0,&\text{其他},
\end{cases}
\qquad
f_Y(y)=
\begin{cases}
1,&0<y<1,\\
0,&\text{其他}.
\end{cases}
\]
所以 \(X,Y\) 相互独立。分布函数可写为
\(
F(x,y)=F_X(x)F_Y(y),
\)
其中
\[
F_X(x)=
\begin{cases}
0,&x\le0,\\
\dfrac{x}{2},&0<x\le2,\\
1,&x>2,
\end{cases}
\qquad
F_Y(y)=
\begin{cases}
0,&y\le0,\\
y,&0<y\le1,\\
1,&y>1.
\end{cases}
\]

(2) 由均匀分布的数字特征，
\[
EX=1,\quad DX=\frac{(2-0)^2}{12}=\frac13,\qquad
EY=\frac12,\quad DY=\frac{(1-0)^2}{12}=\frac1{12}.
\]
又 \(X,Y\) 独立，所以 \(\operatorname{Cov}(X,Y)=0\)。于是
\[
\operatorname{Cov}(\xi,\eta)
=\operatorname{Cov}(X+Y,aX+bY)
=aDX+bDY
=\frac a3+\frac b{12}.
\]
由 \(\xi,\eta\) 不相关，得
\(
b=-4a
\)，
再由
\(
D\eta=D(aX+bY)=a^2DX+b^2DY=\frac{a^2}{3}+\frac{b^2}{12}=1,
\)
代入 \(b=-4a\)，得
\(
\frac{a^2}{3}+\frac{16a^2}{12}=1
\quad\Longrightarrow\quad
\frac{5a^2}{3}=1.
\)
因此
\[
a=\pm\sqrt{\frac35},\qquad b=\mp4\sqrt{\frac35}.
\]
\end{solution}


\begin{exercise}\label{exer:src-2008-01-3a-iv-02}\label{exer:src-paper2-app-02}\label{prob:p2-app-02}
设随机变量 \(X\) 与 \(Y\) 相互独立，\(X,Y\) 分别服从参数为 \(\lambda,\mu\ (\lambda\ne\mu)\) 的指数分布。求
\(
Z=3X+2Y
\)
的密度函数 \(f_Z(z)\)。
\end{exercise}
\begin{solution}
令
\(
U=3X,\qquad V=2Y.
\)
则 \(U,V\) 相互独立，且
\(
f_U(u)=\frac{\lambda}{3}e^{-\lambda u/3},\qquad u>0,
\)
\(
f_V(v)=\frac{\mu}{2}e^{-\mu v/2},\qquad v>0.
\)
因为 \(Z=U+V\)，所以对 \(z>0\)，
\[
\begin{aligned}
f_Z(z)
&=\int_0^z f_U(u)f_V(z-u)\,\mathrm{d}u\\
&=\int_0^z
\frac{\lambda}{3}e^{-\lambda u/3}
\frac{\mu}{2}e^{-\mu(z-u)/2}\,\mathrm{d}u\\
&=\frac{\lambda\mu}{6}e^{-\mu z/2}
\int_0^z e^{-(\lambda/3-\mu/2)u}\,\mathrm{d}u.
\end{aligned}
\]
若 \(\lambda/3\ne\mu/2\)，即 \(2\lambda\ne3\mu\)，则
\[
f_Z(z)=
\frac{\lambda\mu}{6}
\frac{e^{-\mu z/2}-e^{-\lambda z/3}}{\lambda/3-\mu/2},
\qquad z>0.
\]
等价地，
\[
f_Z(z)=
\frac{\lambda\mu}{2\lambda-3\mu}
\left(e^{-\mu z/2}-e^{-\lambda z/3}\right),
\qquad z>0.
\]
若 \(2\lambda=3\mu\)，则两个指数率相同，卷积退化为
\(
f_Z(z)=\frac{\lambda\mu}{6}z e^{-\lambda z/3},\qquad z>0.
\)
综上，
\[
f_Z(z)=
\begin{cases}
\displaystyle
\frac{\lambda\mu}{2\lambda-3\mu}
\left(e^{-\mu z/2}-e^{-\lambda z/3}\right),
& z>0,\ 2\lambda\ne3\mu,\\[10pt]
\displaystyle
\frac{\lambda\mu}{6}z e^{-\lambda z/3},
& z>0,\ 2\lambda=3\mu,\\[8pt]
0, & z\le0.
\end{cases}
\]
\end{solution}



\begin{exercise}\label{exer:src-2008-01-3c-iv-03}\label{exer:src-paper4-iv-03}
设随机变量 \(X\) 在 \((0,a)\) 上服从均匀分布。当观察到 \(X=x\ (0<x<a)\) 时，\(Y\) 在区间 \((x,a)\) 内任一子区间上取值的概率与该子区间长度成正比。求：\par
(1) \((X,Y)\) 的联合密度函数 \(f(x,y)\)；\par
(2) \(Y\) 的密度函数 \(f_Y(y)\)。
\end{exercise}
\begin{solution}
由题意，
\(
f_X(x)=\frac1a,\qquad 0<x<a.
\)
给定 \(X=x\) 后，\(Y\) 在 \((x,a)\) 上均匀分布，因此条件密度为
\(
f_{Y|X}(y|x)=\frac1{a-x},\qquad x<y<a.
\)
于是联合密度为
\(
f(x,y)=f_X(x)f_{Y|X}(y|x)
=\frac1{a(a-x)},\qquad 0<x<y<a,
\)
即
\[
f(x,y)=
\begin{cases}
\dfrac1{a(a-x)}, & 0<x<y<a,\\[8pt]
0, & \text{其他}.
\end{cases}
\]

对 \(Y\) 求边缘密度。固定 \(0<y<a\)，有 \(0<x<y\)，故
\[
f_Y(y)=\int_0^y \frac1{a(a-x)}\,\mathrm{d}x
=\frac1a\left[-\ln(a-x)\right]_0^y
=\frac1a\ln\frac{a}{a-y}.
\]
所以
\[
f_Y(y)=
\begin{cases}
\dfrac1a\ln\dfrac{a}{a-y}, & 0<y<a,\\[10pt]
0, & \text{其他}.
\end{cases}
\]
\end{solution}

\chapter{随机变量的数字特征}


{\small
\[
\begin{cases}
\text{数学期望}
\begin{cases}
\text{计算公式}
\begin{cases}
\text{离散型}
\begin{cases}
\text{一维:}
\begin{cases}
EX = \sum\limits_{i} x_i p_i \\
E\varphi(X) = \sum\limits_{i} \varphi(x_i) p_i
\end{cases} \\
\text{二维: } E[\varphi(X,Y)] = \sum\limits_{i=1}^{m} \sum\limits_{j=1}^{n} \varphi(x_i,y_j) p_{ij}
\end{cases} \\
\text{连续型}
\begin{cases}
\text{一维:}
\begin{cases}
EX = \int_{-\infty}^{+\infty} x f(x) \mathrm{d}x \\
E\varphi(X) = \int_{-\infty}^{+\infty} \varphi(x) f(x) \mathrm{d}x
\end{cases} \\
\text{二维: } E[\varphi(X,Y)] = \int_{-\infty}^{+\infty} \int_{-\infty}^{+\infty} \varphi(x,y) f(x,y) \mathrm{d}x\mathrm{d}y
\end{cases}
\end{cases} \\
\text{性质}
\begin{cases}
(1)\ \text{常数的期望不变: } EC=C,\ \text{特别地: } E(EX)=EX \\
(2)\ \text{线性性质: 如果 } EX \text{ 和 } EY \text{ 都存在, 则 } E(aX+bY)=aEX+bEY \\
(3)\ \text{如果 } EX \text{ 存在, 则 } E(kX)=kEX \\
(4)\ \text{若 } X,Y \text{ 相互独立, 且期望都存在, 则 } E(XY)=EXEY
\end{cases}
\end{cases} \\
\text{方差}
\begin{cases}
\text{计算公式: } DX=E(X-EX)^2=E(X^2)-(EX)^2 \\
\text{性质}
\begin{cases}
(1)\ \text{常数的方差为零: } DC=0,~~~~ (2) DX\geqslant0 \Rightarrow E(X^2)\geqslant (EX)^2 \\
(3)\ D(aX+b)=a^2 DX,~~~ (4)\ D(X\pm Y)=DX+DY\pm 2\text{Cov}(X,Y) \\
(5)\ \text{若 } X,Y \text{ 相互独立, 则 } D(X\pm Y)=DX+DY \\
(6)\ DX=E(X-EX)^2\leqslant E(X-c)^2
\end{cases}
\end{cases} \\
\text{协方差}
\begin{cases}
\text{计算公式: } \text{Cov}(X,Y)=E(X-EX)(Y-EY)=E(XY)-EXEY \\
\text{其中 } E(XY)=
\begin{cases}
\sum\limits_{i}\sum\limits_{j} x_i y_j P\{X=x_i,Y=y_j\} \quad (\text{离散型}) \\
\int_{-\infty}^{+\infty}\int_{-\infty}^{+\infty} xy f(x,y) \mathrm{d}x\mathrm{d}y \quad (\text{连续型})
\end{cases} \\
\text{性质}
\begin{cases}
(1)\ \text{Cov}(X,X)=DX ,~~~~~(2)\ \text{Cov}(X,Y)=\text{Cov}(Y,X) \\
(3)\ \text{Cov}(X,C)=0 \ (C \text{是常数}) \\
(4)\ \text{Cov}(aX+bY,Z)=a\text{Cov}(X,Z)+b\text{Cov}(Y,Z)
\end{cases}
\end{cases} \\
\text{相关系数}
\begin{cases}
\text{计算公式: } \rho_{XY}=\dfrac{\text{Cov}(X,Y)}{\sqrt{DX}\sqrt{DY}},\ |\rho_{XY}|\leqslant1 \\
\text{性质}
\begin{cases}
(1)\ |\rho_{XY}|=1 \Leftrightarrow P\{Y=aX+b\}=1 \\
(2)\ \rho_{XY}=0 \Leftrightarrow X,Y \text{ 不相关} \Leftrightarrow
\begin{cases}
(1)\ \text{Cov}(X,Y)=0 \\
(2)\ E(XY)=EXEY \\
(3)\ D(X\pm Y)=DX+DY
\end{cases} \\
(3)\ \text{若 } (X,Y)\sim N(\mu_1,\mu_2,\sigma_1^2,\sigma_2^2,\rho), \\
\qquad \text{则 } X,Y \text{ 独立与 } X,Y \text{ 不相关等价}
\end{cases}
\end{cases}
\end{cases}
\]
}

\section{数学期望}

\subsection{数学期望的定义}

\begin{definition}[数学期望] \label{def:expectation}
设 \(X\) 是一个随机变量。
\begin{enumerate}
    \item \textbf{离散型}：若 \(X\) 的分布律为 \(P\{X=x_i\} = p_i\ (i=1,2,\dots)\)，且级数 \(\sum\limits_{i=1}^{\infty} x_i p_i\) 绝对收敛，则称 \(X\) 的\textbf{数学期望}存在，其值为
    \(
    EX = \sum_{i=1}^{\infty} x_i p_i.
    \)
    若级数不绝对收敛，则称 \(X\) 的数学期望不存在。
    \item \textbf{连续型}：若 \(X\) 的概率密度为 \(f(x)\)，且积分 \(\int_{-\infty}^{+\infty} x f(x)\,\mathrm{d}x\) \textbf{绝对收敛}，则称 \(X\) 的\textbf{数学期望}存在，其值为
    \(
    EX = \int_{-\infty}^{+\infty} x f(x)\,\mathrm{d}x.
    \)
\end{enumerate}
\end{definition}
\begin{definition}[随机变量函数的期望] \label{def:function-expectation}
设 \(Y = g(X)\)（\(g\) 为连续函数）。
\begin{enumerate}
    \item \textbf{离散型}：若 \(\sum\limits_{i=1}^{\infty} g(x_i) p_i\) 绝对收敛，则
    \(
    E[g(X)] = \sum_{i=1}^{\infty} g(x_i) p_i.
    \)
    \item \textbf{连续型}：若 \(\int_{-\infty}^{+\infty} g(x) f(x)\,\mathrm{d}x\) 绝对收敛，则
    \(
    E[g(X)] = \int_{-\infty}^{+\infty} g(x) f(x)\,\mathrm{d}x.
    \)
\end{enumerate}
\end{definition}

\begin{note}
(1) 数学期望又称为概率平均值，常简称期望或均值。它是描述随机变量平均取值状况的数字特征，刻画了随机变量一切可能值的集中位置。\par
(2) 定义中要求级数（或积分）\textbf{绝对收敛}，这是为了保证期望的值不依赖于求和（或积分）的顺序。\par
(3) 并非所有随机变量都存在数学期望。例如柯西分布 \(X\sim f(x)=\dfrac{1}{\pi(1+x^2)}\)，由于
\[
\int_{-\infty}^{+\infty}\frac{|x|}{\pi(1+x^2)}\,\mathrm{d}x = +\infty,
\]
其数学期望不存在。
\end{note}

\begin{exercise}\label{exer:4-5-poisson}
设随机变量 \(X\) 的概率分布为 \(P\{X=k\} = \dfrac{C}{k!},\ k=0,1,2,\dots\)，则 \(E(|X - EX|) =\,\underline{\hspace{2cm}}\)。
\end{exercise}
\begin{solution}
\[
\begin{aligned}
E(|X - EX|) &= E(|X - 1|) = \sum_{k=0}^{\infty} P(X=k)|k-1| = P(X=0) \cdot 1 + \sum_{k=2}^{\infty} P(X=k)(k-1) \\
&= e^{-1} + \sum_{k=2}^{\infty} \frac{e^{-1}}{k!}(k-1) = e^{-1} + e^{-1}\sum_{j=1}^{\infty}\frac{j}{(j+1)!} \\
&= e^{-1} + e^{-1}\sum_{j=1}^{\infty}\left(\frac{1}{j!} - \frac{1}{(j+1)!}\right) = e^{-1} + e^{-1}\cdot 1 = \frac{2}{e}.
\end{aligned}
\]
\end{solution}


\begin{exercise}\label{exer:4-20-shop}
某商店销售某商品，销量 \(X\) 在 \((50,100)\) 上服从均匀分布。若每销售一单位获利 500 元；如果需求量大于进货量 \(n\)，可以从其他部门调剂，此时每单位获利 300 元；如果有积压，则每单位亏损 200 元。进货量 \(n\) 为多少时平均获利最大？最大为多少？
\end{exercise}
\begin{solution}
设进货量为 \(n\)（\(n \in [50,100]\)），获利为
\[
Y = \begin{cases} 500n + 300(X-n) = 300X + 200n, & X \geqslant n, \\ 500X - 200(n-X) = 700X - 200n, & X < n. \end{cases}
\]
\[
\begin{aligned}
EY &= \frac{1}{50}\left[\int_{50}^{n}(700x - 200n)\,\mathrm{d}x + \int_{n}^{100}(300x + 200n)\,\mathrm{d}x\right] \\
&= \frac{1}{50}\left[350(n^2-2500) - 200n(n-50) + 150(10000-n^2) + 200n(100-n)\right] \\
&= \frac{1}{50}\left[-200n^2 + 30000n + 625000\right] = -4n^2 + 600n + 12500.
\end{aligned}
\]
令 \((EY)' = -8n + 600 = 0\)，得 \(n = 75\)。此时
\(
EY_{\max} = -4 \times 75^2 + 600 \times 75 + 12500 = 35000.
\)
\end{solution}



\begin{exercise}\label{exer:4-basketball}
某决赛采取五局三胜制（三局先胜后比赛终止），双方每局胜率相同。目前比分 1:0（对方领先）。若每场比赛总收入 160 万元，胜方分得 120 万，败方分得 40 万，求对方收入的期望。
\end{exercise}
\begin{solution}
设 \(E(a,b)\) 表示从当前状态出发、对方还差 \(a\) 胜、我方还差 \(b\) 胜时，对方在\textbf{后续比赛}中的收入期望。则有递推关系
\[
E(a,b)=\frac12\bigl(120+E(a-1,b)\bigr)+\frac12\bigl(40+E(a,b-1)\bigr),\qquad E(0,b)=E(a,0)=0.
\]
顺次计算：
\[
E(1,1)=80,\quad E(1,2)=120,\quad E(2,1)=120,\quad E(1,3)=140,\quad E(2,2)=200.
\]
因此
\(
E(2,3)=\frac12(120+140)+\frac12(40+200)=250.
\)
又因为第一局已经打完且对方获胜，所以这 120 万元是确定收入。故对方总收入的期望为
\(
120+250=370\text{（万元）}.
\)
\end{solution}


\begin{exercise}\label{exer:src-paper3-iii-02}\label{prob:p3-iii-02}
设随机变量 \(X\sim B(3,0.1)\)，令
\(
Y=2^X-1.
\)
求 \(EY\)。
\end{exercise}
\begin{solution}
由二项分布的概率母函数，
\(
E(t^X)=(1-p+pt)^n.
\)
这里 \(n=3,\ p=0.1,\ t=2\)，故
\(
E(2^X)=(0.9+0.1\cdot2)^3=1.1^3=1.331.
\)
因此
\(
EY=E(2^X-1)=1.331-1=0.331.
\)
\end{solution}

\begin{exercise}\label{exer:4-s14-cauchymin}
设连续型随机变量 \(X\) 的概率密度为 \(f(x) = \dfrac{1}{\pi(1+x^2)}\ (-\infty < x < +\infty)\)，求 \(E(\min(|X|,1))\)。
\end{exercise}

\begin{solution}
由对称性，
\[
\begin{aligned}
E(\min(|X|,1)) &= 2\int_0^{+\infty}\min(x,1)\,f(x)\,\mathrm{d}x = 2\left[\int_0^1 x f(x)\,\mathrm{d}x + \int_1^{+\infty} f(x)\,\mathrm{d}x\right] \\
&= 2\left(\int_0^1 \frac{x}{\pi(1+x^2)}\,\mathrm{d}x + \int_1^{+\infty}\frac{1}{\pi(1+x^2)}\,\mathrm{d}x\right) \\
&= \frac{1}{\pi}\ln(1+x^2)\Big|_0^1 + \frac{2}{\pi}\arctan x\Big|_1^{+\infty} = \frac{\ln 2}{\pi} + \frac{1}{2}.
\end{aligned}
\]
\end{solution}


\begin{exercise}\label{exer:src-2014-spring-a-05}
设 \(X,Y\) 相互独立，且均服从正态分布 \(N(0,2^2)\)。求随机点 \((X,Y)\) 落在区域
\(
D=\{(x,y)\mid 1\le x^2+y^2\le2\}
\)
内的概率；若 \(Z=\sqrt{X^2+Y^2}\)，求 \(EZ\)。
\end{exercise}
\begin{solution}
由独立性，\((X,Y)\) 的联合密度为
\(
f(x,y)=\frac1{8\pi}\exp\!\left\{-\frac{x^2+y^2}{8}\right\}.
\)
令 \(R=\sqrt{X^2+Y^2}\)，用极坐标可得
\[
\begin{aligned}
P\{(X,Y)\in D\}
&=\int_0^{2\pi}\int_1^{\sqrt2}
\frac1{8\pi}e^{-r^2/8}\,r\,\mathrm dr\,\mathrm d\theta=\frac14\int_1^{\sqrt2}r e^{-r^2/8}\,\mathrm dr
=e^{-1/8}-e^{-1/4}.
\end{aligned}
\]
又 \(R\) 服从参数 \(\sigma=2\) 的瑞利分布，其密度为
\(
f_R(r)=\frac{r}{4}e^{-r^2/8},\qquad r>0.
\)
因此
\[
EZ=\int_0^{+\infty}r\cdot\frac{r}{4}e^{-r^2/8}\,\mathrm dr
=2\sqrt{\frac\pi2}=\sqrt{2\pi}.
\]
\end{solution}



\begin{exercise}\label{exer:4-s9-maxbvn}
已知随机变量 \((X,Y) \sim N(0,0;\,1,1;\,0.5)\)，\(U = \max\{X,Y\}\)，\(V = \min\{X,Y\}\)，求 \(E(U+V)\) 和 \(E(UV)\)。
\end{exercise}
\begin{solution}
根据二维正态分布的参数记号 \((X,Y) \sim N(\mu_1, \mu_2;\, \sigma_1^2, \sigma_2^2;\, \rho)\)，由题意可知：
\begin{itemize}
    \item 期望：$E(X) = 0,\quad E(Y) = 0$
    \item 方差：$D(X) = 1,\quad D(Y) = 1$
    \item 相关系数：$\rho = 0.5$
\end{itemize}
因为 $U = \max\{X,Y\}$，$V = \min\{X,Y\}$，对于任意一对实数取值 $X$ 和 $Y$，无论哪个数较大，$U$ 和 $V$ 必然一个是 $X$ 另一个是 $Y$。因此它们的和与乘积在变换前后保持不变，即有如下代数恒等式：
\(
U + V = X + Y, \quad U \cdot V = X \cdot Y.
\)
\[
E(U+V) = E(X+Y) = E(X) + E(Y) = 0 + 0 = 0.
\]
\[
\mathrm{Cov}(X,Y) = \rho \cdot \sqrt{D(X)} \sqrt{D(Y)} = 0.5 \times \sqrt{1} \times \sqrt{1} = 0.5.
\]
又根据协方差的计算公式 $\mathrm{Cov}(X,Y) = E(XY) - E(X)E(Y)$，可得：
\[
E(UV) = E(XY) = \mathrm{Cov}(X,Y) + E(X)E(Y) = 0.5 + 0 \times 0 = 0.5.
\]
\end{solution}


\begin{exercise}
设随机变量 $X$ 服从参数为 $0.5$ 的泊松分布，试求随机变量 $Y = X/(1+X)$ 的数学期望 $E[Y]$.
\end{exercise}
\begin{solution}
\begin{align*}
E[Y] &= E\left[\frac{X}{1+X}\right] = \sum_{k=0}^{\infty} \frac{k}{(k+1)!}\lambda^k e^{-\lambda} = e^{-\lambda}\sum_{k=0}^{\infty} \left(\frac1{k!} - \frac{1}{(k+1)!}\right)\lambda^k\\
&=e^{-\lambda}\sum_{k=0}^{\infty} \left(\frac{\lambda^k}{k!} - \frac1\lambda\frac{\lambda^{k+1}}{(k+1)!}\right)=e^{-\lambda}(e^{\lambda}-\frac{1}{\lambda}(e^{\lambda}-1))= 1 - \frac{1 - e^{-\lambda}}{\lambda}
\end{align*}
\end{solution}


\subsection{二维随机变量函数的数学期望}

\begin{definition}[二维随机变量函数的期望] \label{def:2d-expectation}
设 \((X,Y)\) 为二维随机变量，\(g(X,Y)\) 为连续函数。
\begin{enumerate}
    \item \textbf{离散型}：若 \((X,Y)\) 的联合分布律为 \(p_{ij} = P\{X=x_i, Y=y_j\}\ (i,j=1,2,\dots)\)，且级数 \(\sum\limits_{i}\sum\limits_{j} |g(x_i,y_j)|\,p_{ij}\) 收敛，则
    \[
    E[g(X,Y)] = \sum_{i=1}^{\infty}\sum_{j=1}^{\infty} g(x_i, y_j)\,p_{ij}.
    \]
    \item \textbf{连续型}：若 \((X,Y)\) 的联合概率密度为 \(f(x,y)\)，且积分 \(\int_{-\infty}^{+\infty}\int_{-\infty}^{+\infty} |g(x,y)|\,f(x,y)\,\mathrm{d}x\mathrm{d}y\) 收敛，则
    \[
    E[g(X,Y)] = \int_{-\infty}^{+\infty}\int_{-\infty}^{+\infty} g(x,y)\,f(x,y)\,\mathrm{d}x\mathrm{d}y.
    \]
\end{enumerate}
若 \(X, Y\) 相互独立，则 \(E(XY) = EX \cdot EY\)。\(E(XY)=EX\cdot EY\)等价于\(\mathrm{Cov}(X,Y)=E(XY)-EX\cdot EY=0.\)这说明 \(X,Y\) \textbf{不相关}，但并不能推出它们独立。
\end{definition}


\begin{exercise}\label{exer:4-6-fx}
设随机变量 \(X\) 的概率密度为 \(f(x) = \begin{cases} \frac12x, & 0 < x < 2, \\ 0, & \text{其他}, \end{cases}\) \(F(x)\) 为 \(X\) 的分布函数，则 \(P\{F(X) > EX - 1\} =\,\underline{\hspace{2cm}}\)。
\end{exercise}

\begin{solution}
\[
EX=\int_0^2 x\cdot \frac12x\,\mathrm dx
=\frac12\int_0^2 x^2\,\mathrm dx
=\frac12\cdot\frac{8}{3}
=\frac43,~~EX-1=\frac43-1=\frac13
\]
对于连续型随机变量，分布函数变换 \(F(X)\) 服从 \(U(0,1)\)，所以令 \(Y=F(X)\)，则 \(Y\sim U(0,1)\)。
\[
P\{F(X)>EX-1\}=P\left(Y>\frac13\right)=1-\frac13=\frac23.
\]
\end{solution}



\begin{exercise}\label{exer:4-8-wait}
甲、乙两人相约于 12:00\textasciitilde 13:00 会面。设 \(X, Y\) 分别是甲、乙到达超过 12:00 的时间（单位：小时），且相互独立，已知
\[
f_X(x) = \begin{cases} 3x^2, & 0 < x < 1, \\ 0, & \text{其他}, \end{cases}\quad
f_Y(y) = \begin{cases} 2y, & 0 < y < 1, \\ 0, & \text{其他}, \end{cases}
\]
求先到达者需要等待的时间的数学期望。
\end{exercise}
\begin{solution}
所求为 \(E(|X-Y|)\)。由独立性，\(f(x,y) = 6x^2y\ (0 < x,y < 1)\)，故
\[
\begin{aligned}
E(|X-Y|) &= \int_0^1\!\!\int_0^1 |x-y| \cdot 6x^2 y\,\mathrm{d}x\mathrm{d}y \\
&= \iint_{x>y} (x-y) \cdot 6x^2 y\,\mathrm{d}x\mathrm{d}y + \iint_{x<y} (y-x) \cdot 6x^2 y\,\mathrm{d}x\mathrm{d}y \\
&= \int_0^1\!\!\int_0^x 6x^3y - 6x^2y^2\,\mathrm{d}y\,\mathrm{d}x + \int_0^1\!\!\int_0^y 6x^2y^2 - 6x^3y\,\mathrm{d}x\,\mathrm{d}y \\
&= \int_0^1 \!\left(3x^5 - 2x^5\right)\mathrm{d}x + \int_0^1 \!\left(2y^5 - \frac{3}{2}y^5\right)\mathrm{d}y = \frac{1}{6} + \frac{1}{12} = \frac{1}{4}.
\end{aligned}
\]
\end{solution}


\begin{exercise}\label{exer:4-two-points}
从长度为 3 的线段上任取两个点，求两点间距离 \(D\) 的数学期望和方差。
\end{exercise}
\begin{solution}
设两个点的位置分别为 \(X_1, X_2\)，独立且均服从 \(U(0,3)\)，则距离 \(D = |X_1 - X_2|\)。单个均匀分布 \( U(0,3) \) 的概率密度为：
\[f_{X_i}(x_i) =\begin{cases}\displaystyle \frac{1}{3}, & 0 < x_i < 3 \\0, & \text{其他}\end{cases}\quad (i=1,2)\]
由独立性，联合密度 = 边缘密度的乘积：
  \[
  f(x_1,x_2) = f_{X_1}(x_1) \cdot f_{X_2}(x_2) =
  \begin{cases}
  \displaystyle \frac{1}{3} \cdot \frac{1}{3} = \frac{1}{9}, & 0 < x_1 < 3,\ 0 < x_2 < 3 \\
  0, & \text{其他}
  \end{cases}
  \]
由对称性，只需计算 \(x_1 > x_2\) 的区域再乘以 2：
\[
ED = 2\int_0^3 \int_0^{x_1} (x_1 - x_2) \cdot \frac{1}{9}\,\mathrm{d}x_2\,\mathrm{d}x_1
= \frac{2}{9}\int_0^3 \frac{x_1^2}{2}\,\mathrm{d}x_1 = 1.
\]
\(
E(X_i) = \frac{3}{2},\quad E(X_i^2) = DX_i + (EX_i)^2 = \frac{3}{4} + \frac{9}{4} = 3.
\)
\(
E(D^2) = E[(X_1 - X_2)^2] = E(X_1^2) - 2EX_1 EX_2 + E(X_2^2) = 3 - 2 \times \frac{9}{4} + 3 = \frac{3}{2},
\)
\(
DD = ED^2 - (ED)^2 = \frac{3}{2} - 1 = \frac{1}{2}.
\)
\end{solution}

\begin{proposition}[均匀分布次序统计量的通用结论]\label{prop:order-stat-uniform}
设相互独立的随机变量 \(X_1, X_2, \dots, X_n\) 均服从 \((0, \theta)\) 上的均匀分布。记次序统计量为
\[
X_{(1)}\le X_{(2)}\le \cdots \le X_{(n)}.
\]
则第 \(k\) 个次序统计量的期望为
\[
E(X_{(k)})=\frac{k}{n+1}\theta,\qquad k=1,2,\dots,n.
\]
特别地，
\[
E(X_{(1)})=\frac{\theta}{n+1},\qquad
E(X_{(n)})=\frac{n\theta}{n+1},
\]
所以极差 \(R=X_{(n)}-X_{(1)}\) 的期望为
\[
E(R)=E(X_{(n)})-E(X_{(1)})=\frac{n-1}{n+1}\theta.
\]
\end{proposition}

\begin{proof}
先证 \(U(0,1)\) 的情形。设 \(U_1,\dots,U_n\) 相互独立且都服从 \(U(0,1)\)，其第 \(k\) 个次序统计量为 \(U_{(k)}\)。要使 \(U_{(k)}\) 落在 \(u\) 附近，需要恰有 \(k-1\) 个样本落在 \(u\) 左边、恰有 \(n-k\) 个样本落在 \(u\) 右边，再有一个样本落在 \([u,u+\mathrm{d}u]\) 内，故
\[
f_{U_{(k)}}(u)=\frac{n!}{(k-1)!(n-k)!}u^{k-1}(1-u)^{n-k},\qquad 0<u<1.
\]
因此
\[
\begin{aligned}
E(U_{(k)})
&=\frac{n!}{(k-1)!(n-k)!}\int_0^1 u^k(1-u)^{n-k}\,\mathrm{d}u=\frac{n!}{(k-1)!(n-k)!}\cdot \frac{k!(n-k)!}{(n+1)!}
=\frac{k}{n+1}.
\end{aligned}
\]
若 \(X_i\sim U(0,\theta)\)，令 \(U_i=X_i/\theta\)，则 \(U_i\sim U(0,1)\)，且 \(X_{(k)}=\theta U_{(k)}\)。所以
\[
E(X_{(k)})=\theta E(U_{(k)})=\frac{k}{n+1}\theta.
\]
\end{proof}


\begin{exercise}
设随机变量 $X_1, X_2, X_3$ 独立同分布，且 $X_i (i=1,2,3)$ 的分布列为：$P(X_i = k) = \frac{1}{3} \quad (k=1,2,3)$，求 $Y = \max\{X_1, X_2, X_3\}$ 的数学期望.
\end{exercise}
\begin{solution}
由于 $X_1, X_2, X_3$ 独立同分布，我们可以先求出 $Y$ 的分布函数（即累积概率）：
$$P(Y \le k) = P(X_1 \le k, X_2 \le k, X_3 \le k) = P(X_1 \le k)P(X_2 \le k)P(X_3 \le k) = [P(X_1 \le k)]^3$$
根据已知条件，$X_i$ 的分布列为 $P(X_i = k) = \frac{1}{3}$（$k=1,2,3$），我们可以算出 $X_i$ 的累积概率：
\begin{itemize}
\item 当 $k=1$ 时，$P(X_1 \le 1) = P(X_1 = 1) = \frac{1}{3}$
\item 当 $k=2$ 时，$P(X_1 \le 2) = P(X_1 = 1) + P(X_1 = 2) = \frac{2}{3}$
\item 当 $k=3$ 时，$P(X_1 \le 3) = P(X_1 = 1) + P(X_1 = 2) + P(X_1 = 3) = 1$
\end{itemize}
\begin{align*}
E[Y] &= \sum_{k=1}^{3} k \cdot P(Y = k) = 1 \cdot \frac{1}{27} + 2 \cdot \frac{7}{27} + 3 \cdot \frac{19}{27} = \frac{8}{3}
\end{align*}
\end{solution}

\section{方差，协方差，相关系数}

\begin{definition}[方差] \label{def:variance}
从方差的定义可以看出，只有当随机变量的期望存在时，其方差才可能存在。但期望存在的随机变量，其方差也不一定存在。设 \(X\) 是随机变量，若 \(E[(X - EX)^2]\) 存在，则称其为 \(X\) 的\textbf{方差}，记为 \(DX\)（或 \(\mathrm{Var}(X)\)），即
\(
DX = E[(X - EX)^2] = E(X^2) - (EX)^2.
\)
称 \(\sqrt{DX}\) 为 \(X\) 的\textbf{标准差}（或均方差），记为 \(\sigma(X)\)。

若 \(DX > 0\)，称 \(X^* = \dfrac{X - EX}{\sqrt{DX}}\) 为 \(X\) 的\textbf{标准化随机变量}，此时 \(EX^* = 0\)，\(DX^* = 1\)。
\end{definition}

\begin{property}\label{prop:variance}
方差具有以下性质（假设所涉及的方差均存在）：
\begin{enumerate}
    \item 对任意两个随机变量 \(X, Y\)，有
    \[
    D(X \pm Y) = DX + DY \pm 2\,\mathrm{Cov}(X,Y).
    \]
    更一般地，
    \[
    D\!\left(\sum_{i=1}^{n} a_i X_i\right) = \sum_{i=1}^{n} a_i^2 DX_i + 2\sum_{1 \leqslant i < j \leqslant n} a_i a_j\,\mathrm{Cov}(X_i,X_j).
    \]
    \item 若 \(X, Y\) 相互独立，则 \(D(X \pm Y) = DX + DY\)，且
\[
\begin{aligned}
D(XY) &= E(X^2 Y^2) - [E(XY)]^2 = E(X^2)E(Y^2) - (EX)^2(EY)^2 \\
&= [DX + (EX)^2][DY + (EY)^2] - (EX)^2(EY)^2 \\
&= DX \cdot DY + (EX)^2 \cdot DY + (EY)^2 \cdot DX.
\end{aligned}
\]
\end{enumerate}
\end{property}

\begin{proposition}[均值的期望与方差]\label{prop:sample-mean-moment}
    若总体 $X$ 的期望为 $\mu$，方差为 $\sigma^2$，$X_1,\dots,X_n$ 是简单随机样本，样本均值为 $\overline{X}$，则：$$\mathbf{E(\overline{X}) = \mu}, \quad \mathbf{D(\overline{X}) = \frac{\sigma^2}{n}}$$
\end{proposition}
\begin{note}
这条结论只要求样本是独立同分布、总体具有有限方差，并不要求总体正态。它常用来说明样本均值是 \(\mu\) 的无偏估计，且样本量越大，样本均值的波动越小；但若题目进一步要求 \(\overline{X}\) 的精确分布，就还需要额外条件，不能只凭这两个矩结论直接下结论。
\end{note}
\begin{proof} 
    \[ E(\overline{X}) = E\!\left(\frac{1}{n}\sum_{i=1}^{n}X_i\right) = \frac{1}{n}\sum_{i=1}^{n}EX_i = \frac{1}{n}\cdot n\mu = \mu. \] 
    由于简单随机样本中 \(X_1, \dots, X_n\) 相互独立，故 \[ D(\overline{X}) = D\!\left(\frac{1}{n}\sum_{i=1}^{n}X_i\right) = \frac{1}{n^2}\sum_{i=1}^{n}DX_i = \frac{1}{n^2}\cdot n\sigma^2 = \frac{\sigma^2}{n}. \] 
\end{proof}

\begin{definition}[协方差] \label{def:covariance}
设随机变量 \(X\) 与 \(Y\) 的方差都存在，称
\(
\mathrm{Cov}(X,Y) = E[(X - EX)(Y - EY)] = E(XY) - EX \cdot EY
\)
为 \(X\) 与 \(Y\) 的\textbf{协方差}。其中
\[
E(XY) = \begin{cases}
\sum\limits_{i}\sum\limits_{j} x_i y_j P\{X=x_i,Y=y_j\} & \text{（离散型）}, \\[8pt]
\int_{-\infty}^{+\infty}\int_{-\infty}^{+\infty} xy\,f(x,y)\,\mathrm{d}x\mathrm{d}y & \text{（连续型）}.
\end{cases}
\]
协方差具有以下性质：\(\mathrm{Cov}(aX, bY) = ab\,\mathrm{Cov}(X,Y)\)。
    \[
    \mathrm{Cov}\!\left(\sum_{i=1}^{n} a_i X_i,\ \sum_{j=1}^{m} b_j Y_j\right) = \sum_{i=1}^{n}\sum_{j=1}^{m} a_i b_j\,\mathrm{Cov}(X_i, Y_j).
    \]
    二元情形：\[\operatorname{Cov}(aX+bY,cU+dV)=ac\operatorname{Cov}(X,U)+ad\operatorname{Cov}(X,V)+bc\operatorname{Cov}(Y,U)+bd\operatorname{Cov}(Y,V).\]
\end{definition}


\begin{definition}[相关系数] \label{def:correlation}
设 \(DX > 0\)，\(DY > 0\)，称
\[
\rho_{XY} = \frac{\mathrm{Cov}(X,Y)}{\sqrt{DX \cdot DY}}
\]
为 \(X\) 与 \(Y\) 的\textbf{相关系数}。若 \(\rho_{XY} = 0\)，则称 \(X\) 与 \(Y\)\textbf{不相关}；若 \(\rho_{XY} > 0\)，称 \(X\) 与 \(Y\)\textbf{正相关}；若 \(\rho_{XY} < 0\)，称 \(X\) 与 \(Y\)\textbf{负相关}。
\(\rho_{XY} = 0\)（只能说明“没有线性关系”，不能说明“没有关系”）的等价条件：
    \begin{itemize}
        \item \(\mathrm{Cov}(X,Y) = 0\)；
        \item \(E(XY) = EX \cdot EY\)；
        \item \(D(X + Y) = DX + DY\)。
    \end{itemize}
独立性与不相关性的关系：
\begin{itemize}
    \item 独立 $\Rightarrow$ 不相关 \checkmark（永远成立）
    \item 不相关 $\Rightarrow$ 独立 $\times$（一般不成立）
    \item 不相关 + 二维正态 $\Rightarrow$ 独立 \checkmark
    \item 不相关 + 各自正态（但未说联合正态）$\Rightarrow$ 独立 $\times$
\end{itemize}
\end{definition}

\begin{exercise}\label{exer:4-4-normal}
设连续型随机变量 \(X\) 的概率密度为
\[
f(x) = \frac{1}{\sqrt{\pi}}\,e^{-x^2 + 2x - 1}, \quad -\infty < x < +\infty,
\]
则 \(EX \cdot \sqrt{DX} =\,\underline{\hspace{2cm}}\)。
\end{exercise}
\begin{solution}
\[
f(x) = \frac{1}{\sqrt{\pi}}\,e^{-x^2+2x-1} = \frac{1}{\sqrt{\pi}}\,e^{-(x-1)^2} = \frac{1}{\sqrt{2\pi}\cdot\frac{1}{\sqrt{2}}}\exp\!\left\{-\frac{(x-1)^2}{2\cdot\left(\frac{1}{\sqrt{2}}\right)^2}\right\}.
\]
即 $X \sim N\!\left(1,\dfrac{1}{2}\right)$。$EX = 1$，$\sqrt{DX} = \dfrac{1}{\sqrt{2}}$，故 $EX \cdot \sqrt{DX} = \dfrac{1}{\sqrt{2}}$。
\end{solution}

\begin{exercise}\label{exer:4-7-square}
设 \(X, Y\) 是随机变量，相关系数为 0.4，\(EX = 2\)，\(DX = 25\)，\(EY = 1\)，\(DY = 16\)，则 \(E(2X - 3Y + 4)^2 =\,\underline{\hspace{2cm}}\)。
\end{exercise}

\begin{solution}
先计算各分量：
\(
EX^2 = DX + (EX)^2 = 29,\quad EY^2 = DY + (EY)^2 = 17,
\)
\[
EXY = EX \cdot EY + \mathrm{Cov}(X,Y) = 2 + 0.4 \times 5 \times 4 = 10.
\]
于是
\[
\begin{aligned}
E(2X-3Y+4)^2 &= E(4X^2 + 9Y^2 - 12XY + 16 + 16X - 24Y) \\
&= 4 \times 29 + 9 \times 17 - 12 \times 10 + 16 + 16 \times 2 - 24 \times 1 \\
&= 116 + 153 - 120 + 16 + 32 - 24 = 173.
\end{aligned}
\]
\end{solution}


\begin{exercise}\label{exer:4-3-abs}
设随机变量 \((X,Y)\) 服从区域 \(G = \{(x,y) \mid 0 \leqslant x \leqslant 2,\ 0 \leqslant y \leqslant 2\}\) 上的均匀分布，\(U = |X-Y|\)。\par
(1) 求 \(U\) 的概率密度 \(f_U(u)\)；\par
(2) 求 \(EU\) 和 \(DU\)。
\end{exercise}
\begin{solution}
(1) \(F_U(u) = P(|X-Y| \leqslant u)\)。当 \(u \leqslant 0\) 时，\(F_U(u) = 0\)；当 \(u \geqslant 2\) 时，\(F_U(u) = 1\)。\\
当 \(0 < u < 2\) 时，区域 \(|x-y| \leqslant u\) 在 \([0,2]^2\) 内的面积为 \(4 - 2 \times \dfrac{1}{2}(2-u)^2 = 4u - u^2\)，故
\(
F_U(u) = \frac{4u - u^2}{4} = \frac{(4-u)u}{4}.
\)
因此
\[
f_U(u) = F_U'(u) = \begin{cases} \dfrac{2-u}{2}, & 0 < u < 2, \\ 0, & \text{其他}. \end{cases}
\]
(2) \(EU = \displaystyle\int_0^2 u \cdot \frac{2-u}{2}\,\mathrm{d}u = \frac{2}{3}\)；
\(EU^2 = \displaystyle\int_0^2 u^2 \cdot \frac{2-u}{2}\,\mathrm{d}u = \frac{2}{3}\)；
\(DU = EU^2 - (EU)^2 = \dfrac{2}{3} - \dfrac{4}{9} = \dfrac{2}{9}\)。
\end{solution}


\begin{exercise}\label{exer:4-sample-cov}
设 \(X_1, X_2, \dots, X_5\) 为总体 \(N(1,4)\) 的简单随机样本，\(X_{(5)} = \max\{X_1, X_2, \dots, X_5\}\)。求：(1) \(P(X_{(5)} < 3)\)；(2) \(\mathrm{Cov}(X_1, \bar{X})\)，其中 \(\bar{X} = \dfrac{1}{5}\sum\limits_{i=1}^{5}X_i\)。
\end{exercise}
\begin{solution}
(1) \(P(X_i < 3) = \Phi\!\left(\dfrac{3-1}{2}\right) = \Phi(1)\)，由样本独立性
\(
P(X_{(5)} < 3) = [P(X_1 < 3)]^5 = [\Phi(1)]^5.
\)
(2) 先把样本均值展开：
\(
\bar X=\frac15(X_1+X_2+X_3+X_4+X_5).
\)
于是
\[
\mathrm{Cov}(X_1,\bar X)
=\mathrm{Cov}\!\left(X_1,\frac15\sum_{i=1}^{5}X_i\right)
=\frac15\sum_{i=1}^{5}\mathrm{Cov}(X_1,X_i).
\]
这里最关键的一步是：由于简单随机样本中的 \(X_1,\dots,X_5\) 相互独立，所以
\(
\mathrm{Cov}(X_1,X_i)=0\qquad (i=2,3,4,5),
\)
只剩下
\(
\mathrm{Cov}(X_1,X_1)=DX_1.
\)
因此
\(
\mathrm{Cov}(X_1,\bar X)=\frac15\,DX_1.
\)
又因为总体 \(N(1,4)\) 的方差为 4，所以 \(DX_1=4\)。故
\(
\mathrm{Cov}(X_1,\bar X)=\frac45.
\)
\end{solution}


\begin{exercise}[多项分布的相关系数] 
袋中有红、白、黑三种颜色球若干，若从袋中任摸一球，摸出的球为红球的概率为 $p_1$，摸出的球为白球的概率为 $p_2$。现从袋中有放回地摸球 $n$ 次，共摸出红球 $X$ 次，摸出白球 $Y$ 次，试求 $X$ 与 $Y$ 的相关系数 $r(X,Y)$。
\end{exercise}

\begin{solution}
由于是有放回抽样，每次摸球是相互独立的，且 $X$ 和 $Y$ 分别服从二项分布：
\(
X \sim B(n, p_1), \quad Y \sim B(n, p_2)
\)，
因此它们的方差分别为：
\(
D(X) = np_1(1-p_1), \quad D(Y) = np_2(1-p_2)
\)，
为了求 $X$ 与 $Y$ 的协方差 $\text{Cov}(X,Y)$，我们引入示性变量（0-1 变量）：
设 
\[
X_i = \begin{cases}
1, & \text{第 } i \text{ 次摸出红球}, \\
0, & \text{其他},
\end{cases}
\qquad
Y_i = \begin{cases}
1, & \text{第 } i \text{ 次摸出白球}, \\
0, & \text{其他},
\end{cases}
\]
其中 $i = 1, 2, \cdots, n$。则 $X = \sum\limits_{i=1}^n X_i$，$Y = \sum\limits_{j=1}^n Y_j$。
利用协方差的双线性性质：
\[
\text{Cov}(X, Y) = \text{Cov}\left(\sum_{i=1}^n X_i, \sum_{j=1}^n Y_j\right) = \sum_{i=1}^n \sum_{j=1}^n \text{Cov}(X_i, Y_j)
\]
分两种情况讨论 $\text{Cov}(X_i, Y_j)$：
\begin{enumerate}
    \item[(1)] 当 $i \neq j$ 时，第 $i$ 次与第 $j$ 次摸球相互独立，因此 $X_i$ 与 $Y_j$ 独立，$\text{Cov}(X_i, Y_j) = 0$。
    \item[(2)] 当 $i = j$ 时，由于同一次摸球不可能既是红球又是白球，故恒有 $X_i Y_i = 0$，从而 $E(X_i Y_i) = 0$。
    此时：
    \(
    \text{Cov}(X_i, Y_i) = E(X_i Y_i) - E(X_i)E(Y_i) = 0 - p_1 p_2 = -p_1 p_2
    \)
\end{enumerate}
因此，协方差为：
\[
\text{Cov}(X,Y) = \sum_{i=1}^n \text{Cov}(X_i, Y_i) = \sum_{i=1}^n (-p_1 p_2) = -n p_1 p_2
\]
最后，求 $X$ 与 $Y$ 的相关系数：
\begin{align*}
r(X,Y) &= \frac{\text{Cov}(X,Y)}{\sqrt{D(X)} \sqrt{D(Y)}} = \frac{-n p_1 p_2}{\sqrt{np_1(1-p_1)} \sqrt{np_2(1-p_2)}} = -\sqrt{\frac{p_1 p_2}{(1-p_1)(1-p_2)}}
\end{align*}
\end{solution}

\begin{exercise}\label{exer:src-2010-0630-a-04}
设二维随机变量 \((X,Y)\) 的概率密度为
\[
f(x,y)=
\begin{cases}
A(x+y), & 0\le x\le1,\ 0\le y\le2,\\
0, & \text{其他}.
\end{cases}
\]
求：\par
(1) 常数 \(A\)；\par
(2) \(E(X)\)、\(\operatorname{Cov}(X,Y)\) 以及 \(X,Y\) 的相关系数；\par
(3) \((X,Y)\) 落入区域
\(
D=\{0\le x\le1,\ y\ge x^2\}
\)
的概率。
\end{exercise}
\begin{solution}
(1) 由归一化条件，
\[
1=A\int_0^1\int_0^2(x+y)\,\mathrm dy\,\mathrm dx
=A\int_0^1(2x+2)\,\mathrm dx=3A,
\]
故
\(
A=\frac13.
\)

(2) 先求所需数字特征：
\[
E(X)=\frac13\int_0^1\int_0^2 x(x+y)\,\mathrm dy\,\mathrm dx
=\frac13\int_0^1(2x^2+2x)\,\mathrm dx
=\frac59.
\]
同理
\[
E(Y)=\frac13\int_0^1\int_0^2 y(x+y)\,\mathrm dy\,\mathrm dx
=\frac{11}{9},
\]
\[
E(XY)=\frac13\int_0^1\int_0^2 xy(x+y)\,\mathrm dy\,\mathrm dx
=\frac23.
\]
因此
\[
\operatorname{Cov}(X,Y)=E(XY)-E(X)E(Y)
=\frac23-\frac59\cdot\frac{11}{9}
=-\frac1{81}.
\]

为求相关系数，再算
\[
E(X^2)=\frac13\int_0^1\int_0^2 x^2(x+y)\,\mathrm dy\,\mathrm dx
=\frac7{18},
\]
\(
D(X)=E(X^2)-[E(X)]^2=\frac7{18}-\frac{25}{81}=\frac{13}{162}.
\)
又
\[
E(Y^2)=\frac13\int_0^1\int_0^2 y^2(x+y)\,\mathrm dy\,\mathrm dx
=\frac{16}{9},
\]
\(
D(Y)=E(Y^2)-[E(Y)]^2=\frac{16}{9}-\frac{121}{81}=\frac{23}{81}.
\)
所以相关系数为
\[
\rho_{XY}
=\frac{\operatorname{Cov}(X,Y)}{\sqrt{D(X)}\sqrt{D(Y)}}
=\frac{-1/81}{\sqrt{(13/162)(23/81)}}
=-\sqrt{\frac{2}{299}}.
\]

(3) 在支撑区域内，\(D\) 的补集为 \(0\le x\le1,\ 0\le y<x^2\)。先算补集概率：
\[
\begin{aligned}
P(Y<x^2)
&=\frac13\int_0^1\int_0^{x^2}(x+y)\,\mathrm dy\,\mathrm dx=\frac13\int_0^1\left(x^3+\frac{x^4}{2}\right)\mathrm dx
=\frac13\left(\frac14+\frac1{10}\right)
=\frac7{60}.
\end{aligned}
\]
故
\(
P\{(X,Y)\in D\}=1-\frac7{60}=\frac{53}{60}.
\)
\end{solution}


\begin{exercise}\label{exer:src-2011-ab-a-03}
随机变量 \(\xi\sim N(0,4)\)，令 \(\eta=2^\xi\)。求：\par
(1) \(\eta\) 的概率密度函数 \(f_\eta(y)\)；\par
(2) \(E\eta\)；\par
(3) \(D\eta\)。
\end{exercise}
\begin{solution}
(1) 因为 \(\eta=2^\xi>0\)，当 \(y>0\) 时，
\(
F_\eta(y)=P(2^\xi\le y)
=P\left(\xi\le\frac{\ln y}{\ln2}\right).
\)
对 \(y\) 求导，得
\[
f_\eta(y)
=\frac{1}{2\sqrt{2\pi}\,y\ln2}
\exp\left\{-\frac{(\ln y)^2}{8\ln^2 2}\right\},
\qquad y>0,
\]
而 \(y\le0\) 时 \(f_\eta(y)=0\)。

(2) 由于 \(2^\xi=e^{(\ln2)\xi}\)，正态变量的矩母函数给出
\[
E\eta=E e^{(\ln2)\xi}
=\exp\left\{\frac12\cdot4\ln^2 2\right\}
=e^{2\ln^2 2}=2^{2\ln2}.
\]

(3)
\[
E\eta^2=E e^{2(\ln2)\xi}
=\exp\left\{\frac12\cdot4\cdot(2\ln2)^2\right\}
=e^{8\ln^2 2}.
\]
因此
\(
D\eta=e^{8\ln^2 2}-e^{4\ln^2 2}
=2^{4\ln2}\left(2^{4\ln2}-1\right).
\)
\end{solution}


\begin{exercise}\label{exer:4-4-sym}
设随机变量 \(X\) 的概率密度为 \(f(x)\)，\(DX = 2\)，而随机变量 \(Y\) 的概率密度为 \(f(-y)\)，它们的相关系数 \(\rho_{XY} = \dfrac{1}{4}\)，记 \(Z = X + Y\)，求 \(EZ\)，\(DZ\)。
\end{exercise}
\begin{solution}
这道题的关键是先读懂密度 \(f(-y)\) 的含义：它表示 \(Y\) 的分布是把 \(X\) 关于原点镜像后的结果，因此 \(Y\) 与 \(-X\) 同分布。
所以\(EZ=EX+EY=EX-EX=0\)，又因为 \(EY=-EX\)，故\(DY=EY^2-(EY)^2=EX^2-(EX)^2=DX=2.\)。题目给的是相关系数，不是协方差，所以还要先换回来：
\[
\mathrm{Cov}(X,Y)=\rho_{XY}\sqrt{DX}\sqrt{DY}=\frac14\times \sqrt{2}\times \sqrt{2}=\frac12.
\]
最后利用 \(D(X+Y)=DX+DY+2\,\mathrm{Cov}(X,Y)\)，得到 \(DZ=2+2+2\times \frac12=5\)。因此 \(EZ=0,\ DZ=5\)。
\end{solution}


\begin{exercise}\label{exer:src-2009-a-09}
设 \((X,Y)\) 的密度函数为
\[
f(x,y)=
\begin{cases}
Axy, & 0<x<1,\ 0<y<1,\\
0, & \text{其他}.
\end{cases}
\]
求：\par
(1) 常数 \(A\)；\par
(2) \(P(X<0.4,\ Y<1.3)\)；\par
(3) \(E(e^{tX+sY})\)；\par
(4) \(EX,\ DX,\ \operatorname{Cov}(X,Y)\)。
\end{exercise}
\begin{solution}
(1) 由归一化条件，
\[
1=\int_0^1\int_0^1 Axy\,\mathrm{d}y\,\mathrm{d}x
=A\left(\int_0^1 x\,\mathrm{d}x\right)\left(\int_0^1 y\,\mathrm{d}y\right)
=\frac{A}{4}.
\]
故
\(
A=4.
\)

(2) 因为 \(Y<1.3\) 在支撑 \(0<Y<1\) 上恒成立，所以
\[
P(X<0.4,Y<1.3)=\int_0^{0.4}\int_0^1 4xy\,\mathrm{d}y\,\mathrm{d}x
=\int_0^{0.4}2x\,\mathrm{d}x=0.16.
\]

(3)
\[
\begin{aligned}
E(e^{tX+sY})
&=\int_0^1\int_0^1 e^{tx+sy}4xy\,\mathrm{d}y\,\mathrm{d}x=4\left(\int_0^1 xe^{tx}\,\mathrm{d}x\right)
\left(\int_0^1 ye^{sy}\,\mathrm{d}y\right).
\end{aligned}
\]
当 \(t\ne0,\ s\ne0\) 时，
\(
\int_0^1 xe^{tx}\,\mathrm{d}x=\frac{e^t}{t}-\frac{e^t}{t^2}+\frac1{t^2},
\)
\(
\int_0^1 ye^{sy}\,\mathrm{d}y=\frac{e^s}{s}-\frac{e^s}{s^2}+\frac1{s^2}.
\)
故
\[
E(e^{tX+sY})
=4\left(\frac{e^t}{t}-\frac{e^t}{t^2}+\frac1{t^2}\right)
\left(\frac{e^s}{s}-\frac{e^s}{s^2}+\frac1{s^2}\right),
\]
其中 \(t=0\) 或 \(s=0\) 的情形可由极限得到。

(4) 由密度可分解为
\(
f(x,y)=2x\cdot2y,\qquad 0<x<1,\ 0<y<1,
\)
故 \(X,Y\) 相互独立，且边缘密度分别为 \(f_X(x)=2x,\ f_Y(y)=2y\)。于是
\(
EX=\int_0^1 2x^2\,\mathrm{d}x=\frac23,
\)
\(
EX^2=\int_0^1 2x^3\,\mathrm{d}x=\frac12,
\)
\(
DX=EX^2-(EX)^2=\frac12-\frac49=\frac1{18}.
\)
又 \(X,Y\) 独立，所以
\(
\operatorname{Cov}(X,Y)=0.
\)
\end{solution}



\begin{exercise}\label{exer:4-19-interval}
在区间 \([0.5, 1]\) 中任取一个值 \(x\)，再在区间 \([-x, x]\) 中任取一个值 \(y\)，构成二维随机变量 \((X,Y)\)。\par
(1) 求 \((X,Y)\) 的联合密度函数 \(f(x,y)\)，关于 \(Y\) 的边缘密度函数 \(f_Y(y)\)，并判断 \(X, Y\) 是否独立；\par
(2) 求 \(\mathrm{Cov}(X,Y)\)，并据此判断 \(X, Y\) 是否相关。
\end{exercise}

\begin{solution}
(1) \((X,Y)\) 的取值区域为 \(0.5 < x < 1\)，\(-x < y < x\)。

\(X \sim U(0.5, 1)\)，故 \(f_X(x) = \dfrac{1}{0.5} = 2\ (0.5 < x < 1)\)。给定 \(X = x\)，\(Y \sim U(-x, x)\)，故 \(f(y|x) = \dfrac{1}{2x}\ (-x < y < x)\)。

联合密度函数：
\(
f(x,y) = f_X(x)\,f(y|x) = 2 \times \frac{1}{2x} = \frac{1}{x},\quad 0.5 < x < 1,\ -x < y < x.
\)

关于 \(Y\) 的边缘密度函数。对固定的 \(y\)，\(x\) 的范围由 \(\max(0.5, |y|) < x < 1\) 确定：
\[
f_Y(y) = \begin{cases}
\displaystyle\int_y^1 \frac{1}{x}\,\mathrm{d}x = -\ln y, & 0.5 \leqslant y < 1, \\[8pt]
\displaystyle\int_{0.5}^1 \frac{1}{x}\,\mathrm{d}x = \ln 2, & -0.5 \leqslant y < 0.5, \\[8pt]
\displaystyle\int_{-y}^1 \frac{1}{x}\,\mathrm{d}x = -\ln(-y), & -1 \leqslant y < -0.5.
\end{cases}
\]

\(f_X(x) = \displaystyle\int_{-x}^{x} \frac{1}{x}\,\mathrm{d}y = 2\)。由于 \(f(x,y) = \dfrac{1}{x} \neq 2\,f_Y(y) = f_X(x)\,f_Y(y)\)，故 \(X\) 与 \(Y\) 不独立。

(2) \(EX = \displaystyle\int_{0.5}^{1} 2x\,\mathrm{d}x = \dfrac{3}{4}\)。由对称性（\(f_Y(y)\) 关于 \(y = 0\) 对称），\(EY = 0\)。

\[
\mathrm{Cov}(X,Y) = EXY - EX \cdot EY = EXY = \int_{0.5}^{1}\int_{-x}^{x} xy \cdot \frac{1}{x}\,\mathrm{d}y\,\mathrm{d}x = \int_{0.5}^{1} x\!\left[\int_{-x}^{x} y\,\mathrm{d}y\right]\mathrm{d}x = 0.
\]

所以 \(X, Y\) 不相关（但由 (1) 知二者不独立）。
\end{solution}

\section{正态分布的数字特征}



\begin{exercise}\label{exer:4-15-absnormal}
设两个随机变量 \(X, Y\) 相互独立，且都服从均值为 0、方差为 0.5 的正态分布，求 \(|X-Y|\) 的数学期望与方差。
\end{exercise}

\begin{solution}
设 \(Z=X-Y\)。由于 \(X,Y\) 相互独立，且都服从均值为 0、方差为 0.5 的正态分布，所以线性组合 \(Z\) 仍服从正态分布，并且 \(EZ=EX-EY=0,\ DZ=DX+DY=0.5+0.5=1\)，因此 \(Z\sim N(0,1)\)。

题目要求的是 \(|X-Y|=|Z|\) 的数学期望与方差，先求 \(E|Z|\)：
\[
E(|X-Y|)=E|Z|=\int_{-\infty}^{+\infty}|z|\,\frac{1}{\sqrt{2\pi}}\,e^{-\frac{z^2}{2}}\,\mathrm{d}z=
E|Z|=\frac{2}{\sqrt{2\pi}}\int_0^{+\infty} z\,e^{-\frac{z^2}{2}}\,\mathrm{d}z
=\sqrt{\frac{2}{\pi}}.
\]

再求方差。注意
\(
|Z|^2=Z^2,
\)
所以
\(
E(|X-Y|^2)=E(|Z|^2)=E(Z^2)=DZ+(EZ)^2=1+0=1.
\)
因此
\(
D(|X-Y|)=D(|Z|)=E(|Z|^2)-[E|Z|]^2=1-\frac{2}{\pi}.
\)

所以
\(
E(|X-Y|)=\sqrt{\frac{2}{\pi}},\qquad D(|X-Y|)=1-\frac{2}{\pi}.
\)
\end{solution}

\begin{exercise}\label{exer:src-2010-0630-a-06}
设随机变量 \(\xi\) 与 \(\eta\) 相互独立，并且有相同的分布 \(N(a,\sigma^2)\)。证明：
\(
E\max\{\xi,\eta\}=a+\frac{\sigma}{\sqrt{\pi}}.
\)
\end{exercise}
\begin{solution}
令
\(
X=\frac{\xi-a}{\sigma},\qquad Y=\frac{\eta-a}{\sigma}.
\)
则 \(X,Y\) 相互独立且均服从标准正态分布，并且
\(
\max\{\xi,\eta\}=a+\sigma\max\{X,Y\}.
\)
因此只需证明 \(E\max\{X,Y\}=1/\sqrt{\pi}\)。

利用恒等式
\(
\max\{X,Y\}=\frac{X+Y+|X-Y|}{2},
\)
得
\(
E\max\{X,Y\}
=\frac12E(X+Y)+\frac12E|X-Y|
=\frac12E|X-Y|.
\)
由于
\(
X-Y\sim N(0,2),
\)
若记 \(Z\sim N(0,1)\)，则 \(X-Y\) 与 \(\sqrt2\,Z\) 同分布。于是
\(
E|X-Y|=\sqrt2\,E|Z|.
\)
又
\[
E|Z|=2\int_0^{+\infty}z\frac1{\sqrt{2\pi}}e^{-z^2/2}\,\mathrm dz
=\sqrt{\frac2\pi}.
\]
所以
\[
E|X-Y|=\sqrt2\sqrt{\frac2\pi}=\frac2{\sqrt{\pi}},
\]
从而
\(
E\max\{X,Y\}=\frac1{\sqrt{\pi}}.
\)
代回原变量，得到
\(
E\max\{\xi,\eta\}=a+\frac{\sigma}{\sqrt{\pi}}.
\)
\end{solution}



\begin{exercise}\label{exer:src-2008-probstat-2a-i-03}\label{exer:src-paper3-ii-03}\label{prob:p3-ii-03}
设随机变量 \(X_1,X_2,\dots,X_n\) 独立同分布，且方差为 \(\sigma^2>0\)。令
\(
Y=\frac1n\sum_{i=1}^n X_i.
\)
则下列结论正确的是（\quad）。\par
A. \(\operatorname{Cov}(X_1,Y)=\dfrac{\sigma^2}{n}\)\quad
B. \(\operatorname{Cov}(X_1,Y)=\sigma^2\)\par
C. \(D(X_1+Y)=\dfrac{(n+2)\sigma^2}{n}\)\quad
D. \(D(X_1-Y)=\dfrac{(n+1)\sigma^2}{n}\)
\end{exercise}
\begin{solution}
应选 A。

由协方差的线性性质和独立性，
\[
\operatorname{Cov}(X_1,Y)
=\operatorname{Cov}\!\left(X_1,\frac1n\sum_{i=1}^nX_i\right)
=\frac1n\operatorname{Cov}(X_1,X_1)
=\frac{\sigma^2}{n}.
\]
因此 A 正确，B 错误。

另外，
\(
DY=\frac{\sigma^2}{n}.
\)
所以
\[
D(X_1+Y)=DX_1+DY+2\operatorname{Cov}(X_1,Y)
=\sigma^2+\frac{\sigma^2}{n}+\frac{2\sigma^2}{n}
=\frac{(n+3)\sigma^2}{n},
\]
\[
D(X_1-Y)=DX_1+DY-2\operatorname{Cov}(X_1,Y)
=\sigma^2+\frac{\sigma^2}{n}-\frac{2\sigma^2}{n}
=\frac{(n-1)\sigma^2}{n}.
\]
故 C、D 也不正确。
\end{solution}


\begin{exercise}\label{exer:src-2006-2a-11}
设随机变量 \(X\) 的密度函数为
\[
f(x)=
\begin{cases}
a\sin x, & 0\le x\le \pi,\\
0, & \text{其他}.
\end{cases}
\]
求：\par
(1) 常数 \(a\) 的值；\par
(2) \(E(X)\) 与 \(D(X)\)；\par
(3) 设事件 \(A=\{X\le 2\pi/3\}\)，\(B=\{X\ge \pi/4\}\)，求 \(P(A|B)\)。
\end{exercise}
\begin{solution}
(1) 由归一化条件，
\[
1=\int_0^\pi a\sin x\,\mathrm{d}x
=a[-\cos x]_0^\pi
=2a,
\]
所以
\(
a=\frac12.
\)

(2) 由 \(f(x)=\dfrac12\sin x\ (0\le x\le\pi)\)，
\[
E(X)=\int_0^\pi x\cdot\frac12\sin x\,\mathrm{d}x
=\frac12[-x\cos x+\sin x]_0^\pi
=\frac{\pi}{2}.
\]
再算二阶矩：
\[
\begin{aligned}
E(X^2)
&=\int_0^\pi x^2\cdot\frac12\sin x\,\mathrm{d}x\\
&=\frac12[-x^2\cos x+2x\sin x+2\cos x]_0^\pi\\
&=\frac{\pi^2-4}{2}.
\end{aligned}
\]
因此
\(
D(X)=E(X^2)-[E(X)]^2
=\frac{\pi^2-4}{2}-\frac{\pi^2}{4}
=\frac{\pi^2}{4}-2.
\)

(3) 由于 \(A\cap B=\{\pi/4\le X\le 2\pi/3\}\)，
\(
P(A|B)=\frac{P(A\cap B)}{P(B)}.
\)
其中
\[
P(A\cap B)=\int_{\pi/4}^{2\pi/3}\frac12\sin x\,\mathrm{d}x
=\frac12\left(\cos\frac{\pi}{4}-\cos\frac{2\pi}{3}\right)
=\frac{\sqrt2+1}{4},
\]
\[
P(B)=\int_{\pi/4}^{\pi}\frac12\sin x\,\mathrm{d}x
=\frac12\left(\cos\frac{\pi}{4}-\cos\pi\right)
=\frac{2+\sqrt2}{4}.
\]
故
\(
P(A|B)=\frac{\sqrt2+1}{\sqrt2+2}.
\)
\end{solution}

\begin{exercise}\label{exer:src-2008-01-3a-i-04}\label{exer:src-paper2-tf-04}\label{prob:p2-tf-04}
设 \(X\) 为离散型随机变量，且存在正数 \(k\)，使得
\(
P(|X|>k)=0.
\)
则 \(X\) 的数学期望 \(E(X)\) 未必存在。判断正误。
\end{exercise}
\begin{solution}
错误。由 \(P(|X|>k)=0\) 可知 \(X\) 几乎必然只在有限区间 \([-k,k]\) 内取值。于是
\(
|E X|\le E|X|\le k<+\infty.
\)
因此 \(EX\) 一定存在。
\end{solution}

\begin{exercise}\label{exer:src-2008-01-3c-i-03}\label{exer:src-paper4-i-03}
若 \(X\) 与 \(Y\) 都是标准正态随机变量，则 \(X+Y\sim N(0,2)\)。判断正误。
\end{exercise}
\begin{solution}
错误。仅知道 \(X,Y\) 都服从标准正态分布，并不能推出 \(X+Y\) 的分布。题目没有给出独立性。
\end{solution}


\chapter{大数定律和中心极限定理}

\section{依概率收敛}

\begin{definition}[依概率收敛] \label{def:convergence-prob}
设 \(X\) 为随机变量，\(\{X_n\}\ (n=1,2,\dots)\) 为随机变量序列。若对任意 \(\varepsilon > 0\)，有
\[
\lim_{n\to\infty} P\{|X_n - X| \geqslant \varepsilon\} = 0
\quad \text{或等价地} \quad
\lim_{n\to\infty} P\{|X_n - X| < \varepsilon\} = 1,
\]
则称序列 \(\{X_n\}\) \textbf{依概率收敛}于 \(X\)，记为
\[
X_n \stackrel{P}{\longrightarrow} X\ (n\to\infty).
\]
特别地，若 \(X = a\) 为常数，则称 \(\{X_n\}\) 依概率收敛于常数 \(a\)。依概率收敛是概率意义下的收敛，并不要求每一次试验中 \(X_n\) 都逐项贴近 \(X\)。它允许偶尔仍出现较大偏差，只要求这类偏差发生的概率趋于 0。依概率收敛是重要的弱收敛方式之一。通常有“几乎必然收敛 \(\Rightarrow\) 依概率收敛 \(\Rightarrow\) 依分布收敛”，但反向一般不成立。本章中的大数定律，描述的正是随机变量序列的算术平均值依概率收敛到其期望这一事实。
\end{definition}

\begin{property}[概率收敛运算性质]\label{prop:conv-prob-operation}
依概率收敛具有以下运算性质。设 \(X_n \stackrel{P}{\longrightarrow} a\)，\(Y_n \stackrel{P}{\longrightarrow} b\)，\(g(x,y)\) 为连续函数，则
\(
g(X_n, Y_n) \stackrel{P}{\longrightarrow} g(a, b).
\)
更一般地，对 \(m\) 元连续函数 \(g(x_1, x_2, \dots, x_m)\)，若 \(X_{n}^{(i)} \stackrel{P}{\longrightarrow} a_i\ (i=1,2,\dots,m)\)，则
\[
g\!\left(X_{n}^{(1)}, X_{n}^{(2)}, \dots, X_{n}^{(m)}\right) \stackrel{P}{\longrightarrow} g(a_1, a_2, \dots, a_m).
\]
\end{property}


\section{马尔可夫不等式和切比雪夫不等式}

\begin{theorem}[马尔可夫不等式] \label{thm:markov}
设随机变量 \(X\ge 0\)，且数学期望 \(EX\) 存在，则对任意 \(a>0\)，有
\(
P\{X\ge a\}\le \frac{EX}{a}.
\)
\end{theorem}


\begin{theorem}[切比雪夫不等式] \label{thm:chebyshev}
切比雪夫不等式其实可以看成把马尔可夫不等式用在非负随机变量 \(Y=(X-EX)^2\) 上。因为事件 \(\{|X-EX|\ge \varepsilon\}\) 等价于\(\{(X-EX)^2\ge \varepsilon^2\},\)
再对 \(Y\) 使用马尔可夫不等式，就得到
\[
P\{|X-EX|\geqslant \varepsilon\}\leqslant \frac{E[(X-EX)^2]}{\varepsilon^2}=\frac{DX}{\varepsilon^2},~~P\{|X - EX| < \varepsilon\} \geqslant 1 - \frac{DX}{\varepsilon^2}
\]
\end{theorem}


\begin{exercise}\label{exer:src-2015-spring-a-ii-06}\label{exer:src-2004-2a-06-chebyshev}\label{exer:src-2004-3a-06-chebyshev}
设 \(X_1,X_2,\dots,X_9\) 相互独立，且 \(EX_i=1,\ DX_i=1\ (i=1,2,\dots,9)\)。则对任意给定的 \(\varepsilon>0\)，下列结论正确的是（\quad）。\par
A. \(\displaystyle P\left\{\left|\sum_{i=1}^9X_i-1\right|<\varepsilon\right\}\ge1-\varepsilon^{-2}\)~~~~~~~
B. \(\displaystyle P\left\{\left|\frac19\sum_{i=1}^9X_i-1\right|<\varepsilon\right\}\ge1-\varepsilon^{-2}\)\par
C. \(\displaystyle P\left\{\left|\sum_{i=1}^9X_i-9\right|<\varepsilon\right\}>1-\varepsilon^{-2}\)~~~~~~~~
D. \(\displaystyle P\left\{\left|\sum_{i=1}^9X_i-9\right|<\varepsilon\right\}\ge1-9\varepsilon^{-2}\)
\end{exercise}
\begin{solution}
应选 D。令
\(
S=\sum_{i=1}^9X_i.
\)
由期望线性性和独立时方差可加，
\(
ES=\sum_{i=1}^9EX_i=9,\qquad
DS=\sum_{i=1}^9DX_i=9.
\)
由切比雪夫不等式的等价形式，
\(
P\{|S-ES|<\varepsilon\}\ge1-\frac{DS}{\varepsilon^2},
\)
得
\[
P\left\{\left|\sum_{i=1}^9X_i-9\right|<\varepsilon\right\}
\ge1-\frac9{\varepsilon^2}.
\]
\end{solution}







\begin{exercise}\label{exer:5-chebyshev-sum}
设 \(EX = -2\)，\(EY = 2\)，\(DX = 1\)，\(DY = 4\)，\(\rho(X,Y) = 0.5\)。求 \(D(X+Y) =\,\underline{\hspace{2cm}}\)，并根据切比雪夫不等式估计 \(P\{|X+Y| \geq 6\}\,\underline{\hspace{2cm}}\)。
\end{exercise}
\begin{solution}
先由相关系数求协方差。因为
\[
\rho(X,Y)=\frac{\mathrm{Cov}(X,Y)}{\sqrt{DX}\sqrt{DY}},
\]
所以
\[
\mathrm{Cov}(X,Y)=\rho(X,Y)\sqrt{DX}\sqrt{DY}
=0.5\times \sqrt{1}\times \sqrt{4}=1.
\]
于是
\(
D(X+Y) = DX + DY + 2\,\mathrm{Cov}(X,Y) = 1 + 4 + 2 = 7.
\)

再看第二问。由期望的线性性质，
\(
E(X+Y)=EX+EY=-2+2=0.
\)
因此事件 \(\{|X+Y|\ge 6\}\) 正好就是 \(\{|(X+Y)-E(X+Y)|\ge 6\}\)，这已经是切比雪夫不等式的标准形式，其中 \(\varepsilon=6\)。所以
\[
P\{|X+Y| \geq 6\} \leq \frac{D(X+Y)}{36} = \frac{7}{36}.
\]
\end{solution}



\begin{exercise}\label{exer:src-2013-a-01-04}\label{exer:src-2004-2a-16}
对随机变量 \(X\) 和 \(Y\)，已知
\[
E(X)=-2,\quad E(Y)=2,\quad D(X)=1,\quad D(Y)=4,\quad \rho_{X,Y}=-0.5.
\]
由切比雪夫不等式所能确定的最小正数 \(c\) 为何值，其中 \(c\) 满足
\(
P\{X+Y\ge 6\}\le c?
\)
\end{exercise}
\begin{solution}
先计算 \(X+Y\) 的期望与方差。由期望线性性，
\(
E(X+Y)=E(X)+E(Y)=-2+2=0.
\)
由相关系数公式，
\[
\mathrm{Cov}(X,Y)=\rho_{X,Y}\sqrt{D(X)}\sqrt{D(Y)}
=-0.5\times 1\times 2=-1.
\]
因此
\(
D(X+Y)=D(X)+D(Y)+2\mathrm{Cov}(X,Y)
=1+4+2(-1)=3.
\)

因为
\(
\{X+Y\ge6\}\subset \{|X+Y|\ge6\},
\)
并且 \(E(X+Y)=0\)，由切比雪夫不等式可得
\(
P\{X+Y\ge6\}
\le P\{|X+Y|\ge6\}
\le \frac{D(X+Y)}{6^2}
=\frac{1}{12}.
\)
所以由切比雪夫不等式能给出的最小上界为
\(
c=\frac1{12}.
\)
\end{solution}

\begin{exercise}\label{exer:5-chebyshev-pdf}
设随机变量 \(X\) 的概率密度为
\[
f(x) = \begin{cases} \dfrac{1}{2}\,x^2\,e^{-x}, & x > 0, \\ 0, & x \leq 0, \end{cases}
\]
利用切比雪夫不等式估计 \(P(1 < X < 5) \geq\,\underline{\hspace{2cm}}\)。
\end{exercise}
\begin{solution}
利用公式 \(\displaystyle\int_0^{+\infty} e^{-x}\,x^k\,\mathrm{d}x = k!\)（即 Gamma 函数）计算期望和方差：
\[
EX = \int_0^{+\infty} \frac{1}{2}\,x^3\,e^{-x}\,\mathrm{d}x = \frac{1}{2}\,\Gamma(4) = 3,\quad
EX^2 = \int_0^{+\infty} \frac{1}{2}\,x^4\,e^{-x}\,\mathrm{d}x = \frac{1}{2}\,\Gamma(5) = 12,
\]
\(
DX = EX^2 - (EX)^2 = 12 - 9 = 3.
\)
注意到 \(P(1 < X < 5) = P(|X - 3| < 2)\)，由切比雪夫不等式：
\(
P(|X - 3| < 2) \geq 1 - \frac{DX}{2^2} = 1 - \frac{3}{4} = \frac{1}{4}.
\)
\end{solution}



\section{大数定律}

\begin{theorem}[切比雪夫大数定律] \label{thm:chebyshev-lln}
设 \(X_1, X_2, \dots, X_n, \dots\) 是相互独立的随机变量序列，就有\(\displaystyle D\overline{X}_n=D\frac1n \sum_{i=1}^n X_i=\frac1{n^2}\sum_{i=1}^n DX_i\)，若方差 \(DX_i\ (i \geq 1)\) 存在且一致有界，就有\(\frac1{n^2}\sum_{i=1}^n DX_i=O(\frac{n}{n^2})=O(\frac1n)\)，即样本平均的方差$D\overline{X}_n$以$O(\frac{1}{n})$速度趋于0。此时再把切比雪夫不等式用于样本平均值和期望的差
\[
\overline{X}_n-\frac1n\sum_{i=1}^n EX_i,
\]
可知对任意 \(\varepsilon>0\)，偏离概率
\[
P\!\left\{\left|\overline{X}_n-\frac1n\sum_{i=1}^n EX_i\right|\ge\varepsilon\right\}
\le \frac{D\overline{X}_n}{\varepsilon^2}\to0.
\]
则\(n\) 个随机变量的算术平均值依概率收敛于其期望的算术平均值。即对任意 \(\varepsilon > 0\)，有
\[
\lim_{n\to\infty} P\!\left\{\left|\frac{1}{n}\sum_{i=1}^{n} X_i - \frac{1}{n}\sum_{i=1}^{n} EX_i\right| < \varepsilon\right\} = 1.
\]它允许 \(X_1,X_2,\dots\) 的期望彼此不同。只有在每个 \(EX_i\) 都等于同一个常数 \(\mu\) 时，结论才简化为
\[
\frac{1}{n}\sum_{i=1}^{n} X_i \stackrel{P}{\longrightarrow} \mu.
\]
\end{theorem}


\begin{exercise}\label{exer:5-lln-exponential}
设 \(X_1, X_2, \dots, X_n, \dots\) 是相互独立的随机变量序列，\(X_n\) 服从参数为 \(n\ (n \geq 1)\) 的指数分布，则下列随机变量序列中\textbf{不}满足切比雪夫大数定律条件的是（\quad）。\par
A. \(X_1, \dfrac{1}{2}X_2, \dots, \dfrac{1}{n}X_n, \dots\)\quad
B. \(X_1, X_2, \dots, X_n, \dots\)\par
C. \(X_1, 2X_2, \dots, nX_n, \dots\)\quad
D. \(X_1, 2^2X_2, \dots, n^2X_n, \dots\)
\end{exercise}

\begin{solution}
切比雪夫大数定律要求：\(\{X_n\}\) 相互独立，方差存在且一致有界（\(DX_n \leq C\)）。
由题设，\(DX_n = \dfrac{1}{n^2}\)（指数分布 \(E(\lambda)\) 的方差为 \(\frac{1}{\lambda^2}\)，此处 \(\lambda = n\)）。
逐一检验：
\begin{itemize}
    \item A: \(D\!\left(\dfrac{1}{n}X_n\right) = \dfrac{1}{n^2}\,DX_n = \dfrac{1}{n^4} \leq 1\) \checkmark
    \item B: \(DX_n = \dfrac{1}{n^2} \leq 1\) \checkmark
    \item C: \(D(nX_n) = n^2\,DX_n = 1 \leq 1\) \checkmark
    \item D: \(D(n^2 X_n) = n^4\,DX_n = n^2\)，当 \(n \to \infty\) 时无上界 \(\times\)
\end{itemize}
\end{solution}

\begin{exercise}\label{exer:src-lln-log-n}
设 \(\{X_n\}\) 是独立的随机变量序列，且假设
\[
P(X_n=\sqrt{\ln n})=P(X_n=-\sqrt{\ln n})=0.5,\qquad n=1,2,\dots.
\]
问 \(\{X_n\}\) 是否服从大数定律。
\end{exercise}
\begin{solution}
先看均值和方差。由对称性，
\(
EX_n=0.
\)
又
\(
X_n^2=\ln n
\)
恒成立，所以
\(
DX_n=EX_n^2-(EX_n)^2=\ln n.
\)
因此 \(\{X_n\}\) 的方差并不一致有界，不能直接套切比雪夫大数定律的充分条件。

但这不代表它不服从大数定律。记
\(
\overline X_n=\frac1n\sum_{k=1}^nX_k.
\)
由于 \(X_1,\dots,X_n\) 相互独立，
\[
D\overline X_n
=\frac1{n^2}\sum_{k=1}^nDX_k
=\frac1{n^2}\sum_{k=1}^n\ln k
\le \frac{n\ln n}{n^2}
=\frac{\ln n}{n}\to0.
\]
于是对任意 \(\varepsilon>0\)，由切比雪夫不等式，
\[
P\{|\overline X_n-0|\ge\varepsilon\}
\le \frac{D\overline X_n}{\varepsilon^2}
\le \frac{\ln n}{n\varepsilon^2}\to0.
\implies
\frac1n\sum_{k=1}^nX_k\stackrel{P}{\longrightarrow}0.
\]
所以 \(\{X_n\}\) 服从大数定律。
\end{solution}

\begin{theorem}[伯努利大数定律] \label{thm:bernoulli-lln}
伯努利大数定律是切比雪夫大数定律的特例（取 \(X_i \sim B(1,p)\)），设 \(\mu_n\) 是 \(n\) 重伯努利试验中事件 \(A\) 发生的次数，\(p\) 是事件 \(A\) 在每次试验中发生的概率（\(0 < p < 1\)），由于 \(\{X_i\}\) 相互独立，方差一致有界（\(\leq 1/4\)），由切比雪夫大数定律直接可得
\(
\overline{X}_n = \frac{\mu_n}{n} \xrightarrow{P} p.
\)，即对任意 \(\varepsilon > 0\)，有
\(
\lim_{n\to\infty} P\!\left\{\left|\frac{\mu_n}{n} - p\right| < \varepsilon\right\} = 1,
\)
即频率 \(\frac{\mu_n}{n}\) 依概率收敛于概率 \(p\)。在数理统计中，经验分布函数
\[
F_n(x) = \frac{\text{样本值中} \leq x \text{的个数}}{n}
\]
正是事件 \(\{X \leq x\}\) 在 \(n\) 次试验中出现的频率。由伯努利大数定律，当 \(n\) 充分大时，\(F_n(x)\) 可作为总体分布函数 \(F(x)\) 的近似，且 \(n\) 越大近似越好。
\end{theorem}

\begin{exercise}\label{exer:5-empirical}
设 \((2, 1, 5, 2, 1, 3, 1)\) 是来自总体 \(X\) 的简单随机样本值，求总体 \(X\) 的经验分布函数 \(F_7(x)\)。
\end{exercise}

\begin{solution}
将样本值按从小到大排列：\(1, 1, 1, 2, 2, 3, 5\)。
\[
F_7(x) = \begin{cases}
0, & x < 1, \\
\dfrac{3}{7}, & 1 \leq x < 2, \\
\dfrac{5}{7}, & 2 \leq x < 3, \\
\dfrac{6}{7}, & 3 \leq x < 5, \\
1, & x \geq 5.
\end{cases}
\]
\end{solution}


\begin{theorem}[辛钦大数定律] \label{thm:khinchin-lln}
\textbf{切比雪夫型大数定律}允许 \(X_1,X_2,\dots\) 的分布彼此不同，所以极限写成\(\frac1n\sum_{i=1}^n EX_i\)，代价是要额外控制方差一致有界。而\textbf{辛钦大数定律}把条件收紧成“独立同分布”，于是极限自然就是这个分布的期望 \(\mu\)；回报是它只要求期望存在，不再要求方差存在。设 \(X_1, X_2, \dots, X_n, \dots\) 是独立同分布的随机变量序列，若数学期望 \(EX_i = \mu\ (i = 1, 2, \dots)\) 存在，则对任意 \(\varepsilon > 0\)，有
\[
\lim_{n\to\infty} P\!\left\{\left|\frac{1}{n}\sum_{i=1}^{n} X_i - \mu\right| < \varepsilon\right\} = 1,
\]
即 \(\dfrac{1}{n}\displaystyle\sum_{i=1}^{n} X_i \stackrel{P}{\longrightarrow} \mu\)。若题目问的不是 \(\frac1n\sum X_i\)，而是
\(
\frac1n\sum h(X_i)
\)
这类“样本函数平均值”的极限，只要 \(X_i\) 独立同分布且 \(E|h(X_1)|<\infty\)，就可以令 \(Y_i=h(X_i)\)，再把辛钦大数定律直接用在 \(Y_i\) 上。进一步地，若 \(X_1, X_2, \dots\) 独立同分布且 \(E|X_1|^k<\infty\)，则把 \(Y_i=X_i^k\) 视为新的独立同分布随机变量，再应用辛钦大数定律可得
\[
\frac{1}{n}\sum_{i=1}^{n} X_i^k \stackrel{P}{\longrightarrow} E(X_1^k).
\]
\end{theorem}

\begin{exercise}\label{exer:5-khinchin-x2}
设随机变量序列 \(X_1, X_2, \dots, X_n, \dots\) 独立同分布，\(X_1\) 的概率密度为
\(
f(x) = \begin{cases} 1 - |x|, & |x| < 1, \\ 0, & \text{其他}, \end{cases}
\)
则当 \(n \to \infty\) 时，\(\dfrac{1}{n}\displaystyle\sum_{i=1}^{n} X_i^2\) 依概率收敛于（\quad）。~~
A. \(\dfrac{1}{8}\)\quad B. \(\dfrac{1}{6}\)\quad C. \(\dfrac{1}{3}\)\quad D. \(\dfrac{1}{2}\)
\end{exercise}
\begin{solution}
令 \(Y_i=X_i^2\)。
由于 \(X_1,X_2,\dots\) 独立同分布，\(Y_1,Y_2,\dots\) 也独立同分布；又因为 \(X_i\) 的取值落在 \([-1,1]\) 上，故 \(E(Y_i)=E(X_i^2)\) 一定存在。于是由辛钦大数定律，
\[
\frac{1}{n}\sum_{i=1}^{n} X_i^2 \stackrel{P}{\longrightarrow} E(X_1^2) = \int_{-1}^{1} x^2(1-|x|)\,\mathrm{d}x = 2\int_0^1 x^2(1-x)\,\mathrm{d}x=\frac{1}{6}.
\]
\end{solution}

\section{中心极限定理}

大数定律告诉我们
\(
\overline{X}_n-\mu \stackrel{P}{\longrightarrow} 0,
\)
但它没有告诉我们：误差通常有多大？怎么近似算概率？中心极限定理补的正是这一块。
设 \(X_1,\dots,X_n\) 独立同分布，它们的和的方差：\(D(S_n)=D\left(\sum_{i=1}^n X_i\right)=\sum_{i=1}^nD(X_i)=n\sigma^2\)。因此 \(D(S_n-n\mu)=D(S_n)=n\sigma^2\)。标准差就是方差开根号：\(\sqrt{D(S_n-n\mu)}=\sqrt{n\sigma^2}=\sqrt n\,\sigma\)。而“典型波动量级”通常就是看标准差的数量级，所以 \(S_n-n\mu\) 的典型波动大约是 \(\sqrt n\)。


\begin{theorem}[列维--林德伯格中心极限定理：要求独立、同分布、期望和方差存在] \label{thm:levy-lindberg}
设 \(X_1, X_2, \dots, X_n, \dots\) 是独立同分布的随机变量序列，且数学期望 \(EX_i = \mu\) 和方差 \(DX_i = \sigma^2 > 0\) 存在 \((i = 1, 2, \dots)\)，则对任意实数 \(x\)，有
\[
\lim_{n\to\infty} P\!\left\{\frac{\sum_{i=1}^{n} X_i - n\mu}{\sqrt{n}\,\sigma} \leq x\right\} = \Phi(x)=\frac{1}{\sqrt{2\pi}}\int_{-\infty}^x e^{-\frac{t^2}{2}}\,\mathrm{d}t.
\]也就是说，当 \(n\) 足够大时，\(\displaystyle\sum_{i=1}^{n} X_i\) 近似服从正态分布 \(N(n\mu,\, n\sigma^2)\)，因此
\[
P\!\left\{a < \sum_{i=1}^{n} X_i < b\right\} \approx \Phi\!\left(\frac{b - n\mu}{\sqrt{n}\,\sigma}\right) - \Phi\!\left(\frac{a - n\mu}{\sqrt{n}\,\sigma}\right).
\]若题目问的是样本均值 \(\overline{X}_n\)，也完全可以使用中心极限定理，因为
\[
\frac{\overline{X}_n-\mu}{\sigma/\sqrt{n}}
=\frac{\sum_{i=1}^n X_i-n\mu}{\sqrt{n}\,\sigma}.
\]
\end{theorem}


\begin{theorem}[棣莫弗--拉普拉斯中心极限定理] \label{thm:demoivre-laplace}
棣莫弗--拉普拉斯定理是列维--林德伯格定理在二项分布上的特例。若将 \(Y_n\) 看作 \(n\) 个独立同分布的 \(B(1,p)\) 随机变量之和，即 \(Y_n = \sum_{i=1}^{n} X_i\)，其中 \(X_i \sim \begin{pmatrix} 1 & 0 \\ p & 1-p \end{pmatrix}\)，则 \(\mu = p\)，\(\sigma^2 = p(1-p)\)，代入列维--林德伯格定理即得——\\
设随机变量 \(Y_n \sim B(n, p)\ (0 < p < 1,\ n \geq 1)\)，则对任意实数 \(x\)，有
\[
\lim_{n\to\infty} P\!\left\{\frac{Y_n - np}{\sqrt{np(1-p)}} \leq x\right\} = \Phi(x).
\]
\end{theorem}

\textbf{二项分布概率计算的三种方法}：设 \(X \sim B(n, p)\)。
\begin{enumerate}
    \item \(n\) 较小（\(n \leq 10\)）：直接用公式 \(P\{X=k\} = \mathrm{C}_n^k\,p^k(1-p)^{n-k}\)。
    \item \(n\) 较大且 \(p\) 较小（\(n > 10,\ p < 0.1,\ np\) 适中）：泊松近似 \(\displaystyle P\{X=k\} \approx \dfrac{\lambda^k}{k!}\,e^{-\lambda}\)，其中 \(\lambda = np\)。
    \item \(n\) 较大且 \(p\) 不太大（\(np \geq 10\)）：正态近似（中心极限定理）
    \[
    P\{a < X < b\} \approx \Phi\!\left(\frac{b-np}{\sqrt{np(1-p)}}\right) - \Phi\!\left(\frac{a-np}{\sqrt{np(1-p)}}\right).
    \]
\end{enumerate}

\begin{exercise}\label{exer:5-clt-condition}
设随机变量 \(X_1, X_2, \dots, X_n\) 相互独立，\(S_n = X_1 + X_2 + \cdots + X_n\)。根据列维--林德伯格定理，当 \(n\) 充分大时，\(S_n\) 近似服从正态分布，只要（\quad）。\par
A. 有相同期望和方差\quad B. 服从同一离散分布\quad
C. 服从同一指数分布\quad D. 服从同一连续分布
\end{exercise}

\begin{solution}
\begin{itemize}
    \item A：有相同的期望和方差 \(\Rightarrow\) 未必同分布，不满足条件。
    \item B：服从同一离散型分布 \(\Rightarrow\) 方差未必存在（反例：若 \(P(X=k)=c/k^3\ (k=1,2,\dots)\)，则 \(EX^2=\sum k^2\cdot c/k^3=\sum c/k\) 发散，方差不存在）。
    \item C：指数分布 \(E(\lambda)\) 的期望和方差分别为 \(1/\lambda\) 和 \(1/\lambda^2\)，均存在，满足全部三个条件 \checkmark
    \item D：服从同一连续型分布 \(\Rightarrow\) 方差未必存在（如柯西分布）。
\end{itemize}
\end{solution}


\begin{exercise}\label{exer:5-clt-box}
生产线生产的产品成箱包装，每箱质量是随机的。假设平均每箱重 50 千克，标准差为 5 千克。若用载重为 5 吨的汽车承运，试用中心极限定理说明每辆车最多可以装多少箱，才能保证不超载的概率大于 0.977。\ (\(\Phi(2) = 0.977\))
\end{exercise}

\begin{solution}
设 \(X_i\ (i=1,2,\dots,n)\) 为第 \(i\) 箱的质量（千克），\(X_i\) 独立同分布，\(EX_i = 50\)，\(DX_i = 5\)。\(n\) 箱总质量 \(T_n = \displaystyle\sum_{i=1}^{n} X_i\)，\(ET_n = 50n\)，\(DT_n = 25n\)，近似服从 \(N(50n,\, 25n)\)：
\[
P(T_n \leq 5000) = P\!\left(\frac{T_n - 50n}{5\sqrt{n}} \leq \frac{5000 - 50n}{5\sqrt{n}}\right) \approx \Phi\!\left(\frac{1000 - 10n}{\sqrt{n}}\right) > 0.977 = \Phi(2),
\]
即 \(\dfrac{1000 - 10n}{\sqrt{n}} > 2\)。故 \(n < 99.0\)。因此最多装 \(\mathbf{98}\) 箱。
\end{solution}


\begin{exercise}\label{exer:5-clt-poisson}
设 \(X_1, X_2, \dots, X_n\) 相互独立，且均服从泊松分布 \(P(2)\)，则由列维--林德伯格中心极限定理
\[
\lim_{n\to\infty} P\!\left(\frac{\sum_{i=1}^n X_i - 2n}{\sqrt{n}} < 2\right) = \text{（\quad）}.
\]
A. \(\Phi(\sqrt{2})\)\quad B. \(\Phi(2)\)\quad C. \(\Phi(1)\)\quad D. \(1\)
\end{exercise}
\begin{solution}
\[
P\!\left(\frac{\sum X_i - 2n}{\sqrt{n}} < 2\right)
= P\!\left(\frac{\sum X_i - 2n}{\sqrt{2n}} < \frac{2}{\sqrt{2}}\right)
= P(Z < \sqrt{2}) = \Phi(\sqrt{2}).
\]
\end{solution}



\begin{exercise}\label{exer:5-clt-moment}
假设 \(X_1, X_2, \dots, X_n\) 是来自总体 \(X\) 的简单随机样本，已知 \(E(X^k) = \alpha_k\ (k = 1,2,3,4)\)。证明：当 \(n\) 充分大时，随机变量
\(
Z_n = \frac{1}{n}\sum_{i=1}^{n} X_i^2
\)
近似服从正态分布，并指出其分布参数。
\end{exercise}
\begin{solution}
由题设，\(X_1, X_2, \dots, X_n\) 独立同分布，故 \(X_1^2, X_2^2, \dots, X_n^2\) 也独立同分布，且
\[
E(X_i^2) = \alpha_2,\quad D(X_i^2) = E(X_i^4) - [E(X_i^2)]^2 = \alpha_4 - \alpha_2^2.
\]
由列维--林德伯格中心极限定理，
\[
\frac{\sum_{i=1}^{n} X_i^2 - n\alpha_2}{\sqrt{n(\alpha_4 - \alpha_2^2)}} \to N(0,1),~~~
\frac{Z_n - \alpha_2}{\sqrt{(\alpha_4 - \alpha_2^2)/n}} \to N(0,1),
\]
因此当 \(n\) 充分大时，\(Z_n\) 近似服从正态分布，参数为
\[
\mu = \alpha_2,\quad \sigma^2 = \frac{\alpha_4 - \alpha_2^2}{n}.
\]
\end{solution}



\begin{exercise}\label{exer:5-clt-tap}
某宿舍楼按 400 名学生的规模建造，每人每天有约 10\% 的时间占用一个水龙头（各人独立）。为能以 95\% 的概率保证用水需要，至少应安装多少个水龙头？\ (\(\Phi(1.645) = 0.95\))
\end{exercise}
\begin{solution}
先把题目建模成二项分布。对每个学生来说，在某一时刻有 \(10\%\) 的概率使用水龙头，又因为 400 名学生彼此独立，所以同一时刻的用水人数 \(X\) 服从 \(X\sim B(400,0.1)\)。因此 \(EX=np=40,\ DX=np(1-p)=400\times 0.1\times 0.9=36\)。题目要求 \(P(X\le N)\ge 0.95\)。
当 \(n=400\) 较大时，可用中心极限定理把 \(X\) 近似为正态分布 \(N(40,36)\)。于是
\[
P(X \leq N) \approx \Phi\!\left(\frac{N - 40}{6}\right) \geq 0.95 = \Phi(1.645),~~~~
\frac{N-40}{6}\ge 1.645 \implies N\ge 40+6\times 1.645=49.87.
\]
由于水龙头数量必须取整数，故至少应安装 \(N=50\) 个水龙头。
\end{solution}


\begin{exercise}\label{exer:src-probstat-test1-06}
某生产线产品合格率为 \(0.8\)。要使一批产品的合格率落在 \(76\%\) 与 \(84\%\) 之间的概率不小于 \(90\%\)，问这批产品至少要生产多少件？分别用以下方法估计：\par
(1) 切比雪夫不等式；\par
(2) 中心极限定理。\par
已知 \(\Phi(1.645)=0.95\)。
\end{exercise}
\begin{solution}
(1) 设生产 \(n\) 件，合格品数为 \(X\)，则
\[
X\sim B(n,0.8),\hat p=\frac Xn\implies
E\hat p=0.8,D\hat p=D\frac{X}{n}=\frac{n0.8(1-0.8)}{n^2}=\frac{0.16}{n}.
\]
题目要求
\[
P(0.76\le \hat p\le0.84)
=P(|\hat p-0.8|\le0.04)\ge0.90.
\]由切比雪夫不等式，
\(
P(|\hat p-0.8|<0.04)
\ge1-\frac{D\hat p}{0.04^2}
=1-\frac{0.16/n}{0.0016}
=1-\frac{100}{n}.
\)
要使下界不小于 \(0.90\)，只需
\(
1-\frac{100}{n}\ge0.90,
\)
即
\(
n\ge1000.
\)\\
(2) 设生产 \(n\) 件，合格数 \(X \sim B(n, 0.8)\)。对于单次实验而言，$\mu=0.8$，方差为 \(p(1-p) = 0.16\)。
\begin{align*}
P\!\left(0.76 \leq \frac{X}{n} \leq 0.84\right) &=P\!\left(0.76 n\leq X \leq 0.84n\right)=P\!\left(-0.04n\leq X-0.8n \leq 0.04n\right)\\
&=P\!\left(|X-0.8n|\leq 0.04n\right)=P\!\left(\left|\frac{X-0.8n}{\sqrt{0.16n}}\right|\leq
\frac{0.04n}{0.4\sqrt n}\right)\\
&\approx P(|Z|\leq0.1\sqrt n)
=2\Phi(0.1\sqrt n)-1.
\end{align*}
要使该概率不小于 \(0.90\)，只需
\(
2\Phi(0.1\sqrt n)-1\ge0.90,
\)
即
\(
\Phi(0.1\sqrt n)\ge0.95.
\)
由 \(\Phi(1.64)=0.95\)，得
\(
0.1\sqrt n\ge1.64,
\)
从而
\(
n\ge(16.4)^2=268.96.
\)
故至少要生产 \(269\) 件。
\end{solution}



\begin{exercise}\label{exer:src-2005-2a-16}
某工厂生产的零件废品率为 \(5\%\)。某人要采购一批零件，他希望以 \(95\%\) 的概率保证其中有 \(2000\) 个及格品，问他至少应购买多少零件？已知 \(\Phi(1.65)=0.9505\)。
\end{exercise}
\begin{solution}
$n$次二项分布，二项分布的期望为$p=0.95$，方差为$0.95\times0.05$，设购买 \(n\) 个零件，其中及格品数为 \(\zeta\)。
\[
P(\zeta\ge 2000)=1-P(\zeta\le 2000)=1-P\left(\frac{\zeta-0.95n}{\sqrt{0.0475n}}\le\frac{2000-0.95n}{\sqrt{0.0475n}}\right)
\approx
1-\Phi\!\left(\frac{2000-0.95n}{\sqrt{0.0475n}}\right).
\]
令该近似值达到 \(0.95\)，即
\(
1-\Phi\!\left(\frac{2000-0.95n}{\sqrt{0.0475n}}\right)=0.95.
\)
由 \(\Phi(1.65)\approx0.95\)，可取
\(
\frac{2000-0.95n}{\sqrt{0.0475n}}=-1.65.
\)
解得 \(n\approx 2122.4\)。因为购买数量必须为整数，且要至少保证要求，故取
\(
n=2123.
\)
所以至少应购买 \(2123\) 个零件。
\end{solution}

\begin{exercise}\label{exer:src-2006-2a-12}
在一家保险公司里有 \(25000\) 辆汽车参加保险，每辆汽车每年付保险费 \(240\) 元。一年内一辆汽车出事故的概率为 \(0.001\)，若一辆车出事故，则保险公司赔付 \(20\) 万元。问：\par
(1) 保险公司亏本的概率为多大；\par
(2) 保险公司每年的利润不小于 \(100\) 万元的概率为多大。\par
（参考数据：\(\Phi(1)=0.8413,\ \Phi(2)=0.9772,\ \Phi(3)=0.9987\)。）
\end{exercise}
\begin{solution}
设一年内出事故的汽车数为 \(X\)，则25000次是否出事的伯努利实验，单次实验期望为$p=0.001$，方差为$0.001\times0.999$\\
(1) 保险公司的保费收入为
\(
25000\times240=6000000\quad(\text{元}),
\)
赔付金额为 \(200000X\)，故利润为
\(
L=6000000-200000X.
\)，亏本即 \(L<0\)，也就是
\[
6000000-200000X<0 \quad\Longleftrightarrow\quad X>30.
\]要有大于30辆车出事故，用正态近似 \(X\approx N(25,25)\)，得
\begin{align*}
P(L<0)&=P(X>30)=P\left(\frac{X-0.001\times25000}{\sqrt{25000}\sqrt{0.001\times0.999}}>\frac{30-0.001\times25000}{\sqrt{25000}\sqrt{0.001\times0.999}}\right)\\
&=1-P\left(\frac{X-0.001\times25000}{\sqrt{25000}\sqrt{0.001\times0.999}}<\frac{30-0.001\times25000}{\sqrt{25000}\sqrt{0.001\times0.999}}\right)\\
&\approx 1-\Phi\left(\frac{30-25}{5}\right)
=1-\Phi(1)
=0.1587.
\end{align*}
(2) 利润不小于 \(100\) 万元，即
\[
6000000-200000X\ge1000000
\quad\Longleftrightarrow\quad
X\le25.
\]
因此
\[
P(L\ge1000000)=P(X\le25)
\approx \Phi\left(\frac{25-25}{5}\right)
=\Phi(0)=0.5.
\]
所以保险公司每年利润不小于 \(100\) 万元的概率近似为 \(0.5\)。
\end{solution}



\begin{exercise}\label{exer:src-2011-ab-a-05}
商场销售某种商品，每周销售量（件数）服从参数 \(\lambda=9\) 的泊松分布，各周销售量相互独立，一年按 \(50\) 个销售周计。每销售一件该商品商场可获得 \(10\) 元利润。求（精确到元）：\par
(1) 一年中商场售出该商品件数在 \(400\) 件到 \(500\) 件之间的概率；\par
(2) 以 \(95\%\) 的把握估计商场一年中能获得的最低利润；\par
(3) 以 \(95\%\) 的把握估计商场一年中能获得的最高利润。
\end{exercise}
\begin{solution}
(1) 设第 \(i\) 周销量为 \(X_i\)，则\(S=\sum_{i=1}^{50}X_i.\)，由中心极限定理，
\[
\begin{aligned}
P(400\le S\le500)
&=P\left\{\frac{400-50\times9}{3\sqrt{50}}\le S\le
\frac{500-50\times9}{3\sqrt{50}}\right\}\\
&=2\Phi\left(\frac{\sqrt{50}}3\right)-1
\approx2\Phi(2.36)-1\approx0.9818.
\end{aligned}
\]
(2) 设最低利润为 \(m\)。要求
\[
0.95\leqslant P(10S\ge m)=P(S\ge \frac{m}{10})=P\left(\frac{S-50\times9}{3\sqrt{50}}\ge \frac{400\frac{m}{10}-50\times9}{3\sqrt{50}}\right)=2\Phi\left(\frac{450-m/10}{3\sqrt{50}}\right)-1
\]
故
\(
m=10\left(450-1.645\cdot3\sqrt{50}\right)\approx4151.
\)
所以最低利润约为 \(4151\) 元。\\
(3) 设最高利润为 \(M\)。要求
\(
P\{10S\le M\}\ge0.95.
\)
同理
\(
\Phi\left(\frac{M/10-450}{3\sqrt{50}}\right)=0.95,
\)
从而
\(
M=10\left(450+1.645\cdot3\sqrt{50}\right)\approx4849.
\)
所以最高利润约为 \(4849\) 元。
\end{solution}



\begin{exercise}\label{exer:src-2013-a-06}
某供电站供应某地区 \(1000\) 户居民用电，各户每日用电量相互独立。已知每户每日用电量（单位：度）在区间 \([0,20]\) 上服从均匀分布。若要以 \(0.99\) 的概率满足该地区居民每日用电需求，问供电站每天至少应向该地区供应多少度电？（已知 \(\Phi(2.33)\approx0.99\)）
\end{exercise}
\begin{solution}
设第 \(k\) 户居民每日用电量为 \(X_k\)，则
\(
X_k\sim U[0,20],\qquad EX_k=10,\qquad DX_k=\frac{20^2}{12}=\frac{100}{3}.
\)
记该地区每日总用电量为
\(
S=\sum_{k=1}^{1000}X_k.
\)
则
\[
ES=1000\cdot10=10000,\qquad
DS=1000\cdot\frac{100}{3}=\frac{100000}{3}.
\]
由中心极限定理，
\(
\frac{S-10000}{\sqrt{100000/3}}\approx N(0,1).
\)
设供电量为 \(L\)，题意要求
\[
P\{S\le L\}\approx
\Phi\!\left(\frac{L-10000}{\sqrt{100000/3}}\right)=0.99.
\]
查表得对应分位数约为 \(2.33\)，故
\(
\frac{L-10000}{\sqrt{100000/3}}\approx2.33.
\)
于是
\[
L\approx10000+2.33\sqrt{\frac{100000}{3}}\approx10425.
\]
因此供电站每天至少应供应约 \(10425\) 度电。
\end{solution}


\begin{exercise}\label{exer:src-2017-a-06}
有一批种子，其中良种占 \(\dfrac16\)。从中任取 \(180\) 粒，问能以 \(0.99\) 的概率保证其中良种的比例与 \(\dfrac16\) 相差多少？已知
\(
\Phi(2.33)=0.99,\qquad \Phi(2.48)=0.995.
\)
\end{exercise}
\begin{solution}
令
\[
X_i=
\begin{cases}
1, & \text{第 }i\text{ 粒为良种},\\
0, & \text{第 }i\text{ 粒不是良种},
\end{cases}
\qquad i=1,2,\dots,180.
\]
则 \(X_i\) 可近似看作独立同分布的伯努利随机变量，且
\(
p=P\{X_i=1\}=\frac16,\qquad q=1-p=\frac56.
\)
样本良种比例为
\(
\hat p=\frac1{180}\sum_{i=1}^{180}X_i.
\)
由中心极限定理，
\(
\frac{\hat p-p}{\sqrt{pq/180}}\approx N(0,1).
\)
因为
\(
\sqrt{\frac{pq}{180}}
=\sqrt{\frac{(1/6)(5/6)}{180}}
=\frac1{36},
\)
故
\(
P\left\{\left|\hat p-\frac16\right|<C\right\}
\approx 2\Phi(36C)-1.
\)
令该概率为 \(0.99\)，得
\(
2\Phi(36C)-1=0.99,
\)
即
\(
\Phi(36C)=0.995.
\)
由 \(\Phi(2.48)=0.995\)，可得
\(
36C=2.48,\qquad C=\frac{2.48}{36}=\frac{31}{450}\approx0.07.
\)
因此能以 \(0.99\) 的概率保证良种比例与 \(\dfrac16\) 的差不超过约 \(0.07\)。
\end{solution}


\begin{exercise}\label{exer:src-2011-ab-b-06-2credit}
机械加工设备加工某种工件的长度 \(X\sim N(100,2.25)\)。在正式出厂前需要试生产 \(100\) 个该种工件。试求在试生产的 \(100\) 个工件中，长度误差不小于 \(3\%\) 的工件个数不少于 \(5\) 件的概率。
\end{exercise}
\begin{solution}
单个工件长度误差不小于 \(3\%\) 的概率为
\[
p=P\{|X-100|\ge3\}
=2\left[1-\Phi\left(\frac3{\sqrt{2.25}}\right)\right]
=2(1-\Phi(2)).
\]
由 \(\Phi(2)=0.9772\)，得 \(p\approx0.0456\)。设 \(Y\) 为 \(100\) 个工件中长度误差不小于 \(3\%\) 的件数，则
\(
Y\sim B(100,0.0456).
\)
用泊松近似，\(\lambda=np=4.56\)。于是
\[
P(Y\ge5)
\approx1-e^{-4.56}\left(1+4.56+\frac{4.56^2}{2!}
+\frac{4.56^3}{3!}+\frac{4.56^4}{4!}\right)
\approx0.479.
\]
\end{solution}


\begin{exercise}\label{exer:src-2007-a-14}
叙述大数定律，并证明下列随机变量序列服从大数定律：
\[
\xi_n\sim
\begin{pmatrix}
-\sqrt n & 0 & \sqrt n\\
\dfrac1n & 1-\dfrac2n & \dfrac1n
\end{pmatrix},
\qquad n=2,3,4,\dots
\]
\end{exercise}
\begin{solution}
切比雪夫大数定律可表述为：若随机变量 \(X_1,X_2,\dots\) 相互独立，且存在常数 \(C\)，使 \(DX_i\le C\)，则对任意 \(\varepsilon>0\)，有
\[
\lim_{n\to\infty}
P\!\left\{\left|\frac1n\sum_{i=1}^n X_i-\frac1n\sum_{i=1}^n EX_i\right|<\varepsilon\right\}=1.
\]

对本题，
\[
E\xi_n=(-\sqrt n)\frac1n+0\left(1-\frac2n\right)+(\sqrt n)\frac1n=0,
\]
并且
\[
D\xi_n=E(\xi_n^2)=n\cdot\frac1n+n\cdot\frac1n=2.
\]
方差一致有界。若 \(\xi_2,\xi_3,\dots\) 相互独立，令
\(
\overline{\xi}_n=\frac1n\sum_{k=2}^{n+1}\xi_k.
\)
则
\[
E\overline{\xi}_n=0,\qquad
D\overline{\xi}_n=\frac1{n^2}\sum_{k=2}^{n+1}D\xi_k=\frac{2}{n}.
\]
由切比雪夫不等式，对任意 \(\varepsilon>0\)，
\[
P\{|\overline{\xi}_n| \ge \varepsilon\}
\le \frac{D\overline{\xi}_n}{\varepsilon^2}
=\frac{2}{n\varepsilon^2}\to0.
\]
故
\[
\lim_{n\to\infty}P\{|\overline{\xi}_n|<\varepsilon\}=1,
\]
即该随机变量序列服从大数定律。
\end{solution}


\begin{exercise}\label{exer:src-2008-probstat-2a-vi-01}\label{exer:src-paper3-iv-04}\label{prob:p3-iv-04}
某厂生产某产品 \(1000\) 件，其价格为 \(2000\) 元/件。该产品使用寿命 \(X\)（单位：天）的密度函数为
\[
f(x)=
\begin{cases}
\dfrac1{20000}e^{-\frac1{20000}(x-365)},&x\ge365,\\
0,&x<365.
\end{cases}
\]
现由某保险公司为其质量进行保险，厂方向保险公司交保费 \(P_0\) 元/件。若每件产品寿命小于 \(1095\) 天（3 年），则保险公司按原价赔偿 \(2000\) 元/件。试利用中心极限定理计算：\par
(1) 若保费 \(P_0=100\) 元/件，保险公司亏本的概率；\par
(2) 确定保费 \(P_0\)，使保险公司亏本的概率不超过 \(1\%\)。\par
已知 \(e^{-0.0365}\approx0.96,\ \Phi(1.61)=0.946,\ \Phi(2.33)=0.99\)。
\end{exercise}
\begin{solution}
先求单件产品需要赔偿的概率。由密度函数可得
\[
F(x)=
\begin{cases}
0,&x<365,\\
1-e^{-\frac1{20000}(x-365)},&x\ge365.
\end{cases}
\]
因此
\(
P(X<1095)=F(1095)=1-e^{-\frac{1095-365}{20000}}
=1-e^{-0.0365}\approx0.04.
\)

设 \(N\) 表示 \(1000\) 件产品中寿命小于 \(1095\) 天的件数，则
\(
N\sim B(1000,0.04).
\)
于是
\[
EN=40,\qquad DN=1000\times0.04\times0.96=38.4,\qquad \sqrt{DN}\approx6.2.
\]

(1) 当 \(P_0=100\) 时，保险公司收取保费 \(1000\times100\) 元，赔偿金额为 \(2000N\) 元。亏本近似对应
\(
2000N\ge100000,\qquad\text{即 }N\ge50.
\)
由中心极限定理，
\[
\begin{aligned}
P(\text{亏本})
&\approx P(N\ge50)\\
&=P\left(\frac{N-40}{6.2}\ge\frac{50-40}{6.2}\right)\\
&\approx1-\Phi(1.61)
=1-0.946
=0.054.
\end{aligned}
\]

(2) 一般地，利润为
\(
1000P_0-2000N.
\)
要求亏本概率不超过 \(1\%\)，即
\(
P(1000P_0-2000N<0)\le0.01,
\)
近似写为
\(
P\left(N\ge\frac{P_0}{2}\right)\le0.01.
\)
由中心极限定理，只需
\(
\frac{P_0/2-40}{6.2}\ge2.33.
\)
故
\(
P_0\ge2(40+6.2\times2.33)=108.892.
\)
因此保费可取
\(
P_0\approx108.9\text{ 元/件},
\)
若要求按整元收费，则至少取 \(109\) 元/件。
\end{solution}

\begin{exercise}\label{exer:src-2009-a-06}\label{exer:src-paper4-iv-04}
某学校北区食堂为提高服务质量，要先对就餐率 \(p\) 进行调查。决定在某天中午随机地对用过午餐的同学进行抽样调查。设调查了 \(n\) 个同学，其中在北区食堂用过餐的学生数为 \(m\)。若要求以大于 \(95\%\) 的概率保证调查所得的就餐频率与 \(p\) 之间的误差上下在 \(10\%\) 以内，问 \(n\) 应取多大？
\end{exercise}
\begin{solution}
把“是否在北区食堂用餐”看作一次伯努利试验，成功概率为 \(p\)。则
\[
m\sim B(n,p),\qquad E\left(\frac mn\right)=p,\qquad D\left(\frac mn\right)=\frac{p(1-p)}{n}.
\]
由中心极限定理，
\(
\frac{m/n-p}{\sqrt{p(1-p)/n}}\approx N(0,1).
\)
题目要求
\(
P\left(\left|\frac mn-p\right|<0.1\right)>0.95.
\)
近似为
\(
2\Phi\left(\frac{0.1}{\sqrt{p(1-p)/n}}\right)-1>0.95.
\)
故需要
\(
\frac{0.1}{\sqrt{p(1-p)/n}}\ge1.96.
\)
也就是
\(
n\ge (19.6)^2p(1-p).
\)
由于 \(p\) 未知，取最保守的
\(
p(1-p)\le\frac14,
\)
于是
\(
n\ge (19.6)^2\cdot\frac14=96.04.
\)
所以应取
\(
n=97.
\)
\end{solution}

\begin{exercise}\label{exer:src-2010-a-09}
证明：如果不独立的随机变量 \(\xi_1,\xi_2,\dots,\xi_n,\dots\) 满足条件
\[
\lim_{n\to\infty}\frac1{n^2}D\!\left(\sum_{i=1}^n \xi_i\right)=0,
\]
则对于任意正数 \(\varepsilon\)，恒有
\[
\lim_{n\to\infty}
P\left(
\left|
\frac1n\sum_{i=1}^n \xi_i
-
\frac1n\sum_{i=1}^n E(\xi_i)
\right|<\varepsilon
\right)=1.
\]
\end{exercise}
\begin{solution}
记
\(
S_n=\sum_{i=1}^n \xi_i.
\)
则
\(
E\left(\frac{S_n}{n}\right)=\frac1n\sum_{i=1}^n E(\xi_i),
\)
且
\(
D\left(\frac{S_n}{n}\right)
=\frac1{n^2}D(S_n)
=\frac1{n^2}D\!\left(\sum_{i=1}^n \xi_i\right).
\)
由切比雪夫不等式，对任意 \(\varepsilon>0\)，
\[
\begin{aligned}
P\left(
\left|
\frac{S_n}{n}
-
E\left(\frac{S_n}{n}\right)
\right|\ge\varepsilon
\right)
&\le
\frac{D(S_n/n)}{\varepsilon^2}\\
&=
\frac1{\varepsilon^2}
\frac1{n^2}D\!\left(\sum_{i=1}^n \xi_i\right).
\end{aligned}
\]
由题设，右端趋于 \(0\)。故
\[
\lim_{n\to\infty}
P\left(
\left|
\frac1n\sum_{i=1}^n \xi_i
-
\frac1n\sum_{i=1}^n E(\xi_i)
\right|\ge\varepsilon
\right)=0.
\]
取对立事件即得
\[
\lim_{n\to\infty}
P\left(
\left|
\frac1n\sum_{i=1}^n \xi_i
-
\frac1n\sum_{i=1}^n E(\xi_i)
\right|<\varepsilon
\right)=1.
\]
\end{solution}


\begin{exercise}\label{exer:5-dml-burn}
设某种电气元件不能承受超负荷试验的概率为 0.05。现对 100 个这样的元件进行超负荷试验，以 \(X\) 表示不能承受试验而烧毁的元件数。根据中心极限定理，求 \(P\{5 \leq X \leq 10\}\) 的近似值。\ (\(\Phi(2.29) = 0.989\))
\end{exercise}
\begin{solution}
\(X \sim B(100, 0.05)\)，\(EX = np = 5\)，\(DX = np(1-p) = 4.75\)。由棣莫弗--拉普拉斯中心极限定理：
\[
\begin{aligned}
P\{5 \leq X \leq 10\}
&= P\!\left\{\frac{5 - 5}{\sqrt{4.75}} \leq \frac{X - 5}{\sqrt{4.75}} \leq \frac{10 - 5}{\sqrt{4.75}}\right\} \\
&= P\!\left\{0 \leq \frac{X - 5}{\sqrt{4.75}} \leq 2.29\right\}
\approx \Phi(2.29) - \Phi(0) = 0.989 - 0.5 = 0.489.
\end{aligned}
\]
\end{solution}


\begin{exercise}\label{exer:5-clt-bernoulli-sum}
设 \(X_1, X_2, \dots, X_{100}\) 为来自总体 \(X\) 的简单随机样本，其中 \(P\{X=0\} = P\{X=1\} = \dfrac{1}{2}\)。利用中心极限定理求 \(P\!\left\{\displaystyle\sum_{i=1}^{100} X_i \leq 55\right\}\) 的近似值（\quad）。\par
A. \(1 - \Phi(1)\)\quad B. \(\Phi(1)\)\quad C. \(1 - \Phi(0.2)\)\quad D. \(\Phi(0.2)\)
\end{exercise}
\begin{solution}
因为 \(P\{X=0\}=P\{X=1\}=\dfrac12\)，所以每个样本变量都服从伯努利分布
\(
X_i\sim B\!\left(1,\frac12\right),
\)
并且 \(X_1,\dots,X_{100}\) 相互独立同分布。

设
\(
S=\sum_{i=1}^{100} X_i.
\)
则
\[
ES=100\times \frac12=50,\qquad DS=100\times \frac14=25.
\]

由于样本量 \(n=100\) 较大，可用中心极限定理把 \(S\) 近似为正态分布 \(N(50,25)\)。于是
\(
P(S \leq 55) = P\!\left(\frac{S - 50}{5} \leq \frac{55 - 50}{5}\right) \approx \Phi(1).
\)，
因此应选 B。
\end{solution}

\begin{exercise}\label{exer:5-climb}
在一次集体登山活动中，每个人意外受伤的概率为 1\%，各人是否受伤相互独立。\par
(1) 为保证没有人受伤的概率大于 0.90，应如何控制参加人数？\par
(2) 若有 100 人参加，求受伤人数不超过 2 人的概率的近似值（用中心极限定理，结果可用 \(\Phi\) 表示）。
\end{exercise}
\begin{solution}
(1) 设人数为 \(n\)，受伤人数 \(X \sim B(n, 0.01)\)。
要求 \(P(X = 0) = 0.99^n > 0.9\)。取对数：
\[
n\ln 0.99 > \ln 0.9 \implies n < \frac{\ln 0.9}{\ln 0.99} \approx 10.48,
\]
故人数控制在 \(\mathbf{10}\) 人及以内。\\
(2) \(n = 100\)，\(X \sim B(100, 0.01)\)，\(EX = 1\)，\(DX = 0.99\)。
\[
P(X \leq 2) = P\!\left(\frac{X-1}{\sqrt{0.99}} \leq \frac{1}{\sqrt{0.99}}\right) \approx \Phi\!\left(\frac{1}{\sqrt{0.99}}\right) = \Phi(\sqrt{100/99}).
\]
\end{solution}


\begin{exercise}\label{exer:5-clt-blood}
某地区某种疾病的患病率为 0.1，现对 10000 人进行血检普查（各人独立）。\par
(1) 用中心极限定理估计患病人数在 \([970, 1030]\) 内的概率（用 \(\Phi\) 表示）；\par
(2) 有两套化验方案：方案一逐一化验（10000 次）；方案二将每 \(k\) 人分为一组，先混合化验，阴性则只需 1 次，阳性则再逐一化验另半血样（共 \(k+1\) 次）。问 \(k\) 满足什么条件时方案二更优？\(k = 4\) 时方案二的平均化验次数是多少？
\end{exercise}
\begin{solution}
(1) 患病人数 \(X \sim B(10000, 0.1)\)，\(EX = 1000\)，\(DX = 900\)，\(\sigma = 30\)。
\[
P(970 \leq X \leq 1030) \approx \Phi\!\left(\frac{30}{30}\right) - \Phi\!\left(\frac{-30}{30}\right) = \Phi(1) - \Phi(-1) = 2\Phi(1) - 1.
\]
\begin{note}
严格使用连续性修正时为 \(P(969.5 < X < 1030.5)\)，因 \(n\) 很大（10000），修正可忽略。
\end{note}
(2) 每组 \(k\) 人全为阴性的概率为 \(0.9^k\)，每组化验次数 \(Y_i\) 满足：
\(
E(Y_i) = 1 \cdot 0.9^k + (k+1)(1-0.9^k) = k + 1 - k \cdot 0.9^k.
\)
总期望 \(E(T_2) = \dfrac{10000}{k}(k+1-k \cdot 0.9^k) = 10000\!\left(1 + \dfrac{1}{k} - 0.9^k\right)\)。\\
方案二更优的条件为 \(E(T_2) < 10000\)，即 \(k \cdot 0.9^k > 1\)。
当 \(k = 4\) 时：
\(
E(Y_i) = 5 - 4 \times 0.9^4 = 5 - 4 \times 0.6561 = 2.3756,
\)
\(
E(T_2) = \frac{10000}{4} \times 2.3756 = 5939.
\)
\end{solution}


\chapter{样本与抽样分布}


{\small
\[
\begin{cases}
\text{基本概念}
\begin{cases}
\text{总体与个体} \\
\text{样本}
\begin{cases}
\text{代表性：与总体具有相同概率分布} \\
\text{独立性：个体之间相互独立}
\end{cases}
\end{cases} \\
\text{统计量}
\begin{cases}
\text{概念：不含其它未知参数的样本函数} \\
\text{常见统计量：}
\begin{cases}
\text{样本均值：}\bar{X}=\displaystyle\frac{1}{n}\sum_{i=1}^n X_i \\
\text{样本方差：}S^2=\displaystyle\frac{1}{n-1}\sum_{i=1}^n (X_i-\bar{X})^2=\frac{1}{n-1}\left(\sum_{i=1}^n X_i^2-n\bar{X}^2\right) \\
\text{样本}k\text{阶原点矩：}A_k=\displaystyle\frac{1}{n}\sum_{i=1}^n X_i^k \\
\text{样本}k\text{阶中心距：}B_k=\displaystyle\frac{1}{n}\sum_{i=1}^n (X_i-\bar{X})^k,k=2,3\cdots
\end{cases} \\
\text{常用统计量的数字特征}
\begin{cases}
E\bar{X}=EX,D\bar{X}=\displaystyle\frac{DX}{n} \\
E(S^2)=DX
\end{cases}
\end{cases} \\
\text{抽样分布}
\begin{cases}
\chi^2\text{分布}
\begin{cases}
\text{形式：}n\text{个独立标准正态分布的平方和，自由度为}n \\
\text{性质：(1)独立可加性；(2)若}X\sim\chi^2(n)\text{,则}EX=n,DX=2n
\end{cases} \\
t\text{分布}
\begin{cases}
\text{形式：}Z=\displaystyle\frac{X}{\sqrt{Y/n}}\text{,其中}X,Y\text{独立，} \\
X\sim N(0,1),Y\sim\chi^2(n) \\
\text{性质：若}Z\sim t(n)\text{,则}EZ=0,DZ=\displaystyle\frac{n}{n-2}
\end{cases} \\
F\text{分布}
\begin{cases}
\text{形式：}F=\displaystyle\frac{X/n_1}{Y/n_2}\sim F(n_1,n_2)\text{,其中}X,Y\text{独立，} \\
X\sim\chi^2(n_1),Y\sim\chi^2(n_2) \\
\text{性质：若}X\sim t(n)\text{,则}X^2\sim F(1,n)
\end{cases} \\
\text{单正态总体抽样}
\begin{cases}
\displaystyle\frac{\bar{X}-\mu}{\sigma/\sqrt{n}}\sim N(0,1),\\
\frac{(n-1)S^2}{\sigma^2}\sim\chi^2(n-1),\quad
\bar{X}\text{与}S^2\text{独立}\\
\displaystyle\sum_{i=1}^n\left(\frac{X_i-\mu}{\sigma}\right)^2\sim\chi^2(n),\quad
\frac{\bar{X}-\mu}{S/\sqrt{n}}\sim t(n-1)
\end{cases} \\
\text{双正态总体抽样}\\
X\sim N(\mu_1,\sigma_1^2),Y\sim N(\mu_2,\sigma_2^2)
\end{cases}
\end{cases}
\]
}
\section{总体与样本，统计量}

\begin{definition}[总体与个体] \label{def:population}
研究对象的全体称为\textbf{总体}，组成总体的每一个元素称为\textbf{个体}。在对总体进行统计研究时，通常将总体与某个随机变量 \(X\) 等同起来，称“总体 \(X\)”，总体的分布即指随机变量 \(X\) 的分布。
\end{definition}

\begin{definition}[简单随机样本的分布函数] \label{def:sample}
设 \(X\) 是具有分布函数 \(F\) 的随机变量，若 \(X_1, X_2, \dots, X_n\) 是与 \(X\) 具有\textbf{同分布}函数 \(F\) 且\textbf{相互独立}的随机变量，则称 \(X_1, X_2, \dots, X_n\) 为从总体 \(X\)（或分布函数 \(F\)）得到的\textbf{容量为 \(n\) 的简单随机样本}，简称\textbf{样本}。它们的观察值 \(x_1, x_2, \dots, x_n\) 称为\textbf{样本值}。设总体 \(X\) 的分布函数为 \(F(x)\)，则样本 \(X_1, X_2, \dots, X_n\) 的联合分布函数为
\[
F(x_1, x_2, \dots, x_n) = \prod_{i=1}^{n} F(x_i).
\]
对于离散型总体，联合分布为 \[P\{X_1 = x_1, \dots, X_n = x_n\} = \prod_{i=1}^{n} P\{X_i = x_i\}\]
对于连续型总体，联合密度为 \[f(x_1, \dots, x_n) = \prod_{i=1}^{n} f(x_i)\]
\end{definition}

\begin{definition}[统计量] \label{def:statistic}
设 \(X_1, X_2, \dots, X_n\) 为来自总体 \(X\) 的一个样本，\(g(X_1, X_2, \dots, X_n)\) 是 \(X_1, X_2, \dots, X_n\) 的函数。若 \(g\) 中\textbf{不含任何未知参数}，则称 \(g(X_1, X_2, \dots, X_n)\) 为一个\textbf{统计量}。 \(\overline{X}\)、\(S^2\)、\(\max X_i\)、\(\min X_i\) 都可以是统计量；关键不在形式，而在是否含未知参数。比如 \(\sum (X_i-\mu)^2\) 在 \(\mu\) 已知时是统计量，但 \(\sum (X_i-\sigma)^2\) 若 \(\sigma\) 未知，就不是统计量。统计量和估计量不是同义词。估计量一定是统计量，因为它必须是不含未知参数的样本函数；但统计量未必专门拿来估参数。
\end{definition}

\begin{definition}[常用统计量] \label{def:common-stat}
设 \(X_1, X_2, \dots, X_n\) 是来自总体 \(X\) 的样本，\(\overline{X}\) 和 \(S^2\) 分别为样本均值与样本方差。
\begin{enumerate}
    \item \textbf{样本均值}：\(\displaystyle\overline{X} = \frac{1}{n}\sum_{i=1}^{n}X_i\)；
    \item \textbf{样本方差}：\[\displaystyle S^2 = \frac{1}{n-1}\sum_{i=1}^{n}(X_i - \overline{X})^2 = \frac{1}{n-1}\!\left(\sum_{i=1}^{n}X_i^2 - n\overline{X}^2\right)\]$n$ 个数据 $X_1,\dots,X_n$ 本来有 $n$ 个”自由度”（可以独立变化的方向）。但一旦你用 $\bar{X}$ 代替了真实均值 $\mu$，就多加了一个约束 $\sum(X_i-\bar{X})=0$，相当于”用掉了一个自由度”。剩下只有 $n-1$ 个独立的偏差，所以除以 $n-1$ 才是正确的”平均”。这也解释了为什么 $\frac{(n-1)S^2}{\sigma^2}\sim\chi^2(n-1)$ 而不是 $\chi^2(n)$——自由度正好反映了被”用掉”的那一个。
    \item \textbf{样本标准差}：\(S = \sqrt{S^2}\)；
    \item \textbf{样本 \(k\) 阶原点矩}：\(\displaystyle A_k = \frac{1}{n}\sum_{i=1}^{n}X_i^k\ (k=1,2,\dots)\)；
    \item \textbf{样本 \(k\) 阶中心矩}：\(\displaystyle B_k = \frac{1}{n}\sum_{i=1}^{n}(X_i - \overline{X})^k\ (k=2,3,\dots)\)。
\end{enumerate}
由大数定律，样本均值和样本方差依概率收敛于总体均值和总体方差：
\(\overline{X} \stackrel{P}{\longrightarrow} \mu\)，\(S^2 \stackrel{P}{\longrightarrow} \sigma^2\)，\(A_k \stackrel{P}{\longrightarrow} EX^k\)。
\end{definition}

\begin{property}[统计量的数字特征] \label{prop:stat-moment}
设总体 \(X\) 的期望为 \(EX = \mu\)，方差为 \(DX = \sigma^2\)，\(X_1, X_2, \dots, X_n\) 为容量为 \(n\) 的样本，则
\[
EX_i = \mu,\quad DX_i = \sigma^2,\quad E\overline{X} = \mu,\quad D\overline{X} = \frac{\sigma^2}{n},\quad E(S^2) = \sigma^2.
\]
\end{property}
\begin{proof}
\(E\overline{X} = \mu\) 和 \(D\overline{X} = \sigma^2/n\) 由期望和方差的性质直接可得。关键在于证明 \(E(S^2) = \sigma^2\)。令 \(Y_i = X_i - \mu\)，则
\[
\sum_{i=1}^{n}(X_i - \overline{X})^2 = \sum_{i=1}^{n}(Y_i - \overline{Y})^2 = \sum_{i=1}^{n} Y_i^2 - n\overline{Y}^2.
\implies
E\!\left[\sum(Y_i - \overline{Y})^2\right] = n\sigma^2 - n\cdot\frac{\sigma^2}{n} = (n-1)\sigma^2.
\]
因此
\[
E(S^2) = E\!\left[\frac{1}{n-1}\sum_{i=1}^{n}(X_i-\overline{X})^2\right] = \frac{(n-1)\sigma^2}{n-1} = \sigma^2.
\]
\end{proof}


\begin{exercise}\label{exer:src-2011-a-06-34credit}
设 \(X_1,X_2,\cdots,X_n\ (n\ge2)\) 为来自总体 \(N(0,\sigma^2)\) 的简单随机样本，其样本均值为 \(\overline X\)。记
\(
Y_i=X_i-\overline X,\qquad i=1,2,\cdots,n.
\)
求：\par
(1) \(Y_i\) 的方差 \(DY_i\ (i=1,2,\cdots,n)\)；\par
(2) \(Y_1\) 与 \(Y_n\) 的协方差 \(\operatorname{Cov}(Y_1,Y_n)\)；\par
(3) 若 \(c(Y_1+Y_n)^2\) 是 \(\sigma^2\) 的无偏估计量，求常数 \(c\)；\par
(4) 若 \(\hat\theta_1\) 与 \(\hat\theta_2\) 相互独立，均为 \(\sigma^2\) 的无偏估计量，且它们的方差存在，试给出一个比它们更有效的无偏估计量。
\end{exercise}
\begin{solution}
(1) 
\[
Y_i=X_i-\overline X
=\left(1-\frac1n\right)X_i-\frac1n\sum_{k\ne i}X_k,~~
DY_i=\left(1-\frac1n\right)^2\sigma^2+(n-1)\frac1{n^2}\sigma^2
=\frac{n-1}{n}\sigma^2.
\]
(2)
\begin{align*}
\operatorname{Cov}(Y_1,Y_n)&=\operatorname{Cov}(X_1-\overline X,\ X_n-\overline X)=\operatorname{Cov}(X_1,X_n)-\operatorname{Cov}(X_1,\overline X)-\operatorname{Cov}(X_n,\overline X)+\operatorname{Cov}(\overline X,\overline X)\\
&=0-\frac{\sigma^2}{n}-\frac{\sigma^2}{n}+\frac{\sigma^2}{n}=-\frac{\sigma^2}{n}.
\end{align*}
(3) 因为 \(E(Y_1+Y_n)=0\)，
\(
E\!\left[c(Y_1+Y_n)^2\right]
=cD(Y_1+Y_n).
\)
由前两问，
\[
D(Y_1+Y_n)
=DY_1+DY_n+2\operatorname{Cov}(Y_1,Y_n)
=\left(\frac{n-1}{n}+\frac{n-1}{n}-\frac2n\right)\sigma^2
=\frac{2(n-2)}{n}\sigma^2.
\]
因此当 \(n>2\) 时，要使 \(c(Y_1+Y_n)^2\) 为 \(\sigma^2\) 的无偏估计量，应有
\(
c\cdot\frac{2(n-2)}{n}=1\implies
c=\frac{n}{2(n-2)}.
\)\\
(4) 设
\(
D\hat\theta_1=V_1,\qquad D\hat\theta_2=V_2.
\)
要把两个无偏估计量合成一个新的无偏估计量，最自然先取线性组合
\(
\hat\theta_3=a\hat\theta_1+(1-a)\hat\theta_2.
\)，因为 \(E\hat\theta_1=E\hat\theta_2=\sigma^2\)，所以
\(
E\hat\theta_3=a\sigma^2+(1-a)\sigma^2=\sigma^2,
\)，无论 \(a\) 取何值，它都是 \(\sigma^2\) 的无偏估计量。接下来只需选 \(a\)，使方差尽量小。由 \(\hat\theta_1,\hat\theta_2\) 相互独立，
\[
D\hat\theta_3=a^2V_1+(1-a)^2V_2.\text{求导：}
2aV_1-2(1-a)V_2=0
\implies a=\frac{V_2}{V_1+V_2}.
\]
于是
\[
\hat\theta_3=\frac{V_2}{V_1+V_2}\hat\theta_1+\frac{V_1}{V_1+V_2}\hat\theta_2,~~
D\hat\theta_3
=\left(\frac{V_2}{V_1+V_2}\right)^2V_1
+\left(\frac{V_1}{V_1+V_2}\right)^2V_2
=\frac{V_1V_2}{V_1+V_2},
\]
当 \(V_1,V_2>0\) 时，它小于 \(V_1,V_2\)，所以 \(\hat\theta_3\) 比原来的两个估计量更有效。
\end{solution}

\begin{exercise}\label{exer:6-sk-sq}
已知总体 \(X\) 的期望 \(EX = 0\)，方差 \(DX = \sigma^2\)，从总体 \(X\) 中抽取容量为 \(n\) 的简单随机样本，其样本均值为 \(\overline{X}\)，样本二阶中心矩为
\[
B_2=\frac1n\sum_{i=1}^n(X_i-\overline X)^2.
\]
记 \(T_k^2=k\overline{X}^2+kB_2\ (k=1,2,3,4)\)，则（\quad）。\par
A. \(E(T_1^2) = \sigma^2\)\quad B. \(E(T_2^2) = \sigma^2\)\quad C. \(E(T_3^2) = \sigma^2\)\quad D. \(E(T_4^2) = \sigma^2\)
\end{exercise}

\begin{solution}
先利用恒等式
\[
\frac1n\sum_{i=1}^nX_i^2
=\frac1n\sum_{i=1}^n\bigl[(X_i-\overline X)+\overline X\bigr]^2
=B_2+\overline X^2.
\]
由于 \(EX=0,\ DX=\sigma^2\)，所以 \(EX^2=DX+(EX)^2=\sigma^2\)。因此
\[
E(\overline X^2+B_2)
=E\left(\frac1n\sum_{i=1}^nX_i^2\right)
=\frac1n\sum_{i=1}^nEX_i^2
=\sigma^2.
\]
于是
\[
E(T_k^2)=kE(\overline X^2+B_2)=k\sigma^2.
\]
要使 \(E(T_k^2)=\sigma^2\)，只能取 \(k=1\)。故选 A。
\end{solution}


\section{三大抽样分布}


\begin{definition}[\texorpdfstring{\(\chi^2\)}{χ²} 分布] \label{def:chi2}
如果总体正态：
\(
X_1,\dots,X_n\sim N(\mu,\sigma^2),
\)，那么标准化后
\(
\frac{X_i-\mu}{\sigma}\sim N(0,1).
\)，如果 \(\mu\) 已知，那么
\(
\sum_{i=1}^n\left(\frac{X_i-\mu}{\sigma}\right)^2\sim \chi^2(n).
\)，这就是 \(n\) 个标准正态平方和。但实际中 \(\mu\) 往往未知，要用 \(\overline X\) 代替 \(\mu\)。这会消耗一个自由度，因为偏差满足约束：
\(
\sum_{i=1}^n (X_i-\overline X)=0.
\)，所以真正独立的偏差只有 \(n-1\) 个，于是有经典结论：
\[
\sum_{i=1}^n(\frac{X_i-\overline X}{\sigma})^2=\frac{(n-1)S^2}{\sigma^2}\sim \chi^2(n-1).
\]
设 \(X_1, X_2, \dots, X_n\) 相互独立且都服从标准正态分布 \(N(0,1)\)，令
\[
\chi^2 = X_1^2 + X_2^2 + \cdots + X_n^2,
\]
则称 \(\chi^2\) 服从自由度为 \(n\) 的 \(\chi^2\) 分布，记为 \(\chi^2 \sim \chi^2(n)\)。
\end{definition}

\begin{property}\label{prop:chi2-property} 
\begin{enumerate}
    \item \(E[\chi^2(n)] = n\)，\(D[\chi^2(n)] = 2n\)。
    \item \textbf{可加性}：若 \(X \sim \chi^2(m)\)，\(Y \sim \chi^2(n)\)，且 \(X, Y\) 相互独立，则 \(X + Y \sim \chi^2(m+n)\)。
\end{enumerate}
\end{property}



\begin{definition}[\texorpdfstring{\(t\)}{t} 分布] \label{def:t-dist}
如果总体正态\(X_1,\dots,X_n\sim N(\mu,\sigma^2),\)并且 \(\sigma\) 已知，那么样本均值可以标准化为\(Z=\frac{n\overline X-n\mu}{\sigma\sqrt n}=\frac{\overline X-\mu}{\sigma/\sqrt n}\sim N(0,1)\)，但现实里 \(\sigma\) 往往未知。用样本标准差 \(S\) 代替 \(\sigma\)，于是构造
\(
T=\frac{\overline X-\mu}{S/\sqrt n}.
\)，
由
\[
\frac{\overline X-\mu}{\sigma/\sqrt n}\sim N(0,1),~~
\frac{(n-1)S^2}{\sigma^2}\sim \chi^2(n-1),
\]
所以
\[
T=\frac{\overline X-\mu}{S/\sqrt n}=\frac{Z}{\sqrt{U/(n-1)}},
~~~Z\sim N(0,1),U\sim \chi^2(n-1),
\]且 \(Z\) 与 \(U\) 独立。这个比值分布就被定义为自由度 \(n-1\) 的 \(t\) 分布：
\[
T\sim t(n-1).
\]
设 \(X \sim N(0,1)\)，\(Y \sim \chi^2(n)\)，且 \(X, Y\) 相互独立，令
\[
t = \frac{X}{\sqrt{Y/n}},
\]
则称 \(t\) 服从自由度为 \(n\) 的 \(t\) 分布，记为 \(t \sim t(n)\)。
\end{definition}

\begin{property}\label{prop:t-dist-property}
\begin{enumerate}
    \item \(t\) 分布的概率密度关于 \(x = 0\) 对称，\(Et = 0\ (n \geq 2)\)，\(Dt = \dfrac{n}{n-2}\ (n > 2)\)。
    \item 当 \(n\) 足够大时，\(t\) 分布近似于标准正态分布；\(n\) 越大，曲线越"瘦"越"高"。
    \item 由对称性：\(t_{1-\alpha}(n) = -t_\alpha(n)\)。
    \item 若 \(t \sim t(n)\)，则 \(t^2 \sim F(1,n)\)。
\end{enumerate}
\end{property}

\begin{definition}[\texorpdfstring{\(F\)}{F} 分布] \label{def:f-dist}
在正态总体下，有经典结论：
\[
\frac{(n_1-1)S_1^2}{\sigma_1^2}\sim \chi^2(n_1-1),~~~
\frac{(n_2-1)S_2^2}{\sigma_2^2}\sim \chi^2(n_2-1),
\]并且如果两组样本独立，那么这两个卡方变量也独立。于是当原假设 \(\sigma_1^2=\sigma_2^2\) 成立时
\[
\frac{S_1^2}{S_2^2}
=
\frac{\chi^2(n_1-1)/(n_1-1)}
{\chi^2(n_2-1)/(n_2-1)}.
\]于是人们就把这种形式定义成 F 分布：
\[
F=\frac{U/m}{V/n},~~~U\sim \chi^2(m), V\sim \chi^2(n)
\]
其中\(U,V\) 相互独立， \(m\) 为第一自由度，\(n\) 为第二自由度。
\end{definition}

\begin{property}\label{prop:f-dist-property}
\begin{enumerate}
    \item 若 \(F \sim F(n_1, n_2)\)，则 \(\dfrac{1}{F} \sim F(n_2, n_1)\)。
    \item \(F_{1-\alpha}(n_1, n_2) = \dfrac{1}{F_\alpha(n_2, n_1)}\)，常用于求 \(F\) 分布表中未列出的分位数。
    \item 若 \(t \sim t(n)\)，则 \(t^2 \sim F(1, n)\)。
\end{enumerate}
\end{property}


\begin{exercise}\label{exer:src-paper1-fill-05}\label{prob:p1-fill-05}
设 \((X_1,X_2,\dots,X_6)\) 是来自正态分布 \(N(0,\sigma^2)\) 的样本，
\[
Y=\left(X_1+X_2+X_3\right)^2+
\left(X_4+X_5+X_6\right)^2.
\]
当 \(c=\underline{\hspace{1.5cm}}\) 时，\(cY\) 服从 \(\chi^2\) 分布，且 \(E(\chi^2)=\underline{\hspace{1.5cm}}\)。
\end{exercise}
\begin{solution}
\[
\begin{cases}X_1+X_2+X_3\sim N(0,3\sigma^2)\\X_4+X_5+X_6\sim N(0,3\sigma^2)\end{cases}\implies\frac{X_1+X_2+X_3}{\sqrt{3}\sigma}\sim N(0,1),\frac{X_4+X_5+X_6}{\sqrt{3}\sigma}\sim N(0,1)
\]
\[(\frac{X_1+X_2+X_3}{\sqrt{3}\sigma})^2+(\frac{X_4+X_5+X_6}{\sqrt{3}\sigma})^2=\frac{Y}{3\sigma^2}\sim \chi^2(2),\text{期望等于自由度}2\]
\end{solution}



\begin{exercise}\label{exer:6-chi2-ab}
设 \(X_1, X_2, X_3, X_4\) 是来自正态总体 \(N(0, 4)\) 的简单随机样本，记
\(
X = a(X_1 - 2X_2)^2 + b(3X_3 - 4X_4)^2,
\)
若统计量 \(X\) 服从自由度为 2 的 \(\chi^2\) 分布，则 \(ab =\,\underline{\hspace{2cm}}\)。
\end{exercise}

\begin{solution}
\[
\begin{cases}
X_1-2X_2\sim N(0,4+(-2)^24)=N(0,20),\\
3X_3-4X_4\sim N(0,3^24+(-4)^24)=N(0,100),
\end{cases}
\implies
\frac{X_1-2X_2}{\sqrt{20}}\sim N(0,1),
\frac{3X_3-4X_4}{10}\sim N(0,1).
\]
\[
\left(\frac{X_1-2X_2}{\sqrt{20}}\right)^2+
\left(\frac{3X_3-4X_4}{10}\right)^2
=\frac{1}{20}(X_1-2X_2)^2+\frac{1}{100}(3X_3-4X_4)^2
\sim \chi^2(2).
\]
因此 \(a=\dfrac1{20},\ b=\dfrac1{100}\)，故
\(
ab=\dfrac1{2000}.
\)
\end{solution}



\begin{exercise}\label{exer:src-2007-a-15}
设随机变量 \(\xi_1,\xi_2,\dots,\xi_6\) 独立同分布，且 \(\xi_i\sim N(\mu,\sigma^2)\)。令
\(
Z=\frac1{\sigma^2}\sum_{i=1}^6(\xi_i-\mu)^2,
\)
则 \(Z\) 的概率密度为（\quad）。\par
A. \(\displaystyle f(z)=
\begin{cases}
\dfrac1{16}z^2e^{z/2}, & z>0,\\
0, & z\le0,
\end{cases}\)~~~~~~~
B. \(\displaystyle f(z)=\dfrac1{\sqrt{12\pi}}e^{z^2/12},\quad -\infty<z<+\infty\)\par
C. \(\displaystyle f(z)=\dfrac1{\sqrt{12\pi}}e^{-z^2/12},\quad -\infty<z<+\infty\)~~~~~~~~
D. \(\displaystyle f(z)=
\begin{cases}
\dfrac1{16}z^2e^{-z/2}, & z>0,\\
0, & z\le0.
\end{cases}\)
\end{exercise}
\begin{solution}
注意$\mu$不是$\overline{X}$，因为
\(
\frac{\xi_i-\mu}{\sigma}\sim N(0,1),
\)
且六个变量相互独立，所以
\(
Z=\sum_{i=1}^6\left(\frac{\xi_i-\mu}{\sigma}\right)^2\sim\chi^2(6).
\)
\(\chi^2(n)\) 的密度为
\(
f(z)=\frac{1}{2^{n/2}\Gamma(n/2)}z^{n/2-1}e^{-z/2},\qquad z>0.
\)
取 \(n=6\)，得
\[
f_Z(z)=\frac{1}{2^3\Gamma(3)}z^2e^{-z/2}
=\frac{1}{8\cdot2}z^2e^{-z/2}
=\frac1{16}z^2e^{-z/2},\qquad z>0.
\]
\end{solution}



\begin{exercise}[\texorpdfstring{\(t\)}{t} 分布的构造] \label{ex:sampling_04}
设 \(X_1,X_2,X_3,X_4\) 为来自 \(N(0,1/2)\) 的简单随机样本，则统计量
\(
Y=\frac{X_1-X_2}{\sqrt{X_3^2+X_4^2}}
\)
服从（\quad）。\par
A. \(N(0,1)\)\qquad B. \(\chi^2(2)\)\qquad C. \(t(2)\)\qquad D. \(F(2,2)\)
\end{exercise}
\begin{solution}
应选 C。\textbf{分子}：\(X_1-X_2\sim N(0,1/2+1/2)=N(0,1)\)，故 \(X_1-X_2\sim N(0,1)\)。\\
\textbf{分母}：
\[\sqrt{X_3^2+X_4^2}=\sqrt{\frac12}\sqrt{\frac{X_3^2}{\frac12}+\frac{X_4^2}{\frac12}}\sim\sqrt{\frac{\chi^2(2)}{2}}.\]
\end{solution}


\begin{exercise}\label{exer:6-t-new}
设 \(X_1, X_2, \dots, X_n\) 为来自正态分布 \(N(\mu, \sigma^2)\) 的简单随机样本，\(\overline{X}\) 和 \(S^2\) 分别为样本均值和样本方差。又设 \(X_{n+1} \sim N(\mu, \sigma^2)\) 且与 \(X_1, \dots, X_n\) 相互独立，则统计量
\(
\frac{X_{n+1} - \overline{X}}{S}\sqrt{\frac{n}{n+1}}
\)
服从\,\underline{\hspace{2cm}}分布（注明自由度）。
\end{exercise}

\begin{solution}
\(\overline{X} \sim N(\mu, \sigma^2/n)\)，\(X_{n+1} \sim N(\mu, \sigma^2)\)，且相互独立，故
\[
X_{n+1} - \overline{X} \sim N\!\left(0,\, \left(1+\frac{1}{n}\right)\sigma^2\right),\quad \frac{X_{n+1}-\overline{X}}{\sigma\sqrt{1+1/n}} \sim N(0,1).
\]

又 \(\dfrac{(n-1)S^2}{\sigma^2} \sim \chi^2(n-1)\)，且与 \(\overline{X}\) 独立，故
\[
\frac{X_{n+1}-\overline{X}}{S}\sqrt{\frac{n}{n+1}} = \frac{(X_{n+1}-\overline{X})\big/\!\left(\sigma\sqrt{1+\frac{1}{n}}\right)}{\sqrt{(n-1)S^2/[\sigma^2(n-1)]}} \sim t(n-1).
\]
\end{solution}

\begin{exercise}\label{exer:6-t-y-sq}
设随机变量 \(X \sim t(n)\)，\(Y \sim F(1,n)\)，给定 \(\alpha\ (0 < \alpha < 0.5)\)，常数 \(c\) 满足 \(P\{X > c\} = \alpha\)，则 \(P\{Y > c^2\} =\,\underline{\hspace{2cm}}\)。\par
A. \(\alpha\)\quad B. \(1-\alpha\)\quad C. \(2\alpha\)\quad D. \(1-2\alpha\)
\end{exercise}
\begin{solution}
由 \(X \sim t(n)\) 知
\(
X^2 \sim F(1,n).
\)
因此 \(Y\) 与 \(X^2\) 同分布，题目所求概率可改写为
\(
P\{Y > c^2\} = P\{X^2 > c^2\}.
\)
\(
\text{而 } X^2>c^2 \Longleftrightarrow X>c \text{ 或 } X<-c.
\)
又因为 \(t\) 分布关于 0 对称，且题设给出 \(P\{X>c\}=\alpha\)，所以
\(
P\{X<-c\}=\alpha.
\)
故
\(
P\{Y > c^2\} = P\{X > c\} + P\{X < -c\} = \alpha + \alpha = 2\alpha.
\)
\end{solution}




\begin{exercise}\label{exer:src-paper1-choice-04}\label{prob:p1-choice-04}
设 \((X_1,X_2,\dots,X_n)\) 为总体 \(X\sim N(0,1)\) 的一个样本，\(\overline X\) 为样本均值，\(S^2\) 为样本方差，则有（\quad）。\par
A. \(\overline X\sim N(0,1)\)\quad
B. \(n\overline X\sim N(0,1)\)\par
C. \(\overline X/S\sim t(n-1)\)\quad
D. \(\dfrac{(n-1)X_1^2}{\sum_{i=2}^{n}X_i^2}\sim F(1,n-1)\)
\end{exercise}
\begin{solution}
应选 D。

因为 \(X_1\sim N(0,1)\)，所以 \(X_1^2\sim\chi^2(1)\)。又由于样本相互独立，
\(
\sum_{i=2}^{n}X_i^2\sim\chi^2(n-1),
\)
并且 \(X_1^2\) 与 \(\sum_{i=2}^{n}X_i^2\) 相互独立。因此
\[
\frac{X_1^2/1}{\left(\sum_{i=2}^{n}X_i^2\right)/(n-1)}
=\frac{(n-1)X_1^2}{\sum_{i=2}^{n}X_i^2}
\sim F(1,n-1).
\]
A、B 错在 \(\overline X\sim N(0,1/n)\)，而 \(n\overline X\sim N(0,n)\)。C 中若使用样本标准差，应有 \(\sqrt n\,\overline X/S\sim t(n-1)\)。
\end{solution}

\section{正态总体下的抽样分布}

\subsection{单个正态总体}

\begin{theorem}[单正态总体抽样分布] \label{thm:normal-single}
设总体 \(X \sim N(\mu, \sigma^2)\)，\(X_1, X_2, \dots, X_n\) 为来自 \(X\) 的样本，\(\overline{X}, S^2\) 分别为样本均值和样本方差，则
\begin{enumerate}
    \item \(\displaystyle\frac{\overline{X} - \mu}{\sigma/\sqrt{n}} \sim N(0,1)\)，即 \(\overline{X} \sim N\!\left(\mu, \frac{\sigma^2}{n}\right)\)；
    \item \(\displaystyle\frac{(n-1)S^2}{\sigma^2} = \sum_{i=1}^{n}\!\left(\frac{X_i - \overline{X}}{\sigma}\right)^{\!2} \sim \chi^2(n-1)\)，且 \(\overline{X}\) 与 \(S^2\) 相互独立；
    \item \(\displaystyle\frac{\overline{X} - \mu}{S/\sqrt{n}} \sim t(n-1)\)（当 \(\sigma\) 未知时，用 \(S\) 替代 \(\sigma\)）。
\end{enumerate}
\end{theorem}

\begin{note}
这个定理是后面所有 \(t\) 检验、\(\chi^2\) 检验、区间估计公式的出发点，因此要特别清楚它\textbf{什么时候能直接用}：前提必须是总体本身服从正态分布。若总体不正态，则“\(\overline{X}\) 与 \(S^2\) 独立”“\((n-1)S^2/\sigma^2\sim\chi^2(n-1)\)”这两条一般不能直接照搬。

定理中结论(2)的证明涉及矩阵知识，不要求掌握，但其结论必须牢记。\textbf{注意自由度}：\(\sum(X_i-\overline{X})^2/\sigma^2\) 的自由度为 \(n-1\)（而非 \(n\)），因为 \(\overline{X}\) 已经消耗了一个自由度。进一步有 \(\displaystyle\frac{n(\overline{X}-\mu)^2}{S^2} \sim F(1,n-1)\)。常见误用是把“总体均值未知”和“总体均值恰好等于 0”混成一回事，从而把自由度写错。
\textbf{自学追问：} 为什么这个定理几乎”统治”第 6 章？因为区间估计、\(t\) 检验、\(\chi^2\) 检验、\(F\) 检验的统计量都由它直接导出。
\end{note}

\begin{note}
\textbf{$\bar X$ 与 $S^2$ 独立的直觉解释：}

这个结论违反直觉——$S^2$ 的公式里明明含有 $\bar X$，怎么可能和 $\bar X$ 独立？

\textbf{几何视角：} 把样本 $(X_1,\dots,X_n)$ 看作 $\mathbb{R}^n$ 中的一个点。$\bar X$ 是该点在”全1方向” $\mathbf{1}=(1,\dots,1)$ 上的投影长度（除以 $\sqrt{n}$）；而 $S^2$ 度量的是该点到这条直线的距离。在正态分布下，球对称性保证了”沿某方向的投影”与”到该方向的距离”相互独立。

\textbf{为什么仅正态成立？} 因为只有正态分布的联合密度具有球对称性（等密度面是球面）。换成其他分布（如均匀分布），等密度面不再是球面，投影和距离就不再独立。

\textbf{考试意义：} 凡是需要构造 $t=\dfrac{\bar X-\mu}{S/\sqrt{n}}$ 的地方，都隐含使用了”分子的 $\bar X$ 与分母的 $S$ 独立”这一前提。若题目总体不是正态，这个 $t$ 统计量就不服从 $t$ 分布。
\end{note}

\begin{exercise}\label{exer:6-s2-chi2}
设 \(X_1, X_2, \dots, X_n\) 是总体 \(X \sim N(\mu, \sigma^2)\) 的简单随机样本，\(\overline{X}\) 为样本均值，则下列结论正确的是（\quad）。\par
A. \(\displaystyle\frac{1}{\sigma^2}\sum_{i=1}^{n}(X_i-\overline{X})^2 \sim \chi^2(n-1)\)\qquad B. \(\displaystyle\frac{1}{n}\sum_{i=1}^{n}(X_i-\overline{X})^2 \sim \chi^2(n-1)\)\par
C. \(\displaystyle\frac{1}{n-1}\sum_{i=1}^{n}(X_i-\overline{X})^2 \sim \chi^2(n-1)\)\qquad D. \(\displaystyle\frac{1}{\sigma^2}\sum_{i=1}^{n}(X_i-\overline{X})^2 \sim \chi^2(n)\)
\end{exercise}
\begin{solution}
应选 A。

这道题的核心判据只有一句：
\[
\frac{(n-1)S^2}{\sigma^2}
=\frac{1}{\sigma^2}\sum_{i=1}^{n}(X_i-\overline{X})^2
\sim \chi^2(n-1).
\]

于是 A 与定理完全一致，正确。

B 错在漏掉了 \(\sigma^2\)。\(\dfrac{1}{n}\sum(X_i-\overline{X})^2\) 本身只是一个样本方差型统计量，不会直接服从 \(\chi^2(n-1)\)。

C 也错在没有除以 \(\sigma^2\)。\(\dfrac{1}{n-1}\sum(X_i-\overline{X})^2=S^2\)，它不是卡方分布，而是 \(\chi^2(n-1)\) 再乘上 \(\dfrac{\sigma^2}{n-1}\) 得到的。

D 错在自由度写成了 \(n\)。只要出现 \(\overline{X}\)，就意味着损失了一个自由度，因此应为 \(n-1\)。

所以正确选项是 A。
\end{solution}

\begin{exercise}[样本方差的期望] \label{ex:statistic_05}
设总体 \(X\) 的概率密度为 \(f(x)=\dfrac{1}{2}e^{-|x|}\)（\(-\infty<x<+\infty\)），\(X_1,X_2,\dots,X_n\) 为简单随机样本，样本方差为 \(S^2\)，则 \(E(S^2)=\)\underline{\qquad}。
\end{exercise}
\begin{solution}
由 \(E(S^2)=DX\)，先计算总体矩：
\[
EX=\int_{-\infty}^{+\infty}x\cdot\frac{1}{2}e^{-|x|}\,\mathrm{d}x=0 \quad(\text{奇函数}).
\]
\[
EX^2=\int_{-\infty}^{+\infty}x^2\cdot\frac{1}{2}e^{-|x|}\,\mathrm{d}x
=2\int_0^{+\infty}x^2\cdot\frac{1}{2}e^{-x}\,\mathrm{d}x
=\int_0^{+\infty}x^2\,e^{-x}\,\mathrm{d}x=\Gamma(3)=2.
\]
\(
DX=EX^2-(EX)^2=2,\qquad E(S^2)=DX=2.
\)
\begin{note}
这里利用了拉普拉斯分布 \(f(x)=\frac{1}{2}e^{-|x|}\) 的矩公式，其本质是 \(\text{Exp}(1)\) 分布的偶函数延拓。
\end{note}
\end{solution}

\begin{exercise}\label{exer:src-2008-01-3c-iii-05}\label{exer:src-paper4-iii-05}
设 \(X_1,X_2,\dots,X_n\) 是来自总体 \(\chi^2(n)\) 分布的样本，\(\overline X\) 是样本均值。求
\(
E(\overline X),\qquad D(\overline X).
\)
\end{exercise}
\begin{solution}
若 \(X_i\sim\chi^2(n)\)，则
\(
E(X_i)=n,\qquad D(X_i)=2n.
\)
由于 \(X_1,\dots,X_n\) 为样本，相互独立同分布，所以
\(
E(\overline X)=E(X_1)=n,
\)
并且
\(
D(\overline X)=\frac{D(X_1)}{n}
=\frac{2n}{n}=2.
\)
故
\(
E(\overline X)=n,\qquad D(\overline X)=2.
\)
\end{solution}


\begin{exercise}\label{exer:6-normal-judge}
设 \(X_1, X_2, \dots, X_n\) 是来自总体 \(X \sim N(\mu, \sigma^2)\) 的简单随机样本，\(\overline{X}\) 与 \(S^2\) 分别为其样本均值和样本方差，则下列结论正确的是（\quad）。\par
A. \(2X_2 - X_1 \sim N(\mu, \sigma^2)\)\quad B. \(\dfrac{n(\overline{X}-\mu)^2}{S^2} \sim F(1,n-1)\)\par
C. \(\dfrac{S^2}{\sigma^2} \sim \chi^2(n-1)\)\quad D. \(\dfrac{\overline{X}-\mu}{S}\sqrt{n-1} \sim t(n-1)\)
\end{exercise}
\begin{solution}
应选 B。

这类题的关键，是先回忆正态总体抽样分布的三个标准结论：
\[
\overline{X}\sim N\!\left(\mu,\frac{\sigma^2}{n}\right),\qquad
\frac{(n-1)S^2}{\sigma^2}\sim\chi^2(n-1),\qquad
\frac{\sqrt{n}(\overline{X}-\mu)}{S}\sim t(n-1).
\]
然后逐项判断。

对 A，\(2X_2-X_1\) 仍服从正态分布，但
\(
E(2X_2-X_1)=2\mu-\mu=\mu,
\)
\(
D(2X_2-X_1)=4DX_2+DX_1=4\sigma^2+\sigma^2=5\sigma^2.
\)
所以它应服从 \(N(\mu,5\sigma^2)\)，不是 \(N(\mu,\sigma^2)\)，故 A 错。

对 B，由
\(
\frac{\sqrt{n}(\overline{X}-\mu)}{S}\sim t(n-1),
\)
两边平方可得
\(
\frac{n(\overline{X}-\mu)^2}{S^2}\sim F(1,n-1),
\)
所以 B 对。

对 C，正确的卡方结论应为
\(
\frac{(n-1)S^2}{\sigma^2}\sim\chi^2(n-1),
\)
题中少了系数 \(n-1\)，故 C 错。

对 D，\(t\) 统计量的正确形式应为
\(
\frac{\sqrt{n}(\overline{X}-\mu)}{S}\sim t(n-1),
\)
前面的系数应是 \(\sqrt{n}\)，不是 \(\sqrt{n-1}\)，故 D 错。

综上，正确选项是 B。
\end{solution}



\begin{exercise}\label{exer:6-esum-sq}
设总体 \(X \sim N(1, 4)\)，\(X_1, X_2, \dots, X_n\) 为来自总体 \(X\) 的样本，\(\overline{X}\) 为样本均值。则
\[
E\!\left(\sum_{i=1}^{n}(X_i-\overline{X})^2\right)=\underline{\hspace{2cm}}.
\]
A. \(4n\)\quad B. \(4(n-1)\)\quad C. \(2n\)\quad D. \(2(n-1)\)
\end{exercise}

\begin{solution}
应选 B。

题目要求的是
\(
E\!\left(\sum_{i=1}^{n}(X_i-\overline{X})^2\right),
\)
而不是 \(E(S^2)\)。最方便的做法是直接套正态总体下的抽样分布结论：
\[
\frac{1}{\sigma^2}\sum_{i=1}^{n}(X_i-\overline{X})^2 \sim \chi^2(n-1).
\]

因为题中总体 \(X\sim N(1,4)\)，所以 \(\sigma^2=4\)。两边取期望，利用 \(\chi^2(n-1)\) 的期望是 \(n-1\)，得
\(
E\!\left[\frac{1}{4}\sum_{i=1}^{n}(X_i-\overline{X})^2\right]=n-1.
\)
于是
\(
E\!\left(\sum_{i=1}^{n}(X_i-\overline{X})^2\right)=4(n-1).
\)

因此应选 B。
\end{solution}

\begin{exercise}\label{exer:6-y-t-dist}
设总体 \(X \sim N(0, 3^2)\)，\(X_1, X_2, \dots, X_8\) 是来自总体 \(X\) 的简单随机样本，则统计量
\(
Y = \frac{X_1 + X_2 + X_3 + X_4}{\sqrt{X_5^2 + X_6^2 + X_7^2 + X_8^2}}
\)
服从（\quad）。\par
A. \(\chi^2(4)\)\quad B. \(t(4)\)\quad C. \(F(4,4)\)\quad D. \(N(0, 4^2)\)
\end{exercise}

\begin{solution}
应选 B。

要判断 \(Y\) 的分布，关键是把它整理成
\(
\frac{U}{\sqrt{V/\nu}}
\)
的形式，其中 \(U\sim N(0,1)\)、\(V\sim\chi^2(\nu)\)，且 \(U,V\) 相互独立。

\(
X_1+X_2+X_3+X_4\sim N(0,4\times 9)=N(0,36),
\)
\[
U = \frac{X_1+X_2+X_3+X_4}{\sqrt{4 \times 9}} = \frac{X_1+X_2+X_3+X_4}{6} \sim N(0,1),
\]
\(
V = \frac{X_5^2+X_6^2+X_7^2+X_8^2}{9} \sim \chi^2(4),
\)
这是因为 \(\dfrac{X_i}{3}\sim N(0,1)\ (i=5,6,7,8)\)，四个标准正态平方和服从自由度为 4 的卡方分布。

又由于 \(U\) 只由 \(X_1,\dots,X_4\) 构成，\(V\) 只由 \(X_5,\dots,X_8\) 构成，而原样本彼此独立，所以 \(U,V\) 相互独立。故
\[
Y = \frac{(X_1+X_2+X_3+X_4)/6}{\sqrt{(X_5^2+\cdots+X_8^2)/4}} = \frac{U}{\sqrt{V/4}} \sim t(4).
\]

因此 \(Y\sim t(4)\)，应选 B。
\end{solution}


\begin{exercise}\label{exer:6-cov-t}
设 \(X_1, X_2\) 是取自正态总体 \(N(\mu, \sigma^2)\) 的简单随机样本，\(Y_1=X_1+X_2\)，\(Y_2=X_1-X_2\)。则 \(\mathrm{Cov}(Y_1, Y_2) =\,\underline{\hspace{2cm}}\)；已知 \((Y_1, Y_2)\) 服从二维正态分布，若 \(c\) 为非零常数，则当 \(c =\,\underline{\hspace{2cm}}\) 时，\(\dfrac{c(Y_1-2\mu)}{|Y_2|}\) 服从自由度为\,\underline{\hspace{2cm}}\,的\,\underline{\hspace{2cm}}\,分布。
\end{exercise}
\begin{solution}
第一空填 \(0\)：
\[
\mathrm{Cov}(Y_1, Y_2) = \mathrm{Cov}(X_1+X_2,\, X_1-X_2) = DX_1 - DX_2 = \sigma^2 - \sigma^2 = 0.
\]
第二空填 \(1\)，第三空填 \(1\)，第四空填 \(t\)。由 \((Y_1,Y_2)\) 服从二维正态分布且协方差为 0，知 \(Y_1, Y_2\) 相互独立。\(Y_1 \sim N(2\mu, 2\sigma^2)\)，故 \(\dfrac{Y_1-2\mu}{\sqrt{2}\,\sigma} \sim N(0,1)\)；\(Y_2 \sim N(0, 2\sigma^2)\)，故 \(\left(\dfrac{Y_2}{\sqrt{2}\,\sigma}\right)^{\!2} \sim \chi^2(1)\)。
构造 \(t\) 分布：
\[
t = \frac{(Y_1-2\mu)/(\sqrt{2}\,\sigma)}{\sqrt{(Y_2/(\sqrt{2}\,\sigma))^2}} = \frac{Y_1-2\mu}{|Y_2|} \sim t(1),
\]
故 \(c = 1\)。
\end{solution}


\begin{exercise}\label{exer:6-f-xbar}
设 \((X_1, X_2, \dots, X_{10})\) 是取自正态总体 \(N(0, 4)\) 的一个样本，\(\overline{X}\) 为样本均值，\(S^2\) 为样本方差。当 \(c =\,\underline{\hspace{2cm}}\) 时，\(c\,\dfrac{\overline{X}^2}{S^2}\) 服从 \(F\) 分布。
\end{exercise}
\begin{solution}
应填 \(10\)。

因为总体 \(X\sim N(0,4)\)，样本量 \(n=10\)，所以
\(
\overline{X}\sim N\!\left(0,\frac{4}{10}\right).
\)
标准化得
\(
\frac{\sqrt{10}\,\overline{X}}{2}\sim N(0,1),
\)
因此
\(
\frac{10\overline{X}^2}{4}\sim\chi^2(1).
\)

另一方面，在正态总体下，
\(
\frac{(n-1)S^2}{\sigma^2}=\frac{9S^2}{4}\sim\chi^2(9),
\)
并且 \(\overline{X}\) 与 \(S^2\) 相互独立。

故
\[
F = \frac{\chi^2(1)/1}{\chi^2(9)/9} = \frac{10\overline{X}^2/4}{S^2/4} = 10\,\frac{\overline{X}^2}{S^2} \sim F(1, 9).
\]

所以所求常数为
\(
c=10.
\)
\end{solution}

\begin{exercise}\label{exer:6-t-prob}
设 \(X_1, X_2, \dots, X_8\) 是来自正态总体 \(N(0, 4)\) 的简单随机样本。常数 \(c =\,\underline{\hspace{2cm}}\) 时，统计量 \(Y = \dfrac{c(X_1-X_2)}{\sqrt{\sum_{i=3}^{8}X_i^2}} \sim\,\underline{\hspace{2cm}}\)；\(P\!\left\{|\overline{X}| \leq \dfrac{\sqrt{2}}{2}\right\} =\,\underline{\hspace{2cm}}\)。
\end{exercise}
\begin{solution}
第一空填 \(\sqrt{3}\)，第二空填 \(t(6)\)，第三空填 \(2\Phi(1)-1\)。

先看统计量 \(Y\)。因为 \(X_1,X_2\) 独立，且都服从 \(N(0,4)\)，所以
\(
X_1-X_2\sim N(0,4+4)=N(0,8),
\)
从而
\(
\frac{X_1-X_2}{\sqrt{8}}\sim N(0,1).
\)

再看分母。对 \(i=3,\dots,8\)，有 \(\dfrac{X_i}{2}\sim N(0,1)\)，故
\(
\sum_{i=3}^{8}X_i^2
=4\sum_{i=3}^{8}\left(\frac{X_i}{2}\right)^2
\sim 4\chi^2(6).
\)
又因为分子只含 \(X_1,X_2\)，分母只含 \(X_3,\dots,X_8\)，所以它们相互独立。

由 \(t\) 分布定义：
\[
t = \frac{(X_1-X_2)/\sqrt{8}}{\sqrt{\sum_{i=3}^{8}X_i^2/(4 \times 6)}} = \sqrt{\frac{24}{8}}\,\frac{X_1-X_2}{\sqrt{\sum_{i=3}^{8}X_i^2}} = \sqrt{3}\,\frac{X_1-X_2}{\sqrt{\sum_{i=3}^{8}X_i^2}} \sim t(6).
\]
故
\(
c=\sqrt{3},\qquad Y\sim t(6).
\)

再算概率。因为
\(
\overline{X}\sim N\!\left(0,\frac{4}{8}\right)=N\!\left(0,\frac12\right),
\)
所以
\(
\sqrt{2}\,\overline{X}\sim N(0,1).
\)
\[
P\!\left\{|\overline{X}| \leq \frac{\sqrt{2}}{2}\right\} = P\!\left\{-1 \leq \sqrt{2}\,\overline{X} \leq 1\right\} = 2\Phi(1)-1.
\]
\end{solution}

\begin{exercise}\label{exer:6-xbar-s2-prob}
设总体 \(X \sim N(1, \sigma^2)\)，\(X_1, X_2, \dots, X_6\) 是取自 \(X\) 的简单随机样本，\(\overline{X}, S\) 分别为样本均值与样本标准差。求 \(P\{\overline{X} > 1\}\) 和 \(P\{\overline{X} < 1,\ S^2 < 1.8472\sigma^2\}\)。\ (\(\chi^2_{0.05}(5) = 9.236\))
\end{exercise}
\begin{solution}
\(
\overline{X}\sim N\!\left(1,\frac{\sigma^2}{6}\right).
\)
由于正态分布关于均值 1 对称，所以 \(P\{\overline{X}>1\}=0.5\)。

又因为总体正态时 \(\overline{X}\) 与 \(S^2\) 相互独立，所以
\[
P\{\overline{X} < 1,\ S^2 < 1.8472\sigma^2\} = P\{\overline{X} < 1\}\,P\{S^2 < 1.8472\sigma^2\}.
\]
其中 \(P\{\overline{X}<1\}=0.5\)。

再由
\(
\frac{5S^2}{\sigma^2}\sim\chi^2(5),
\)
以及 \(\chi^2_{0.05}(5)=9.236\) 表示 \(P\{\chi^2(5)>9.236\}=0.05\)，
可得
\(
P\{S^2 < 1.8472\sigma^2\} = P\!\left\{\frac{5S^2}{\sigma^2} < 9.236\right\} = 1 - 0.05 = 0.95.
\)

因此
\[
P\{\overline{X} < 1,\ S^2 < 1.8472\sigma^2\}=0.5\times 0.95=0.475.
\]

故所求两个概率分别为 \(0.5\) 和 \(0.475\)。
\end{solution}


\subsection{两个正态总体}

\begin{theorem}[双正态总体抽样分布] \label{thm:normal-two}
设总体 \(X \sim N(\mu_1, \sigma_1^2)\)，\(Y \sim N(\mu_2, \sigma_2^2)\)，\(X_1, \dots, X_{n_1}\) 和 \(Y_1, \dots, Y_{n_2}\) 分别为来自 \(X\) 和 \(Y\) 的样本，两组样本相互独立。记 \(\overline{X}, S_1^2\) 和 \(\overline{Y}, S_2^2\) 分别为两组的样本均值与样本方差，则
\begin{enumerate}
    \item \(\displaystyle\frac{S_1^2/\sigma_1^2}{S_2^2/\sigma_2^2} \sim F(n_1-1, n_2-1)\)；
    \item 当 \(\sigma_1^2 = \sigma_2^2 = \sigma^2\) 时，
    \[
    \frac{(\overline{X}-\overline{Y})-(\mu_1-\mu_2)}{S_w\sqrt{\frac{1}{n_1}+\frac{1}{n_2}}} \sim t(n_1+n_2-2),
    \]
    其中 \(S_w^2 = \dfrac{(n_1-1)S_1^2+(n_2-1)S_2^2}{n_1+n_2-2}\) 称为\textbf{合并样本方差}。
\end{enumerate}
\end{theorem}

\begin{note}
这一定理能直接使用的前提有两层：一是两个总体都必须正态，二是两组样本必须彼此独立。第 (1) 条主要用于方差比问题，第 (2) 条只在“\(\sigma_1^2=\sigma_2^2\) 相等但未知”时成立。最常见误用是：方差不相等时仍强行套合并方差 \(S_w^2\) 的 \(t\) 统计量；或者两组样本并非独立却仍按双样本 \(t/F\) 分布处理。
\end{note}

\begin{exercise}[样本均值差的分布] \label{ex:sampling_02}
设总体 \(X\) 与 \(Y\) 相互独立且都服从 \(N(\mu,\sigma^2)\)，\(\overline{X},\overline{Y}\) 分别是来自 \(X,Y\)（容量均为 \(n\)）的样本均值，则当 \(n\) 固定时，\(P\{|\overline{X}-\overline{Y}|>\sigma\}\) 随 \(\sigma\) 的增大而（\quad）。\par
A. 单调增大\qquad B. 单调减小\qquad C. 保持不变\qquad D. 增减不定
\end{exercise}
\begin{solution}
应选 C。由独立性，\(\overline{X}\sim N\!\left(\mu,\frac{\sigma^2}{n}\right)\)，\(\overline{Y}\sim N\!\left(\mu,\frac{\sigma^2}{n}\right)\)，故
\(
\overline{X}-\overline{Y}\sim N\!\left(0,\frac{2\sigma^2}{n}\right).
\)
标准化后：
\[
P\{|\overline{X}-\overline{Y}|>\sigma\}
= P\!\left\{\left|\frac{\overline{X}-\overline{Y}}{\sigma\sqrt{2/n}}\right|>\sqrt{\frac{n}{2}}\right\}
= 2\!\left[1-\Phi\!\left(\sqrt{\frac{n}{2}}\right)\right],
\]
该值仅与 \(n\) 有关，与 \(\sigma\) 无关，故选 C。
\end{solution}

\begin{exercise}\label{exer:src-2015-spring-a-i-06}
设总体 \(X\) 和 \(Y\) 相互独立，\(X\sim N(0,4)\)，\(Y\sim N(0,9)\)。记
\[
\overline X=\frac1{10}\sum_{i=1}^{10}X_i,\qquad
\overline Y=\frac1{15}\sum_{i=1}^{15}Y_i,
\]
其中 \(X_1,X_2,\dots,X_{10}\) 与 \(Y_1,Y_2,\dots,Y_{15}\) 分别来自总体 \(X\) 与 \(Y\) 的简单随机样本。求 \(E|\overline X-\overline Y|\)。
\end{exercise}
\begin{solution}
由样本均值的分布，
\[
\overline X\sim N\!\left(0,\frac4{10}\right),\qquad
\overline Y\sim N\!\left(0,\frac9{15}\right).
\]
又两组样本相互独立，所以
\[
\overline X-\overline Y\sim N\!\left(0,\frac4{10}+\frac9{15}\right)=N(0,1).
\]
若 \(Z\sim N(0,1)\)，则
\[
E|Z|=2\int_0^{+\infty}z\frac1{\sqrt{2\pi}}e^{-z^2/2}\,\mathrm dz
=\sqrt{\frac2\pi}.
\]
因此
\(
E|\overline X-\overline Y|=\sqrt{\frac2\pi}.
\)
\end{solution}


\begin{exercise}\label{exer:6-y-f-dist}
设 \(X_1, X_2, \dots, X_{15}\) 是来自正态总体 \(N(0, 9)\) 的简单随机样本，则统计量
\[
Y = \frac{2X_1^2 + X_2^2 + \cdots + X_{10}^2}{X_{11}^2 + X_{12}^2 + \cdots + X_{15}^2}
\]
服从（\quad）。\par
A. \(\chi^2(10)\)\quad B. \(t(15)\)\quad C. \(F(10,5)\)\quad D. \(F(5,10)\)
\end{exercise}

\begin{solution}
应选 C。
\(
\chi_1^2 = \frac{2X_1^2+X_2^2+\cdots+X_{10}^2}{9}
\)
注意到 \(2X_1^2/9 = 2(X_1/3)^2\)，而 \(X_1/3 \sim N(0,1)\)，故 \((X_1/3)^2 \sim \chi^2(1)\)，从而 \(2(X_1/3)^2\) 服从自由度为 2 的 \(\chi^2\) 分布的尺度变换。更直接的方法是：
\[
Y = \frac{(2X_1^2+X_2^2+\cdots+X_{10}^2)/9}{(X_{11}^2+\cdots+X_{15}^2)/9} = \frac{U/10}{V/5},
\]
其中
\[
U = \dfrac{2X_1^2+\cdots+X_{10}^2}{9},\qquad
V = \dfrac{X_{11}^2+\cdots+X_{15}^2}{9}.
\]
分子等价于 10 个独立 \(N(0,1)\) 变量的平方和，所以 \(U \sim \chi^2(10)\)；同理 \(V \sim \chi^2(5)\)。

故 \(Y = \dfrac{U/10}{V/5} \sim F(10, 5)\)。

\begin{note}
\(X_1/3 \sim N(0,1)\)，故 \((X_1/3)^2 \sim \chi^2(1)\)。由伽马分布的可加性，\(2(X_1/3)^2 \sim \chi^2(2)\)（因为 \(\chi^2(1)\) 即 \(\text{Ga}(1/2,1/2)\)，两倍服从 \(\text{Ga}(1,1/2)=\chi^2(2)\)）。因此 \(2X_1^2/9\) 贡献 2 个自由度，分子共 10 个自由度。
\end{note}
\end{solution}

\begin{exercise}\label{exer:src-2015-spring-a-i-04}
设总体 \(X\sim N(0,4)\)，而 \(X_1,X_2,\dots,X_{15}\) 为取自该总体的简单随机样本，则统计量
\[
Y=\frac{X_1^2+X_2^2+\cdots+X_{10}^2}
{2\left(X_{11}^2+X_{12}^2+\cdots+X_{15}^2\right)}
\]
服从什么分布？
\end{exercise}
\begin{solution}
因为 \(X_i\sim N(0,4)\)，所以 \(X_i/2\sim N(0,1)\)。令
\(
U=\sum_{i=1}^{10}\left(\frac{X_i}{2}\right)^2,\qquad
V=\sum_{i=11}^{15}\left(\frac{X_i}{2}\right)^2.
\)
则 \(U\sim\chi^2(10)\)，\(V\sim\chi^2(5)\)，且 \(U,V\) 相互独立。题中统计量可化为
\(
Y=\frac{4U}{2\cdot4V}
=\frac{U}{2V}
=\frac{U/10}{V/5}.
\)
因此
\(
Y\sim F(10,5).
\)
\end{solution}

\begin{exercise}\label{exer:src-2017-a-i-05}
设 \(Y_1,Y_2,Y_3\) 独立且均服从分布 \(\chi^2(n)\)，则
\(
\frac{2Y_1}{Y_2+Y_3}
\)
服从什么分布？
\end{exercise}
\begin{solution}
由 \(\chi^2\) 分布的可加性，
\(
Y_2+Y_3\sim\chi^2(2n).
\)
又 \(Y_1\) 与 \(Y_2+Y_3\) 相互独立。题中统计量可写成
\(
\frac{2Y_1}{Y_2+Y_3}
=\frac{Y_1/n}{(Y_2+Y_3)/(2n)}.
\)
因此
\(
\frac{2Y_1}{Y_2+Y_3}\sim F(n,2n).
\)
\end{solution}

\subsection*{历年试题补充：2008-01 3学分 A 卷}

\begin{exercise}\label{exer:src-2008-01-3a-ii-05}\label{exer:src-2017-a-ii-06}\label{exer:src-paper2-choice-05}\label{prob:p2-choice-05}
设 \((X_1,X_2,\dots,X_n)\) 为总体 \(N(1,2^2)\) 的一个样本，\(\overline X\) 为样本均值，则下列结论中正确的是（\quad）。\par
A. \(\displaystyle \frac{\overline X-1}{2/\sqrt n}\sim t(n)\)\quad
B. \(\displaystyle \frac14\sum_{i=1}^n(X_i-1)^2\sim F(n,1)\)\par
C. \(\displaystyle \frac{\overline X-1}{\sqrt2/\sqrt n}\sim N(0,1)\)\quad
D. \(\displaystyle \frac14\sum_{i=1}^n(X_i-1)^2\sim \chi^2(n)\)
\end{exercise}
\begin{solution}
应选 D。

因为 \(X_i\sim N(1,4)\)，所以
\(
\frac{X_i-1}{2}\sim N(0,1).
\)
于是
\[
\frac14\sum_{i=1}^n(X_i-1)^2
=\sum_{i=1}^n\left(\frac{X_i-1}{2}\right)^2
\sim\chi^2(n),
\]
故 D 正确。

另外，
\(
\overline X\sim N\left(1,\frac4n\right),
\)
所以正确的标准化形式应为
\(
\frac{\overline X-1}{2/\sqrt n}\sim N(0,1),
\)
不是 \(t(n)\)，也不是用 \(\sqrt2/\sqrt n\) 作分母。
\end{solution}

\begin{exercise}\label{exer:src-2009-07-3a-i-05}\label{exer:src-probstat-test1-i-05}
设 \(X_1,X_2,\dots,X_n\) 是来自总体 \(X\sim N(\mu,\sigma^2)\) 的样本，则下列说法中正确的是（\quad）。\par
A. \(\overline X\sim N(\mu,\sigma^2)\)\quad
B. \(\displaystyle\frac1{\sigma^2}\sum_{i=1}^{n}(X_i-\overline X)^2\sim\chi^2(n)\)\par
C. \(\displaystyle\frac{\sqrt n(\overline X-\mu)}{S}\sim t(n)\)\quad
D. \(\displaystyle\frac1{\sigma^2}\sum_{i=1}^{n}(X_i-\mu)^2\sim\chi^2(n)\)
\end{exercise}
\begin{solution}
应选 D。

因为
\(
\frac{X_i-\mu}{\sigma}\sim N(0,1),\qquad i=1,2,\dots,n,
\)
且样本相互独立，所以
\[
\frac1{\sigma^2}\sum_{i=1}^{n}(X_i-\mu)^2
=\sum_{i=1}^{n}\left(\frac{X_i-\mu}{\sigma}\right)^2
\sim\chi^2(n).
\]
故 D 正确。

A 错在样本均值的方差应为 \(\sigma^2/n\)，即
\(
\overline X\sim N\left(\mu,\frac{\sigma^2}{n}\right).
\)
B 错在用样本均值 \(\overline X\) 替代真实均值会损失一个自由度，应为
\[
\frac1{\sigma^2}\sum_{i=1}^{n}(X_i-\overline X)^2\sim\chi^2(n-1).
\]
C 中若 \(S\) 为样本标准差，则
\(
\frac{\sqrt n(\overline X-\mu)}{S}\sim t(n-1),
\)
自由度也不是 \(n\)。
\end{solution}

\begin{exercise}\label{exer:src-2008-01-3a-iii-05}\label{exer:src-paper2-fill-05}\label{prob:p2-fill-05}
设 \(X\sim t(m)\)，则随机变量
\(
Y=X^2
\)
服从什么分布？请写出自由度。
\end{exercise}
\begin{solution}
由 \(t\) 分布定义，若 \(X\sim t(m)\)，则可表示为
\(
X=\frac{U}{\sqrt{V/m}},
\)
其中 \(U\sim N(0,1)\)，\(V\sim\chi^2(m)\)，且 \(U,V\) 相互独立。于是
\(
X^2=\frac{U^2/1}{V/m}.
\)
因为 \(U^2\sim\chi^2(1)\)，所以
\(
Y=X^2\sim F(1,m).
\)
\end{solution}


\begin{exercise}\label{exer:src-paper3-ii-01}\label{prob:p3-ii-01}
设随机变量 \(X\sim N(0,1)\)。对给定的 \(\alpha\ (0<\alpha<1)\)，数 \(z_\alpha\) 满足
\(
P(X>z_\alpha)=\alpha.
\)
若
\(
P(|X|<c)=\alpha,
\)
求 \(c\)。
\end{exercise}
\begin{solution}
由标准正态分布的对称性，
\(
P(|X|<c)=\alpha
\)
等价于两侧尾概率之和为 \(1-\alpha\)，且左右两侧各为
\(
\frac{1-\alpha}{2}.
\)
因此
\(
P(X>c)=\frac{1-\alpha}{2}.
\)
按上侧分位数定义，
\(
c=z_{\frac{1-\alpha}{2}}.
\)
\end{solution}

\begin{exercise}\label{exer:src-paper3-ii-04}\label{prob:p3-ii-04}
设 \(X_1,X_2,\dots,X_n\) 是来自正态总体 \(N(\mu,1)\) 的一个简单随机样本，\(\overline X,S^2\) 分别为样本均值与样本方差。判断下列结论中正确的是（\quad）。\par
A. \(\overline X\sim N(0,1)\)\quad
B. \(\displaystyle\sum_{i=1}^n (X_i-\overline X)^2\sim\chi^2(n-1)\)\par
C. \(\displaystyle\sum_{i=1}^n (X_i-\mu)^2\sim\chi^2(n-1)\)\quad
D. \(\displaystyle\frac{\overline X}{S/\sqrt{n-1}}\sim t(n-1)\)。
\end{exercise}
\begin{solution}
应选 B。

因为总体方差为 \(1\)，所以
\[
\sum_{i=1}^n (X_i-\overline X)^2
=\frac{1}{1}\sum_{i=1}^n (X_i-\overline X)^2
\sim\chi^2(n-1).
\]

A 错在 \(\overline X\sim N(\mu,1/n)\)，不一定是 \(N(0,1)\)。C 中若用真均值 \(\mu\)，则
\(
\sum_{i=1}^n (X_i-\mu)^2\sim\chi^2(n),
\)
自由度是 \(n\) 而不是 \(n-1\)。D 的 \(t\) 统计量应为
\(
\frac{\overline X-\mu}{S/\sqrt n}\sim t(n-1),
\)
题中既没有减去 \(\mu\)，分母也写错。
\end{solution}

\begin{exercise}\label{exer:src-paper3-iii-06}\label{prob:p3-iii-06}
设 \(X_1,X_2,X_3,X_4\) 是来自正态总体 \(N(0,9)\) 的一个简单随机样本，令
\(
\xi=\frac{(X_2+X_3+X_4)^2}{3X_1^2}.
\)
求 \(\xi\) 的分布，并写出自由度。
\end{exercise}
\begin{solution}
因为 \(X_i\sim N(0,9)\)，且相互独立，所以
\(
X_2+X_3+X_4\sim N(0,27).
\)
于是
\(
\frac{X_2+X_3+X_4}{3\sqrt3}\sim N(0,1),
\)
从而
\(
\frac{(X_2+X_3+X_4)^2}{27}\sim\chi^2(1).
\)
同理
\(
\frac{X_1^2}{9}\sim\chi^2(1).
\)
且两者相互独立。因此
\(
\xi
=\frac{(X_2+X_3+X_4)^2}{3X_1^2}
=\frac{\dfrac{(X_2+X_3+X_4)^2}{27}}{\dfrac{X_1^2}{9}}
\sim F(1,1).
\)
\end{solution}

\section{分布的分位数}

\begin{note}
\textbf{自学导读：分位数是查表的"接口"。}

后面区间估计和假设检验中，你会反复看到 $z_{\alpha/2}$、$t_{\alpha/2}(n-1)$、$\chi^2_{\alpha/2}(n-1)$ 这类符号。它们的含义统一且简单：\textbf{上 $\alpha$ 分位数 $=$ 右尾面积恰好为 $\alpha$ 的那个临界点}。

直觉：想象密度曲线下的面积为 1。从右边切掉面积 $\alpha$ 的那个"刀口"位置，就是上 $\alpha$ 分位数。

\textbf{易错点：} 不同教材可能用"上侧"或"下侧"定义。本书统一用\textbf{上侧}：$P(X>a_\alpha)=\alpha$。做题时先确认教材/试卷用哪种定义，再查表。
\end{note}

\begin{definition}[上 \texorpdfstring{\(\alpha\)}{α} 分位数] \label{def:quantile}
设 \(X\) 为随机变量，\(0 < \alpha < 1\)。若实数 \(a_\alpha\) 满足 \(P\{X > a_\alpha\} = \alpha\)（即上侧尾部面积为 \(\alpha\)），则称 \(a_\alpha\) 为 \(X\) 的\textbf{上 \(\alpha\) 分位数}。
\end{definition}

三大抽样分布的上 \(\alpha\) 分位数分别为：
\begin{enumerate}
    \item \(\chi^2\) 分布：\(P\{\chi^2 > \chi_\alpha^2(n)\} = \alpha\)；
    \item \(t\) 分布：\(P\{t > t_\alpha(n)\} = \alpha\)，由对称性 \(t_{1-\alpha}(n) = -t_\alpha(n)\)；
    \item \(F\) 分布：\(P\{F > F_\alpha(m,n)\} = \alpha\)，且 \(F_{1-\alpha}(m,n) = \dfrac{1}{F_\alpha(n,m)}\)。
\end{enumerate}

\begin{note}
不同教材对分位数的定义可能不同：有的用上侧（\(P\{X > a\} = \alpha\)），有的用下侧（\(P\{X \leq a\} = \alpha\)）。使用时需注意对应关系：
\begin{itemize}
    \item 上侧 \(\alpha\) 分位数 \(a_\alpha\) 对应下侧 \(1-\alpha\) 分位数；
    \item 两种定义下，\(\chi^2\) 分布满足 \(\chi^2_{\alpha,\text{上侧}}(n) = \chi^2_{1-\alpha,\text{下侧}}(n)\)。
\end{itemize}
\end{note}

\begin{note}
\textbf{使用提醒：} 查分位数表前，建议固定做三步检查：先认分布族（标准正态、\(t\)、\(\chi^2\) 还是 \(F\)），再认自由度，最后认清教材用的是上侧定义还是下侧定义。尤其 \(F\) 分布除了尾部方向外，还要额外检查第一、第二自由度的顺序，次序一换，分位数就会跟着变。
\end{note}

\subsection*{历年试题补充：2009.07 3学分 A 卷}

\begin{exercise}\label{exer:src-2009-07-3a-i-03}
设连续随机变量 \(X\) 的密度函数满足 \(f(x)=f(-x)\)，\(x_\alpha\) 是 \(X\) 的上 \(\alpha\) 分位数，则
\(
P(|X|>x_\alpha)=(\quad).
\)
A. \(2-\alpha\)\quad B. \(1-2\alpha\)\par
C. \(2\alpha-1\)\quad D. \(2\alpha\)
\end{exercise}
\begin{solution}
应选 D。

由 \(f(x)=f(-x)\) 可知分布关于 \(0\) 对称。又 \(x_\alpha\) 是上 \(\alpha\) 分位数，所以
\(
P(X>x_\alpha)=\alpha.
\)
由对称性，
\(
P(X<-x_\alpha)=\alpha.
\)
因此
\(
P(|X|>x_\alpha)=P(X>x_\alpha)+P(X<-x_\alpha)=2\alpha.
\)
\end{solution}

\begin{exercise}\label{exer:src-probstat-test1-i-03}
设连续随机变量 \(X\) 的密度函数满足 \(f(x)=f(-x)\)，\(x_\alpha\) 是 \(X\) 的 \(\alpha\) 分位数，且 \(x_\alpha>0\)，则
\(
P(|X|\le x_\alpha)=(\quad).
\)
A. \(2-\alpha\)\quad B. \(1-2\alpha\)\quad C. \(2\alpha-1\)\quad D. \(2\alpha\)
\end{exercise}
\begin{solution}
应选 C。

这里的 \(\alpha\) 分位数按下侧概率理解，即
\(
P(X\le x_\alpha)=\alpha.
\)
又由 \(f(x)=f(-x)\) 知 \(X\) 的分布关于 \(0\) 对称，因此
\(
P(X<-x_\alpha)=P(X>x_\alpha)=1-\alpha.
\)
于是
\[
P(|X|\le x_\alpha)
=P(-x_\alpha\le X\le x_\alpha)
=1-2(1-\alpha)
=2\alpha-1.
\]
\end{solution}

\begin{exercise}\label{exer:6-quantile-s2}
设总体 \(X \sim N(\mu, 16)\)，\(X_1, X_2, \dots, X_{10}\) 是取自 \(X\) 的样本，\(S^2\) 为样本方差。求常数 \(a =\,\underline{\hspace{2cm}}\)，使 \(P(S^2 \geq a) = 0.1\) 成立。\par
A. \(\dfrac{16}{9}\chi_{0.1}^2(9)\)\quad B. \(\dfrac{16}{9}\chi_{0.05}^2(10)\)\quad C. \(\chi_{0.1}^2(10)\)\quad D. \(\chi_{0.05}^2(9)\)
\end{exercise}

\begin{solution}
应选 A。

本题先识别样本方差的抽样分布。由于总体 \(X \sim N(\mu,16)\)，样本容量 \(n=10\)，所以
\(
\frac{(n-1)S^2}{\sigma^2}=\frac{9S^2}{16}\sim\chi^2(9).
\)
题目要求
\(
P(S^2 \geq a)=0.1,
\)
等价于
\(
P\!\left(\frac{9}{16}S^2 \geq \frac{9}{16}a\right)=0.1.
\)
而上 \(0.1\) 分位数 \(\chi_{0.1}^2(9)\) 的定义正是右尾概率等于 \(0.1\) 的点，因此必须有
\(
\frac{9}{16}a = \chi_{0.1}^2(9)\implies a = \frac{16}{9}\chi_{0.1}^2(9).
\)
故应选 A。
\end{solution}

\begin{exercise}\label{exer:6-quantile-range}
设随机变量 \(X \sim \chi^2(2)\)，用 \(\chi_\alpha^2(2)\) 表示自由度为 2 的 \(\chi^2\) 分布的上 \(\alpha\) 分位数。已知 \(P(x < X < y) = 0.95\)，\(P(X > y) = 0.02\)，请用分位数表示 \(x\) 和 \(y\)。
\end{exercise}

\begin{solution}
先看右端点 \(y\)。题目已给 \(P(X>y)=0.02\)。
而 \(\chi^2_\alpha(2)\) 的定义正是“使右尾概率等于 \(\alpha\) 的点”，所以 \(y=\chi^2_{0.02}(2)\)。

再求左端点 \(x\)。由 \(P(x<X<y)=0.95\) 以及右尾概率 \(P(X>y)=0.02\)，可知剩下的左侧概率为
\(
P(X\le x)=1-0.95-0.02=0.03.
\)
这一步最容易算错成 \(0.05\) 或 \(0.02\)，本质上是在把总概率 1 分成三段：左边、中间、右边。

于是 \(P(X>x)=1-0.03=0.97\)，按上 \(\alpha\) 分位数的定义，满足 \(P(X>x)=0.97\) 的点就是 \(x=\chi^2_{0.97}(2)\)。

所以
\(
x=\chi^2_{0.97}(2),\qquad y=\chi^2_{0.02}(2).
\)
\end{solution}


\section{综合练习}

\begin{exercise}\label{exer:6-multi-normal}
设 \(X \sim N(0,1)\)，\(X_1, X_2, \dots, X_n\) 是取自 \(X\) 的样本，\(\overline{X}\) 为样本均值，\(S^2\) 为样本方差，则（\quad）。\par
A. \(\overline{X} \sim N(0,1)\)\qquad B. \(n\overline{X} \sim N(0,1)\)\par
C. \(\displaystyle\sum_{i=1}^{n}X_i^2 \sim \chi^2(n)\)\qquad D. \(\dfrac{\overline{X}}{S} \sim t(n-1)\)
\end{exercise}

\begin{solution}
应选 C。

\textbf{这类综合判断题的做法是逐项对照单正态总体抽样分布的三个标准结论。}

A: 由 \(\overline{X}\sim N(0,1/n)\)，可知它不是 \(N(0,1)\)，故 A 错。

B: 因为 \(n\overline{X}=\sum_{i=1}^n X_i\)，所以
\(
n\overline{X}\sim N(0,n),
\)
也不是 \(N(0,1)\)，故 B 错。

C: 由于 \(X_1,\dots,X_n\) 相互独立且都服从 \(N(0,1)\)，所以
\(
\sum_{i=1}^n X_i^2\sim \chi^2(n),
\)
这正是 \(\chi^2\) 分布的定义，故 C 对。

D: \(t\) 分布的标准形式应为
\(
\frac{\overline{X}-\mu}{S/\sqrt{n}}\sim t(n-1).
\)
题中只写 \(\overline{X}/S\)，既少了中心化的 \(\mu\)，也少了 \(\sqrt{n}\) 的缩放，因此 D 错。
\end{solution}

\begin{exercise}\label{exer:6-multi-harbin}
设 \(X_1, X_2, \dots, X_n\) 是来自总体 \(X \sim N(0,1)\) 的简单随机样本，则下列统计量的分布不正确的是（\quad）。\par
A. \(\displaystyle\sum_{i=1}^{n}X_i^2 \sim \chi^2(n)\)\qquad B. \(\displaystyle\frac{\sqrt{n-1}\,X_n}{\sqrt{\sum_{i=1}^{n-1}X_i^2}} \sim t(n-1)\)\par
C. \(\displaystyle\frac{1}{n}\sum_{i=1}^{n}X_i \sim N(0,1)\)\qquad D. \(\displaystyle\frac{(n-1)\sum_{i=1}^{2}X_i^2}{2\sum_{i=3}^{n}X_i^2} \sim F(2,n-2)\)
\end{exercise}

\begin{solution}
应选 C。

\textbf{仍按“定义或标准形式”逐项核对。}

A: 因为 \(X_1,\dots,X_n\) 独立同分布于 \(N(0,1)\)，所以
\(
\sum_{i=1}^{n}X_i^2 \sim \chi^2(n),
\)
故 A 正确。

B: \(X_n \sim N(0,1)\)，且
\(
\sum_{i=1}^{n-1}X_i^2 \sim \chi^2(n-1).
\)
又因 \(X_n\) 与 \(X_1,\dots,X_{n-1}\) 独立，所以
\(
\frac{X_n}{\sqrt{\sum_{i=1}^{n-1}X_i^2/(n-1)}} \sim t(n-1),
\)
即 B 正确。

C: \(\dfrac{1}{n}\sum X_i=\overline{X}\sim N(0,1/n)\)，不是 \(N(0,1)\)，所以 C 错。

D: 前两项平方和与后 \(n-2\) 项平方和分别服从
\[
\sum_{i=1}^{2}X_i^2\sim \chi^2(2),\qquad \sum_{i=3}^{n}X_i^2\sim \chi^2(n-2),
\]
且二者独立，因此
\(
\frac{\sum_{i=1}^{2}X_i^2/2}{\sum_{i=3}^{n}X_i^2/(n-2)} \sim F(2,n-2),
\)
故 D 正确。

因此分布不正确的是 C。
\end{solution}


\begin{exercise}\label{exer:6-multi-dist}
设 \(X_1, X_2, \dots, X_n\) 是总体 \(N(0,1)\) 的简单随机样本，则 \(a =\,\underline{\hspace{2cm}}\) 时 \(a(X_1-X_2)^2 \sim\,\underline{\hspace{2cm}}\)；\(\dfrac{2(X_1-X_2)}{\sqrt{\sum_{i=3}^{n}X_i^2}} \sim\,\underline{\hspace{2cm}}\)。
\end{exercise}
\begin{solution}
第一空填 \(\dfrac{1}{2}\)，第二空填 \(\chi^2(1)\)，第三空填 \(t(8)\)。

先处理前两空。因为 \(X_1,X_2\) 相互独立且都服从 \(N(0,1)\)，所以
\(
X_1-X_2\sim N(0,1+1)=N(0,2).
\)
把它标准化：
\(
\frac{X_1-X_2}{\sqrt2}\sim N(0,1).
\)
对标准正态平方，就得到自由度为 1 的卡方分布：
\(
\left(\frac{X_1-X_2}{\sqrt2}\right)^2=\frac{(X_1-X_2)^2}{2}\sim\chi^2(1).
\)
因此
\(
a=\frac12,\qquad a(X_1-X_2)^2\sim\chi^2(1).
\)

再看第三空。分子已经知道
\(
\frac{X_1-X_2}{\sqrt2}\sim N(0,1).
\)
分母中的
\(
\sum_{i=3}^{n}X_i^2
\)
是 \(n-2\) 个独立标准正态平方和，所以
\(
\sum_{i=3}^{n}X_i^2\sim\chi^2(n-2).
\)
而且它与分子由不同样本项构成，因此相互独立。于是按 \(t\) 分布定义，
\[
\frac{(X_1-X_2)/\sqrt2}{\sqrt{\sum_{i=3}^{n}X_i^2/(n-2)}}\sim t(n-2).
\]
把左式整理成题目给定形式：
\[
\frac{(X_1-X_2)/\sqrt2}{\sqrt{\sum_{i=3}^{n}X_i^2/(n-2)}}
=\sqrt{\frac{n-2}{2}}\cdot\frac{X_1-X_2}{\sqrt{\sum_{i=3}^{n}X_i^2}}.
\]
题目系数是 2，所以令
\(
\sqrt{\frac{n-2}{2}}=2 \implies n-2=8,\qquad n=10.
\)
因此该统计量服从 \(t(n-2)=t(8)\)。

综上，答案依次为 \(\dfrac12\)、\(\chi^2(1)\)、\(t(8)\)。
\end{solution}


\begin{exercise}\label{exer:6-ET-DT}
设总体 \(X \sim N(\mu, \sigma^2)\)，\(X_1, X_2, \dots, X_n\) 是取自 \(X\) 的简单随机样本，\(\overline{X}, S^2\) 分别为样本均值和样本方差。记 \(T = \overline{X}^2 - \dfrac{1}{n}S^2\)。\par
(1) 证明 \(ET = \mu^2\)；\par
(2) 当 \(\mu = 0\)，\(\sigma = 1\) 时，求 \(DT\)。
\end{exercise}
\begin{solution}
(1) \(E(\overline{X}) = \mu\)，\(D(\overline{X}) = \sigma^2/n\)，\(E(S^2) = \sigma^2\)。
\[
E(\overline{X}^2) = D(\overline{X}) + [E(\overline{X})]^2 = \frac{\sigma^2}{n} + \mu^2.
\]
\[
ET = E\!\left(\overline{X}^2 - \frac{1}{n}S^2\right) = \left(\frac{\sigma^2}{n} + \mu^2\right) - \frac{\sigma^2}{n} = \mu^2. 
\]
(2) 当 \(\mu=0\)，\(\sigma=1\) 时，\(X \sim N(0,1)\)。\(\sqrt{n}\,\overline{X} \sim N(0,1)\)，故 \(n\overline{X}^2 \sim \chi^2(1)\)。\((n-1)S^2 \sim \chi^2(n-1)\)。由 \(\overline{X}\) 与 \(S^2\) 独立：
\[
D(\overline{X}^2) = D\!\left(\frac{\chi^2(1)}{n}\right) = \frac{2}{n^2},\quad D\!\left(\frac{S^2}{n}\right) = \frac{1}{n^2}\cdot\frac{2(n-1)}{(n-1)^2} = \frac{2}{n^2(n-1)}.
\]
\(
DT = D(\overline{X}^2) + D\!\left(\frac{S^2}{n}\right) = \frac{2}{n^2} + \frac{2}{n^2(n-1)} = \frac{2n}{n^2(n-1)} = \frac{2}{n(n-1)}.
\)
\end{solution}


\chapter{点估计}


{\small
\[
\text{点估计}
\begin{cases}
\text{具体内容: 总体分布未知, 寻求适当的统计量估计分布中的未知参数} \\
\text{估计方法}
\begin{cases}
\text{矩估计}
\begin{cases}
\text{思想: 样本矩替换总体矩} \\
\text{掌握: 一阶矩、二阶矩}
\end{cases} \\
\text{最大似然估计}
\begin{cases}
\text{思想: 样本似然函数取最大值时对应的未知参数} \\
\text{掌握: 最大似然估计步骤}
\end{cases}
\end{cases} \\
\text{估计量的评选标准}
\begin{cases}
\text{无偏性: } E(\hat{\theta}) = \theta \\
\text{有效性: 若 } \hat{\theta}_1,\hat{\theta}_2 \text{ 为 } \theta \text{ 的无偏估计量,} \\
\qquad D(\hat{\theta}_1) \leq D(\hat{\theta}_2), \text{ 则 } \hat{\theta}_1 \text{ 比 } \hat{\theta}_2 \text{ 有效} \\
\text{一致性(相合性): } \lim\limits_{n \to \infty} P\left\{|\hat{\theta} - \theta| < \varepsilon\right\} = 1, \\
\qquad \text{也可记为 } \hat{\theta} \xrightarrow{P} \theta \\
\text{均方误差: } \hat{\theta} \text{ 为 } \theta \text{ 的估计量,} \\
\qquad \text{称 } M(\hat{\theta}) = E(\hat{\theta} - \theta)^2 \text{ 为 } \hat{\theta} \text{ 关于 } \theta \text{ 的均方误差}, \\
\quad \text{若 } \hat{\theta}_1,\hat{\theta}_2 \text{ 为 } \theta \text{ 的无偏估计量, 且 } M(\hat{\theta}_1) \leq M(\hat{\theta}_2), \\
\quad \text{则 } \hat{\theta}_1 \text{ 在均方误差下比 } \hat{\theta}_2 \text{ 有效}
\end{cases}
\end{cases}
\]
}


\section{点估计的基本概念}

\begin{definition}[点估计] \label{def:point-est}
设总体 \(X\) 的分布函数形式已知，但含有一个或多个未知参数 \(\theta\)。利用总体 \(X\) 的一个样本 \(X_1, X_2, \dots, X_n\)
构造一个适当的统计量 \(\hat{\theta}(X_1, X_2, \dots, X_n)\) 
来估计未知参数 \(\theta\)，则称其为 \(\theta\) 的\textbf{估计量}。将样本观测值 \(x_1, x_2, \dots, x_n\)代入估计量，得到的数值 \(\hat{\theta}(x_1, x_2, \dots, x_n)\) 称为 \(\theta\) 的\textbf{估计值}。
\end{definition}

\begin{note}
估计量是随机变量（样本的函数），估计值是具体的数值。常用的点估计方法有\textbf{矩估计法}和\textbf{最大似然估计法}。
\end{note}

\begin{note}
\textbf{自学追问：} 为什么要把“估计量”和“估计值”分开？因为前者是随机变量，后者是具体数值；不分开就没法谈无偏性、方差和一致性。
\end{note}

\section{矩估计法（MME）} \label{subsec:mme}

\begin{definition}[矩] \label{def:moments}
设 \(X, Y\) 为随机变量，\(k, l\) 为正整数。
\begin{enumerate}
    \item 称 \(E(X^k)\) 为 \(X\) 的\textbf{\(k\) 阶原点矩}。
    \item 称 \(E[(X - EX)^k]\) 为 \(X\) 的\textbf{\(k\) 阶中心矩}。
    \item 称 \(E(X^k Y^l)\) 为 \(X\) 与 \(Y\) 的\textbf{\(k+l\) 阶混合原点矩}。
    \item 称 \(E[(X - EX)^k (Y - EY)^l]\) 为 \(X\) 与 \(Y\) 的\textbf{\(k+l\) 阶混合中心矩}。
\end{enumerate}
\end{definition}

\begin{note}
从矩的定义可以看出：数学期望 \(EX\) 是 \(X\) 的一阶原点矩，方差 \(DX\) 是 \(X\) 的二阶中心矩，协方差 \(\mathrm{Cov}(X,Y)\) 是 \(X\) 与 \(Y\) 的二阶混合中心矩。
\end{note}

矩估计法的核心思想是\textbf{用样本矩替换总体矩}：由大数定律，样本矩依概率收敛于总体矩，因此可以用样本矩来估计总体矩，从而解出未知参数。

\begin{property}[矩估计的步骤] \label{prop:mme-steps}
设总体 \(X\) 含有 \(k\) 个未知参数 \(\theta_1, \theta_2, \dots, \theta_k\)，\(X_1, X_2, \dots, X_n\) 为样本。
\begin{enumerate}
    \item 计算总体的前 \(k\) 阶矩 \(\mu_j = E(X^j)\ (j=1,2,\dots,k)\)，它们是 \(\theta_1, \dots, \theta_k\) 的函数。
    \item 计算对应的样本矩 \(A_j = \dfrac{1}{n}\displaystyle\sum_{i=1}^{n}X_i^j\ (j=1,2,\dots,k)\)。
    \item 令 \(\mu_j = A_j\ (j=1,2,\dots,k)\)，建立 \(k\) 个方程。
    \item 解方程组，得到 \(\hat{\theta}_1, \hat{\theta}_2, \dots, \hat{\theta}_k\) 作为矩估计量。
\end{enumerate}
\end{property}

\begin{note}
求矩估计时，以低阶为原则——取能解出参数的最小 \(j\) 值。例如，若一个参数可用一阶矩解出，就不必使用二阶矩。对于两个参数，通常用 \(EX\) 和 \(EX^2\)（或等价地 \(EX\) 和 \(DX\)）建立两个方程。
\end{note}

\begin{exercise}[矩估计量的方差] \label{ex:mom_03}
设总体 \(X\) 的概率密度为
\[
f(x)=
\begin{cases}
\dfrac{6x}{\theta^3}(\theta-x), & 0<x<\theta,\\
0, & \text{其他},
\end{cases}
\]
\(X_1,X_2,\dots,X_n\) 是简单随机样本。求：(1) \(\theta\) 的矩估计量；(2) 该矩估计量的方差。
\end{exercise}
\begin{solution}
(1) 计算总体期望：
\[
EX=\int_0^\theta\frac{6x^2}{\theta^3}(\theta-x)\,\mathrm{d}x=\frac{6}{\theta^3}\!\left[\theta\cdot\frac{\theta^3}{3}-\frac{\theta^4}{4}\right]=\frac{6}{\theta^3}\cdot\frac{\theta^4}{12}=\frac{\theta}{2}.
\]
令 \(\theta/2=\overline{X}\)，得 \(\hat{\theta}=2\overline{X}\)。\\
(2) 先求 \(DX\)：
\[
EX^2=\int_0^\theta\frac{6x^3}{\theta^3}(\theta-x)\,\mathrm{d}x=\frac{6}{\theta^3}\!\left[\theta\cdot\frac{\theta^4}{4}-\frac{\theta^5}{5}\right]=\frac{6}{\theta^3}\cdot\frac{\theta^5}{20}=\frac{3\theta^2}{10}.
\]
\[
DX=EX^2-(EX)^2=\frac{3\theta^2}{10}-\frac{\theta^2}{4}=\frac{6\theta^2-5\theta^2}{20}=\frac{\theta^2}{20}.
\]
因此
\[
D\hat{\theta}=D(2\overline{X})=4D\overline{X}=\frac{4}{n}DX=\frac{4}{n}\cdot\frac{\theta^2}{20}=\frac{\theta^2}{5n}.
\]
\end{solution}



\section{最大似然估计法（MLE）} \label{subsec:mle}

最大似然估计的基本思想：在参数可能的取值范围内，选取使\textbf{样本似然函数}最大的参数值作为估计值。对离散型分布，可以理解为“让当前观测结果出现的概率尽可能大”；对连续型分布，则应理解为“让观测值对应的密度值尽可能大”，不能机械地把它说成概率本身。

\begin{property}[最大似然估计的步骤] \label{prop:mle-steps}
设总体 \(X\) 的分布（离散型为 \(p(x;\theta)\)，连续型为 \(f(x;\theta)\)）含有未知参数 \(\theta\)，样本观测值为 \(x_1, x_2, \dots, x_n\)。
\begin{enumerate}
    \item 写出\textbf{似然函数}：
    \[
    L(\theta) = \prod_{i=1}^{n}p(x_i;\theta)\quad\text{（离散型）}\qquad\text{或}\qquad L(\theta) = \prod_{i=1}^{n}f(x_i;\theta)\quad\text{（连续型）}.
    \]
    \item 取对数：\(\ln L(\theta)\)（将乘积化为求和，便于求导）。
    \item 对 \(\theta\) 求导并令导数为零，解出 \(\hat{\theta}\)：
    \begin{itemize}
        \item \textbf{有驻点}：令 \(\dfrac{\mathrm{d}\ln L}{\mathrm{d}\theta} = 0\)，解出 \(\hat{\theta}\)；
        \item \textbf{单调型（无驻点）}：\(\hat{\theta}\) 在参数范围的\textbf{边界}取得（通常是 \(\max\{X_i\}\) 或 \(\min\{X_i\}\)）；
        \item \textbf{常数型}：\(\hat{\theta}\) 不唯一。
    \end{itemize}
\end{enumerate}
\end{property}

\begin{note}
\textbf{易错点：} 最大似然估计不一定都来自“令导数等于 0”。若似然函数在参数范围内单调，极大值往往直接落在边界点（如样本最大值或最小值）；若参数还影响支持区间，甚至应先写出参数约束，再谈求导。做 MLE 题时，\textbf{先看参数范围和似然是否单调，再看有没有驻点}，通常更不容易漏解。
\end{note}



\begin{property}[最大似然估计的不变性] \label{prop:mle-inv}
设 \(\hat{\theta}\) 是 \(\theta\) 的最大似然估计，\(u = u(\theta)\) 具有反函数，则 \(\hat{u} = u(\hat{\theta})\) 是 \(u(\theta)\) 的最大似然估计。例如，若 \(\hat{\theta}\) 是 \(\theta\) 的 MLE，则 \(e^{\hat{\theta}}\) 是 \(e^\theta\) 的 MLE。
\end{property}


\begin{exercise}[最大似然估计的不变性] \label{ex:mle_04}
设总体 \(X\) 的概率密度为
\[
f(x)=
\begin{cases}
\dfrac{2x}{\alpha^2}, & 0\leq x\leq\alpha,\\
0, & \text{其他},
\end{cases}
\]
其中 \(\alpha>1\) 未知。\(X_1,X_2,\dots,X_n\) 为简单随机样本。\par
(1) 求 \(p=P\{0<X<\sqrt{\alpha}\}\)；\par
(2) 求 \(p\) 的最大似然估计量 \(\hat{p}\)。
\end{exercise}
\begin{solution}
(1)
\[
p=P\{0<X<\sqrt{\alpha}\}=\int_0^{\sqrt{\alpha}}\frac{2x}{\alpha^2}\,\mathrm{d}x=\frac{1}{\alpha^2}\cdot(\sqrt{\alpha})^2=\frac{1}{\alpha}.
\]
(2) 当 \(0\leq x_i\leq\alpha\ (i=1,\dots,n)\) 时，
\[
L(\alpha)=\prod_{i=1}^{n}\frac{2x_i}{\alpha^2}=\frac{2^n\prod_{i=1}^{n}x_i}{\alpha^{2n}}.
\]
\(L(\alpha)\) 关于 \(\alpha\) 单调递减，约束为 \(\alpha\geq\max\{X_1,\dots,X_n\}\)，故
\(
\hat{\alpha}=\max\{X_1,X_2,\dots,X_n\}.
\)
由 MLE 的\textbf{不变性}，\(p=g(\alpha)=1/\alpha\) 的 MLE 为
\[
\hat{p}=\frac{1}{\hat{\alpha}}=\frac{1}{\max\{X_1,X_2,\dots,X_n\}}.
\]
\end{solution}

\begin{note}
\textbf{MLE 单调型的处理步骤详解：} 上题是典型的"似然函数单调"情形。遇到这类题时，按以下步骤操作：
\begin{enumerate}
    \item 写出似然函数 $L(\theta)$，注意\textbf{参数的约束范围}（通常来自密度函数的支撑集，如 $0\leqslant x\leqslant\theta$ 意味着 $\theta\geqslant\max\{x_i\}$）；
    \item 判断 $L(\theta)$（或 $\ln L(\theta)$）关于 $\theta$ 是单调递增还是递减；
    \item 若单调递减：$\theta$ 越小 $L$ 越大，但 $\theta$ 不能小于约束下界 $\to$ $\hat\theta = \max\{X_i\}$；
    \item 若单调递增：$\theta$ 越大 $L$ 越大，但 $\theta$ 不能超过约束上界 $\to$ $\hat\theta = \min\{X_i\}$。
\end{enumerate}
\textbf{口诀：} 递减取 max，递增取 min。
\end{note}

\begin{note}
\textbf{图解"支撑含参数"型 MLE（以 $U(0,\theta)$ 为例）：}

设 $X_1,\dots,X_n\overset{\text{iid}}{\sim}U(0,\theta)$，似然函数为
\[
L(\theta)=\frac{1}{\theta^n}\cdot\mathbf{1}\{\theta\geq X_{(n)}\}.
\]
\textbf{关键理解：} 当 $\theta < X_{(n)}$ 时，$L(\theta)=0$（不可能产生比 $\theta$ 还大的观测值）。当 $\theta\geq X_{(n)}$ 时，$L(\theta)=1/\theta^n$ 关于 $\theta$ 严格递减。

想象 $L(\theta)$ 的图像：$\theta$ 轴上，左半段（$\theta<X_{(n)}$）为零，右半段（$\theta\geq X_{(n)}$）是下降曲线。最大值恰在"断崖边缘" $\theta=X_{(n)}$ 处取得。

\textbf{做题模板：} 凡是密度中出现 $0\leq x\leq\theta$ 或 $\theta\leq x$ 这类条件的：
\begin{enumerate}
    \item 写出似然函数时，\textbf{必须带上示性函数}（约束条件）；
    \item 在约束允许的范围内判断单调性；
    \item 最大值在约束边界取得 $\Rightarrow$ $\hat\theta=X_{(n)}$ 或 $X_{(1)}$。
\end{enumerate}
\end{note}

\begin{exercise}\label{exer:src-2017-a-07}
(1) 设总体 \(X\) 等可能地取值 \(1,2,\dots,N\)，其中 \(N\) 是未知正整数。\(X_1,\dots,X_n\) 是来自该总体的简单随机样本，求 \(N\) 的最大似然估计量。\par
(2) 某单位自行车棚内存放了 \(N\) 辆自行车，编号为 \(1,2,\dots,N\)。假定取走自行车是等可能的。某人连续 \(12\) 天记录观察到的取走的第一辆自行车编号为
\[
12,\ 203,\ 23,\ 7,\ 239,\ 45,\ 73,\ 189,\ 95,\ 112,\ 73,\ 159.
\]
求在上述样本观测值下 \(N\) 的最大似然估计值。
\end{exercise}
\begin{solution}
(1) 对样本观测值 \(x_1,\dots,x_n\)，似然函数为
\(
L(N)=\prod_{i=1}^n P\{X=x_i\}
=\left(\frac1N\right)^n,
\)
但必须满足
\(
N\ge \max\{x_1,\dots,x_n\}=x_{(n)}.
\)
在约束 \(N\ge x_{(n)}\) 下，\(L(N)=N^{-n}\) 随 \(N\) 增大而减小，因此最大值在最小可取的 \(N\) 处达到。故
\(
\widehat N=X_{(n)}=\max\{X_1,\dots,X_n\}.
\)

(2) 观测值中的最大值为
\(
\max\{12,203,23,7,239,45,73,189,95,112,73,159\}=239.
\)
因此最大似然估计值为
\(
\widehat N=239.
\)
\end{solution}

\begin{exercise}[离散分布的最大似然估计] \label{ex:mle_01}
设总体 \(X\) 的概率分布为
\(
P\{X=1\}=\frac{1-\theta}{2},\quad P\{X=2\}=P\{X=3\}=\frac{1+\theta}{4}.
\)
利用来自总体 \(X\) 的样本值 \(1,3,2,2,1,3,1,2\)，可得 \(\theta\) 的最大似然估计值为（\quad）。\par
A. \(\dfrac{1}{4}\)\qquad B. \(\dfrac{3}{8}\)\qquad C. \(\dfrac{1}{2}\)\qquad D. \(\dfrac{5}{8}\)
\end{exercise}
\begin{solution}
应选 A。样本中 \(X=1\) 出现 3 次，\(X=2\) 出现 3 次，\(X=3\) 出现 2 次。似然函数：
\(
L(\theta)=\left(\frac{1-\theta}{2}\right)^{\!3}\left(\frac{1+\theta}{4}\right)^{\!5}.
\)
对数似然函数：
\[
\ln L(\theta)=-3\ln 2-5\ln 4+3\ln(1-\theta)+5\ln(1+\theta).
\]
令导数为零：
\[
\frac{\mathrm{d}}{\mathrm{d}\theta}\ln L(\theta)=-\frac{3}{1-\theta}+\frac{5}{1+\theta}=0
\implies -3(1+\theta)+5(1-\theta)=0
\implies \theta=\frac{1}{4}.
\]
\end{solution}

\begin{exercise}[矩估计与最大似然估计]\label{ex:est_02}\label{exer:src-2017-201207-a-08-34credit}
已知总体 \(X\) 的概率密度为
\[
f(x)=
\begin{cases}
(1+\theta)x^\theta, & 0<x<1,\\
0, & \text{其他},
\end{cases}
\]
其中 \(\theta>-1\) 未知。\(X_1,X_2,\dots,X_n\) 为简单随机样本。求：\par
(1) \(\theta\) 的矩估计量；\par
(2) \(\theta\) 的最大似然估计量；\par
(3) \(E(X)\) 的最大似然估计量。
\end{exercise}
\begin{solution}
\textbf{(1) 矩估计。}
\[
\mu_1=EX=\int_0^1(1+\theta)x^{\theta+1}\,\mathrm{d}x=(1+\theta)\left[\frac{x^{\theta+2}}{\theta+2}\right]_0^1=\frac{1+\theta}{2+\theta}.
\]
令 \(\mu_1=\overline{X}\)，即 \(\dfrac{1+\theta}{2+\theta}=\overline{X}\)，解得
\[
1+\theta=(2+\theta)\overline{X}\implies \theta-\theta\overline{X}=2\overline{X}-1\implies \hat{\theta}_{\text{矩}}=\frac{2\overline{X}-1}{1-\overline{X}}.
\]
\textbf{(2) 最大似然估计。}似然函数
\[
L(\theta)=\prod_{i=1}^{n}(1+\theta)x_i^\theta=(1+\theta)^n\left(\prod_{i=1}^{n}x_i\right)^{\!\theta},\quad 0<x_i<1.
\]
\[
\ln L(\theta)=n\ln(1+\theta)+\theta\sum_{i=1}^{n}\ln x_i.
\]
\[
\frac{\mathrm{d}}{\mathrm{d}\theta}\ln L(\theta)=\frac{n}{1+\theta}+\sum_{i=1}^{n}\ln x_i=0
\implies \hat{\theta}_{\text{MLE}}=-\frac{n}{\displaystyle\sum_{i=1}^{n}\ln X_i}-1.
\]

\textbf{(3) \(E(X)\) 的最大似然估计。}由
\(
E(X)=\frac{\theta+1}{\theta+2},
\)
代入 \(\hat\theta_{\text{MLE}}\)。记
\[
\overline L=\frac1n\sum_{i=1}^n\ln X_i,
\qquad
\hat\theta_{\text{MLE}}=-1-\frac1{\overline L}.
\]
则
\[
\widehat{E(X)}_{\text{MLE}}
=\frac{\hat\theta_{\text{MLE}}+1}{\hat\theta_{\text{MLE}}+2}
=\frac{1}{1-\overline L}
=\frac{1}{1-\frac1n\sum_{i=1}^n\ln X_i}.
\]
\end{solution}


\begin{exercise}\label{exer:7-poisson-bacteria}
为检验某种自然水消毒设备的效果，从消毒后的水中随机抽取 50 升，化验每升水中大肠杆菌的个数（服从泊松分布），结果如下：
\begin{center}
\begin{tabular}{c|ccccccc}
大肠杆菌数/升 & 0 & 1 & 2 & 3 & 4 & 5 & 6 \\
\hline
升数 & 17 & 20 & 10 & 2 & 1 & 0 & 0
\end{tabular}
\end{center}
问：平均每升水中大肠杆菌的个数为多少时，出现上述情况的概率最大？
\end{exercise}
\begin{solution}
设每升水中大肠杆菌个数为 \(X_i\sim P(\lambda)\)，且各升之间可视为相互独立。题目所问“平均每升多少个时，出现当前观测结果的概率最大”，本质上就是要求参数 \(\lambda\) 的最大似然估计。

对于泊松总体，最大似然估计量是样本均值，因此只需根据频数表计算
\[
\overline{x} = \frac{0\times17+1\times20+2\times10+3\times2+4\times1}{50} = \frac{50}{50} = 1.
\]
所以
\(
\hat{\lambda}=1.
\)
也就是说，当平均每升大肠杆菌个数取 \(1\) 时，观察到题目所给结果的似然最大。
\end{solution}

\begin{exercise}\label{exer:7-normal-sigma2}
设 \((X_1, X_2, \dots, X_n)\) 为取自总体 \(N(0, \sigma^2)\) 的样本，则 \(\sigma^2\) 的最大似然估计是（\quad）。\par
A. \(S^2\)\qquad B. \(\dfrac{n-1}{n}S^2\)\qquad C. \(\overline{X}^2\)\qquad D. \(\dfrac{1}{n}\displaystyle\sum_{i=1}^{n}X_i^2\)
\end{exercise}
\begin{solution}
应选 D。关键在于本题中总体均值 \(\mu=0\) 是已知的，因此似然函数应直接围绕 \(\sum x_i^2\) 展开，而不是围绕样本方差 \(S^2\) 展开。

样本的似然函数为
\[
L(\sigma^2) = \frac{1}{(\sqrt{2\pi\sigma^2})^n}\exp\!\left(-\frac{1}{2\sigma^2}\sum_{i=1}^{n}x_i^2\right)~~~~~
\ln L = -\frac{n}{2}\ln(2\pi) - \frac{n}{2}\ln\sigma^2 - \frac{1}{2\sigma^2}\sum_{i=1}^{n}x_i^2.
\]
令 \(\dfrac{\partial\ln L}{\partial\sigma^2} = -\dfrac{n}{2\sigma^2} + \dfrac{1}{2(\sigma^2)^2}\displaystyle\sum_{i=1}^{n}x_i^2 = 0\)，得
\(
\hat{\sigma}^2 = \frac{1}{n}\sum_{i=1}^{n}X_i^2.
\)
因此最大似然估计量正是选项 D。
\end{solution}


\begin{proposition}\label{prop:mle-mme-equiv} 
    对于某些单参数指数族模型，如果其概率质量函数或密度函数可写成
    \(
    P(X=x; \theta)=h(x)c(\theta)\exp\{w(\theta)t(x)\},
    \)
    其中 \(t(x)\) 体现了样本对参数的“自然统计量”结构。若它还能化简为 \(t(x)=x\)，那么一阶矩方程与最大似然方程往往会写成同一个形式，因此矩估计与最大似然估计常常一致。几何分布与本文前面的二项分布就是这样的典型例子。
\end{proposition}
\begin{note}
这个结论\textbf{什么时候可以直接用}：当模型确实能化成单参数指数族，且自然统计量正好就是 \(t(x)=x\) 时，可以迅速判断“一阶矩方程”和“似然方程”会落到同一条等式上。不要把它误推广成“任何单参数分布都满足 MME=MLE”，更不能在多参数模型、\(t(x)\neq x\) 或边界型似然问题里机械套用。
\end{note}
\begin{proof}
\begin{enumerate}
    \item 写出对数似然函数：$$\ln L(\theta) = \sum_{i=1}^n \ln P(X_i; \theta) = \sum_{i=1}^n \ln h(X_i) + n\ln c(\theta) + w(\theta)\sum_{i=1}^n X_i$$
    \item 求导寻找最大似然估计（MLE）：将对数似然函数对 $\theta$ 求导并令其等于0：$$\frac{\mathrm{d}\ln L}{\mathrm{d}\theta} = n\frac{c'(\theta)}{c(\theta)} + w'(\theta)\sum_{i=1}^n X_i = 0$$等式两边同除以 $n \cdot w'(\theta)$，我们得到 MLE 的核心方程：\(-\frac{c'(\theta)}{c(\theta)w'(\theta)} = \frac{1}{n}\sum_{i=1}^n X_i = \overline{X}\)
    \item 再来看看理论期望（总体一阶矩）：在数理统计中有一个基础定理（正则条件下对数似然导数的期望为0），即 $E\left[ \frac{\mathrm{d}\ln P(X; \theta)}{\mathrm{d}\theta} \right] = 0$。代入单样本的导数 $\frac{c'(\theta)}{c(\theta)} + w'(\theta)X$，取期望得到：$$\frac{c'(\theta)}{c(\theta)} + w'(\theta)E(X) = 0 \implies E(X) = -\frac{c'(\theta)}{c(\theta)w'(\theta)}$$
    \item 将第 2 步中的 \(\overline{X}\) 与第 3 步得到的 \(E(X)\) 对应起来，可以发现两者都归结为同一个方程：\(E(X)=\overline{X}\)这正是一阶矩估计的基本方程。
\end{enumerate}
结论：在本题这类“单参数指数族且 \(t(x)=x\)”的模型中，最大似然方程会化成 \(\overline{X}\) 与总体一阶矩的对应关系，因此 MLE 与 MME 往往会得到同一个估计量。但这只是这类模型下的典型现象，不能推广到所有分布。
\end{proof}

\begin{exercise}\label{exer:7-geometric}
设总体 \(X\) 的分布律为 \(P(X=x) = p(1-p)^{x-1}\ (x=1,2,\dots)\)，其中 \(p\ (0<p<1)\) 未知。\(X_1, \dots, X_n\) 为样本，求 \(p\) 的矩估计量和最大似然估计量。
\end{exercise}
\begin{solution}
几何分布的数学期望
\begin{align*}
E(X) = \sum_{x=1}^\infty x \cdot p(1-p)^{x-1} = p \sum_{x=1}^\infty x(1-p)^{x-1} = p \cdot \frac{1}{[1-(1-p)]^2} = p \cdot \frac{1}{p^2} = \frac{1}{p}
\end{align*}
矩估计的核心思想是用样本矩代替总体矩，建立方程求解未知参数。对于单参数问题，我们用一阶原点矩（期望）来求解
即：
$$E(X) = \overline{X},~~~\text{其中}\overline{X} = \frac{1}{n}\sum_{i=1}^n X_i\implies \frac1p=\overline{X}\implies \hat{p}_{\text{矩}} = \frac1{\overline{X}}$$
\textbf{最大似然估计}是想找到一个参数值 $\hat{p}$，使得当前样本出现的概率最大（即似然函数取最大值），似然函数是样本的联合概率（离散型），即所有样本观测值概率的乘积：
\begin{align*}
L(p) &= \prod_{i=1}^n P(X=x_i) = \prod_{i=1}^n \left[ p(1-p)^{x_i-1} \right]=  p^n \cdot (1-p)^{\left( \sum_{i=1}^n x_i \right) - n}\\
\ln L(p) &= n\ln p + \left( \sum_{i=1}^n X_i - n \right) \ln(1-p)\\
\frac{\mathrm{d}}{\mathrm{d}p} \ln L(p) &= \frac{n}{p} - \frac{\sum_{i=1}^n X_i - n}{1-p} = 0\implies n(1-p) = p\left( \sum_{i=1}^n X_i - n \right),~n = p\sum_{i=1}^n X_i
\end{align*}
故 \(\hat{p} = \frac{1}{\overline{X}}\)
\end{solution}

\begin{exercise}\label{exer:7-discrete-3val}
设总体 \(X\) 的概率分布为 \(P\{X=0\} = \theta^2\)，\(P\{X=1\} = 2\theta(1-\theta)\)，\(P\{X=2\} = (1-\theta)^2\)，其中 \(\theta\ (0 < \theta < 1/2)\) 是未知参数。利用样本值 2, 0, 2, 2, 1, 0, 2, 2，求 \(\theta\) 的矩估计值和最大似然估计值。
\end{exercise}
\begin{solution}
\textbf{矩估计}：
\(
EX = 2\theta(1-\theta) + 2(1-\theta)^2 = 2-2\theta.
\)
\(\overline{x} = \frac{11}{8}\)，令 \(2-2\theta = \frac{11}{8}\)，解得 \(\hat{\theta}_{\text{矩}} = \frac{5}{16}\)。\\
\textbf{最大似然估计}：样本值中 \(X=0\) 出现2次，\(X=1\) 出现1次，\(X=2\) 出现5次。
\[
L(\theta) = \theta^4\cdot2\theta(1-\theta)\cdot(1-\theta)^{10} = 2\theta^5(1-\theta)^{11}.
\]
\[
\ln L = \ln 2 + 5\ln\theta + 11\ln(1-\theta),\quad \frac{\mathrm{d}\ln L}{\mathrm{d}\theta} = \frac{5}{\theta} - \frac{11}{1-\theta} = 0.
\]
解得 \(\hat{\theta}_{\text{MLE}} = \frac{5}{16}\)。
\end{solution}

\begin{exercise}\label{exer:7-discrete-123}
设总体 \(X\) 的分布列如下：
\begin{center}
\begin{tabular}{c|ccc}
\(X\) & 1 & 2 & 3 \\
\hline
\(P\) & \(\theta^2\) & \(2\theta(1-\theta)\) & \((1-\theta)^2\)
\end{tabular}
\end{center}
其中 \(\theta\) 为未知参数，样本观测值为 \(X_1=1, X_2=2, X_3=1\)。求 \(\theta\) 的矩估计值和最大似然估计值。
\end{exercise}
\begin{solution}
\textbf{矩估计}：
\[
EX = 1\cdot\theta^2 + 2\cdot2\theta(1-\theta) + 3(1-\theta)^2 = 3-2\theta.
\]
\(\overline{x} = \frac{4}{3}\)，令 \(3-2\theta = \frac{4}{3}\)，得 \(\hat{\theta}_{\text{矩}} = \frac56\)。\\
\textbf{最大似然估计}：
\[
L(\theta) = \theta^2 \cdot 2\theta(1-\theta) \cdot \theta^2 = 2\theta^5(1-\theta).
\]
\(\ln L = \ln 2 + 5\ln\theta + \ln(1-\theta)\)，令 \(\dfrac{5}{\theta} - \dfrac{1}{1-\theta} = 0\)，得 \(\hat{\theta}_{\text{MLE}} = \frac56\)。
\end{solution}

\begin{note}
    几何分布和二项分布都属于指数族分布。下面给出一些别的例子。
\end{note}

\begin{exercise}\label{exer:7-discrete-4val}\label{exer:src-paper1-app-05}\label{prob:p1-app-05}
设总体 \(X\) 的概率分布为
\begin{center}
\begin{tabular}{c|cccc}
\(X\) & 0 & 1 & 2 & 3 \\
\hline
\(P\) & \(\theta^2\) & \(2\theta(1-\theta)\) & \(\theta^2\) & \(1-2\theta\)
\end{tabular}
\end{center}
其中 \(\theta\ (0 < \theta < 1/2)\) 是未知参数。从总体 \(X\) 中抽取容量为 8 的样本，其样本值为 3, 1, 3, 0, 3, 1, 2, 3。求 \(\theta\) 的矩估计值和最大似然估计值。
\end{exercise}
\begin{solution}
\textbf{矩估计}：总体仅含一个参数，用一阶矩。
\[
EX = 0\cdot\theta^2 + 1\cdot2\theta(1-\theta) + 2\cdot\theta^2 + 3(1-2\theta) = 3-4\theta.
\]
令 \(3-4\theta = \overline{X}\)，\(\overline{x} = \dfrac{3+1+3+0+3+1+2+3}{8} = 2\)，故$\hat{\theta}_{\text{矩}} = \dfrac{3-2}{4} = \dfrac{1}{4}$.\\
\textbf{最大似然估计}：样本值中 \(X=0\) 出现1次，\(X=1\) 出现2次，\(X=2\) 出现1次，\(X=3\) 出现4次。
\[
L(\theta) = \theta^2\cdot[2\theta(1-\theta)]^2\cdot\theta^2\cdot(1-2\theta)^4 = 4\theta^6(1-\theta)^2(1-2\theta)^4.
\]
\[
\ln L(\theta) = \ln 4 + 6\ln\theta + 2\ln(1-\theta) + 4\ln(1-2\theta).
\]
令 \(\dfrac{\mathrm{d}\ln L}{\mathrm{d}\theta} = \dfrac{6}{\theta} - \dfrac{2}{1-\theta} - \dfrac{8}{1-2\theta} = 0\)，化简得 \(12\theta^2 - 14\theta + 3 = 0\)。解得$\theta = \dfrac{14\pm\sqrt{196-144}}{24} = \dfrac{7\pm\sqrt{13}}{12}$因为 \(\dfrac{7+\sqrt{13}}{12} > \dfrac{1}{2}\)（不合题意），故
\(
\hat{\theta}_{\text{MLE}} = \frac{7-\sqrt{13}}{12}.
\)
\end{solution}


\begin{exercise}\label{exer:7-discrete-0123}
设总体 \(X\) 的分布列为
\begin{center}
\begin{tabular}{c|cccc}
\(X\) & 0 & 1 & 2 & 3 \\
\hline
\(P\) & \(\dfrac{\theta(1-\theta)}{2}\) & \(\theta^2\) & \(\dfrac{\theta(1-\theta)}{2}\) & \(1-\theta\)
\end{tabular}
\end{center}
其中 \(\theta\ (0 < \theta < 1)\) 未知。样本观测值为 0, 0, 1, 2, 2, 3, 3, 3, 0, 2。求 \(\theta\) 的矩估计值与最大似然估计值。
\end{exercise}
\begin{solution}
频数统计：\(X=0\) 出现3次，\(X=1\) 出现1次，\(X=2\) 出现3次，\(X=3\) 出现3次。\(\overline{x} = 1.6\)。
\textbf{矩估计}：
\[
EX = \theta^2 + \theta(1-\theta) + 3(1-\theta) = 3-2\theta.
\]
令 \(3-2\theta = 1.6\)，得 \(\hat{\theta}_{\text{矩}} = 0.7\)。\\
\textbf{最大似然估计}：
\[
L(\theta) = \left[\frac{\theta(1-\theta)}{2}\right]^{\!3}\theta^2\!\left[\frac{\theta(1-\theta)}{2}\right]^{\!3}(1-\theta)^3 = \frac{1}{64}\,\theta^8(1-\theta)^9.
\]
\(\ln L = 8\ln\theta + 9\ln(1-\theta) - \ln 64\)，令 \(\dfrac{8}{\theta} - \dfrac{9}{1-\theta} = 0\)，得 \(\hat{\theta}_{\text{MLE}} = \dfrac{8}{17}\)。
\end{solution}

\begin{exercise}\label{exer:src-2004-3a-08}
设总体 \(X\) 服从几何分布，即
\(
P\{X=k\}=(1-p)^{k-1}p,\qquad k=1,2,\dots,
\)
其中 \(0<p<1\)。现从 \(X\) 中抽得容量为 \(n\) 的样本，其一组观察值为 \(x_1,\dots,x_n\)，求参数 \(p\) 的最大似然估计。
\end{exercise}
\begin{solution}
样本联合概率为
\[
L(p)=\prod_{i=1}^n (1-p)^{x_i-1}p
=p^n(1-p)^{\sum_{i=1}^n x_i-n},\qquad 0<p<1.
\]
取对数似然函数：
\(
\ell(p)=n\ln p+\left(\sum_{i=1}^n x_i-n\right)\ln(1-p).
\)
求导得
\(
\ell'(p)=\frac{n}{p}-\frac{\sum_{i=1}^n x_i-n}{1-p}.
\)
令 \(\ell'(p)=0\)，有
\(
n(1-p)=p\left(\sum_{i=1}^n x_i-n\right),
\)
即
\(
n=p\sum_{i=1}^n x_i.
\)
所以
\(
\hat p=\frac{n}{\sum_{i=1}^n x_i}=\frac{1}{\overline X}.
\)
由于
\(
\ell''(p)=-\frac{n}{p^2}-\frac{\sum_{i=1}^n x_i-n}{(1-p)^2}<0,
\)
该驻点给出最大值。因此 \(p\) 的最大似然估计量为
\(
\boxed{\hat p=\dfrac1{\overline X}}.
\)
\end{solution}

\begin{remark}
以上是离散型的，我们再来看一下连续型的。
\end{remark}

\begin{exercise}\label{exer:7-poisson-tri}
设总体 \(X\) 服从泊松分布 \(P(\theta)\)，总体 \(Y\) 具有密度函数 \(f(y,\lambda) = \dfrac{6y}{\lambda^3}(\lambda-y)\ (0 < y < \lambda)\)。\(X_1, \dots, X_n\) 及 \(Y_1, \dots, Y_n\) 分别为来自 \(X\) 和 \(Y\) 的样本。\par
(1) 求 \(\theta\) 的最大似然估计 \(\hat{\theta}\)；~~(2) 求 \(\lambda\) 的矩估计 \(\hat{\lambda}\)；~~(3) 求 \(\hat{\lambda}\) 的方差。
\end{exercise}
\begin{solution}
(1) 即 \(P(X=x) = e^{-\theta}\dfrac{\theta^x }{x!},\ \theta > 0\)，泊松分布的似然函数：
\[
L(\theta) = \prod_{i=1}^{n}\frac{\theta^{x_i}e^{-\theta}}{x_i!} = e^{-n\theta}\,\theta^{\sum x_i}\prod_{i=1}^{n}\frac{1}{x_i!}.
\]
\[
\ln L = -n\theta + \sum x_i\ln\theta + C,\quad \frac{\partial\ln L}{\partial\theta} = -n + \frac{\sum x_i}{\theta} = 0\implies \hat{\theta} = \overline{X}
\]
(2) \[EY = \displaystyle\int_0^\lambda y\cdot\frac{6y}{\lambda^3}(\lambda-y)\,\mathrm{d}y = \frac{6}{\lambda^3}\!\left[\lambda\int_0^\lambda y^2\,\mathrm{d}y - \int_0^\lambda y^3\,\mathrm{d}y\right] = \frac{6}{\lambda^3}\!\left(\frac{\lambda^4}{3}-\frac{\lambda^4}{4}\right) = \frac{\lambda}{2}\]
令 \(\dfrac{\lambda}{2} = \overline{Y}\)，得 \(\hat{\lambda} = 2\overline{Y}\)。\\
(3) \[EY^2 = \displaystyle\int_0^\lambda \frac{6y^3}{\lambda^3}(\lambda-y)\,\mathrm{d}y = \frac{6}{\lambda^3}\!\left(\frac{\lambda^5}{4}-\frac{\lambda^5}{5}\right) = \frac{3\lambda^2}{10}\]
\(DY = EY^2 - (EY)^2 = \dfrac{3\lambda^2}{10} - \dfrac{\lambda^2}{4} = \dfrac{\lambda^2}{20}\)。
\(D(\hat{\lambda}) = D(2\overline{Y}) = \dfrac{4}{n}\,DY = \dfrac{\lambda^2}{5n}\)。
\end{solution}

\begin{note}
对于单调型似然函数（\(\ln L\) 关于 \(\theta\) 单调），\(\hat{\theta}\) 通常在参数范围的边界取得。常见模式：若 \(L(\theta)\) 关于 \(\theta\) 单调递减，则 \(\hat{\theta} = \max\{X_i\}\)；若单调递增，则 \(\hat{\theta} = \min\{X_i\}\)。
\end{note}

\begin{exercise}[单减取最大，单增取最小]\label{exer:7-unif-theta}
设总体 \(X\) 的概率密度为
\[
f(x;\theta) = \begin{cases} \dfrac{1}{1-\theta}, & \theta \leq x \leq 1, \\ 0, & \text{其他}, \end{cases}
\]
其中 \(\theta\) 为未知参数，\(X_1, X_2, \dots, X_n\) 为简单随机样本。\par
(1) 求 \(\theta\) 的矩估计量；\par
(2) 求 \(\theta\) 的最大似然估计量。
\end{exercise}
\begin{solution}
(1) 总体 \(X \sim U(\theta, 1)\)，\(EX = \dfrac{1+\theta}{2}\)。令 \(\dfrac{1+\theta}{2} = \overline{X}\)，得 \(\hat{\theta}_{\text{矩}} = 2\overline{X}-1\)。\\
(2) 似然函数为\footnote{密度为 0 的区间直接被过滤了，否则整个乘积为 0，似然函数取最小值，不可能对应最大似然估计。}
\[
L(\theta) = \frac{1}{(1-\theta)^n},\quad \text{约束条件：}\theta \leq \min\{x_1, \dots, x_n\},\ x_i \leq 1.
\]
\(L(\theta)\) 关于 \(\theta\) 单调递增（分母 \(1-\theta\) 递减），故当 \(\theta = \min\{X_1, \dots, X_n\}\) 时 \(L\) 最大：
\(
\hat{\theta}_{\text{MLE}} = \min\{X_1, X_2, \dots, X_n\}.
\)
\end{solution}

\begin{exercise}\label{exer:src-2010-0630-a-08-34credit}
设总体 \(X\) 的密度函数为
\[
f(x;\theta)=
\begin{cases}
\dfrac1{|\theta|}, & \theta<x<\theta+|\theta|,\\
0, & \text{其他}.
\end{cases}
\]
分别在以下两种情形下求参数 \(\theta\) 的最大似然估计量：\par
(1) \(\theta<0\)；\par
(2) \(\theta>0\)。
\end{exercise}
\begin{solution}
设样本为 \(X_1,\dots,X_n\)，记
\[
X_{(1)}=\min_{1\le i\le n}X_i,\qquad
X_{(n)}=\max_{1\le i\le n}X_i.
\]

(1) 当 \(\theta<0\) 时，\(|\theta|=-\theta\)，总体的支撑区间为
\(
\theta<x<0.
\)
因此似然函数为
\[
L(\theta)=\left(-\frac1\theta\right)^n
\mathbf{1}\{\theta<X_{(1)},\ X_{(n)}<0\},\qquad \theta<0.
\]
在允许范围 \(\theta<X_{(1)}\) 内，\(\left(-1/\theta\right)^n\) 随 \(\theta\) 增大而增大，故最大值在边界处取得：
\(
\hat\theta=X_{(1)}=\min\{X_1,\dots,X_n\}.
\)

(2) 当 \(\theta>0\) 时，\(|\theta|=\theta\)，总体的支撑区间为
\(
\theta<x<2\theta.
\)
因此样本必须满足
\(
\theta<X_{(1)},\qquad X_{(n)}<2\theta,
\)
即
\(
\frac{X_{(n)}}2<\theta<X_{(1)}.
\)
似然函数为
\[
L(\theta)=\theta^{-n}
\mathbf{1}\left\{\frac{X_{(n)}}2<\theta<X_{(1)}\right\}.
\]
在允许范围内，\(\theta^{-n}\) 随 \(\theta\) 增大而减小，故最大值在左端边界处取得：
\[
\hat\theta=\frac12X_{(n)}=\frac12\max\{X_1,\dots,X_n\}.
\]
\end{solution}

\begin{exercise}\label{exer:src-2012-a-06-34credit}
设总体 \(X\) 的分布律为
\[
\begin{array}{c|ccc}
X & -1 & 0 & 3\\
\hline
P & \theta & 3\theta & 1-4\theta
\end{array}
\]
其中 \(0<\theta<\dfrac14\) 未知。设 \(X_1,X_2,\dots,X_n\) 为来自该总体的简单随机样本。求：\par
(1) \(\theta\) 的矩估计量 \(\hat\theta_1\)；\par
(2) \(\theta\) 的极大似然估计量 \(\hat\theta_2\)；\par
(3) 讨论 \(\hat\theta_1,\hat\theta_2\) 是否为无偏估计量。
\end{exercise}
\begin{solution}
(1) 先求总体均值：
\[
EX=(-1)\theta+0\cdot3\theta+3(1-4\theta)=3-13\theta.
\]
令 \(EX=\overline X\)，得矩估计量
\(
\hat\theta_1=\frac{3-\overline X}{13}.
\)

(2) 设样本中 \(-1,0,3\) 出现的次数分别为 \(r_1,r_2,r_3\)，则 \(r_1+r_2+r_3=n\)。似然函数为
\[
L(\theta)=\theta^{r_1}(3\theta)^{r_2}(1-4\theta)^{r_3}
=3^{r_2}\theta^{r_1+r_2}(1-4\theta)^{r_3}.
\]
对数似然函数为
\[
\ell(\theta)=r_2\ln3+(r_1+r_2)\ln\theta+r_3\ln(1-4\theta).
\]
若 \(r_1+r_2>0\)，令导数为零：
\(
\ell'(\theta)=\frac{r_1+r_2}{\theta}-\frac{4r_3}{1-4\theta}=0.
\)
解得
\(
\hat\theta_2=\frac{r_1+r_2}{4n}
=\frac{1-r_3/n}{4}.
\)
若样本全为 \(3\)，则上式给出边界极限值 \(0\)，可理解为似然在闭区间扩展下的边界估计。

(3) 对矩估计量，
\[
E\hat\theta_1=\frac{3-E\overline X}{13}
=\frac{3-(3-13\theta)}{13}=\theta,
\]
故 \(\hat\theta_1\) 是无偏估计量。

对极大似然估计量，由于 \(r_3\sim B(n,1-4\theta)\)，
\(
E\left(\frac{r_3}{n}\right)=1-4\theta.
\)
于是
\[
E\hat\theta_2
=E\left(\frac{1-r_3/n}{4}\right)
=\frac{1-(1-4\theta)}4=\theta.
\]
因此 \(\hat\theta_2\) 也是无偏估计量。
\end{solution}

\begin{exercise}\label{exer:src-2012-b-06-34credit}
设总体 \(X\) 的分布律为
\[
\begin{array}{c|ccc}
X & -1 & 1 & 2\\
\hline
P & \theta & 2\theta & 1-3\theta
\end{array}
\]
其中 \(0<\theta<\dfrac13\) 未知。设 \(X_1,X_2,\dots,X_n\) 为来自该总体的简单随机样本。求：\par
(1) \(\theta\) 的矩估计量 \(\hat\theta_1\)；\par
(2) \(\theta\) 的极大似然估计量 \(\hat\theta_2\)；\par
(3) 讨论 \(\hat\theta_1,\hat\theta_2\) 是否为无偏估计量。
\end{exercise}
\begin{solution}
(1) 先求总体均值：
\[
EX=(-1)\theta+1\cdot2\theta+2(1-3\theta)=2-5\theta.
\]
令 \(EX=\overline X\)，得矩估计量
\(
\hat\theta_1=\frac{2-\overline X}{5}.
\)

(2) 设样本中 \(-1,1,2\) 出现的次数分别为 \(r_1,r_2,r_3\)，且 \(r_1+r_2+r_3=n\)。似然函数为
\[
L(\theta)=\theta^{r_1}(2\theta)^{r_2}(1-3\theta)^{r_3}
=2^{r_2}\theta^{r_1+r_2}(1-3\theta)^{r_3}.
\]
对数似然函数为
\[
\ell(\theta)=r_2\ln2+(r_1+r_2)\ln\theta+r_3\ln(1-3\theta).
\]
若 \(r_1+r_2>0\)，令导数为零：
\(
\ell'(\theta)=\frac{r_1+r_2}{\theta}-\frac{3r_3}{1-3\theta}=0.
\)
由于 \(r_1+r_2=n-r_3\)，解得
\(
\hat\theta_2=\frac13\left(1-\frac{r_3}{n}\right).
\)
若样本全为 \(2\)，则上式给出边界极限值 \(0\)，可理解为似然在闭区间扩展下的边界估计。

(3) 对矩估计量，
\[
E\hat\theta_1=\frac{2-E\overline X}{5}
=\frac{2-(2-5\theta)}{5}=\theta,
\]
故 \(\hat\theta_1\) 是无偏估计量。

对极大似然估计量，\(r_3\sim B(n,1-3\theta)\)，所以
\(
E\left(\frac{r_3}{n}\right)=1-3\theta.
\)
于是
\[
E\hat\theta_2
=\frac13\left(1-E\frac{r_3}{n}\right)
=\frac13\{1-(1-3\theta)\}=\theta.
\]
因此 \(\hat\theta_2\) 也是无偏估计量。
\end{solution}

\begin{exercise}\label{exer:src-2011-ab-a-08-34credit}
设总体 \(X\) 的概率密度为
\[
f(x;\theta)=
\begin{cases}
2e^{-2(x-\theta)}, & x>\theta,\\
0, & x\le\theta,
\end{cases}
\]
其中 \(\theta>0\) 是未知参数。从总体 \(X\) 中抽取简单随机样本
\(X_1,X_2,\dots,X_n\)。求：\par
(1) \(\theta\) 的矩估计量 \(\hat\theta\)；\par
(2) 讨论 \(\hat\theta\) 是否具有无偏性。
\end{exercise}
\begin{solution}
令 \(Y=X-\theta\)，则 \(Y\) 服从参数为 \(2\) 的指数分布，因此
\(
EX=\theta+EY=\theta+\frac12.
\)
(1) 用一阶矩作矩估计，令 \(EX=\overline X\)，得
\(
\hat\theta=\overline X-\frac12.
\)

(2) 因为
\[
E\hat\theta=E\overline X-\frac12
=\left(\theta+\frac12\right)-\frac12
=\theta,
\]
所以 \(\hat\theta=\overline X-\dfrac12\) 是 \(\theta\) 的无偏估计量。
\end{solution}

\begin{exercise}\label{exer:src-2011-ab-b-08-34credit}
设总体 \(X\) 的概率密度为
\[
f(x;\theta)=
\begin{cases}
2\theta e^{-2\theta x}, & x>0,\\
0, & x\le0,
\end{cases}
\]
其中 \(\theta>0\) 是未知参数。从总体 \(X\) 中抽取简单随机样本
\(X_1,X_2,\dots,X_n\)。求：\par
(1) \(\theta\) 的极大似然估计 \(\hat\theta_1\) 和矩估计 \(\hat\theta_2\)；\par
(2) 讨论 \(\hat\theta_1\) 是否具有无偏性。
\end{exercise}
\begin{solution}
(1) 似然函数为
\[
L(\theta)=\prod_{i=1}^n 2\theta e^{-2\theta X_i}
=(2\theta)^n\exp\left(-2\theta\sum_{i=1}^n X_i\right).
\]
对数似然函数为
\[
\ell(\theta)=n\ln(2\theta)-2\theta\sum_{i=1}^n X_i.
\]
令
\(
\ell'(\theta)=\frac{n}{\theta}-2\sum_{i=1}^n X_i=0,
\)
得
\[
\hat\theta_1=\frac{n}{2\sum_{i=1}^n X_i}
=\frac1{2\overline X}.
\]
另一方面，\(X\) 服从参数为 \(2\theta\) 的指数分布，所以
\(
EX=\frac1{2\theta}.
\)
令 \(EX=\overline X\)，得矩估计量
\(
\hat\theta_2=\frac1{2\overline X}.
\)

(2) 记 \(S=\sum_{i=1}^n X_i\)。因为 \(X_i\) 独立同分布且服从参数为 \(2\theta\) 的指数分布，所以
\(
S\sim \Gamma(n,2\theta)
\)
（这里第二个参数为率参数）。当 \(n>1\) 时，
\(
E\left(\frac1S\right)=\frac{2\theta}{n-1}.
\)
于是
\[
E\hat\theta_1
=E\left(\frac{n}{2S}\right)
=\frac n2\cdot\frac{2\theta}{n-1}
=\frac{n}{n-1}\theta.
\]
因此 \(\hat\theta_1\) 不是 \(\theta\) 的无偏估计量；若需修正为无偏估计，可取
\(
\frac{n-1}{n}\hat\theta_1.
\)
\end{solution}




\begin{exercise}\label{exer:7-rayleigh}
设某种产品的使用寿命 \(T\) 的分布函数 \(F(t)\) 满足方程
\(
F'(t) + \frac{2t}{\theta^2}[F(t)-1] = 0,\quad t \geq 0,\quad F(0) = 0,
\)
其中 \(\theta > 1\) 为未知参数。产品性能 \(Q(\theta) = \dfrac{\theta^2}{2}\ln\theta - \dfrac{3}{4}\theta^2 + \theta\)。\par
(1) 求 \(P\{T > t\}\) 与 \(P\{T > s+t \mid T > s\}\)（\(s,t > 0\)）；\par
(2) 任取 \(n\) 个产品做寿命试验，测得寿命 \(t_1, \dots, t_n\)，求 \(\theta\) 的最大似然估计值 \(\hat{\theta}\)；\par
(3) 求 \(Q\) 的最大似然估计值 \(\hat{Q}\)。
\end{exercise}

\begin{solution}
\textbf{先求 \(F(t)\)}。微分方程 \(F'(t) + \dfrac{2t}{\theta^2}F(t) = \dfrac{2t}{\theta^2}\) 是一阶线性 ODE，积分因子为 \(e^{t^2/\theta^2}\)：
\[
F(t)\,e^{t^2/\theta^2} = \int e^{t^2/\theta^2}\,\mathrm{d}\!\left(\frac{t^2}{\theta^2}\right) + C = e^{t^2/\theta^2} + C.
\]
由 \(F(0)=0\) 得 \(C=-1\)，故
\(
F(t) = 1 - e^{-t^2/\theta^2}\quad(t \geq 0).
\)
(1) \(P\{T > t\} = e^{-t^2/\theta^2}\)。
\[
P\{T > s+t \mid T > s\} = \frac{P\{T > s+t\}}{P\{T > s\}} = \frac{e^{-(s+t)^2/\theta^2}}{e^{-s^2/\theta^2}} = e^{-(t^2+2st)/\theta^2}.
\]
(2) 密度 \(f(t;\theta) = F'(t) = \dfrac{2t}{\theta^2}\,e^{-t^2/\theta^2}\ (t > 0)\)。
\[
L(\theta) = \theta^{-2n}\cdot 2^n\!\left(\prod_{i=1}^{n}t_i\right)e^{-\sum t_i^2/\theta^2},\quad
\ln L = n\ln 2 - 2n\ln\theta + \sum\ln t_i - \frac{1}{\theta^2}\sum t_i^2.
\]
令 \(\dfrac{\mathrm{d}\ln L}{\mathrm{d}\theta} = -\dfrac{2n}{\theta} + \dfrac{2}{\theta^3}\sum t_i^2 = 0\)，得 \(\hat{\theta} = \displaystyle\left(\dfrac{1}{n}\sum_{i=1}^{n}t_i^2\right)^{\!1/2}\)。\\
(3) 由不变性原则，\(Q'(\theta) = \theta\ln\theta - \theta + 1>0\)，故 \(Q(\theta)\) 有反函数。
\[
\hat{Q} = Q(\hat{\theta}) = \frac{\hat{\theta}^2}{2}\ln\hat{\theta} - \frac{3}{4}\hat{\theta}^2 + \hat{\theta}.
\]
\end{solution}



\begin{exercise}\label{exer:src-paper3-i-06}\label{prob:p3-i-06}\label{exer:src-paper1-tf-05}\label{prob:p1-tf-05}
样本均值的平方 \(\overline X^2\) 是总体期望平方 \(\mu^2\) 的无偏估计。判断正误。
\end{exercise}
\begin{solution}
错误。

若总体方差为 \(\sigma^2\)，则
\[
E\overline X=\mu,\qquad D\overline X=\frac{\sigma^2}{n}.
\]
因此
\[
E(\overline X^2)=D\overline X+(E\overline X)^2
=\frac{\sigma^2}{n}+\mu^2.
\]
一般情况下 \(\sigma^2>0\)，所以
\(
E(\overline X^2)\ne \mu^2.
\)
故 \(\overline X^2\) 不是 \(\mu^2\) 的无偏估计量。
\end{solution}

\begin{exercise}\label{exer:src-paper3-iv-05}\label{prob:p3-iv-05}
已知随机变量 \(X\) 的密度函数为
\[
f(x)=
\begin{cases}
(\theta+1)(x-5)^\theta, & 5<x<6,\\
0, & \text{其他},
\end{cases}
\qquad \theta>0.
\]
其中 \(\theta\) 为未知参数。求 \(\theta\) 的矩估计量与极大似然估计量。
\end{exercise}
\begin{solution}
先求一阶矩。令 \(u=x-5\)，则 \(0<u<1\)，且密度为 \((\theta+1)u^\theta\)。于是
\[
E(X)=5+E(U)
=5+\int_0^1 u(\theta+1)u^\theta\,\mathrm{d}u
=5+\frac{\theta+1}{\theta+2}
=6-\frac1{\theta+2}.
\]
令 \(E(X)=\overline X\)，得
\(
\overline X=6-\frac1{\theta+2},
\)
故矩估计量为
\(
\hat\theta_M=\frac1{6-\overline X}-2.
\)

设样本为 \(X_1,\dots,X_n\)。当 \(5<X_i<6\) 时，似然函数为
\(
L(\theta)=(\theta+1)^n\prod_{i=1}^n (X_i-5)^\theta.
\)
对数似然为
\[
\ell(\theta)=n\ln(\theta+1)+\theta\sum_{i=1}^n\ln(X_i-5).
\]
求导：
\(
\ell'(\theta)=\frac{n}{\theta+1}+\sum_{i=1}^n\ln(X_i-5).
\)
令 \(\ell'(\theta)=0\)，得
\(
\hat\theta_L
=-\frac{n}{\sum_{i=1}^n\ln(X_i-5)}-1.
\)
由于 \(0<X_i-5<1\)，故 \(\ln(X_i-5)<0\)，该估计量在通常样本下有意义；并且
\(
\ell''(\theta)=-\frac{n}{(\theta+1)^2}<0,
\)
所以它给出极大值点。
\end{solution}

\section{估计量的评价标准} \label{subsec:criteria}

\begin{definition}[无偏性] \label{def:unbiased}
若 \(E\hat{\theta} = \theta\)，则称 \(\hat{\theta}\) 为 \(\theta\) 的\textbf{无偏估计量}。也就是说，从重复抽样的长期平均效果看，这个估计量既不系统性偏大，也不系统性偏小。
\end{definition}

\begin{note}
无偏性评价的是“平均是否对准真值”，不是“每次估得是否都接近真值”。一个无偏估计量仍可能波动很大；一个有偏估计量也可能因为方差更小而在均方误差意义下更好。因此无偏只是评价标准之一，不是唯一标准。
\end{note}

\begin{exercise}\label{exer:src-paper1-choice-05}\label{prob:p1-choice-05}
设 \((X_1,X_2,\dots,X_n)\) 为总体 \(N(\mu,\sigma^2)\) 的一个样本，\(\mu\) 已知，\(\overline X\) 为样本均值。则在总体方差 \(\sigma^2\) 的下列估计量中，为无偏估计量的是（\quad）。\par
A. \(\hat\sigma_1^2=\dfrac1n\sum_{i=1}^{n}(X_i-\overline X)^2\)\quad
B. \(\hat\sigma_2^2=\dfrac1{n-1}\sum_{i=1}^{n}(X_i-\overline X)^2\)\par
C. \(\hat\sigma_3^2=\dfrac1n\sum_{i=1}^{n}(X_i-\mu)^2\)\quad
D. \(\hat\sigma_4^2=\dfrac1{n-1}\sum_{i=1}^{n}(X_i-\mu)^2\)
\end{exercise}
\begin{solution}
应选 C。

由于 \(\mu\) 已知，直接围绕真实均值展开时有
\(
E\sum_{i=1}^{n}(X_i-\mu)^2
=\sum_{i=1}^{n}E(X_i-\mu)^2
=n\sigma^2.
\)
因此
\[
E\hat\sigma_3^2
=E\left[\frac1n\sum_{i=1}^{n}(X_i-\mu)^2\right]
=\sigma^2,
\]
所以 C 为无偏估计量。

A 的期望为 \(\dfrac{n-1}{n}\sigma^2\)，有偏；B 是在 \(\mu\) 未知、使用 \(\overline X\) 代替 \(\mu\) 时的无偏估计；D 的期望为 \(\dfrac{n}{n-1}\sigma^2\)，有偏。
\end{solution}

\begin{exercise}[无偏估计的构造与验证] \label{ex:unbiased_05}
设 \(X_1,X_2,\dots,X_n\) 是来自 \(N(\mu,\sigma^2)\) 的简单随机样本，记
\[
\overline{X}=\frac{1}{n}\sum_{i=1}^{n}X_i,\quad S^2=\frac{1}{n-1}\sum_{i=1}^{n}(X_i-\overline{X})^2,\quad T=\overline{X}^2-\frac{1}{n}S^2.
\]
证明：\(T\) 是 \(\mu^2\) 的无偏估计量。
\end{exercise}
\begin{solution}
要证明 \(T\) 是 \(\mu^2\) 的无偏估计量，只需验证 \(ET=\mu^2\)。

先对 \(T\) 逐项取期望：
\(
ET=E\!\left(\overline{X}^2-\frac{1}{n}S^2\right)
=E(\overline{X}^2)-\frac{1}{n}E(S^2).
\)

其中
\[
E(\overline{X})=\mu,\qquad D(\overline{X})=\frac{\sigma^2}{n},\qquad E(S^2)=\sigma^2.
\]
因此
\[
E(\overline{X}^2)=D(\overline{X})+[E(\overline{X})]^2
=\frac{\sigma^2}{n}+\mu^2.
\]
代回去可得
\(
ET=\left(\frac{\sigma^2}{n}+\mu^2\right)-\frac{1}{n}\sigma^2=\mu^2.
\)

这就说明 \(T\) 的期望恰好等于待估参数 \(\mu^2\)，所以 \(T\) 是 \(\mu^2\) 的无偏估计量。

\begin{note}
此题的技巧是：对 \(\overline{X}^2\) 使用 \(E(\overline{X}^2)=D(\overline{X})+(E\overline{X})^2\) 展开，再利用 \(E(S^2)=\sigma^2\) 消去方差项。此结论在数学三中也常以"求 \(ET\)"的形式出现。
\end{note}
\end{solution}



\begin{proposition}[系数和为1]\label{prop:unbiased-coeff}
若估计量形式为：
\(
\hat \mu = a_1 X_1 + a_2 X_2 + \dots + a_n X_n
\)
则
\[
\hat\mu \text{ 是 } \mu \text{ 的无偏估计} \iff a_1+a_2+\dots+a_n=1
\]
\end{proposition}
\begin{proof}
\[
E(\hat{\mu}) = E\!\left(\sum_{i=1}^{n}a_i X_i\right) = \sum_{i=1}^{n}a_i\,EX_i = \sum_{i=1}^{n}a_i\,\mu = \left(\sum_{i=1}^{n}a_i\right)\mu.
\]
故 \(E(\hat{\mu})=\mu \iff \displaystyle\sum_{i=1}^{n}a_i = 1\)。
\end{proof}

\begin{exercise}\label{exer:7-unbiased-judge}
设 \(X_1, X_2, X_3\) 是总体的样本，则下列统计量不是总体均值 \(\mu\) 的无偏估计量的是（\quad）。\par
A. \(\dfrac{1}{3}(X_1+X_2+X_3)\)~~~ B. \(\dfrac{1}{3}X_1+\dfrac{1}{6}X_2+\dfrac{1}{2}X_3\)~~~
C. \(\dfrac{1}{6}X_1+\dfrac{1}{3}X_2+\dfrac{1}{2}X_3\)~~~ D. \(\dfrac{1}{6}X_1+\dfrac{1}{6}X_2+\dfrac{1}{3}X_3\)
\end{exercise}
\begin{solution}
应选 D。

对于形如
\(
\hat\mu=a_1X_1+a_2X_2+a_3X_3
\)
的估计量，要成为总体均值 \(\mu\) 的无偏估计，充要条件是系数和等于 1，即 \(a_1+a_2+a_3=1\)。

逐项检查：

A 的系数和是 \(\frac13+\frac13+\frac13=1\)，所以无偏。

B 的系数和是 \(\frac13+\frac16+\frac12=1\)，所以无偏。

C 的系数和是 \(\frac16+\frac13+\frac12=1\)，所以无偏。

D 的系数和是 \(\frac16+\frac16+\frac13=\frac23\ne1\)，因此它不是 \(\mu\) 的无偏估计量。

故应选 D。
\end{solution}

\begin{exercise}\label{exer:src-2017-a-i-06}
设总体 \(X\sim N(\mu,\sigma^2)\)，\(X_1,X_2,X_3\) 为来自 \(X\) 的样本。若
\(
\widehat\mu=\frac14X_1+aX_2+\frac12X_3
\)
是未知参数 \(\mu\) 的无偏估计，则常数 \(a=\underline{\hspace{2cm}}\)。
\end{exercise}
\begin{solution}
应填 \(\dfrac14\)。因为
\(
E\widehat\mu
=\left(\frac14+a+\frac12\right)\mu.
\)
要使 \(\widehat\mu\) 是 \(\mu\) 的无偏估计量，必须有 \(E\widehat\mu=\mu\)，即
\(
\frac14+a+\frac12=1.
\)
解得
\(
a=\frac14.
\)
\end{solution}

\begin{exercise}\label{exer:7-unbiased-c}
设 \(X_1, X_2, \dots, X_n\) 是来自分布 \(f(x;\theta) = \dfrac{3x^2}{\theta^3},0 \leq x \leq \theta\) 的简单随机样本。若统计量 \(g(X_1, \dots, X_n) = c\overline{X}\) 为 \(\theta\) 的无偏估计，则常数 \(c\)为？
\end{exercise}
\begin{solution}
应填 \(\dfrac{4}{3}\)。

要使 \(g(X_1,\dots,X_n)=c\overline{X}\) 成为 \(\theta\) 的无偏估计，必须满足
\(
E(c\overline{X})=\theta.
\)
因此先求总体期望 \(EX\)。

由密度函数
\(
f(x;\theta)=\frac{3x^2}{\theta^3},\qquad 0\le x\le \theta,
\)
可得
\[
EX=\int_0^\theta x\cdot \frac{3x^2}{\theta^3}\,\mathrm dx
=\frac{3}{\theta^3}\int_0^\theta x^3\,\mathrm dx
=\frac{3}{\theta^3}\cdot\frac{\theta^4}{4}
=\frac{3\theta}{4}.
\]

由于 \(E(\overline{X})=EX=\dfrac{3\theta}{4}\)，于是
\(
E(c\overline{X})=cE(\overline{X})=c\cdot\frac{3\theta}{4}.
\)
令它等于 \(\theta\)，得
\(
c\cdot\frac{3\theta}{4}=\theta\implies c=\frac{4}{3}.
\)

所以常数 \(c=\dfrac{4}{3}\)。
\end{solution}

\begin{exercise}\label{exer:7-exp-mle}
设总体 \(X\) 具有概率密度
\[
f(x,\theta) = \begin{cases} \dfrac{1}{\theta}\,e^{-x/\theta}, & x > 0\ (\theta > 0), \\ 0, & x \leq 0, \end{cases}
\]
其中 \(\theta\) 未知，\(X_1, \dots, X_n\) 为样本。求 \(\theta\) 的最大似然估计 \(\hat{\theta}\)，并证明 \(\hat{\theta}\) 是 \(\theta\) 的无偏估计。
\end{exercise}
\begin{solution}
先求最大似然估计。

由样本独立性，似然函数为
\[
L(\theta)=\prod_{i=1}^n \frac{1}{\theta}e^{-X_i/\theta}
=\theta^{-n}\exp\!\left(-\frac{1}{\theta}\sum_{i=1}^n X_i\right),\qquad \theta>0.
\]

取对数得
\(
\ln L(\theta)=-n\ln\theta-\frac{1}{\theta}\sum_{i=1}^n X_i.
\)

对 \(\theta\) 求导：
\[
\frac{\mathrm{d}\ln L}{\mathrm{d}\theta}
=-\frac{n}{\theta}+\frac{1}{\theta^2}\sum_{i=1}^n X_i.
\]
令导数为 0，得
\[
-\frac{n}{\theta}+\frac{1}{\theta^2}\sum_{i=1}^n X_i=0
\quad\Longrightarrow\quad
\hat{\theta}=\frac{1}{n}\sum_{i=1}^n X_i=\overline{X}.
\]

再验证它确实对应极大值。二阶导数为
\[
\frac{\mathrm{d}^2\ln L}{\mathrm{d}\theta^2}
=\frac{n}{\theta^2}-\frac{2}{\theta^3}\sum_{i=1}^n X_i.
\]
在 \(\theta=\overline{X}\) 处，
\[
\left.\frac{\mathrm{d}^2\ln L}{\mathrm{d}\theta^2}\right|_{\theta=\overline{X}}
=\frac{n}{\overline{X}^2}-\frac{2n\overline{X}}{\overline{X}^3}
=-\frac{n}{\overline{X}^2}<0,
\]
所以 \(\hat{\theta}=\overline{X}\) 的确是最大似然估计。

下面证明无偏性。指数分布满足
\(
EX=\theta,
\)
因此
\[
E(\hat{\theta})=E(\overline{X})=\frac{1}{n}\sum_{i=1}^n E(X_i)=\frac{1}{n}\cdot n\theta=\theta.
\]

故 \(\hat{\theta}=\overline{X}\) 是 \(\theta\) 的无偏估计。
\end{solution}


\begin{exercise}\label{exer:7-unbiased-diff}
设 \(X_1, X_2, \dots, X_n\) 是来自正态总体 \(X \sim N(\mu, \sigma^2)\) 的简单随机样本，为使
\(
D = k\sum_{i=1}^{n-1}(X_{i+1}-X_i)^2
\)
成为总体方差 \(\sigma^2\) 的无偏估计量，则 \(k =\,\underline{\hspace{2cm}}\)。
\end{exercise}
\begin{solution}
要使 \(D\) 成为 \(\sigma^2\) 的无偏估计，只需令 \(ED=\sigma^2\)。

先看每一项 \((X_{i+1}-X_i)^2\) 的期望。由于 \(X_i,X_{i+1}\) 相互独立，且都服从 \(N(\mu,\sigma^2)\)，所以
\(
E(X_{i+1}-X_i)=\mu-\mu=0,
\)
\(
D(X_{i+1}-X_i)=DX_{i+1}+DX_i=2\sigma^2.
\)
从而
\[
E\bigl[(X_{i+1}-X_i)^2\bigr]
=D(X_{i+1}-X_i)+\bigl(E(X_{i+1}-X_i)\bigr)^2
=2\sigma^2.
\]

于是
\[
ED
=k\sum_{i=1}^{n-1}E\bigl[(X_{i+1}-X_i)^2\bigr]
=k(n-1)\cdot 2\sigma^2
=2k(n-1)\sigma^2.
\]

令它等于 \(\sigma^2\)，得到
\[
2k(n-1)\sigma^2=\sigma^2 \quad\Longrightarrow\quad k=\frac{1}{2(n-1)}.
\]

故所求 \(k=\frac{1}{2(n-1)}\)。
\end{solution}




\begin{exercise}\label{exer:7-two-normal-mle}
设 \(X_1, \dots, X_m\) 是取自总体 \(N(1, \sigma^2)\) 的简单随机样本，\(Y_1, \dots, Y_n\) 是取自总体 \(N(2, \sigma^2)\) 的简单随机样本。\par
(1) 求联合密度函数；\par
(2) 求 \(\sigma^2\) 的最大似然估计 \(\hat{\sigma}^2\)；\par
(3) \(\hat{\sigma}^2\) 是否为 \(\sigma^2\) 的无偏估计？
\end{exercise}
\begin{solution}
(1) 设 \(T = \displaystyle\sum_{i=1}^{m}(X_i-1)^2 + \sum_{j=1}^{n}(Y_j-2)^2\)，则
\(
f = (2\pi\sigma^2)^{-(m+n)/2}\,\exp\!\left\{-\frac{T}{2\sigma^2}\right\}.
\)
(2) \(\ln f = -\dfrac{m+n}{2}\ln(2\pi\sigma^2) - \dfrac{T}{2\sigma^2}\)，令 \(\dfrac{\partial\ln f}{\partial\sigma^2} = -\dfrac{m+n}{2\sigma^2}+\dfrac{T}{2(\sigma^2)^2}=0\)，得
\(
\hat{\sigma}^2 = \frac{1}{m+n}\!\left[\sum_{i=1}^{m}(X_i-1)^2+\sum_{j=1}^{n}(Y_j-2)^2\right].
\)
(3) 是。\(E\!\left[\sum(X_i-1)^2\right] = m\sigma^2\)，\(E\!\left[\sum(Y_j-2)^2\right] = n\sigma^2\)，故
\(
E(\hat{\sigma}^2) = \frac{(m+n)\sigma^2}{m+n} = \sigma^2. 
\)
\end{solution}


\begin{exercise}\label{exer:7-pareto}\label{exer:src-2015-spring-a-07}
设总体 \(X\) 的分布函数为
\[
F(x,\alpha,\beta) = \begin{cases} 1-\left(\dfrac{\alpha}{x}\right)^{\!\beta}, & x > \alpha, \\ 0, & x \leq \alpha, \end{cases}
\]
其中 \(\alpha > 0\)，\(\beta > 1\)，\(X_1, \dots, X_n\) 为样本。\par
(1) 当 \(\alpha=1\) 时，求 \(\beta\) 的矩估计和最大似然估计；\par
(2) 当 \(\beta=2\) 时，求 \(\alpha\) 的最大似然估计并讨论它的无偏性。
\end{exercise}
\begin{solution}
密度 \(f(x) = \beta\,\alpha^\beta\,x^{-(\beta+1)}\ (x > \alpha)\)。\\
(1) \(\alpha=1\) 时，\(f(x) = \beta\,x^{-(\beta+1)}\ (x > 1)\)。
\[
EX = \int_1^\infty x\cdot\beta\,x^{-(\beta+1)}\,\mathrm{d}x = \frac{\beta}{\beta-1}.
\]
矩估计：令 \(\dfrac{\beta}{\beta-1} = \overline{X}\)，得 \(\hat{\beta}_{\text{矩}} = \dfrac{\overline{X}}{\overline{X}-1}\)。\\
似然函数 \[L(\beta) = \beta^n\!\left(\prod X_i\right)^{-(\beta+1)},~~~\ln L = n\ln\beta - (\beta+1)\sum\ln X_i\]
令 \(n/\beta - \sum\ln X_i = 0\)，得 \(\hat{\beta}_{\text{MLE}} = \dfrac{n}{\sum\ln X_i}\)。\\
(2) \(\beta=2\) 时，\(f(x) = 2\alpha^2 x^{-3}\ (x > \alpha)\)。约束 \(\alpha \leq X_{(1)}\)。\\
\(\ln L = n\ln 2 + 2n\ln\alpha - 3\sum\ln X_i\)，\(\dfrac{\partial\ln L}{\partial\alpha} = \dfrac{2n}{\alpha} > 0\)，故 \(L\) 关于 \(\alpha\) 递增，\(\hat{\alpha} = X_{(1)}\)。\\
对于 Pareto(\(\alpha\),2)，\(E(X_{(1)}) = \dfrac{2n\alpha}{2n-1} \neq \alpha\)，故 \(\hat{\alpha}\) 是有偏估计。
\end{solution}


\begin{exercise}\label{exer:7-mono-3x2}
设总体 \(X\) 的概率密度为 \(f(x,\theta) = \dfrac{3x^2}{\theta^3}\ (0 \leq x \leq \theta)\)，其中 \(\theta > 0\) 未知，\(X_1, \dots, X_n\) 为简单随机样本。\par
(1) 求 \(\theta\) 的矩估计量 \(\hat{\theta}_1\) 和最大似然估计量 \(\hat{\theta}_2\)；\par
(2) 判断 \(\hat{\theta}_2\) 是否是 \(\theta\) 的无偏估计量。
\end{exercise}
\begin{solution}
(1) \(EX = \displaystyle\int_0^\theta \frac{3x^3}{\theta^3}\,\mathrm{d}x = \frac{3\theta}{4}\)。令 \(\dfrac{3\theta}{4} = \overline{X}\)，得 \(\hat{\theta}_1 = \dfrac{4}{3}\overline{X}\)。似然函数 \(L(\theta) = \dfrac{3^n\prod x_i^2}{\theta^{3n}}\) 关于 \(\theta\) 单减，故
\(
\hat{\theta}_2 = \max\{X_1, X_2, \dots, X_n\}.
\)
(2) 由于样本同分布，$P(X_i \le x) = P(X \le x) = F_X(x)$，对于样本最大值，分布函数公式为：
\begin{align*}
F_{\max}(x) &= P\left(\max\{X_1,X_2,\dots,X_n\} \leq x\right) = P\left(X_1 \leq x,\ X_2 \leq x,\ \dots,\ X_n \leq x\right)\\
&= \prod_{i=1}^n P(X_i \leq x) = \prod_{i=1}^n F_X(x) = \left[F_X(x)\right]^n
\end{align*}
求总体 $X$ 的分布函数 $F_X(x)$
$$F_X(x) =
\begin{cases}
0, & x<0 \\
\displaystyle \int_{0}^{x} \frac{3t^2}{\theta^3}\,\mathrm{d}t = \frac{x^3}{\theta^3}, & 0\leq x\leq\theta \\
1, & x>\theta
\end{cases}$$
代入最大值分布公式，当 $0\leq x\leq\theta$ 时：
\(F_{\hat{\theta}_2}(x) = \left(F_X(x)\right)^n = \left(\frac{x^3}{\theta^3}\right)^n = \left(\frac{x}{\theta}\right)^{3n}\)
当 $x<0$ 时 $F=0$，$x>\theta$ 时 $F=1$。密度是分布函数的导数：
$$f_{\hat{\theta}_2}(x) = F'_{\hat{\theta}_2}(x) =\begin{cases}\displaystyle \frac{3n x^{3n-1}}{\theta^{3n}}, & 0\leq x\leq\theta \\0, & \text{其他}\end{cases}$$
根据期望公式 $E(\hat{\theta}_2) = \int_{-\infty}^{+\infty} x f_{\hat{\theta}_2}(x)\,\mathrm{d}x$：
$$\begin{aligned}
E(\hat{\theta}_2) &= \int_{0}^{\theta} x \cdot \frac{3n x^{3n-1}}{\theta^{3n}}\,\mathrm{d}x = \frac{3n}{\theta^{3n}} \int_{0}^{\theta} x^{3n}\,\mathrm{d}x = \frac{3n}{\theta^{3n}} \cdot \left. \frac{x^{3n+1}}{3n+1} \right|_{0}^{\theta} = \frac{3n}{\theta^{3n}} \cdot \frac{\theta^{3n+1}}{3n+1} = \frac{3n}{3n+1} \theta
\end{aligned}$$
故 \(\hat{\theta}_2\) 是\textbf{有偏}估计量。
\begin{note}
作为对比，\(\hat{\theta}_1 = \dfrac{4}{3}\overline{X}\) 满足 \(E(\hat{\theta}_1) = \dfrac{4}{3}\cdot\dfrac{3\theta}{4} = \theta\)，是无偏的。
\end{note}
\end{solution}



\begin{exercise}\label{exer:7-exp4-shift}
设总体 \(X\) 的密度函数为
\[
f_X(x,\theta) = \begin{cases} 4\,e^{-4(x-\theta)}, & x \geq \theta, \\ 0, & \text{其他}, \end{cases}
\]
\(X_1, \dots, X_n\) 为来自总体 \(X\) 的样本。\par
(1) 求 \(\theta\) 的最大似然估计 \(\hat{\theta}\)；\par
(2) \(\hat{\theta}\) 是否为无偏估计？如果不是，修正为无偏估计。
\end{exercise}
\begin{solution}
(1) 构造似然函数：
$$L(\theta) = \prod_{i=1}^n 4 e^{-4(X_i-\theta)} = 4^n e^{-4\sum_{i=1}^n X_i + 4n\theta}, \quad \theta \leq X_{(1)} = \min\{X_i\}$$
$L(\theta)$ 关于 $\theta$ 严格单调递增，因此在定义域右端点取最大值，故最大似然估计量：
\(\hat{\theta} = X_{(1)} = \min\{X_1,X_2,\dots,X_n\}\)
(2) 我们先做一个变量平移：令 $Y_i = X_i - \theta$，则当 $X_i \geq \theta$ 时，$Y_i \geq 0$，代入密度函数：
  \(  f_{Y_i}(y) = f_X(y+\theta,\theta) = 4 e^{-4y}, \quad y \geq 0  \)
  这正是参数为 $\lambda=4$ 的指数分布：$Y_i \sim \text{Exp}(4)$，且 $Y_1,Y_2,\dots,Y_n$ 独立同分布。
样本最小值 $X_{(1)} = \min\{X_i\} = \min\{Y_i+\theta\} = \theta + \min\{Y_i\}$，因此：
  \(  X_{(1)} - \theta = \min\{Y_1,Y_2,\dots,Y_n\}  \)
  即：$\mathbf{X_{(1)} - \theta}$ 是 $n$ 个独立 $\text{Exp}(4)$ 变量的最小值。单个 $\text{Exp}(4)$ 的分布函数：$F_Y(y) = 1 - e^{-4y},\ y\geq0$，$n$ 个独立变量的最小值的分布函数：
  $$  F_{\min Y}(y) = P(\min Y_i \leq y) = 1 - P(\min Y_i > y) = 1 - [1-F_Y(y)]^n = 1 - e^{-4ny}  $$
  这正是参数为 $\lambda=4n$ 的指数分布：$\min Y_i \sim \text{Exp}(4n)$。根据指数分布的期望公式：若 $Z \sim \text{Exp}(\lambda)$，则 $E(Z) = \frac{1}{\lambda}$。因此：
  \(  E\left(\min Y_i\right) = \frac{1}{4n}  \)
由 $X_{(1)} = \theta + \min Y_i$，两边取期望：
\begin{align*}
E(\hat{\theta}) = E(X_{(1)}) = E\left(\theta + \min Y_i\right) = \theta + E\left(\min Y_i\right) = \theta + \frac{1}{4n}\neq \theta
\end{align*}
因此 $\hat{\theta} = X_{(1)}$ 是有偏估计量，偏差为 $\frac{1}{4n}$，我们只需要把偏差减掉，构造新的估计量：
\(
\hat{\theta}_{\text{unbiased}} = X_{(1)} - \frac{1}{4n}
\)
验证无偏性：
$$
E\left(\hat{\theta}_{\text{unbiased}}\right) = E\left(X_{(1)} - \frac{1}{4n}\right) = \left(\theta + \frac{1}{4n}\right) - \frac{1}{4n} = \theta
$$
修正后的估计量是 $\theta$ 的无偏估计。
\end{solution}
\begin{note}
题目中的 $f_X(x,\theta)$ 是位置参数型指数分布（也叫平移指数分布）：标准指数分布 $\text{Exp}(\lambda)$ 是 $f(x) = \lambda e^{-\lambda x},\ x\geq0$；平移 $\theta$ 后，分布变为 $f(x) = \lambda e^{-\lambda(x-\theta)},\ x\geq\theta$，其中 $\theta$ 是位置参数（分布的左端点）。在这类“下端点未知、似然函数随 \(\theta\) 单调变化”的典型边界参数模型中，极大似然估计往往落在样本最小值处；本题正是这一类代表情形。
\end{note}

\begin{proposition}[样本方差的无偏性]\label{prop:sample-variance-unbiased}
    样本方差 $S^2$ 是总体方差 $\sigma^2$ 的无偏估计，这个结论对“任意”具有有限方差的总体分布都成立。
\end{proposition}
\begin{note}
这条结论的使用范围非常广：\textbf{不要求总体正态}，只要求总体方差存在即可，因此它比卡方分布结论更稳。不要误把它和“\((n-1)S^2/\sigma^2\sim\chi^2(n-1)\)”混为一谈；后者必须建立在总体正态的前提下。
\end{note}
\begin{proof}
首先给出样本方差的定义
\(
S^2=\frac{1}{n-1}\sum_{i=1}^n (X_i-\overline{X})^2.
\)
将 \(X_i-\overline{X}\) 改写为 \((X_i-\mu)-(\overline{X}-\mu)\)，则
\[
\sum_{i=1}^n (X_i-\overline{X})^2
=\sum_{i=1}^n (X_i-\mu)^2-n(\overline{X}-\mu)^2.
\]
于是取期望得
\[
E\!\left[\sum_{i=1}^n (X_i-\overline{X})^2\right]
=\sum_{i=1}^n E(X_i-\mu)^2-nE(\overline{X}-\mu)^2.
\]
由于 \(E(X_i-\mu)^2=\sigma^2\)，且
\(
E(\overline{X}-\mu)^2=D(\overline{X})=\frac{\sigma^2}{n},
\)
所以
\[
E\!\left[\sum_{i=1}^n (X_i-\overline{X})^2\right]
=n\sigma^2-n\cdot\frac{\sigma^2}{n}=(n-1)\sigma^2.
\]
因此
\(
E(S^2)=\frac{1}{n-1}E\!\left[\sum_{i=1}^n (X_i-\overline{X})^2\right]=\sigma^2.
\)
\end{proof}

\begin{exercise}\label{exer:7-sigma2-unbiased}
设 \(X_1, \dots, X_n\) 和 \(Y_1, \dots, Y_m\) 分别是来自总体 \(X \sim N(\mu_1, \sigma^2)\) 和 \(Y \sim N(\mu_2, \sigma^2)\) 的简单随机样本，样本方差分别为 \(S_X^2, S_Y^2\)。则 \(\sigma^2\) 的无偏估计是（\quad）。\par
A. \(S_X^2 + S_Y^2\)\quad B. \((m-1)S_X^2 + (n-1)S_Y^2\)\quad C. \(\dfrac{S_X^2+S_Y^2}{m+n-2}\)\quad D. \(\dfrac{(m-1)S_X^2+(n-1)S_Y^2}{m+n-2}\)
\end{exercise}
\begin{solution}
应选 D。

判断哪个是 \(\sigma^2\) 的无偏估计，最直接的方法就是分别求期望。

由于样本方差对总体方差都是无偏的，所以
\(
E(S_X^2)=\sigma^2,\qquad E(S_Y^2)=\sigma^2.
\)

于是各选项对应的期望分别为
\(
E(S_X^2+S_Y^2)=2\sigma^2,
\)
\(
E\!\bigl((m-1)S_X^2+(n-1)S_Y^2\bigr)=(m+n-2)\sigma^2,
\)
\(
E\!\left(\frac{S_X^2+S_Y^2}{m+n-2}\right)=\frac{2\sigma^2}{m+n-2},
\)
\(
E\!\left(\frac{(m-1)S_X^2+(n-1)S_Y^2}{m+n-2}\right)
=\frac{(m-1)\sigma^2+(n-1)\sigma^2}{m+n-2}.
\)

最后一个式子化简为
\(
E\!\left(\frac{(m-1)S_X^2+(n-1)S_Y^2}{m+n-2}\right) = \frac{(m-1)\sigma^2+(n-1)\sigma^2}{m+n-2} = \sigma^2.
\)

因此只有 D 的期望恰好等于 \(\sigma^2\)，故 D 是 \(\sigma^2\) 的无偏估计。

这其实就是“无偏估计量做加权平均仍可保持无偏”的一个简单例子。
\end{solution}


\begin{proposition}[泊松分布无偏估计的不唯一性]\label{prop:poisson-unbiased}
    设 $X_1, X_2, \dots, X_n$ 是来自泊松总体 $P(\lambda)$ 的简单随机样本，那么样本均值 $\overline{X}$ 和样本方差 $S^2$ 都是未知参数 $\lambda$ 的无偏估计。
\end{proposition}
\begin{note}
这条命题的重点不只是“会算”，更在于提醒我们：\textbf{无偏估计一般不唯一}。以后做选择题时，如果题目只问“哪个是无偏估计”，不要先入为主地认为答案只能有样本均值这一种；但若进一步比较优劣，还需要看方差或均方误差，不能只看无偏性。
\end{note}
\begin{proof}
泊松总体 \(P(\lambda)\) 满足
\(
E(X)=\lambda,\qquad D(X)=\lambda.
\)
因此对样本均值 \(\overline{X}\)，有
\(
E(\overline{X})=\frac{1}{n}\sum_{i=1}^n E(X_i)=\lambda,
\)
故 \(\overline{X}\) 是 \(\lambda\) 的无偏估计量。

另一方面，由前述“样本方差的无偏性”，可知
\(
E(S^2)=D(X)=\lambda.
\)
所以 \(S^2\) 也是 \(\lambda\) 的无偏估计量。

这说明：对同一个参数 \(\lambda\)，无偏估计量并不唯一。
\end{proof}

\begin{exercise}
    设 $X_1,X_2,\dots,X_n$ 是取自总体 $P(\lambda)$ 的样本，$\overline{X}$ 为样本均值，则 $\lambda^2$ 的一个无偏估计为（\quad）。\\
A. $\overline{X}^2 - \overline{X}$\qquad B. $\overline{X}^2 + \overline{X}$\qquad C. $\overline{X}^2 - \dfrac{1}{n}\overline{X}$\qquad D. $\overline{X}^2 + \dfrac{1}{n}\overline{X}$
\end{exercise}
\begin{solution}
应选 C。

要求的是 \(\lambda^2\) 的无偏估计，因此关键是先把 \(E(\overline{X}^2)\) 算出来。

因为总体服从泊松分布 \(P(\lambda)\)，所以
\(
EX=\lambda,\qquad DX=\lambda.
\)
从而样本均值满足
\[
E(\overline{X})=\lambda,\qquad D(\overline{X})=\frac{\lambda}{n}.
\]

利用恒等式 \(E(T^2)=DT+[E(T)]^2\)，取 \(T=\overline{X}\)，得
\[
E(\overline{X}^2)=D(\overline{X})+[E(\overline{X})]^2
=\frac{\lambda}{n}+\lambda^2.
\]

现在只需把多出来的 \(\dfrac{\lambda}{n}\) 消掉。又因为
\(
E\!\left(\frac{1}{n}\overline{X}\right)=\frac{1}{n}E(\overline{X})=\frac{\lambda}{n},
\)
故
\[
E\!\left(\overline{X}^2-\frac{1}{n}\overline{X}\right)
=\left(\frac{\lambda}{n}+\lambda^2\right)-\frac{\lambda}{n}
=\lambda^2.
\]

因此
\(
\overline{X}^2-\frac{1}{n}\overline{X}
\)
是 \(\lambda^2\) 的无偏估计，故应选 C。
\end{solution}

\begin{proposition}[均匀分布参数的无偏估计]\label{prop:uniform-unbiased}
    设 $X_1, X_2, \dots, X_n$ 是来自均匀总体 $U(0, \theta)$ 的简单随机样本，次序统计量最大值 \(X_{(n)} = \max(X_1, \dots, X_n)\)是 $\theta$ 的极大似然估计量，但它是有偏的。经过修正后的统计量 $T = \frac{n+1}{n}X_{(n)}$ 才是 $\theta$ 的无偏估计。
\end{proposition}
\begin{note}
这个结论适用于“支撑区间上端点未知”的均匀分布模型。直接把 \(X_{(n)}\) 当作无偏估计会系统性偏小，因为样本最大值通常达不到真实上界 \(\theta\)。考试里一旦看到 \(U(0,\theta)\) 或类似边界参数模型，就要警惕“MLE 往往在边界上，但常常有偏”。
\end{note}
\begin{proof}
    首先求出最大值 $X_{(n)}$ 的分布。因为样本相互独立，所以 $X_{(n)} \le x$ 等价于所有的 $X_i \le x$：$$F_{X_{(n)}}(x) = P(X_{(n)} \le x) = \prod_{i=1}^n P(X_i \le x) = \left(\frac{x}{\theta}\right)^n \quad (0 < x < \theta)$$对其求导，得到 $X_{(n)}$ 的概率密度函数：\(f_{X_{(n)}}(x) = \frac{n}{\theta^n}x^{n-1} \quad (0 < x < \theta)\)接下来求 $X_{(n)}$ 的数学期望：$$E(X_{(n)}) = \int_0^\theta x \cdot \frac{n}{\theta^n}x^{n-1}\,\mathrm{d}x = \frac{n}{\theta^n} \int_0^\theta x^n\,\mathrm{d}x = \frac{n}{\theta^n} \cdot \frac{\theta^{n+1}}{n+1} = \frac{n}{n+1}\theta$$由于 $E(X_{(n)}) \ne \theta$，所以极大似然估计量 $X_{(n)}$ 是有偏的（它总是倾向于低估总体的上限 $\theta$）。
    为了消除偏差，我们在等式两边同乘 $\frac{n+1}{n}$：\(E\left(\frac{n+1}{n}X_{(n)}\right) = \frac{n+1}{n} \cdot \frac{n}{n+1}\theta = \theta\)因此，修正后的估计量 $T = \frac{n+1}{n}X_{(n)}$ 是总体参数 $\theta$ 的无偏估计量。
\end{proof}


\begin{definition}[有效性] \label{def:efficiency}
若 \(\hat{\theta}_1, \hat{\theta}_2\) 都是 \(\theta\) 的无偏估计量，且 \(D\hat{\theta}_1 < D\hat{\theta}_2\)，则称 \(\hat{\theta}_1\) 比 \(\hat{\theta}_2\) \textbf{有效}。直观上说，在两者都“不偏”的前提下，谁波动更小，谁就更值得优先采用。
\end{definition}

\begin{note}
有效性比较有一个前提：\textbf{只能在无偏估计量之间比较}。如果某个估计量本身有偏，就不能直接拿“方差更小”来宣布它更有效，而应转去比较均方误差。很多同学会把“有效”误理解成“任何意义下都更好”，这是不准确的。
\end{note}





\begin{exercise}\label{exer:7-efficiency}
设总体 \(X \sim N(\mu, \sigma^2)\)，\(X_1, X_2, X_3\) 是来自总体 \(X\) 的简单随机样本，记
\[
\hat{\mu}_1 = \frac{1}{5}X_1 + \frac{3}{10}X_2 + \frac{1}{2}X_3,\quad
\hat{\mu}_2 = \frac{1}{3}X_1 + \frac{1}{4}X_2 + \frac{5}{12}X_3,\quad
\hat{\mu}_3 = \frac{1}{3}X_1 + \frac{1}{6}X_2 + \frac{1}{2}X_3.
\]
证明 \(\hat{\mu}_1, \hat{\mu}_2, \hat{\mu}_3\) 都是 \(\mu\) 的无偏估计，并指出哪一个最有效。
\end{exercise}

\begin{solution}
先验证无偏性。线性估计量是否对 \(\mu\) 无偏，只需看各系数之和是否为 1。

对 \(\hat{\mu}_1\)：
\(
E\hat{\mu}_1
=\left(\frac15+\frac{3}{10}+\frac12\right)\mu
=\mu.
\)

对 \(\hat{\mu}_2\)：
\(
E\hat{\mu}_2
=\left(\frac13+\frac14+\frac{5}{12}\right)\mu
=\mu.
\)

对 \(\hat{\mu}_3\)：
\[
E\hat{\mu}_3
=\left(\frac13+\frac16+\frac12\right)\mu
=\mu.
\]

所以 \(\hat{\mu}_1,\hat{\mu}_2,\hat{\mu}_3\) 都是 \(\mu\) 的无偏估计。

下面比较谁最有效。由于 \(X_1,X_2,X_3\) 相互独立，且方差都等于 \(\sigma^2\)，所以各线性组合的方差等于“系数平方和乘 \(\sigma^2\)”。

\[
D\hat{\mu}_1
=\left[\left(\frac15\right)^2+\left(\frac{3}{10}\right)^2+\left(\frac12\right)^2\right]\sigma^2
=\left(\frac{1}{25}+\frac{9}{100}+\frac{1}{4}\right)\sigma^2
=\frac{38}{100}\sigma^2,
\]
\[
D\hat{\mu}_2
=\left[\left(\frac13\right)^2+\left(\frac14\right)^2+\left(\frac{5}{12}\right)^2\right]\sigma^2
=\left(\frac{1}{9}+\frac{1}{16}+\frac{25}{144}\right)\sigma^2
=\frac{50}{144}\sigma^2,
\]
\[
D\hat{\mu}_3
=\left[\left(\frac13\right)^2+\left(\frac16\right)^2+\left(\frac12\right)^2\right]\sigma^2
=\left(\frac{1}{9}+\frac{1}{36}+\frac{1}{4}\right)\sigma^2
=\frac{14}{36}\sigma^2.
\]

把它们化成小数：
\[
\frac{38}{100}=0.380,\qquad
\frac{50}{144}\approx 0.347,\qquad
\frac{14}{36}\approx 0.389.
\]
所以
\(
\frac{50}{144}<\frac{38}{100}<\frac{14}{36}.
\)

因此 \(\hat{\mu}_2\) 的方差最小，所以 \(\hat{\mu}_2\) 最有效。
\end{solution}

\begin{exercise}\label{exer:7-eff-N01}
设总体 \(X \sim N(0,1)\)，\(X_1, X_2\) 是从总体抽取的一个样本，下列四个无偏估计量中最有效的是（\quad）。\par
A. \(\mu_1 = X_1\)\qquad B. \(\mu_2 = \dfrac{2}{3}X_1+\dfrac{1}{3}X_2\)\qquad C. \(\mu_3 = \dfrac{1}{4}X_1+\dfrac{3}{4}X_2\)\qquad D. \(\mu_4 = \dfrac{1}{2}X_1+\dfrac{1}{2}X_2\)
\end{exercise}
\begin{solution}
应选 D。

题目已经说明这四个估计量都是无偏的，因此比较“最有效”只需比较方差大小。

又因为总体 \(X\sim N(0,1)\)，所以 \(DX_1=DX_2=1\)，
且 \(X_1,X_2\) 相互独立。于是线性组合的方差等于各项系数平方乘方差后相加。

对 A：\(D\mu_1=DX_1=1\)。

对 B：
\[
D\mu_2=D\!\left(\frac23X_1+\frac13X_2\right)
=\left(\frac23\right)^2+\left(\frac13\right)^2
=\frac59.
\]

对 C：
\(
D\mu_3=\left(\frac14\right)^2+\left(\frac34\right)^2
=\frac{1}{16}+\frac{9}{16}
=\frac58.
\)

对 D：
\(
D\mu_4=\left(\frac12\right)^2+\left(\frac12\right)^2=\frac12.
\)

比较可得
\(
\frac12<\frac59<\frac58<1.
\)
所以 \(\mu_4\) 的方差最小，也就是最有效。

故应选 D。
\end{solution}


\begin{exercise}\label{exer:7-u0-compare}
设样本 \(X_1, X_2, X_3\) 来自总体 \(U(0,\theta)\)，则估计量 \(\hat{\theta}_1=2\overline{X}\)，\(\hat{\theta}_2=X_1+X_2\)，\(\hat{\theta}_3=\dfrac{2}{5}X_1+\dfrac{4}{5}(X_2+X_3)\) 中，哪一个较优？
\end{exercise}
\begin{solution}
先判断三个估计量是否都能作为 \(\theta\) 的无偏估计量来比较。对总体 \(U(0,\theta)\)，有
\(
EX=\frac{\theta}{2},\qquad DX=\frac{\theta^2}{12}.
\)
\begin{note}
对 \(U(0,\theta)\) 来说，\(EX=\theta/2\)，所以线性估计量 \(a_1X_1+a_2X_2+a_3X_3\) 无偏的条件不是“系数和为 1”，而是 \(a_1+a_2+a_3=2\)。这正是这道题最容易误判的地方。
\end{note}
因此
\(
E(\hat{\theta}_1)=E(2\overline{X})=2EX=\theta,
\)
\(
E(\hat{\theta}_2)=E(X_1+X_2)=2EX=\theta,
\)
\[
E(\hat{\theta}_3)=\left(\frac25+\frac45+\frac45\right)EX=2EX=\theta.
\]
所以三者都对 \(\theta\) 无偏，可以继续比较方差。
\[
D\hat{\theta}_1 = 4D(\overline{X}) = \frac{4\theta^2}{12\cdot3} = \frac{\theta^2}{9},
\quad
D\hat{\theta}_2 = 2\cdot\frac{\theta^2}{12} = \frac{\theta^2}{6},
\]
\[
D\hat{\theta}_3 = \left(\frac{4}{25}+\frac{16}{25}+\frac{16}{25}\right)\frac{\theta^2}{12} = \frac{36\theta^2}{300} = \frac{3\theta^2}{25}.
\]
\(
\frac{\theta^2}{9}<\frac{3\theta^2}{25}<\frac{\theta^2}{6}.
\)
故 \(\hat{\theta}_1=2\overline{X}\) 的方差最小，因此它在这三个无偏估计量中最有效，也就是较优的估计量。
\end{solution}


\begin{exercise}\label{exer:src-2009-3a-07}
设 \(X_1,X_2,\dots,X_n\) 为总体的一个样本，\(X\) 服从均匀分布：
\[
f(x;\theta)=
\begin{cases}
\dfrac1\theta, & 0<x<\theta,\\
0, & \text{其它}.
\end{cases}
\]
求参数 \(\theta\) 的极大似然估计 \(\hat\theta_1\)；求参数 \(\theta\) 的矩估计 \(\hat\theta_2\)；讨论 \(\hat\theta_1,\hat\theta_2\) 的无偏性。
\end{exercise}
\begin{solution}
似然函数为
\[
L(\theta)=\prod_{i=1}^n \frac1\theta \mathbf 1_{(0,\theta)}(X_i)
=\theta^{-n}\mathbf 1_{\{\theta\ge X_{(n)}\}},
\]
其中 \(X_{(n)}=\max\{X_1,\dots,X_n\}\)。因为 \(L(\theta)\) 在 \([\;X_{(n)},+\infty)\) 上单调递减，故
\(
\hat\theta_1=X_{(n)}.
\)
又因 \(EX=\theta/2\)，矩估计满足 \(\overline X=\theta/2\)，故
\(
\hat\theta_2=2\overline X.
\)
关于无偏性，
\(
E(\hat\theta_2)=2E(\overline X)=2EX=\theta,
\)
所以 \(\hat\theta_2\) 是无偏估计。

而
\(
E(\hat\theta_1)=E(X_{(n)})=\frac{n}{n+1}\theta,
\)
故 \(\hat\theta_1\) 有偏。
\end{solution}


\begin{definition}[一致性（相合性）] \label{def:consistent}
若 \(\hat{\theta}\xrightarrow{P}\theta\)（即对任意 \(\varepsilon > 0\)，\(\lim_{n\to\infty}P\{|\hat{\theta}-\theta| < \varepsilon\} = 1\)），则称 \(\hat{\theta}\) 为 \(\theta\) 的\textbf{一致估计量}（或\textbf{相合估计量}）。这表示样本量越来越大时，估计量偏离真值的概率会越来越小。
\end{definition}

\begin{note}
验证一致性常用的两种方法：
\begin{enumerate}
    \item \textbf{切比雪夫不等式}：若 \(D\hat{\theta} \to 0\)（当 \(n\to\infty\)），则 \(P\{|\hat{\theta}-\theta| \geq \varepsilon\} \leq D\hat{\theta}/\varepsilon^2 \to 0\)。
    \item \textbf{辛钦大数定律}：若 \(\hat{\theta}\) 是样本值的函数且可表示为独立同分布随机变量的均值，则可直接得出依概率收敛。
\end{enumerate}
\textbf{不要混淆无偏与一致}：无偏看的是每个固定 \(n\) 下期望是否等于真值；一致看的是 \(n\to\infty\) 时估计量本身是否逼近真值。二者彼此都不能自动推出对方。
\end{note}


\begin{exercise}\label{exer:7-laplace}\label{exer:src-2013-a-08-mle}\label{exer:src-probstat-test1-07}
设总体 \(X\) 的概率密度为 \(f(x) = \dfrac{1}{2\theta}\,e^{-|x|/\theta}\ (-\infty < x < \infty)\)，其中 \(\theta > 0\) 未知，\(X_1, \dots, X_n\) 为样本。\par
(1) 求 \(\theta\) 的矩估计量 \(\hat{\theta}_M\)；\par
(2) 求 \(\theta\) 的最大似然估计量 \(\hat{\theta}_L\)；\par
(3) 证明 \(\hat{\theta}_L\) 是 \(\theta\) 的无偏估计量，且 \(\hat{\theta}_L\) 是 \(\theta\) 的一致估计量。
\end{exercise}
\begin{solution}
(1) 由于该分布关于 \(0\) 对称，\(EX=0\)，不能用一阶矩解出参数。改用二阶原点矩：
\[
EX^2
=\frac{1}{2\theta}\int_{-\infty}^{+\infty}x^2e^{-|x|/\theta}\,\mathrm{d}x
=\frac1\theta\int_0^\infty x^2e^{-x/\theta}\,\mathrm{d}x.
\]
令 \(t=x/\theta\)，得
\[
EX^2
=\theta^2\int_0^\infty t^2e^{-t}\,\mathrm{d}t
=2\theta^2.
\]
令 \(EX^2\) 等于样本二阶原点矩，得到
\(
2\theta^2=\frac1n\sum_{i=1}^{n}X_i^2,
\)
故
\(
\hat{\theta}_M=\sqrt{\frac{1}{2n}\sum_{i=1}^{n}X_i^2}.
\)

(2) 由样本独立性，似然函数为
\[
L(\theta)=\prod_{i=1}^{n}\frac{1}{2\theta}e^{-|X_i|/\theta}
=(2\theta)^{-n}e^{-\sum|X_i|/\theta}.
\]
取对数得
\[
\ln L(\theta)=-n\ln(2\theta)-\frac{1}{\theta}\sum_{i=1}^{n}|X_i|.
\]
对 \(\theta\) 求导并令其为 0：
\[
\frac{\mathrm d}{\mathrm d\theta}\ln L(\theta)
=-\frac{n}{\theta}+\frac{1}{\theta^2}\sum_{i=1}^{n}|X_i|=0,
\]
故
\(
\hat{\theta}_L=\frac{1}{n}\sum_{i=1}^{n}|X_i|.
\)

(3) 为证明无偏性，先计算单个样本的 \(E|X|\)：
\[
E|X|
=\frac{1}{2\theta}\int_{-\infty}^{+\infty}|x|e^{-|x|/\theta}\,\mathrm{d}x
=\frac{1}{\theta}\int_0^\infty xe^{-x/\theta}\,\mathrm{d}x.
\]
令 \(t=x/\theta\)，则 \(x=\theta t,\ \mathrm dx=\theta\,\mathrm dt\)，于是
\[
E|X|=\theta\int_0^\infty te^{-t}\,\mathrm{d}t=\theta\,\Gamma(2)=\theta.
\]
因此
\[
E(\hat{\theta}_L)=E\!\left(\frac1n\sum_{i=1}^{n}|X_i|\right)=\frac1n\sum_{i=1}^{n}E|X_i|=\theta,
\]
所以 \(\hat{\theta}_L\) 是 \(\theta\) 的无偏估计量。

由于 \(|X_1|,|X_2|,\dots,|X_n|\) 独立同分布，且 \(E|X_i|=\theta\)，由辛钦大数定律：
\[
\hat{\theta}_L = \frac{1}{n}\sum|X_i| \xrightarrow{P} E|X| = \theta. 
\]
故 \(\hat{\theta}_L\) 是 \(\theta\) 的一致估计量。
\end{solution}


\begin{exercise}\label{exer:src-2013-a-08-mme}
设总体 \(X\) 的概率密度为
\[
f(x;\theta)=\frac{1}{2\theta}e^{-|x|/\theta},\qquad -\infty<x<+\infty,\quad \theta>0.
\]
\(X_1,\dots,X_n\) 为来自总体 \(X\) 的样本。求参数 \(\theta\) 的矩估计量。
\end{exercise}
\begin{solution}
由于该分布关于 \(0\) 对称，一阶矩为
\(
EX=0,
\)
不能用来解出 \(\theta\)。改用二阶原点矩：
\[
EX^2
=\frac{1}{2\theta}\int_{-\infty}^{+\infty}x^2e^{-|x|/\theta}\,\mathrm{d}x
=\frac{1}{\theta}\int_0^\infty x^2e^{-x/\theta}\,\mathrm{d}x.
\]
令 \(t=x/\theta\)，则
\[
EX^2
=\frac{1}{\theta}\int_0^\infty \theta^2t^2e^{-t}\theta\,\mathrm{d}t
=\theta^2\Gamma(3)=2\theta^2.
\]
令总体二阶矩等于样本二阶矩，
\(
2\theta^2=\frac1n\sum_{i=1}^n X_i^2,
\)
得到矩估计量
\(
\hat\theta_M=\sqrt{\frac{1}{2n}\sum_{i=1}^n X_i^2}.
\)
\end{solution}

\begin{exercise}\label{exer:7-consistent}
设 \(X_1, X_2, \dots, X_n\) 是来自总体 \(X\) 的简单随机样本，\(EX = \mu\)，\(DX = \sigma^2\)。证明统计量
\(
Y = \frac{2}{n(n+1)}\sum_{i=1}^{n}iX_i
\)
是 \(\mu\) 的无偏相合估计量。
\end{exercise}
\begin{solution}
\textbf{无偏性}：
\(
EY = \frac{2}{n(n+1)}\sum_{i=1}^{n}i\,EX_i = \frac{2\mu}{n(n+1)}\cdot\frac{n(n+1)}{2} = \mu.
\)
\textbf{一致性}：由 \(\sum i^2 = \dfrac{n(n+1)(2n+1)}{6}\)，
\[
DY = \left[\frac{2}{n(n+1)}\right]^{\!2}\sigma^2\sum_{i=1}^{n}i^2 = \frac{4\sigma^2}{n^2(n+1)^2}\cdot\frac{n(n+1)(2n+1)}{6} = \frac{2(2n+1)}{3n(n+1)}\,\sigma^2.
\]
由切比雪夫不等式，对任意 \(\varepsilon > 0\)：
\[
P\{|Y-\mu| < \varepsilon\} \geq 1 - \frac{DY}{\varepsilon^2} = 1 - \frac{2(2n+1)\sigma^2}{3n(n+1)\varepsilon^2}.
\]
令 \(n\to\infty\)，右端 \(\to 1\)，故 \(Y\xrightarrow{P}\mu\)，\(Y\) 是 \(\mu\) 的相合估计量。
\end{solution}



\begin{definition}[均方误差] \label{def:mse}
设 \(\hat{\theta}\) 为参数 \(\theta\) 的估计量，称
\[
\mathrm{MSE}(\hat{\theta}) = E(\hat{\theta} - \theta)^2
\]
为 \(\hat{\theta}\) 关于 \(\theta\) 的\textbf{均方误差}（Mean Squared Error, MSE）。
\end{definition}

\begin{note}
均方误差把“偏差”和“波动”合在一起衡量，因此它适合比较\textbf{不一定无偏}的估计量。实际使用时，一个估计量即使略有偏差，只要方差降得足够多，也可能拥有更小的均方误差。这正是“有偏但更稳定”的估计量依然有价值的原因。
\end{note}

\begin{note}
\textbf{使用提醒：} 比较估计量时可以按这个顺序想：若题目强调“无偏”，先筛掉有偏候选量，再比较方差看有效性；若候选量本身可能有偏，就转去比较均方误差；若题目关心“大样本下会不会越来越贴近真值”，再看一致性。把这些评价标准放回各自的使用场景，才不会出现“明明在比 MSE，却只盯着无偏性”的误判。
\end{note}

\begin{property}[均方误差的分解] \label{prop:mse-decomp}
均方误差可分解为偏差的平方与方差之和：
\[
\mathrm{MSE}(\hat{\theta}) = [E(\hat{\theta}) - \theta]^2 + D(\hat{\theta}) = \mathrm{Bias}^2(\hat{\theta}) + D(\hat{\theta}).
\]
\begin{proof}
\[
E(\hat{\theta}-\theta)^2 = E\!\left[(\hat{\theta}-E\hat{\theta})+(E\hat{\theta}-\theta)\right]^2 = D(\hat{\theta}) + [E(\hat{\theta})-\theta]^2.
\]
交叉项 \(2E[(\hat{\theta}-E\hat{\theta})(E\hat{\theta}-\theta)] = 0\)。
\end{proof}
\end{property}

\begin{note}
均方误差综合考虑了估计量的\textbf{偏差}和\textbf{方差}：
\begin{itemize}
    \item 若 \(\hat{\theta}\) 是无偏估计（\(E\hat{\theta}=\theta\)），则 \(\mathrm{MSE}(\hat{\theta}) = D(\hat{\theta})\)，此时均方误差最小等价于方差最小（即有效性）。
    \item 若 \(\hat{\theta}_1, \hat{\theta}_2\) 为 \(\theta\) 的估计量（不一定无偏），且 \(\mathrm{MSE}(\hat{\theta}_1) \leq \mathrm{MSE}(\hat{\theta}_2)\)，则称 \(\hat{\theta}_1\) 在均方误差意义下比 \(\hat{\theta}_2\) 有效。
    \item 一个有偏但方差很小的估计量，其均方误差可能\textbf{优于}无偏但方差较大的估计量。
\end{itemize}
\end{note}


\begin{exercise}\label{exer:7-shifted-exp}
设总体 \(X\) 的概率密度为
\[
f(x) = \begin{cases} e^{\theta-x}, & x > \theta, \\ 0, & x \leq \theta, \end{cases}
\]
\(\theta\) 未知，\(X_1, \dots, X_n\) 为样本。\par
(1) 求 \(\theta\) 的矩估计量 \(\hat{\theta}_M\) 和最大似然估计量 \(\hat{\theta}_L\)；\par
(2) 比较 \(\hat{\theta}_M\) 与 \(\hat{\theta}_L\) 的优劣。
\end{exercise}
\begin{solution}
(1) \(EX = \displaystyle\int_\theta^\infty x\,e^{\theta-x}\,\mathrm{d}x = \theta + 1\)。令 \(\theta+1 = \overline{X}\)，得 \(\hat{\theta}_M = \overline{X}-1\)。
\(L(\theta) = e^{n\theta-\sum X_i}\)，关于 \(\theta\) 递增（\(\theta \leq X_{(1)}\)），故 \(\hat{\theta}_L = X_{(1)}\)。\\
(2) \(\hat{\theta}_M\) 无偏。\(DX = E(X^2)-(EX)^2 = \displaystyle\int_\theta^\infty(x-\theta-1)^2e^{\theta-x}\,\mathrm{d}x\)。令 \(t=x-\theta\)：\(\displaystyle\int_0^\infty(t-1)^2e^{-t}\,\mathrm{d}t = \Gamma(3)-2\Gamma(2)+\Gamma(1) = 2-2+1 = 1\)，故 \(D(\hat{\theta}_M) = D(\overline{X}) = 1/n\)，\(\mathrm{MSE}_M = 1/n\)。\\
\(X_{(1)}-\theta\) 为 \(n\) 个 \(\mathrm{Exp}(1)\) 的最小值，\(E(X_{(1)}-\theta) = 1/n\)，\(D(X_{(1)}-\theta) = 1/n^2\)。故 \(E(\hat{\theta}_L) = \theta + 1/n\)，\(D(\hat{\theta}_L) = 1/n^2\)，
\(
\mathrm{MSE}_L = \left(\frac{1}{n}\right)^{\!2} + \frac{1}{n^2} = \frac{2}{n^2}.
\)
当 \(n > 2\) 时，\(2/n^2 < 1/n\)，\(\hat{\theta}_L\) 的均方误差更小，更优。
\end{solution}

\begin{exercise}\label{exer:src-2014-spring-a-08}
设总体 \(X\) 的概率密度为
\[
f(x;\theta)=
\begin{cases}
2e^{-2(x-\theta)}, & x>\theta,\\
0, & x\le\theta,
\end{cases}
\qquad \theta>0.
\]
从该总体中抽取简单随机样本 \(X_1,\dots,X_n\)，令 \(\hat\theta=\min\{X_1,\dots,X_n\}\)。求：\par
(1) \(X\) 的分布函数 \(F(x)\)；\par
(2) \(\hat\theta\) 的分布函数；\par
(3) 若用 \(\hat\theta\) 估计 \(\theta\)，判断其是否无偏；\par
(4) \(D\hat\theta\)。
\end{exercise}
\begin{solution}
(1) 当 \(x\le\theta\) 时，\(F(x)=0\)。当 \(x>\theta\) 时，
\[
F(x)=\int_{\theta}^{x}2e^{-2(t-\theta)}\,\mathrm dt
=1-e^{-2(x-\theta)}.
\]
故
\[
F(x)=
\begin{cases}
0, & x\le\theta,\\
1-e^{-2(x-\theta)}, & x>\theta.
\end{cases}
\]

(2) 记 \(X_{(1)}=\min\{X_1,\dots,X_n\}\)。当 \(x\le\theta\) 时，\(P\{X_{(1)}\le x\}=0\)；当 \(x>\theta\) 时，
\[
\begin{aligned}
F_{X_{(1)}}(x)
&=P\{X_{(1)}\le x\}\\
&=1-P\{X_1>x,\dots,X_n>x\}\\
&=1-[1-F(x)]^n
=1-e^{-2n(x-\theta)}.
\end{aligned}
\]
因此
\[
F_{\hat\theta}(x)=
\begin{cases}
0, & x\le\theta,\\
1-e^{-2n(x-\theta)}, & x>\theta.
\end{cases}
\]

(3) 由上式可知 \(\hat\theta-\theta\sim E(2n)\)，于是
\(
E\hat\theta=\theta+\frac1{2n}.
\)
故 \(\hat\theta\) 不是 \(\theta\) 的无偏估计量，而是有偏估计量。

(4) 由于 \(\hat\theta-\theta\sim E(2n)\)，指数分布方差为参数平方的倒数，故
\[
D\hat\theta=D(\hat\theta-\theta)=\frac1{(2n)^2}=\frac1{4n^2}.
\]
\end{solution}

\begin{note}
\textbf{自学检查点——点估计：}
\begin{enumerate}
    \item 矩估计的核心思想是什么？（答：用样本矩代替总体矩，解方程）
    \item MLE 的四步是什么？（答：写似然函数 $\to$ 取对数 $\to$ 求导令为零 $\to$ 解方程）
    \item 似然函数单调递增时，MLE 怎么取？（答：取参数允许范围的上界，通常是 $X_{(n)}$ 或 $X_{(1)}$）
    \item 无偏性的定义？（答：$E(\hat\theta)=\theta$）
    \item $S^2$ 是 $\sigma^2$ 的无偏估计，那 $S$ 是 $\sigma$ 的无偏估计吗？（答：不是！$E(S)\neq\sigma$，因为期望和开根号不能交换）
\end{enumerate}
\end{note}


\begin{note}
\textbf{【高频题型模板】MLE 完整解题流程：}
\begin{enumerate}
    \item \textbf{写似然函数：} $L(\theta)=\displaystyle\prod_{i=1}^n f(x_i;\theta)$（连续型）或 $\prod P(X=x_i;\theta)$（离散型）。
    \item \textbf{取对数：} $\ln L(\theta)=\sum_{i=1}^n \ln f(x_i;\theta)$。
    \item \textbf{判断类型：}
    \begin{itemize}
        \item 若 $\ln L$ 对 $\theta$ 可导且有驻点 $\Rightarrow$ 令 $\dfrac{\mathrm{d}\ln L}{\mathrm{d}\theta}=0$ 解出 $\hat\theta$；
        \item 若 $\ln L$ 单调（支撑含参数）$\Rightarrow$ 在约束边界取极值（递减取 max，递增取 min）。
    \end{itemize}
    \item \textbf{验证：} 确认是极大值（二阶导 $<0$ 或单调性判断）。
\end{enumerate}
\textbf{常见追问：} 求出 $\hat\theta$ 后，题目可能继续问 $g(\theta)$ 的 MLE $\Rightarrow$ 直接用不变性：$\widehat{g(\theta)}=g(\hat\theta)$。
\end{note}

\subsection*{历年试题补充：2008-01 3学分 A 卷}

\begin{exercise}\label{exer:src-2008-01-3a-iii-07}\label{exer:src-paper2-fill-07}\label{prob:p2-fill-07}
设 \(X\) 的分布律为
\[
\begin{array}{c|ccc}
X & 1 & 2 & 3\\
\hline
P & \theta^2 & 2\theta(1-\theta) & (1-\theta)^2
\end{array}
\]
已知一个样本值
\(
(x_1,x_2,x_3)=(1,2,1),
\)
求参数 \(\theta\) 的极大似然估计值。
\end{exercise}
\begin{solution}
由样本值可知，\(X=1\) 出现 \(2\) 次，\(X=2\) 出现 \(1\) 次，\(X=3\) 出现 \(0\) 次。似然函数为
\[
L(\theta)=(\theta^2)^2\cdot[2\theta(1-\theta)]
=2\theta^5(1-\theta),\qquad 0<\theta<1.
\]
取对数，
\(
\ln L(\theta)=\ln2+5\ln\theta+\ln(1-\theta).
\)
求导并令其为 \(0\)，得
\(
\frac{5}{\theta}-\frac{1}{1-\theta}=0.
\)
于是
\(
5(1-\theta)=\theta,
\)
解得
\(
\hat\theta=\frac56.
\)
\end{solution}

\begin{exercise}\label{exer:src-2008-01-3c-iv-05}\label{exer:src-paper4-iv-05}
设总体 \(X\) 的密度函数为
\[
f(x;\theta)=\frac1{2\theta}e^{-|x|/\theta},\qquad \theta>0,\ -\infty<x<+\infty,
\]
其中 \(\theta\) 未知。设 \(X_1,X_2,\dots,X_n\) 为来自 \(X\) 的样本，求 \(\theta\) 的：\par
(1) 矩估计量；\par
(2) 极大似然估计量。
\end{exercise}
\begin{solution}
该分布关于 \(0\) 对称，因此 \(EX=0\)，一阶矩不能用于估计 \(\theta\)。注意到
\[
EX^2=\int_{-\infty}^{+\infty}x^2\frac1{2\theta}e^{-|x|/\theta}\,\mathrm{d}x
=\frac1{\theta}\int_0^{+\infty}x^2e^{-x/\theta}\,\mathrm{d}x.
\]
令 \(t=x/\theta\)，则
\[
EX^2=\theta^2\int_0^{+\infty}t^2e^{-t}\,\mathrm{d}t=2\theta^2.
\]
按原点矩作矩估计，令
\(
\frac1n\sum_{i=1}^{n}X_i^2=2\theta^2,
\)
得矩估计量
\(
\hat\theta_M=\sqrt{\frac1{2n}\sum_{i=1}^{n}X_i^2}.
\)

再求极大似然估计。似然函数为
\[
L(\theta)=\prod_{i=1}^{n}\frac1{2\theta}e^{-|X_i|/\theta}
=\frac1{(2\theta)^n}\exp\!\left(-\frac{\sum_{i=1}^{n}|X_i|}{\theta}\right),
\qquad \theta>0.
\]
对数似然为
\[
\ell(\theta)=-n\ln2-n\ln\theta-\frac{\sum_{i=1}^{n}|X_i|}{\theta}.
\]
求导得
\[
\ell'(\theta)=-\frac n\theta+\frac{\sum_{i=1}^{n}|X_i|}{\theta^2}.
\]
令 \(\ell'(\theta)=0\)，得
\(
\hat\theta_L=\frac1n\sum_{i=1}^{n}|X_i|.
\)
二阶导数在该点为负，故为极大值点。
\end{solution}

\begin{exercise}\label{exer:src-2008-01-3a-iv-04}\label{exer:src-paper2-app-04}\label{prob:p2-app-04}
设总体 \(X\sim N(\mu,\sigma^2)\)，\((X_1,X_2,\dots,X_n)\) 为总体 \(X\) 的一个样本。求常数 \(k\)，使
\(
k\sum_{i=1}^{n}|X_i-\overline X|
\)
为 \(\sigma\) 的无偏估计量。
\end{exercise}
\begin{solution}
记
\(
Y_i=X_i-\overline X.
\)
由于正态样本的线性组合仍为正态变量，且
\(
EY_i=0.
\)
计算方差：
\[
D(Y_i)=D(X_i-\overline X)
=DX_i+D\overline X-2\operatorname{Cov}(X_i,\overline X).
\]
其中
\(
DX_i=\sigma^2,\qquad D\overline X=\frac{\sigma^2}{n},
\)
\[
\operatorname{Cov}(X_i,\overline X)
=\operatorname{Cov}\left(X_i,\frac1n\sum_{j=1}^n X_j\right)
=\frac{\sigma^2}{n}.
\]
所以
\(
D(Y_i)=\sigma^2+\frac{\sigma^2}{n}-2\frac{\sigma^2}{n}
=\frac{n-1}{n}\sigma^2.
\)
即
\(
Y_i\sim N\left(0,\frac{n-1}{n}\sigma^2\right).
\)
若 \(Z\sim N(0,\tau^2)\)，则
\(
E|Z|=\tau\sqrt{\frac2\pi}.
\)
因此
\[
E|X_i-\overline X|
=\sigma\sqrt{\frac{n-1}{n}}\sqrt{\frac2\pi}.
\]
从而
\[
E\left[k\sum_{i=1}^{n}|X_i-\overline X|\right]
=kn\sigma\sqrt{\frac{n-1}{n}}\sqrt{\frac2\pi}.
\]
令其等于 \(\sigma\)，得
\(
k=\frac1{n}\sqrt{\frac{n}{n-1}}\sqrt{\frac{\pi}{2}}
=\sqrt{\frac{\pi}{2n(n-1)}}.
\)
\end{solution}


\chapter{区间估计}

\begin{note}
\textbf{自学导读：从"一个数"到"一个区间"。}

第 7 章的点估计给出一个数 $\hat\mu = \bar{x}$，但你无法知道这个数离真值 $\mu$ 有多近。区间估计的想法是：给出一个区间 $[\bar{x}-\Delta,\ \bar{x}+\Delta]$，并保证"真值 $\mu$ 落在这个区间内的概率 $\geqslant 1-\alpha$"。

\textbf{核心逻辑（三步走）：}
\begin{enumerate}
    \item 找一个包含未知参数 $\theta$ 和样本的统计量（叫\textbf{枢轴量}），其分布已知且不含未知参数；
    \item 根据置信水平 $1-\alpha$，从分布表中查出临界值；
    \item 解不等式，把 $\theta$ 夹在中间，得到置信区间。
\end{enumerate}

\textbf{自学提醒：} 不要死记公式表。记住"$\sigma$ 已知用 $Z$，$\sigma$ 未知用 $t$，估方差用 $\chi^2$"这三条规则，再按三步走推导，公式自然出来。

\textbf{与第 6 章的联系：} 枢轴量的分布正是第 6 章学过的抽样分布定理——$\frac{\bar{X}-\mu}{\sigma/\sqrt{n}}\sim N(0,1)$，$\frac{\bar{X}-\mu}{S/\sqrt{n}}\sim t(n-1)$，$\frac{(n-1)S^2}{\sigma^2}\sim\chi^2(n-1)$。本章只是把这些结论"反过来用"。
\end{note}


% 单个正态总体（完整保留所有内容，优化行距和排版）
{\small
\[
\begin{cases}
\text{单个正态 } X\sim N(\mu,\sigma^2) \\
\text{置信水平为 } 1-\alpha
\end{cases}
\Rightarrow
\begin{cases}
\text{总体均值 } \mu:
\begin{cases}
\sigma^2 \text{ 已知}: & \left( \overline{X} - \dfrac{\sigma}{\sqrt{n}} z_{\frac{\alpha}{2}}, \overline{X} + \dfrac{\sigma}{\sqrt{n}} z_{\frac{\alpha}{2}} \right)\\
\sigma^2 \text{ 未知}: & \left( \overline{X} - \dfrac{S}{\sqrt{n}} t_{\frac{\alpha}{2}}(n-1), \overline{X} + \dfrac{S}{\sqrt{n}} t_{\frac{\alpha}{2}}(n-1) \right)
\end{cases}\\[8pt]
\text{总体方差 } \sigma^2:
\begin{cases}
\mu \text{ 已知}: & \left( \dfrac{\sum_{i=1}^n (X_i-\mu)^2}{\chi_{\frac{\alpha}{2}}^2(n)}, \dfrac{\sum_{i=1}^n (X_i-\mu)^2}{\chi_{1-\frac{\alpha}{2}}^2(n)} \right)\\
\mu \text{ 未知}: & \left( \dfrac{(n-1)S^2}{\chi_{\frac{\alpha}{2}}^2(n-1)}, \dfrac{(n-1)S^2}{\chi_{1-\frac{\alpha}{2}}^2(n-1)} \right)
\end{cases}
\end{cases}
\]
}

\vspace{1em}
\[
\begin{aligned}
&\text{两个正态 } \\
&X_1 \sim N(\mu_1,\sigma_1^2);\\
& X_2 \sim N(\mu_2,\sigma_2^2) \\
&\text{置信水平为 }1-\alpha
\end{aligned}
\;
\left\{
\begin{aligned}
&\text{总体均值差 } (\mu_1-\mu_2): \\
&\qquad
\begin{cases}
\sigma_1^2,\sigma_2^2 \text{ 已知}: \\[2pt]
\quad (\bar X_1-\bar X_2) \pm z_{\alpha/2} \sqrt{\sigma_1^2/n_1+\sigma_2^2/n_2};\\[8pt]
\sigma_1^2=\sigma_2^2=\sigma^2,\;\sigma^2 \text{ 未知}: \\[2pt]
\quad(\bar X_1-\bar X_2) \pm t_{\alpha/2}(n_1+n_2-2)\,S_w\sqrt{1/n_1+1/n_2};
\end{cases}
\\[12pt]
&\text{总体方差比 } \bigl( \sigma_1^2/\sigma_2^2 \bigr): \\
&\qquad
\begin{cases}
\mu_1,\mu_2 \text{ 已知}: \\[2pt]
\quad\Bigl( \frac{\frac1{n_1}\sum_{i=1}^{n_1}(X_i-\mu_1)^2}{\frac1{n_2}\sum_{j=1}^{n_2}(X_j-\mu_2)^2}
\cdot\frac1{F_{\alpha/2}(n_1,n_2)},\;
\frac{\frac1{n_1}\sum_{i=1}^{n_1}(X_i-\mu_1)^2}{\frac1{n_2}\sum_{j=1}^{n_2}(X_j-\mu_2)^2}
\cdot\frac1{F_{1-\alpha/2}(n_1,n_2)} \Bigr) \\[8pt]
\mu_1,\mu_2 \text{ 未知}: \\[2pt]
\quad\Bigl( \frac{S_1^2}{S_2^2}\cdot\frac1{F_{\alpha/2}(n_1-1,n_2-1)},\;
\frac{S_1^2}{S_2^2}\cdot\frac1{F_{1-\alpha/2}(n_1-1,n_2-1)} \Bigr)
\end{cases}
\end{aligned}
\right.
\]


\begin{introduction}
  \item 区间估计的基本概念：双侧置信区间与单侧置信区间
  \item 构造置信区间的基本原理与枢轴量法
  \item 正态总体参数（均值、方差）的区间估计
\end{introduction}

\begin{note}
\textbf{使用提醒：} 区间估计题的选型可以按“估谁 + 总体条件 + 单侧/双侧”三步走。先认是均值、方差、均值差还是方差比；再看总体是否正态、方差已知还是未知；最后根据“至少/至多/是否等于”判断该构造单侧区间还是双侧区间。把这三步先理顺，再去选 \(z/t/\chi^2/F\) 会稳很多。
\end{note}

\section{区间估计的基本概念}

对于未知参数 $\theta$，除了求出它的点估计外，我们往往还希望给出一个合理的范围，并知道这个范围包含参数 $\theta$ 真值的可信程度。这样的范围通常以区间的形式给出，这种估计形式被称为\textbf{区间估计}。

\begin{definition}[置信区间与置信水平] \label{def:confidence_interval}
设 $\theta$ 是总体 $X$ 的分布函数中的一个未知参数，对于给定的显著性水平 $\alpha\ (0<\alpha<1)$，如果由样本 $X_1,X_2,\dots,X_n$ 确定的两个统计量 $\underline{\theta}=\underline{\theta}(X_1,X_2,\dots,X_n)$ 和 $\overline{\theta}=\overline{\theta}(X_1,X_2,\dots,X_n)$，满足 $\underline{\theta} < \overline{\theta}$，且对于任意的 $\theta \in \Theta$ 均有：
\begin{equation}
P\{\underline{\theta}<\theta<\overline{\theta}\} \ge 1-\alpha
\end{equation}
则称随机区间 $(\underline{\theta},\overline{\theta})$ 是 $\theta$ 的置信度（或置信水平）为 $1-\alpha$ 的\textbf{双侧置信区间}。其中，$\underline{\theta}$ 和 $\overline{\theta}$ 分别称为置信度为 $1-\alpha$ 的\textbf{置信下限}和\textbf{置信上限}。

特别地，如果满足 $P\{\theta < \underline{\theta}\} = P\{\theta > \overline{\theta}\} = \frac{\alpha}{2}$，则称这种置信区间为\textbf{等尾置信区间}。
\end{definition}

\begin{note}
置信区间最容易混淆的地方是“概率到底作用在谁身上”。参数 \(\theta\) 是固定但未知的，真正有随机性的是由样本算出来的区间端点 \(\underline{\theta},\overline{\theta}\)。因此“置信水平 \(1-\alpha\)”的含义是：\textbf{在重复抽样中，这种构造方法给出的区间有 \(1-\alpha\) 的比例能覆盖真值}，而不是“对当前这一个已经算出来的区间，\(\theta\) 落在其中的概率是 \(1-\alpha\)”。
\end{note}

\begin{note}
\textbf{自学追问：} 为什么每道区间题都先找枢轴量？因为要先构造一个“分布已知、又含参数”的量，再把不等式反解成参数区间。换句话说，置信区间的计算本质上是“先标准化，再反解”。
\end{note}

在实际应用中，有时我们只关心参数的下界或上界（例如零件的寿命越长越好，废品率越低越好），此时就需要使用单侧置信区间。

\begin{definition}[单侧置信区间] \label{def:one_sided_confidence_interval}
对于给定值 $\alpha\ (0<\alpha<1)$，若由样本确定的统计量 $\underline{\theta}$ 满足 $P\{\theta > \underline{\theta}\} \ge 1-\alpha$，则称 $(\underline{\theta},+\infty)$ 是 $\theta$ 的置信水平为 $1-\alpha$ 的\textbf{单侧置信区间}，$\underline{\theta}$ 称为\textbf{单侧置信下限}。

同理，若由样本确定的统计量 $\overline{\theta}$ 满足 $P\{\theta < \overline{\theta}\} \ge 1-\alpha$，则称 $(-\infty,\overline{\theta})$ 是 $\theta$ 的置信水平为 $1-\alpha$ 的\textbf{单侧置信区间}，$\overline{\theta}$ 称为\textbf{单侧置信上限}。
\end{definition}

\begin{note}
若总体 $X$ 为连续型随机变量，我们通常要求上述概率严格等于 $1-\alpha$；若 $X$ 为离散型随机变量，则使其尽可能地接近（且不小于） $1-\alpha$。单侧区间适用于“只关心下界”或“只关心上界”的场景，例如寿命至少多大、次品率至多多大。不要把单侧区间中的 \(\alpha\) 仍按双侧区间那样平均分到两端。
\end{note}

\begin{note}
\textbf{自学追问：} 什么时候该选下限，什么时候该选上限？若题目关心“参数至少有多大”“不能太小”，就构造单侧下限；若关心“参数至多有多大”“不能太大”，就构造单侧上限。方向由实际问题中的“至少/至多、下不低于/上不超过”这些关键词决定。
\end{note}

\begin{exercise}\label{exer:src-paper1-tf-06}\label{prob:p1-tf-06}
在给定的置信度 \(1-\alpha\) 下，被估参数的置信区间不一定唯一。判断正误。
\end{exercise}
\begin{solution}
正确。

置信区间是一种由样本构造随机区间的方法，只要求覆盖概率达到给定置信度。例如对连续分布的参数，可以构造等尾区间，也可以把尾部概率不平均地分配到左右两端，只要整体覆盖概率仍为 \(1-\alpha\)，就可能得到不同的置信区间。因此同一置信度下置信区间不一定唯一。
\end{solution}

\section{构造置信区间的基本原理}

构造置信区间最常用的方法是\textbf{枢轴量法}。枢轴量是一个由样本和未知参数共同构成的函数，其最大的特点是\textbf{它的分布完全已知，不依赖于任何未知参数}。下面我们以正态分布均值的区间估计为例，说明枢轴量法的基本步骤。

\begin{note}
\textbf{易错点：} 枢轴量不是“把未知参数消掉后的纯样本函数”，而是“虽然还含参数，但它的分布已经与未知参数无关”的量。若你构造出来的统计量里，分布形式仍然要靠未知参数才能写清楚，那它就还不能直接拿来反解置信区间。
\end{note}

利用上述枢轴量法，我们推导出了现实中最常见的正态总体（包括单个和两个正态总体）在不同情况下的参数置信区间。为了便于查阅，我们将相关公式归纳为表~\ref{tab:confidence_intervals}。

\begin{table}[htbp]
\centering
\caption{常见的正态总体参数置信区间（置信度为 $1-\alpha$）}
\label{tab:confidence_intervals}
\resizebox{\textwidth}{!}{
\begin{tabular}{lllc}
\toprule
\textbf{总体条件} & \textbf{待估参数} & \textbf{枢轴量及分布} & \textbf{置信区间} \\
\midrule
\multicolumn{4}{c}{\textbf{单个正态总体 $X \sim N(\mu, \sigma^2)$}} \\
\midrule
$\sigma^2$ 已知 & $\mu$ & $\frac{\overline{X}-\mu}{\sigma/\sqrt{n}} \sim N(0,1)$ & $\left(\overline{X}-\frac{\sigma}{\sqrt{n}}z_{\alpha/2},\ \overline{X}+\frac{\sigma}{\sqrt{n}}z_{\alpha/2}\right)$ \\
$\sigma^2$ 未知 & $\mu$ & $\frac{\overline{X}-\mu}{S/\sqrt{n}} \sim t(n-1)$ & $\left(\overline{X}-\frac{S}{\sqrt{n}}t_{\alpha/2}(n-1),\ \overline{X}+\frac{S}{\sqrt{n}}t_{\alpha/2}(n-1)\right)$ \\
$\mu$ 未知 & $\sigma^2$ & $\frac{(n-1)S^2}{\sigma^2} \sim \chi^2(n-1)$ & $\left(\frac{(n-1)S^2}{\chi^2_{\alpha/2}(n-1)},\ \frac{(n-1)S^2}{\chi^2_{1-\alpha/2}(n-1)}\right)$ \\
$\mu$ 已知 & $\sigma^2$ & $\frac{\sum(X_i-\mu)^2}{\sigma^2} \sim \chi^2(n)$ & $\left(\frac{\sum(X_i-\mu)^2}{\chi^2_{\alpha/2}(n)},\ \frac{\sum(X_i-\mu)^2}{\chi^2_{1-\alpha/2}(n)}\right)$ \\
\midrule
\multicolumn{4}{c}{\textbf{两个正态总体 $X \sim N(\mu_1, \sigma_1^2), Y \sim N(\mu_2, \sigma_2^2)$}} \\
\midrule
$\sigma_1^2,\sigma_2^2$ 已知 & $\mu_1-\mu_2$ & 正态分布 $N(0,1)$ & $\left((\overline{X}-\overline{Y}) \pm z_{\alpha/2}\sqrt{\frac{\sigma_1^2}{n_1}+\frac{\sigma_2^2}{n_2}}\right)$ \\
$\sigma_1^2=\sigma_2^2$ 未知 & $\mu_1-\mu_2$ & $t$ 分布 $t(n_1+n_2-2)$ & $\left((\overline{X}-\overline{Y}) \pm t_{\alpha/2}(n_1+n_2-2)S_w\sqrt{\frac{1}{n_1}+\frac{1}{n_2}}\right)$ \\
$\mu_1,\mu_2$ 已知 & $\frac{\sigma_1^2}{\sigma_2^2}$ & $F$ 分布 $F(n_1,n_2)$ & $\left(\frac{\sum(X_i-\mu_1)^2/n_1}{\sum(Y_j-\mu_2)^2/n_2}\frac{1}{F_{\alpha/2}(n_1,n_2)},\ \frac{\sum(X_i-\mu_1)^2/n_1}{\sum(Y_j-\mu_2)^2/n_2} F_{\alpha/2}(n_2,n_1)\right)$ \\
$\mu_1,\mu_2$ 未知 & $\frac{\sigma_1^2}{\sigma_2^2}$ & $F$ 分布 $F(n_1-1,n_2-1)$ & $\left(\frac{S_1^2}{S_2^2}\frac{1}{F_{\alpha/2}(n_1-1,n_2-1)},\ \frac{S_1^2}{S_2^2}F_{\alpha/2}(n_2-1,n_1-1)\right)$ \\
\bottomrule
\end{tabular}
}
\end{table}

\begin{remark}
上述表格中公式的推导核心在于”寻找正确的枢轴量”，因此读者不应当采用死记硬背的方式来记忆公式，而是应当通过枢轴量和分位数的定义，自行推导得出。
\end{remark}

\begin{note}
\textbf{枢轴量选取的决策流程（做题三问）：}
\begin{enumerate}
    \item \textbf{估什么？} 确定待估参数是 $\mu$ 还是 $\sigma^2$（或两总体的差/比）。
    \item \textbf{知什么？} 看题目条件中哪些参数已知、哪些未知：
    \begin{itemize}
        \item 估 $\mu$ 且 $\sigma^2$ 已知 $\Longrightarrow$ 用 $Z=\dfrac{\bar X-\mu}{\sigma/\sqrt{n}}\sim N(0,1)$；
        \item 估 $\mu$ 且 $\sigma^2$ 未知 $\Longrightarrow$ 用 $T=\dfrac{\bar X-\mu}{S/\sqrt{n}}\sim t(n-1)$；
        \item 估 $\sigma^2$ 且 $\mu$ 未知 $\Longrightarrow$ 用 $\chi^2=\dfrac{(n-1)S^2}{\sigma^2}\sim\chi^2(n-1)$；
        \item 估 $\sigma^2$ 且 $\mu$ 已知 $\Longrightarrow$ 用 $\chi^2=\dfrac{\sum(X_i-\mu)^2}{\sigma^2}\sim\chi^2(n)$。
    \end{itemize}
    \item \textbf{反解。} 对枢轴量写出 $P(a\le W\le b)=1-\alpha$，再把不等式解成”参数 $\in$ (下界, 上界)”的形式。
\end{enumerate}
记住口诀：\textbf{估谁含谁，知啥用啥，分布定型后反解}。
\end{note}

\begin{exercise}[正态总体均值的区间估计] \label{ex:pivotal_quantity}
设总体 $X\sim N(\mu,\sigma^2)$，其中方差 $\sigma^2$ 已知，均值 $\mu$ 未知。现有一组来自 $X$ 的样本 $X_1,X_2,\dots,X_n$，求 $\mu$ 的置信水平为 $1-\alpha$ 的置信区间。
\end{exercise}

\begin{solution}
这道题是枢轴量法最标准的起点。思路分三步：先找枢轴量，再写出中心概率式，最后把不等式解成关于 \(\mu\) 的区间。

第一步，构造关于 \(\mu\) 的枢轴量。因为总体正态且 \(\sigma^2\) 已知，所以
\(
\overline{X}\sim N\!\left(\mu,\frac{\sigma^2}{n}\right).
\)
把它标准化可得
\[
W=\frac{\overline{X}-\mu}{\sigma/\sqrt{n}}\sim N(0,1).
\]
这里 \(W\) 的分布已经完全已知，不含未知参数，因此它就是所需的枢轴量。

第二步，根据置信水平 \(1-\alpha\) 写出等尾概率式：
\[
P\left\{-z_{\alpha/2}<\frac{\overline{X}-\mu}{\sigma/\sqrt{n}}<z_{\alpha/2}\right\}=1-\alpha.
\]

第三步，把中间的不等式对 \(\mu\) 进行变形。两边同乘 \(\sigma/\sqrt{n}\)，再移项，得到
\[
P\left\{\overline{X}-\frac{\sigma}{\sqrt{n}}z_{\alpha/2}<\mu<\overline{X}+\frac{\sigma}{\sqrt{n}}z_{\alpha/2}\right\}=1-\alpha.
\]

因此，\(\mu\) 的置信水平为 \(1-\alpha\) 的置信区间为
\[
\left(\overline{X}-\frac{\sigma}{\sqrt{n}}z_{\alpha/2},\ \overline{X}+\frac{\sigma}{\sqrt{n}}z_{\alpha/2}\right).
\]

\begin{note}
本题之所以用 \(z_{\alpha/2}\)，关键在于\(\sigma\) 已知且采用双侧等尾区间。以后若 \(\sigma\) 未知，就不能继续使用这个正态枢轴量，而要改用 \(t\) 分布。
\end{note}
\end{solution}

\begin{exercise}\label{exer:src-paper1-fill-06}\label{prob:p1-fill-06}
设某种清漆干燥时间 \(X\sim N(\mu,\sigma^2)\)（单位：小时），取 \(n=9\) 的样本，得样本均值和方差分别为 \(\overline X=6,\ S^2=0.33\)。则 \(\mu\) 的置信度为 \(95\%\) 的单侧置信区间上限为 \(\underline{\hspace{2cm}}\)。已知 \(t_{0.05}(8)=1.8595\)。
\end{exercise}
\begin{solution}
应填 \(6.356\)。

总体正态且 \(\sigma^2\) 未知，使用
\(
T=\frac{\overline X-\mu}{S/\sqrt n}\sim t(n-1)=t(8).
\)
要构造置信度为 \(95\%\) 的单侧上限 \(U\)，取
\(
P\left\{T\ge -t_{0.05}(8)\right\}=0.95,
\)
即
\(
P\left\{\mu\le \overline X+t_{0.05}(8)\frac{S}{\sqrt n}\right\}=0.95.
\)
因此
\(
U=6+1.8595\cdot\frac{\sqrt{0.33}}{3}
\approx 6.356.
\)
\end{solution}

\subsection*{历年试题补充：2011 试卷 A 卷}

\begin{exercise}\label{exer:src-2011-a-07-34credit}
假设 \(0.50,1.25,0.80,2.00\) 是来自总体 \(X\) 的简单随机样本值。已知 \(Y=\ln X\) 服从正态分布 \(N(\mu,1)\)。求：\par
(1) \(X\) 的数学期望 \(EX\)（记 \(EX\) 为 \(b\)）；\par
(2) \(\mu\) 的置信度为 \(0.95\) 的置信区间；\par
(3) \(b\) 的置信度为 \(0.95\) 的置信区间。
\end{exercise}
\begin{solution}
(1) 因为 \(Y=\ln X\sim N(\mu,1)\)，所以 \(X=e^Y\)。于是
\(
b=EX=E(e^Y)=e^{\mu+\frac12}.
\)

(2) 对样本作对数变换，得到 \(Y_1,\dots,Y_4\) 的样本。由于 \(Y\sim N(\mu,1)\)，
\(
\overline Y\sim N\left(\mu,\frac14\right).
\)
样本均值为
\[
\overline y
=\frac14(\ln0.50+\ln1.25+\ln0.80+\ln2.00)
=\frac14\ln(0.50\times1.25\times0.80\times2.00)=0.
\]
置信度为 \(0.95\) 时，\(z_{0.025}=1.96\)，故 \(\mu\) 的置信区间为
\(
\left(\overline y-\frac{1.96}{2},\ \overline y+\frac{1.96}{2}\right)
=(-0.98,\ 0.98).
\)

(3) 由 \(b=e^{\mu+1/2}\)，且指数函数严格单调递增，将上面的区间整体作变换：
\[
-0.98<\mu<0.98
\quad\Longrightarrow\quad
-0.48<\mu+\frac12<1.48.
\]
因此 \(b\) 的 \(0.95\) 置信区间为
\(
\left(e^{-0.48},\ e^{1.48}\right).
\)
\end{solution}

\begin{exercise}\label{exer:src-2011-ab-a-06-34credit}
某地某种商品在一家商场中的月消费额 \(X\sim N(\mu,\sigma^2)\)，且已知 \(\sigma=100\) 元。现商业部门要对该商品在商场中的平均月消费额 \(\mu\) 进行估计，并要求以不小于 \(95\%\) 的把握保证估计结果的误差不超过 \(20\) 元。问至少需要随机调查多少家商场？
\end{exercise}
\begin{solution}
因为总体标准差 \(\sigma=100\) 已知，要求
\(
P\{|\overline X-\mu|\le20\}\ge0.95.
\)
由
\(
\frac{\overline X-\mu}{\sigma/\sqrt n}\sim N(0,1),
\)
可得
\(
P\{|\overline X-\mu|\le20\}
=2\Phi\left(\frac{20\sqrt n}{100}\right)-1.
\)
令其不小于 \(0.95\)，即
\(
\frac{20\sqrt n}{100}\ge z_{0.025}=1.96.
\)
于是
\(
n\ge\left(\frac{1.96\times100}{20}\right)^2=96.04.
\)
所以至少需要随机调查 \(97\) 家商场。
\end{solution}

\begin{exercise}\label{exer:src-2011-ab-b-06-34credit}
某地某种商品在一家商场中的月消费额 \(X\sim N(\mu,\sigma^2)\)，且已知 \(\sigma=200\) 元。现商业部门要对该商品在商场中的平均月消费额 \(\mu\) 进行估计，并要求以不小于 \(95\%\) 的把握保证估计结果的误差不超过 \(80\) 元。问至少需要随机调查多少家商场？
\end{exercise}
\begin{solution}
同样使用正态总体均值在 \(\sigma\) 已知时的误差控制公式：
\(
P\{|\overline X-\mu|\le80\}
=2\Phi\left(\frac{80\sqrt n}{200}\right)-1.
\)
要求该概率不小于 \(0.95\)，故
\(
\frac{80\sqrt n}{200}\ge1.96.
\)
于是
\(
n\ge\left(\frac{1.96\times200}{80}\right)^2=24.01.
\)
所以至少需要随机调查 \(25\) 家商场。
\end{solution}

除了上述 $\sigma^2$ 已知时 $\mu$ 的区间估计外，我们还可以利用类似的方法推导其他常见参数的置信区间。

\begin{exercise}[正态总体均值的区间估计（方差未知）]\label{exer:src-paper3-iii-05}\label{prob:p3-iii-05}\label{ex:ci_mu_unknown_var}
设总体 $X\sim N(\mu,\sigma^2)$，其中均值 $\mu$ 和方差 $\sigma^2$ 均未知。现有一组来自 $X$ 的独立样本 $X_1,X_2,\dots,X_n$，样本均值为 $\overline{X}$，样本方差为 $S^2$。求 $\mu$ 的置信水平为 $1-\alpha$ 的置信区间。
\end{exercise}

\begin{solution}
与上一题相比，唯一的关键变化是：\(\sigma^2\) 现在未知，所以不能再直接用
\(
\frac{\overline{X}-\mu}{\sigma/\sqrt{n}}.
\)
这时要改用 \(t\) 统计量。

因为总体正态，所以有经典结论
\(
W=\frac{\overline{X}-\mu}{S/\sqrt{n}}\sim t(n-1).
\)
这里之所以能用 \(t\) 分布，靠的是“总体正态”以及“\(\overline{X}\) 与 \(S^2\) 独立”这两个前提。

在置信水平 \(1-\alpha\) 下，取双侧等尾区间：
\[
P\left\{-t_{\alpha/2}(n-1)<\frac{\overline{X}-\mu}{S/\sqrt{n}}<t_{\alpha/2}(n-1)\right\}=1-\alpha.
\]

把不等式解出 \(\mu\)，得到
\[
P\left\{\overline{X}-\frac{S}{\sqrt{n}}t_{\alpha/2}(n-1)<\mu<\overline{X}+\frac{S}{\sqrt{n}}t_{\alpha/2}(n-1)\right\}=1-\alpha.
\]

因此 \(\mu\) 的置信区间为
\[
\left(\overline{X}-\frac{S}{\sqrt{n}}t_{\alpha/2}(n-1),\ \overline{X}+\frac{S}{\sqrt{n}}t_{\alpha/2}(n-1)\right).
\]

\begin{note}
这道题最容易错的地方有两个：一是把 \(z_{\alpha/2}\) 误写成 \(t_{\alpha/2}(n-1)\) 之外的正态分位点；二是把自由度写成 \(n\)。只要样本方差 \(S^2\) 是围绕 \(\overline{X}\) 构造的，\(t\) 分布自由度就是 \(n-1\)。
\end{note}
\end{solution}

\begin{exercise}[正态总体方差的区间估计（均值未知）] \label{ex:ci_var_unknown_mu}
设总体 $X\sim N(\mu,\sigma^2)$，其中均值 $\mu$ 和方差 $\sigma^2$ 均未知。现有一组来自 $X$ 的独立样本 $X_1,X_2,\dots,X_n$，样本方差为 $S^2$。求方差 $\sigma^2$ 的置信水平为 $1-\alpha$ 的置信区间。
\end{exercise}

\begin{solution}
这一题要估计的是方差 \(\sigma^2\)，所以应优先寻找与样本方差有关的枢轴量。

在总体正态且总体均值未知的条件下，有
\(
W=\frac{(n-1)S^2}{\sigma^2}\sim\chi^2(n-1).
\)
这个统计量的分布不再含未知参数，因此可以直接用来构造区间。

由于 \(\chi^2\) 分布不是对称分布，双侧区间不能像正态和 \(t\) 分布那样写成“正负同一个分位点”的形式，而应分别取左右尾各 \(\alpha/2\)：
\[
P\left\{\chi^2_{1-\alpha/2}(n-1)<\frac{(n-1)S^2}{\sigma^2}<\chi^2_{\alpha/2}(n-1)\right\}=1-\alpha.
\]

现在解这个不等式。注意中间量与 \(\sigma^2\) 在分母上，变形时要格外小心方向：
\[
P\left\{\frac{(n-1)S^2}{\chi^2_{\alpha/2}(n-1)}<\sigma^2<\frac{(n-1)S^2}{\chi^2_{1-\alpha/2}(n-1)}\right\}=1-\alpha.
\]

因此，\(\sigma^2\) 的置信区间为
\[
\left(\frac{(n-1)S^2}{\chi^2_{\alpha/2}(n-1)},\ \frac{(n-1)S^2}{\chi^2_{1-\alpha/2}(n-1)}\right).
\]

\begin{note}
本题最容易错在上下界次序。因为 \(\chi^2_{\alpha/2}(n-1)\) 比 \(\chi^2_{1-\alpha/2}(n-1)\) 大，所以它反而出现在分母较大的那一侧，最终形成较小的下界和较大的上界。若机械照搬“左分位写左边、右分位写右边”，就很容易写反。
\end{note}
\end{solution}

对于两个正态总体 $X \sim N(\mu_1, \sigma_1^2)$ 和 $Y \sim N(\mu_2, \sigma_2^2)$，我们同样可以利用类似的方法构造枢轴量，来估计它们的均值差或方差比。

\begin{exercise}[两正态总体均值差的区间估计（方差相等且未知）] \label{ex:ci_diff_mu_unknown_var}
设有两个独立的正态总体 $X \sim N(\mu_1, \sigma^2)$ 和 $Y \sim N(\mu_2, \sigma^2)$，方差相等但具体数值未知。分别抽取容量为 $n_1$ 和 $n_2$ 的独立样本，求均值差 $\mu_1-\mu_2$ 的置信水平为 $1-\alpha$ 的置信区间。
\end{exercise}

\begin{solution}
我们要估计的是均值差 \(\mu_1-\mu_2\)，最自然的点估计量是
\(
\overline{X}-\overline{Y}.
\)

题目特别强调“两总体方差相等但未知”，这时不能分别用 \(S_1,S_2\) 生搬硬套，而要先构造\textbf{合并样本方差}
\(
S_w^2=\frac{(n_1-1)S_1^2+(n_2-1)S_2^2}{n_1+n_2-2}.
\)

在两个总体正态、两组样本独立且 \(\sigma_1^2=\sigma_2^2\) 的前提下，有
\[
W=\frac{(\overline{X}-\overline{Y})-(\mu_1-\mu_2)}{S_w\sqrt{\frac{1}{n_1}+\frac{1}{n_2}}}\sim t(n_1+n_2-2).
\]

于是按双侧等尾原则写出
\[
P\left\{-t_{\alpha/2}(n_1+n_2-2)<\frac{(\overline{X}-\overline{Y})-(\mu_1-\mu_2)}{S_w\sqrt{\frac{1}{n_1}+\frac{1}{n_2}}}<t_{\alpha/2}(n_1+n_2-2)\right\}=1-\alpha.
\]

解出 \(\mu_1-\mu_2\)，便得到区间
\[
\left((\overline{X}-\overline{Y})-t_{\alpha/2}(n_1+n_2-2)S_w\sqrt{\frac{1}{n_1}+\frac{1}{n_2}},\ 
(\overline{X}-\overline{Y})+t_{\alpha/2}(n_1+n_2-2)S_w\sqrt{\frac{1}{n_1}+\frac{1}{n_2}}\right).
\]

\begin{note}
这里最关键的适用条件不是“方差未知”，而是“\textbf{两总体方差相等且未知}”。若方差不相等，就不能直接用这一套合并方差 \(t\) 公式。
\end{note}
\end{solution}

\begin{exercise}[两正态总体方差比的区间估计（均值未知）] \label{ex:ci_ratio_var_unknown_mu}
设有两个独立的正态总体 $X \sim N(\mu_1, \sigma_1^2)$ 和 $Y \sim N(\mu_2, \sigma_2^2)$，其均值均未知。求方差比 $\sigma_1^2/\sigma_2^2$ 的置信水平为 $1-\alpha$ 的置信区间。
\end{exercise}

\begin{solution}
要估计方差比 \(\sigma_1^2/\sigma_2^2\)，最自然的样本对应物是 \(S_1^2/S_2^2\)。在两个总体正态且均值未知的条件下，有
\[
\frac{S_1^2/\sigma_1^2}{S_2^2/\sigma_2^2}
=\frac{S_1^2}{S_2^2}\cdot\frac{\sigma_2^2}{\sigma_1^2}
\sim F(n_1-1,n_2-1).
\]
这就是所需的枢轴量。

在置信水平 \(1-\alpha\) 下，写成双侧形式：
\[
P\left\{F_{1-\alpha/2}(n_1-1,n_2-1)<\frac{S_1^2}{S_2^2}\cdot\frac{\sigma_2^2}{\sigma_1^2}<F_{\alpha/2}(n_1-1,n_2-1)\right\}=1-\alpha.
\]

接下来把它整理成关于 \(\sigma_1^2/\sigma_2^2\) 的不等式。变形后得到
\[
P\left\{\frac{S_1^2}{S_2^2}\cdot\frac{1}{F_{\alpha/2}(n_1-1,n_2-1)}
<\frac{\sigma_1^2}{\sigma_2^2}
<\frac{S_1^2}{S_2^2}\cdot\frac{1}{F_{1-\alpha/2}(n_1-1,n_2-1)}\right\}=1-\alpha.
\]

再利用 \(F\) 分布分位点的倒数关系
\(
F_{1-\alpha/2}(n_1-1,n_2-1)=\frac{1}{F_{\alpha/2}(n_2-1,n_1-1)},
\)
可把上界写得更常用一些。于是最终区间为
\[
\left(\frac{S_1^2}{S_2^2}\cdot\frac{1}{F_{\alpha/2}(n_1-1,n_2-1)},\ 
\frac{S_1^2}{S_2^2}\cdot F_{\alpha/2}(n_2-1,n_1-1)\right).
\]

\begin{note}
这题常错在两个地方：一是把自由度误写成 \(n_1,n_2\)；二是忘记 \(F\) 分布左右分位点并不对称，需要借助倒数关系来整理上界。
\end{note}
\end{solution}

\begin{exercise}[正态总体方差的区间估计（均值已知）] \label{ex:ci_var_known_mu}
设总体 $X\sim N(\mu,\sigma^2)$，其中均值 $\mu$ 已知，方差 $\sigma^2$ 未知。现有一组来自 $X$ 的独立样本 $X_1,X_2,\dots,X_n$。求方差 $\sigma^2$ 的置信水平为 $1-\alpha$ 的置信区间。
\end{exercise}

\begin{solution}
由于均值 $\mu$ 已知，我们在构造关于方差 $\sigma^2$ 的枢轴量时，不应再使用样本均值 $\overline{X}$ 带来的自由度损失，而是直接使用真实均值 $\mu$。

首先，对每个样本进行标准化处理。因为 $X_i \sim N(\mu, \sigma^2)$，所以：
\begin{equation*}
\frac{X_i-\mu}{\sigma} \sim N(0,1), \quad i=1,2,\dots,n
\end{equation*}

根据卡方分布的定义，有限个相互独立的标准正态随机变量的平方和服从 $\chi^2$ 分布。将上述 $n$ 个相互独立的随机变量平方后求和，得到我们所需的枢轴量 $W$：
\begin{equation*}
W = \sum_{i=1}^n \left(\frac{X_i-\mu}{\sigma}\right)^2 = \frac{\sum_{i=1}^n (X_i-\mu)^2}{\sigma^2} \sim \chi^2(n)
\end{equation*}
此时 $W$ 服从自由度为 $n$ 的卡方分布（注意：因为使用的是已知的总体均值 $\mu$ 而不是样本均值 $\overline{X}$，所以自由度没有减少，仍为 $n$），且 $W$ 的分布不依赖于未知参数 $\sigma^2$，是一个理想的枢轴量。

利用 $\chi^2(n)$ 分布的上 $\alpha/2$ 分位点 $\chi^2_{\alpha/2}(n)$ 和上 $1-\alpha/2$ 分位点 $\chi^2_{1-\alpha/2}(n)$，在等尾概率的情况下有：
\begin{equation*}
P\left\{\chi^2_{1-\alpha/2}(n) < \frac{\sum_{i=1}^n (X_i-\mu)^2}{\sigma^2} < \chi^2_{\alpha/2}(n)\right\}=1-\alpha
\end{equation*}

接下来从不等式中解出 $\sigma^2$。对各项取倒数（此时不等号反向）：
\begin{equation*}
P\left\{\frac{1}{\chi^2_{\alpha/2}(n)} < \frac{\sigma^2}{\sum_{i=1}^n (X_i-\mu)^2} < \frac{1}{\chi^2_{1-\alpha/2}(n)}\right\}=1-\alpha
\end{equation*}

同时乘以 $\sum_{i=1}^n (X_i-\mu)^2$，得到：
\begin{equation*}
P\left\{\frac{\sum_{i=1}^n (X_i-\mu)^2}{\chi^2_{\alpha/2}(n)} < \sigma^2 < \frac{\sum_{i=1}^n (X_i-\mu)^2}{\chi^2_{1-\alpha/2}(n)}\right\}=1-\alpha
\end{equation*}

由此得到方差 $\sigma^2$ 的置信水平为 $1-\alpha$ 的置信区间为：
\begin{equation*}
\left(\frac{\sum_{i=1}^n(X_i-\mu)^2}{\chi^2_{\alpha/2}(n)},\ \frac{\sum_{i=1}^n(X_i-\mu)^2}{\chi^2_{1-\alpha/2}(n)}\right)
\end{equation*}
\begin{note}
与“\(\mu\) 未知”时相比，这里自由度是 \(n\) 而不是 \(n-1\)。原因不在于题型换了，而在于我们没有再用样本去估计总体均值，因此没有损失那一个自由度。
\end{note}
\end{solution}


\begin{exercise}[两正态总体均值差的区间估计（方差已知）] \label{ex:ci_diff_mu_known_var}
设有两个独立的正态总体 $X \sim N(\mu_1, \sigma_1^2)$ 和 $Y \sim N(\mu_2, \sigma_2^2)$，其中方差 $\sigma_1^2, \sigma_2^2$ 均已知，但均值 $\mu_1, \mu_2$ 未知。分别抽取容量为 $n_1$ 和 $n_2$ 的独立样本。求均值差 $\mu_1-\mu_2$ 的置信水平为 $1-\alpha$ 的置信区间。
\end{exercise}

\begin{solution}
我们已知样本均值的分布分别为：
\begin{equation*}
\overline{X} \sim N\left(\mu_1, \frac{\sigma_1^2}{n_1}\right), \quad \overline{Y} \sim N\left(\mu_2, \frac{\sigma_2^2}{n_2}\right)
\end{equation*}

要估计均值差 $\mu_1-\mu_2$，最自然的无偏估计量是两样本均值之差 $\overline{X}-\overline{Y}$。由于两样本相互独立，根据正态分布的可加性，$\overline{X}-\overline{Y}$ 也服从正态分布，其数学期望和方差分别为：
\begin{align*}
E(\overline{X}-\overline{Y}) &= E(\overline{X}) - E(\overline{Y}) = \mu_1 - \mu_2 \\
D(\overline{X}-\overline{Y}) &= D(\overline{X}) + D(\overline{Y}) = \frac{\sigma_1^2}{n_1} + \frac{\sigma_2^2}{n_2}
\end{align*}
因此有：
\begin{equation*}
\overline{X}-\overline{Y} \sim N\left(\mu_1-\mu_2,\ \frac{\sigma_1^2}{n_1} + \frac{\sigma_2^2}{n_2}\right)
\end{equation*}

对其进行标准化处理，便可构造出完全不含未知参数的枢轴量 $W$：
\begin{equation*}
W = \frac{(\overline{X}-\overline{Y}) - (\mu_1-\mu_2)}{\sqrt{\frac{\sigma_1^2}{n_1} + \frac{\sigma_2^2}{n_2}}} \sim N(0,1)
\end{equation*}

利用标准正态分布的对称性，以及上 $\alpha/2$ 分位点 $z_{\alpha/2}$，可以建立概率等式：
\begin{equation*}
P\left\{-z_{\alpha/2} < \frac{(\overline{X}-\overline{Y}) - (\mu_1-\mu_2)}{\sqrt{\frac{\sigma_1^2}{n_1} + \frac{\sigma_2^2}{n_2}}} < z_{\alpha/2}\right\} = 1-\alpha
\end{equation*}

解此不等式，分离出待估参数 $\mu_1-\mu_2$：
\begin{equation*}
-z_{\alpha/2}\sqrt{\frac{\sigma_1^2}{n_1} + \frac{\sigma_2^2}{n_2}} < (\overline{X}-\overline{Y}) - (\mu_1-\mu_2) < z_{\alpha/2}\sqrt{\frac{\sigma_1^2}{n_1} + \frac{\sigma_2^2}{n_2}}
\end{equation*}
同时乘以 $-1$ 并加上 $\overline{X}-\overline{Y}$，由于对称性，正负边界的形式依然保持：
\begin{equation*}
P\left\{(\overline{X}-\overline{Y}) - z_{\alpha/2}\sqrt{\frac{\sigma_1^2}{n_1} + \frac{\sigma_2^2}{n_2}} < \mu_1-\mu_2 < (\overline{X}-\overline{Y}) + z_{\alpha/2}\sqrt{\frac{\sigma_1^2}{n_1} + \frac{\sigma_2^2}{n_2}}\right\} = 1-\alpha
\end{equation*}

故 $\mu_1-\mu_2$ 的置信水平为 $1-\alpha$ 的置信区间为：
\begin{equation*}
\left((\overline{X}-\overline{Y}) - z_{\alpha/2}\sqrt{\frac{\sigma_1^2}{n_1}+\frac{\sigma_2^2}{n_2}},\ (\overline{X}-\overline{Y}) + z_{\alpha/2}\sqrt{\frac{\sigma_1^2}{n_1}+\frac{\sigma_2^2}{n_2}}\right)
\end{equation*}
\end{solution}


\begin{exercise}[两正态总体方差比的区间估计（均值已知）] \label{ex:ci_ratio_var_known_mu}
设有两个独立的正态总体 $X \sim N(\mu_1, \sigma_1^2)$ 和 $Y \sim N(\mu_2, \sigma_2^2)$，其中均值 $\mu_1, \mu_2$ 均已知，但方差 $\sigma_1^2, \sigma_2^2$ 未知。分别抽取容量为 $n_1$ 和 $n_2$ 的独立样本。求方差比 $\sigma_1^2/\sigma_2^2$ 的置信水平为 $1-\alpha$ 的置信区间。
\end{exercise}

\begin{solution}
由于总体的真实均值 $\mu_1, \mu_2$ 均已知，我们可以充分利用它们构造 $\chi^2$ 统计量，而无需借助于消耗自由度的样本方差 $S^2$。

根据前面的推导，对于总体 $X$ 和总体 $Y$，我们分别有：
\begin{equation*}
U = \frac{\sum_{i=1}^{n_1}(X_i-\mu_1)^2}{\sigma_1^2} \sim \chi^2(n_1) \quad \text{和} \quad V = \frac{\sum_{j=1}^{n_2}(Y_j-\mu_2)^2}{\sigma_2^2} \sim \chi^2(n_2)
\end{equation*}

根据 $F$ 分布的定义，由两个独立的 $\chi^2$ 分布除以它们各自的自由度后求比值，即可得到服从 $F$ 分布的随机变量。因此，我们构造如下枢轴量：
\begin{equation*}
W = \frac{U/n_1}{V/n_2} = \frac{\frac{1}{n_1\sigma_1^2}\sum_{i=1}^{n_1}(X_i-\mu_1)^2}{\frac{1}{n_2\sigma_2^2}\sum_{j=1}^{n_2}(Y_j-\mu_2)^2} \sim F(n_1, n_2)
\end{equation*}
为了方便书写和解不等式，我们将 $W$ 重写为含有未知参数 $\sigma_1^2/\sigma_2^2$ 和纯样本数值的形式：
\begin{equation*}
W = \frac{\frac{1}{n_1}\sum_{i=1}^{n_1}(X_i-\mu_1)^2}{\frac{1}{n_2}\sum_{j=1}^{n_2}(Y_j-\mu_2)^2} \cdot \frac{\sigma_2^2}{\sigma_1^2} \sim F(n_1, n_2)
\end{equation*}
此时 $W$ 的分布完全确定，不含任何未知参数，是一个合适的枢轴量。

利用 $F$ 分布的分位点，我们有：
\begin{equation*}
P\left\{F_{1-\alpha/2}(n_1, n_2) < \frac{\frac{1}{n_1}\sum_{i=1}^{n_1}(X_i-\mu_1)^2}{\frac{1}{n_2}\sum_{j=1}^{n_2}(Y_j-\mu_2)^2} \cdot \frac{\sigma_2^2}{\sigma_1^2} < F_{\alpha/2}(n_1, n_2)\right\}=1-\alpha
\end{equation*}

对上述不等式的每一项取倒数（中间项的 $\sigma_2^2/\sigma_1^2$ 变为 $\sigma_1^2/\sigma_2^2$，且不等号方向反向）：
\begin{equation*}
P\left\{\frac{1}{F_{\alpha/2}(n_1, n_2)} < \frac{\frac{1}{n_2}\sum_{j=1}^{n_2}(Y_j-\mu_2)^2}{\frac{1}{n_1}\sum_{i=1}^{n_1}(X_i-\mu_1)^2} \cdot \frac{\sigma_1^2}{\sigma_2^2} < \frac{1}{F_{1-\alpha/2}(n_1, n_2)}\right\}=1-\alpha
\end{equation*}

同时乘以 $\frac{\frac{1}{n_1}\sum_{i=1}^{n_1}(X_i-\mu_1)^2}{\frac{1}{n_2}\sum_{j=1}^{n_2}(Y_j-\mu_2)^2}$ 使得中间只留下待估的方差比：
\begin{equation*}
P\left\{\frac{\frac{1}{n_1}\sum(X_i-\mu_1)^2}{\frac{1}{n_2}\sum(Y_j-\mu_2)^2} \cdot \frac{1}{F_{\alpha/2}(n_1, n_2)} < \frac{\sigma_1^2}{\sigma_2^2} < \frac{\frac{1}{n_1}\sum(X_i-\mu_1)^2}{\frac{1}{n_2}\sum(Y_j-\mu_2)^2} \cdot \frac{1}{F_{1-\alpha/2}(n_1, n_2)}\right\}=1-\alpha
\end{equation*}

最后，利用 $F$ 分布的倒数性质 $F_{1-\alpha/2}(n_1, n_2) = \frac{1}{F_{\alpha/2}(n_2, n_1)}$，可以将上限的右侧转化为更直观的形式。最终得到 $\sigma_1^2/\sigma_2^2$ 的置信区间为：
\begin{equation*}
\left(\frac{\frac{1}{n_1}\sum_{i=1}^{n_1}(X_i-\mu_1)^2}{\frac{1}{n_2}\sum_{j=1}^{n_2}(Y_j-\mu_2)^2}\frac{1}{F_{\alpha/2}(n_1,n_2)},\ \frac{\frac{1}{n_1}\sum_{i=1}^{n_1}(X_i-\mu_1)^2}{\frac{1}{n_2}\sum_{j=1}^{n_2}(Y_j-\mu_2)^2} F_{\alpha/2}(n_2,n_1)\right)
\end{equation*}
\end{solution}


\section{典型例题}

\begin{exercise} \label{ex:ci_length}
设总体 $X \sim N(\mu,\sigma^2)$（$\sigma^2$ 已知），则在给定样本容量 $n$ 及置信度 $1-\alpha$ 的情况下，未知参数 $\mu$ 的置信区间长度随着样本均值 $\overline{X}$ 的增加而（\qquad）。
\begin{enumerate}
  \item[(A)] 增加
  \item[(B)] 减少
  \item[(C)] 不变
  \item[(D)] 不能确定增或减
\end{enumerate}
\end{exercise}

\begin{solution}
应选 (C)。

在给定条件下，由表~\ref{tab:confidence_intervals} 可知，未知参数 $\mu$ 的置信区间为：
\begin{equation*}
\left( \overline{X} - \frac{\sigma}{\sqrt{n}} z_{\alpha/2},\ \overline{X} + \frac{\sigma}{\sqrt{n}} z_{\alpha/2} \right)
\end{equation*}
该置信区间的长度为
\(
L = 2\frac{\sigma}{\sqrt{n}} z_{\alpha/2}.
\)
显然，区间长度 \(L\) 只与总体标准差 \(\sigma\)、样本容量 \(n\) 以及对应置信水平的分位点 \(z_{\alpha/2}\) 有关，而与样本均值 \(\overline{X}\) 毫无关系。

也就是说，\(\overline{X}\) 的变化只会让整个区间左右平移，不会改变区间的宽度。因此，随着 \(\overline{X}\) 的变化，置信区间的长度保持不变。
\end{solution}

\begin{exercise} \label{ex:ci_large_sample}
设总体 $X$ 的方差为 $1$，根据来自 $X$ 的容量为 $100$ 的简单随机样本，测得样本均值为 $5$。已知标准正态分布函数值 $\Phi(1.96)=0.975$。求 $X$ 的数学期望 $\mu$ 的置信度近似等于 $0.95$ 的置信区间。
\end{exercise}

\begin{solution}
题目中并未说明总体 $X$ 是否为正态分布，但是由于样本容量 $n=100 \ge 30$，属于大样本。根据中心极限定理，样本均值 $\overline{X}$ 近似服从正态分布，即：
\begin{equation*}
Z = \frac{\overline{X}-\mu}{\sigma/\sqrt{n}} = \frac{\overline{X}-\mu}{1/\sqrt{100}} \stackrel{\text{近似}}{\sim} N(0,1)
\end{equation*}
又因为 $\Phi(1.96)=0.975$，可知 $z_{0.025} = 1.96$，故 $P\{-1.96 < Z < 1.96\} \approx 0.95$。

由此可得 $\mu$ 的近似置信区间为：
\begin{equation*}
\left( \overline{X} - \frac{1}{10} \times 1.96,\ \overline{X} + \frac{1}{10} \times 1.96 \right)
\end{equation*}
代入已知数据 $\overline{X}=5$，计算出置信区间的下限为 $4.804$，上限为 $5.196$。

所以 $\mu$ 的置信区间近似为 $(4.804,\ 5.196)$。
\end{solution}


\begin{exercise} \label{prob:ci_1}
对于正态总体 $X \sim N(\mu, \sigma^2)$，$\sigma^2$ \textbf{已知}，记 $L$ 是 $\mu$ 的置信度为 $0.95$ 的置信区间的长度，那么当样本容量 $n$ \underline{\qquad\qquad} 时，$L \leq \sigma$。\\
（参考数据：$z_{0.05}=1.645$，$z_{0.025}=1.96$，$t_{0.025}(15)=1.753$，$t_{0.025}(16)=2.12$）
\end{exercise}
\begin{solution}
因为 $\sigma^2$ 已知，使用正态分布作为枢轴量。置信区间长度为：
\begin{equation*}
L = 2 \cdot \frac{\sigma}{\sqrt{n}} z_{0.025}
\end{equation*}
代入 $z_{0.025}=1.96$，并令 $L \leq \sigma$：
\begin{equation*}
2 \cdot 1.96 \cdot \frac{\sigma}{\sqrt{n}} \leq \sigma \implies \frac{3.92}{\sqrt{n}} \leq 1 \implies \sqrt{n} \geq 3.92 \implies n \geq 15.3664
\end{equation*}
由于样本容量 $n$ 必须为正整数，故 $n \ge 16$。
\end{solution}

\begin{exercise} \label{prob:ci_2}
生产一个零件所需时间 $X \sim N(\mu, \sigma^2)$。现观察了 $25$ 个零件，测得样本均值 $\overline{X}=5.5$，样本标准差 $S=1.73$。求 $\mu$ 的置信度为 $0.95$ 的置信区间 \underline{\qquad\qquad}。\\
（参考数据：$t_{0.05}(24)=1.7109,\ t_{0.025}(24)=2.0639,\ \Phi(1.96)=0.975$）
\end{exercise}
\begin{solution}
由于总体方差 $\sigma^2$ 未知，需要使用 $t$ 分布。自由度 $df = n-1 = 24$，查表得 $t_{0.025}(24)=2.0639$。
置信区间公式为：
\begin{equation*}
\left( \overline{X}-\frac{S}{\sqrt{n}}t_{0.025}(24),\ \overline{X}+\frac{S}{\sqrt{n}}t_{0.025}(24) \right)
\end{equation*}
计算边际误差：
\begin{equation*}
\frac{1.73}{5} \times 2.0639 \approx 0.7141
\end{equation*}
下限为 $5.5 - 0.7141 = 4.7859$，上限为 $5.5 + 0.7141 = 6.2141$。故置信区间为 $(4.7859,\ 6.2141)$。
\end{solution}

\begin{exercise} \label{prob:ci_3}
设总体 $X \sim N(\mu, \sigma^2)$，有一个容量为 $10$ 的样本，测得样本方差 $S^2=2$。求 $\sigma^2$ 的置信水平为 $0.95$ 的置信区间 \underline{\qquad\qquad}。\\
（参考数据：$\chi^2_{0.975}(9)=2.70$，$\chi^2_{0.025}(9)=19.023$）
\end{exercise}
\begin{solution}
由于 $\mu$ 未知，求 $\sigma^2$ 的置信区间使用 $\chi^2$ 分布，自由度 $df = 10-1 = 9$。
方差的置信区间公式为：
\begin{equation*}
\left( \frac{(n-1)S^2}{\chi^2_{0.025}(9)},\ \frac{(n-1)S^2}{\chi^2_{0.975}(9)} \right)
\end{equation*}
计算分子：$(n-1)S^2 = 9 \times 2 = 18$。代入分母得到：
\begin{equation*}
\frac{18}{19.023} \approx 0.9462,\quad \frac{18}{2.70} \approx 6.6667
\end{equation*}
故置信区间为 $(0.9462,\ 6.6667)$。
\end{solution}

\begin{exercise} \label{prob:ci_4}
已知 $X\sim N(\mu,\sigma^2)$，样本容量 $n=9$，样本均值 $\overline{X}=6$，求 $\mu$ 在置信水平 $0.95$ 下的单侧置信下限 \underline{\qquad\qquad}。\\
（参考数据：$\Phi(1.96)=0.975,\ \Phi(1.645)=0.95$）
\end{exercise}
\begin{solution}
题目没有给出 \(\sigma\) 的具体数值，因此结果只能保留含 \(\sigma\) 的形式。又因为参考数据给的是标准正态分位点，所以这里对应的是“\(\sigma\) 已知时 \(\mu\) 的单侧置信区间”。

单侧置信下限的要求是
\(
P\{\mu \ge \underline{\mu}\}=0.95.
\)
利用标准正态分布的枢轴量，有：
\begin{equation*}
P\left( \frac{\overline{X}-\mu}{\sigma/\sqrt{n}} \le z_{0.05} \right) = 0.95
\end{equation*}
解得 $\mu \ge \overline{X} - z_{0.05}\frac{\sigma}{\sqrt{n}}$。代入数据 $\overline{X}=6, n=9, z_{0.05}=1.645$：
\begin{equation*}
\underline{\mu} = 6 - 1.645 \frac{\sigma}{3} \approx 6 - 0.5483\sigma
\end{equation*}
\end{solution}

\begin{exercise} \label{prob:ci_5}
对 $\alpha\in(0,1)$，由样本 $X_1,\dots,X_n$ 得到统计量 $\underline{\theta}$，若对一切参数 $\theta$ 均有 \underline{\qquad\qquad}，则称 $(\underline{\theta},+\infty)$ 为 $\theta$ 的置信水平为 $1-\alpha$ 的单侧置信区间，$\underline{\theta}$ 称为单侧置信下限。
\end{exercise}
\begin{solution}
应填
\(
P\{\theta>\underline{\theta}\}\ge 1-\alpha.
\)

理由直接来自单侧置信下限的定义。区间 \((\underline{\theta},+\infty)\) 要成为参数 \(\theta\) 的置信水平为 \(1-\alpha\) 的单侧置信区间，意思就是：由样本算出的下限 \(\underline{\theta}\) 不应把真值“压到区间外面去”，并且这种覆盖成功的概率至少要达到 \(1-\alpha\)。

换成事件语言，就是
\(
P\{\theta>\underline{\theta}\}\ge 1-\alpha.
\)
如果教材专门讨论连续型情形，也常把它写成等号
\(
P\{\theta>\underline{\theta}\}=1-\alpha,
\)
但从定义本身看，填空的标准形式应为“\(\ge 1-\alpha\)”。
\end{solution}



\subsection*{历年试题补充：2005 级概率试卷 A 卷}

\begin{exercise}\label{exer:src-2005-a-02}
测量某冶炼炉内的温度，重复测量 \(5\) 次，数据如下（单位：\({}^\circ\mathrm{C}\)）：
\[
1820,\quad 1834,\quad 1831,\quad 1816,\quad 1824.
\]
假定重复测量所得温度 \(X\sim N(\mu,\sigma^2)\)，并取 \(\sigma=10\)。求总体温度真值 \(\mu\) 的置信度为 \(0.95\) 的置信区间。\par
（参考数据：\(\Phi(1.96)=0.975,\ \Phi(1.65)=0.95\)。）
\end{exercise}
\begin{solution}
总体服从正态分布，且标准差 \(\sigma=10\) 已知，因此使用标准正态枢轴量
\[
Z=\frac{\overline{X}-\mu}{\sigma/\sqrt{n}}\sim N(0,1).
\]
由 \(\Phi(1.96)=0.975\)，得 \(z_{0.025}=1.96\)。样本均值为
\(
\overline{x}
=\frac{1820+1834+1831+1816+1824}{5}
=\frac{9125}{5}=1825.
\)
于是 \(\mu\) 的 \(0.95\) 置信区间为
\[
\left(
\overline{x}-1.96\frac{\sigma}{\sqrt n},
\overline{x}+1.96\frac{\sigma}{\sqrt n}
\right).
\]
代入 \(n=5,\ \sigma=10\)，得误差限
\(
1.96\frac{10}{\sqrt5}\approx 8.77.
\)
故所求置信区间为
\(
[1825-8.77,\ 1825+8.77]=[1816.23,\ 1833.77].
\)
所以总体温度真值 \(\mu\) 的 \(0.95\) 置信区间为
\(
[1816.23,\ 1833.77].
\)
\end{solution}

\begin{exercise} \label{prob:ci_6}
设总体 $X\sim N(\mu,\sigma^2)$，有一个容量 $n=9$ 的样本，已知 $\sum_{i=1}^9 X_i=198$，$\sum_{i=1}^9 X_i^2=4372$。求 $\mu$ 与 $\sigma$ 的置信水平为 $0.95$ 的双侧置信区间（结果保留四位小数）。\\
（参考数据：$t_{0.025}(8)=2.306$，$\chi^2_{0.975}(8)=2.180$，$\chi^2_{0.025}(8)=17.535$）
\end{exercise}
\begin{solution}
首先计算样本均值和样本方差：
\begin{equation*}
\overline{X} = \frac{1}{n}\sum_{i=1}^9 X_i = \frac{198}{9} = 22
\end{equation*}
\begin{equation*}
S^2 = \frac{1}{n-1}\left(\sum_{i=1}^9 X_i^2 - n\overline{X}^2\right) = \frac{1}{8}\left(4372 - 9 \times 22^2\right) = \frac{16}{8} = 2
\end{equation*}
于是样本标准差 $S = \sqrt{2} \approx 1.4142$。

求 $\mu$ 的置信区间（方差未知，用 $t$ 分布）：
自由度 $df=8$，分位点 $t_{0.025}(8)=2.306$。边际误差为 $\frac{\sqrt{2}}{3} \times 2.306 \approx 1.0871$。
下限：$22 - 1.0871 = 20.9129$，上限：$22 + 1.0871 = 23.0871$。
即 $\mu$ 的置信区间为 $(20.9129,\ 23.0871)$。

求 $\sigma$ 的置信区间：
先求出 $\sigma^2$ 的置信区间：
\begin{equation*}
\left( \frac{(n-1)S^2}{\chi^2_{0.025}(8)},\ \frac{(n-1)S^2}{\chi^2_{0.975}(8)} \right) = \left( \frac{16}{17.535},\ \frac{16}{2.180} \right) \approx (0.9125,\ 7.3395)
\end{equation*}
对上下限开平方，得到标准差 $\sigma$ 的置信区间：
\begin{equation*}
(\sqrt{0.9125},\ \sqrt{7.3395}) \approx (0.9552,\ 2.7091)
\end{equation*}
综上，$\mu$ 的置信区间为 $(20.9129,\ 23.0871)$，$\sigma$ 的置信区间为 $(0.9552,\ 2.7091)$。
\end{solution}

\begin{exercise}\label{exer:src-2009-3a-05}
假设一枚硬币抛了 \(400\) 次，结果只有 \(175\) 次正面。求正面的 \(95\%\) 置信区间，并在水平 \(\alpha=0.05\) 下说明这是否为均匀硬币。
\end{exercise}
\begin{solution}
设正面概率为 \(p\)，则样本正面频率
\(
\hat p=\frac{175}{400}=0.4375.
\)
由于样本容量较大，采用正态近似，取 \(z_{0.025}=1.96\)，则
\[
p\in\left(\hat p-1.96\sqrt{\frac{\hat p(1-\hat p)}{400}},\ \hat p+1.96\sqrt{\frac{\hat p(1-\hat p)}{400}}\right).
\]
代入 \(\hat p=0.4375\)，得
\(
\sqrt{\frac{0.4375\times0.5625}{400}}\approx0.0248,
\)
所以置信区间约为
\(
(0.4375-0.0486,\ 0.4375+0.0486)=(0.3889,\ 0.4861).
\)
因为 \(0.5\) 不在区间内，故在 \(\alpha=0.05\) 下可认为这不是均匀硬币。
\end{solution}

\begin{exercise}\label{exer:src-2009-3b-05}
假设一枚硬币抛了 \(600\) 次，结果只有 \(275\) 次正面。求正面的 \(95\%\) 置信区间，并在水平 \(\alpha=0.05\) 下说明这是否为均匀硬币。
\end{exercise}
\begin{solution}
设正面概率为 \(p\)，则样本正面频率
\(
\hat p=\frac{275}{600}\approx0.4583.
\)
由于样本容量较大，采用正态近似，取 \(z_{0.025}=1.96\)，得
\[
p\in\left(\hat p-1.96\sqrt{\frac{\hat p(1-\hat p)}{600}},\ \hat p+1.96\sqrt{\frac{\hat p(1-\hat p)}{600}}\right).
\]
代入 \(\hat p=275/600\)，有
\(
\sqrt{\frac{\hat p(1-\hat p)}{600}}\approx0.0203,
\)
所以置信区间约为
\(
(0.4583-0.0399,\ 0.4583+0.0399)=(0.4185,\ 0.4982).
\)
因为 \(0.5\) 不在区间内，故在 \(\alpha=0.05\) 下可认为这不是均匀硬币。
\end{solution}



\subsection*{历年试题补充：2008-01 3学分 A 卷}

\begin{exercise}\label{exer:src-2008-01-3a-iii-06}\label{exer:src-paper2-fill-06}\label{prob:p2-fill-06}
设某种保险丝熔化时间
\(
X\sim N(\mu,\sigma^2)
\)
（单位：秒），取 \(n=16\) 的样本，得样本均值和方差分别为
\(
\overline X=15,\qquad S^2=0.36.
\)
求 \(\mu\) 的置信度为 \(95\%\) 的单侧置信区间上限。
\end{exercise}
\begin{solution}
总体正态且 \(\sigma^2\) 未知，使用
\(
T=\frac{\overline X-\mu}{S/\sqrt n}\sim t(n-1).
\)
要求置信度为 \(0.95\) 的单侧上限 \(\overline\mu\)，即
\(
P(\mu<\overline\mu)=0.95.
\)
由
\(
P\left\{\frac{\overline X-\mu}{S/\sqrt n}>-t_{0.05}(15)\right\}=0.95
\)
可得
\(
\mu<\overline X+t_{0.05}(15)\frac{S}{\sqrt n}.
\)
题中
\(
S=\sqrt{0.36}=0.6,\qquad n=16.
\)
查表 \(t_{0.05}(15)=1.7531\)，故单侧置信上限为
\[
\overline\mu
=15+1.7531\cdot\frac{0.6}{4}
=15+0.2630
=15.2630.
\]
因此 \(\mu\) 的 \(95\%\) 单侧置信区间可写为
\(
(-\infty,\,15.2630).
\)
\end{solution}


\chapter{假设检验}

\begin{note}
\textbf{自学导读：区间估计的"反面"——用数据做决策。}

区间估计回答"参数大概在哪个范围"；假设检验回答"参数是否等于某个特定值"。两者用的数学工具完全相同（同样的枢轴量、同样的分布），只是问法不同。

\textbf{核心逻辑：}
\begin{enumerate}
    \item 先假设 $H_0$（原假设）为真——比如"$\mu=\mu_0$"；
    \item 在 $H_0$ 为真的前提下，算出检验统计量应该落在什么范围（接受域）；
    \item 如果实际算出的统计量落在接受域\textbf{之外}（拒绝域），说明"如果 $H_0$ 为真，出现这么极端的数据概率 $<\alpha$"——小概率事件发生了，所以拒绝 $H_0$。
\end{enumerate}

\textbf{自学提醒：}
\begin{itemize}
    \item 假设检验和区间估计是\textbf{一体两面}：$\mu_0$ 落在置信区间内 $\iff$ 不拒绝 $H_0:\mu=\mu_0$。
    \item 选统计量的规则和第 8 章完全一样：$\sigma$ 已知用 $Z$，未知用 $t$，检验方差用 $\chi^2$，比较两组方差用 $F$。
    \item "不拒绝 $H_0$"$\neq$"$H_0$ 为真"。只是说数据不够极端，无法推翻 $H_0$。
\end{itemize}
\end{note}

{\small
\[
\begin{cases}
\text{假设检验}
\begin{cases}
\text{基本思想: 随机试验中小概率事件一般不会发生} \\
\text{基本步骤: 提出假设} \Rightarrow \text{构建统计量} \Rightarrow \text{写出拒绝域} \Rightarrow \text{检验} \\
\text{主要检验: 双边检验或单边检验} \\
\text{显著性检验: 仅考虑控制第一类错误的假设检验，限定第一类错误发生的最大概率为} \alpha
\end{cases} \\
\text{单个正态总体}
\begin{cases}
\text{总体均值}
\begin{cases}
\text{总体方差已知, } Z\text{检验，也称} U\text{检验} \\
\text{总体方差未知, } t\text{检验}
\end{cases} \\
\text{总体方差}
\begin{cases}
\text{总体均值已知, } \chi^2(n)\text{检验} \\
\text{总体均值未知, } \chi^2(n-1)\text{检验}
\end{cases}
\end{cases} \\
\text{两个正态总体} \\
(X \sim N(\mu_1,\sigma_1^2), Y \sim N(\mu_2,\sigma_2^2))
\begin{cases}
\text{总体均值差} (\mu_1-\mu_2)
\begin{cases}
\sigma_1^2=\sigma_2^2=\sigma^2\text{已知, } Z\text{检验} \\
\sigma_1^2=\sigma_2^2=\sigma^2\text{未知, } t(n_1+n_2-2)\text{检验}
\end{cases} \\
\text{总体方差比} \left( \frac{\sigma_1^2}{\sigma_2^2} \right)
\begin{cases}
\mu_1,\mu_2\text{已知, } F(n_1,n_2)\text{检验} \\
\mu_1,\mu_2\text{未知, } F(n_1-1,n_2-1)\text{检验}
\end{cases}
\end{cases}
\end{cases}
\]
}

\begin{introduction}
    \item 原假设与备择假设~\ref{def:hypothesis-test}
    \item 两类错误与显著性水平~\ref{def:two-errors}
    \item 假设检验的基本步骤~\ref{sec:ht-steps}
    \item 正态总体参数的假设检验~\ref{sec:ht-normal}
    \item 置信区间与假设检验的关系~\ref{sec:ht-ci}
\end{introduction}

\begin{note}
\textbf{使用提醒：} 假设检验选方法时，不要先背拒绝域，先按四个问题定位：检验哪个参数、单样本还是双样本、总体方差已知还是未知、备择假设是左尾右尾还是双尾。前两问决定该用哪类统计量，第三问常常决定是 \(z\) 还是 \(t\)、是 \(\chi^2\) 还是 \(F\)，最后一问才决定拒绝域落在哪个尾部。
\end{note}

\section{基本概念}

\begin{definition}[假设检验] \label{def:hypothesis-test}
关于总体分布中的未知参数或分布类型的每一种论断称为\textbf{统计假设}。根据样本信息去推断该假设是否成立的统计推断问题称为\textbf{假设检验}。其中需要检验的假设称为\textbf{原假设}（或零假设），记为 \(H_0\)；与 \(H_0\) 对立的假设称为\textbf{备择假设}，记为 \(H_1\)。
\end{definition}

\begin{note}
假设检验的本质不是“证明 \(H_0\) 真”，而是“根据样本判断是否有足够证据拒绝 \(H_0\) ”。因此最后的结论通常只能说“拒绝 \(H_0\)”或“不能拒绝 \(H_0\)”，而不宜表述成“接受 \(H_0\) 一定正确”。
\end{note}

\begin{note}
\textbf{自学追问：} 为什么总要先设 \(H_0\)？因为检验需要在 \(H_0\) 下给出可计算的分布，从而控制第一类错误。结论也要只写“拒绝”或“不能拒绝”，不要把“不能拒绝”误写成“证明了 \(H_0\) 真”。
\end{note}

\begin{definition}[两类错误] \label{def:two-errors}
\begin{enumerate}
    \item \textbf{第一类错误（弃真）}：\(H_0\) 为真时，错误地拒绝了 \(H_0\)。
    \(
    \alpha = P\{\text{拒绝}\, H_0 \mid H_0\text{ 为真}\}.
    \)
    \item \textbf{第二类错误（取伪）}：\(H_0\) 不真时，错误地未拒绝 \(H_0\)。
    \(
    \beta = P\{\text{未拒绝}\, H_0 \mid H_1\text{ 为真}\}.
    \)
\end{enumerate}
\end{definition}

\begin{note}
两类错误不能同时都压到很小而又不付出代价。样本量固定时，如果把拒绝域设得更苛刻以减小第一类错误 \(\alpha\)，往往会增大第二类错误 \(\beta\)。所以假设检验本质上是在不同错误风险之间做权衡。
\end{note}

\begin{exercise}\label{exer:src-2008-01-3a-i-05}\label{exer:src-paper2-tf-05}\label{prob:p2-tf-05}
在一个确定的假设检验中，当样本容量确定时，犯第一类错误的概率与犯第二类错误的概率不能同时减少。判断正误。
\end{exercise}
\begin{solution}
正确。

在样本容量固定时，检验统计量的波动规模已经确定。若通过缩小拒绝域来减小第一类错误概率 \(\alpha\)，通常会增大在备择假设成立时未落入拒绝域的概率，即第二类错误概率 \(\beta\)。若扩大拒绝域以减小 \(\beta\)，又会增大 \(\alpha\)。

要在通常意义下同时降低两类错误概率，需要增加样本容量，使原假设和备择假设下统计量分布的重叠部分减小；仅在样本容量固定时调整检验规则，二者不能同时减少。
\end{solution}

\begin{definition}[显著性水平与显著性检验] \label{def:significance-level}
在假设检验中，只控制犯第一类错误的概率不超过 \(\alpha\)（\(\alpha\) 通常取 0.1, 0.05, 0.01），而不考虑第二类错误的概率，这种检验称为\textbf{显著性检验}，\(\alpha\) 称为\textbf{显著性水平}。
\end{definition}

\begin{note}
\(\alpha\) 不是“原假设为真的概率”，也不是“犯任何错误的总概率”。它专指：\textbf{在 \(H_0\) 真的前提下，却把 \(H_0\) 错误拒绝掉的最大容许概率}。做题时只要看到“显著性水平”，就要立刻联想到第一类错误。
\end{note}

\begin{note}
\textbf{使用提醒：} 把 \(\alpha\) 取得更小，意味着检验更谨慎、更不容易拒绝 \(H_0\)；这通常会降低第一类错误，却常常提高第二类错误。它不意味着”结论更真”，只意味着”拒绝 \(H_0\) 的门槛更高”。
\end{note}

\begin{note}
\textbf{$\alpha$ 与 $\beta$ 的跷跷板关系（直觉图景）：}

想象两条钟形曲线：一条以 $\mu_0$（$H_0$ 真时）为中心，另一条以 $\mu_1$（$H_1$ 真时）为中心，两条曲线有重叠区域。拒绝域是一条竖线（临界值）右边的区域。
\begin{itemize}
    \item 把竖线\textbf{右移}（$\alpha\downarrow$，更难拒绝）$\Longrightarrow$ 左边曲线尾巴变小（$\alpha$ 减小），但右边曲线被截掉的面积变大（$\beta$ 增大）。
    \item 把竖线\textbf{左移}（$\alpha\uparrow$，更容易拒绝）$\Longrightarrow$ $\alpha$ 增大，$\beta$ 减小。
\end{itemize}
\textbf{唯一能同时减小 $\alpha$ 和 $\beta$ 的方法：增大样本量 $n$}——这会让两条曲线都变窄（方差 $\sigma^2/n$ 减小），重叠区域缩小。

\textbf{考试常考判断题：}
\begin{itemize}
    \item “$\alpha$ 越小检验越好” $\times$（$\beta$ 可能很大）
    \item “增大 $n$ 可同时减小 $\alpha$ 和 $\beta$” \checkmark
    \item “$\alpha+\beta=1$” $\times$（两者没有这种简单关系）
\end{itemize}
\end{note}

\begin{note}
\textbf{小概率原理}是假设检验的理论基础：概率很小的事件在一次试验中认为不会发生。若在 \(H_0\) 成立的条件下，样本观测值导致某个概率 \(\leq \alpha\) 的事件发生，则拒绝 \(H_0\)。
\end{note}

\begin{definition}[双边检验与单边检验]\label{def:ht-types}
\begin{itemize}
    \item \textbf{双边检验}：\(H_0\colon \mu = \mu_0\)，\(H_1\colon \mu \neq \mu_0\)。
    \item \textbf{右边检验}：\(H_0\colon \mu \leq \mu_0\)，\(H_1\colon \mu > \mu_0\)。
    \item \textbf{左边检验}：\(H_0\colon \mu \geq \mu_0\)，\(H_1\colon \mu < \mu_0\)。
\end{itemize}
\end{definition}

\begin{note}
原假设 \(H_0\) 与备择假设 \(H_1\) 的地位不对等：\(H_0\) 是受保护的（控制弃真概率），因此选择 \(H_0\) 需要谨慎，一般遵循以下原则：
\begin{enumerate}
    \item 将后果较严重的错误设为第一类错误。例如检验药品真假，\(H_0\colon\) 药品为假，\(H_1\colon\) 药品为真（误判假药为真的后果更严重）。
    \item 若无严重后果，常选 \(H_0\) 为"维持现状"。例如检验新技术是否提高收益，\(H_0\colon\) 新技术未提高收益。
    \item 用\textbf{反证法}思想：若要验证 \(x\) 显著高于 \(y\)，则设 \(H_0\colon x \leq y\)，拒绝 \(H_0\) 后便有充分理由认为 \(x > y\)。
\end{enumerate}
\textbf{检验类型由备择假设决定}：若 \(H_1\) 写”\(\neq\)”就是双边检验，写”\(>\)”就是右边检验，写”\(<\)”就是左边检验。不要只盯着 \(H_0\) 看方向。
\end{note}

\begin{note}
\textbf{考试中如何快速确定 $H_0$ 和 $H_1$（三步法）：}
\begin{enumerate}
    \item \textbf{找”想证明的结论”。} 题目中”是否显著提高/降低/不同”所指向的方向，就是 $H_1$。因为假设检验的逻辑是反证法——把想证的放在 $H_1$，拒绝 $H_0$ 后才能下结论。
    \item \textbf{$H_0$ 取等号。} $H_0$ 永远包含”$=$”（$=$、$\le$、$\ge$），因为只有在等号成立时才能精确算出统计量的分布。
    \item \textbf{定方向。}
    \begin{itemize}
        \item 题目说”是否有显著差异/变化” $\Longrightarrow$ 双边：$H_0\colon\theta=\theta_0$，$H_1\colon\theta\neq\theta_0$；
        \item 题目说”是否显著提高/大于” $\Longrightarrow$ 右边：$H_0\colon\theta\le\theta_0$，$H_1\colon\theta>\theta_0$；
        \item 题目说”是否显著降低/小于” $\Longrightarrow$ 左边：$H_0\colon\theta\ge\theta_0$，$H_1\colon\theta<\theta_0$。
    \end{itemize}
\end{enumerate}
\textbf{常见陷阱：} 题目说”检验均值是否为 $\mu_0$”，很多同学误以为 $H_0\colon\mu\neq\mu_0$。记住：$H_0$ 永远是”没变化”的那一方，$H_1$ 才是”有变化”。
\end{note}

\section{假设检验的步骤} \label{sec:ht-steps}

\begin{property}[假设检验的四步流程]\label{prop:ht-steps}
\begin{enumerate}
    \item 根据问题提出 \(H_0\) 和 \(H_1\)，确定检验类型（双边/单边）。
    \item 给定显著性水平 \(\alpha\) 和样本容量 \(n\)。
    \item 选择\textbf{检验统计量}，确定\textbf{拒绝域}。
    \item 由样本计算统计量的观测值，判断是否落入拒绝域，作出结论。
\end{enumerate}
\end{property}

\begin{note}
\textbf{拒绝域的形式与备择假设 \(H_1\) 一致}，便于记忆：
\begin{itemize}
    \item \(H_1\colon \theta > \theta_0\)（右边检验）\(\Longrightarrow\) 拒绝域在\textbf{右尾}；
    \item \(H_1\colon \theta < \theta_0\)（左边检验）\(\Longrightarrow\) 拒绝域在\textbf{左尾}；
    \item \(H_1\colon \theta \neq \theta_0\)（双边检验）\(\Longrightarrow\) 拒绝域在\textbf{双尾}。
\end{itemize}
不能拒绝 \(H_0\) 并不意味着 \(H_0\) 一定为真，只是没有足够证据拒绝它。
\end{note}

\begin{note}
\textbf{假设检验与区间估计的统一视角：} 两者用的数学工具完全相同——同样的枢轴量、同样的分布表。区别仅在于"问法"：
\begin{itemize}
    \item 区间估计问"参数在哪个范围？" $\Longrightarrow$ 反解不等式得区间；
    \item 假设检验问"参数是否等于某值？" $\Longrightarrow$ 把该值代入枢轴量，看观测值是否落入小概率区域。
\end{itemize}
\textbf{快速判断法：} 若 $\mu_0$ 落在 $\mu$ 的 $1-\alpha$ 置信区间内，则在显著性水平 $\alpha$ 下不拒绝 $H_0\colon\mu=\mu_0$；反之拒绝。这是两章内容的桥梁。
\end{note}

\begin{exercise}\label{exer:src-paper3-i-07}\label{prob:p3-i-07}
在假设检验中，拒绝域的形式是根据备择假设 \(H_1\) 而确定的。判断正误。
\end{exercise}
\begin{solution}
正确。

检验统计量的类型由总体分布、参数和已知条件决定，而拒绝域落在哪一侧由备择假设决定：
\(
H_1\colon \theta>\theta_0
\)
对应右尾拒绝域；
\(
H_1\colon \theta<\theta_0
\)
对应左尾拒绝域；
\(
H_1\colon \theta\ne\theta_0
\)
对应双侧拒绝域。
\end{solution}

\begin{exercise}\label{exer:src-paper3-ii-05}\label{prob:p3-ii-05}
在以 \(H_0\) 为原假设、\(H_1\) 为备择假设的假设检验中，若显著性水平为 \(\alpha\)，则下列说法正确的是（\quad）。\par
A. \(P(\text{接受 }H_0\mid H_0\text{ 成立})=\alpha\)\quad
B. \(P(\text{接受 }H_1\mid H_1\text{ 成立})=\alpha\)\par
C. \(P(\text{接受 }H_1\mid H_0\text{ 成立})=\alpha\)\quad
D. \(P(\text{接受 }H_0\mid H_1\text{ 成立})=\alpha\)。
\end{exercise}
\begin{solution}
应选 C。

显著性水平 \(\alpha\) 是第一类错误概率，即在 \(H_0\) 成立时拒绝 \(H_0\) 的概率：
\(
\alpha=P(\text{拒绝 }H_0\mid H_0\text{ 成立}).
\)
在二择一的检验表述中，“接受 \(H_1\)”等价于“拒绝 \(H_0\)”，因此
\(
P(\text{接受 }H_1\mid H_0\text{ 成立})=\alpha.
\)
\end{solution}


\section{正态总体参数的假设检验} \label{sec:ht-normal}

\begin{exercise}[检验统计量的表达式] \label{ex:ht_stat}
设 \((X_1,\dots,X_n)\) 是来自 \(N(\mu,\sigma^2)\) 的样本（\(\mu,\sigma^2\) 未知），\(\overline{X}=\dfrac{1}{n}\displaystyle\sum_{i=1}^{n}X_i\)，\(Q^2=\displaystyle\sum_{i=1}^{n}(X_i-\overline{X})^2\)。检验 \(H_0\colon\mu=0\) 时，应选取的检验统计量为\underline{\qquad}（用 \(\overline{X},Q^2\) 表示）。
\end{exercise}
\begin{solution}
\textbf{先判题型。} 总体是正态总体，但 \(\sigma^2\) 未知，而题目检验的是均值 \(\mu\)，因此这里不能用 \(z\) 检验，而应使用单样本 \(t\) 检验。

在原假设 \(H_0\colon \mu=\mu_0\) 下，标准统计量应写成
\(
T=\frac{\overline{X}-\mu_0}{S/\sqrt{n}}.
\)
本题中 \(\mu_0=0\)，又由
\(
S^2=\frac{Q^2}{n-1},\qquad S=\frac{Q}{\sqrt{n-1}},
\)
代入即可把统计量改写成题目要求的 \(\overline{X},Q^2\) 形式：
\[
T=\frac{\overline{X}-\mu_0}{S/\sqrt{n}}\bigg|_{\mu_0=0}
=\frac{\overline{X}}{Q/\sqrt{n(n-1)}}=\frac{\overline{X}\sqrt{n(n-1)}}{Q}\sim t(n-1).
\]
\(
\text{故应选取的检验统计量为}\quad T=\frac{\overline{X}\sqrt{n(n-1)}}{Q}.
\)
\end{solution}


以下汇总了单正态总体与双正态总体在不同条件下的常用检验统计量及其拒绝域。
\begin{table}[htbp]
\centering
\caption{正态总体均值与方差的假设检验（显著性水平为 $\alpha$）} \label{tab:hypothesis_tests}
\resizebox{\textwidth}{!}{
\begin{tabular}{lllll}
\toprule
\textbf{参数} & \textbf{条件} & \textbf{原假设 $H_0$} & \textbf{检验统计量} & \textbf{拒绝域（对应右单边、左单边、双边）} \\
\midrule
\multirow{3}{*}{$\mu$} & \multirow{3}{*}{$\sigma^2$ 已知} & $\mu \leq \mu_0$ ($H_1: \mu > \mu_0$) & \multirow{3}{*}{$Z = \frac{\overline{X}-\mu_0}{\sigma/\sqrt{n}} \sim N(0,1)$} & $z \geq z_\alpha$ \\
& & $\mu \geq \mu_0$ ($H_1: \mu < \mu_0$) & & $z \leq -z_\alpha$ \\
& & $\mu = \mu_0$ ($H_1: \mu \neq \mu_0$) & & $|z| \geq z_{\frac{\alpha}{2}}$ \\
\midrule
\multirow{3}{*}{$\mu$} & \multirow{3}{*}{$\sigma^2$ 未知} & $\mu \leq \mu_0$ ($H_1: \mu > \mu_0$) & \multirow{3}{*}{$t = \frac{\overline{X}-\mu_0}{S/\sqrt{n}} \sim t(n-1)$} & $t \geq t_\alpha(n-1)$ \\
& & $\mu \geq \mu_0$ ($H_1: \mu < \mu_0$) & & $t \leq -t_\alpha(n-1)$ \\
& & $\mu = \mu_0$ ($H_1: \mu \neq \mu_0$) & & $|t| \geq t_{\frac{\alpha}{2}}(n-1)$ \\
\midrule
\multirow{3}{*}{$\mu_1-\mu_2$} & \multirow{3}{*}{$\sigma_1^2,\sigma_2^2$ 已知} & $\mu_1-\mu_2 \leq \delta$ & \multirow{3}{*}{$Z = \frac{(\overline{X}-\overline{Y})-\delta}{\sqrt{\frac{\sigma_1^2}{n_1}+\frac{\sigma_2^2}{n_2}}} \sim N(0,1)$} & $z \geq z_\alpha$ \\
& & $\mu_1-\mu_2 \geq \delta$ & & $z \leq -z_\alpha$ \\
& & $\mu_1-\mu_2 = \delta$ & & $|z| \geq z_{\frac{\alpha}{2}}$ \\
\midrule
\multirow{3}{*}{$\mu_1-\mu_2$} & \multirow{3}{*}{$\sigma_1^2=\sigma_2^2$ 未知} & $\mu_1-\mu_2 \leq \delta$ & \multirow{3}{*}{$t = \frac{(\overline{X}-\overline{Y})-\delta}{S_w\sqrt{\frac{1}{n_1}+\frac{1}{n_2}}} \sim t(n_1+n_2-2)$} & $t \geq t_\alpha(n_1+n_2-2)$ \\
& & $\mu_1-\mu_2 \geq \delta$ & & $t \leq -t_\alpha(n_1+n_2-2)$ \\
& & $\mu_1-\mu_2 = \delta$ & & $|t| \geq t_{\frac{\alpha}{2}}(n_1+n_2-2)$ \\
\midrule
\multirow{3}{*}{$\sigma^2$} & \multirow{3}{*}{$\mu$ 未知} & $\sigma^2 \leq \sigma_0^2$ & \multirow{3}{*}{$\chi^2 = \frac{(n-1)S^2}{\sigma_0^2} \sim \chi^2(n-1)$} & $\chi^2 \geq \chi^2_\alpha(n-1)$ \\
& & $\sigma^2 \geq \sigma_0^2$ & & $\chi^2 \leq \chi^2_{1-\alpha}(n-1)$ \\
& & $\sigma^2 = \sigma_0^2$ & & $\chi^2 \geq \chi^2_{\frac{\alpha}{2}}(n-1)$ 或 $\chi^2 \leq \chi^2_{1-\frac{\alpha}{2}}(n-1)$ \\
\midrule
\multirow{3}{*}{$\sigma_1^2 / \sigma_2^2$} & \multirow{3}{*}{$\mu_1,\mu_2$ 未知} & $\sigma_1^2 \leq \sigma_2^2$ & \multirow{3}{*}{$F = \frac{S_1^2}{S_2^2} \sim F(n_1-1,n_2-1)$} & $F \geq F_\alpha(n_1-1,n_2-1)$ \\
& & $\sigma_1^2 \geq \sigma_2^2$ & & $F \leq F_{1-\alpha}(n_1-1,n_2-1)$ \\
& & $\sigma_1^2 = \sigma_2^2$ & & $F \geq F_{\frac{\alpha}{2}}(n_1-1,n_2-1)$ 或 $F \leq F_{1-\frac{\alpha}{2}}(n_1-1,n_2-1)$ \\
\bottomrule
\end{tabular}
}
\end{table}
\begin{remark} \label{rem:rejection_region}
\textbf{记忆捷径}：仔细观察上表可知，拒绝域的不等号方向总是与备择假设$H_1$的形式（不等号方向）保持一致。
\end{remark}

\begin{exercise}\label{exer:src-paper1-tf-07}\label{prob:p1-tf-07}
在参数的假设检验中，拒绝域的形式是根据备择假设 \(H_1\) 而确定的。判断正误。
\end{exercise}
\begin{solution}
正确。

显著性水平 \(\alpha\) 决定拒绝域的概率大小，而备择假设 \(H_1\) 决定拒绝域放在哪个方向：若 \(H_1\) 为“参数大于某值”，拒绝域在右尾；若 \(H_1\) 为“小于某值”，拒绝域在左尾；若 \(H_1\) 为“不等于某值”，拒绝域在两端。因此拒绝域形式要依据 \(H_1\) 来定。
\end{solution}

\begin{note}
\textbf{检验统计量选取的决策流程（与区间估计完全对称）：}
\begin{enumerate}
    \item \textbf{检什么？} 检验 $\mu$ 还是 $\sigma^2$（或两总体的差/比）。
    \item \textbf{知什么？} 与区间估计相同的判断：
    \begin{itemize}
        \item 检 $\mu$ 且 $\sigma^2$ 已知 $\Longrightarrow$ $Z$ 检验；
        \item 检 $\mu$ 且 $\sigma^2$ 未知 $\Longrightarrow$ $t$ 检验；
        \item 检 $\sigma^2$（$\mu$ 未知）$\Longrightarrow$ $\chi^2$ 检验；
        \item 检两总体方差比 $\Longrightarrow$ $F$ 检验。
    \end{itemize}
    \item \textbf{代 $H_0$ 的值。} 把 $H_0$ 中的参数值代入枢轴量，得到检验统计量。
    \item \textbf{看 $H_1$ 定拒绝域方向。}
\end{enumerate}
口诀：\textbf{区间估计反着来——代入 $\theta_0$ 后看落不落在尾巴上}。
\end{note}

\subsection{单个正态总体均值的检验}

\begin{property}[Z 检验（\(\sigma^2\) 已知）]\label{prop:z-test}
检验统计量 \(Z = \dfrac{\overline{X}-\mu_0}{\sigma/\sqrt{n}} \sim N(0,1)\)。

\begin{center}
\begin{tabular}{ccc}
\toprule
\(H_1\) & 拒绝域 \\
\midrule
\(\mu > \mu_0\) & \(z \geq z_\alpha\) \\
\(\mu < \mu_0\) & \(z \leq -z_\alpha\) \\
\(\mu \neq \mu_0\) & \(|z| \geq z_{\alpha/2}\) \\
\bottomrule
\end{tabular}
\end{center}
\end{property}

\begin{exercise}[辛烷等级的左边 \texorpdfstring{\(Z\)}{Z} 检验] \label{ex:ht_z_left}
一种燃料的辛烷等级服从正态分布，平均值为 98，标准差为 0.8。从一批新燃料中随机抽取 25 桶，样本均值为 97.7。假定新燃料辛烷等级标准差仍为 0.8，问新燃料辛烷等级平均值是否比原燃料偏低（\(\alpha=0.05\)）？
（\(z_{0.05}=1.64\)）
\end{exercise}
\begin{solution}
\(H_0\colon\mu\geq98\)，\(H_1\colon\mu<98\)（左边检验，\(\alpha=0.05\)）。
\(\sigma=0.8\) 已知，使用 \(Z\) 检验：
\[
Z=\frac{\overline{X}-98}{0.8/\sqrt{25}}\sim N(0,1),\quad\text{拒绝域}\ z\leq -z_{0.05}=-1.64.
\]
\(
z=\frac{97.7-98}{0.8/5}=\frac{-0.3}{0.16}=-1.875<-1.64.
\)
落入拒绝域，拒绝 \(H_0\)，认为新燃料辛烷等级显著偏低。
\end{solution}



\begin{exercise}\label{exer:src-2004-3a-09}
在正态总体 \(N(\mu,1)\) 中抽取容量为 \(100\) 的样本，经计算得样本均值的观测值 \(\overline{x}=5.32\)。试在显著性水平 \(\alpha=0.01\) 下，检验假设
\(
H_0:\mu=5,\qquad H_1:\mu\ne 5,
\)
其中 \(\Phi(2.57)=0.995,\ \Phi(2.33)=0.99\)。
\end{exercise}
\begin{solution}
这是单个正态总体均值的双边检验，且总体方差已知为 \(\sigma^2=1\)。在 \(H_0:\mu=5\) 成立时，检验统计量为
\(
Z=\frac{\overline{X}-5}{1/\sqrt{100}}\sim N(0,1).
\)

显著性水平 \(\alpha=0.01\)，双边检验的拒绝域为
\(
|Z|\ge z_{\alpha/2}=z_{0.005}.
\)
由 \(\Phi(2.57)=0.995\)，得
\(
z_{0.005}=2.57.
\)

代入样本均值：
\(
z=\frac{5.32-5}{1/\sqrt{100}}
=\frac{0.32}{0.1}
=3.2.
\)
由于
\(
|z|=3.2>2.57,
\)
统计量落入拒绝域，因此拒绝 \(H_0\)。在显著性水平 \(0.01\) 下，可以认为总体均值 \(\mu\) 与 \(5\) 有显著差异。
\end{solution}

\begin{exercise}\label{exer:src-2008-01-3a-iv-05a}\label{exer:src-paper2-app-05a}\label{prob:p2-app-05a}
某厂生产的金属丝折断力 \(X\sim N(\mu,\sigma^2)\)，已知 \(\sigma=8\) kg。现抽取 \(10\) 根，测得平均折断力 \(\overline{x}=575.2\) kg。问在显著性水平 \(\alpha=0.05\) 下，能否认为这批金属丝的平均折断力为 \(570\) kg？
（\(z_{0.025}=1.96\)）
\end{exercise}
\begin{solution}
检验问题为
\[
H_0\colon \mu=570,\qquad H_1\colon \mu\neq570.
\]
总体标准差 \(\sigma=8\) 已知，采用双边 \(Z\) 检验。在 \(H_0\) 成立时，
\(
Z=\frac{\overline{X}-570}{8/\sqrt{10}}\sim N(0,1).
\)
显著性水平为 \(0.05\)，拒绝域为
\(
|z|\ge z_{0.025}=1.96.
\)
代入样本均值：
\[
z=\frac{575.2-570}{8/\sqrt{10}}
=\frac{5.2\sqrt{10}}{8}
\approx 2.055.
\]
由于 \(2.055>1.96\)，统计量落入拒绝域，故拒绝 \(H_0\)。在 \(5\%\) 显著性水平下，不能认为这批金属丝的平均折断力为 \(570\) kg。
\end{solution}

\begin{exercise}\label{exer:9-z-test}
某工厂生产灯泡的寿命 \(X \sim N(\mu, 200^2)\)。过去灯泡平均寿命为 1500 小时，采用新工艺后抽取 25 只，测得平均寿命为 1675 小时。问在显著性水平 \(\alpha = 0.05\) 下，新工艺是否显著提高了灯泡寿命？（\(z_{0.05}=1.65\)）
\end{exercise}
\begin{solution}
\textbf{第一步：提出假设。}问题问"是否显著提高"，故用右边检验：
\[
H_0\colon \mu \leq 1500,\quad H_1\colon \mu > 1500.
\]
\textbf{第二步：选择检验统计量。}总体方差 \(\sigma^2=200^2\) 已知，使用 \(Z\) 检验：
\[
Z = \frac{\overline{X}-\mu_0}{\sigma/\sqrt{n}} = \frac{\overline{X}-1500}{200/\sqrt{25}} \sim N(0,1).
\]
\textbf{第三步：确定拒绝域。}右边检验，\(\alpha=0.05\)，拒绝域为 \(z \geq z_{0.05}=1.65\)。\\
\textbf{第四步：计算并判断。}将样本数据代入：
\(
z = \frac{1675-1500}{200/5} = \frac{175}{40} = 4.375.
\)
\(4.375 > 1.65\)，\(z\) 落入拒绝域，拒绝 \(H_0\)，即在 5\% 显著性水平下认为新工艺显著提高了灯泡寿命。
\end{solution}


\begin{property}[t 检验（\(\sigma^2\) 未知）]\label{prop:t-test}
检验统计量 \(t = \dfrac{\overline{X}-\mu_0}{S/\sqrt{n}} \sim t(n-1)\)。
\begin{center}
\begin{tabular}{ccc}
\toprule
\(H_1\) & 拒绝域 \\
\midrule
\(\mu > \mu_0\) & \(t \geq t_\alpha(n-1)\) \\
\(\mu < \mu_0\) & \(t \leq -t_\alpha(n-1)\) \\
\(\mu \neq \mu_0\) & \(|t| \geq t_{\alpha/2}(n-1)\) \\
\bottomrule
\end{tabular}
\end{center}
\end{property}


\begin{exercise}\label{exer:9-t-test-two}
某厂生产灯泡寿命 \(X \sim N(\mu, \sigma^2)\)（\(\mu, \sigma^2\) 均未知），抽取 25 只测得 \(\overline{x}=1700\)，\(s=490\)。在 \(\alpha=0.01\) 下，能否认为平均寿命为 2000 小时？
（\(\;t_{0.005}(24)=2.7969\)）
\end{exercise}
\begin{solution}
\textbf{提出假设。}问题问"能否认为平均寿命为 2000"，属于双边检验：
\[
H_0\colon \mu = 2000,\quad H_1\colon \mu \neq 2000.
\]
\textbf{选择检验统计量。}\(\sigma^2\) 未知，使用 \(t\) 检验：
\[
t = \frac{\overline{X}-\mu_0}{S/\sqrt{n}} \sim t(n-1) = t(24).
\]
\textbf{确定拒绝域。}双边检验，\(\alpha=0.01\)，拒绝域为 \(|t| \geq t_{\alpha/2}(24) = t_{0.005}(24)=2.7969\)。\\
\textbf{计算并判断。}已知 \(\overline{x}=1700,\ s=490,\ n=25\)：
\[
t = \frac{1700-2000}{490/\sqrt{25}} = \frac{-300}{98} = -3.061,\quad |t| = 3.061.
\]
\(|t| = 3.061 > 2.7969\)，落入拒绝域，拒绝 \(H_0\)，即在 1\% 显著性水平下不能认为平均寿命为 2000 小时。
\end{solution}


\begin{exercise}\label{exer:9-t-test-right}
某高校某专业 25 名学生高数成绩服从正态分布，样本均值 \(\overline{x}=69.5\)，标准差 \(s=15\)。在 \(\alpha=0.05\) 下，检验该专业平均成绩是否\textbf{显著高于} 70 分。（\(t_{0.05}(24)=1.711\)）
\end{exercise}
\begin{solution}
\textbf{提出假设。}问题问"是否\textbf{显著高于} 70 分"，故用右边检验：
\(
H_0\colon \mu \leq 70,\quad H_1\colon \mu > 70.
\)
\textbf{选择检验统计量。}\(\sigma^2\) 未知，使用 \(t\) 检验：
\(
t = \frac{\overline{X}-\mu_0}{S/\sqrt{n}} \sim t(24).
\)
\textbf{确定拒绝域。}右边检验，\(\alpha=0.05\)，拒绝域为 \(t \geq t_{0.05}(24)=1.711\)。\\
\textbf{计算并判断。}已知 \(\overline{x}=69.5,\ s=15,\ n=25\)：
\(
t = \frac{69.5-70}{15/\sqrt{25}} = \frac{-0.5}{3} = -0.167.
\)
\(-0.167 < 1.711\)，未落入拒绝域（实际上统计量为负，远小于正的临界值），故不能拒绝 \(H_0\)，即在 5\% 显著性水平下没有充分证据认为平均成绩显著高于 70 分。
\begin{note}
从直觉上看，样本均值 \(\overline{x}=69.5\) 反而\textbf{低于} 70 分，自然不可能得出"显著高于"的结论。这提醒我们：在使用单边检验时，应先看样本数据的方向是否与 \(H_1\) 一致；若方向相反，通常就已经说明样本没有支持 \(H_1\) 的趋势。
\end{note}
\end{solution}

\begin{exercise}[考试成绩均值的 \texorpdfstring{\(t\)}{t} 区间与检验] \label{ex:ht_t_two}\label{exer:src-probstat-test1-04}
设某次考试成绩服从正态分布 \(N(\mu,\sigma^2)\)，随机抽取 \(36\) 位考生，算得 \(\overline{x}=66.5\) 分，修正样本标准差 \(s=15\)。求：\par
(1) \(\mu\) 的 \(95\%\) 置信区间；\par
(2) 在 \(\alpha=0.05\) 下，是否可以认为全体考生平均成绩为 \(70\) 分？\par
已知 \(t_{0.025}(35)=2.0301\)。
\end{exercise}
\begin{solution}
(1) 总体正态且 \(\sigma^2\) 未知，使用
\(
T=\frac{\overline X-\mu}{S/\sqrt n}\sim t(n-1).
\)
这里 \(n=36,\ \overline x=66.5,\ s=15\)，故 \(\mu\) 的 \(95\%\) 置信区间为
\(
\overline x\pm t_{0.025}(35)\frac{s}{\sqrt n}.
\)
代入数值得误差限
\[
2.0301\cdot\frac{15}{\sqrt{36}}
=2.0301\cdot2.5
=5.07525.
\]
因此
\[
\mu\in(66.5-5.07525,\ 66.5+5.07525)
\approx(61.42,\ 71.58).
\]

(2) 检验
\[
H_0\colon\mu=70,\qquad H_1\colon\mu\neq70
\]
（双边 \(t\) 检验，\(\alpha=0.05\)）。在 \(H_0\) 成立时，
\[
t=\frac{\overline{X}-70}{S/\sqrt{n}}\sim t(35),\quad\text{拒绝域}\ |t|\geq t_{0.025}(35)=2.0301.
\]
\(
t=\frac{66.5-70}{15/\sqrt{36}}=\frac{-3.5}{2.5}=-1.4,\quad |t|=1.4<2.0301.
\)
未落入拒绝域，故不能拒绝 \(H_0\)；在该显著性水平下，没有充分证据认为平均成绩偏离 70 分。
\end{solution}

\begin{exercise}[电池寿命方差的双边 \texorpdfstring{\(\chi^2\)}{χ²} 检验] \label{ex:ht_chi2_battery}
某厂电池寿命 \(X\sim N(\mu,\sigma^2)\)，正常生产时 \(\sigma^2=5000\)。现因生产条件波动较大，抽取 26 只电池测得样本方差 \(s^2=7200\)。能否判断寿命波动性有显著变化（\(\alpha=0.02\)）？
（\(\chi^2_{0.01}(25)=44.314\)，\(\chi^2_{0.99}(25)=11.524\)）
\end{exercise}
\begin{solution}
\(H_0\colon\sigma^2=5000\)，\(H_1\colon\sigma^2\neq5000\)（双边 \(\chi^2\) 检验，\(\alpha=0.02\)）。
\(
\chi^2=\frac{(n-1)s^2}{\sigma_0^2}=\frac{25\times7200}{5000}=36.
\)
\textbf{不拒绝域}为
\(
[\,\chi^2_{0.99}(25),\ \chi^2_{0.01}(25)\,]=[11.524,\ 44.314].
\)
\(11.524<36<44.314\)，未落入拒绝域，故不能拒绝 \(H_0\)，即暂时没有充分证据判断波动性有显著变化。
\end{solution}


\begin{exercise}[机床零件标准差的 \texorpdfstring{\(\chi^2\)}{χ²} 检验] \label{ex:ht_chi2_two}
某机床加工零件尺寸 \(X\sim N(\mu,\sigma^2)\)，抽取 10 个零件测得长度为
\(
5.5,\ 6.0,\ 6.2,\ 5.8,\ 5.6,\ 6.0,\ 6.4,\ 6.4,\ 6.0,\ 6.1.
\)
能否认为标准差为 0.4（\(\alpha=0.05\)）？
（\(\chi^2_{0.025}(9)=19.023\)，\(\chi^2_{0.975}(9)=2.700\)）
\end{exercise}
\begin{solution}
\(H_0\colon\sigma^2=0.16\)，\(H_1\colon\sigma^2\neq0.16\)（双边 \(\chi^2\) 检验）。\\
\(\overline{x}=60.0/10=6.0\)，离差平方和：
\(
\sum_{i=1}^{10}(x_i-6.0)^2=0.25+0+0.04+0.04+0.16+0+0.16+0.16+0+0.01=0.82.
\)
\(
\chi^2=\frac{(n-1)s^2}{\sigma_0^2}=\frac{0.82}{0.16}=5.125.
\)
拒绝域：\(\chi^2\leq 2.700\) 或 \(\chi^2\geq 19.023\)。
\(2.700<5.125<19.023\)，未落入拒绝域，故不能拒绝 \(H_0\)；在该显著性水平下，没有充分证据认为标准差偏离 \(0.4\)。
\end{solution}

\begin{exercise}\label{exer:src-paper1-app-06}\label{prob:p1-app-06}
某冶金实验室对锰的熔化点作了四次试验，结果分别为
\[
1269^\circ\mathrm C,\quad 1271^\circ\mathrm C,\quad 1263^\circ\mathrm C,\quad 1265^\circ\mathrm C.
\]
设数据服从正态分布 \(N(\mu,\sigma^2)\)，以 \(\alpha=5\%\) 的水平作如下检验：\par
(1) 这些结果是否符合公布的数字 \(1260^\circ\mathrm C\)？\par
(2) 测定值的标准差是否不超过 \(2^\circ\mathrm C\)？\par
已知 \(t_{0.025}(3)=3.1824,\ \chi^2_{0.05}(3)=7.815\)。
\end{exercise}
\begin{solution}
样本均值为
\(
\overline x=\frac{1269+1271+1263+1265}{4}=1267.
\)
修正样本方差为
\[
S^2=\frac1{3}\sum_{i=1}^{4}(x_i-\overline x)^2
=\frac{40}{3},
\qquad
S=\sqrt{\frac{40}{3}}\approx3.65.
\]

(1) 检验
\(
H_0:\mu=1260,\qquad H_1:\mu\ne1260.
\)
总体正态且 \(\sigma^2\) 未知，采用双边 \(t\) 检验：
\(
T=\frac{\overline X-1260}{S/\sqrt4}\sim t(3).
\)
代入样本值，
\(
|T_0|=\frac{1267-1260}{3.65/2}\approx3.836.
\)
由于
\(
3.836>t_{0.025}(3)=3.1824,
\)
故拒绝 \(H_0\)。在 \(5\%\) 显著性水平下，不能认为这些结果符合公布的 \(1260^\circ\mathrm C\)。

(2) 检验
\(
H_0:\sigma\le2,\qquad H_1:\sigma>2.
\)
取边界值 \(\sigma_0=2\)，检验统计量为
\(
\chi^2=\frac{(n-1)S^2}{\sigma_0^2}\sim\chi^2(3).
\)
样本值为
\(
\chi_0^2=\frac{3\cdot(40/3)}{4}=10.
\)
右边检验的拒绝域为
\(
\chi^2>\chi^2_{0.05}(3)=7.815.
\)
由于 \(10>7.815\)，故拒绝 \(H_0\)。在 \(5\%\) 显著性水平下，有充分证据认为测定值的标准差超过 \(2^\circ\mathrm C\)。
\end{solution}


\subsection{两个正态总体均值差的检验}

\begin{property}[两个正态总体均值差的检验]\label{prop:two-sample-mean-test}
设 \(X_1,\dots,X_{n_1}\) 来自 \(N(\mu_1,\sigma_1^2)\)，\(Y_1,\dots,Y_{n_2}\) 来自 \(N(\mu_2,\sigma_2^2)\)，两组独立。\\
\textbf{情形1}：\(\sigma_1^2, \sigma_2^2\) 已知。检验统计量
\[
Z = \frac{(\overline{X}-\overline{Y})-\delta}{\sqrt{\sigma_1^2/n_1+\sigma_2^2/n_2}} \sim N(0,1).
\]
\textbf{情形2}：\(\sigma_1^2=\sigma_2^2=\sigma^2\) 未知。检验统计量
\[
t = \frac{(\overline{X}-\overline{Y})-\delta}{S_w\sqrt{1/n_1+1/n_2}} \sim t(n_1+n_2-2),
\]
其中 \(S_w^2 = \dfrac{(n_1-1)S_1^2+(n_2-1)S_2^2}{n_1+n_2-2}\) 为\textbf{合并样本方差}。\\
两种情形下，\(H_1\colon \mu_1-\mu_2 > \delta\) 对应右尾拒绝域，\(H_1\colon \mu_1-\mu_2 < \delta\) 对应左尾，\(H_1\colon \mu_1-\mu_2 \neq \delta\) 对应双尾。
\end{property}

\subsection{单个正态总体方差的检验}

\begin{property}[$\chi^2$ 检验]\label{prop:chi2-test}
\textbf{\(\mu\) 未知}时，检验统计量
\(
\chi^2 = \frac{(n-1)S^2}{\sigma_0^2} \sim \chi^2(n-1).
\)
\textbf{\(\mu\) 已知}时，检验统计量
\[
\chi^2 = \frac{\sum_{i=1}^{n}(X_i-\mu)^2}{\sigma_0^2} \sim \chi^2(n).
\]
\begin{center}
\begin{tabular}{cc}
\toprule
\(H_1\) & 拒绝域（\(\mu\) 未知） \\
\midrule
\(\sigma^2 > \sigma_0^2\) & \(\chi^2 \geq \chi^2_\alpha(n-1)\) \\
\(\sigma^2 < \sigma_0^2\) & \(\chi^2 \leq \chi^2_{1-\alpha}(n-1)\) \\
\(\sigma^2 \neq \sigma_0^2\) & \(\chi^2 \geq \chi^2_{\alpha/2}(n-1)\) 或 \(\chi^2 \leq \chi^2_{1-\alpha/2}(n-1)\) \\
\bottomrule
\end{tabular}
\end{center}
\end{property}

\begin{exercise}\label{exer:9-chi2-two}
某厂铜线折断力 \(X \sim N(\mu, \sigma^2)\)，抽查 10 根，\(\overline{x}=575.2\)，\(s^2=68.16\)。在 \(\alpha=0.05\) 下检验 \(H_0\colon \sigma^2=64\)，\(H_1\colon \sigma^2\neq64\)。
（\(\chi^2_{0.025}(9)=19.0\)，\(\chi^2_{0.975}(9)=2.7\)）
\end{exercise}
\begin{solution}
\textbf{提出假设。}检验方差是否为 64，使用双边检验：
\[
H_0\colon \sigma^2 = 64,\quad H_1\colon \sigma^2 \neq 64.
\]
\textbf{选择检验统计量。}\(\mu\) 未知，使用 \(\chi^2\) 检验：
\(
\chi^2 = \frac{(n-1)S^2}{\sigma_0^2} = \frac{(10-1)S^2}{64} \sim \chi^2(9).
\)
\textbf{确定拒绝域。}双边检验，\(\alpha=0.05\)，拒绝域为：
\[
\chi^2 \leq \chi^2_{0.975}(9) = 2.7 \quad\text{或}\quad \chi^2 \geq \chi^2_{0.025}(9) = 19.0.
\]
不拒绝域为 \((2.7,\ 19.0)\)。\\
\textbf{计算并判断。}
\(
\chi^2 = \frac{9 \times 68.16}{64} = \frac{613.44}{64} = 9.585.
\)
\(2.7 < 9.585 < 19.0\)，未落入拒绝域，故不能拒绝 \(H_0\)，即在 5\% 显著性水平下没有充分证据认为这批铜线折断力的方差偏离 64。
\end{solution}

\begin{exercise}\label{exer:src-2008-01-3a-iv-05b}\label{exer:src-paper2-app-05b}\label{prob:p2-app-05b}
在正常条件下，某种维尼纶纤度服从 \(N(\mu,0.048^2)\)。现抽取 \(5\) 根测得纤度为
\(
1.31,\ 1.55,\ 1.34,\ 1.40,\ 1.45.
\)
问在显著性水平 \(\alpha=0.10\) 下，能否认为该纤度的方差仍属正常？
（\(\chi^2_{0.05}(4)=9.488,\ \chi^2_{0.95}(4)=0.711\)）
\end{exercise}
\begin{solution}
要检验方差是否仍为正常值 \(0.048^2\)，作双边检验：
\[
H_0\colon \sigma^2=0.048^2,\qquad H_1\colon \sigma^2\neq0.048^2.
\]
总体均值未知，检验统计量为
\[
\chi^2=\frac{(n-1)S^2}{\sigma_0^2}
=\frac{\sum_{i=1}^{n}(X_i-\overline{X})^2}{0.048^2}
\sim \chi^2(4).
\]
样本均值
\(
\overline{x}=\frac{1.31+1.55+1.34+1.40+1.45}{5}=1.41,
\)
离差平方和为
\[
\sum_{i=1}^{5}(x_i-\overline{x})^2
=0.0100+0.0196+0.0049+0.0001+0.0016
=0.0362.
\]
故
\(
\chi^2=\frac{0.0362}{0.048^2}\approx15.7118.
\)
双边检验显著性水平 \(\alpha=0.10\)，拒绝域为
\[
\chi^2\le \chi^2_{0.95}(4)=0.711
\quad\text{或}\quad
\chi^2\ge \chi^2_{0.05}(4)=9.488.
\]
由于 \(15.7118>9.488\)，统计量落入拒绝域，拒绝 \(H_0\)。因此在 \(10\%\) 显著性水平下，不能认为该纤度的方差仍属正常。
\end{solution}


\begin{exercise}\label{exer:src-2008-01-3c-ii-05}\label{exer:src-paper4-ii-05}
设总体 \(X\sim N(\mu,\sigma^2)\)，\(\mu\) 为未知参数，样本 \(X_1,X_2,\dots,X_n\) 的样本方差为 \(S^2\)。对假设
\[
H_0\colon \sigma\ge2,\qquad H_1\colon \sigma<2,
\]
显著性水平为 \(\alpha\) 的拒绝域是（\quad）。\par
A. \(\chi^2\le \chi^2_{1-\alpha/2}(n-1)\)\quad
B. \(\chi^2\le \chi^2_{1-\alpha}(n-1)\)\par
C. \(\chi^2\le \chi^2_{1-\alpha/2}(n)\)\quad
D. \(\chi^2\le \chi^2_{1-\alpha}(n)\)
\end{exercise}
\begin{solution}
应选 B。

这是单个正态总体方差的左边检验。由于 \(\mu\) 未知，检验统计量为
\(
\chi^2=\frac{(n-1)S^2}{\sigma_0^2}
=\frac{(n-1)S^2}{4}\sim\chi^2(n-1)
\)
在 \(H_0\) 的边界 \(\sigma_0=2\) 下成立。

备择假设为 \(\sigma<2\)，对应左尾拒绝域。因此显著性水平为 \(\alpha\) 时，
\(
\chi^2\le \chi^2_{1-\alpha}(n-1),
\)
其中本书采用上侧分位数记号 \(P\{\chi^2>\chi^2_\gamma(n-1)\}=\gamma\)。
\end{solution}

\begin{exercise}\label{exer:9-chi2-left}
测定溶液水分，10 个测定值 \(s=0.037\%\)，总体服从正态分布。在 \(\alpha=0.05\) 下检验 \(H_0\colon \sigma \geq 0.04\%\)，\(H_1\colon \sigma < 0.04\%\)。
（\(\chi^2_{0.95}(9)=3.325\)）
\end{exercise}
\begin{solution}
\textbf{提出假设。}题目给出 \(H_0\colon \sigma \geq 0.04\%\)，需等价转化为关于 \(\sigma^2\) 的假设：
\[
H_0\colon \sigma^2 \geq (0.04\%)^2 = 1.6\times10^{-6},\quad H_1\colon \sigma^2 < (0.04\%)^2.
\]
这是\textbf{左边检验}。\\
\textbf{选择检验统计量。}\(\mu\) 未知，使用 \(\chi^2\) 检验：
\(
\chi^2 = \frac{(n-1)S^2}{\sigma_0^2} = \frac{9\,S^2}{(0.04\%)^2} \sim \chi^2(9).
\)
\textbf{确定拒绝域。}左边检验，\(\alpha=0.05\)，拒绝域为 \(\chi^2 \leq \chi^2_{1-\alpha}(9) = \chi^2_{0.95}(9) = 3.325\)。\\
\textbf{计算并判断。}已知 \(s=0.037\%\)：
\[
\chi^2 = \frac{9 \times (0.037\%)^2}{(0.04\%)^2} = 9 \times \left(\frac{0.037}{0.04}\right)^{\!2} = 9 \times 0.8556 = 7.701.
\]
\(7.701 > 3.325\)，未落入拒绝域，故不能拒绝 \(H_0\)，即在 5\% 显著性水平下没有充分证据认为 \(\sigma < 0.04\%\)。
\end{solution}

\begin{exercise}\label{exer:src-2017-201207-a-07-34credit}
某糖厂用自动打印机装糖。已知每袋糖的质量（单位：kg）服从正态分布 \(N(\mu,\sigma^2)\)。现随机抽取 \(9\) 袋，称得样本均值 \(\overline x=48.5\)，样本标准差 \(S=2.5\)。在显著性水平 \(\alpha=0.05\) 下，对
\(
H_0\colon \mu=50
\)
分别在下列两种情形中进行检验：\par
(1) 已知 \(\sigma^2=4\)；\par
(2) \(\sigma^2\) 未知。
\end{exercise}
\begin{solution}
(1) 已知 \(\sigma=2\)，采用双边 \(u\) 检验。检验统计量为
\(
u=\frac{\overline X-\mu_0}{\sigma/\sqrt n}\sim N(0,1).
\)
代入数据：
\(
u=\frac{48.5-50}{2/\sqrt9}=-2.25.
\)
当 \(\alpha=0.05\) 时，双边拒绝域为
\(
|u|\ge u_{0.025}=1.96.
\)
由于 \(|u|=2.25>1.96\)，拒绝 \(H_0\)。在已知 \(\sigma^2=4\) 时，样本显示平均质量与 \(50\) 有显著差异。

(2) 当 \(\sigma^2\) 未知时，采用双边 \(t\) 检验：
\(
t=\frac{\overline X-\mu_0}{S/\sqrt n}\sim t(n-1).
\)
代入数据：
\(
t=\frac{48.5-50}{2.5/\sqrt9}=-1.8.
\)
自由度为 \(8\)，显著性水平 \(\alpha=0.05\) 的双边拒绝域为
\(
|t|\ge t_{0.025}(8)=2.31.
\)
由于 \(|t|=1.8<2.31\)，不能拒绝 \(H_0\)。在 \(\sigma^2\) 未知时，没有充分证据认为平均质量与 \(50\) 有显著差异。
\end{solution}

\begin{exercise}\label{exer:src-2017-201207-a-09-34credit}
设某种油漆的 \(9\) 个样本，其干燥时间（单位：h）分别为
\(
6.0,\ 5.7,\ 5.8,\ 6.5,\ 7.0,\ 6.3,\ 5.6,\ 6.1,\ 5.0.
\)
设干燥时间总体服从正态分布 \(N(\mu,\sigma^2)\)，求 \(\mu\) 的置信度为 \(95\%\) 的置信区间：\par
(1) 若已知 \(\sigma=0.6\)；\par
(2) 若 \(\sigma\) 未知。
\end{exercise}
\begin{solution}
样本容量 \(n=9\)，样本均值为
\(
\overline x=\frac{6.0+5.7+\cdots+5.0}{9}=6.
\)

(1) 当 \(\sigma=0.6\) 已知时，
\(
\overline X\sim N\!\left(\mu,\frac{\sigma^2}{n}\right),
\)
故 \(\mu\) 的 \(95\%\) 置信区间为
\(
\overline x\pm u_{0.025}\frac{\sigma}{\sqrt n}.
\)
代入 \(u_{0.025}=1.96\)，得
\(
6\pm1.96\cdot\frac{0.6}{3}
=6\pm0.392,
\)
即
\(
(5.608,\ 6.392).
\)

(2) 当 \(\sigma\) 未知时，先计算修正样本方差
\[
S^2=\frac1{n-1}\sum_{i=1}^9(x_i-\overline x)^2
=\frac{2.63}{8}\approx0.33.
\]
此时
\(
\frac{\overline X-\mu}{S/\sqrt n}\sim t(8).
\)
由 \(t_{0.025}(8)=2.306\)，置信区间为
\[
\overline x\pm t_{0.025}(8)\frac{S}{\sqrt n}
\approx6\pm2.306\frac{\sqrt{0.33}}{3}.
\]
因此
\(
\mu\in(5.558,\ 6.442).
\)
\end{solution}



\subsection{两个正态总体方差比的检验}

\begin{property}[F 检验]\label{prop:f-test}
\(\mu_1, \mu_2\) 未知时，检验统计量
\(
F = \frac{S_1^2}{S_2^2} \sim F(n_1-1, n_2-1).
\)
\begin{center}
\begin{tabular}{cc}
\toprule
\(H_1\) & 拒绝域 \\
\midrule
\(\sigma_1^2 > \sigma_2^2\) & \(F \geq F_\alpha(n_1-1, n_2-1)\) \\
\(\sigma_1^2 < \sigma_2^2\) & \(F \leq F_{1-\alpha}(n_1-1, n_2-1)\) \\
\(\sigma_1^2 \neq \sigma_2^2\) & \(F \geq F_{\alpha/2}\) 或 \(F \leq F_{1-\alpha/2}\) \\
\bottomrule
\end{tabular}
\end{center}
其中 \(F_{1-\alpha}(n_1-1,n_2-1) = \dfrac{1}{F_\alpha(n_2-1,n_1-1)}\)。
\end{property}

\begin{exercise}\label{exer:9-f-test}
两个小麦品种从播种到抽穗天数如下：
\begin{center}
\begin{tabular}{ccccccccccc|c}
\toprule
品种1 & 101 & 100 & 99 & 99 & 98 & 100 & 98 & 99 & 99 & 99 & 行和 992 \\
\midrule
品种2 & 100 & 98 & 100 & 99 & 98 & 99 & 98 & 98 & 99 & 100 & 行和 989 \\
\bottomrule
\end{tabular}
\end{center}
在 \(\alpha=0.05\) 下：(1) 检验方差是否有显著差异；(2) 检验均值是否有显著差异。
（\(F_{0.025}(9,9)=4.03\)，\(F_{0.975}(9,9)=0.248\)，\(t_{0.025}(18)=2.101\)）
\end{exercise}
\begin{solution}
\textbf{准备工作}：\(n_1=n_2=10\)，\(\overline{x}_1=992/10=99.2\)，\(\overline{x}_2=989/10=98.9\)。
计算两组的离差平方和：
\[
\sum_{i=1}^{10}(x_{1i}-99.2)^2 = (101{-}99.2)^2\!\times\!1 + (100{-}99.2)^2\!\times\!2 + (99{-}99.2)^2\!\times\!5 + (98{-}99.2)^2\!\times\!2 = 7.6,
\]
\[
\sum_{j=1}^{10}(x_{2j}-98.9)^2 = (100{-}98.9)^2\!\times\!3 + (99{-}98.9)^2\!\times\!3 + (98{-}98.9)^2\!\times\!4 = 6.9.
\]
\(
S_1^2 = \frac{7.6}{9} \approx 0.844,\quad S_2^2 = \frac{6.9}{9} \approx 0.767.
\)
\textbf{(1) 检验方差是否有显著差异（F 检验）。}提出假设：\(H_0\colon \sigma_1^2=\sigma_2^2\)，\(H_1\colon \sigma_1^2\neq\sigma_2^2\)（双边检验）。
检验统计量 \(F = S_1^2/S_2^2 \sim F(9,9)\)（在 \(H_0\) 成立时）。拒绝域：\(F \geq F_{0.025}(9,9)=4.03\) 或 \(F \leq F_{0.975}(9,9)=0.248\)。
\(
F = \frac{0.844}{0.767} \approx 1.100.
\)
\(0.248 < 1.100 < 4.03\)，未落入拒绝域，故不能拒绝 \(H_0\)，即没有充分证据认为两品种方差存在显著差异。因此后续可以用等方差的 pooled \(t\) 检验。\\
\textbf{(2) 检验均值是否有显著差异（t 检验）。}提出假设：\(H_0\colon \mu_1=\mu_2\)，\(H_1\colon \mu_1\neq\mu_2\)（双边检验）。由(1)知方差相等，使用 pooled \(t\) 检验。先计算合并样本方差：
\[
S_w^2 = \frac{(n_1-1)S_1^2+(n_2-1)S_2^2}{n_1+n_2-2} = \frac{9\times0.844+9\times0.767}{18} = \frac{7.6+6.9}{18} = \frac{14.5}{18} \approx 0.806.
\]
检验统计量：
\[
t = \frac{\overline{X}_1-\overline{X}_2}{S_w\sqrt{1/n_1+1/n_2}} = \frac{99.2-98.9}{\sqrt{0.806\times(0.1+0.1)}} = \frac{0.3}{\sqrt{0.1612}} = \frac{0.3}{0.401} \approx 0.747.
\]
拒绝域：\(|t| \geq t_{0.025}(18)=2.101\)。\(|0.747| < 2.101\)，未落入拒绝域，故不能拒绝 \(H_0\)，即没有充分证据认为两品种均值存在显著差异。
\begin{note}
在实际应用中，检验两总体均值差时，通常\textbf{先}做 F 检验判断方差是否相等，再根据结果选择合适的 \(t\) 检验方法。
\end{note}
\end{solution}

\begin{exercise}[绿地面积的 \texorpdfstring{\(t\)}{t} 检验与 \texorpdfstring{\(\chi^2\)}{χ²} 检验] \label{ex:ht_combo}
为改建绿地，5 位学生独立测量面积（单位：km\(^2\)）：
\(
1.23,\ 1.22,\ 1.20,\ 1.26,\ 1.23.
\)
设测量误差服从正态分布（\(\alpha=0.05\)）\\
(1) 以前面积 \(\mu_0=1.23\) km\(^2\)，是否有必要修改？\\
(2) 若要求标准差不超过 \(\sigma=0.015\)，能否认为标准差显著偏大？
（\(t_{0.025}(4)=2.7764\)，\(\chi^2_{0.05}(4)=9.488\)）
\end{exercise}
\begin{solution}
\(\overline{x}=1.228\)，离差平方和 \(\sum(x_i-\overline{x})^2=0.00188\)，\(s^2=0.000470\)，\(s\approx0.0217\)。\\
(1) \(H_0\colon\mu=1.23\)，\(H_1\colon\mu\neq1.23\)（双边 \(t\) 检验）。
\[
t=\frac{1.228-1.23}{0.0217/\sqrt{5}}\approx-0.206,\quad |t|<2.7764.
\]
故不能拒绝 \(H_0\)，在该显著性水平下没有必要修改。\\
(2) \(H_0\colon\sigma\leq0.015\)，\(H_1\colon\sigma>0.015\)（右边 \(\chi^2\) 检验）。
\[
\chi^2=\frac{4\times0.000470}{0.015^2}=\frac{0.00188}{0.000225}\approx8.356<9.488.
\]
故不能拒绝 \(H_0\)，不能认为标准差显著偏大。
\end{solution}

\begin{exercise}\label{exer:src-2009-3b-08}
自动包装机加工袋装食盐，每袋盐的净重 \(X\sim N(\mu,\sigma^2)\)，其中 \(\mu,\sigma^2\) 均未知。按规定每袋盐的标准重量为 \(100\) 克，标准差不能超过 \(5\) 克。一天，为检查机器的工作情况，随机地抽取 \(6\) 袋，测得样本均值 \(\overline x=99.3\) 克，样本标准差 \(s=6.24\) 克。问：分别通过检验期望 \(\mu\) 和方差 \(\sigma^2\) 来判断包装机该天的工作是否正常(\(\alpha=0.05\))？
\end{exercise}
\begin{solution}
分别检验均值与方差。

\textbf{(1) 检验均值。}作双边检验
\[
H_0\colon \mu=100,\qquad H_1\colon \mu\neq100.
\]
由于 \(\sigma^2\) 未知，采用 \(t\) 检验：
\(
t=\frac{\overline X-100}{S/\sqrt n}\sim t(5).
\)
代入数据得
\(
t=\frac{99.3-100}{6.24/\sqrt6}\approx -0.275.
\)
双边拒绝域为 \(|t|\ge t_{0.025}(5)=2.5706\)，故不能拒绝 \(H_0\)。

\textbf{(2) 检验方差。}规定标准差不能超过 \(5\) 克，作右边检验
\[
H_0\colon \sigma\le5,\qquad H_1\colon \sigma>5.
\]
等价地检验 \(\sigma^2\le25\) 对 \(\sigma^2>25\)。统计量为
\(
\chi^2=\frac{(n-1)S^2}{25}=\frac{5\times6.24^2}{25}\approx7.800.
\)
右边检验拒绝域为 \(\chi^2\ge\chi^2_{0.05}(5)=11.071\)，故不能拒绝 \(H_0\)。

综上，在显著性水平 \(0.05\) 下，可认为该天包装机工作正常。
\end{solution}

\subsection*{历年试题补充：2011 评分标准 A/B 卷}

\begin{exercise}\label{exer:src-2011-ab-a-07-34credit}
自动包装机将水泥装袋，每袋的标称重量为 \(100\) 千克，实际重量 \(X\sim N(\mu,\sigma^2)\)，其中 \(\mu,\sigma^2\) 未知，标准差不能超过 \(2\) 千克。为检查机器的工作情况，随机抽取 \(10\) 袋，测得样本均值 \(\overline x=98.2\) 千克，样本均方差 \(s=2.25\) 千克。试通过检验期望 \(\mu\) 和方差 \(\sigma^2\) 来判断包装机的工作是否正常（\(\alpha=0.05\)）。
\end{exercise}
\begin{solution}
分别检验均值是否等于标称重量，以及标准差是否超过规定上限。

\textbf{(1) 检验均值。}作双边检验
\[
H_0\colon \mu=100,\qquad H_1\colon \mu\ne100.
\]
总体正态且方差未知，采用 \(t\) 检验：
\(
t=\frac{\overline X-100}{S/\sqrt n}\sim t(n-1)=t(9).
\)
代入数据，
\(
t=\frac{98.2-100}{2.25/\sqrt{10}}\approx -2.53.
\)
双边拒绝域为
\(
|t|\ge t_{0.025}(9)=2.26.
\)
由于 \(|-2.53|>2.26\)，故拒绝 \(H_0\)，认为平均装袋量与 \(100\) 千克有显著差异，且样本均值偏低。

\textbf{(2) 检验方差。}规定标准差不能超过 \(2\) 千克，即 \(\sigma^2\le4\)。作右边检验
\[
H_0\colon \sigma^2\le4,\qquad H_1\colon \sigma^2>4.
\]
在边界 \(\sigma_0^2=4\) 下，
\(
\chi^2=\frac{(n-1)S^2}{\sigma_0^2}\sim\chi^2(9).
\)
代入得
\(
\chi^2=\frac{9\times2.25^2}{4}\approx11.39.
\)
右边检验拒绝域为
\(
\chi^2\ge \chi^2_{0.05}(9).
\)
由试卷给出的双侧分位点可见 \(2.7<11.39<19\)，故不拒绝 \(\sigma^2=4\) 的基准判断，也没有充分证据认为标准差超过 \(2\) 千克。

综上，方差方面未见异常，但均值检验拒绝 \(H_0\)，且装袋量偏低，因此包装机工作不正常。
\end{solution}

\subsection*{历年试题补充：2012 试卷 A 卷}

\begin{exercise}\label{exer:src-2012-a-07-34credit}
某种钢材的强度服从正态分布。已知旧工艺下平均强度为 \(52.00\)。改用新工艺后，随机抽取 \(7\) 个试件，测得强度为
\[
52.45,\quad 48.51,\quad 56.02,\quad 51.53,\quad
49.02,\quad 53.38,\quad 54.04.
\]
试在显著性水平 \(\alpha=0.05\) 下检验新工艺是否使平均强度有显著提高，并求平均强度 \(\mu\) 的 \(95\%\) 置信区间。
\end{exercise}
\begin{solution}
由样本计算得
\[
n=7,\qquad \overline{x}\approx52.1357,\qquad s\approx2.6954.
\]

先检验平均强度是否有显著提高。建立假设
\(
H_0\colon \mu=52.00,\qquad H_1\colon \mu>52.00.
\)
总体正态且方差未知，采用单样本右尾 \(t\) 检验：
\(
t=\frac{\overline X-52.00}{S/\sqrt n}\sim t(6).
\)
代入样本值，
\(
t=\frac{52.1357-52.00}{2.6954/\sqrt7}\approx0.133.
\)
显著性水平 \(\alpha=0.05\) 的右尾拒绝域为
\(
t\ge t_{0.05}(6)\approx1.943.
\)
由于 \(0.133<1.943\)，故不能拒绝 \(H_0\)。在 \(0.05\) 显著性水平下，没有充分证据认为新工艺使平均强度显著提高。

再求 \(95\%\) 置信区间。因 \(\sigma^2\) 未知，
\[
\mu\in\left(\overline{x}-t_{0.025}(6)\frac{s}{\sqrt n},
\overline{x}+t_{0.025}(6)\frac{s}{\sqrt n}\right).
\]
取 \(t_{0.025}(6)\approx2.447\)，得
\[
\mu\in
\left(52.1357-2.447\frac{2.6954}{\sqrt7},
52.1357+2.447\frac{2.6954}{\sqrt7}\right)
\approx(49.64,\ 54.63).
\]
\end{solution}

\subsection*{历年试题补充：2011 试卷 A 卷}

\begin{exercise}\label{exer:src-2011-a-08-34credit}
某市居民上月平均伙食费为 \(235.5\) 元。随机抽取 \(49\) 个居民，他们本月的伙食费平均为 \(236.5\) 元，由这 \(49\) 个样本算出的标准差 \(S=3.5\) 元。假设该市居民月伙食费 \(X\) 服从正态分布，试在 \(\alpha=0.01\) 时，检验“本月该市居民平均伙食费较之上个月无变化”的假设。
\end{exercise}
\begin{solution}
需要检验
\[
H_0\colon \mu=235.5,\qquad H_1\colon \mu\ne235.5.
\]
总体正态，方差未知，采用双边 \(t\) 检验：
\(
t=\frac{\sqrt n(\overline X-\mu_0)}{S}\sim t(n-1).
\)
代入 \(n=49,\ \overline x=236.5,\ S=3.5,\ \mu_0=235.5\)，得
\(
|t|
=\left|\frac{\sqrt{49}(236.5-235.5)}{3.5}\right|
=2.
\)
自由度 \(48>30\)，可用标准正态分位点近似，显著性水平 \(\alpha=0.01\) 的双边拒绝域近似为
\(
|t|\ge u_{0.005}=2.58.
\)
由于
\(
2<2.58,
\)
故不能拒绝 \(H_0\)。在显著性水平 \(0.01\) 下，没有充分证据认为本月该市居民平均伙食费较上个月有显著变化。
\end{solution}

\begin{exercise}\label{exer:src-2008-01-3c-iv-06}\label{exer:src-paper4-iv-06}
自动包装机加工袋装食盐，每袋盐的净重 \(X\sim N(\mu,\sigma^2)\)，其中 \(\mu,\sigma^2\) 均未知。按规定每袋盐的标准重量为 \(500\) 克，标准差不能超过 \(10\) 克。某天随机抽取 \(6\) 袋，测得样本均值
\(
\overline{x}=495.3\ \text{克},
\)
样本标准差
\(
s=13.74\ \text{克}.
\)
在显著性水平 \(\alpha=0.05\) 下，通过检验期望 \(\mu\) 和方差 \(\sigma^2\) 判断包装机该天工作是否正常。
（\(t_{0.025}(5)=2.5706,\ \chi^2_{0.05}(5)=11.071,\ \chi^2_{0.025}(5)=12.833\)）
\end{exercise}
\begin{solution}
需要分别检验均值是否偏离标准重量，以及标准差是否超过规定上限。

\textbf{(1) 检验均值。}作双边检验
\[
H_0\colon \mu=500,\qquad H_1\colon \mu\neq500.
\]
总体方差未知，采用 \(t\) 检验：
\(
t=\frac{\overline X-500}{S/\sqrt n}\sim t(5).
\)
代入数据，
\(
t=\frac{495.3-500}{13.74/\sqrt6}
\approx \frac{-4.7}{5.609}
\approx -0.838.
\)
双边检验拒绝域为
\(
|t|\ge t_{0.025}(5)=2.5706.
\)
由于 \(|t|=0.838<2.5706\)，故不能拒绝 \(H_0\)。没有充分证据认为平均净重偏离 \(500\) 克。

\textbf{(2) 检验方差。}规定“标准差不能超过 \(10\) 克”，故作右边检验
\[
H_0\colon \sigma\le10,\qquad H_1\colon \sigma>10.
\]
等价地检验 \(\sigma^2\le100\) 对 \(\sigma^2>100\)。统计量为
\(
\chi^2=\frac{(n-1)S^2}{100}\sim\chi^2(5)
\)
在边界 \(\sigma_0=10\) 下成立。代入数据：
\[
\chi^2=\frac{5\times 13.74^2}{100}
=\frac{5\times188.7876}{100}
\approx9.439.
\]
右边检验拒绝域为
\(
\chi^2\ge\chi^2_{0.05}(5)=11.071.
\)
由于 \(9.439<11.071\)，故不能拒绝 \(H_0\)。没有充分证据认为标准差超过 \(10\) 克。

综上，在显著性水平 \(0.05\) 下，均值和标准差两方面都未显示异常，可认为包装机该天工作正常。
\end{solution}

\begin{exercise}\label{exer:src-paper3-iv-06}\label{prob:p3-iv-06}
机器自动包装食盐，设每袋盐的净重服从正态分布，规定每袋盐的标准重量为 \(500\) 克，标准差不能超过 \(10\) 克。某天开工后，为了检验机器是否正常工作，从已经包装好的食盐中随机取 \(9\) 袋，测得
\(
\overline X=499,\qquad S=16.03.
\)
在显著性水平 \(\alpha=0.05\) 下检验：\par
(1) \(H_0\colon \mu=500,\ H_1\colon \mu\ne500\)；\par
(2) \(H_0\colon \sigma^2\le100,\ H_1\colon \sigma^2>100\)。\par
已知 \(t_{0.025}(8)=2.306,\ \chi^2_{0.05}(8)=15.507\)。
\end{exercise}
\begin{solution}
(1) 检验均值。总体正态且方差未知，采用
\(
T=\frac{\overline X-500}{S/\sqrt n}\sim t(8).
\)
双边检验的拒绝域为
\(
|T|>t_{0.025}(8)=2.306.
\)
代入数据：
\(
T_0=\frac{499-500}{16.03/3}
=-\frac3{16.03}
\approx -0.187.
\)
由于 \(|T_0|<2.306\)，故不能拒绝 \(H_0\)。均值方面没有显著异常。

(2) 检验方差。由于备择假设为 \(\sigma^2>100\)，采用右尾检验。统计量为
\(
\chi^2=\frac{(n-1)S^2}{100}
=\frac{8S^2}{100},
\)
在边界 \(\sigma^2=100\) 下服从 \(\chi^2(8)\)。拒绝域为
\(
\frac{8S^2}{100}>\chi^2_{0.05}(8)=15.507.
\)
代入数据：
\(
\frac{8\times16.03^2}{100}
\approx20.56.
\)
由于 \(20.56>15.507\)，故拒绝 \(H_0\)，认为标准差超过规定上限。

综上，均值检验未显示异常，但方差检验显示波动过大，因此这天自动包装机工作不正常。
\end{solution}


\begin{exercise}[\texorpdfstring{\(F\)}{F} 检验与合并 \texorpdfstring{\(t\)}{t} 检验] \label{ex:ht_f_pooled}
在施肥与不施肥的两块苗床上各抽取 6 株苗木，苗高数据（cm）如下：
\begin{center}
\begin{tabular}{ccccccc|c}
\toprule
施肥 & 77.3 & 79.1 & 81.0 & 79.1 & 82.1 & 77.3 & 行和 475.9 \\
\midrule
不施肥 & 75.5 & 76.2 & 78.1 & 72.4 & 77.4 & 76.7 & 行和 456.3 \\
\bottomrule
\end{tabular}
\end{center}
设苗高服从正态分布（\(\alpha=0.05\)）。\\
(1) 检验两组方差是否相同；(2) 检验施肥是否显著高于不施肥。
（\(F_{0.025}(5,5)=7.15\)，\(F_{0.975}(5,5)=0.14\)，\(t_{0.05}(10)=1.812\)）
\end{exercise}
\begin{solution}
\(n_1=n_2=6\)，\(\overline{x}_1=475.9/6\approx79.317\)，\(\overline{x}_2=456.3/6=76.05\)。
离差平方和：
\[
\sum(x_{1i}-\overline{x}_1)^2\approx18.808,\quad S_1^2=\frac{18.808}{5}=3.762;
\]
\[
\sum(x_{2j}-\overline{x}_2)^2=20.095,\quad S_2^2=\frac{20.095}{5}=4.019.
\]
\textbf{(1)} \(H_0\colon\sigma_1^2=\sigma_2^2\)，\(H_1\colon\sigma_1^2\neq\sigma_2^2\)。
\(
F=\frac{S_1^2}{S_2^2}=\frac{3.762}{4.019}\approx0.936.
\)
拒绝域：\(F\geq7.15\) 或 \(F\leq0.14\)。\(0.14<0.936<7.15\)，未落入拒绝域，故不能拒绝 \(H_0\)，即方差无显著差异。\\
\textbf{(2)} 方差齐性成立，用 pooled \(t\) 检验。\(H_0\colon\mu_1\leq\mu_2\)，\(H_1\colon\mu_1>\mu_2\)。
\(
S_w^2=\frac{18.808+20.095}{10}=3.890,\quad S_w\approx1.972.
\)
\[
t=\frac{79.317-76.05}{1.972\sqrt{1/6+1/6}}=\frac{3.267}{1.972/\sqrt{3}}\approx\frac{3.267}{1.139}\approx2.869.
\]
拒绝域 \(t\geq t_{0.05}(10)=1.812\)。\(2.869>1.812\)，拒绝 \(H_0\)，施肥苗高显著高于不施肥。
\end{solution}

\begin{exercise}\label{exer:src-2015-spring-a-08}
(1) 已知多名实习生相互独立地测量同一块土地的面积。设每名实习生得到的测量数据 \(X\)（平方米）服从正态分布 \(N(\mu,\sigma^2)\)，从这些测量数据中随机抽取 \(7\) 个，经计算其平均面积为 \(125\) 平方米，标准差为 \(2.71\) 平方米。求 \(\mu\) 的置信度为 \(90\%\) 的置信区间。（已知 \(t_{0.95}(6)=1.943\)）\par
(2) 甲乙两厂生产的灯泡，其寿命 \(X\) 和 \(Y\) 分别服从 \(N(\mu_1,84^2)\) 和 \(N(\mu_2,96^2)\)。现从两厂生产的灯泡中各取 \(60\) 只，测得平均寿命甲厂为 \(\overline x=1295\) 小时，乙厂为 \(\overline y=1230\) 小时。问在显著性水平 \(\alpha=0.05\) 下，能否认为两厂生产的灯泡寿命无显著差异？
\end{exercise}
\begin{solution}
(1) 总体正态且方差未知，使用 \(t\) 分布。按源卷给出的样本标准差记号，置信度为 \(1-\alpha=0.90\) 的置信区间为
\[
\left(
\overline X-t_{1-\alpha/2}(n-1)\frac{S}{\sqrt{n-1}},
\overline X+t_{1-\alpha/2}(n-1)\frac{S}{\sqrt{n-1}}
\right).
\]
代入
\[
\overline X=125,\qquad S=2.71,\qquad n=7,\qquad t_{0.95}(6)=1.943,
\]
得误差限
\(
1.943\cdot\frac{2.71}{\sqrt6}\approx2.15.
\)
故 \(\mu\) 的 \(90\%\) 置信区间约为
\(
(125-2.15,\ 125+2.15)=(122.85,\ 127.15),
\)
按一位小数写为
\(
(122.9,\ 127.1).
\)

(2) 检验
\[
H_0\colon\mu_1=\mu_2,\qquad H_1\colon\mu_1\ne\mu_2.
\]
两总体方差已知，采用双总体均值差的 \(u\) 检验。在 \(H_0\) 下，
\[
u=
\frac{(\overline X-\overline Y)-0}
{\sqrt{\sigma_1^2/n+\sigma_2^2/m}}
\sim N(0,1).
\]
代入数据：
\[
u=
\frac{1295-1230}
{\sqrt{84^2/60+96^2/60}}
\approx3.95.
\]
显著性水平 \(\alpha=0.05\) 的双边拒绝域为
\(
|u|\ge u_{0.975}=1.96.
\)
因为 \(3.95>1.96\)，所以拒绝 \(H_0\)。因此在显著性水平 \(0.05\) 下，应认为两厂生产的灯泡寿命有显著差异。
\end{solution}

\begin{exercise}\label{exer:src-2017-a-08}
已知多名实习生相互独立地测量同一块土地的面积。设每名实习生得到的测量数据 \(X\)（平方米）服从正态分布 \(N(\mu,\sigma^2)\)。随机抽取 \(7\) 个测量数据，经计算其平均面积为 \(125\) 平方米，修正样本标准差为 \(S^*=2.71\) 平方米。\par
(1) 求 \(\mu\) 的置信度为 \(90\%\) 的置信区间；\par
(2) 在显著性水平 \(\alpha=0.1\) 下，能否认为这块土地的平均面积为 \(124\) 平方米？\par
已知 \(t_{0.95}(6)=1.943\)。
\end{exercise}
\begin{solution}
(1) 总体正态且方差未知，使用
\(
T=\frac{\overline X-\mu}{S^*/\sqrt n}\sim t(n-1).
\)
这里 \(n=7,\ \overline x=125,\ S^*=2.71\)，故置信度 \(90\%\) 的置信区间为
\(
\overline x\pm t_{0.95}(6)\frac{S^*}{\sqrt7}.
\)
代入得误差限
\(
1.943\cdot\frac{2.71}{\sqrt7}\approx1.99.
\)
因此
\(
\mu\in(125-1.99,\ 125+1.99)=(123.01,\ 126.99),
\)
按源卷取整写为
\(
(123,\ 127).
\)

(2) 作双边检验
\[
H_0\colon \mu=124,\qquad H_1\colon \mu\ne124.
\]
检验统计量为
\(
T=\frac{\overline X-\mu_0}{S^*/\sqrt n}
=\frac{125-124}{2.71/\sqrt7}
\approx0.976.
\)
显著性水平 \(\alpha=0.1\) 的双边拒绝域为
\(
|T|\ge t_{0.95}(6)=1.943.
\)
由于
\(
0.976<1.943,
\)
故不能拒绝 \(H_0\)。因此在显著性水平 \(0.1\) 下，可以认为这块土地的平均面积为 \(124\) 平方米。
\end{solution}


\section{置信区间与假设检验的关系} \label{sec:ht-ci}

\begin{property}[置信区间与假设检验的等价性]\label{prop:ci-test-equiv}
设 \((\theta_L, \theta_U)\) 为参数 \(\theta\) 的置信水平为 \(1-\alpha\) 的置信区间，则对于双边检验
\[
H_0\colon \theta=\theta_0,\quad H_1\colon \theta\neq\theta_0,
\]
“不能拒绝 \(H_0\)”等价于 “\(\theta_0\) 落在置信区间 \((\theta_L, \theta_U)\) 内”。对于单边检验也有类似的对应关系：若做右边检验 \(H_0\colon \theta\leq\theta_0\) 并构造下单侧置信区间 \((\underline{\theta},+\infty)\)，则拒绝域 \(z\geq z_\alpha\) 等价于 \(\theta_0<\underline{\theta}\)；若做左边检验 \(H_0\colon \theta\geq\theta_0\) 并构造上单侧置信区间 \((-\infty,\overline{\theta})\)，则拒绝域对应于 \(\theta_0>\overline{\theta}\)。
\end{property}

\begin{note}
\textbf{易错点：} “置信区间和假设检验等价”有严格前提：参数相同、显著性水平匹配、而且尾部形式也要匹配。双侧 \(1-\alpha\) 置信区间对应双边显著性水平为 \(\alpha\) 的检验；右边检验要配下单侧区间，左边检验要配上单侧区间，不能拿双侧区间去机械替代单边检验。
\end{note}

\begin{exercise}\label{exer:9-ci-test}
零件长度 \(X \sim N(\mu, \sigma^2)\)（\(\mu, \sigma^2\) 均未知），\(n=16\)，\(\overline{x}=10\)，\(s^2=0.16\)。\\
(1) 求 \(\mu\) 的置信水平为 0.95 的置信区间；\\
(2) 在 \(\alpha=0.05\) 下检验 \(H_0\colon \mu=9.7\)，\(H_1\colon \mu\neq9.7\)。（\(t_{0.025}(15)=2.132\)）
\end{exercise}
\begin{solution}
\textbf{(1) 求置信区间。}\(\mu\) 未知，\(\sigma^2\) 未知，使用 \(t\) 分布。
已知 \(n=16\)，\(\overline{x}=10\)，\(s=\sqrt{0.16}=0.4\)，\(t_{0.025}(15)=2.132\)。置信区间为：
\[
\left(\overline{x} - \frac{s}{\sqrt{n}}\,t_{\alpha/2}(n{-}1),\quad \overline{x} + \frac{s}{\sqrt{n}}\,t_{\alpha/2}(n{-}1)\right).
\]
代入计算：
\[
\frac{s}{\sqrt{n}}\,t_{\alpha/2}(15) = \frac{0.4}{4}\times2.132 = 0.2132.
\]
\(
\text{置信区间} = (10-0.2132,\ 10+0.2132) = (9.7868,\ 10.2132).
\)
\textbf{(2) 假设检验。}提出假设：\(H_0\colon \mu=9.7\)，\(H_1\colon \mu\neq9.7\)（双边检验，\(\alpha=0.05\)）。
检验统计量 \(t = \dfrac{\overline{X}-9.7}{S/\sqrt{n}} \sim t(15)\)，拒绝域 \(|t| \geq t_{0.025}(15)=2.132\)。
\(
t = \frac{10-9.7}{0.4/4} = \frac{0.3}{0.1} = 3.0 > 2.132.
\)
落入拒绝域，拒绝 \(H_0\)，认为 \(\mu\neq9.7\)。
\textbf{等价性验证}：由命题~\ref{prop:ci-test-equiv}，双边检验中“不能拒绝 \(H_0\)”等价于 \(\mu_0\) 落在置信区间内。此处 \(\mu_0=9.7\)，而置信区间为 \((9.7868,\ 10.2132)\)，\(9.7 < 9.7868\)，\(\mu_0\) 在区间外，故应拒绝 \(H_0\)，与上述检验结论一致。
\end{solution}

\begin{exercise}\label{exer:src-2014-spring-a-07}
(1) 设零件长度 \(X\sim N(\mu,\sigma^2)\)，其中 \(\mu,\sigma^2\) 均未知。现抽取容量为 \(16\) 的样本，测得 \(\overline x=10,\ s^2=0.16\)。求 \(\mu\) 的 \(95\%\) 置信区间。（已知 \(t_{0.025}(15)=2.132\)）\par
(2) 某维尼纶纤维在正常生产条件下纤度服从 \(N(1.405,0.048^2)\)。现抽取 \(5\) 根纤维，测得纤度为
\(
1.32,\quad 1.55,\quad 1.36,\quad 1.40,\quad 1.44.
\)
试在显著性水平 \(\alpha=0.05\) 下检验生产方差是否正常。（已知 \(\chi^2_{0.025}(5)=12.833,\ \chi^2_{0.975}(5)=0.831\)）
\end{exercise}
\begin{solution}
(1) 因为总体正态且 \(\sigma^2\) 未知，使用 \(t\) 分布。由 \(s=\sqrt{0.16}=0.4,\ n=16\)，得
\[
\mu\in
\left(
\overline x-t_{0.025}(15)\frac{s}{\sqrt n},
\overline x+t_{0.025}(15)\frac{s}{\sqrt n}
\right).
\]
代入数值，
\(
t_{0.025}(15)\frac{s}{\sqrt n}
=2.132\cdot\frac{0.4}{4}=0.2132.
\)
故
\(
\mu\in(9.7868,\ 10.2132).
\)

(2) 检验
\[
H_0\colon \sigma^2=0.048^2,\qquad
H_1\colon \sigma^2\ne0.048^2.
\]
由于题中正常生产均值 \(1.405\) 已知，故在 \(H_0\) 下
\(
\chi^2=\sum_{i=1}^{5}\frac{(X_i-1.405)^2}{0.048^2}\sim\chi^2(5).
\)
样本平方离差和为
\[
\begin{aligned}
\sum_{i=1}^{5}(x_i-1.405)^2
&=(-0.085)^2+0.145^2+(-0.045)^2+(-0.005)^2+0.035^2\\
&=0.031525.
\end{aligned}
\]
因此
\(
\chi^2=\frac{0.031525}{0.048^2}\approx13.683.
\)
双边检验的拒绝域为
\[
\chi^2>\chi^2_{0.025}(5)=12.833
\quad\text{或}\quad
\chi^2<\chi^2_{0.975}(5)=0.831.
\]
由于 \(13.683>12.833\)，故拒绝 \(H_0\)，认为生产方差不正常。
\end{solution}

\subsection*{历年试题补充：2010.06.30 试卷 A 卷}

\begin{exercise}\label{exer:src-2010-0630-a-07-34credit}
化肥厂用自动打包机包装化肥，某日测得 \(8\) 包化肥的重量（斤）如下：
\[
98.7,\quad 100.5,\quad 101.2,\quad 98.3,\quad
99.7,\quad 99.5,\quad 101.4,\quad 100.5.
\]
已知各包重量服从正态分布 \(N(\mu,\sigma^2)\)。求：\par
(1) 在显著性水平 \(\alpha=0.05\) 下，是否可以认为每包平均重量为 \(100\) 斤？\par
(2) 参数 \(\sigma^2\) 的 \(90\%\) 置信区间。\par
可用分位点：\(t_{0.025}(7)\approx2.365\)，\(\chi^2_{0.05}(7)=14.067\)，\(\chi^2_{0.95}(7)=2.167\)。
\end{exercise}
\begin{solution}
由样本数据得
\(
\overline{x}=\frac{799.8}{8}=99.975,
\)
样本离差平方和为
\(
\sum_{i=1}^8(x_i-\overline{x})^2=8.815.
\)
因此
\(
S^2=\frac{8.815}{7}\approx1.2593,\qquad S\approx1.1222.
\)

(1) 检验
\[
H_0\colon \mu=100,\qquad H_1\colon \mu\ne100.
\]
总体正态且 \(\sigma^2\) 未知，使用 \(t\) 检验：
\(
t=\frac{\overline{X}-100}{S/\sqrt n}\sim t(7).
\)
代入样本值，
\(
t=\frac{99.975-100}{1.1222/\sqrt8}\approx -0.063.
\)
双边检验的拒绝域为
\(
|t|\ge t_{0.025}(7)\approx2.365.
\)
由于
\(
|-0.063|<2.365,
\)
故不能拒绝 \(H_0\)。在显著性水平 \(0.05\) 下，可以认为每包平均重量为 \(100\) 斤。

(2) 由于总体正态且均值未知，
\(
\frac{(n-1)S^2}{\sigma^2}\sim\chi^2(n-1)=\chi^2(7).
\)
\(\sigma^2\) 的 \(90\%\) 置信区间为
\(
\left(
\frac{(n-1)S^2}{\chi^2_{0.05}(7)},
\frac{(n-1)S^2}{\chi^2_{0.95}(7)}
\right).
\)
代入 \((n-1)S^2=8.815\)，得
\(
\left(
\frac{8.815}{14.067},
\frac{8.815}{2.167}
\right)
\approx(0.6268,\ 4.068).
\)
因此 \(\sigma^2\) 的 \(90\%\) 置信区间约为
\(
(0.627,\ 4.068).
\)
\end{solution}

\begin{exercise}[置信区间与假设检验的综合] \label{ex:ci_ht_01}
根据环保条例，排放工业废水中某有害物质含量不得超过 0.5。设该含量 \(X\sim N(\mu,\sigma^2)\)，现取 5 份水样，测得含量分别为
\(
0.530,\ 0.542,\ 0.510,\ 0.495,\ 0.515.
\)
(1) 求 \(\mu\) 的置信度为 95\% 的置信区间；\\
(2) 能否据此说明有害物质含量超过了规定标准（\(\alpha=0.05\)）？
（\(t_{0.05}(4)=2.132\)，\(t_{0.025}(4)=2.776\)）
\end{exercise}
\begin{solution}
\textbf{数据预处理。} \(n=5\)
\(
\overline{x}=\frac{0.530+0.542+0.510+0.495+0.515}{5}=0.5184.
\)
\[
\sum_{i=1}^{5}(x_i-0.5184)^2=0.0116^2+0.0236^2+0.0084^2+0.0234^2+0.0034^2=0.001321.
\]
\(
s^2=\frac{0.001321}{4}=0.000330,\quad s\approx 0.0182.
\)
(1) \(\sigma^2\) 未知，\(\mu\) 的 95\% 置信区间为
\begin{align*}
&\left(\overline{x}-\frac{s}{\sqrt{n}}\,t_{0.025}(4),\;\overline{x}+\frac{s}{\sqrt{n}}\,t_{0.025}(4)\right)\\
=&\left(0.5184-\frac{0.0182}{\sqrt{5}}\times2.776,\;0.5184+\frac{0.0182}{\sqrt{5}}\times2.776\right)=(0.4958,\;0.5410).
\end{align*}
(2) 右边检验：\(H_0\colon\mu\leq 0.5\)，\(H_1\colon\mu>0.5\)。
\[
t=\frac{0.5184-0.5}{0.0182/\sqrt{5}}=\frac{0.0184}{0.00813}\approx 2.264.
\]
拒绝域 \(t\geq t_{0.05}(4)=2.132\)。\(2.264>2.132\)，落入拒绝域，拒绝 \(H_0\)，即认为有害物质含量超过了规定标准。
\begin{note}
也可以利用(1)的结论：\(\mu_0=0.5\) 不在置信区间 \((0.4958,0.5410)\) 内，但此区间为\textbf{双侧}置信区间，对应\textbf{双边}检验。而本题(2)是\textbf{单边}检验，二者不完全等价，需单独计算。
\end{note}
\end{solution}


\section{两类错误}

\begin{note}
在固定样本容量 \(n\) 的条件下，\(\alpha\) 与 \(\beta\) 存在此消彼长的关系：\(\alpha\) 小则 \(\beta\) 大，\(\beta\) 小则 \(\alpha\) 大。\textbf{它们不满足} \(\beta=1-\alpha\)。
实际中，总是在控制 \(\alpha\)（使弃真概率较小）的条件下尽量使 \(\beta\) 小，因为人们通常认为错误地拒绝 \(H_0\) 比错误地未能拒绝 \(H_0\) 的后果更严重。
\end{note}

\begin{exercise}[离散总体第一类错误的计算] \label{ex:error_discrete}
设总体 \(X\) 的分布律为 \(P\{X=1\}=\theta\)，\(P\{X=0\}=1-\theta\)（\(0<\theta<1\)），\(X_1,X_2,X_3\) 为样本。对于检验
\[
H_0\colon\theta=0.1,\quad H_1\colon\theta>0.1,
\]
拒绝域为 \(\{X_1=1,\ X_2=0,\ X_3=1\}\)。求犯第一类错误的概率。
\end{exercise}
\begin{solution}
第一类错误指的是：\textbf{\(H_0\) 为真时却拒绝了 \(H_0\)} 的概率。

本题原假设为 \(H_0:\theta=0.1\)，而拒绝域给成 \(\{X_1=1,\ X_2=0,\ X_3=1\}\)。
因此犯第一类错误的概率，就是在 \(\theta=0.1\) 成立时，样本恰好落入这个拒绝域的概率。

当 \(\theta=0.1\) 时，
\(
P(X_i=1)=0.1,\qquad P(X_i=0)=0.9.
\)
又因为 \(X_1,X_2,X_3\) 是样本，相互独立，所以
\[
\alpha=P\{X_1=1,X_2=0,X_3=1\mid\theta=0.1\}=0.1\times0.9\times0.1=0.009.
\]

故犯第一类错误的概率为
\(
\alpha=0.009.
\)
\end{solution}


\begin{exercise}[正态总体两类错误的计算] \label{ex:error_normal}
设 \(X_1,X_2,\dots,X_{20}\) 为总体 \(N(\mu,1)\) 的样本，考虑假设检验
\[
H_0\colon\mu=2\quad\text{vs}\quad H_1\colon\mu=3,
\]
拒绝域为 \(W=\{\overline{X}\geq 2.6\}\)。求犯第一类错误的概率 \(\alpha\) 和犯第二类错误的概率 \(\beta\)。
\end{exercise}
\begin{solution}
\(\overline{X}\sim N(\mu,\,1/n)\)，此处 \(n=20\)，\(\overline{X}\sim N(\mu,\,1/20)\)。\\
\textbf{第一类错误}：\(\alpha=P\{\overline{X}\geq2.6\mid\mu=2\}\)。
标准化：令 \(Z=\dfrac{\overline{X}-2}{1/\sqrt{20}}\sim N(0,1)\)，
\[
\alpha=P\!\left\{Z\geq\frac{2.6-2}{1/\sqrt{20}}\right\}=P\{Z\geq 0.6\sqrt{20}\}=P\{Z\geq 1.2\sqrt{5}\}=1-\Phi(1.2\sqrt{5}).
\]
\textbf{第二类错误}：\(\beta=P\{\overline{X}<2.6\mid\mu=3\}\)。
令 \(Z'=\dfrac{\overline{X}-3}{1/\sqrt{20}}\sim N(0,1)\)，
\[
\beta=P\!\left\{Z'<\frac{2.6-3}{1/\sqrt{20}}\right\}=P\{Z'<-0.4\sqrt{20}\}=P\{Z'<-0.8\sqrt{5}\}=\Phi(-0.8\sqrt{5}).
\]
\end{solution}


\begin{exercise}\label{exer:9-error-calc}
设连续型随机变量 \(U\) 的密度分别为：
\[
H_0\colon f(x)=\begin{cases}1/2, & 0\leq x\leq2,\\0, & \text{其他},\end{cases}
\quad
H_1\colon f(x)=\begin{cases}x/2, & 0\leq x\leq2,\\0, & \text{其他}.\end{cases}
\]
检验规则：当 \(U>3/2\) 时拒绝 \(H_0\)。求犯第一类错误的概率 \(\alpha\) 和犯第二类错误的概率 \(\beta\)。\par
A. \(3/4,\;7/16\)\qquad B. \(7/16,\;3/4\)\qquad C. \(9/16,\;1/4\)\qquad D. \(1/4,\;9/16\)
\end{exercise}
\begin{solution}
应选 D。根据两类错误的定义：
\begin{itemize}
    \item \textbf{第一类错误}\(\alpha = P\{\text{拒绝}\,H_0 \mid H_0\text{ 为真}\} = P\!\left\{U>\tfrac{3}{2}\,\middle|\,H_0\right\}\)。
    在 \(H_0\) 下，\(U\) 服从 \([0,2]\) 上的均匀分布，密度 \(f(x)=1/2\)：
    \[
    \alpha = \int_{3/2}^{2}\frac{1}{2}\,\mathrm{d}x = \frac{1}{2}\times\left(2-\frac{3}{2}\right) = \frac{1}{2}\times\frac{1}{2} = \frac{1}{4}.
    \]
    \item \textbf{第二类错误}\(\beta = P\{\text{未拒绝}\,H_0 \mid H_1\text{ 为真}\} = P\!\left\{U\leq\tfrac{3}{2}\,\middle|\,H_1\right\}\)。
    在 \(H_1\) 下，\(U\) 的密度 \(f(x)=x/2\)（\(0\leq x\leq2\)）：
    \[
    \beta = \int_0^{3/2}\frac{x}{2}\,\mathrm{d}x = \frac{1}{2}\times\frac{x^2}{2}\Bigg|_0^{3/2} = \frac{1}{4}\times\frac{9}{4} = \frac{9}{16}.
    \]
\end{itemize}
故 \(\alpha=1/4\)，\(\beta=9/16\)，选 D。
\end{solution}


\begin{exercise}\label{exer:9-error-concept1}
在假设检验中，\(H_0\) 不成立时，检验统计量值未落入拒绝域，从而没有拒绝 \(H_0\)，这种错误称为第\textbf{二}类错误。
\end{exercise}
\begin{solution}
第二类错误（取伪）：\(H_0\) 不成立（\(H_1\) 为真）时，检验统计量值未落入拒绝域，从而错误地未拒绝 \(H_0\)。
\begin{note}
要区分两类错误的关键：
\begin{itemize}
    \item 第一类错误：\(H_0\) 本来为真，却被错误拒绝；
    \item 第二类错误：\(H_0\) 本来不真，却没有被拒绝。
\end{itemize}
\end{note}
\end{solution}


\begin{exercise}\label{exer:9-error-concept2}
显著性水平 \(\alpha\) 的含义是（\quad）。\par
A. \(H_0\) 成立，经检验被拒绝的概率\qquad B. \(H_0\) 成立，经检验不能拒绝的概率\par
C. \(H_0\) 不成立，经检验被拒绝的概率\qquad D. \(H_0\) 不成立，经检验不能拒绝的概率
\end{exercise}
\begin{solution}
应选 A。显著性水平 \(\alpha\) 的定义是：在 \(H_0\) 成立（为真）的条件下，检验结果拒绝 \(H_0\) 的概率，即
\(
\alpha = P\{\text{拒绝}\,H_0 \mid H_0\text{ 为真}\}.
\)
逐项分析：
\begin{itemize}
    \item A: \(H_0\) 成立，经检验被拒绝 \(\longrightarrow\) 恰好是 \(\alpha\) 的定义 \(\checkmark\)。
    \item B: \(H_0\) 成立，不能拒绝 \(\longrightarrow\) 这是 \(1-\alpha\)，不是 \(\alpha\) \(\times\)。
    \item C: \(H_0\) 不成立，被拒绝 \(\longrightarrow\) 这是\textbf{检验功效} \(1-\beta\) \(\times\)。
    \item D: \(H_0\) 不成立，不能拒绝 \(\longrightarrow\) 这是第二类错误概率 \(\beta\) \(\times\)。
\end{itemize}
\end{solution}

\chapter{历年真题}

\section{2018-2019-2学期《概率论与数理统计》A卷}

\paragraph*{一、选择题（共6题，每题3分，共18分）}
\begin{exercise}
\par
设 \(X_1,X_2,\cdots,X_n\ (n\ge2)\) 为来自总体 \(N(\mu,1)\) 的简单随机样本，记 \(\bar{X}=\frac{1}{n}\sum_{i=1}^n X_i\)。则下列结论中不正确的是（ ）。

（A）\(\sum_{i=1}^n (X_i-\mu)^2\) 服从 \(\chi^2\) 分布

（B）\(2(X_n-X_1)^2\) 服从 \(\chi^2\) 分布

（C）\(\sum_{i=1}^n (X_i-\bar{X})^2\) 服从 \(\chi^2\) 分布

（D）\(n(\bar{X}-\mu)^2\) 服从 \(\chi^2\) 分布
\end{exercise}

\begin{solution}
\par
由 \(X_i\sim N(\mu,1)\) 知 \(\sum_{i=1}^n (X_i-\mu)^2\sim \chi^2_n\)，
\(\sum_{i=1}^n (X_i-\bar{X})^2\sim \chi^2_{n-1}\)，
且 \(n(\bar{X}-\mu)^2\sim \chi^2_1\)。
又 \(X_n-X_1\sim N(0,2)\)，故 \(\dfrac{(X_n-X_1)^2}{2}\sim \chi^2_1\)，而 \(2(X_n-X_1)^2\) 不是 \(\chi^2\) 分布。
故选 B。
\end{solution}
\begin{exercise}
\par
设随机变量 \(X\sim P(3)\)，\(Y\sim B\!\left(8,\frac13\right)\)，且 \(X,Y\) 相互独立，则 \(\mathrm{Var}(X-3Y-4)=（ ）\)。

A. \(-13\)\qquad B. \(15\)\qquad C. \(19\)\qquad D. \(23\)
\end{exercise}

\begin{solution}
\par
\(\mathrm{Var}(X-3Y-4)=\mathrm{Var}(X)+9\mathrm{Var}(Y)\)。
对泊松分布 \(P(3)\)，\(\mathrm{Var}(X)=3\)；
对二项分布 \(B(8,\frac13)\)，
\(\mathrm{Var}(Y)=8\cdot\frac13\cdot\frac23=\frac{16}{9}\)。
因此
\[
  \mathrm{Var}(X-3Y-4)=3+9\cdot\frac{16}{9}=19.
\]
故选 C。
\end{solution}
\begin{exercise}
\par
设一个盒子中有 5 件产品，其中有 3 件正品，2 件次品。从盒子中任取两件，则取出的两件产品中至少有一件次品的概率为（ ）。

A. \(\frac{3}{10}\)\qquad B. \(\frac{5}{10}\)\qquad C. \(\frac{7}{10}\)\qquad D. \(\frac{1}{5}\)
\end{exercise}

\begin{solution}
\par
设事件 \(A\) 表示“至少有一件次品”。
其对立事件为“取出的两件全是正品”。
因此
\[
  P(A)=1-\frac{\binom32}{\binom52}=1-\frac{3}{10}=\frac{7}{10}.
\]
故选 C。
\end{solution}
\begin{exercise}
\par
随机变量 \(X\sim N(0,1)\)。对于给定的 \(\alpha\ (0<\alpha<1)\)，数 \(u_\alpha\) 满足 \(P(u>u_\alpha)=\alpha\)。若 \(P(|X|<c)=\alpha\)，则 \(c=\)（ ）。

A. \(u_{\alpha/2}\)\qquad B. \(u_{(1-\alpha)/2}\)\qquad C. \(u_{1-\alpha}\)\qquad D. \(u_{1-\alpha/2}\)
\end{exercise}

\begin{solution}
\par
由 \(P(|X|<c)=\alpha\) 知两侧尾概率各为 \(\dfrac{1-\alpha}{2}\)。
又 \(u_\beta\) 表示上侧 \(\beta\) 分位点，即 \(P(X>u_\beta)=\beta\)。
故
\[
  c=u_{(1-\alpha)/2}.
\]
选 B。
\end{solution}
\begin{exercise}
\par
设随机变量 \(X\) 与 \(Y\) 均服从正态分布，\(X\sim N(\mu,4^2)\)，\(Y\sim N(\mu,5^2)\)，而 \(p_1=P\{X\le \mu-4\}\)，\(p_2=P\{Y\ge \mu+5\}\)，则（ ）。

A. 对任何实数 \(\mu\)，都有 \(p_1=p_2\)．

B. 对任何实数 \(\mu\)，都有 \(p_1<p_2\)．

C. 只对 \(\mu\) 的个别值，才有 \(p_1=p_2\)．

D. 对任何实数 \(\mu\)，都有 \(p_1>p_2\)．
\end{exercise}

\begin{solution}
\par
标准化得
\[
  p_1=P\!\left\{Z\le \frac{-4}{4}\right\}=P(Z\le -1)=\Phi(-1),
\]
\[
  p_2=P\!\left\{Z\ge \frac{5}{5}\right\}=P(Z\ge 1)=\Phi(-1).
\]
故对任意 \(\mu\)，都有 \(p_1=p_2\)。选 A。
\end{solution}
\begin{exercise}
\par
设 \(X\sim P(\lambda)\)，且 \(E[(X-1)(X-2)]=1\)，则 \(\lambda=\)（ ）。

A. \(1\)\qquad B. \(2\)\qquad C. \(3\)\qquad D. \(0\)
\end{exercise}

\begin{solution}
\par
由 \(E[X]=E[X^2]-E[X]\) 的思路直接算更快：
\[
  E[(X-1)(X-2)]
  =E(X^2-3X+2)
  =E(X^2)-3E(X)+2.
\]
对泊松分布 \(P(\lambda)\)，有 \(E(X)=\lambda\)，\(\mathrm{Var}(X)=\lambda\)，故
\[
  E(X^2)=\mathrm{Var}(X)+[E(X)]^2=\lambda+\lambda^2.
\]
于是
\[
  E[(X-1)(X-2)]=\lambda^2-2\lambda+2.
\]
令其等于 \(1\)，得 \((\lambda-1)^2=0\)，故 \(\lambda=1\)。选 A。
\end{solution}

\paragraph*{二、填空题（共6题，每题3分，共18分）}
\begin{exercise}
\par
设 \(X_1,X_2\) 是来自总体 \(X\) 的样本，\(\hat{\mu}_1=\frac13X_1+\frac23X_2\)，\(\hat{\mu}_2=\frac12X_1+\frac12X_2\) 为总体均值 \(\mu\) 的无偏估计，则 \(\hat{\mu}_1,\hat{\mu}_2\) 中较有效的是 \underline{\hspace{4cm}}。
\end{exercise}

\begin{solution}
\par
无偏估计中，方差较小者较有效。设总体方差为 \(\sigma^2\)，则
\[
  \mathrm{Var}(\hat{\mu}_1)=\left(\frac19+\frac49\right)\sigma^2=\frac59\sigma^2,
  \qquad
  \mathrm{Var}(\hat{\mu}_2)=\left(\frac14+\frac14\right)\sigma^2=\frac12\sigma^2.
\]
因为 \(\frac12<\frac59\)，故较有效的是 \(\hat{\mu}_2\)。
\end{solution}
\begin{exercise}
\par
设随机变量 \(X\) 与 \(Y\) 独立且都服从 \([0,3]\) 上的均匀分布，则 \(P[\min(X,Y)\ge2]=\underline{\hspace{4cm}}\)。
\end{exercise}

\begin{solution}
\par
\[
  P[\min(X,Y)\ge2]=P(X\ge2,\ Y\ge2)=P(X\ge2)P(Y\ge2).
\]
对 \([0,3]\) 上的均匀分布，\(P(X\ge2)=\frac13\)，故
\[
  P[\min(X,Y)\ge2]=\frac13\cdot\frac13=\frac19.
\]
\end{solution}
\begin{exercise}
\par
设随机变量 \(X\) 的分布函数为 \(F(x)=0.5\Phi(x)+0.5\Phi\!\left(\frac{x-4}{2}\right)\)，其中 \(\Phi(x)\) 为标准正态分布函数，则 \(EX=\underline{\hspace{4cm}}\)。
\end{exercise}

\begin{solution}
\par
该分布函数对应两个正态分布的混合：
\(X\sim 0.5\,N(0,1)+0.5\,N(4,2^2)\)。
故
\[
  EX=0.5\cdot0+0.5\cdot4=2.
\]
\end{solution}
\begin{exercise}
\par
设随机变量 \(X\) 和 \(Y\) 的数学期望分别为 \(-2\) 和 \(2\)，方差分别为 \(1\) 和 \(4\)，而相关系数为 \(-0.5\)，则根据契比雪夫不等式 \(P\{|X+Y|\ge6\}\le\underline{\hspace{4cm}}\)。
\end{exercise}

\begin{solution}
\par
先算
\[
  E(X+Y)=-2+2=0.
\]
相关系数为 \(-0.5\)，故协方差
\[
  \mathrm{Cov}(X,Y)=\rho_{XY}\sigma_X\sigma_Y=-0.5\cdot1\cdot2=-1.
\]
于是
\[
  D(X+Y)=1+4+2(-1)=3.
\]
由契比雪夫不等式，
\[
  P\{|X+Y|\ge6\}\le \frac{D(X+Y)}{6^2}=\frac{3}{36}=\frac{1}{12}.
\]
\end{solution}
\begin{exercise}
\par
设随机变量 \(\xi\sim B(n,p)\)，\(E(\xi)=3\)，\(Var(\xi)=1.2\)，则 \(n=\underline{\hspace{4cm}}\)。
\end{exercise}

\begin{solution}
\par
由 \(E(\xi)=np=3\)，\(Var(\xi)=np(1-p)=1.2\)。
两式相除得
\[
  1-p=\frac{1.2}{3}=0.4,\qquad p=0.6.
\]
再由 \(np=3\) 得
\[
  n=\frac{3}{0.6}=5.
\]
\end{solution}
\begin{exercise}
\par
设总体 \(X\) 服从正态分布 \(N(0,4)\)，而 \(X_1,X_2,\cdots,X_{15}\) 是来自总体 \(X\) 的简单随机样本，则
\[
Y=\frac{X_1^2+X_2^2+\cdots+X_{10}^2}{2(X_{11}^2+X_{12}^2+\cdots+X_{15}^2)}
\]
服从的分布是 \underline{\hspace{4cm}}（包括分布参数）。
\end{exercise}

\begin{solution}
\par
因 \(X_i\sim N(0,4)\)，故 \(\frac{X_i}{2}\sim N(0,1)\)，于是
\[
  \frac{X_1^2+\cdots+X_{10}^2}{4}\sim \chi^2(10),\qquad
  \frac{X_{11}^2+\cdots+X_{15}^2}{4}\sim \chi^2(5).
\]
记前者为 \(U\)，后者为 \(V\)，则
\[
  Y=\frac{4U}{2\cdot4V}=\frac{U}{2V}
  =\frac{U/10}{V/5}.
\]
故 \(Y\sim F(10,5)\)。
\end{solution}

\paragraph*{三、计算题（10分）}
\begin{exercise}
\par
某保险公司把被保险人分为三类：谨慎的、一般的、冒失的，统计资料表明，上述三种人在一年内发生事故的概率依次为 \(0.05\)、\(0.15\) 和 \(0.30\)。如果谨慎的占总的被保人数的 \(20\%\)，一般的占 \(50\%\)，冒失的占 \(30\%\)。

(1) 求某被保人在一年内发生事故的概率；

(2) 若此人在一年内发生事故，则他是谨慎的客户的概率是多少。
\end{exercise}

\begin{solution}
\par
设 \(A_1,A_2,A_3\) 分别表示谨慎的、一般的、冒失的被保险人，设 \(B\) 表示“一年内发生事故”。
由全概率公式，
\[
  P(B)=\sum_{i=1}^3 P(B|A_i)P(A_i)
  =0.05\cdot0.2+0.15\cdot0.5+0.30\cdot0.3=0.175.
\]
再由贝叶斯公式，
\[
  P(A_1|B)=\frac{P(B|A_1)P(A_1)}{P(B)}
  =\frac{0.05\cdot0.2}{0.175}
  =\frac{2}{35}\approx0.0571.
\]
故 (1) 为 \(0.175\)，(2) 为 \(\dfrac{2}{35}\approx0.0571\)。
\end{solution}

\paragraph*{四、计算题（10分）}
\begin{exercise}
\par
设某次概率统计考试考生的成绩 \(X\sim N(\mu,\sigma^2)\)，从中随机地抽取 36 位考生的成绩，算得平均成绩为 \(\bar{x}=66.5\) 分，修正标准差为 \(S^*=15\) 分。

(1) 在置信度为 \(0.95\) 时，求考生成绩数学期望 \(\mu\) 的置信区间；

(2) 在显著性水平 \(\alpha=0.05\) 下，检验是否可以认为这次考试的平均成绩为 70 分。

已知：\(t_{0.975}(35)=2.0301,\ t_{0.975}(36)=2.0281,\ t_{0.95}(35)=1.6896,\ t_{0.95}(36)=1.6883\)
\end{exercise}

\begin{solution}
\par
(1) 因 \(\sigma\) 未知，采用 \(t\) 区间：
\[
  \frac{\bar{X}-\mu}{S^*/\sqrt{n}}\sim t(35).
\]
所以
\[
  \mu\in \bar{x}\pm t_{0.975}(35)\frac{S^*}{\sqrt{n}}
  =66.5\pm2.0301\cdot\frac{15}{6}.
\]
计算得置信区间为
\[
  (61.42,71.58).
\]

(2) 提出假设
\[
  H_0:\mu=70,\qquad H_1:\mu\ne70.
\]
检验统计量为
\[
  t=\frac{\bar{x}-70}{S^*/\sqrt{n}}
  =\frac{66.5-70}{15/6}=-1.4.
\]
因为 \(|t|=1.4<t_{0.975}(35)=2.0301\)，故不拒绝原假设。
可以认为这次考试的平均成绩为 70 分。
\end{solution}

\paragraph*{五、计算题（10分）}
\begin{exercise}
\par
某单位有 400 部内线电话，每时刻每部电话打外线的概率为 \(10\%\)，设各电话使用外线与否是相互独立的，估计在任一时刻有 30～50 部电话同时使用外线的概率。

已知：\(\Phi(1.67)=0.9525,\ \Phi(1.60)=0.9452,\ \Phi(1.52)=0.9357,\ \Phi(1.36)=0.9131\)
\end{exercise}

\begin{solution}
\par
设 \(X\) 为同时使用外线的电话数，则
\[
  X\sim B(400,0.1),\qquad EX=np=40,\qquad D X=np(1-p)=36.
\]
用正态近似，并作连续性修正：
\[
  P(30\le X\le50)\approx P\!\left\{-\frac{10}{6}\le \frac{X-40}{6}\le\frac{10}{6}\right\}.
\]
标准化得
\[
  P(30\le X\le50)\approx \Phi(1.67)-\Phi(-1.67).
\]
利用对称性，
\[
  \Phi(1.67)-\Phi(-1.67)=2\Phi(1.67)-1=0.905.
\]
故概率约为 \(0.905\)。
\end{solution}

\paragraph*{六、计算题（10分）}
\begin{exercise}
\par
设 \(X_1,X_2,\dots,X_n\) 是取自双参数指数分布总体的一组样本，密度函数为
\[
f(x;\theta,\mu)=
\begin{cases}
\dfrac{1}{\theta}e^{-\frac{x-\mu}{\theta}}, & x>\mu,\\
0, & \text{其它},
\end{cases}
\]
其中 \(\theta,\mu\) 是未知参数，\(x_1,x_2,\dots,x_n\) 是一组样本值，求 \(\theta,\mu\) 的最大似然估计量。
\end{exercise}

\begin{solution}
\par
似然函数为
\[
  L(\mu,\theta)=\frac{1}{\theta^n}\exp\!\left\{-\frac{1}{\theta}\left(\sum_{i=1}^n x_i-n\mu\right)\right\},
  \qquad \mu\le \min\{x_1,\dots,x_n\}.
\]
对数似然
\[
  \ell(\mu,\theta)=-n\ln\theta-\frac{1}{\theta}\left(\sum_{i=1}^n x_i-n\mu\right).
\]
对 \(\mu\) 求偏导，
\[
  \frac{\partial \ell}{\partial\mu}=\frac{n}{\theta}>0,
\]
故 \(\ell\) 随 \(\mu\) 单调增，取
\[
  \hat{\mu}_L=X_{(1)}=\min\{X_1,\dots,X_n\}.
\]
再对 \(\theta\) 求偏导并令其为零，
\[
  \hat{\theta}_L=\bar{X}-X_{(1)}.
\]
故最大似然估计量为
\[
  \hat{\mu}_L=X_{(1)},\qquad \hat{\theta}_L=\bar{X}-X_{(1)}.
\]
\end{solution}

\paragraph*{七、计算题（12分）}
\begin{exercise}
\par
设随机变量 \(X\) 和 \(Y\) 的联合分布在以点 \((0,1),(1,0),(1,1)\) 为顶点的三角形区域上服从均匀分布，求随机变量 \(U=X+Y\) 的方差。
\end{exercise}

\begin{solution}
\par
联合分布的支持区域为
\[
  G=\{(x,y):0\le x\le1,\ 0\le y\le1,\ x+y\ge1\},
\]
其面积为 \(\frac12\)，故联合密度为
\[
  f(x,y)=
  \begin{cases}
    2, & (x,y)\in G,\\
    0, & (x,y)\notin G.
  \end{cases}
\]
先求边缘密度。对 \(0<x<1\)，有
\[
  f_X(x)=\int_{1-x}^1 2\,\mathrm{d}y=2x.
\]
因此
\[
  EX=\int_0^1 x\cdot2x\,\mathrm{d}x=\frac23,\qquad
  EX^2=\int_0^1 x^2\cdot2x\,\mathrm{d}x=\frac12,
\]
从而
\[
  DX=\frac12-\left(\frac23\right)^2=\frac1{18}.
\]
同理 \(EY=\frac23,\ DY=\frac1{18}\)。

再算
\[
  EXY=\iint_G 2xy\,\mathrm{d}x\mathrm{d}y
  =2\int_0^1 x\left(\int_{1-x}^1 y\,\mathrm{d}y\right)\mathrm{d}x
  =\frac{5}{12}.
\]
于是
\[
  \mathrm{Cov}(X,Y)=EXY-EX\cdot EY
  =\frac{5}{12}-\frac49=-\frac1{36}.
\]
所以
\[
  D(U)=D(X+Y)=DX+DY+2\mathrm{Cov}(X,Y)
  =\frac1{18}+\frac1{18}-\frac{2}{36}
  =\frac1{18}.
\]
\end{solution}

\paragraph*{八、计算题（12分）}
\begin{exercise}
\par
设随机变量 \(X,Y\) 相互独立，且 \(X\) 的概率分布为 \(P\{X=0\}=P\{X=1\}=\frac12\)，\(Y\) 的概率密度为
\[
f(y)=
\begin{cases}
2y, & 0<y<1,\\
0, & \text{其他}.
\end{cases}
\]
（1）求 \(P\{Y\le EY\}\)；

（2）求 \(Z=X+Y\) 的概率密度。
\end{exercise}

\begin{solution}
\par
(1) 先求 \(EY\)：
\[
  EY=\int_0^1 y\cdot2y\,\mathrm{d}y=\frac23.
\]
因此
\[
  P\{Y\le EY\}=P\!\left\{Y\le\frac23\right\}
  =\int_0^{2/3}2y\,\mathrm{d}y=\frac49.
\]

(2) 由 \(X\) 取 \(0,1\) 的概率各为 \(\frac12\)，得
\[
  F_Z(z)=P(Z\le z)
  =\frac12P(Y\le z)+\frac12P(Y\le z-1)
  =\frac12F_Y(z)+\frac12F_Y(z-1).
\]
对 \(F_Z\) 求导，得
\[
  f_Z(z)=\frac12f_Y(z)+\frac12f_Y(z-1)
  =
  \begin{cases}
    z, & 0<z<1,\\
    z-1, & 1\le z<2,\\
    0, & \text{其他}.
  \end{cases}
\]
故所求密度如上。
\end{solution}



\section{2019-2020-2学期《概率论与数理统计》A卷}
\paragraph*{一、选择题}
\begin{exercise}
\par
若 \(P(AB)=0\)，则下述命题中正确的是（\( \boldsymbol{B} \)）。
A. \(A\) 与 \(B\) 互不相容；
B. \(P(A\cup B)=P(A)+P(B)\)；
C. \(A\) 与 \(B\) 相互独立；
D. \(P(A)=0\) 或 \(P(B)=0\)。
\end{exercise}

\begin{solution}
\par
由
\[
P(A\cup B)=P(A)+P(B)-P(AB)=P(A)+P(B),
\]
故选 B。
\end{solution}

\begin{exercise}
\par
设 \(X\sim N(0,2),Y\sim N(1,2)\)，且 \(X\) 与 \(Y\) 相互独立，则（\( \boldsymbol{D} \)）。
A. \(P(X+Y\le0)=\frac12\)；
B. \(P(X-Y\le1)=\frac12\)；
C. \(P(X-Y\le0)=\frac12\)；
D. \(P(X+Y\le1)=\frac12\)。
\end{exercise}

\begin{solution}
\par
\(X+Y\sim N(1,4)\)，故
\[
P(X+Y\le1)=\frac12.
\]
故选 D。
\end{solution}

\begin{exercise}
\par
设总体 \(X\sim N(0,1)\)，\(X_1,X_2\) 是总体 \(X\) 的一个样本，令 \(Y=\frac{X_1}{|X_2|}\)，则 \(Y\) 服从（\( \boldsymbol{C} \)）。
A. 正态分布； B. \(\chi^2\) 分布； C. \(t\) 分布； D. \(F\) 分布。
\end{exercise}

\begin{solution}
\par
因 \(X_1\sim N(0,1)\)、\(X_2^2\sim\chi^2(1)\)，且
\[
Y=\frac{X_1}{\sqrt{X_2^2/1}},
\]
故 \(Y\sim t(1)\)，即 \(t\) 分布。故选 C。
\end{solution}

\begin{exercise}
\par
设 \(X_1,X_2,X_3\) 独立同分布，分布函数为 \(F(x)\)，令 \(Y=\min\{X_1,X_2,X_3\}\)，则 \(P(Y\le1)=\)（\( \boldsymbol{B} \)）。
A. \(1-(F(1))^3\)； B. \(1-(1-F(1))^3\)； C. \((F(1))^3\)； D. \((1-F(1))^3\)。
\end{exercise}

\begin{solution}
\par
\[
P(Y\le1)=1-P(X_1>1,X_2>1,X_3>1)
=1-(1-F(1))^3.
\]
故选 B。
\end{solution}

\begin{exercise}
\par
设随机变量 \(X\) 的密度函数为
\[
f(x)=\frac{1}{2\sqrt{\pi}}e^{-\frac{x^2}{4}},
\]
则 \(Y=2X\) 的密度函数为（\( \boldsymbol{B} \)）。
A. \(\frac{1}{2\sqrt{\pi}}e^{-\frac{y^2}{16}}\)； 
B. \(\frac{1}{4\sqrt{\pi}}e^{-\frac{y^2}{16}}\)； 
C. \(\frac{1}{2\sqrt{\pi}}e^{-\frac{y^2}{8}}\)； 
D. \(\frac{1}{2\sqrt{2\pi}}e^{-\frac{y^2}{8}}\)。
\end{exercise}

\begin{solution}
\par
由变量变换 \(x=y/2\)，得
\[
f_Y(y)=f_X\!\left(\frac y2\right)\cdot\frac12
=\frac{1}{4\sqrt{\pi}}e^{-\frac{y^2}{16}}.
\]
故选 B。
\end{solution}

\begin{exercise}
\par
设随机变量 \(X\) 服从 \(t\) 分布 \(t(n)\)，对给定的 \(\alpha(0<\alpha<1)\)，数 \(t_\alpha(n)\) 满足 \(P(X>t_\alpha(n))=\alpha\)，若 \(P(|X|\le x)=\alpha\)，则 \(x\) 等于（\( \boldsymbol{A} \)）。
A. \(t_{\frac{1-\alpha}{2}}(n)\)； B. \(t_{1-\frac{\alpha}{2}}(n)\)； C. \(t_{\frac{\alpha}{2}}(n)\)； D. \(t_{1-\alpha}(n)\)。
\end{exercise}

\begin{solution}
\par
由对称性，\(P(|X|\le x)=\alpha\) 等价于两侧尾概率各为 \(\frac{1-\alpha}{2}\)，故
\[
x=t_{\frac{1-\alpha}{2}}(n).
\]
故选 A。
\end{solution}

\paragraph*{二、填空题}
\begin{exercise}
\par
设随机变量 \(X\) 和 \(Y\) 的期望分别为 \(-1\) 和 \(1\)，方差分别为 \(1\) 和 \(2\)，且 \(X\) 和 \(Y\) 相互独立，用切比雪夫不等式估计 \(P(|X+Y|\ge3)\le \underline{\hspace{2cm}}\)。
\end{exercise}

\begin{solution}
\par
\[
E(X+Y)=0,\qquad D(X+Y)=1+2=3.
\]
由切比雪夫不等式，
\[
P(|X+Y|\ge3)\le \frac{D(X+Y)}{3^2}=\frac13.
\]
\end{solution}

\begin{exercise}
\par
设随机变量 \(X\) 的分布函数为
\[
F(x)=
\begin{cases}
0, & x<1,\\
0.4, & 1\le x<3,\\
1, & x\ge3,
\end{cases}
\]
则 \(P(X<2\mid X\neq1)=\underline{\hspace{2cm}}\)。
\end{exercise}

\begin{solution}
\par
\(X\) 只可能取 \(1\) 或 \(3\)，故在条件 \(X\neq1\) 下必有 \(X=3\)，于是
\[
P(X<2\mid X\neq1)=0.
\]
\end{solution}

\begin{exercise}
\par
若二维随机向量 \((X,Y)\sim N(1,2,4,9,0)\)，则 \(Var[X-2Y+1]=\underline{\hspace{2cm}}\)。
\end{exercise}

\begin{solution}
\par
由题意可读出 \(D(X)=4,\ D(Y)=9,\ \rho_{XY}=0\)，故
\[
D(X-2Y+1)=D(X)+4D(Y)=4+36=40.
\]
\end{solution}

\begin{exercise}
\par
设 \(X_1,X_2,\dots,X_n\) 是来自总体 \(X\sim N(0,\sigma^2)\) 的样本，若 \(\hat{\sigma}^2=a\sum_{i=1}^n X_i^2\) 是 \(\sigma^2\) 的无偏估计，则 \(a=\underline{\hspace{2cm}}\)。
\end{exercise}

\begin{solution}
\par
\[
E\!\left(\sum_{i=1}^n X_i^2\right)=n\sigma^2,
\]
故无偏性要求 \(an\sigma^2=\sigma^2\)，于是
\[
a=\frac1n.
\]
\end{solution}

\begin{exercise}
\par
设 \(X_1,\dots,X_4\) 是来自总体 \(X\sim N(0,2)\) 的样本，令
\[
Y=\frac12(X_1-X_2)^2+\frac12(X_3+X_4)^2,
\]
则 \(E[Y]=\underline{\hspace{2cm}}\)。
\end{exercise}

\begin{solution}
\par
\[
E[(X_1-X_2)^2]=D(X_1-X_2)=2+2=4,
\qquad
E[(X_3+X_4)^2]=D(X_3+X_4)=2+2=4.
\]
故
\[
E[Y]=\frac12\cdot4+\frac12\cdot4=4.
\]
\end{solution}

\begin{exercise}
\par
已知 \(X,Y\) 满足 \(E[X]=1,Var[X]=9,E[Y]=0,Var[Y]=16\)，相关系数 \(r(X,Y)=-\frac12\)，设随机变量 \(Z=\frac{X}{3}+\frac{Y}{2}\)，则 \(X\) 与 \(Z\) 的协方差 \(\mathrm{Cov}(X,Z)=\underline{\hspace{2cm}}\)。
\end{exercise}

\begin{solution}
\par
\[
\mathrm{Cov}(X,Y)=r(X,Y)\sigma_X\sigma_Y=-\frac12\cdot3\cdot4=-6.
\]
因此
\[
\mathrm{Cov}(X,Z)=\frac13\mathrm{Var}(X)+\frac12\mathrm{Cov}(X,Y)
=\frac13\cdot9+\frac12(-6)=0.
\]
\end{solution}

\paragraph*{三、计算题}
\begin{exercise}
\par
在房间里有 10 个人，分别佩戴从 1 号到 10 号的纪念章，任选 3 人记录其纪念章的号码。试求：
(1) 最大号码为 5 的概率；
(2) 如果已知记录的最大号码为 5，求记录的三个号码中有 3 的概率？
\end{exercise}

\begin{solution}
\par
设 \(A=\{\)记录的最大号码为 5\(\}\)，则
\[
P(A)=\frac{\binom42}{\binom{10}3}=\frac1{20}.
\]
设 \(B=\{\)记录的号码中有 3\(\}\)，则在事件 \(A\) 下，满足条件的三元组共有
\(\{3,5,\cdot\}\) 中从 \(\{1,2,4\}\) 选 1 个，共 3 种，而 \(A\) 共 6 种，故
\[
P(B\mid A)=\frac12.
\]
\end{solution}

\begin{exercise}
\par
设某面粉厂采用自动流水线灌装面粉，装袋重量 \(X\sim N(\mu,\sigma^2)\)。从中随机地抽取 36 袋，经计算得平均重量为 \(\bar{x}=24.92\mathrm{kg}\)，修正标准差 \(s_n^*=1.974\mathrm{kg}\)。
(1) 在置信度为 0.95 时，求出面粉重量平均值 \(\mu\) 的置信区间。
(2) 在显著性水平 \(\alpha=0.05\) 下，检验是否可以认为该流水线面粉的平均重量为每袋 25 kg。
\end{exercise}

\begin{solution}
\par
(1) 因 \(\sigma\) 未知，取 \(t\) 区间：
\[
\mu\in \bar{x}\pm t_{0.975}(35)\frac{s_n^*}{\sqrt{36}}
=24.92\pm2.0301\cdot\frac{1.974}{6}.
\]
故置信区间为
\[
[24.252,25.588].
\]

(2) 提出假设：
\[
H_0:\mu=25,\qquad H_1:\mu\ne25.
\]
统计量
\[
T=\frac{\bar{X}-25}{s_n^*/\sqrt{36}}\sim t(35).
\]
观测值为
\[
t=\frac{24.92-25}{1.974/6}\approx-0.2432.
\]
由于 \(|t|<2.0301\)，故不拒绝 \(H_0\)，可以认为平均重量为每袋 25 kg。
\end{solution}

\begin{exercise}
\par
抽样检查产品质量时，如果发现有多于 10 个的次品，则拒绝接受这批产品。设某批产品的次品率为 10\%，试用中心极限定理来判断，至少应抽取多少个产品来检查，才能保证拒绝接受该产品的概率达到 0.95？
\end{exercise}

\begin{solution}
\par
设抽取 \(n\) 个产品，次品数为 \(X\sim B(n,0.1)\)。由中心极限定理，
\[
X\approx N(0.1n,0.09n).
\]
要求
\[
P(X>10)\approx 0.95,
\]
于是
\[
\frac{0.1n-10}{0.3\sqrt n}\approx1.645.
\]
解得 \(n\approx163.01\)，故至少应抽取
\[
n=164
\]
个产品。
\end{solution}

\begin{exercise}
\par
设二维随机向量 \((X,Y)\) 服从矩形区域 \(D=\{(x,y)\mid 0\le x\le2,0\le y\le1\}\) 上的均匀分布，且
\[
U=\begin{cases}
-1, & X\ge Y,\\
1, & X<Y,
\end{cases}
\qquad
V=\begin{cases}
0, & X\le2Y,\\
1, & X>2Y.
\end{cases}
\]
求：(1) \(P(U=-1)\)；(2) \(U,V\) 的联合分布；(3) \(W=UV\) 的分布列。
\end{exercise}

\begin{solution}
\par
区域面积为 2，故联合密度为 \(f(x,y)=\frac12\)（在 \(D\) 上）。

(1)
\[
P(U=-1)=P(X\ge Y)=\frac34.
\]

(2) 分四块计算：
\[
P(U=-1,V=0)=\frac14,\quad
P(U=-1,V=1)=0,
\]
\[
P(U=1,V=0)=\frac14,\quad
P(U=1,V=1)=\frac12.
\]
故联合分布为
\[
\begin{array}{c|cc}
U\backslash V & 0 & 1\\ \hline
-1 & \frac14 & 0\\
1 & \frac14 & \frac12
\end{array}
\]

(3) \(W=UV\) 只取 \(0,1\)，且
\[
P(W=0)=\frac12,\qquad P(W=1)=\frac12.
\]
\end{solution}

\begin{exercise}
\par
设随机变量 \((X,Y)\) 的概率密度函数为
\[
f(x,y)=
\begin{cases}
Ae^{-(x+2y)}, & 0<x<1,\ y>0,\\
0, & \text{其它},
\end{cases}
\]
（1）求 \(A\) 的值；（2）判断 \(X,Y\) 是否独立；（3）期望 \(E[2XY]\)。
\end{exercise}

\begin{solution}
\par
(1) 由归一性，
\[
1=\int_0^1\int_0^\infty Ae^{-(x+2y)}\,\mathrm{d}y\mathrm{d}x
=\frac A2(1-e^{-1}),
\]
故
\[
A=\frac{2}{1-e^{-1}}.
\]

(2) 边缘密度为
\[
f_X(x)=\frac{e^{-x}}{1-e^{-1}},\quad 0<x<1,
\qquad
f_Y(y)=2e^{-2y},\quad y>0.
\]
并且 \(f(x,y)=f_X(x)f_Y(y)\)，故 \(X,Y\) 相互独立。

(3)
\[
E[X]=\frac{1-2e^{-1}}{1-e^{-1}},\qquad E[Y]=\frac12,
\]
于是
\[
E[2XY]=2E[X]E[Y]=\frac{1-2e^{-1}}{1-e^{-1}}.
\]
\end{solution}

\paragraph*{四、证明/计算题}
\begin{exercise}
\par
设总体 \(X\) 具有分布律：
\[
\begin{array}{c|ccc}
X & 1 & 2 & 3\\ \hline
P & \theta^2 & 2\theta(1-\theta) & (1-\theta)^2
\end{array}
\]
其中 \(\theta(0<\theta<1)\) 为未知参数。
(1) 求 \(\theta\) 的矩估计量 \(\hat{\theta}_M\)，并讨论 \(\hat{\theta}_M\) 的无偏性；
(2) 当样本观测值为 \(x_1=1,x_2=2,x_3=1,x_4=1,x_5=3\) 时，求 \(\theta\) 的极大似然估计值。
\end{exercise}

\begin{solution}
\par
(1) 先求数学期望：
\[
E[X]=1\cdot\theta^2+2\cdot2\theta(1-\theta)+3(1-\theta)^2=3-2\theta.
\]
令 \(\bar X=3-2\theta\)，得矩估计量
\[
\hat\theta_M=\frac{3-\bar X}{2}.
\]
又因 \(E[\bar X]=E[X]=3-2\theta\)，故
\[
E(\hat\theta_M)=\theta,
\]
所以 \(\hat\theta_M\) 是无偏的。

(2) 样本中 \(1,2,3\) 的个数分别为 \(3,1,1\)，故似然函数为
\[
L(\theta)=\theta^{6}\cdot[2\theta(1-\theta)]\cdot(1-\theta)^2
=2\theta^7(1-\theta)^3.
\]
取对数并求极值，得
\[
\frac{7}{\theta}-\frac{3}{1-\theta}=0,
\]
解得
\[
\hat\theta_{MLE}=\frac{7}{10}.
\]
\end{solution}


\end{document}
